\documentclass[12pt,b5paper,twoside,openany]{book}

\usepackage[T1]{fontenc}
\usepackage[english]{babel}
\usepackage{anyfontsize}
\usepackage{newpxtext}
\usepackage{newpxmath}
\usepackage{mathtools}
\usepackage{bm}
\usepackage{array}
\usepackage{booktabs}
\usepackage{longtable}
\usepackage{enumitem}
\usepackage{manyfoot}
\usepackage{xcolor}

\makeatletter
\renewcommand\normalsize{%
  \@setfontsize\normalsize{13pt}{15.6pt}%
  \abovedisplayskip 12pt plus 3pt minus 7pt%
  \abovedisplayshortskip 0pt plus 3pt%
  \belowdisplayshortskip 7pt plus 3pt minus 4pt%
  \belowdisplayskip \abovedisplayskip
  \let\@listi\@listI
}
\renewcommand\small{\@setfontsize\small{11pt}{13.2pt}}
\renewcommand\footnotesize{\@setfontsize\footnotesize{10pt}{12pt}}
\renewcommand\scriptsize{\@setfontsize\scriptsize{9pt}{10.8pt}}
\renewcommand\tiny{\@setfontsize\tiny{7.5pt}{9pt}}
\renewcommand\large{\@setfontsize\large{15pt}{18pt}}
\renewcommand\Large{\@setfontsize\Large{18pt}{21.6pt}}
\renewcommand\LARGE{\@setfontsize\LARGE{21pt}{25.2pt}}
\renewcommand\huge{\@setfontsize\huge{24pt}{28.8pt}}
\renewcommand\Huge{\@setfontsize\Huge{28pt}{33.6pt}}
\makeatother

\allowdisplaybreaks[4]
\linespread{1.1}
\normalsize

\usepackage[
  paperwidth=176mm,
  paperheight=250mm,
  % The text block is inset by an additional 1 mm so that italic and
  % mathematical glyph overhangs still leave the requested 7 mm of ink margin.
  left=8mm,
  right=8mm,
  top=9.1mm,
  % The extra depth compensates for descenders, display limits, and footnotes,
  % keeping the visible text at least 9.1 mm above the lower paper edge.
  bottom=11.4mm,
  includehead,
  headheight=6.1mm,
  headsep=1mm,
  footskip=5pt
]{geometry}

\usepackage{fancyhdr}
\fancyhf{}
\fancyheadoffset[LE,RO]{0.95mm}
\fancyhead[LE]{\raisebox{1.36mm}[0pt][0pt]{\normalfont\normalsize\thepage}}
\fancyhead[RO]{\raisebox{1.36mm}[0pt][0pt]{\normalfont\normalsize\thepage}}
\renewcommand{\headrulewidth}{0pt}
\renewcommand{\footrulewidth}{0pt}
\pagestyle{fancy}
\fancypagestyle{plain}{%
  \fancyhf{}
  \fancyheadoffset[LE,RO]{0.95mm}
  \fancyhead[LE]{\raisebox{1.36mm}[0pt][0pt]{\normalfont\normalsize\thepage}}
  \fancyhead[RO]{\raisebox{1.36mm}[0pt][0pt]{\normalfont\normalsize\thepage}}
  \renewcommand{\headrulewidth}{0pt}
  \renewcommand{\footrulewidth}{0pt}
}

\definecolor{LinkBlue}{HTML}{1F4E8C}
\definecolor{URLTeal}{HTML}{006D77}
\usepackage[
  unicode=true,
  colorlinks=true,
  hypertexnames=false,
  linktoc=all,
  linkcolor=LinkBlue,
  citecolor=LinkBlue,
  urlcolor=URLTeal,
  filecolor=URLTeal,
  pdfborder={0 0 0},
  bookmarksnumbered=true,
  pdfpagelabels=true,
  pdftitle={C*-Algebras and Their Representations, Second Edition},
  pdfauthor={Jacques Dixmier}
]{hyperref}

\setcounter{tocdepth}{1}
\setcounter{secnumdepth}{-1}
\makeatletter
\renewcommand*\l@chapter[2]{%
  \ifnum \c@tocdepth >\m@ne
    \addpenalty{-\@highpenalty}%
    \vskip 1.0em \@plus\p@
    \setlength\@tempdima{2em}%
    \begingroup
      \parindent \z@ \rightskip \@pnumwidth
      \parfillskip -\@pnumwidth
      \leavevmode \bfseries
      \advance\leftskip\@tempdima
      \hskip -\leftskip
      #1\nobreak\hfil
      \nobreak\hb@xt@\@pnumwidth{\hss #2\kern-\p@\kern\p@}\par
      \penalty\@highpenalty
    \endgroup
  \fi}
\renewcommand*\l@section{\@dottedtocline{1}{0.8em}{2.8em}}
\makeatother
\setlength{\parindent}{1.5em}
\setlength{\parskip}{0pt}
\setlist{nosep,leftmargin=2.4em}
\raggedbottom
\emergencystretch=2em
\clubpenalty=10000
\widowpenalty=10000
\displaywidowpenalty=10000

\DeclareNewFootnote{T}[alph]
\makeatletter
\newcommand{\translatornote}[1]{%
  \footnoteT{\hypertarget{\Hy@footnote@currentHref}{}%
    \textit{Translator's note.}\enspace #1}%
}
\makeatother

\newcommand{\ResultHeading}[2]{%
  \par\addvspace{.55\baselineskip}%
  \hypertarget{result:#1}{}%
  \noindent\textbf{#2}\enspace\ignorespaces
}
\newcommand{\ResultLabel}[1]{%
  \par\addvspace{.55\baselineskip}%
  \noindent\textbf{#1}\enspace\ignorespaces
}
\newcommand{\Proof}{%
  \par\addvspace{.35\baselineskip}%
  \noindent\textit{Proof.}\enspace\ignorespaces
}
\newcommand{\ProofOf}[1]{%
  \par\addvspace{.35\baselineskip}%
  \noindent\textit{#1}\enspace\ignorespaces
}
\newcommand{\XRef}[2]{\hyperlink{result:#1}{#2}}
\newcommand{\SectionRef}[2]{\hyperlink{section:#1}{#2}}
\newcommand{\AvNTarget}{\hypertarget{reference:AvN}{[A v N]}}
\newcommand{\AvNRef}{\cite{refAvN}}
\newcommand{\PdfOrTeXTitle}[2]{%
  \if\relax\detokenize{#1}\relax
    #2%
  \else
    \texorpdfstring{#2}{#1}%
  \fi
}
\DeclareRobustCommand{\RaggedTitleBreak}{\hfill\break}

\newcommand{\BookSection}[3][]{%
  \clearpage
  \phantomsection
  \hypertarget{section:#2}{}%
  \addcontentsline{toc}{chapter}{\protect\numberline{\S~#2}\PdfOrTeXTitle{#1}{#3}}%
  \begin{center}
    \Large\bfseries \S~#2. --- #3
  \end{center}
  \vspace{.25\baselineskip}
}
\newcommand{\BookSectionHere}[3][]{%
  \phantomsection
  \hypertarget{section:#2}{}%
  \addcontentsline{toc}{chapter}{\protect\numberline{\S~#2}\PdfOrTeXTitle{#1}{#3}}%
  \begin{center}
    \Large\bfseries \S~#2. --- #3
  \end{center}
  \vspace{.25\baselineskip}
}
\newcommand{\BookSubsection}[3][]{%
  \section*{\texorpdfstring{%
    \raisebox{1.5\baselineskip}[0pt][0pt]{\hypertarget{section:#2}{}}%
  }{}#2. #3}%
  \addcontentsline{toc}{section}{\protect\numberline{#2}\PdfOrTeXTitle{#1}{#3}}%
}
\newcommand{\BookSubsectionToc}[4][]{%
  \section*{\texorpdfstring{%
    \raisebox{1.5\baselineskip}[0pt][0pt]{\hypertarget{section:#2}{}}%
  }{}#2. #3}%
  \addcontentsline{toc}{section}{\protect\numberline{#2}\PdfOrTeXTitle{#1}{#4}}%
}
\newcommand{\BookSubsubsection}[2]{%
  \subsection*{\texorpdfstring{%
    \raisebox{1.5\baselineskip}[0pt][0pt]{\hypertarget{section:#1}{}}%
  }{}#1. #2}%
}
\newcommand{\BookDivision}[3][]{%
  \clearpage
  \phantomsection
  \hypertarget{division:#2}{}%
  \addcontentsline{toc}{part}{\PdfOrTeXTitle{#1}{#3}}%
  \begin{center}
    \LARGE\bfseries #3
  \end{center}
  \vspace{.5\baselineskip}
}
\newcommand{\BookAppendix}[3][]{%
  \clearpage
  \phantomsection
  \hypertarget{section:#2}{}%
  \addcontentsline{toc}{chapter}{Appendix #2: \PdfOrTeXTitle{#1}{#3}}%
  \begin{center}
    \Large\bfseries Appendix #2\\[.35\baselineskip]#3
  \end{center}
  \vspace{.25\baselineskip}
}
\newcommand{\BookBackMatter}[2]{%
  \clearpage
  \phantomsection
  \hypertarget{backmatter:#1}{}%
  \addcontentsline{toc}{chapter}{#2}%
  \begin{center}
    \Large\bfseries #2
  \end{center}
  \vspace{.25\baselineskip}
}

\makeatletter
\newenvironment{DixmierBibliography}[1]
  {\list{\@biblabel{\@arabic\c@enumiv}}%
    {\settowidth\labelwidth{\@biblabel{#1}}%
     \leftmargin\labelwidth
     \advance\leftmargin\labelsep
     \@openbib@code
     \usecounter{enumiv}%
     \let\p@enumiv\@empty
     \renewcommand\theenumiv{\@arabic\c@enumiv}}%
   \sloppy
   \hbadness=2000
   \clubpenalty4000
   \@clubpenalty\clubpenalty
   \widowpenalty4000%
   \sfcode`\.\@m}
  {\def\@noitemerr{\@latex@warning{Empty `DixmierBibliography' environment}}%
   \endlist}
\makeatother

\pdfstringdefDisableCommands{%
  \def\S{Section}%
  \def\ast{*}%
  \def\mathrm#1{#1}%
  \def\mathbf#1{#1}%
}

\begin{document}
\pagenumbering{arabic}
\setcounter{page}{1}
% ===== ASSEMBLED TRANSLATION CHUNK 01: chunk_001_029.tex =====
% Source PDF physical page 001
\begin{titlepage}
\thispagestyle{empty}
\begin{center}
  {\Large\bfseries CAHIERS SCIENTIFIQUES\par}
  \vspace{0.4em}
  {\small Published under the direction of M. Gaston Julia\par}
  \vspace{2em}
  {\large\bfseries FASCICULE XXIX\par}
  \vfill
  {\Huge\bfseries LES $C^{*}$-ALG\`EBRES\par}
  \vspace{0.35em}
  {\Large\bfseries ET\par}
  \vspace{0.35em}
  {\Huge\bfseries LEURS REPR\'ESENTATIONS\par}
  \vspace{1.5em}
  {\LARGE\bfseries $C^{*}$-ALGEBRAS\par}
  \vspace{0.35em}
  {\large\bfseries AND THEIR REPRESENTATIONS\par}
  \vspace{2em}
  {\small BY\par}
  \vspace{0.6em}
  {\Large\bfseries Jacques DIXMIER\par}
  \vspace{0.3em}
  {\small Professor in the Faculty of Sciences of Paris\par}
  \vspace{2em}
  {\large\bfseries SECOND EDITION\par}
  \vfill
  {\large PARIS\par}
  {\large GAUTHIER-VILLARS, PUBLISHER\par}
  {\large 1969\par}
\end{center}
\end{titlepage}

% English hyperlinked table of contents, inserted immediately after the title page.
\clearpage
\setcounter{page}{2}
\pdfbookmark[0]{Contents}{bookmark:contents}
\begingroup
\sloppy
\tableofcontents
\endgroup

% Source PDF physical page 002
\clearpage
\thispagestyle{fancy}
\vspace*{\fill}
\begin{center}
  {\large TO MARC, \'{E}DITH AND FRAN\c{C}OISE}
\end{center}
\vspace*{\fill}

% Source PDF physical page 003

% Source PDF physical page 004
\chapter*{Preface to the Second Edition}
\phantomsection
\addcontentsline{toc}{chapter}{Preface to the Second Edition}

This second edition was produced essentially by photographic means. This naturally ruled out any major revision. I have confined myself to making a few indispensable corrections (some of them due to J. Dieudonn\'{e}), removing the problems solved since 1964, and supplementing the bibliography.

% Source PDF physical page 005

% Source PDF physical page 006
\chapter*{Introduction}
\phantomsection
\addcontentsline{toc}{chapter}{Introduction}

Let $H$ be a Hilbert space, and let $\mathcal L(H)$ be the set of continuous linear operators on $H$. Consider a subset $A$ of $\mathcal L(H)$ that is stable under addition, multiplication, multiplication by scalars, and adjunction; suppose that $A$ is closed in the operator norm. Then $A$ is an involutive Banach algebra of a special kind. Such an algebra is called a $C^{*}$-algebra.

The theory began in 1943, when Gelfand and Naimark showed that, among involutive Banach algebras, $C^{*}$-algebras can be characterized by simple axioms. It was subsequently realized that $C^{*}$-algebras play a universal role in the study of representations of a very broad class of involutive Banach algebras: for every algebra $B$ in this class, one can construct a $C^{*}$-algebra $A$ such that the representations of $B$ on a Hilbert space can be identified with the representations of $A$. For many questions (especially those involving ideals), $A$ is easier to handle than $B$. In particular, this construction applies when $B$ is taken to be the algebra of integrable functions on a locally compact group $G$. The study of the unitary representations of $G$ is thereby reduced to that of the representations of a certain $C^{*}$-algebra, called the $C^{*}$-algebra of $G$.

The study of $C^{*}$-algebras occupies nearly four-fifths of this book. We present the principal results, due especially to the work of Fell, Glimm, Kadison, Kaplansky, Mackey, Segal, and others. It seemed a pity not to take advantage of the material thus accumulated, together with the material contained in my book on von Neumann algebras (\emph{Cahiers Scientifiques}, fasc. XXV; cited as \cite{refAvN}), to say a few words about unitary representations of groups. This was all the more appropriate because group theory supplies the most interesting examples of $C^{*}$-algebras. The final pages of the book, however, by no means constitute a treatise on group representations. To gain an idea of the questions that ought to be treated in such a work, the reader may consult Mackey's lecture, \emph{Infinite dimensional
% Source PDF physical page 007
group representations} (\emph{Bull. Amer. Math. Soc.}, vol.~69, 1963, pp.~629--686); of the sixty pages of that lecture, only the first twelve concern the questions treated in the present book. Moreover, Gelfand and the Russian school on the one hand, and Harish-Chandra on the other, are scarcely cited here; this indicates how incomplete my account is as regards groups.

Although von Neumann algebras are examples of $C^{*}$-algebras, the theory of $C^{*}$-algebras is in fact dependent on that of von Neumann algebras, as will be seen repeatedly. To spare the reader from having to refer constantly to \cite{refAvN}, the indispensable results on von Neumann algebras have been collected in Appendix~A. I admit that I took the opportunity to reformulate certain statements so as to make them more readily usable. Almost always this reformulation is very easy; in the exceptional cases where it requires some work, the proof is indicated.

Results of various kinds used in the course of the text are, on the other hand, collected in Appendix~B. That appendix is, of course, highly miscellaneous.

When writing \cite{refAvN}, it seemed to me that most of the theorems presented there had a more or less definitive form. By contrast, the theory of $C^{*}$-algebras seems to me far from having reached a stable state.

\begin{center}
$*\ *\ *$
\end{center}

Each section ends with supplements, given without proof or with a sketch of proof. Some of these easy supplements may serve as exercises, but it was difficult to separate them sharply from the others; I have merely marked with an asterisk those supplements whose proof is genuinely long (for example, more than three pages). The reader is asked not to be overly exacting about this classification. The supplements are not used in subsequent sections.

The bibliography contains the essentials as regards $C^{*}$-algebras. It is, however, very incomplete as regards groups.

As with \cite{refAvN}, reading the present book presupposes a good knowledge of general topology, topological vector spaces, and integration. Proofs concerning $C^{*}$-algebras are given in fairly full detail; those concerning groups are more condensed and assume that the reader is familiar with the properties of convolution. It has also been assumed that the reader already knows the theory of locally compact
% Source PDF physical page 008
abelian groups: although, in principle, that theory is not invoked (except for occasional points of detail), reference to the abelian case sheds much light on the problems studied.

The results are numbered lexicographically. A reference such as \XRef{4.7.5}{4.7.5} is self-explanatory. References such as \XRef{A.8}{A~8} and \XRef{B.15}{B~15} concern Appendices A and B.

\begin{center}
$*\ *\ *$
\end{center}

Finally, a few words on the relation between this book and those of Naimark and Rickart. The treatises of Naimark and Rickart contain, among other things, the general theory of normed algebras, which will not be found here. On the other hand, on the much more specialized subject of $C^{*}$-algebras, the present book is more detailed. There are of course developments common to all three books, but they are fairly limited. The general results on normed algebras used here are few in number and are cited in Appendix~B.

% Source PDF physical page 009

% Source PDF physical page 010
\clearpage
\begin{center}
  {\Large\bfseries $C^{*}$-Algebras and Their Representations\par}
  \vspace{1.2em}
  {\large\bfseries $C^{*}$-Algebras\par}
\end{center}
\phantomsection
\hypertarget{division:cstar}{}
\addcontentsline{toc}{part}{\texorpdfstring{$C^{*}$-Algebras}{C*-Algebras}}

\BookSectionHere{1}{Involutive Normed Algebras}

\BookSubsection{1.1}{Involutive Algebras}

\ResultHeading{1.1.1}{1.1.1. Definition.}
\begin{itshape}
Let $A$ be an algebra over the field $\mathbb C$ of complex numbers. An \emph{involution} on $A$ is a map $x\mapsto x^{*}$ from $A$ into itself such that
\begin{align*}
\textup{(i)}\quad &(x^{*})^{*}=x,\\
\textup{(ii)}\quad &(x+y)^{*}=x^{*}+y^{*},\\
\textup{(iii)}\quad &(\lambda x)^{*}=\overline{\lambda}\,x^{*},\\
\textup{(iv)}\quad &(xy)^{*}=y^{*}x^{*},
\end{align*}
for all $x,y\in A$ and $\lambda\in\mathbb C$. An algebra over $\mathbb C$ equipped with an involution is called an \emph{involutive algebra}.
\end{itshape}

One often says that $x^{*}$ is the adjoint of $x$. A subset of $A$ stable under the involution is called self-adjoint.

Property (i) implies that an involution on $A$ is necessarily a bijection of $A$ onto itself.

\ResultHeading{1.1.2}{1.1.2. Examples.}
\begin{enumerate}[label=\textup{(\arabic*)}]
\item On $A=\mathbb C$, the map $z\mapsto\overline z$ (the complex conjugate of $z$) is an involution that makes $A$ a commutative involutive algebra.

\item Let $X$ be a locally compact space, and let $A$ be the algebra of continuous complex-valued functions on $X$ that vanish at infinity. Equipped with the map $f\mapsto\overline f$, it is a commutative involutive algebra. When $X$ consists of a single point, this gives Example~(1).

\item Let $H$ be a Hilbert space, and let $A=\mathcal L(H)$ be the algebra of continuous endomorphisms of $H$. Equipped with the usual adjunction operation, it is an involutive algebra. Examples~(2) and (3) will play a
% Source PDF physical page 011
fundamental role. (Throughout this book, ``Hilbert space'' means ``complex Hilbert space''.)

\item Let $G$ be a unimodular locally compact group, and let $A$ be the algebra $L^{1}(G)$ (with convolution). For every $f\in L^{1}(G)$, put
\[
f^{*}(s)=\overline{f(s^{-1})}\qquad(s\in G).
\]
Equipped with the map $f\mapsto f^{*}$, $A$ is an involutive algebra.
\end{enumerate}

\ResultHeading{1.1.3}{1.1.3.}
We now introduce terminology inspired by Example~(3) above. Let $A$ be an involutive algebra. An element $x\in A$ is called \emph{Hermitian} if $x^{*}=x$, and \emph{normal} if $xx^{*}=x^{*}x$. A Hermitian idempotent is called a \emph{projection}. Every Hermitian element is normal. The set of Hermitian elements is a real vector subspace of $A$. If $x$ and $y$ are Hermitian, then $(xy)^{*}=y^{*}x^{*}=yx$, and hence $xy$ is Hermitian if $x$ and $y$ commute. For every $x\in A$, the elements $xx^{*}$ and $x^{*}x$ are Hermitian [but a Hermitian element does not in general have either of these forms, as Example~(1) above shows].

\ResultHeading{1.1.4}{1.1.4.}
Every $x\in A$ can be written uniquely as $x_1+ix_2$, with $x_1,x_2$ Hermitian. [In Example~(1) above, this is the decomposition of a complex number into its real and imaginary parts.] Indeed, if we put
\[
x_1=\frac12(x+x^{*}),\qquad x_2=\frac1{2i}(x-x^{*}),
\]
then $x_1$ and $x_2$ are Hermitian and $x=x_1+ix_2$. Conversely, if $x=x_1+ix_2$ with $x_1$ and $x_2$ Hermitian, then $x^{*}=x_1-ix_2$, so
\[
x_1=\frac12(x+x^{*}),\qquad x_2=\frac1{2i}(x-x^{*}),
\]
which proves the assertion. Observe that
\begin{align*}
xx^{*}&=x_1^2+x_2^2+i(x_2x_1-x_1x_2),\\
x^{*}x&=x_1^2+x_2^2-i(x_2x_1-x_1x_2);
\end{align*}
hence $x$ is normal if and only if $x_1$ and $x_2$ commute.

\ResultHeading{1.1.5}{1.1.5.}
If $A$ has a left identity $1$, then, for every $x\in A$,
\[
x1^{*}=(1x^{*})^{*}=x^{**}=x;
\]
thus $1^{*}$ is a right identity, and consequently $1=1^{*}$ is an identity element of $A$. If $x$ is invertible in $A$, then
\[
(x^{-1})^{*}x^{*}=(xx^{-1})^{*}=1^{*}=1,
\qquad
x^{*}(x^{-1})^{*}=(x^{-1}x)^{*}=1^{*}=1;
\]
hence $x^{*}$ is invertible and $(x^{*})^{-1}=(x^{-1})^{*}$. Conversely, if $x^{*}$ is invertible, then $x^{**}=x$ is invertible. Since
\[
(x-\lambda 1)^{*}=x^{*}-\overline\lambda 1
\]
for every $\lambda\in\mathbb C$, it follows that
\[
\operatorname{Sp}_{A}x^{*}=\overline{\operatorname{Sp}_{A}x}.
\]
% Source PDF physical page 012
(In any unital algebra $B$, the \emph{spectrum} of an element $x$, denoted by $\operatorname{Sp}_{B}x$ or $\operatorname{Sp}x$, is the set of scalars $\lambda$ such that $x-\lambda 1$ is not invertible.)

An element $x\in A$ is called \emph{unitary} if $xx^{*}=x^{*}x=1$, or equivalently if $x$ is invertible and $x=(x^{*})^{-1}$. [In Example~(1) above, the unitary elements are the complex numbers of modulus~1.] The unitary elements of $A$ form a group under multiplication, the \emph{unitary group} of $A$; indeed, if $x$ and $y$ are unitary elements of $A$, then
\[
(xy)^{*-1}=(y^{*}x^{*})^{-1}=x^{*-1}y^{*-1}=xy,
\]
so $xy$ is unitary, and
\[
(x^{-1})^{*-1}=(x^{*-1})^{-1}=x^{-1},
\]
so $x^{-1}$ is unitary.

\ResultHeading{1.1.6}{1.1.6.}
Let $A$ be an involutive algebra. If $\widetilde A$ is the algebra obtained from $A$ by adjoining an identity element, there is a unique involution on $\widetilde A$ extending that on $A$, as is readily seen; it is defined by
\[
(\lambda,x)^{*}=(\overline\lambda,x^{*})
\qquad(\lambda\in\mathbb C,\ x\in A).
\]
For every $x\in A$,
\[
\operatorname{Sp}'_{A}x=\operatorname{Sp}_{\widetilde A}x.
\]
(In any algebra $B$, whether or not it has an identity, $\operatorname{Sp}'_{B}x$ or $\operatorname{Sp}'x$ denotes $\operatorname{Sp}_{\widetilde B}x$, where $\widetilde B$ is the algebra obtained from $B$ by adjoining an identity. One has $0\in\operatorname{Sp}'_{B}x$ for every $x\in B$.)

\ResultHeading{1.1.7}{1.1.7.}
Let $A$ and $B$ be two involutive algebras. A \emph{morphism} (respectively, an \emph{isomorphism}) from $A$ into $B$ is a map (respectively, a bijection) $\varphi$ from $A$ into $B$ such that
\begin{gather*}
\varphi(x+y)=\varphi(x)+\varphi(y),\qquad
\varphi(\lambda x)=\lambda\varphi(x),\\
\varphi(xy)=\varphi(x)\varphi(y),\qquad
\varphi(x^{*})=\varphi(x)^{*}
\end{gather*}
for all $x,y\in A$ and $\lambda\in\mathbb C$. In particular, $\varphi$ is a morphism from the algebra underlying $A$ into the algebra underlying $B$. When there is a risk of confusion, we shall specify ``morphism for the involutive-algebra structure'' or ``morphism for the algebra structure.''

\ResultHeading{1.1.8}{1.1.8.}
Let $A$ be an involutive algebra. An \emph{involutive subalgebra} of $A$ is a subalgebra of $A$ stable under the involution. Every intersection of involutive subalgebras is an involutive subalgebra; hence, if $M$ is any subset of $A$, there is a smallest involutive subalgebra of $A$ containing $M$, namely the intersection of all involutive subalgebras of $A$ containing $M$. It is called the involutive subalgebra of $A$ generated by $M$; it is the set of linear combinations of elements of the form $x_1x_2\cdots x_n$, where $x_1,x_2,\ldots,x_n\in M\cup M^{*}$. When $M$ consists of a single element $x$, this subalgebra is commutative if and only if $x$ is normal.

Let $A$ be an involutive algebra, and let $B$ be a self-adjoint two-sided ideal of $A$. The involution of $A$ induces, on passage to the quotient, an involution on the algebra $A/B$, and the canonical map from $A$ onto $A/B$ is a morphism.

% Source PDF physical page 013
Every product of involutive algebras is in the evident way an involutive algebra.

The algebra opposite to an involutive algebra, equipped with the same involution, is an involutive algebra.

\ResultHeading{1.1.9}{1.1.9.}
Let $A$ be an involutive algebra and let $M$ be a self-adjoint subset of $A$. Then the commutant $M'$ of $M$ in $A$ is an involutive subalgebra of $A$. Consequently, the bicommutant $M''$ of $M$ in $A$ is an involutive subalgebra of $A$ containing $M$; in general it differs from the involutive subalgebra of $A$ generated by $M$ (for example, when $A$ is commutative). If the elements of $M$ commute pairwise, then $M\subset M'$, hence $M'\supset M''$, and therefore $M''$ is commutative. Let $x\in A$ and $M=\{x,x^{*}\}$; then $M''$ is commutative if and only if $x$ is normal.

\ResultHeading{1.1.10}{1.1.10.}
Let $A$ be an involutive algebra. If $f$ is a linear form on $A$, the function
\[
x\longmapsto\overline{f(x^{*})}
\]
on $A$ is a linear form on $A$, denoted by $f^{*}$ and called the adjoint of $f$. One has
\[
f^{**}=f,\qquad (f+f')^{*}=f^{*}+f'^{*},\qquad
(\lambda f)^{*}=\overline\lambda f^{*}\quad(\lambda\in\mathbb C).
\]
The form $f$ is called Hermitian if $f=f^{*}$. Every linear form $f$ on $A$ can be written uniquely as $f_1+if_2$, with $f_1,f_2$ Hermitian; namely,
\[
f_1=\frac12(f+f^{*}),\qquad f_2=\frac1{2i}(f-f^{*}).
\]
A linear form $f$ on $A$ is Hermitian if and only if it is real-valued on the set $A_h$ of Hermitian elements of $A$. The map $f\mapsto f|_{A_h}$ is an isomorphism from the real vector space of Hermitian forms onto the dual of the real vector space $A_h$. If $A$ is commutative and $\chi$ is a character of $A$, then $\chi^{*}$ is a character of $A$.

Bibliography: \cite{ref108,ref126}.

\BookSubsection{1.2}{Involutive Normed Algebras}

\ResultHeading{1.2.1}{1.2.1. Definition.}
\begin{itshape}
An \emph{involutive normed algebra} is a normed algebra $A$ equipped with an involution $x\mapsto x^{*}$ such that $\|x^{*}\|=\|x\|$ for every $x\in A$. If, in addition, $A$ is complete, it is called an \emph{involutive Banach algebra}.
\end{itshape}

\ResultHeading{1.2.2}{1.2.2. Examples.}
The four examples in \XRef{1.1.2}{1.1.2} are examples of involutive Banach algebras when the norms are defined as follows: in Example~(1), take $\|z\|=|z|$ for every $z\in\mathbb C$; in Example~(2), take $\|f\|=\sup_{t\in X}|f(t)|$ for every $f\in A$; in Example~(3), take the usual norm on $\mathcal L(H)$; in Example~(4), take
\[
\|f\|=\int_G |f(g)|\,dg
\]
for every $f\in L^{1}(G)$.

% Source PDF physical page 014
\ResultHeading{1.2.3}{1.2.3.}
Let $A$ be an involutive normed algebra, and let $\widetilde A$ be the involutive algebra obtained from $A$ by adjoining an identity element. There are norms on $\widetilde A$ extending that of $A$ and making $\widetilde A$ an involutive normed algebra (for example, one may put $\|(\lambda,x)\|=|\lambda|+\|x\|$ for $\lambda\in\mathbb C$, $x\in A$). Every involutive normed algebra obtained in this way is called an involutive normed algebra obtained from $A$ by adjoining an identity element.

\ResultHeading{1.2.4}{1.2.4.}
Let $A$ and $B$ be two involutive normed algebras. A morphism from $A$ into $B$ will simply mean a morphism of the underlying involutive algebras (with no condition on the norms). By contrast, an isomorphism will mean an isomorphism of the underlying involutive algebras that preserves the norms.

\ResultHeading{1.2.5}{1.2.5.}
Let $A$ be an involutive normed algebra. The closure of an involutive subalgebra is an involutive subalgebra. If $M\subset A$, the smallest closed involutive subalgebra $B$ containing $M$ is called the \emph{closed involutive subalgebra generated by $M$}; it is the closure of the involutive subalgebra generated by $M$. If $M$ consists of a normal element, $B$ is commutative.

The quotient of an involutive normed algebra by a closed self-adjoint two-sided ideal, the product of finitely many involutive normed algebras, the opposite of an involutive normed algebra, and the completion of an involutive normed algebra are naturally involutive normed algebras.

\ResultHeading{1.2.6}{1.2.6.}
Let $A$ be an involutive normed algebra. If $f$ is a continuous linear form on $A$, then $f^{*}$ is continuous and $\|f^{*}\|=\|f\|$, because the unit ball of $A$ is self-adjoint. The set $A_h$ of Hermitian elements of $A$ is a real normed vector space. Let $f$ be a continuous Hermitian linear form on $A$, and put $g=f|_{A_h}$. Then $\|f\|=\|g\|$. Indeed, it is clear that $\|f\|\geq\|g\|$. On the other hand, for every $\varepsilon>0$, there exists $x\in A$ such that $\|x\|\leq1$ and $|f(x)|\geq\|f\|-\varepsilon$; multiplying $x$ by a scalar of modulus~1, we may suppose that $f(x)\geq0$. Then
\[
\left|g\!\left(\frac12(x+x^{*})\right)\right|
=\frac12\left|f(x)+\overline{f(x)}\right|
=f(x)\geq\|f\|-\varepsilon.
\]
Since $\left\|\frac12(x+x^{*})\right\|\leq1$, it follows that $\|g\|\geq\|f\|-\varepsilon$, proving the assertion. We shall therefore identify continuous Hermitian linear forms on $A$ with continuous real linear forms on $A_h$.

Bibliography: \cite{ref108,ref126}.

% Source PDF physical page 015
\BookSubsection[C*-Algebras]{1.3}{$C^{*}$-Algebras}

\ResultHeading{1.3.1}{1.3.1. Definition.}
\begin{itshape}
A $C^{*}$-algebra is an involutive Banach algebra $A$ such that
\[
\|x\|^{2}=\|x^{*}x\|
\]
for every $x\in A$.
\end{itshape}

\ResultHeading{1.3.2}{1.3.2.}
Examples~(1), (2), and (3) of \XRef{1.1.2}{1.1.2} and \XRef{1.2.2}{1.2.2} are examples of $C^{*}$-algebras. By contrast, Example~(4) is not in general a $C^{*}$-algebra.

\ResultHeading{1.3.3}{1.3.3.}
If $A$ is a $C^{*}$-algebra, every closed involutive subalgebra of $A$ is a $C^{*}$-algebra. In particular, if $H$ is a Hilbert space, every closed involutive subalgebra of $\mathcal L(H)$ is a $C^{*}$-algebra. We shall see later (\XRef{2.6.1}{2.6.1}) that every $C^{*}$-algebra is isomorphic to a $C^{*}$-algebra of this kind; it was in connection with this example that the theory of $C^{*}$-algebras was created.

Let $(A_i)_{i\in I}$ be a family of $C^{*}$-algebras. Let $A$ be the set of families $(x_i)_{i\in I}$ such that $x_i\in A_i$ for every $i\in I$ and
\[
\sup_{i\in I}\|x_i\|<+\infty.
\]
With the operations
\begin{align*}
(x_i)+(y_i)&=(x_i+y_i), & \lambda(x_i)&=(\lambda x_i),\\
(x_i)(y_i)&=(x_iy_i), & (x_i)^{*}&=(x_i^{*}),
\end{align*}
and the norm
\[
\|(x_i)\|=\sup_i\|x_i\|,
\]
one verifies at once that $A$ is a $C^{*}$-algebra, called the product $C^{*}$-algebra of the $A_i$. Notice that the set $A$ is not the full Cartesian product of the $A_i$.

Let $A$ be a $C^{*}$-algebra. Retain all the operations and the norm of $A$, except for the multiplication $(x,y)\mapsto xy$, which is replaced by $(x,y)\mapsto yx$. In other words, consider the involutive normed algebra $A^{\mathrm o}$ opposite to $A$. It is immediate that $A^{\mathrm o}$ is a $C^{*}$-algebra.

The question of quotient $C^{*}$-algebras is more delicate (see \XRef{1.8.2}{1.8.2}).

\ResultHeading{1.3.4}{1.3.4.}
Let $A$ be a Banach algebra equipped with an involution such that
\[
\|x\|^{2}\leq\|x^{*}x\|.
\]
It follows that $\|x\|^{2}\leq\|x^{*}\|\,\|x\|$, whence $\|x\|\leq\|x^{*}\|$; replacing $x$ by $x^{*}$ shows that $\|x^{*}\|=\|x\|$. The hypothesis then implies
\[
\|x\|^{2}\leq\|x^{*}x\|\leq\|x^{*}\|\,\|x\|=\|x\|^{2},
\]
so $A$ is a $C^{*}$-algebra.

\ResultHeading{1.3.5}{1.3.5.}
Let $A$ be a $C^{*}$-algebra. For every $x\in A$,
\[
\|x\|=\sup_{\|x'\|\leq1}\|xx'\|.
\]
% Source PDF physical page 016
Indeed, it is clear that $\|x'\|\leq1$ implies $\|xx'\|\leq\|x\|$. To prove the reverse inequality, we may suppose that $\|x\|=1$; then $\|x^{*}\|=1$, and
\[
\sup_{\|x'\|\leq1}\|xx'\|\geq\|xx^{*}\|=\|x\|^{2}=1.
\]

\ResultHeading{1.3.6}{1.3.6.}
Let $A$ be a unital $C^{*}$-algebra. One has
\[
\|1\|^{2}=\|1^{*}1\|=\|1\|,
\]
and therefore $\|1\|=0$ or $1$. Thus, if $A\neq0$, then $\|1\|=1$. It follows that every unitary element $u$ of $A$ satisfies
\[
\|u\|=\|u^{*}u\|^{1/2}=1
\]
when $A\neq0$. 

Recall that no continuity condition was imposed in the definition of morphisms of involutive normed algebras. We shall see that, for $C^{*}$-algebras, every morphism is automatically continuous. More precisely:

\ResultHeading{1.3.7}{1.3.7. Proposition.}
\begin{itshape}
Let $A$ be an involutive Banach algebra, let $B$ be a $C^{*}$-algebra, and let $\pi$ be a morphism from $A$ into $B$. Then
\[
\|\pi(x)\|\leq\|x\|
\]
for every $x\in A$.
\end{itshape}

\Proof
For every Hermitian element $y$ of $B$, one has
\[
\|y^2\|=\|y^{*}y\|=\|y\|^2,
\]
and hence, by induction, $\|y^{2^n}\|^{2^{-n}}=\|y\|$. As $n\to+\infty$, the left-hand side tends to the spectral radius $\rho(y)$ of $y$ (see \XRef{B.1}{B~1}); hence
\begin{equation}
\rho(y)=\|y\|.
\tag{1}\label{eq:1.3.7-1}
\end{equation}
Now, for every $x\in A$,
\[
\operatorname{Sp}'_{B}\pi(x)\subset\operatorname{Sp}'_{A}x,
\qquad
\rho(\pi(x))\leq\rho(x)\leq\|x\|.
\]
Consequently, in view of \eqref{eq:1.3.7-1},
\[
\|\pi(x)\|^{2}
=\|\pi(x^{*}x)\|
=\rho(\pi(x^{*}x))
\leq\|x^{*}x\|
\leq\|x^{*}\|\,\|x\|
=\|x\|^{2}.
\]

The following proposition will make it possible, in almost every case, to reduce the study of $C^{*}$-algebras to that of $C^{*}$-algebras with an identity element.

\ResultHeading{1.3.8}{1.3.8. Proposition.}
\begin{itshape}
Let $A$ be a $C^{*}$-algebra, and let $\widetilde A$ be the involutive algebra obtained from $A$ by adjoining an identity element. There is a unique norm on $\widetilde A$ extending that of $A$ and making $\widetilde A$ a $C^{*}$-algebra.
\end{itshape}

\Proof
Uniqueness follows from \XRef{1.3.7}{1.3.7}. We prove existence. Suppose first that $A$ has an identity element $e$, and denote the identity element of $\widetilde A$ by $1$. In $\widetilde A$, the spaces $A$ and $\mathbb C(1-e)$ are supplementary self-adjoint two-sided ideals; consequently there is an isomorphism of the involutive algebra $\widetilde A$ onto the involutive algebra $\mathbb C\times A$ that carries $A$ onto $\{0\}\times A$. We may therefore equip $\mathbb C\times A$ with its product $C^{*}$-algebra structure. Suppose
% Source PDF physical page 017
from now on that $A$ has no identity element. Every element $x$ of $\widetilde A$ defines an operator $L_x$ of left multiplication on the two-sided ideal $A$ of $\widetilde A$; put
\[
\|x\|=\|L_x\|.
\]
For $x\in A$, this gives back the given norm on $A$, by \XRef{1.3.5}{1.3.5}. It is also clear that $x\mapsto\|x\|$ is a seminorm on $\widetilde A$ and that $\|xy\|\leq\|x\|\,\|y\|$. This seminorm is a norm. Indeed, let $x=\lambda1-x'$ ($\lambda\in\mathbb C$, $x'\in A$) be an element of $\widetilde A$ such that $\|L_x\|=0$, that is, $xy=0$ for every $y\in A$, and let us prove that $x=0$. If $\lambda\neq0$, then for every $y\in A$,
\[
0=\lambda^{-1}xy=y-\lambda^{-1}x'y;
\]
thus $\lambda^{-1}x'$ is a left identity of $A$, and $A$ has an identity element by \XRef{1.1.5}{1.1.5}, contrary to the hypothesis. Therefore $\lambda=0$, so $x\in A$, and $\|x\|=0$ implies $x=0$. We have thus proved that $x\mapsto\|x\|$ is a norm on $\widetilde A$. Since $A$ is complete and has codimension one in $\widetilde A$, the latter is complete. It remains to prove, by \XRef{1.3.4}{1.3.4}, that
\[
\|x\|^{2}\leq\|x^{*}x\|
\]
for every $x\in\widetilde A$, and it suffices to do so when $\|x\|=1$. For every $r<1$, there is a $y\in A$ such that $\|y\|\leq1$ and $\|xy\|^{2}\geq r$; then, since $xy\in A$,
\[
\|x^{*}x\|
\geq\|y^{*}(x^{*}x)y\|
=\|(xy)^{*}(xy)\|
=\|xy\|^{2}
\geq r,
\]
and hence $\|x^{*}x\|\geq1$.

\ResultHeading{1.3.9}{1.3.9. Proposition.}
\begin{itshape}
Let $A$ be a $C^{*}$-algebra.
\begin{enumerate}[label=\textup{(\roman*)}]
\item If $h$ is a Hermitian element of $A$, every point of $\operatorname{Sp}'h$ is real.
\item If $A$ has an identity element and $u$ is a unitary element of $A$, every point of $\operatorname{Sp}u$ has modulus~$1$.
\end{enumerate}
\end{itshape}

\Proof
By \XRef{1.3.8}{1.3.8}, in proving both parts we may suppose that $A$ has an identity element (and that $A\neq0$). By \XRef{1.3.6}{1.3.6},
\[
\|u\|=\|u^{-1}\|=1;
\]
thus $\rho(u)\leq1$ and $\rho(u^{-1})\leq1$, so both $\operatorname{Sp}u$ and $\operatorname{Sp}(u^{-1})=(\operatorname{Sp}u)^{-1}$ lie in the closed unit disk of $\mathbb C$, which proves (ii). On the other hand,
\[
(\exp(ih))^{*}
=\left(1+ih+\frac{i^{2}h^{2}}{2!}+\cdots\right)^{*}
=1+(-i)h+\frac{(-i)^{2}h^{2}}{2!}+\cdots
=\exp(-ih),
\]
so $\exp(ih)$ is unitary. Hence, if $z\in\operatorname{Sp}h$, then $|\exp(iz)|=1$ (\XRef{B.4}{B~4}), and therefore $z\in\mathbb R$. (Throughout the book, $\mathbb R$ denotes the set of real numbers.)

The following proposition shows that, in the case of $C^{*}$-algebras, there is no need to be concerned with the ambient algebra when defining the spectrum of an element.

\ResultHeading{1.3.10}{1.3.10. Proposition.}
\begin{itshape}
Let $A$ be a $C^{*}$-algebra, let $B$ be a $C^{*}$-subalgebra, and let $x\in B$.
\begin{enumerate}[label=\textup{(\roman*)}]
\item $\operatorname{Sp}'_{A}x=\operatorname{Sp}'_{B}x$.
\item If $A$ has an identity element belonging to $B$, then $\operatorname{Sp}_{A}x=\operatorname{Sp}_{B}x$.
\end{enumerate}
\end{itshape}

\Proof
By adjoining an identity element, one sees that (i) follows from (ii). We prove (ii). If $x$ is Hermitian, then $\operatorname{Sp}_{B}x\subset\mathbb R$ by \XRef{1.3.9}{1.3.9}, and hence $\operatorname{Sp}_{B}x=\operatorname{Sp}_{A}x$ by \XRef{B.2}{B~2}. In the general case, if $x\in B$ is invertible in $A$, then $xx^{*}$ is invertible in $A$, hence in $B$ by what precedes, and therefore $x$ is right-invertible in $B$. Similarly, $x$ is left-invertible in $B$, and is thus invertible in $B$. Applying this result to $x-\lambda1$, where $\lambda\in\mathbb C$, proves (ii).

$C^{*}$-algebras are called completely regular algebras in \cite{ref108}, and $B^{*}$-algebras by many authors.

Bibliography: \cite{ref36,ref108,ref123,ref125,ref126}.

% Source PDF physical page 018
\BookSubsection[Commutative C*-Algebras]{1.4}{Commutative $C^{*}$-Algebras}

Their theory is completed by the following theorem:

\ResultHeading{1.4.1}{1.4.1. Theorem.}
\begin{itshape}
Let $A$ be a commutative $C^{*}$-algebra, let $S$ be its spectrum (which is a locally compact space), and let $B$ be the $C^{*}$-algebra of continuous complex-valued functions on $S$ that vanish at infinity.
\begin{enumerate}[label=\textup{(\roman*)}]
\item Every character of $A$ is Hermitian.
\item The Gelfand transform is an isomorphism of the $C^{*}$-algebra $A$ onto the $C^{*}$-algebra $B$.
\end{enumerate}
\end{itshape}

\Proof
Let $\chi$ be a character of $A$. If $y$ is a Hermitian element of $A$, then
\[
\chi(y)\in\operatorname{Sp}'y\subset\mathbb R
\]
by \XRef{1.3.9}{1.3.9}. If $x$ is any element of $A$, write $x=x_1+ix_2$ with $x_1,x_2$ Hermitian; then
\[
\chi(x^{*})=\chi(x_1-ix_2)
=\chi(x_1)-i\chi(x_2)
=\overline{\chi(x)},
\]
which proves (i).

In other words, if $\mathcal G$ denotes the Gelfand transform, then
\[
\mathcal G(x^{*})=\overline{\mathcal G(x)}
\]
for every $x\in A$. Moreover, $\mathcal G(A)$ separates the points of $S$, and at every point of $S$ there is a function in $\mathcal G(A)$ that does not vanish (\XRef{B.3}{B~3}). The Stone--Weierstrass theorem therefore shows that $\mathcal G(A)$ is dense in $B$. To complete the proof, it remains only to show that $\mathcal G$ is isometric. But $\|\mathcal G(y)\|=\|y\|$ for Hermitian $y$ by \XRef{1.3.7}{1.3.7}, formula~\eqref{eq:1.3.7-1}; hence, for every $x\in A$,
\[
\|x\|^{2}
=\|x^{*}x\|
=\|\mathcal G(x^{*}x)\|
=\|\overline{\mathcal G(x)}\,\mathcal G(x)\|
=\|\mathcal G(x)\|^{2}.
\]

\ResultHeading{1.4.2}{1.4.2.}
Retain the preceding notation. If $f\in B$, then $g(f)$ has a meaning for every continuous complex-valued function $g$ on
\[
f(S)\cup\{0\}=\operatorname{Sp}'_{B}f
\]
that vanishes at $0$, and it represents an element of $B$. Thus, if $x\in A$, one may define $h(x)\in A$ for every continuous complex-valued function $h$ on $\operatorname{Sp}'_{A}x$ that vanishes at $0$. In \XRef{1.5.1}{1.5.1} we shall see that such a ``functional calculus'' can be defined for normal elements of an arbitrary $C^{*}$-algebra (not necessarily commutative). It will be a very useful tool in what follows.

% Source PDF physical page 019
\ResultHeading{1.4.3}{1.4.3. Proposition.}
\begin{itshape}
Let $A$ be a unital commutative $C^{*}$-algebra, let $S$ be its spectrum, and let $x\in A$. Suppose that the $C^{*}$-subalgebra of $A$ generated by $1$ and $x$ is equal to $A$. Then
\[
\chi\longmapsto\chi(x)
\]
is a homeomorphism from $S$ onto $\operatorname{Sp}_{A}x$.
\end{itshape}

\Proof
This map is continuous and has image $\operatorname{Sp}_{A}x$ (\XRef{B.3}{B~3}). Let $\chi,\chi'\in S$ satisfy $\chi(x)=\chi'(x)$. Since $\chi$ and $\chi'$ are continuous and Hermitian, the set $A'$ of all $y\in A$ such that $\chi(y)=\chi'(y)$ is a $C^{*}$-subalgebra of $A$ containing $x$. Thus $A'=A$ and $\chi'=\chi$. The map in question is injective and hence is a homeomorphism, since $S$ is compact.

Bibliography: \cite{ref1,ref41,ref108,ref126}.

\BookSubsection[Functional Calculus in C*-Algebras]{1.5}{Functional Calculus in $C^{*}$-Algebras}

\ResultHeading{1.5.1}{1.5.1. Theorem.}
\begin{itshape}
Let $A$ be a unital $C^{*}$-algebra, let $x$ be a normal element of $A$, let $S=\operatorname{Sp}_{A}x$, and let $A'$ be the $C^{*}$-algebra of continuous complex-valued functions on $S$. There is a unique morphism $\varphi$ from $A'$ into $A$ such that
\[
\varphi(1)=1,\qquad \varphi(\iota)=x,
\]
where $\iota$ denotes the function $z\mapsto z$ on $S$. This morphism is isometric. Its image is the $C^{*}$-subalgebra of $A$ generated by $1$ and $x$, and hence consists of normal elements.
\end{itshape}

\Proof
The polynomials in $z$ and $\overline z$ are dense in $A'$, and every morphism from $A'$ into $A$ is continuous by \XRef{1.3.7}{1.3.7}; this proves uniqueness of $\varphi$. Let $B$ be the commutative $C^{*}$-subalgebra of $A$ generated by $1$ and $x$, let $T$ be its spectrum, let $C$ be the $C^{*}$-algebra of continuous complex-valued functions on $T$, and let $\mathcal G_B$ be the Gelfand transform for $B$, which is an isomorphism of $B$ onto $C$ by \XRef{1.4.1}{1.4.1}. Proposition~\XRef{1.4.3}{1.4.3} gives a homeomorphism from $T$ onto $\operatorname{Sp}_{B}x=S$ by \XRef{1.3.10}{1.3.10}, and hence an isomorphism $\psi:A'\to C$ carrying $\iota$ to $\mathcal G_Bx$ [for every $\chi\in T$,
\[
(\mathcal G_Bx)(\chi)=\chi(x)=\iota(\chi(x)).
\]
] Consider the composite isomorphism
\[
A'\xrightarrow{\ \psi\ }C\xrightarrow{\ \mathcal G_B^{-1}\ }B.
\]
Composing it with the canonical injection of $B$ into $A$ gives a morphism from $A'$ into $A$ with the asserted properties.

\ResultHeading{1.5.2}{1.5.2. Definition.}
\begin{itshape}
Let $A$ be a unital $C^{*}$-algebra. If $x$ is a normal element of $A$ and $f$ is a continuous complex-valued function on $\operatorname{Sp}_{A}x$, the element $\varphi(f)$ of Theorem~\XRef{1.5.1}{1.5.1} is denoted by $f(x)$.
\end{itshape}

The fact that $\varphi$ is an isometric isomorphism is expressed by the following formulas, where $f$ and $g$ are continuous complex-valued functions on $\operatorname{Sp}_{A}x$:
\begin{align}
(f+g)(x)&=f(x)+g(x),
  \tag{1}\label{eq:1.5.2-1}\\
(fg)(x)&=f(x)g(x),
  \tag{2}\label{eq:1.5.2-2}\\
\overline f(x)&=f(x)^{*},
  \tag{3}\label{eq:1.5.2-3}\\
\|f(x)\|&=\|f\|.
  \tag{4}\label{eq:1.5.2-4}
\end{align}

% Source PDF physical page 020
If $f$ is the restriction to $\operatorname{Sp}_{A}x$ of a polynomial $z\mapsto P(z,\overline z)$ in $z$ and $\overline z$, then
\[
f(x)=P(x,x^{*}),
\]
where the symbol $P(x,x^{*})$ has its usual algebraic meaning (recall that $xx^{*}=x^{*}x$).

With the notation above,
\[
\operatorname{Sp}_{A}f(x)
=\operatorname{Sp}_{B}f(x)
=\operatorname{Sp}_{A'}f
=f(S),
\]
or, equivalently,
\begin{equation}
\operatorname{Sp}_{A}f(x)=f(\operatorname{Sp}_{A}x).
\tag{5}\label{eq:1.5.2-5}
\end{equation}

\ResultHeading{1.5.3}{1.5.3. Proposition.}
\begin{itshape}
Let $A$ and $B$ be two unital $C^{*}$-algebras, let $\varphi$ be a morphism from $A$ into $B$ carrying $1$ to $1$, and let $x$ be a normal element of $A$, so that $\varphi(x)$ is normal in $B$. Let $f$ be a continuous complex-valued function on $\operatorname{Sp}_{A}x$. Denoting its restriction to $\operatorname{Sp}_{B}\varphi(x)$ again by $f$, one has
\[
\varphi(f(x))=f(\varphi(x)).
\]
\end{itshape}

\Proof
Let $C$ be the $C^{*}$-algebra of continuous complex-valued functions on $\operatorname{Sp}_{A}x$. The maps
\[
f\longmapsto\varphi(f(x)),
\qquad
f\longmapsto f(\varphi(x))
\]
are morphisms from $C$ into $B$ and take the same values when $f$ is one of the functions $z\mapsto1$, $z\mapsto z$, or $z\mapsto\overline z$. Since the $C^{*}$-subalgebra of $C$ generated by these functions is all of $C$, the two morphisms are equal.

\ResultHeading{1.5.4}{1.5.4. Corollary.}
\begin{itshape}
Let $A$ be a unital commutative $C^{*}$-algebra, let $x\in A$, let $\mathcal G$ be the Gelfand transform for $A$, and let $f$ be a continuous complex-valued function on $\operatorname{Sp}_{A}x$. Then
\[
\mathcal G(f(x))=f\circ\mathcal G(x).
\]
\end{itshape}

This follows, for example, from \XRef{1.5.3}{1.5.3} applied to $\varphi=\mathcal G$.

\ResultHeading{1.5.5}{1.5.5. Corollary.}
\begin{itshape}
Let $A$ be a unital $C^{*}$-algebra, let $x$ be a normal element of $A$, let $C$ be the $C^{*}$-algebra of continuous complex-valued functions on $\operatorname{Sp}x$, let $f\in C$, let $C'$ be the $C^{*}$-algebra of continuous complex-valued functions on
\[
\operatorname{Sp}f(x)=f(\operatorname{Sp}x),
\]
and let $g\in C'$. Then $g\circ f\in C$ and
\[
(g\circ f)(x)=g(f(x)).
\]
\end{itshape}

\Proof
The map $g\mapsto(g\circ f)(x)$ is a morphism from $C'$ into $A$ that carries $1$ to $1$ and the function $z\mapsto z$ to $f(x)$. By the uniqueness assertion in Theorem~\XRef{1.5.1}{1.5.1},
\[
(g\circ f)(x)=g(f(x)).
\]

\ResultHeading{1.5.6}{1.5.6. Proposition.}
\begin{itshape}
Let $A$ be a $C^{*}$-algebra, let $x\in A$ be normal, let $S=\operatorname{Sp}'_{A}x$, and let $A'$ be the $C^{*}$-algebra of continuous complex-valued functions on $S$ that vanish at $0$. There is a unique morphism $\varphi$ from $A'$ into $A$ such that
\[
\varphi(\iota)=x,
\]
where $\iota$ denotes the function $z\mapsto z$ on $S$. This morphism is isometric. Its image is the $C^{*}$-subalgebra of $A$ generated by $x$, and hence consists of normal elements.
\end{itshape}

\Proof
Since the polynomials in $z$ and $\overline z$ without constant terms are dense in $A'$, uniqueness of $\varphi$ is immediate. Existence follows from Theorem~\XRef{1.5.1}{1.5.1} upon adjoining an identity element to $A$.

% Source PDF physical page 021
\ResultHeading{1.5.7}{1.5.7.}
By adjoining an identity element, all the results of this subsection extend immediately, with the evident modifications, to $C^{*}$-algebras without an identity element. We shall therefore refer to these results even when the existence of an identity element has not been assumed.

In particular, let $x$ be a Hermitian element of the $C^{*}$-algebra $A$. Its spectrum is real. Consider the continuous functions of a real variable
\[
t\longmapsto f_1(t)=\sup(t,0),\qquad
t\longmapsto f_2(t)=\sup(-t,0),\qquad
t\longmapsto f_3(t)=|t|.
\]
Put
\[
x^{+}=f_1(x),\qquad x^{-}=f_2(x),\qquad |x|=f_3(x).
\]
These are Hermitian elements of $A$ (indeed, of the $C^{*}$-subalgebra of $A$ generated by $x$), because $\overline f_1=f_1$, $\overline f_2=f_2$, and $\overline f_3=f_3$. Since $f_1,f_2,f_3$ take nonnegative values, one has
\begin{equation}
\operatorname{Sp}'(x^{+})\geq0,
\qquad
\operatorname{Sp}'(x^{-})\geq0,
\qquad
\operatorname{Sp}'(|x|)\geq0.
\tag{1}\label{eq:1.5.7-1}
\end{equation}
Since
\[
f_1(t)-f_2(t)=t,
\qquad
f_1(t)+f_2(t)=f_3(t),
\qquad
f_1(t)f_2(t)=0,
\]
one has
\begin{equation}
x=x^{+}-x^{-},
\qquad
|x|=x^{+}+x^{-},
\qquad
x^{+}x^{-}=x^{-}x^{+}=0.
\tag{2}\label{eq:1.5.7-2}
\end{equation}
The norm of $|x|$ is equal to that of $x$, and the norms of $x^{+}$ and $x^{-}$ are at most that of $x$. Observe, for \XRef{1.6.4}{1.6.4}, that $x^{+}$ and $x^{-}$ are the squares of two Hermitian elements whose product is zero: it suffices to consider the functions $\sqrt{f_1}$ and $\sqrt{f_2}$.

\ResultHeading{1.5.8}{1.5.8.}
Let $A$ be a $C^{*}$-algebra. For every integer $n>0$, one has $A=A^n$ (the set of linear combinations of products of $n$ elements of $A$). Indeed, it suffices to prove that every Hermitian element $x$ of $A$ is a product of $n$ elements of $A$. If $f_1,\ldots,f_n$ are continuous real-valued functions of a real variable such that
\[
f_1(t)f_2(t)\cdots f_n(t)=t,
\qquad
f_1(0)=\cdots=f_n(0)=0,
\]
then
\[
x=f_1(x)f_2(x)\cdots f_n(x).
\]

Bibliography: \cite{ref41,ref108,ref126}.

\BookSubsection[Positive Elements in C*-Algebras]{1.6}{Positive Elements in $C^{*}$-Algebras}

\ResultHeading{1.6.1}{1.6.1. Proposition.}
\begin{itshape}
Let $A$ be a $C^{*}$-algebra, and let $x$ be a Hermitian element of $A$. The following conditions are equivalent:
\begin{enumerate}[label=\textup{(\roman*)}]
\item $\operatorname{Sp}'_{A}x\geq0$.
\item $x$ has the form $yy^{*}$ for some $y\in A$.
\item $x$ has the form $h^2$ for some Hermitian $h\in A$.
\end{enumerate}
Moreover, the set $P$ of elements satisfying these conditions is a closed convex cone such that
\[
P\cap(-P)=\{0\}.
\]
\end{itshape}

\Proof
For the proof, let $P$ denote the set of Hermitian elements of $A$ satisfying condition~(i).

(i)$\Rightarrow$(iii): if $\operatorname{Sp}'_{A}x\geq0$, one may form $h=x^{1/2}$, which is a Hermitian element of $A$, and $x=h^2$.

(iii)$\Rightarrow$(i): if $x=h^2$ with $h$ Hermitian, then
\[
\operatorname{Sp}'_{A}x=(\operatorname{Sp}'_{A}h)^2\geq0,
\]
since $\operatorname{Sp}'_{A}h$ is real.

(iii)$\Rightarrow$(ii) is evident.

To prove (ii)$\Rightarrow$(iii) and the final assertion of the proposition, we shall need the following lemmas.

% Source PDF physical page 022
\ResultHeading{1.6.2}{1.6.2. Lemma.}
\begin{itshape}
Suppose that $A$ has an identity element. If $x\in A$ is Hermitian and $\|1-x\|\leq1$, then $x\in P$. If $x\in P$ and $\|x\|\leq1$, then $\|1-x\|\leq1$.
\end{itshape}

\Proof
Considering the $C^{*}$-subalgebra of $A$ generated by $x$, we reduce first to the case where $A$ is commutative and then, by \XRef{1.4.1}{1.4.1}, to the case where $A$ is the $C^{*}$-algebra of continuous complex-valued functions on a compact space. The lemma is then evident.

\ResultHeading{1.6.3}{1.6.3. Lemma.}
\begin{itshape}
Suppose that $A$ has an identity element. Let $x$ be a Hermitian element of $A$. Then $x\in P$ if and only if
\[
\|\,\|x\|1-x\|\leq\|x\|.
\]
\end{itshape}

\Proof
We may suppose that $x\neq0$ and hence, by scaling, that $\|x\|=1$. The assertion then follows at once from Lemma~\XRef{1.6.2}{1.6.2}.

\ResultHeading{1.6.4}{1.6.4.}
Return to the setting of \XRef{1.6.1}{1.6.1}. To prove that $P$ is a closed convex cone with $P\cap(-P)=\{0\}$, we may suppose, by \XRef{1.3.8}{1.3.8}, that $A$ has an identity element. Lemma~\XRef{1.6.3}{1.6.3} implies that $P$ is closed. It is clear that $x\in P$ and $\lambda\geq0$ imply $\lambda x\in P$. Let $x,y\in P$, and let us show that $x+y\in P$. We may suppose that $\|x\|\leq1$ and $\|y\|\leq1$. Then $\|1-x\|\leq1$ and $\|1-y\|\leq1$ by \XRef{1.6.2}{1.6.2}; hence
\[
\left\|1-\frac12(x+y)\right\|
=\frac12\|1-x+1-y\|
\leq\frac12\bigl(\|1-x\|+\|1-y\|\bigr)
\leq1.
\]
Thus $\frac12(x+y)\in P$ by \XRef{1.6.2}{1.6.2}, and $x+y\in P$. If $x\in P\cap(-P)$, then $\operatorname{Sp}x=\{0\}$, so $\rho(x)=0$ and hence $x=0$ by \XRef{1.3.7}{1.3.7}, formula~\eqref{eq:1.3.7-1}.

Finally, without assuming that $A$ has an identity element, let us prove (ii)$\Rightarrow$(iii) in \XRef{1.6.1}{1.6.1}. Let $y\in A$. Write
\[
(yy^{*})^{+}=u^2,
\qquad
(yy^{*})^{-}=v^2,
\]
where $u,v$ are Hermitian elements of $A$ with $uv=0$, as in \XRef{1.5.7}{1.5.7}. Then
\[
(vy)(vy)^{*}=v(yy^{*})v=vu^2v-v^4=-v^4\in-P,
\]
% Source PDF physical page 023
because (iii)$\Rightarrow$(i). Write $vy=s+it$ with $s,t$ Hermitian. Then
\begin{align*}
(vy)^{*}(vy)
&=-(vy)(vy)^{*}+(s+it)(s-it)+(s-it)(s+it)\\
&=-(vy)(vy)^{*}+2s^2+2t^2\in P,
\end{align*}
since $-(vy)(vy)^{*}\in P$, $s^2\in P$, $t^2\in P$, and $P$ is a convex cone. Therefore $(vy)(vy)^{*}\in P$, because in every algebra the spectrum of a product does not depend on the order of its factors (\XRef{B.26}{B~26}). Thus
\[
(vy)(vy)^{*}\in P\cap(-P)=\{0\}.
\]
It follows that $v^4=0$, hence $v=0$ and $yy^{*}=u^2$. Proposition~\XRef{1.6.1}{1.6.1} is proved.

\ResultHeading{1.6.5}{1.6.5. Definition.}
\begin{itshape}
Let $A$ be a $C^{*}$-algebra. An element $x$ of $A$ is called \emph{positive}, and one writes $x\geq0$, if it is Hermitian and satisfies the three equivalent conditions of Proposition~\XRef{1.6.1}{1.6.1}. The set of positive elements of $A$ is denoted by $A^{+}$.
\end{itshape}

If $B$ is a $C^{*}$-subalgebra of $A$ and $x\in B$, then, by \XRef{1.3.10}{1.3.10}, the condition $x\geq0$ has the same meaning relative to $A$ as it does relative to $B$.

Since $A^{+}$ is a convex cone and $A^{+}\cap(-A^{+})=\{0\}$, the relation $x-y\geq0$ is an order relation on $A_h$, compatible with its real vector-space structure; this relation is denoted by $x\geq y$ or $y\leq x$. If $z$ is a Hermitian element of $A$, then $z^{+}\geq0$, $z^{-}\geq0$, and $|z|\geq0$ by \XRef{1.5.7}{1.5.7}, formulas~\eqref{eq:1.5.7-1}; $z^{+}$ (respectively, $z^{-}$) is called the positive (respectively, negative) part of $z$. By \XRef{1.5.7}{1.5.7}, formulas~\eqref{eq:1.5.7-2}, every Hermitian element of $A$ is the difference of two elements of $A^{+}$.

\ResultHeading{1.6.6}{1.6.6.}
Let $A$ be the $C^{*}$-algebra of continuous complex-valued functions vanishing at infinity on a locally compact space $T$, and let $f\in A$. The relation $f\geq0$ in $A$ has its usual meaning, by condition~(i) of Proposition~\XRef{1.6.1}{1.6.1} and the fact that
\[
\operatorname{Sp}'f=f(T)\cup\{0\}.
\]

\ResultHeading{1.6.7}{1.6.7.}
Let $H$ be a Hilbert space, let $A$ be the $C^{*}$-algebra $\mathcal L(H)$, and let $x\in A$. We show that $x\geq0$ is equivalent to
\[
\langle x\xi,\xi\rangle\geq0
\qquad\text{for every }\xi\in H,
\]
that is, to the usual definition of operators $\geq0$. If $x=y^{*}y$ with $y\in A$, then
\[
\langle x\xi,\xi\rangle
=\langle y^{*}y\xi,\xi\rangle
=\|y\xi\|^{2}\geq0
\qquad(\xi\in H).
\]
Conversely, suppose that $\langle x\xi,\xi\rangle\geq0$ for every $\xi\in H$. For every $\eta\in H$,
\begin{align*}
0&\leq\langle x(x^{-}\eta),x^{-}\eta\rangle\\
 &=\langle(x^{+}-x^{-})(x^{-}\eta),x^{-}\eta\rangle\\
 &=-\langle(x^{-})^{2}\eta,x^{-}\eta\rangle
  =-\langle(x^{-})^{3}\eta,\eta\rangle.
\end{align*}
Since $(x^{-})^{3}\geq0$, one also has $\langle(x^{-})^{3}\eta,\eta\rangle\geq0$; hence this inner product is zero for every $\eta$, so $(x^{-})^{3}=0$, $x^{-}=0$, and finally $x=x^{+}\geq0$.

\ResultHeading{1.6.8}{1.6.8.}
Let $A$ be a $C^{*}$-algebra, and let $a,b,x\in A$. If $a\leq b$, then
\[
x^{*}ax\leq x^{*}bx.
\]
Indeed, there is a $y\in A$ such that $b-a=y^{*}y$, whence
\[
x^{*}bx-x^{*}ax=x^{*}(y^{*}y)x=(yx)^{*}(yx)\geq0.
\]

% Source PDF physical page 024
Suppose that $A$ has an identity element, and let $b\in A$ satisfy $b\geq1$ (so that $b$ is invertible). Applying what precedes with $x=b^{-1/2}$, one obtains $1\geq b^{-1}$. More generally, let $a$ and $b$ be invertible elements of $A^{+}$ such that $0\leq a\leq b$. Then $a^{-1}\geq b^{-1}$; indeed, by what precedes, successively,
\[
1\leq a^{-1/2}ba^{-1/2},
\qquad
1\geq a^{1/2}b^{-1}a^{1/2},
\qquad
a^{-1}\geq b^{-1}.
\]

\ResultHeading{1.6.9}{1.6.9.}
Let $A$ be a $C^{*}$-algebra, and let $x,y\in A^{+}$ satisfy $y\leq x$. Then
\[
\|y\|\leq\|x\|.
\]
Indeed, we may suppose that $A$ has an identity element. By \XRef{1.4.1}{1.4.1}, for example, $x\leq\|x\|1$. Hence
\[
0\leq y\leq\|x\|1,
\]
and therefore $\|y\|\leq\|x\|$, again by \XRef{1.4.1}{1.4.1}.

\ResultHeading{1.6.10}{1.6.10.}
Let $A$ and $B$ be two $C^{*}$-algebras, and let $\varphi$ be a morphism from $A$ into $B$. Clearly,
\[
\varphi(A^{+})\subset\varphi(A)\cap B^{+}.
\]
Conversely, let $y\in\varphi(A)\cap B^{+}$; there is an $x\in A$ such that $y=\varphi(x)$, and
\[
y=(y^{*}y)^{1/2}=\varphi\bigl((x^{*}x)^{1/2}\bigr),
\]
so $y\in\varphi(A^{+})$. Therefore
\[
\varphi(A^{+})=\varphi(A)\cap B^{+}.
\]

Bibliography: \cite{ref36,ref80,ref108,ref126,ref132,ref142}.

\BookSubsection[Approximate Identities in C*-Algebras]{1.7}{Approximate Identities in $C^{*}$-Algebras}

\ResultHeading{1.7.1}{1.7.1.}
Let $A$ be a $C^{*}$-algebra. We shall say that an approximate identity $(u_\lambda)$ of $A$ (\XRef{B.29}{B~29}) is \emph{increasing and directed} if $u_\lambda\geq0$ for every $\lambda$, and if $\lambda\leq\mu$ implies $u_\lambda\leq u_\mu$.

\ResultHeading{1.7.2}{1.7.2. Proposition.}
\begin{itshape}
Let $A$ be a $C^{*}$-algebra, and let $m$ be a dense two-sided ideal of $A$. There exists an increasing directed approximate identity of $A$ consisting of elements of $m$. If $A$ is separable, this approximate identity may moreover be indexed by $\{1,2,\ldots\}$.
\end{itshape}

\Proof
Let $\widetilde A$ be the $C^{*}$-algebra obtained from $A$ by adjoining an identity element. Let $\Lambda$ be the set of finite subsets of $m$, ordered by inclusion. For
\[
\lambda=\{x_1,\ldots,x_n\}\in\Lambda,
\]
put
\[
v_\lambda=x_1x_1^{*}+\cdots+x_nx_n^{*}\in m,
\qquad
u_\lambda=v_\lambda\left(\frac1n+v_\lambda\right)^{-1}
\]
(the element $u_\lambda$ is calculated in $\widetilde A$, but in fact $u_\lambda\in m$). Since the real-valued function
\[
t\longmapsto t\left(\frac1n+t\right)^{-1}
\]
lies between $0$ and $1$ for $t\geq0$, one has
\[
0\leq u_\lambda\leq1.
\]
Moreover,
\[
\sum_{i=1}^{n}[(u_\lambda-1)x_i][(u_\lambda-1)x_i]^{*}
=(u_\lambda-1)v_\lambda(u_\lambda-1)
=\frac1{n^2}v_\lambda\left(\frac1n+v_\lambda\right)^{-2}.
\]
% Source PDF physical page 025
The real-valued function
\[
t\longmapsto t\left(\frac1n+t\right)^{-2}
\]
is at most $n/4$. Therefore
\[
\sum_{i=1}^{n}[(u_\lambda-1)x_i][(u_\lambda-1)x_i]^{*}
\leq\frac1{4n}.
\]
For $i=1,2,\ldots,n$, it follows that
\[
[(u_\lambda-1)x_i][(u_\lambda-1)x_i]^{*}
\leq\frac1{4n},
\]
and hence, by \XRef{1.6.9}{1.6.9},
\[
\|(u_\lambda-1)x_i\|^{2}\leq\frac1{4n}.
\]
Thus $\|(u_\lambda-1)x\|\to0$ for every $x\in m$, and consequently for every $x\in A$, since $\overline m=A$ and $\|u_\lambda\|\leq1$; moreover,
\[
\|xu_\lambda-x\|=\|u_\lambda x^{*}-x^{*}\|\longrightarrow0.
\]
Thus $(u_\lambda)$ is an approximate identity.

Let $\lambda,\mu\in\Lambda$ with $\lambda\leq\mu$. Write
\[
\lambda=\{x_1,\ldots,x_n\},
\qquad
\mu=\{x_1,\ldots,x_p\},
\qquad p\geq n;
\]
then $v_\lambda\leq v_\mu$. Hence, by \XRef{1.6.8}{1.6.8},
\[
\left(\frac1n+v_\lambda\right)^{-1}
\geq
\left(\frac1n+v_\mu\right)^{-1}.
\]
For every real number $t\geq0$,
\[
\frac1n\left(\frac1n+t\right)^{-1}
\geq
\frac1p\left(\frac1p+t\right)^{-1};
\]
therefore
\[
\frac1n\left(\frac1n+v_\mu\right)^{-1}
\geq
\frac1p\left(\frac1p+v_\mu\right)^{-1}.
\]
Altogether this gives
\[
1-\frac1n\left(\frac1n+v_\lambda\right)^{-1}
\leq
1-\frac1n\left(\frac1n+v_\mu\right)^{-1}
\leq
1-\frac1p\left(\frac1p+v_\mu\right)^{-1},
\]
that is, $u_\lambda\leq u_\mu$. The approximate identity $(u_\lambda)$ is increasing and directed.

Suppose that $A$ is separable. There is a sequence $(y_1,y_2,\ldots)$ dense in $m$. Put
\[
u_n=u_{\{y_1,y_2,\ldots,y_n\}}.
\]
The preceding argument shows that, for every $i$, $\|u_ny_i-y_i\|\to0$ as $n\to+\infty$. Since $\|u_n\|\leq1$, it follows that $u_nx\to x$ for every $x\in A$, and the proof is completed as above.

\ResultHeading{1.7.3}{1.7.3.}
Let $A$ be a $C^{*}$-algebra, and let $I$ be a right ideal of $A$. There is a family $(u_\lambda)$ in $I\cap A^{+}$, indexed by a directed set, such that:
\begin{enumerate}[label=\textup{\arabic*\textsuperscript{o}}]
\item $\|u_\lambda\|\leq1$;
\item $\lambda\leq\mu$ implies $u_\lambda\leq u_\mu$;
\item for every $x\in\overline I$, one has $\|u_\lambda x-x\|\to0$.
\end{enumerate}
The proof of \XRef{1.7.2}{1.7.2} applies word for word.

Bibliography: \cite{ref16,ref134}.

\BookSubsection[Quotient of a C*-Algebra]{1.8}{Quotient of a $C^{*}$-Algebra}

\ResultHeading{1.8.1}{1.8.1. Proposition.}
\begin{itshape}
Let $A$ be a $C^{*}$-algebra, let $B$ be an involutive normed algebra, and let $\varphi$ be an injective morphism from $A$ into $B$. Then
\[
\|\varphi(x)\|\geq\|x\|
\]
for every $x\in A$.
\end{itshape}

% Source PDF physical page 026
\Proof
Let $x\in A$. If we show that
\[
\|\varphi(x^{*}x)\|\geq\|x^{*}x\|,
\]
then
\[
\|x\|^{2}
=\|x^{*}x\|
\leq\|\varphi(x^{*}x)\|
=\|\varphi(x)^{*}\varphi(x)\|
\leq\|\varphi(x)\|^{2},
\]
which proves the proposition. We may therefore suppose from now on that $x$ is Hermitian. Replacing $\varphi$ by its restriction to the $C^{*}$-subalgebra of $A$ generated by $x$, we may suppose that $A$ is commutative. Replacing $B$ by $\varphi(A)$, we may suppose that $B$ is commutative, and then replace $B$ by its completion. We may also adjoin identity elements to $A$ and $B$. In short, we are reduced to the case where $A$ and $B$ are commutative, complete, and unital.

Let $S$ and $T$ be the spectra of $A$ and $B$, which are compact spaces. For every $\chi\in T$, the composite $\chi\circ\varphi$ is a character of $A$, that is, an element of $S$, which we denote by $\varphi'(\chi)$. If $x\in A$, then
\[
\varphi'(\chi)(x)=\chi(\varphi(x))
\]
is a continuous function of $\chi$, so $\varphi'$ is a continuous map from $T$ into $S$, and $\varphi'(T)$ is a compact subset of $S$. If $\varphi'(T)\neq S$, there are continuous complex-valued functions $f,g$ on $S$ such that
\[
fg=0,\qquad f\neq0,\qquad g=1\text{ on }\varphi'(T).
\]
By \XRef{1.4.1}{1.4.1}, $f$ and $g$ are the Gelfand transforms of two elements $x,y\in A$ such that
\[
xy=0,\qquad x\neq0,\qquad \chi(\varphi(y))=1
\quad\text{for every }\chi\in T.
\]
Then $\varphi(y)$ is invertible in $B$ by \XRef{B.3}{B~3}, while $\varphi(x)\varphi(y)=0$ and $\varphi(x)\neq0$, a contradiction. Hence $\varphi'(T)=S$. It follows that, for every $x\in A$,
\begin{align*}
\|x\|
&=\sup_{\chi\in S}|\chi(x)|
 =\sup_{\chi\in T}|\varphi'(\chi)(x)|\\
&=\sup_{\chi\in T}|\chi(\varphi(x))|
 \leq\|\varphi(x)\|.
\end{align*}

\ResultHeading{1.8.2}{1.8.2. Proposition.}
\begin{itshape}
Let $A$ be a $C^{*}$-algebra, and let $I$ be a closed two-sided ideal of $A$. Then $I$ is self-adjoint, and $A/I$, equipped with its natural involutive-algebra structure and the quotient norm, is a $C^{*}$-algebra.
\end{itshape}

\Proof
Let $(u_\lambda)$ be a family with the properties of \XRef{1.7.3}{1.7.3}. If $x\in I$, then
\[
\|x^{*}u_\lambda-x^{*}\|
=\|u_\lambda x-x\|\longrightarrow0,
\]
and $x^{*}u_\lambda\in I$, so $x^{*}\in\overline I=I$; hence $I=I^{*}$.

It is known that $A/I$ is an involutive algebra satisfying the axioms of a Banach algebra. Denote the canonical map from $A$ onto $A/I$ by $x\mapsto\dot x$. To prove that $A/I$ is a $C^{*}$-algebra, it is enough, by \XRef{1.3.4}{1.3.4}, to prove that
\[
\|\dot x\|^{2}\leq\|\dot x\dot x^{*}\|.
\]
One has
\begin{equation}
\|\dot x\|=\lim_\lambda\|x-u_\lambda x\|.
\tag{1}\label{eq:1.8.2-1}
\end{equation}
Indeed, if $y\in I$, then $u_\lambda y\to y$, and therefore
\begin{align*}
\lim_\lambda\|x-u_\lambda x\|
&=\lim_\lambda\|x-u_\lambda x+y-u_\lambda y\|\\
&=\lim_\lambda\|(1-u_\lambda)(x+y)\|
\leq\|x+y\|,
\end{align*}
% Source PDF physical page 027
where the calculation is made in $\widetilde A$. Thus
\[
\|\dot x\|
\leq\lim_\lambda\|x-u_\lambda x\|
\leq\inf_{y\in I}\|x+y\|
=\|\dot x\|,
\]
which proves \eqref{eq:1.8.2-1}. Now, for every $z\in I$,
\begin{align*}
\|\dot x\|^{2}
&=\lim_\lambda\|x-u_\lambda x\|^{2}\\
&=\lim_\lambda\|(x-u_\lambda x)(x-u_\lambda x)^{*}\|\\
&=\lim_\lambda\|xx^{*}-u_\lambda xx^{*}-xx^{*}u_\lambda
  +u_\lambda xx^{*}u_\lambda\|\\
&=\lim_\lambda\|xx^{*}+z-u_\lambda z-u_\lambda xx^{*}
  -xx^{*}u_\lambda-zu_\lambda
  +u_\lambda zu_\lambda+u_\lambda xx^{*}u_\lambda\|\\
&=\lim_\lambda\|(1-u_\lambda)(xx^{*}+z)(1-u_\lambda)\|\\
&\leq\|xx^{*}+z\|.
\end{align*}
Therefore
\[
\|\dot x\|^{2}\leq\|\dot x\dot x^{*}\|.
\]

\ResultHeading{1.8.3}{1.8.3. Corollary.}
\begin{itshape}
Let $A$ and $B$ be $C^{*}$-algebras, let $\varphi$ be a morphism from $A$ into $B$, and let $I$ be the kernel of $\varphi$. Consider the canonical factorization
\[
A\longrightarrow A/I\xrightarrow{\ \psi\ }\varphi(A)\longrightarrow B.
\]
Then $I$ is closed in $A$, $\varphi(A)$ is closed in $B$, and $\psi$ is an isometric isomorphism of the $C^{*}$-algebra $A/I$ onto the $C^{*}$-algebra $\varphi(A)$.
\end{itshape}

\Proof
Since $\varphi$ is continuous by \XRef{1.3.7}{1.3.7}, $I$ is closed; and $A/I$ is a $C^{*}$-algebra by \XRef{1.8.2}{1.8.2}. The morphism $A/I\to B$ induced by $\varphi$ on passing to the quotient is injective, and hence isometric by \XRef{1.3.7}{1.3.7} and \XRef{1.8.1}{1.8.1}. Thus $\varphi(A)$ is complete and consequently closed in $B$.

\ResultHeading{1.8.4}{1.8.4. Corollary.}
\begin{itshape}
Let $A$ be a $C^{*}$-algebra, let $B$ be a $C^{*}$-subalgebra of $A$, and let $I$ be a closed two-sided ideal of $A$. Then $B+I$ is a $C^{*}$-subalgebra of $A$, and the $C^{*}$-algebras
\[
(B+I)/I
\qquad\text{and}\qquad
B/(B\cap I)
\]
are canonically isomorphic.
\end{itshape}

\Proof
Let $\varphi$ be the canonical morphism $A\to A/I$, and let $\psi$ be the restriction of $\varphi$ to $B$. Then
\[
\psi(B)=(B+I)/I
\]
is closed in $A/I$ by \XRef{1.8.3}{1.8.3}, so $B+I$ is closed in $A$. Consider the canonical factorization of $\psi$:
\[
B\longrightarrow B/(B\cap I)\longrightarrow\psi(B)\longrightarrow A/I.
\]
By \XRef{1.8.3}{1.8.3}, the morphism
\[
B/(B\cap I)\longrightarrow\psi(B)=(B+I)/I
\]
is an isomorphism of $C^{*}$-algebras.

\ResultHeading{1.8.5}{1.8.5. Proposition.}
\begin{itshape}
Let $A$ be a $C^{*}$-algebra, let $I$ be a closed two-sided ideal of $A$, and let $J$ be a closed two-sided ideal of $I$. Then $J$ is a closed two-sided ideal of $A$.
\end{itshape}

\Proof
By \XRef{1.5.8}{1.5.8}, $J=J^{3}$, and hence
\[
AJA=AJ^{3}A\subset IJI\subset J.
\]

Bibliography: \cite{ref41,ref74,ref108,ref126,ref135}. The proofs of \XRef{1.8.2}{1.8.2} and \XRef{1.7.3}{1.7.3} were communicated to me orally by F. Combes.

\BookSubsection{1.9}{Supplements}

\ResultHeading{1.9.1}{1.9.1.}
Let $A$ be a $C^{*}$-algebra. Retain the norm and all the operations of $A$ except scalar multiplication $(\lambda,x)\mapsto\lambda x$, which is replaced by $(\lambda,x)\mapsto\overline\lambda x$. The resulting $C^{*}$-algebra $\overline A$ is called the conjugate of $A$.

% Source PDF physical page 028
\ResultHeading{1.9.2}{1.9.2.}
Let $A$ be a $C^{*}$-algebra. If $x\in A$ is normal, then
\[
\|x\|=\rho(x).
\]
(Use \XRef{1.4.1}{1.4.1}.) \cite{ref126}.

\ResultHeading{1.9.3}{1.9.3.}
Let $A$ be a $C^{*}$-algebra.
\begin{enumerate}[label=\textup{\alph*.}]
\item If every maximal commutative $C^{*}$-subalgebra of $A$ has an identity element, then $A$ has an identity element.
\item If every maximal commutative $C^{*}$-subalgebra of $A$ is finite-dimensional, then $A$ is finite-dimensional. \cite{ref116}
\end{enumerate}

\ResultHeading{1.9.4}{1.9.4.}
Let $A$ be a $C^{*}$-algebra. If the conditions $x\in A^{+}$, $y\in A^{+}$, and $x\geq y$ imply $x^{2}\geq y^{2}$, then $A$ is commutative. \cite{ref117}

\ResultHeading{1.9.5}{* 1.9.5.}
Let $A$ be a unital Banach algebra equipped with an involution such that
\[
\|x^{*}x\|=\|x^{*}\|\,\|x\|
\]
for every $x\in A$. Then $A$ is a $C^{*}$-algebra. \cite{ref49,ref120}

\ResultHeading{1.9.6}{1.9.6. Problem.}
Does a unital $C^{*}$-algebra $A$ of dimension $>1$, whose only closed two-sided ideals are $0$ and $A$, have projections distinct from $0$ and $1$? \cite{ref79}

\ResultHeading{1.9.7}{1.9.7.}
Let $A$ be a unital $C^{*}$-algebra, and let $x\in A$ be not left-invertible. Then $x^{*}x$ is not invertible in the $C^{*}$-subalgebra $B$ generated by $1$ and $x^{*}x$; consequently there exist $y_1,y_2,\ldots\in B$ such that
\[
\|y_n\|=1,
\qquad
\|y_nx^{*}x\|\longrightarrow0
\]
(use \XRef{1.4.1}{1.4.1}). Thus, if $x\in A$ is not invertible, then $x$ is a topological zero divisor. \cite{ref126}

\ResultHeading{1.9.8}{1.9.8.}
Let $A$ be a unital $C^{*}$-algebra, let $N$ be the set of normal elements of $A$, let $x\in N$, and let $V$ be a neighborhood of $0$ in $\mathbb C$. There is a neighborhood $U$ of $x$ in $N$ such that, for every $y\in U$,
\[
\operatorname{Sp}y\subset\operatorname{Sp}x+V,
\qquad
\operatorname{Sp}x\subset\operatorname{Sp}y+V.
\]
\cite{ref126}

\ResultHeading{1.9.9}{* 1.9.9.}
Let $A$ be a unital $C^{*}$-algebra, let $H$ be a Hilbert space, and let $\varphi$ be a linear map from $A$ into $\mathcal L(H)$ such that
\[
\varphi(A^{+})\subset\mathcal L(H)^{+}.
\]
Then, for every Hermitian $x\in A$,
\[
\varphi(x^{2})\geq\varphi(x)^{2}.
\]
\cite{ref66}

\ResultHeading{1.9.10}{* 1.9.10.}
Let $A$ and $B$ be two unital $C^{*}$-algebras.
\begin{enumerate}[label=\textup{\alph*.}]
\item Let $\rho:A\to B$ be a bijective linear map such that
\[
\rho(x^{*})=\rho(x)^{*},
\qquad
\rho(x^{n})=\rho(x)^{n}
\]
for Hermitian $x$ and every integer $n>0$. Then $\rho(A^{+})=B^{+}$, $\rho$ is isometric, $\rho$ carries commuting pairs to commuting pairs, and $\rho(1)=1$.

\item Let $\tau:A\to B$ be a bijective linear isometry. Then $\tau$ is the product of a map $\rho$ having the properties in (a) and left multiplication by the element $\tau(1)$ of $B$.

\item Let $\tau:A\to B$ be a linear isometry such that $\tau(1)=1$. Then
\[
\tau(x^{*})=\tau(x)^{*}
\]
for every $x\in A$. \cite{ref3,ref65}
\end{enumerate}

\ResultHeading{1.9.11}{* 1.9.11.}
Let $A$ be a $C^{*}$-algebra.
\begin{enumerate}[label=\textup{\alph*.}]
\item Every derivation of $A$ is continuous.

\item Let $D$ be a derivation of $A$, and let $x$ be a normal element of $A$. If
\[
x(Dx)=(Dx)x,
\]
then $Dx=0$. In particular, every derivation of a
% Source PDF physical page 029
commutative $C^{*}$-algebra is zero. It follows that a closed commutative two-sided ideal in a $C^{*}$-algebra is central.

\item Let $D$ be a derivation of $A$, and let $x$ be a normal element of $A$ such that $Dx=0$. Then $Dx^{*}=0$. \cite{ref79,ref129,ref130}

\item Let $D$ be a derivation of $A$, and let $I$ be a closed two-sided ideal of $A$. Then
\[
D(I)\subset I.
\]
(Use \XRef{1.5.8}{1.5.8}.)

\item Let $H$ be an infinite-dimensional Hilbert space, let $A$ be the $C^{*}$-algebra of compact operators on $H$, and let $x\in\mathcal L(H)$ with $x\notin A+\mathbb C1$. The derivation
\[
y\longmapsto xy-yx
\]
of $A$ is not inner.
\end{enumerate}

\ResultHeading{1.9.12}{1.9.12.}
Let $A$ be a $C^{*}$-algebra.
\begin{enumerate}[label=\textup{\alph*.}]
\item Let $I,J$ be two closed two-sided ideals of $A$. The product ideal $IJ$ is equal to $I\cap J$. (If $x\in(I\cap J)^{+}$, then $x^{1/2}\in(I\cap J)^{+}$.) \cite{ref21}

One has
\[
(I+J)^{+}=I^{+}+J^{+}.
\]
\cite{ref293,ref365}

\item If $I\cap J=0$, the canonical map
\[
I+J\longrightarrow I\times J
\]
is an isomorphism of $C^{*}$-algebras.

\item Let $(I_\lambda)$ be a family of closed two-sided ideals of $A$ with intersection $I$. Let $\omega_\lambda$ (respectively, $\omega$) be the canonical morphism of $A$ onto $A/I_\lambda$ (respectively, $A/I$). For every $x\in A$,
\[
\|\omega(x)\|=\sup_\lambda\|\omega_\lambda(x)\|.
\]
\cite{ref29}
\end{enumerate}

\ResultHeading{1.9.13}{1.9.13.}
Let $A$ be a $C^{*}$-algebra. A \emph{$C^{*}$-seminorm} on $A$ is a seminorm $N$ such that
\[
N(x)\leq\|x\|,
\qquad
N(xy)\leq N(x)N(y),
\qquad
N(x^{*}x)=N(x)^{2}
\]
for all $x,y\in A$. The set $\mathcal N(A)$ of $C^{*}$-seminorms on $A$ is compact for the topology of pointwise convergence. If $I$ is a closed two-sided ideal of $A$ and $x\in A$, let $N_I(x)$ be the norm of the canonical image of $x$ in $A/I$. Then
\[
I\longmapsto N_I
\]
is a bijection from the set $\mathcal J(A)$ of closed two-sided ideals of $A$ onto $\mathcal N(A)$. If $I,J\in\mathcal J(A)$, then
\[
N_{I\cap J}=\sup(N_I,N_J)
\]
(use \XRef{1.9.12}{1.9.12}(b)). A $C^{*}$-seminorm $N$ is called \emph{extremal} if it cannot be the supremum of two $C^{*}$-seminorms without being equal to one of them. A nonzero $C^{*}$-seminorm $N_I$ is extremal if and only if $I$ is prime. [In an arbitrary algebra $R$, a two-sided ideal $J$ is called prime if $J\neq R$ and the relation $J'J''\subset J$, where $J',J''$ are two-sided ideals of $R$, implies $J'\subset J$ or $J''\subset J$.] If $A$ is separable, the set $E$ of extremal $C^{*}$-seminorms is a $G_\delta$ in $\mathcal N(A)$. [The set of pairs
\[
(N,N')\in\mathcal N(A)\times\mathcal N(A)
\]
such that $N\leq N'$ or $N'\leq N$ is compact. Its complement is a countable union of compact sets; its image in $\mathcal N(A)$ under the continuous map
\[
(N,N')\longmapsto\sup(N,N')
\]
is the complement of $E$.] \cite{ref20,ref29}

\ResultHeading{1.9.14}{1.9.14.}
Let $(A_i)_{i\in I}$ be a family of $C^{*}$-algebras. Let $A$ be the set of
\[
x=(x_i)\in\prod_{i\in I}A_i
\]
such that, for every $\varepsilon>0$, one has $\|x_i\|\leq\varepsilon$ except for finitely many indices $i$. With
\[
\|x\|=\sup_i\|x_i\|,
\]
$A$ becomes naturally a $C^{*}$-algebra, called the restricted product of the $A_i$.

\ResultHeading{1.9.15}{1.9.15.}
The radical $R$ of a $C^{*}$-algebra is zero. [If $x\in R$, then $1+\lambda x^{*}x$ is invertible in $\widetilde A$ for every $\lambda\in\mathbb C$, so $\operatorname{Sp}'(x^{*}x)=\{0\}$ and $x=0$.] \cite{ref133}
% ===== ASSEMBLED TRANSLATION CHUNK 02: chunk_030_059.tex =====
% Source PDF physical page 030
\BookSection{2}{Positive Functionals and Representations}

One of the most important notions in this book is that of a representation.
Given an involutive Banach algebra \(A\), it would be difficult to establish
directly the existence of representations of \(A\). We shall instead define a
correspondence between representations of \(A\) and positive functionals on
\(A\), and in particular between irreducible representations and pure positive
functionals. The classical tools of functional analysis (the Hahn--Banach and
Krein--Milman theorems) will then make it possible to prove the existence of
positive functionals, and even of pure positive functionals. This is the
essential idea of this section.

In the various statements, varying hypotheses are imposed on the involutive
algebras under consideration. On a first reading, one may suppose throughout
that they are \(C^{*}\)-algebras.

\BookSubsection{2.1}{Positive Functionals}

\ResultHeading{2.1.1}{2.1.1. Definition.}
\begin{itshape}
Let \(A\) be an involutive algebra. A linear functional \(f\) on \(A\) is said
to be \emph{positive} if \(f(x^{\ast}x)\geq 0\) for every \(x\in A\). If \(A\)
is an involutive normed algebra, a continuous positive linear functional \(f\)
on \(A\) such that \(\lVert f\rVert=1\) is called a \emph{state} of \(A\).
\end{itshape}

Let \(A\) be an involutive algebra, and let \(f\) and \(g\) be two linear
functionals on \(A\). We say that \(f\) \emph{majorizes} \(g\), and write
\(f\geq g\) or \(g\leq f\), if \(f-g\) is positive. This defines on the dual
of \(A\) a preorder relation compatible with its real vector-space structure.
If \(A\) is a \(C^{*}\)-algebra, this relation is in fact an order relation: for
if \(f\geq0\) and \(f\leq0\), then \(f\) vanishes on \(A^{+}\), hence on the
set of hermitian elements of \(A\) (\XRef{1.6.5}{1.6.5}), and consequently
\(f=0\).

Let \(\Omega\) be a locally compact space, and let \(A\) be the
\(C^{*}\)-algebra of continuous complex-valued functions on \(\Omega\) that
vanish at infinity. A continuous linear functional on \(A\) is nothing other
than a bounded measure \(\mu\) on \(\Omega\). To say that this functional is
positive means exactly that the measure \(\mu\) is positive.

\ResultHeading{2.1.2}{2.1.2.}
Let \(A\) be an involutive algebra and \(f\) a positive functional on \(A\).
For \(x,y\in A\), put
\[
 (x\mid y)=f(y^{\ast}x).
\]
This inner product is linear in \(x\), conjugate-linear in \(y\), and
\((x\mid x)\geq0\) for every \(x\). Thus \(A\) acquires the structure of a
pre-Hilbert space.

In particular, this proves that
\begin{align}
 \overline{f(y^{\ast}x)}&=f(x^{\ast}y)
   &&(x\in A,\ y\in A), \tag{1}\label{eq:2.1.2-1}\\
 |f(y^{\ast}x)|^{2}&\leq f(x^{\ast}x)f(y^{\ast}y)
   &&(x\in A,\ y\in A). \tag{2}\label{eq:2.1.2-2}
\end{align}

% Source PDF physical page 031
If \(A\) has an identity element, then, on setting \(y=1\) in
\eqref{eq:2.1.2-1} and \eqref{eq:2.1.2-2}, we obtain
\begin{align}
 f(x^{\ast})&=\overline{f(x)}, \tag{3}\label{eq:2.1.2-3}\\
 |f(x)|^{2}&\leq f(1)f(x^{\ast}x). \tag{4}\label{eq:2.1.2-4}
\end{align}

Let \(H\) be the separated pre-Hilbert space canonically obtained from the
pre-Hilbert space \(A\): thus \(H=A/N\), where \(N\) is the set of all
\(x\in A\) such that \(f(x^{\ast}x)=0\). By \eqref{eq:2.1.2-2}, \(N\) is
also the set of all \(x\in A\) such that \(f(yx)=0\) for every \(y\in A\).
Hence \(N\) is a left ideal of \(A\).

[If one put \((x\mid y)=f(xy^{\ast})\), one would obtain another pre-Hilbert
space structure on \(A\), with analogous properties.]

\ResultHeading{2.1.3}{2.1.3. Lemma.}
\begin{itshape}
Let \(A\) be a Banach algebra with identity, let \(x'\in A\) satisfy
\(\lVert x'\rVert\leq1\), and set \(x=1+x'\). The series
\begin{multline*}
 1+\frac12x'+\frac1{2!}\frac12\left(\frac12-1\right)x'^{\,2}
 +\cdots \\
 {}+\frac1{n!}\frac12\left(\frac12-1\right)\cdots
 \left(\frac12-n+1\right)x'^{\,n}+\cdots
\end{multline*}
converges to an element \(y\in A\) such that \(y^{2}=x\). If \(A\) is
equipped with an isometric involution and \(x\) is hermitian, then \(y\) is
hermitian.
\end{itshape}

The series
\begin{multline*}
 1+\frac12\lVert x'\rVert
 +\frac1{2!}\left|\frac12\left(\frac12-1\right)\right|
   \lVert x'\rVert^{2}+\cdots \\
 {}+\frac1{n!}\left|\frac12\left(\frac12-1\right)\cdots
 \left(\frac12-n+1\right)\right|\lVert x'\rVert^{n}+\cdots
\end{multline*}
converges, which proves the existence of \(y\). If \(y^{2}\) is calculated,
one obtains a power series in \(x'\) whose coefficients are known from the
classical case \(A=\mathbb C\); hence \(y^{2}=1+x'=x\). If \(A\) is equipped
with an isometric involution and \(x\) is hermitian, \(y\) appears as a limit
of hermitian elements and is therefore hermitian, since the involution is
continuous.

\ResultHeading{2.1.4}{2.1.4. Proposition.}
\begin{itshape}
Let \(A\) be an involutive Banach algebra having an identity \(1\) such that
\(\lVert1\rVert=1\). If \(f\) is a positive linear functional on \(A\), then
\(f\) is continuous and \(\lVert f\rVert=f(1)\).
\end{itshape}

If \(x\in A\) is hermitian and \(\lVert x\rVert\leq1\), then \(1-x\) can be
written in the form \(y^{\ast}y\) (\XRef{2.1.3}{Lemma 2.1.3}); hence
\(f(1-x)\geq0\) and \(f(x)\leq f(1)\). If \(x'\in A\) and
\(\lVert x'\rVert\leq1\), then \(\lVert x'^{\ast}x'\rVert\leq1\), and
therefore, by \XRef{2.1.2}{2.1.2}, formula \eqref{eq:2.1.2-4},
\[
 |f(x')|^{2}\leq f(1)f(x'^{\ast}x')\leq f(1)^{2}.
\]
This proves that \(f\) is continuous and that \(\lVert f\rVert\leq f(1)\).
It is clear that
\[
 f(1)\leq\lVert f\rVert,
 \qquad\text{and hence}\qquad
 \lVert f\rVert=f(1).
\]

\ResultHeading{2.1.5}{2.1.5. Proposition.}
\begin{itshape}
Let \(A\) be an involutive Banach algebra admitting an approximate identity
(\XRef{B.29}{B 29}), let \(\widetilde A\) be the involutive algebra obtained
from \(A\)
% Source PDF physical page 032
by adjoining an identity element, and let \(f\) be a continuous positive
linear functional on \(A\).
\begin{enumerate}[label=\textup{(\roman*)}]
 \item \(f(x^{\ast})=\overline{f(x)}\) and
 \(\lvert f(x)\rvert^{2}\leq\lVert f\rVert f(x^{\ast}x)\) for every
 \(x\in A\).
 \item \(\lvert f(y^{\ast}xy)\rvert\leq\lVert x\rVert f(y^{\ast}y)\) for
 all \(x,y\in A\).
 \item \(\displaystyle \lVert f\rVert=
 \sup_{x\in A,\,\lVert x\rVert\leq1}f(x^{\ast}x)\).
 \item If \((x_i)_{i\in I}\) is a family of elements of \(A\), indexed by a
 directed set, such that \(\lVert x_i\rVert\leq1\) and
 \(f(x_i)\to\lVert f\rVert\), then
 \(f(x_i^{\ast}x_i)\to\lVert f\rVert\).
 \item If \((u_i)_{i\in I}\) is an approximate identity of \(A\), then
 \(f(u_i)\to\lVert f\rVert\) and
 \(f(u_i^{\ast}u_i)\to\lVert f\rVert\).
 \item \(f\) extends uniquely to a positive functional \(\widetilde f\) on
 \(\widetilde A\) such that \(\widetilde f(1)=\lVert f\rVert\). Every
 positive functional on \(\widetilde A\) extending \(f\) majorizes
 \(\widetilde f\).
 \item With the notation of \textup{(iv)}, \(x_i\to1\) for the pre-Hilbert
 space structure defined by \(\widetilde f\) on \(\widetilde A\);
 consequently, \(A\) is dense in \(\widetilde A\) for this structure.
\end{enumerate}
\end{itshape}

For every \(x\in A\), we have
\begin{align*}
 \overline{f(x^{\ast})}
 &=\lim_i\overline{f(x^{\ast}u_i)}
  =\lim_i\overline{f((u_i^{\ast}x)^{\ast})}
  =\lim_i f(u_i^{\ast}x)=f(x),\\
 |f(x)|^{2}
 &=\lim_i|f(xu_i)|^{2}
 \leq f(x^{\ast}x)\lim_i f(u_i^{\ast}u_i)
 \leq\lVert f\rVert f(x^{\ast}x).
\end{align*}
This proves (i), and the implications (i) \(\Rightarrow\) (iv)
\(\Rightarrow\) (iii) are clear. Let us prove (vi). The uniqueness of
\(\widetilde f\) is immediate. We prove its existence. For every element
\((\lambda,x)=\lambda+x\) of \(\widetilde A\)
\((\lambda\in\mathbb C,\ x\in A)\), put
\[
 \widetilde f(\lambda+x)=\lambda\lVert f\rVert+f(x).
\]
Then \(\widetilde f\) is a linear functional on \(\widetilde A\) extending
\(f\), and, by (i),
\begin{align*}
 \widetilde f((\lambda+x)^{\ast}(\lambda+x))
 &=f(x^{\ast}x+\overline\lambda x+\lambda x^{\ast})
   +|\lambda|^{2}\lVert f\rVert\\
 &=f(x^{\ast}x)+2\operatorname{Re}(\overline\lambda f(x))
   +|\lambda|^{2}\lVert f\rVert\\
 &\geq f(x^{\ast}x)-2|\lambda|\lVert f\rVert^{1/2}f(x^{\ast}x)^{1/2}
   +|\lambda|^{2}\lVert f\rVert\\
 &=\left(f(x^{\ast}x)^{1/2}
   -|\lambda|\lVert f\rVert^{1/2}\right)^{2}\geq0.
\end{align*}
Thus \(\widetilde f\) is positive. Clearly,
\(\widetilde f(1)=\lVert f\rVert\). Let \(g\) be a positive functional on
\(\widetilde A\) extending \(f\). One may give \(\widetilde A\) the
structure of an involutive Banach algebra extending that of \(A\) and
satisfying \(\lVert1\rVert=1\). By \XRef{2.1.4}{2.1.4},
\(g(1)=\lVert g\rVert\geq\lVert f\rVert\), whence
\begin{align*}
 g((\lambda+x)^{\ast}(\lambda+x))
 &=f(x^{\ast}x+\overline\lambda x+\lambda x^{\ast})
   +|\lambda|^{2}g(1)\\
 &\geq f(x^{\ast}x+\overline\lambda x+\lambda x^{\ast})
   +|\lambda|^{2}\lVert f\rVert\\
 &=\widetilde f((\lambda+x)^{\ast}(\lambda+x)).
\end{align*}
Hence \(g\geq\widetilde f\), and (vi) is proved. If \(y\in A\), the
functional \(x\mapsto\widetilde f(y^{\ast}xy)\) on \(\widetilde A\) is
positive, because
\(\widetilde f(y^{\ast}(x^{\ast}x)y)
=\widetilde f((xy)^{\ast}(xy))\geq0\). By \XRef{2.1.4}{2.1.4}, its norm
is \(f(y^{\ast}y)\), which proves (ii).

% Source PDF physical page 033
With the notation of (iv),
\begin{align*}
 \widetilde f((x_i-1)^{\ast}(x_i-1))
 &=f(x_i^{\ast}x_i)-\overline{f(x_i)}-f(x_i)+\lVert f\rVert\\
 &\longrightarrow\lVert f\rVert-\lVert f\rVert-\lVert f\rVert
   +\lVert f\rVert=0,
\end{align*}
which proves (vii).

Denote by \(\lVert\cdot\rVert'\) the seminorm of the pre-Hilbert space
\(\widetilde A\). Let \(\varepsilon>0\). By (vii), there exists \(x\in A\)
such that \(\lVert x\rVert\leq1\) and
\(\lVert x-1\rVert'<\varepsilon\). Next there exists \(j_0\in I\) such
that \(j\geq j_0\) implies \(\lVert u_jx-x\rVert<\varepsilon\), whence
\[
 \lVert u_jx-x\rVert'
 =f((u_jx-x)^{\ast}(u_jx-x))^{1/2}
 \leq\lVert f\rVert^{1/2}\lVert u_jx-x\rVert
 \leq\varepsilon\lVert f\rVert^{1/2}.
\]
Finally, using (ii) for \(\widetilde f\), we have
\[
 \lVert u_jx-u_j\rVert'
 =\widetilde f((x-1)^{\ast}u_j^{\ast}u_j(x-1))^{1/2}
 \leq\lVert u_j^{\ast}u_j\rVert^{1/2}\lVert x-1\rVert'
 \leq\varepsilon.
\]
Thus \(j\geq j_0\) implies
\(\lVert u_j-1\rVert'\leq\varepsilon(2+\lVert f\rVert^{1/2})\). It follows
that \(f(u_j)=(u_j\mid1)\) tends to
\((1\mid1)=\widetilde f(1)=\lVert f\rVert\), and hence, by (iv), that
\(f(u_j^{\ast}u_j)\) tends to \(\lVert f\rVert\). This completes the proof.

\ResultHeading{2.1.6}{2.1.6.}
With the notation of \XRef{2.1.5}{2.1.5}, \(\widetilde f\) is called the
\emph{canonical extension} of \(f\) to \(\widetilde A\).

\ResultLabel{Corollary.}
\begin{itshape}
Let \(f\) and \(g\) be two continuous positive functionals on \(A\), and let
\(\widetilde f\) and \(\widetilde g\) be their canonical extensions to
\(\widetilde A\). Then
\[
 \lVert f+g\rVert=\lVert f\rVert+\lVert g\rVert,
 \qquad
 \widetilde{f+g}=\widetilde f+\widetilde g.
\]
\end{itshape}
The first formula follows from \XRef{2.1.5}{2.1.5(v)}, and the second follows
from the first.

Consequently, the set of states of \(A\) is a convex subset of the dual of
\(A\).

\ResultHeading{2.1.7}{2.1.7.}
Retain the notation of \XRef{2.1.5}{2.1.5}. Let \(Q\) be the set of
continuous positive functionals on \(A\), and let \(\widetilde Q\) be the
set of positive functionals \(g\) on \(\widetilde A\) such that
\(g(1)=\lVert g|_A\rVert\). Then \(f\mapsto\widetilde f\) is a bijection
from \(Q\) onto \(\widetilde Q\) that preserves order and addition. If a
positive functional on \(\widetilde A\) is majorized by an element of
\(\widetilde Q\), it belongs to \(\widetilde Q\). Indeed, suppose
\(g=g_1+g_2\), where \(g\in\widetilde Q\) and \(g_1,g_2\) are positive
functionals on \(\widetilde A\), and let \(f,f_1,f_2\) be the restrictions
of \(g,g_1,g_2\) to \(A\). We have
\[
 g_1+g_2=g=\widetilde f=\widetilde f_1+\widetilde f_2
\]
and
\[
 g_1(1)\geq\widetilde f_1(1),
 \qquad
 g_2(1)\geq\widetilde f_2(1),
\]
whence
\[
 g_1(1)=\widetilde f_1(1),
 \qquad
 g_2(1)=\widetilde f_2(1).
\]
Thus \(g_1,g_2\in\widetilde Q\).

% Source PDF physical page 034
\ResultHeading{2.1.8}{2.1.8.}
Let \(A\) be a \(C^{*}\)-algebra and \(f\) a positive functional on \(A\).
Then \(f\) is continuous. Indeed, let \((x_1,x_2,\ldots)\) be a sequence
of elements of \(A^{+}\) such that \(\lVert x_i\rVert\leq1\), and let us
show that the numbers \(f(x_i)\) are bounded. For every sequence
\((\lambda_1,\lambda_2,\ldots)\) of nonnegative numbers such that
\(\sum_{i=1}^{\infty}\lambda_i<+\infty\), the series
\(\sum_{i=1}^{\infty}\lambda_ix_i\) converges to an element \(x\in A\).
For every integer \(n\geq1\),
\[
 \sum_{i=n+1}^{\infty}\lambda_ix_i\geq0
\]
by \XRef{1.6.1}{1.6.1}, and hence
\[
 \sum_{i=1}^{n}\lambda_if(x_i)
 =f\left(\sum_{i=1}^{n}\lambda_ix_i\right)\leq f(x).
\]
Thus \(\sum_{i=1}^{\infty}\lambda_if(x_i)<+\infty\). Since the sequence
\((\lambda_i)\) is arbitrary, this proves that the \(f(x_i)\) are bounded.
It follows that
\[
 M=\sup_{x\in A^{+},\,\lVert x\rVert\leq1}f(x)<+\infty.
\]
If \(x\) is hermitian and has norm at most \(1\), then
\[
 |f(x)|\leq|f(x^{+})|+|f(x^{-})|\leq2M.
\]
If \(x\) is arbitrary and has norm at most \(1\), then
\[
 |f(x)|\leq
 \left|f\left(\frac12(x+x^{\ast})\right)\right|
 +\left|f\left(\frac1{2i}(x-x^{\ast})\right)\right|
 \leq4M.
\]
Therefore \(f\) is continuous.

In what follows, we shall constantly speak of positive functionals on
\(C^{*}\)-algebras. The reader should bear in mind that the continuity of
these functionals is automatic.

\ResultHeading{2.1.9}{2.1.9. Proposition.}
\begin{itshape}
Let \(A\) be a unital \(C^{*}\)-algebra and \(f\) a continuous linear
functional on \(A\). Then \(f\) is positive if and only if
\(\lVert f\rVert=f(1)\).
\end{itshape}

If \(f\geq0\), then \(\lVert f\rVert=f(1)\) by \XRef{2.1.4}{2.1.4}.
Conversely, suppose that \(\lVert f\rVert=f(1)\). Let \(x\in A^{+}\), and
suppose that \(f(x)\) is not nonnegative. We may at once reduce to the case
\(f(1)=1\). There exists a closed disk
\(\lvert z-z_0\rvert\leq\rho\) in \(\mathbb C\) that contains
\(\operatorname{Sp}x\) but does not contain \(f(x)\). The spectrum of the
normal element \(x-z_0\) is then contained in the disk
\(\lvert z\rvert\leq\rho\), so that \(\lVert x-z_0\rVert\leq\rho\).
Hence
\[
 |f(x)-z_0|=|f(x)-z_0f(1)|=|f(x-z_0)|
 \leq\lVert f\rVert\lVert x-z_0\rVert\leq\rho,
\]
which is a contradiction. (Another proof can be based on
\XRef{B.28}{B 28}.)

% Source PDF physical page 035
The word \emph{state} is borrowed from physics.

Bibliography: \cite{ref3,ref41,ref42,ref52,ref105,ref108,ref126,ref134}.
The result of \XRef{2.1.8}{2.1.8} has been generalized to involutive Banach
algebras admitting an approximate identity \cite{ref210}.

\BookSubsection{2.2}{Representations}

\ResultHeading{2.2.1}{2.2.1. Definition.}
\begin{itshape}
Let \(A\) be an involutive algebra and \(H\) a Hilbert space. A
\emph{representation} of \(A\) in \(H\) is a homomorphism of the involutive
algebra \(A\) into the involutive algebra \(\mathcal L(H)\).
\end{itshape}

In other words, a representation of \(A\) in \(H\) is a map \(\pi\) from
\(A\) into \(\mathcal L(H)\) such that
\begin{align*}
 \pi(x+y)&=\pi(x)+\pi(y), & \pi(\lambda x)&=\lambda\pi(x),\\
 \pi(xy)&=\pi(x)\pi(y), & \pi(x^{\ast})&=\pi(x)^{\ast}
\end{align*}
for \(x,y\in A\) and \(\lambda\in\mathbb C\).

The Hilbert-space dimension of \(H\) is called the dimension of \(\pi\) and
is denoted by \(\dim\pi\). The space \(H\) is called the space of \(\pi\)
and is denoted by \(H_{\pi}\).

Two representations \(\pi\) and \(\pi'\) of \(A\) in \(H\) and \(H'\) are
said to be \emph{equivalent}, and we write \(\pi\sim\pi'\), if there exists
a Hilbert-space isomorphism \(U:H\to H'\) that carries \(\pi(x)\) into
\(\pi'(x)\) for every \(x\in A\). This gives the notion of a class of
representations. (For convenience of language, we shall often identify
representations with classes of representations.)

\ResultHeading{2.2.2}{2.2.2.}
An \emph{intertwining operator} for \(\pi\) and \(\pi'\) is a continuous
linear operator \(T:H\to H'\) such that \(T\pi(x)=\pi'(x)T\) for every
\(x\in A\). The operator \(U\) in the preceding definition is an
intertwining operator. The set of intertwining operators for \(\pi\) and
\(\pi'\) is a vector space whose dimension is called the
\emph{intertwining number} of \(\pi\) and \(\pi'\). It is equal to the
intertwining number of \(\pi'\) and \(\pi\). Indeed, let \(T:H\to H'\) be
an intertwining operator for \(\pi\) and \(\pi'\). Then \(T^{\ast}:H'\to H\)
is an intertwining operator for \(\pi'\) and \(\pi\), since
\[
 T^{\ast}\pi'(x)=(\pi'(x^{\ast})T)^{\ast}
 =(T\pi(x^{\ast}))^{\ast}=\pi(x)T^{\ast}.
\]
It follows that
\[
 T^{\ast}T\pi(x)=T^{\ast}\pi'(x)T=\pi(x)T^{\ast}T.
\]
Thus \(\lvert T\rvert=(T^{\ast}T)^{1/2}\) commutes with \(\pi(A)\). Let
\(T=U\lvert T\rvert\) be the polar decomposition of \(T\). For every
\(x\in A\),
\begin{equation}
 U\pi(x)|T|=U|T|\pi(x)=T\pi(x)=\pi'(x)T
 =\pi'(x)U|T|. \tag{1}\label{eq:2.2.2-1}
\end{equation}
If \(\ker T=0\), then \(\lvert T\rvert(H)\) is dense in \(H\), and
\eqref{eq:2.2.2-1} implies
\begin{equation}
 U\pi(x)=\pi'(x)U. \tag{2}\label{eq:2.2.2-2}
\end{equation}
If, in addition, \(T(H)=H'\), or equivalently \(\ker T^{\ast}=0\), then
\(U\) is an isomorphism of \(H\) onto \(H'\), and
\eqref{eq:2.2.2-2} proves that \(\pi\) and \(\pi'\) are equivalent.

% Source PDF physical page 036
\ResultHeading{2.2.3}{2.2.3.}
Let \((\pi_i)_{i\in I}\) be a family of representations of \(A\) in Hilbert
spaces \(H_i\), and let \(H\) be the Hilbert direct sum of the \(H_i\). If,
for every \(x\in A\), the numbers \(\lVert\pi_i(x)\rVert\) are bounded
(this is the case, by \XRef{1.3.7}{1.3.7}, when \(A\) is an involutive
Banach algebra), one can form the continuous linear operator \(\pi(x)\) on
\(H\) whose restriction to each \(H_i\) is \(\pi_i(x)\). Then
\(x\mapsto\pi(x)\) is a representation of \(A\) in \(H\), called the
\emph{Hilbert direct sum} of the \(\pi_i\), and denoted by
\(\bigoplus_{i\in I}\pi_i\), or by
\(\pi_1\oplus\pi_2\oplus\cdots\oplus\pi_n\) for a family
\((\pi_1,\pi_2,\ldots,\pi_n)\) of representations.

If \((\pi_i)_{i\in I}\) is a family of representations of \(A\), all equal
to a representation \(\pi\), and if \(\operatorname{Card}I=c\), then
\(\bigoplus\pi_i\) is denoted by \(c\pi\). Any representation equivalent to
a representation of this form is called a \emph{multiple} of \(\pi\).

\ResultHeading{2.2.4}{2.2.4.}
Let \(\pi\) be a representation of \(A\) in \(H\). If a closed vector
subspace \(K\) of \(H\) is invariant under \(\pi(A)\), the restrictions of
the \(\pi(x)\) to \(K\) define a \emph{subrepresentation} \(\pi'\) of
\(A\) in \(K\), denoted by \(\pi_K\), or by \(\pi_E\) if \(E=P_K\). Then
\(H\ominus K\) is invariant under \(\pi(A)\), for if \(\xi\in K\) and
\(\eta\in H\ominus K\), then \(\pi(x^{\ast})\xi\in K\) for every
\(x\in A\), and hence
\[
 (\pi(x)\eta\mid\xi)=(\eta\mid\pi(x)^{\ast}\xi)
 =(\eta\mid\pi(x^{\ast})\xi)=0.
\]
Thus \(\pi(x)\eta\in H\ominus K\), and \(P_K\) commutes with \(\pi(A)\).
If \(\pi''\) is the subrepresentation of \(\pi\) defined by
\(H\ominus K\), then \(\pi\sim\pi'\oplus\pi''\).

Let \(\rho,\rho'\) be two representations of \(A\). If \(\rho'\) is
equivalent to a subrepresentation of \(\rho\), we say that \(\rho'\) is
\emph{contained in} \(\rho\), or that \(\rho\) \emph{contains} \(\rho'\),
and write \(\rho'\leq\rho\) or \(\rho\geq\rho'\).

\ResultHeading{2.2.5}{2.2.5.}
Let \(\pi\) be a representation of \(A\) in \(H\), and let \(\xi\in H\).
The closure of \(\pi(A)\xi\) is a closed vector subspace of \(H\) invariant
under \(\pi(A)\). If this subspace is \(H\), then \(\xi\) is said to be
\emph{cyclic} for \(\pi\).

\ResultHeading{2.2.6}{2.2.6. Proposition.}
\begin{itshape}
Let \(\pi\) be a representation of \(A\) in \(H\). Let \(K\) be the closed
vector subspace of \(H\) spanned by the vectors
\(\pi(x)\xi\) (\(x\in A,\ \xi\in H\)). Let \(K'\) be the closed vector
subspace of \(H\) formed by the vectors \(\xi\in H\) annihilated by all
\(\pi(x)\). Then \(K\) and \(K'\) are invariant under \(\pi(A)\) and
\(\pi(A)'\), are orthogonal, and have sum \(H\).
\end{itshape}

It is clear that \(K\) and \(K'\) are invariant under \(\pi(A)\). An
operator commuting with \(\pi(x)\) leaves invariant the kernel and range of
\(\pi(x)\), and hence \(K\) and \(K'\) are invariant under \(\pi(A)'\).
Since \(\pi(A)(K)\subset K\),
\[
 \pi(A)(H\ominus K)\subset H\ominus K;
\]
but \(\pi(A)(H)\subset K\), so \(\pi(A)(H\ominus K)=0\), and consequently
\(H\ominus K\subset K'\). If \(\xi\in K'\), \(x\in A\), and
\(\eta\in H\), then
\[
 (\xi\mid\pi(x)\eta)=(\pi(x^{\ast})\xi\mid\eta)=0,
\]
and therefore \(\xi\in H\ominus K\). This proves that \(K'=H\ominus K\).

% Source PDF physical page 037
The space \(K\) is called the \emph{essential subspace} of \(\pi\). The
representation \(\pi\) is said to be \emph{nondegenerate} if \(K=H\).
What precedes proves that every representation of \(A\) is, in a unique
way, the Hilbert direct sum of a zero representation and a nondegenerate
representation.

\ResultHeading{2.2.7}{2.2.7.}
It is clear that a Hilbert direct sum of nondegenerate representations is
nondegenerate, and that a representation admitting a cyclic vector is
nondegenerate. Conversely:

\ResultLabel{Proposition.}
\begin{itshape}
Every nondegenerate representation of \(A\) is a Hilbert direct sum of
representations admitting cyclic vectors.
\end{itshape}

Let \(\pi\) be a nondegenerate representation of \(A\) in \(H\), with
\(H\neq0\). It is enough to prove that \(\pi\) has a subrepresentation on a
nonzero space that admits a cyclic vector (for one then applies Zorn's
theorem). Let \(\xi\) be a nonzero element of \(H\). The closure \(K\) of
\(\pi(A)\xi\) is nonzero and invariant under \(\pi\). Put
\(L=H\ominus K\), and write \(\xi=\xi_1+\xi_2\), where
\(\xi_1\in K\) and \(\xi_2\in L\). If \(x\in A\), then
\(\pi(x)\xi_1\in K\), \(\pi(x)\xi_2\in L\), and
\(\pi(x)\xi_1+\pi(x)\xi_2=\pi(x)\xi\in K\); hence
\(\pi(x)\xi_2=0\), and consequently \(\pi(x)\xi=\pi(x)\xi_1\).
Thus \(K\) is the closure of \(\pi(A)\xi_1\).

\ResultHeading{2.2.8}{2.2.8.}
Let \(\pi\) be a representation of \(A\) in \(H\), let \(\overline H\) be
the conjugate Hilbert space of \(H\) (\(\overline H=H\), but on passing
from \(H\) to \(\overline H\), multiplication by
\(\lambda\in\mathbb C\) is replaced by multiplication by
\(\overline\lambda\), and inner products are replaced by their complex
conjugates), and let \(A^{\mathrm o}\) be the opposite involutive algebra of
\(A\). It is readily verified that the map
\(x\mapsto\overline{\pi(x^{\ast})}\) is a representation of \(A^{\mathrm o}\)
in \(\overline H\), denoted by \(\bar\pi^{\mathrm o}\).

\ResultHeading{2.2.9}{2.2.9.}
Let \(\widetilde A\) be the involutive algebra obtained from \(A\) by
adjoining an identity element, and let \(\pi\) be a representation of \(A\)
in \(H\). Then \(\pi\) extends uniquely to a representation
\(\widetilde\pi\) of \(\widetilde A\) in \(H\) such that
\(\widetilde\pi(1)=1\). This extension is called canonical.

\ResultHeading{2.2.10}{2.2.10.}
Let \(A\) be an involutive Banach algebra and \(\pi\) a representation of
\(A\) in \(H\). Recall from \XRef{1.3.7}{1.3.7} that
\(\lVert\pi(x)\rVert\leq\lVert x\rVert\) for every \(x\in A\). If \(A\)
has an approximate identity \((u_i)\) and \(\pi\) is nondegenerate, then
\(\pi(u_i)\) converges strongly to \(1\). Indeed, for every \(y\in A\) and
\(\xi\in H\), \(\pi(u_i)(\pi(y)\xi)=\pi(u_iy)\xi\) converges strongly to
\(\pi(y)\xi\), since
\[
 \lVert\pi(u_iy)-\pi(y)\rVert\leq\lVert u_iy-y\rVert\longrightarrow0;
\]
the vectors \(\pi(y)\xi\) form a total set in \(H\), while
\(\lVert\pi(u_i)\rVert\leq\lVert u_i\rVert\leq1\). Hence
\(\pi(u_i)\eta\) converges strongly to \(\eta\) for every \(\eta\in H\).

Bibliography: \cite{ref41,ref108,ref126}.

% Source PDF physical page 038
\BookSubsection{2.3}{Topologically Irreducible Representations}

\ResultHeading{2.3.1}{2.3.1. Proposition.}
\begin{itshape}
Let \(A\) be an involutive algebra, \(H\) a Hilbert space, and \(\pi\) a
representation of \(A\) in \(H\). The following conditions are equivalent:
\begin{enumerate}[label=\textup{(\roman*)}]
 \item The only closed vector subspaces of \(H\) invariant under \(\pi(A)\)
 are \(0\) and \(H\).
 \item The commutant of \(\pi(A)\) in \(\mathcal L(H)\) consists only of
 scalar operators.
 \item Every nonzero vector of \(H\) is cyclic for \(\pi\), or else \(\pi\)
 is the zero representation of dimension \(1\).
\end{enumerate}
\end{itshape}

(i) \(\Rightarrow\) (iii): suppose (i) holds, and let
\(\xi\in H\), \(\xi\neq0\). If \(\pi(A)\xi\) is not dense in \(H\), its
closure is \(0\) by (i), and hence \(\pi(A)\xi=0\). Then
\(\mathbb C\xi\) is invariant under \(\pi(A)\), so \(H=\mathbb C\xi\) and
\(\pi\) is the zero representation of dimension \(1\).

(iii) \(\Rightarrow\) (i): suppose (iii) holds. Let \(K\neq0\) be a closed
vector subspace of \(H\) invariant under \(\pi(A)\). We must prove that
\(K=H\). This is clear if \(\dim\pi=1\). Otherwise every nonzero vector of
\(H\) is cyclic for \(\pi\). Let \(\xi\in K\), \(\xi\neq0\). Then
\[
 \pi(A)\xi\subset K
 \qquad\text{and}\qquad
 \overline{\pi(A)\xi}=H,
\]
so \(K=H\).

(ii) \(\Rightarrow\) (i): suppose (ii) holds. Let \(K\) be a closed vector
subspace of \(H\) invariant under \(\pi(A)\). Then \(P_K\) commutes with
\(\pi(A)\) by \XRef{2.2.4}{2.2.4}; hence \(P_K\) is scalar and therefore
\(P_K=0\) or \(1\), that is, \(K=0\) or \(H\).

(i) \(\Rightarrow\) (ii): suppose (i) holds. Let \(T\in\mathcal L(H)\)
commute with \(\pi(A)\), and let us prove that \(T\) is scalar. Since
\(T+T^{\ast}\) and \(i(T-T^{\ast})\) commute with \(\pi(A)\), it is enough
to consider the case in which \(T\) is hermitian. The spectral projections
of \(T\) then commute with \(\pi(A)\), and hence, by (i), are all \(0\) or
\(1\). Thus \(T\) is scalar.

\ResultHeading{2.3.2}{2.3.2. Definition.}
\begin{itshape}
Let \(A\) be an involutive algebra, \(H\) a Hilbert space, and \(\pi\) a
representation of \(A\) in \(H\). The representation \(\pi\) is said to be
\emph{topologically irreducible} if \(H\neq0\) and if \(\pi\) satisfies the
equivalent conditions of \XRef{2.3.1}{2.3.1}.
\end{itshape}

Such a representation is either the zero representation of dimension \(1\),
or it is nonzero and nondegenerate. We denote by \(\widehat A\) the set of
classes of nonzero topologically irreducible representations of \(A\).

Recall that irreducibility of \(\pi\) in the algebraic sense means that the
only vector subspaces of \(H\) invariant under \(\pi(A)\) are
\(\{0\}\) and \(H\). If \(\dim\pi=+\infty\), this is a much stronger
condition than topological irreducibility. We shall see, however, in
\XRef{2.8.4}{2.8.4} that, if \(A\) is a \(C^{*}\)-algebra, the two notions
are equivalent.

% Source PDF physical page 039
A finite-dimensional representation \(\pi\) of \(A\) is a Hilbert direct
sum of topologically irreducible representations
(\XRef{2.3.5}{2.3.5}). We shall see in \XRef{8.5.2}{8.5.2} that this result
extends, to a certain extent, to infinite-dimensional representations, but
in a much subtler form. This is one reason why we shall be particularly
interested below in topologically irreducible representations. Given an
involutive algebra \(A\), an essential problem is to list all its
topologically irreducible representations (up to equivalence). It may very
well happen that there is no nonzero topologically irreducible
representation of \(A\) (or even no nonzero representation at all). We
shall nevertheless see in \XRef{2.7.3}{2.7.3} that every
\(C^{*}\)-algebra has sufficiently many topologically irreducible
representations.

\ResultHeading{2.3.3}{2.3.3.}
Let \(A\) be a separable involutive normed algebra, and let \(\pi\) be a
topologically irreducible representation of \(A\) in a Hilbert space \(H\).
Then \(H\) is separable. This is clear if \(\pi\) is the zero representation
of dimension \(1\). Otherwise, let \(\xi\) be a nonzero vector of \(H\);
if \((x_n)\) is a dense sequence in \(A\), then the vectors
\(\pi(x_n)\xi\) are dense in \(H\). The same reasoning shows that, for every
involutive normed algebra \(B\), the dimensions of the topologically
irreducible representations of \(B\) are bounded above by a fixed cardinal.

\ResultHeading{2.3.4}{2.3.4.}
Let \(A\) be an involutive algebra, let \(\pi,\pi'\) be topologically
irreducible representations of \(A\), and let \(n\) be their intertwining
number. If \(n=0\), then \(\pi\) and \(\pi'\) are inequivalent. If \(n>0\),
let \(T:H_{\pi}\to H_{\pi'}\) be a nonzero intertwining operator. By
\XRef{2.2.2}{2.2.2}, \(T^{\ast}T\) and \(TT^{\ast}\) are nonzero scalar
operators, and \(\pi,\pi'\) are equivalent.

\ResultHeading{2.3.5}{2.3.5.}
Let \(A\) be an involutive algebra. We analyze its finite-dimensional
representations. What follows may be regarded both as a special case of
later theorems and as a special case of theorems of pure algebra; but, to
make the reader's task easier, we prefer to give a direct analysis now.

Let \(\pi\) be a finite-dimensional representation of \(A\). Then
\[
 \pi=\pi_1\oplus\cdots\oplus\pi_n,
\]
where the \(\pi_i\) are irreducible. This is clear if \(\dim\pi=0\) (with
\(n=0\)). Suppose \(\dim\pi=q\) and the assertion proved when
\(\dim\pi<q\). If \(\pi\) is irreducible, the assertion is again clear.
Otherwise \(\pi=\pi'\oplus\pi''\), with
\(\dim\pi'<q\) and \(\dim\pi''<q\), and it is enough to apply the induction
hypothesis. The decomposition
\(\pi=\pi_1\oplus\cdots\oplus\pi_n\) is not unique (for example, the
representation \(\lambda\mapsto\lambda\,1\) of \(\mathbb C\) in
\(\mathbb C^n\) admits, for \(n>1\), infinitely many decompositions into
one-dimensional representations). We shall nevertheless establish a
uniqueness result. Let \(\rho_1,\rho_2\) be two
% Source PDF physical page 040
irreducible subrepresentations of \(\pi\), and let \(P_1,P_2\) be the
projections of \(H_{\pi}\) onto \(H_{\rho_1}\) and \(H_{\rho_2}\). They
commute with \(\pi(A)\). Hence the restriction of \(P_2\) to
\(H_{\rho_1}\) is an intertwining operator for \(\rho_1\) and \(\rho_2\).
Thus, if \(H_{\rho_1}\) and \(H_{\rho_2}\) are not orthogonal, then
\(\rho_1\sim\rho_2\) by \XRef{2.3.4}{2.3.4}. This proves that every
irreducible subrepresentation of \(\pi\) is equivalent to one of the
\(\pi_i\). Regrouping the \(\pi_i\), we see that
\[
 \pi=\nu_1\oplus\cdots\oplus\nu_m,
\]
where each \(\nu_i\) is a multiple \(p_i\nu_i'\) of an irreducible
representation \(\nu_i'\), and the \(\nu_i'\) are pairwise inequivalent.
If \(\rho\) is an irreducible subrepresentation of \(\pi\), the preceding
argument shows that \(H_{\rho}\) is orthogonal to every \(H_{\nu_j}\)
except one. Hence \(H_{\rho}\) is contained in that one \(H_{\nu_i}\).
This proves that each subspace \(H_{\nu_i}\) is uniquely determined: it is
the subspace of \(H_{\pi}\) spanned by the spaces of the subrepresentations
of \(\pi\) equivalent to \(\nu_i'\).

Thus, in a decomposition
\(\pi=p_1\nu_1'\oplus\cdots\oplus p_m\nu_m'\) of \(\pi\), with the
\(\nu_i'\) irreducible and pairwise inequivalent, the integers \(p_i\), the
classes of the \(\nu_i'\), and the spaces of the \(p_i\nu_i'\) are uniquely
determined.

Bibliography: \cite{ref108,ref126}.

\BookSubsection{2.4}{Positive Functionals and Representations}

\ResultHeading{2.4.1}{2.4.1. Proposition.}
\begin{itshape}
Let \(A\) be an involutive algebra.
\begin{enumerate}[label=\textup{(\roman*)}]
 \item If \(\pi\) is a representation of \(A\) in \(H\), and if
 \(\xi\in H\), then \(x\mapsto(\pi(x)\xi\mid\xi)\) is a positive functional
 on \(A\).
 \item Let \(\pi,\pi'\) be representations of \(A\) in \(H,H'\), and let
 \(\xi\) (respectively \(\xi'\)) be cyclic for \(\pi\) (respectively
 \(\pi'\)). If
 \[
  (\pi(x)\xi\mid\xi)=(\pi'(x)\xi'\mid\xi')
  \qquad\text{for every }x\in A,
 \]
 there exists a unique isomorphism from \(H\) onto \(H'\) that carries
 \(\pi\) into \(\pi'\) and \(\xi\) into \(\xi'\).
\end{enumerate}
\end{itshape}

We have
\[
 (\pi(x^{\ast}x)\xi\mid\xi)
 =(\pi(x)^{\ast}\pi(x)\xi\mid\xi)
 =\lVert\pi(x)\xi\rVert^{2}\geq0,
\]
which proves (i). Now assume the hypotheses of (ii). For arbitrary
\(x,y\in A\),
\[
 (\pi(x)\xi\mid\pi(y)\xi)
 =(\pi(y^{\ast}x)\xi\mid\xi)
 =(\pi'(y^{\ast}x)\xi'\mid\xi')
 =(\pi'(x)\xi'\mid\pi'(y)\xi').
\]
Since the \(\pi(x)\xi\) (respectively the \(\pi'(x)\xi'\)) are dense in
\(H\) (respectively in \(H'\)), there exists an isomorphism \(U:H\to H'\)
such that \(U(\pi(x)\xi)=\pi'(x)\xi'\) for every \(x\in A\). Let us show
that \(U\) carries \(\pi\) into \(\pi'\), that is,
\(U\pi(x)=\pi'(x)U\) for every \(x\in A\). For every \(y\in A\),
\[
 (U\pi(x))(\pi(y)\xi)=U\pi(xy)\xi=\pi'(xy)\xi'
 =\pi'(x)(\pi'(y)\xi')=(\pi'(x)U)(\pi(y)\xi).
\]
% Source PDF physical page 041
Since the vectors \(\pi(y)\xi\) are dense in \(H\), this indeed gives
\(U\pi(x)=\pi'(x)U\). Moreover, for every \(x\in A\),
\[
 (\xi'\mid\pi'(x)\xi')=(\xi\mid\pi(x)\xi)
 =(U\xi\mid U\pi(x)\xi)=(U\xi\mid\pi'(x)\xi'),
\]
and therefore \(\xi'=U\xi\). Finally, the asserted uniqueness of \(U\) is
immediate, since its values are prescribed on \(\pi(A)\xi\).

\ResultHeading{2.4.2}{2.4.2.}
Retain the preceding notation. The functional
\(x\mapsto(\pi(x)\xi\mid\xi)\) on \(A\) is called the
\emph{functional defined by \(\pi\) and \(\xi\)}. With \(\pi\) fixed and
\(\xi\) varying in \(H\), one obtains the \emph{functionals associated with
\(\pi\)}. If \(S\) is a set of representations of \(A\), the functionals
associated with the various elements of \(S\) are called the functionals
associated with \(S\).

Let \(H\) be a Hilbert space, \(B\) an involutive subalgebra of
\(\mathcal L(H)\), and \(\xi\in H\). Write \(\omega_{\xi}\) for the positive
functional on \(B\) defined by the identity representation of \(B\) and
\(\xi\), that is, \(x\mapsto(x\xi\mid\xi)\). A positive functional on \(B\)
is called a \emph{vector functional} if it is equal to \(\omega_{\xi}\)
for a suitable \(\xi\in H\).

\ResultHeading{2.4.3}{2.4.3.}
Let \(A\) be an involutive Banach algebra admitting an approximate identity
\((u_i)\), let \(\pi\) be a nondegenerate representation of \(A\) in \(H\),
let \(\xi\in H\), and let \(f\) be the positive functional defined by
\(\pi\) and \(\xi\). Then
\(\lVert f\rVert=(\xi\mid\xi)\). Indeed, by \XRef{2.1.5}{2.1.5(v)},
\[
 \lVert f\rVert=\lim_i f(u_i)
 =\lim_i(\pi(u_i)\xi\mid\xi),
\]
and \(\pi(u_i)\) converges strongly to \(1\) by \XRef{2.2.10}{2.2.10}.
It follows that, if \(\widetilde A\) denotes the involutive algebra obtained
from \(A\) by adjoining an identity, and if \(\widetilde f\) and
\(\widetilde\pi\) are the canonical extensions of \(f\) and \(\pi\) to
\(\widetilde A\), then
\[
 \widetilde f(x)=(\widetilde\pi(x)\xi\mid\xi)
 \qquad\text{for every }x\in\widetilde A.
\]
In particular, if \(\pi\) is the nondegenerate identity representation of
a \(C^{*}\)-subalgebra \(A\) of \(\mathcal L(H)\) that does not contain the
identity operator, then the canonical extension to
\(\widetilde A=A+\mathbb C1\) of \(\omega_{\xi}|_A\) is
\(\omega_{\xi}|_{\widetilde A}\).

\ResultHeading{2.4.4}{2.4.4. Proposition.}
\begin{itshape}
Let \(A\) be an involutive Banach algebra having an approximate identity,
let \(\widetilde A\) be the involutive algebra obtained from \(A\) by
adjoining an identity element, let \(f\) be a continuous positive
functional on \(A\), let \(\widetilde f\) be its canonical extension to
\(\widetilde A\), let \(N\) be the left ideal of \(\widetilde A\) formed by
the \(x\in\widetilde A\) such that \(\widetilde f(x^{\ast}x)=0\), let
\(A_f'\) be the separated pre-Hilbert space \(\widetilde A/N\), and let
\(A_f\) be the Hilbert-space completion of \(A_f'\). For every
\(x\in\widetilde A\), let \(\pi'(x)\) be the operator on
\(\widetilde A/N\) induced on passing to the quotient by left multiplication
by \(x\) in \(\widetilde A\). Let \(\xi\) be the canonical image of \(1\)
in \(A_f'\).
\begin{enumerate}[label=\textup{(\roman*)}]
 \item Every \(\pi'(x)\) extends uniquely to a continuous linear operator
 \(\pi(x)\) on \(A_f\).
% Source PDF physical page 042
 \item The map \(x\mapsto\pi(x)\), \(x\in A\), is a representation of
 \(A\) in \(A_f\).
 \item \(\xi\) is cyclic for \(\pi(A)\).
 \item \(f(x)=(\pi(x)\xi\mid\xi)\) for every \(x\in A\).
\end{enumerate}
\end{itshape}

By \XRef{2.1.5}{2.1.5(ii)}, for \(x,y\in\widetilde A\),
\begin{align*}
 (\pi'(x)\pi'(y)\xi\mid\pi'(x)\pi'(y)\xi)
 &=\widetilde f(y^{\ast}x^{\ast}xy)\\
 &\leq\lVert x^{\ast}x\rVert\,\widetilde f(y^{\ast}y)\\
 &=\lVert x^{\ast}x\rVert
   (\pi'(y)\xi\mid\pi'(y)\xi),
\end{align*}
which proves (i). It is clear that the \(\pi'(x)\), and hence the
\(\pi(x)\), are representations as far as the algebraic structure is
concerned. For \(x,y,z\in\widetilde A\),
\begin{align*}
 (\pi(x)\pi(y)\xi\mid\pi(z)\xi)
 &=\widetilde f(z^{\ast}xy)
 =\widetilde f((x^{\ast}z)^{\ast}y)\\
 &=(\pi(y)\xi\mid\pi(x^{\ast})\pi(z)\xi),
\end{align*}
and therefore \(\pi(x)^{\ast}=\pi(x^{\ast})\), which proves (ii). The set
\(\pi(A)\xi\) is the canonical image of \(A\) in \(A_f'\), and hence is
dense in \(A_f\) by \XRef{2.1.5}{2.1.5(vii)}; this proves (iii). Finally,
for every \(x\in\widetilde A\),
\[
 (\pi(x)\xi\mid\xi)=\widetilde f(1^{\ast}x1)=\widetilde f(x),
\]
which proves (iv).

The representation \(\pi\) and vector \(\xi\) are said to be
\emph{defined by \(f\)}. We denote them by \(\pi_f\) and \(\xi_f\).

\ResultHeading{2.4.5}{2.4.5.}
Retain the preceding notation. Let \(M\) be the left ideal of \(A\) formed
by the \(x\in A\) such that \(f(x^{\ast}x)=0\). The canonical image of \(A\)
in \(A_f'\) is naturally identified with \(A/M\), and is dense in \(A_f\).
Thus \(A_f\) may be defined as the completion of the separated pre-Hilbert
space \(A/M\), and, for \(x\in A\), \(\pi(x)\) may be defined as the
continuous extension of left multiplication by \(x\) on \(A/M\). This
avoids introducing \(\widetilde A\) and \(\widetilde f\). From this point
of view, however, the definition of \(\xi\) would be less easy.

\ResultHeading{2.4.6}{2.4.6.}
Let \(A\) be an involutive Banach algebra having an approximate identity,
let \(f\) be a continuous positive functional on \(A\), and let \(\pi,\xi\)
be the representation and vector defined by \(f\). Then the positive
functional defined by \(\pi,\xi\) is \(f\), by
\XRef{2.4.4}{2.4.4(iv)}.

Conversely, start with a representation \(\pi\) of \(A\) in a Hilbert space
and a cyclic vector \(\xi\) for \(\pi\). Let \(f\) be the positive
functional defined by \(\pi,\xi\); it is continuous because \(\pi\) is
continuous. Let \(\pi',\xi'\) be the representation and vector defined by
\(f\). We have
\[
 (\pi(x)\xi\mid\xi)=f(x)=(\pi'(x)\xi'\mid\xi')
 \qquad\text{for every }x\in A,
\]
and \(\xi'\) is cyclic for \(\pi'\). By \XRef{2.4.1}{2.4.1(ii)}, there is a
unique isomorphism from \(H_{\pi}\) onto \(H_{\pi'}\) carrying \(\pi\) into
\(\pi'\) and \(\xi\) into \(\xi'\).

In particular, let \(\pi\) be a nonzero topologically irreducible
representation of \(A\). Every nonzero vector of \(H_{\pi}\) is cyclic for
\(\pi\); hence \(\pi\) is defined by any nonzero functional associated
with \(\pi\).

% Source PDF physical page 043
\ResultHeading{2.4.7}{2.4.7. Proposition.}
\begin{itshape}
Let \(A\) be an involutive Banach algebra having an approximate identity,
and let \(\pi\) be a representation of \(A\). For every \(x\in A\),
\[
 \lVert\pi(x)\rVert^{2}
 =\sup f(x^{\ast}x),
\]
where \(f\) ranges over the positive functionals associated with \(\pi\)
such that \(\lVert f\rVert\leq1\).
\end{itshape}

By \XRef{2.2.6}{2.2.6}, we may reduce to the case in which \(\pi\) is
nondegenerate. By \XRef{2.4.3}{2.4.3}, the positive functionals associated
with \(\pi\) and of norm at most \(1\) are the functionals
\(\omega_{\xi}\circ\pi\), where \(\xi\in H_{\pi}\) and
\(\lVert\xi\rVert\leq1\). But
\begin{align*}
 \lVert\pi(x)\rVert^{2}
 &=\sup_{\lVert\xi\rVert\leq1}
   (\pi(x)\xi\mid\pi(x)\xi)\\
 &=\sup_{\lVert\xi\rVert\leq1}
   (\omega_{\xi}\circ\pi)(x^{\ast}x).
\end{align*}

\ResultHeading{2.4.8}{2.4.8. Proposition.}
\begin{itshape}
Let \(A\) be an involutive Banach algebra having an approximate identity,
let \(f\) be a continuous positive functional on \(A\), and let
\(\pi,\xi\) be the representation and vector defined by \(f\).
\begin{enumerate}[label=\textup{(\roman*)}]
 \item If \(x_0\in A\), then the functional
 \(x\mapsto f(x_0^{\ast}xx_0)\) is associated with \(\pi\).
 \item If \(f'\) is a positive functional associated with \(\pi\), then
 \(f'\) is a norm limit of functionals
 \(x\mapsto f(x_0^{\ast}xx_0)\), where \(x_0\in A\).
\end{enumerate}
\end{itshape}

If \(x_0\in A\), then
\[
 f(x_0^{\ast}xx_0)
 =(\pi(x_0^{\ast}xx_0)\xi\mid\xi)
 =(\pi(x)\pi(x_0)\xi\mid\pi(x_0)\xi),
\]
which proves (i).

Let \(\xi'\in H_{\pi}\) and \(f'=\omega_{\xi'}\circ\pi\). For every
\(\varepsilon>0\), there exists \(x_0\in A\) such that
\(\lVert\pi(x_0)\xi-\xi'\rVert\leq\varepsilon\). Then, for every \(x\in A\),
\begin{align*}
 |f'(x)-f(x_0^{\ast}xx_0)|
 &=|(\pi(x)\xi'\mid\xi')
   -(\pi(x)\pi(x_0)\xi\mid\pi(x_0)\xi)|\\
 &\leq\lVert\pi(x)\xi'\rVert\,
       \lVert\xi'-\pi(x_0)\xi\rVert\\
 &\quad+\lVert\pi(x)\xi'-\pi(x)\pi(x_0)\xi\rVert\,
       \lVert\pi(x_0)\xi\rVert\\
 &\leq\lVert x\rVert\,\lVert\xi'\rVert\varepsilon
  +\lVert x\rVert\varepsilon(\lVert\xi'\rVert+\varepsilon)\\
 &=\lVert x\rVert(2\varepsilon\lVert\xi'\rVert+\varepsilon^{2}),
\end{align*}
and \(2\varepsilon\lVert\xi'\rVert+\varepsilon^{2}\) is arbitrarily small.

\ResultHeading{2.4.9}{2.4.9.}
Let \(A\) be a \(C^{*}\)-algebra, let \(I\) be a closed two-sided ideal of
\(A\), let \(B=A/I\), and let \(\omega:A\to B\) be the canonical
homomorphism. If \(f\) is a positive functional on \(A\) such that
\(f(I)=0\), then \(f\) induces by passage to the quotient a positive
functional \(g\) on \(B\). For every \(x\in A\),
\[
 (\pi_g(\omega(x))\xi_g\mid\xi_g)
 =g(\omega(x))=f(x),
\]
and hence \(\pi_g\circ\omega\) and \(\xi_g\) identify, respectively, with
\(\pi_f\) and \(\xi_f\) by \XRef{2.4.1}{2.4.1}.

\ResultLabel{Proposition.}
\begin{itshape}
Let \(A\) be a \(C^{*}\)-algebra, \(f\) a positive functional on \(A\), and
\(I\) a closed two-sided ideal of \(A\). Then \(f(I)=0\) if and only if
\(\pi_f(I)=0\).
\end{itshape}

% Source PDF physical page 044
We have \(\ker\pi_f\subset\ker f\), so \(\pi_f(I)=0\) implies \(f(I)=0\).
If \(f(I)=0\), then, with the preceding notation,
\(\pi_f\sim\pi_g\circ\omega\), and consequently \(\pi_f(I)=0\).

\ResultHeading{2.4.10}{2.4.10. Corollary.}
\begin{itshape}
\(\ker\pi_f\) is the largest closed two-sided ideal of \(A\) contained in
\(\ker f\).
\end{itshape}

\ResultHeading{2.4.11}{2.4.11. Corollary.}
\begin{itshape}
Let \(A\) be a \(C^{*}\)-algebra and let \(f,g\) be two positive
functionals on \(A\). The following conditions are equivalent:
\begin{enumerate}[label=\textup{(\roman*)}]
 \item \(\ker\pi_f\subset\ker\pi_g\).
 \item \(g\) vanishes on \(\ker\pi_f\).
\end{enumerate}
\end{itshape}

Bibliography: \cite{ref41,ref42,ref52,ref105,ref108,ref126,ref134}.

\BookSubsection{2.5}{Pure Functionals and Irreducible Representations}

We know how to associate representations with positive functionals. When
does this procedure yield topologically irreducible representations? We
shall now solve this problem.

\ResultHeading{2.5.1}{2.5.1. Proposition.}
\begin{itshape}
Let \(A\) be an involutive algebra, let \(\pi\) be a representation of \(A\)
in \(H\), let \(\xi\in H\), and let \(f\) be the positive functional on
\(A\) defined by \(\pi\) and \(\xi\).
\begin{enumerate}[label=\textup{(\roman*)}]
 \item If \(T\) is a hermitian operator on \(H\), commuting with the
 \(\pi(x)\), and satisfying \(0\leq T\leq1\), then the functional
 \[
  x\longmapsto(\pi(x)T\xi\mid T\xi)
  =(\pi(x)\xi\mid T^{2}\xi)
 \]
 on \(A\) is a positive functional \(f_T\) majorized by \(f\).
 \item If \(\xi\) is cyclic for \(\pi\), the map \(T\mapsto f_T\) is
 injective.
 \item If \(A\) is an involutive Banach algebra having an approximate
 identity, then every continuous positive functional on \(A\) majorized by
 \(f\) is of the form \(f_T\).
\end{enumerate}
\end{itshape}

Since \(f_T=\omega_{T\xi}\circ\pi\), we have \(f_T\geq0\). If \(x\in A\),
\begin{align*}
 f_T(x^{\ast}x)
 &=(\pi(x^{\ast}x)T\xi\mid T\xi)
  =\lVert\pi(x)T\xi\rVert^{2}\\
 &=\lVert T\pi(x)\xi\rVert^{2}
 \leq\lVert\pi(x)\xi\rVert^{2}=f(x^{\ast}x),
\end{align*}
so \(f_T\leq f\). This proves (i).

If \(f_T=f_{T'}\), then
\[
 (\pi(x)\xi\mid T^{2}\xi)
 =(\pi(x)\xi\mid T'^{\,2}\xi)
\]
for every \(x\in A\), and hence \(T^{2}\xi=T'^{\,2}\xi\) if \(\xi\) is
cyclic for \(\pi(A)\). Since \(\xi\) is then separating for the commutant
of \(\pi(A)\) (\XRef{A.14}{A 14}), it follows that \(T^{2}=T'^{\,2}\),
and hence \(T=T'\), because \(T,T'\geq0\). This proves (ii).

Let \(g\) be a positive functional on \(A\) majorized by \(f\). For
\(x,y\in A\), we have
\[
 |g(y^{\ast}x)|^{2}
 \leq g(x^{\ast}x)g(y^{\ast}y)
 \leq f(x^{\ast}x)f(y^{\ast}y)
 =\lVert\pi(x)\xi\rVert^{2}\lVert\pi(y)\xi\rVert^{2}.
\]
% Source PDF physical page 045
Thus, on the vector subspace \(\pi(A)\xi\) of \(H\), the formula
\[
 (\pi(x)\xi\mid\pi(y)\xi)_0=g(y^{\ast}x)
\]
defines a unique continuous sesquilinear form, plainly hermitian and
positive. Hence there exists on \(X=\overline{\pi(A)\xi}\) a hermitian
operator \(T_0\) such that \(0\leq T_0\leq1\) and
\[
 g(y^{\ast}x)=(\pi(x)\xi\mid T_0\pi(y)\xi).
\]
For \(x,y,z\in A\),
\begin{align*}
 (\pi(y)\xi\mid T_0\pi(z)\pi(x)\xi)
 &=g((zx)^{\ast}y)=g(x^{\ast}(z^{\ast}y))\\
 &=(\pi(z^{\ast}y)\xi\mid T_0\pi(x)\xi)\\
 &=(\pi(y)\xi\mid\pi(z)T_0\pi(x)\xi),
\end{align*}
and therefore \(T_0\pi(z)=\pi(z)T_0\) on \(X\). Moreover, \(X\) is
invariant under \(\pi(A)\), so \(P_X\) commutes with \(\pi(A)\). Thus the
operator \(T_0P_X\), extended by zero on \(H\ominus X\), is a hermitian
operator on \(H\), commuting with every \(\pi(z)\). Let \(T\) be its
positive square root. We have \(0\leq T\leq1\), and
\begin{align*}
 g(y^{\ast}x)
 &=(\pi(x)\xi\mid T^{2}\pi(y)\xi)
  =(\pi(x)T\xi\mid\pi(y)T\xi)\\
 &=(\pi(y^{\ast}x)T\xi\mid T\xi)
  =f_T(y^{\ast}x).
\end{align*}
Finally, suppose that \(A\) is an involutive Banach algebra admitting an
approximate identity \((u_i)\), and that \(g\) is continuous. Then
\[
 g(y^{\ast})=\lim_i g(y^{\ast}u_i),
 \qquad
 f_T(y^{\ast})=\lim_i f_T(y^{\ast}u_i),
 \qquad\text{so}\qquad g=f_T.
\]

\ResultHeading{2.5.2}{2.5.2. Definition.}
\begin{itshape}
Let \(A\) be an involutive normed algebra and \(f\) a continuous positive
functional on \(A\). The functional \(f\) is said to be \emph{pure} if
\(f\neq0\) and if every continuous positive functional on \(A\) majorized
by \(f\) is of the form \(\lambda f\), where \(0\leq\lambda\leq1\). We
denote by \(P(A)\) the set of pure states of \(A\).
\end{itshape}

Let \(\Omega\) be a locally compact space, and let \(A\) be the
\(C^{*}\)-algebra of continuous complex-valued functions on \(\Omega\)
vanishing at infinity. The pure positive functionals on \(A\) identify with
the positive measures on \(\Omega\) whose support consists of one point,
that is, with the measures
\[
 f\longmapsto\lambda f(\omega)
 \qquad(\lambda>0,\ \omega\text{ a fixed point of }\Omega).
\]
It follows that the pure states of a commutative \(C^{*}\)-algebra are the
characters of that algebra.

\ResultHeading{2.5.3}{2.5.3.}
Let \(A\) be an involutive Banach algebra admitting an approximate identity,
let \(\widetilde A\) be an involutive Banach algebra obtained from \(A\) by
adjoining an identity element, let \(f\) be a continuous positive functional
on \(A\), and let \(\widetilde f\) be its canonical extension to
\(\widetilde A\). Then \(f\) is pure if and only if \(\widetilde f\) is
pure. Indeed, first, the conditions \(f=0\) and \(\widetilde f=0\) are
equivalent. Further, as \(g\) ranges over the continuous positive
functionals on \(A\) majorized by \(f\), \(\widetilde g\) ranges over the
positive functionals
% Source PDF physical page 046
on \(\widetilde A\) majorized by \(\widetilde f\), by
\XRef{2.1.7}{2.1.7}. Finally, \(g\) is of the form \(\lambda f\), with
\(0\leq\lambda\leq1\), if and only if
\(\widetilde g=\lambda\widetilde f\).

\ResultHeading{2.5.4}{2.5.4. Proposition.}
\begin{itshape}
Let \(A\) be an involutive Banach algebra having an approximate identity,
let \(f\) be a continuous positive functional on \(A\), and let \(\pi\) be
the representation of \(A\) defined by \(f\). Then \(\pi\) is nonzero and
topologically irreducible if and only if \(f\) is pure.
\end{itshape}

Let \(\xi\) be the vector of \(H_{\pi}\) defined by \(f\). Suppose \(f\) is
pure. Let \(E\) be a projection on \(H_{\pi}\) commuting with \(\pi(A)\).
The functional
\[
 x\longmapsto(\pi(x)E\xi\mid E\xi)
\]
on \(A\) is continuous, positive, and majorized by \(f\), by
\XRef{2.5.1}{2.5.1(i)}; it is therefore equal to \(\lambda f\), with
\(0\leq\lambda\leq1\). Hence, for every \(x\in A\),
\[
 (\pi(x)E\xi\mid E\xi)
 =(\pi(x)\lambda^{1/2}\xi\mid\lambda^{1/2}\xi).
\]
By \XRef{2.5.1}{2.5.1(ii)}, \(E=\lambda^{1/2}1\), and therefore \(E=0\)
or \(1\). On the other hand, there exist \(x\in A\) with \(f(x)\neq0\),
and hence \((\pi(x)\xi\mid\xi)\neq0\). This proves that \(\pi\) is nonzero
and topologically irreducible.

Conversely, suppose that \(\pi\) is nonzero and topologically irreducible.
Then there exist \(x\in A\) such that \((\pi(x)\xi\mid\xi)\neq0\), so
\(f\neq0\). Let \(g\) be a continuous positive functional on \(A\)
majorized by \(f\). By \XRef{2.5.1}{2.5.1(iii)}, there exists a hermitian
operator \(T\in\pi(A)'\) such that \(0\leq T\leq1\) and
\[
 g(x)=(\pi(x)T\xi\mid T\xi)
\]
for every \(x\in A\). Since \(\pi\) is topologically irreducible, we have
\(T=\mu1\), with \(0\leq\mu\leq1\), whence \(g=\mu^{2}f\). Thus \(f\) is
pure, and the proposition is proved.

Proposition \XRef{2.5.4}{2.5.4} defines a canonical map
\[
 P(A)\longrightarrow\widehat A.
\]
By \XRef{2.4.6}{2.4.6} and \XRef{2.5.4}{2.5.4}, this map is surjective. If
\(\pi\in\widehat A\), its inverse image in \(P(A)\) is the set of states
associated with \(\pi\) (on this point, see \XRef{2.5.7}{2.5.7}).

\ResultHeading{2.5.5}{2.5.5. Proposition.}
\begin{itshape}
Let \(A\) be an involutive Banach algebra having an approximate identity,
and let \(B\) be the set of continuous positive functionals on \(A\) of
norm at most \(1\).
\begin{enumerate}[label=\textup{(\roman*)}]
 \item \(B\) is convex and compact for the weak topology
 \(\sigma(A',A)\) of the dual \(A'\) of \(A\).
 \item The extreme points of \(B\) are \(0\) and the pure states.
 \item \(B\) is the weakly closed convex hull of \(0\) and the set of pure
 states.
\end{enumerate}
\end{itshape}

The set \(B\) is a weakly closed convex subset of the unit ball of \(A'\),
and that ball is weakly compact. This proves (i).

% Source PDF physical page 047
Let us show that \(0\) is an extreme point of \(B\). If \(f\in B\) and
\(-f\in B\), then \(f(x^{\ast}x)=0\) for every \(x\in A\), and hence
\[
 |f(x)|^{2}\leq\lVert f\rVert f(x^{\ast}x)=0
\]
by \XRef{2.1.5}{2.1.5}; thus \(f=0\).

Let \(f\) be a pure state. Suppose
\[
 f=\lambda f_1+(1-\lambda)f_2,
 \qquad 0<\lambda<1,\quad f_1,f_2\in B.
\]
Then \(\lambda f_1\) is majorized by \(f\), so
\(\lambda f_1=\mu f\), with \(0\leq\mu\leq1\). Since
\[
 1=\lVert f\rVert
 =\lambda\lVert f_1\rVert+(1-\lambda)\lVert f_2\rVert
\]
and \(\lVert f_1\rVert\leq1\), \(\lVert f_2\rVert\leq1\), necessarily
\(\lVert f_1\rVert=\lVert f_2\rVert=1\), and consequently
\(\lambda=\mu\) and \(f_1=f\), \(f_2=f\). Thus \(f\) is an extreme point
of \(B\).

Conversely, let \(f\) be an extreme point of \(B\), distinct from \(0\).
Clearly \(\lVert f\rVert=1\). Suppose \(f=f_1+f_2\), where \(f_1,f_2\)
are nonzero continuous positive functionals. Put
\(\lVert f_1\rVert=\lambda\), so that
\(\lVert f_2\rVert=1-\lambda\), and set
\[
 g_1=\lambda^{-1}f_1,\qquad
 g_2=(1-\lambda)^{-1}f_2.
\]
Then \(f=\lambda g_1+(1-\lambda)g_2\), where \(g_1,g_2\in B\). Since \(f\)
is extreme, \(f=g_1=g_2\). Hence \(f_1=\lambda f\) and
\(f_2=(1-\lambda)f\), so \(f\) is a pure state. This proves (ii).

Finally, (iii) follows from (i), (ii), and the Krein--Milman theorem.

\ResultHeading{2.5.6}{2.5.6.}
Retain the preceding notation, and suppose in addition that \(A\) has an
identity element. Then \(B\) is the set of positive functionals \(f\) on
\(A\) such that \(f(1)\leq1\), and the state space \(E(A)\) of \(A\) is
the set of positive functionals \(f\) on \(A\) such that \(f(1)=1\). It
follows that \(E(A)\) is convex and weakly compact, and that the set of
extreme points of \(E(A)\) is the set of extreme points of \(B\) belonging
to \(E(A)\), namely \(P(A)\). Thus \(E(A)\) is the closed convex hull of
\(P(A)\).

\ResultHeading{2.5.7}{2.5.7. Proposition.}
\begin{itshape}
Let \(A\) be an involutive algebra, let \(\pi\) be a nonzero topologically
irreducible representation of \(A\), let \(\xi_1,\xi_2\) be two vectors of
\(H_{\pi}\), and let \(f_1,f_2\) be the positive functionals defined by
\((\pi,\xi_1)\) and \((\pi,\xi_2)\). Then \(f_1=f_2\) if and only if there
exists a complex number \(\lambda\), with \(|\lambda|=1\), such that
\(\xi_2=\lambda\xi_1\).
\end{itshape}

If \(\xi_2=\lambda\xi_1\) with \(|\lambda|=1\), it is clear that
\(f_1=f_2\). Conversely, suppose \(f_1=f_2\). Since \(\xi_1,\xi_2\) are
cyclic by \XRef{2.3.1}{2.3.1}, there exists by \XRef{2.4.1}{2.4.1(ii)} an
automorphism \(U\) of \(H_{\pi}\), commuting with the \(\pi(x)\), such
that \(U\xi_1=\xi_2\). But \(U\) is a scalar operator \(\lambda\) by
\XRef{2.3.1}{2.3.1}, whence \(\xi_2=\lambda\xi_1\) and
\(|\lambda|=1\).

In particular, the canonical map \(P(A)\to\widehat A\) is bijective if and
only if every topologically irreducible representation of \(A\) has
dimension \(1\). When \(A\) is a \(C^{*}\)-algebra, Theorem
\XRef{2.7.3}{2.7.3} will imply that this condition holds if and only if
\(A\) is commutative.

Bibliography: \cite{ref41,ref42,ref52,ref105,ref108,ref126,ref134}.

% Source PDF physical page 048
\BookSubsection[Existence of Representations of C*-Algebras]{2.6}{Existence of Representations of \(C^{*}\)-Algebras}

\ResultHeading{2.6.1}{2.6.1. Theorem.}
\begin{itshape}
Let \(A\) be a \(C^{*}\)-algebra. There exists an isometric representation
of \(A\) in a Hilbert space.
\end{itshape}

Let \(A_h\) be the real Banach space formed by the hermitian elements of
\(A\), and let \(x\) be a nonzero element of \(A\). Then
\(x^{\ast}x\in A^{+}\setminus\{0\}\) by \XRef{1.6.1}{1.6.1}, and
\(A^{+}\) is a closed convex cone by \XRef{1.6.1}{1.6.1}. Hence there
exists a continuous real linear functional \(f_x\) on \(A_h\) such that
\(f_x(y)\geq0\) for \(y\in A^{+}\) and
\(f_x(-x^{\ast}x)<0\), by \XRef{B.5}{B 5}. Extend \(f_x\) to a hermitian
functional on \(A\). Then \(f_x\) is a positive functional on \(A\), and
\(f_x(x^{\ast}x)>0\); consequently, the representation \(\pi_x\) defined
by \(f_x\), in accordance with \XRef{2.4.4}{2.4.4}, satisfies
\(\pi_x(x)\neq0\). Let \(\pi\) be the Hilbert direct sum of the
\(\pi_x\), for \(x\in A\), \(x\neq0\). Then \(\pi\) is injective, and
hence isometric by \XRef{1.3.7}{1.3.7} and \XRef{1.8.1}{1.8.1}. This
proves the theorem.

Thus, as announced, the closed involutive subalgebras of an algebra
\(\mathcal L(H)\), with \(H\) a Hilbert space, give the most general
examples of \(C^{*}\)-algebras.

\ResultHeading{2.6.2}{2.6.2. Proposition.}
\begin{itshape}
Let \(A\) be a \(C^{*}\)-algebra and \(x\in A\). The following conditions
are equivalent:
\begin{enumerate}[label=\textup{(\roman*)}]
 \item \(x\geq0\).
 \item For every representation \(\pi\) of \(A\), the operator \(\pi(x)\)
 is nonnegative.
 \item For every positive functional \(f\) on \(A\), one has \(f(x)\geq0\).
\end{enumerate}
\end{itshape}

(i) \(\Rightarrow\) (iii) is clear.

(iii) \(\Rightarrow\) (ii): let \(\pi\) be a representation of \(A\) and
\(\xi\in H_{\pi}\). The functional
\(y\mapsto(\pi(y)\xi\mid\xi)\) on \(A\) is positive. Hence, if (iii)
holds, then \((\pi(x)\xi\mid\xi)\geq0\), so \(\pi(x)\geq0\).

(ii) \(\Rightarrow\) (i): let \(\pi\) be an isometric representation of
\(A\), furnished by \XRef{2.6.1}{2.6.1}. If \(\pi(x)\geq0\), then
\(\pi(x)\) is positive relative to \(\pi(A)\) by \XRef{1.6.5}{1.6.5};
hence \(x\) is positive relative to \(A\).

\ResultHeading{2.6.3}{2.6.3. Corollary.}
\begin{itshape}
Let \(A\) be a \(C^{*}\)-algebra, let \(A_h\) be the set of hermitian
elements of \(A\), let \(B\) be the set of positive functionals on \(A\)
of norm at most \(1\), and let \(P\subset B\) be the set of pure states of
\(A\). Give \(B\) and \(P\) the weak topology. Let \(C(B)\) and \(C(P)\)
be the sets of continuous real-valued functions on \(B\) and \(P\). For
each \(x\in A_h\), let \(F_x\) be the continuous real-valued function
\[
 f\longmapsto\langle f,x\rangle
\]
on \(B\), and let \(G_x\) be its restriction to \(P\). The map
\(x\mapsto F_x\) (respectively \(x\mapsto G_x\)) is an isometric
isomorphism of the ordered Banach space \(A_h\) onto a vector subspace of
the ordered Banach space \(C(B)\) (respectively \(C(P)\)).
\end{itshape}

It is clear that these maps are linear, and that
\[
 x\geq0\ \Longrightarrow\ F_x\geq0\ \Longrightarrow\ G_x\geq0.
\]
Conversely, suppose that \(G_x\geq0\), that is, \(f(x)\geq0\) for every
\(f\in P\). By \XRef{2.5.5}{2.5.5}, we then have \(f(x)\geq0\) for every
\(f\in B\), and therefore \(x\geq0\)
% Source PDF physical page 049
by \XRef{2.6.2}{2.6.2}. On the other hand, realize \(A\) on a Hilbert
space \(H\) by \XRef{2.6.1}{2.6.1}, and let \(x\in A_h\) satisfy
\(\lVert x\rVert=1\). There exist unit vectors \(\xi\in H\) such that
\(|(x\xi\mid\xi)|\) is arbitrarily close to \(1\), hence functionals
\(f\in B\) such that \(|f(x)|\) is arbitrarily close to \(1\), and hence,
by \XRef{2.5.5}{2.5.5}, functionals \(f\in P\) such that \(|f(x)|\) is
arbitrarily close to \(1\). Consequently,
\[
 1\leq\lVert G_x\rVert\leq\lVert F_x\rVert\leq1.
\]
The corollary follows.

When \(A\) is commutative, the map \(x\mapsto G_x\) is just the Gelfand
isomorphism, restricted to \(A_h\). Another generalization of the Gelfand
isomorphism to the noncommutative case, more satisfactory in certain
respects, will be studied later in \XRef{10.5.4}{10.5.4}.

\ResultHeading{2.6.4}{2.6.4. Corollary.}
\begin{itshape}
Let \(A\) be a \(C^{*}\)-algebra and \(g\) a continuous hermitian linear
functional on \(A\). There exist positive functionals \(f,f'\) on \(A\)
such that
\[
 g=f-f',
 \qquad
 \lVert g\rVert=\lVert f\rVert+\lVert f'\rVert.
\]
\end{itshape}

Retain the notation of \XRef{2.6.3}{2.6.3}, and identify \(A_h\) with a
vector subspace of \(C(B)\). By the Hahn--Banach theorem, \(g\) extends to
a linear functional \(\mu\) on \(C(B)\) such that
\(\lVert\mu\rVert=\lVert g\rVert\) (cf. \XRef{1.2.6}{1.2.6}). The measure
\(\mu\) on \(B\) can be written \(\mu=\mu^{+}-\mu^{-}\), where
\(\mu^{+},\mu^{-}\) are positive measures such that
\[
 \lVert\mu\rVert=\lVert\mu^{+}\rVert+\lVert\mu^{-}\rVert.
\]
Let \(f=\mu^{+}|_{A_h}\) and \(f'=\mu^{-}|_{A_h}\). Then \(g=f-f'\), and
\[
 \lVert g\rVert\leq\lVert f\rVert+\lVert f'\rVert
 \leq\lVert\mu^{+}\rVert+\lVert\mu^{-}\rVert
 =\lVert\mu\rVert=\lVert g\rVert,
\]
whence
\[
 \lVert g\rVert=\lVert f\rVert+\lVert f'\rVert.
\]

We shall see later in \XRef{12.3.4}{12.3.4} that the decomposition in
\XRef{2.6.4}{2.6.4} is unique.

Bibliography:
\cite{ref41,ref42,ref55,ref70,ref105,ref108,ref126,ref134,ref150}.

\BookSubsection[Enveloping C*-Algebra of an Involutive Banach Algebra]{2.7}{Enveloping \(C^{*}\)-Algebra of an Involutive Banach Algebra}

\ResultHeading{2.7.1}{2.7.1. Proposition.}
\begin{itshape}
Let \(A\) be an involutive Banach algebra having an approximate identity.
Let \(R\) be the set of representations of \(A\), \(R'\) the set of
topologically irreducible representations of \(A\), \(B\) the set of
continuous positive functionals on \(A\) of norm at most \(1\), and \(P\)
the set of pure states of \(A\). For every \(x\in A\),
\begin{equation}
 \sup_{\pi\in R}\lVert\pi(x)\rVert
 =\sup_{\pi\in R'}\lVert\pi(x)\rVert
 =\sup_{f\in B}f(x^{\ast}x)^{1/2}
 =\sup_{f\in P}f(x^{\ast}x)^{1/2}. \tag{1}\label{eq:2.7.1-1}
\end{equation}
Denote by \(\lVert x\rVert'\) the common value of these four expressions.
Then \(\lVert x\rVert'\leq\lVert x\rVert\), and
\(x\mapsto\lVert x\rVert'\) is a seminorm on \(A\) such that
\[
 \lVert xy\rVert'\leq\lVert x\rVert'\lVert y\rVert',
 \qquad
 \lVert x^{\ast}\rVert'=\lVert x\rVert',
 \qquad
 \lVert x^{\ast}x\rVert'=\lVert x\rVert'^{\,2}
\]
for all \(x,y\in A\).
\end{itshape}

% Source PDF physical page 050
Denote by \(a,b,c,d\), in succession, the four numbers appearing in
\eqref{eq:2.7.1-1}.

\(d\leq b\): let \(f\in P\). Then \(f\) is associated with a
representation \(\pi\in R'\) by \XRef{2.5.4}{2.5.4}, and
\[
 f(x^{\ast}x)\leq\lVert\pi(x)\rVert^{2}
\]
by \XRef{2.4.7}{2.4.7}.

\(b\leq a\) is clear.

\(a\leq c\) follows from \XRef{2.4.7}{2.4.7}.

\(c\leq d\): let \(f\in B\). By \XRef{2.5.5}{2.5.5},
\(f(x^{\ast}x)\) belongs to the closed interval generated by \(0\) and the
numbers \(g(x^{\ast}x)\), where \(g\in P\). Hence
\[
 f(x^{\ast}x)\leq\sup_{g\in P}g(x^{\ast}x).
\]

By \XRef{1.3.7}{1.3.7}, \(\lVert\pi(x)\rVert\leq\lVert x\rVert\) for
every \(\pi\in R\), so \(\lVert x\rVert'\leq\lVert x\rVert\). Moreover,
each function \(x\mapsto\lVert\pi(x)\rVert\) is a seminorm on \(A\);
hence \(x\mapsto\lVert x\rVert'\) is a seminorm on \(A\). We have
\[
 \lVert\pi(x^{\ast})\rVert=\lVert\pi(x)\rVert,
 \qquad
 \lVert\pi(x^{\ast}x)\rVert=\lVert\pi(x)\rVert^{2}
\]
for every \(\pi\in R\), and therefore
\[
 \lVert x^{\ast}\rVert'=\lVert x\rVert',
 \qquad
 \lVert x^{\ast}x\rVert'=\lVert x\rVert'^{\,2}.
\]
Finally, for \(x,y\in A\) and every \(\pi\in R\),
\[
 \lVert\pi(xy)\rVert\leq\lVert\pi(x)\rVert\lVert\pi(y)\rVert
 \leq\lVert x\rVert'\lVert y\rVert',
\]
and hence \(\lVert xy\rVert'\leq\lVert x\rVert'\lVert y\rVert'\).

\ResultHeading{2.7.2}{2.7.2.}
Let \(I\) be the set of \(x\in A\) such that \(\lVert x\rVert'=0\). It is
a closed self-adjoint two-sided ideal of \(A\). The map
\(x\mapsto\lVert x\rVert'\) induces, on passing to the quotient, a norm on
\(A/I\). With this norm, \(A/I\) satisfies all the axioms of a
\(C^{*}\)-algebra except that in general it is not complete. The completion
\(B\) of \(A/I\) is a \(C^{*}\)-algebra called the
\emph{enveloping \(C^{*}\)-algebra} of \(A\). The canonical map from \(A\)
into \(B\) is a norm-decreasing homomorphism of involutive algebras and has
dense image in \(B\).

\ResultHeading{2.7.3}{2.7.3.}
When \(A\) is a \(C^{*}\)-algebra, \(\lVert x\rVert'=\lVert x\rVert\) by
\XRef{2.6.1}{2.6.1}, and \(A\) identifies with its enveloping
\(C^{*}\)-algebra. In this case, Proposition \XRef{2.7.1}{2.7.1} gives,
as announced, the existence of sufficiently many topologically
irreducible representations, and considerably strengthens Theorem
\XRef{2.6.1}{2.6.1}:

\ResultLabel{Theorem.}
\begin{itshape}
Let \(A\) be a \(C^{*}\)-algebra. There exists a family \((\pi_i)\) of
topologically irreducible representations of \(A\) such that
\[
 \lVert x\rVert=\sup_i\lVert\pi_i(x)\rVert
\]
for every \(x\in A\).
\end{itshape}

\ResultHeading{2.7.4}{2.7.4. Proposition.}
\begin{itshape}
Let \(A\) be an involutive Banach algebra having an approximate identity,
let \(B\) be the enveloping \(C^{*}\)-algebra of \(A\), and let
\(\tau:A\to B\) be the canonical map.
\begin{enumerate}[label=\textup{(\roman*)}]
 \item If \(\pi\) is a representation of \(A\), there exists one and only
 one representation \(\rho\) of \(B\) such that
 \(\pi=\rho\circ\tau\); moreover, \(\rho(B)\) is the \(C^{*}\)-algebra
 generated by \(\pi(A)\).
 \item The map \(\pi\mapsto\rho\) is a bijection from the set of
 representations of \(A\) onto the set of representations of \(B\).
% Source PDF physical page 051
 \item The representation \(\rho\) is nondegenerate (respectively
 topologically irreducible) if and only if \(\pi\) is nondegenerate
 (respectively topologically irreducible).
\end{enumerate}
\end{itshape}

Let \(\pi\) be a representation of \(A\). With the notation of
\XRef{2.7.2}{2.7.2}, \(\pi\) vanishes on \(I\) and induces, on passing to
the quotient, a representation \(\pi'\) of \(A/I\) such that
\(\lVert\pi'(z)\rVert\leq\lVert z\rVert\) for every \(z\in A/I\), where
\(\lVert\cdot\rVert\) denotes the norm in \(B\). Hence \(\pi'\) extends to
a representation \(\rho\) of \(B\), and \(\pi=\rho\circ\tau\). The
uniqueness assertion in (i) follows from the density of \(\tau(A)\) in
\(B\). This also implies that \(\pi(A)\) is norm-dense in \(\rho(B)\).
Since \(\rho(B)\) is a \(C^{*}\)-algebra by \XRef{1.8.3}{1.8.3}, it is the
\(C^{*}\)-algebra generated by \(\pi(A)\). Assertions (ii) and (iii) are
immediate.

Proposition \XRef{2.7.4}{2.7.4} shows that \(B\) solves a universal
problem. It also shows that, in most questions concerning representations
of involutive Banach algebras having an approximate identity, one may
restrict oneself to the case of \(C^{*}\)-algebras.

\ResultHeading{2.7.5}{2.7.5. Proposition.}
\begin{itshape}
Retain the notation \(A,B,\tau\) of \XRef{2.7.4}{2.7.4}.
\begin{enumerate}[label=\textup{(\roman*)}]
 \item If \(f\) is a continuous positive functional on \(A\), there exists
 one and only one positive functional \(g\) on \(B\) such that
 \(f=g\circ\tau\). Moreover, \(\lVert g\rVert=\lVert f\rVert\).
 \item The map \(f\mapsto g\) is a bijection from the set of continuous
 positive functionals on \(A\) onto the set of positive functionals on
 \(B\).
 \item Let \(M\) be a bounded set of continuous positive functionals on
 \(A\). The map \(f\mapsto g\), restricted to \(M\), is bicontinuous for
 the weak topologies \(\sigma(A',A)\) and \(\sigma(B',B)\), where \(A'\)
 and \(B'\) denote the duals of \(A\) and \(B\).
\end{enumerate}
\end{itshape}

Let \(f\) be a continuous positive functional on \(A\). For every
\(x\in A\), by \XRef{2.1.5}{2.1.5(i)},
\[
 |f(x)|\leq\lVert f\rVert^{1/2}f(x^{\ast}x)^{1/2}
 \leq\lVert f\rVert\,\lVert x\rVert'.
\]
It follows that there exists a continuous linear functional \(g\) on \(B\)
such that \(f=g\circ\tau\) and \(\lVert g\rVert\leq\lVert f\rVert\). If
\(y\in B\), there exists a sequence \((x_n)\) in \(A\) such that
\(\tau(x_n)\to y\), and hence
\[
 g(y^{\ast}y)=\lim_n f(x_n^{\ast}x_n)\geq0;
\]
thus \(g\) is positive. If \(x\in A\) and \(\lVert x\rVert\leq1\), then
\[
 |f(x)|=|g(\tau(x))|
 \leq\lVert g\rVert\,\lVert\tau(x)\rVert
 \leq\lVert g\rVert\,\lVert x\rVert
 \leq\lVert g\rVert,
\]
so \(\lVert f\rVert\leq\lVert g\rVert\). The uniqueness of \(g\) follows
from the density of \(\tau(A)\) in \(B\). This proves (i), and (ii) is
immediate. Let \(M\subset A'\) be a set of continuous positive functionals
on \(A\), and let \(N\subset B'\) be its image under \(f\mapsto g\). It is
clear that \(f\mapsto g\) is bicontinuous from \(M\) onto \(N\) for the
weak topologies \(\sigma(A',A)\) and \(\sigma(B',\tau(A))\). If \(M\) is
bounded, then \(N\) is bounded by (i), and consequently
\(\sigma(B',\tau(A))\) and \(\sigma(B',B)\) coincide on \(N\), since
\(\tau(A)\) is dense in \(B\).

% Source PDF physical page 052
\ResultHeading{2.7.6}{2.7.6.}
Let \(A\) be an involutive Banach algebra admitting an approximate
identity, and let \(Q\) be the set of continuous positive functionals on
\(A\). For every \(f\in Q\), let \(\pi_f\) be the representation defined
by \(f\). Then
\[
 \pi=\bigoplus_{f\in Q}\pi_f
\]
is called the \emph{universal representation} of \(A\). It is
nondegenerate. By \XRef{2.2.7}{2.2.7} and \XRef{2.4.6}{2.4.6}, every
nondegenerate representation of \(A\) is a Hilbert direct sum of
representations \(\pi_f\), and hence
\(\lVert x\rVert'=\lVert\pi(x)\rVert\) for every \(x\in A\). Thus, if
\(B\) is the enveloping \(C^{*}\)-algebra of \(A\), the representation of
\(B\) corresponding to \(\pi\) is an isomorphism of \(B\) onto
\(\overline{\pi(A)}\).

Bibliography: \cite{ref27,ref42,ref52,ref105,ref108,ref126,ref134}.

\BookSubsection{2.8}{A Transitivity Theorem}

\ResultHeading{2.8.1}{2.8.1. Lemma.}
\begin{itshape}
Let \(H\) be a Hilbert space, let
\((\xi_1,\ldots,\xi_n)\) be an orthonormal system in \(H\), and let
\(\eta_1,\ldots,\eta_n\) be vectors of \(H\) of norm at most \(r\). There
exists \(b\in\mathcal L(H)\), of norm at most
\((2n)^{1/2}r\), such that
\[
 b\xi_1=\eta_1,\ldots,b\xi_n=\eta_n.
\]
If there exists a hermitian element \(h\in\mathcal L(H)\) such that
\(h\xi_1=\eta_1,\ldots,h\xi_n=\eta_n\), then \(b\) may also be chosen
hermitian.
\end{itshape}

Let \(K\) be the vector subspace of \(H\) spanned by
\(\xi_1,\ldots,\xi_n,\eta_1,\ldots,\eta_n\). We define \(b\) to leave
\(K\) invariant and to vanish on \(H\ominus K\). Thus we may reduce to the
case \(K=H\). Let
\((\xi_1,\ldots,\xi_n,\xi_{n+1},\ldots,\xi_m)\) be an orthonormal basis of
\(H\). Let \(b\) be the operator on \(H\) whose matrix relative to
\((\xi_1,\ldots,\xi_m)\) is chosen as follows:
\begin{enumerate}[label=\alph*.]
 \item its first \(n\) columns are the coordinates of
 \(\eta_1,\ldots,\eta_n\) relative to the \(\xi_i\);
 \item if \(h\) exists, its first \(n\) rows are the conjugates of its
 first \(n\) columns (this is possible because the existence of \(h\)
 implies that the square matrix formed by the first \(n\) rows and first
 \(n\) columns is hermitian);
 \item the entries not yet chosen are zero.
\end{enumerate}
Then
\[
 b\xi_1=\eta_1,\ldots,b\xi_n=\eta_n
\]
and
\[
 \lVert b\rVert^{2}\leq\operatorname{Tr}(b^{\ast}b)
 \leq2(\lVert\eta_1\rVert^{2}+\cdots+\lVert\eta_n\rVert^{2})
 \leq2nr^{2}.
\]
If \(h\) exists, \(b\) is hermitian by construction.

\ResultHeading{2.8.2}{2.8.2. Lemma.}
\begin{itshape}
Let \(H_1,\ldots,H_p\) be Hilbert spaces, let
\(H=H_1\oplus\cdots\oplus H_p\), and let \(A\) be a \(C^{*}\)-algebra of
operators on \(H\), strongly dense in the von Neumann algebra
\[
 \mathcal L(H_1)\times\cdots\times\mathcal L(H_p).
\]
Let \((\xi_1,\ldots,\xi_n)\) be an orthonormal system in \(H\), and let
\(\eta_1,\ldots,\eta_n\) be vectors of \(H\) of norm at most \(r\).
Suppose that, for \(i=1,\ldots,n\), the vectors \(\xi_i\) and \(\eta_i\)
belong to the same summand \(H_{j_i}\). There exists \(b\in A\), of norm at
most \(3n^{1/2}r\),
% Source PDF physical page 053
such that
\[
 b\xi_1=\eta_1,\ldots,b\xi_n=\eta_n.
\]
If there exists a hermitian element \(h\in\mathcal L(H)\) such that
\(h\xi_1=\eta_1,\ldots,h\xi_n=\eta_n\), then \(b\) may also be chosen
hermitian in \(A\).
\end{itshape}

By the Kaplansky density theorem, the unit ball of \(A\) (respectively the
unit ball of the hermitian part of \(A\)) is strongly dense in the unit
ball of
\(\mathcal L(H_1)\times\cdots\times\mathcal L(H_p)\) (respectively its
hermitian part). Thus there first exists
\(y_0\in\mathcal L(H_1)\times\cdots\times\mathcal L(H_p)\) such that
\[
 y_0\xi_1=\eta_1,\ldots,y_0\xi_n=\eta_n,
 \qquad
 \lVert y_0\rVert\leq(2n)^{1/2}r,
\]
by applying \XRef{2.8.1}{2.8.1} successively in
\(H_1,\ldots,H_p\); and then there exists \(x_0\in A\) such that
\[
 \lVert x_0\xi_i-y_0\xi_i\rVert\leq\frac12r
 \quad(i=1,\ldots,n),
 \qquad
 \lVert x_0\rVert\leq(2n)^{1/2}r.
\]
Next there exists
\(y_1\in\mathcal L(H_1)\times\cdots\times\mathcal L(H_p)\) such that
\[
 y_1\xi_i=\eta_i-x_0\xi_i\quad(i=1,\ldots,n),
 \qquad
 \lVert y_1\rVert\leq\frac12(2n)^{1/2}r,
\]
by \XRef{2.8.1}{2.8.1}; and then there exists \(x_1\in A\) such that
\[
 \lVert x_1\xi_i-y_1\xi_i\rVert\leq\frac14r
 \quad(i=1,\ldots,n),
 \qquad
 \lVert x_1\rVert\leq\frac12(2n)^{1/2}r.
\]
Inductively, we construct elements
\[
 y_k\in\mathcal L(H_1)\times\cdots\times\mathcal L(H_p),
 \qquad x_k\in A,
\]
such that
\begin{align*}
 y_k\xi_i&=\eta_i-x_0\xi_i-\cdots-x_{k-1}\xi_i,
 &\lVert y_k\rVert&\leq\frac1{2^k}(2n)^{1/2}r,\\
 \lVert x_k\xi_i-y_k\xi_i\rVert
 &\leq\frac1{2^{k+1}}r,
 &\lVert x_k\rVert&\leq\frac1{2^k}(2n)^{1/2}r
\end{align*}
for \(i=1,\ldots,n\). (If \(h\) exists, the \(y_k\) and \(x_k\) may be
chosen hermitian.) Then \(x_0+x_1+\cdots\) converges in norm to an element
\(b\in A\) such that
\[
 b\xi_1=\eta_1,\ldots,b\xi_n=\eta_n,
\]
and
\[
 \lVert b\rVert
 \leq\left(1+\frac12+\frac14+\cdots\right)(2n)^{1/2}r
 \leq3n^{1/2}r.
\]

\ResultHeading{2.8.3}{2.8.3. Theorem.}
\begin{itshape}
Let \(A\) be a \(C^{*}\)-algebra, let \(\widetilde A\) be the
\(C^{*}\)-algebra obtained from \(A\) by adjoining an identity element,
let \(\pi_1,\ldots,\pi_p\) be representations of \(A\) in Hilbert spaces
\(H_1,\ldots,H_p\), and let
\(\widetilde\pi_1,\ldots,\widetilde\pi_p\) be their canonical extensions
to \(\widetilde A\). Suppose that the \(\pi_j\) are nonzero,
topologically irreducible, and pairwise inequivalent.
\begin{enumerate}[label=\textup{(\roman*)}]
 \item Let \(T_1\in\mathcal L(H_1),\ldots,T_p\in\mathcal L(H_p)\), and let
 \(K_1,\ldots,K_p\) be finite-dimensional vector subspaces of
 \(H_1,\ldots,H_p\). There exists \(x\in A\) such that
\[
 \pi_j(x)|_{K_j}=T_j|_{K_j}
 \qquad(j=1,\ldots,p).
\]
% Source PDF physical page 054
 \item Let \(T_1\in\mathcal L(H_1),\ldots,T_p\in\mathcal L(H_p)\) be
 hermitian operators, and let \(K_1,\ldots,K_p\) be finite-dimensional
 vector subspaces of \(H_1,\ldots,H_p\). There exists a hermitian element
 \(x\in A\) such that
 \[
  \pi_j(x)|_{K_j}=T_j|_{K_j}
  \qquad(j=1,\ldots,p).
 \]
 \item Let \(T_1\in\mathcal L(H_1),\ldots,T_p\in\mathcal L(H_p)\) be
 unitary operators, and let \(K_1,\ldots,K_p\) be finite-dimensional
 vector subspaces of \(H_1,\ldots,H_p\). There exists a unitary element
 \(x\in\widetilde A\) such that
 \[
  \widetilde\pi_j(x)|_{K_j}=T_j|_{K_j}
  \qquad(j=1,\ldots,p).
 \]
\end{enumerate}
\end{itshape}

Let
\[
 H=H_1\oplus\cdots\oplus H_p,
 \qquad
 \pi=\pi_1\oplus\cdots\oplus\pi_p.
\]
By \XRef{1.8.3}{1.8.3}, \(\pi(A)\) is a \(C^{*}\)-subalgebra of
\(\mathcal L(H)\) commuting with the projections \(E_j=P_{H_j}\). Let
\(B\) be the von Neumann algebra generated by \(\pi(A)\). We have
\[
 \pi(A)\subset
 \mathcal L(H_1)\times\cdots\times\mathcal L(H_p),
 \qquad\text{and hence}\qquad
 B\subset
 \mathcal L(H_1)\times\cdots\times\mathcal L(H_p).
\]
Since each \(\pi_j\) is nondegenerate, \(\pi\) is nondegenerate; hence
\(\pi(A)\) is strongly dense in \(B\). Since \(\pi_j\) is topologically
irreducible, \(E_j\) is a minimal projection of
\(\pi(A)'=B'\), and the von Neumann algebra induced by \(B\) on \(H_j\)
is \(\mathcal L(H_j)\). Let \(F_j\in B\cap B'\) be the central support of
\(E_j\) (\XRef{A.10}{A 10}). Let \(j,k\) be distinct indices. If \(F_j\)
and \(F_k\) are not orthogonal, there exist nonzero projections
\(E_j',E_k'\in B'\), respectively dominated by \(E_j,E_k\), and
equivalent relative to \(B'\), by \XRef{A.44}{A 44}. Since \(E_j,E_k\)
are minimal projections of \(B'\), we have \(E_j'=E_j\) and \(E_k'=E_k\).
Hence there exists a partial isometry in \(B'\) having \(E_j,E_k\) as its
initial and final projections; but then \(\pi_j\) and \(\pi_k\) are
equivalent, a contradiction. Thus the \(F_j\) are pairwise orthogonal.
Now \(F_j\geq E_j\) for every \(j\), and consequently
\(E_j=F_j\in B\). Since \(B_{E_j}=\mathcal L(H_j)\), it follows that
\[
 B\supset
 \mathcal L(H_1)\times\cdots\times\mathcal L(H_p),
\]
and therefore, finally,
\[
 B=\mathcal L(H_1)\times\cdots\times\mathcal L(H_p).
\]

Let \(T_j\in\mathcal L(H_j)\), and let \(K_j\) be a finite-dimensional
vector subspace of \(H_j\), for \(j=1,\ldots,p\). By
\XRef{2.8.2}{2.8.2}, there exists \(x\in A\) such that
\(\pi_j(x)|_{K_j}=T_j|_{K_j}\) for all \(j\). If the \(T_j\) are
hermitian, \(x\) may be chosen so that \(\pi(x)\) is hermitian, by
\XRef{2.8.2}{2.8.2}. Since
\[
 \pi(x)=\pi\left(\frac12(x+x^{\ast})\right),
\]
\(x\) itself may be chosen hermitian.

Suppose that the \(T_j\) are unitary. For \(j=1,\ldots,p\), there is a
finite-dimensional vector subspace \(K_j'\supset K_j\) of \(H_j\), and a
unitary operator \(T_j'\in\mathcal L(H_j)\), leaving \(K_j'\) invariant,
such that \(T_j'|_{K_j}=T_j|_{K_j}\). There is then a hermitian operator
\(T_j''\in\mathcal L(H_j)\), leaving \(K_j'\) invariant, such that
\[
 \exp(iT_j'')|_{K_j'}=T_j'|_{K_j'}.
\]
By what precedes, there exists a hermitian element \(y\in\widetilde A\)
such that
\(\widetilde\pi_j(y)|_{K_j'}=T_j''|_{K_j'}\) for every \(j\). Then
\(x=\exp(iy)\) is a unitary element of \(\widetilde A\), and
\(\widetilde\pi_j(x)|_{K_j}=T_j|_{K_j}\) for every \(j\).

\ResultHeading{2.8.4}{2.8.4. Corollary.}
\begin{itshape}
Every topologically irreducible representation of a \(C^{*}\)-algebra is
algebraically irreducible.
\end{itshape}

% Source PDF physical page 055
It is enough to apply Theorem \XRef{2.8.3}{2.8.3} with \(p=1\) and
\(\dim K_1=1\).

Henceforth we shall simply speak of irreducible representations of a
\(C^{*}\)-algebra, without further qualification.

\ResultHeading{2.8.5}{2.8.5. Corollary.}
\begin{itshape}
Let \(A\) be a \(C^{*}\)-algebra, let \(f\) be a pure positive functional
on \(A\), and let \(N\) be the left ideal formed by the \(x\in A\) such
that \(f(x^{\ast}x)=0\). Then \(A/N\), equipped with the inner product
induced by \(f(y^{\ast}x)\), is complete, and hence is equal to the space
of the representation defined by \(f\).
\end{itshape}

Indeed, let \(\pi\) be this representation; it is topologically
irreducible by \XRef{2.5.4}{2.5.4}. By construction of \(\pi\), \(A/N\)
is a vector subspace of \(H_{\pi}\) invariant under \(\pi\). But
\(A/N\neq0\), and therefore \(A/N=H_{\pi}\) by \XRef{2.8.4}{2.8.4}.

\ResultHeading{2.8.6}{2.8.6. Corollary.}
\begin{itshape}
Let \(A\) be a \(C^{*}\)-algebra, let \(\widetilde A\) be the
\(C^{*}\)-algebra obtained from \(A\) by adjoining an identity element,
and let \(f_1,f_2\) be two pure states of \(A\). The representations
\(\pi_1,\pi_2\) defined by \(f_1,f_2\) are equivalent if and only if there
exists a unitary element \(u\in\widetilde A\) such that
\[
 f_2(x)=f_1(u^{\ast}xu)
 \qquad\text{for every }x\in A.
\]
\end{itshape}

Let \(\xi_1\) be the vector of \(H_{\pi_1}\) defined by \(f_1\). The states
of \(A\) that define representations equivalent to \(\pi_1\) are the
states associated with \(\pi_1\), by \XRef{2.4.6}{2.4.6}; that is, the
functionals
\[
 x\longmapsto(\pi_1(x)\xi\mid\xi),
\]
where \(\xi\) is a unit vector of \(H_{\pi_1}\), by
\XRef{2.4.3}{2.4.3}. But the unit vectors of \(H_{\pi_1}\) are precisely
the vectors \(\widetilde\pi_1(u)\xi_1\), where \(u\) is unitary in
\(\widetilde A\), by \XRef{2.8.3}{2.8.3}; here
\(\widetilde\pi_1\) denotes the canonical extension of \(\pi_1\) to
\(\widetilde A\). Finally,
\begin{align*}
 (\pi_1(x)\widetilde\pi_1(u)\xi_1
 \mid\widetilde\pi_1(u)\xi_1)
 &=(\widetilde\pi_1(u^{\ast}xu)\xi_1\mid\xi_1)\\
 &=f_1(u^{\ast}xu).
\end{align*}

Bibliography: \cite{ref49,ref69,ref126}.

\BookSubsection[Ideals in C*-Algebras]{2.9}{Ideals in \(C^{*}\)-Algebras}

\ResultHeading{2.9.1}{2.9.1. Proposition.}
\begin{itshape}
Let \(A\) be a \(C^{*}\)-algebra, let \(f\) be a positive functional on
\(A\), let \(M\) be its kernel, and let \(N\) be the set of \(x\in A\)
such that \(f(x^{\ast}x)=0\). Then
\[
 M\supset N+N^{\ast},
\]
and \(M=N+N^{\ast}\) if \(f\) is pure.
\end{itshape}

Since
\[
 |f(x)|^{2}\leq\lVert f\rVert f(x^{\ast}x),
\]
we have \(N\subset M\), and hence
\(N^{\ast}\subset M^{\ast}=M\). Thus \(N+N^{\ast}\subset M\).

Now suppose \(f\) is pure. Let \(\widetilde A\) be the
\(C^{*}\)-algebra obtained from \(A\) by adjoining an identity, and let
\(\widetilde f\) be the canonical extension of \(f\) to \(\widetilde A\).
Let \(\pi,\xi\) be the irreducible representation and vector defined by
\(f\). Let \(b\in M\), and let \(\eta\) be its canonical image in
\(H_{\pi}\). By construction of \(\pi,\xi\),
\[
 (\eta\mid\xi)=\widetilde f(1^{\ast}b)=0.
\]
Thus there exists a hermitian operator on \(H_{\pi}\) that annihilates
\(\xi\) and fixes \(\eta\). By \XRef{2.8.3}{2.8.3}, there exists a
hermitian element \(a\in A\) such that
\(\pi(a)\xi=0\) and \(\pi(a)\eta=\eta\); in other words,
\(a\in N\) and \(b-ab\in N\). Then
\[
 b^{\ast}=(b-ab)^{\ast}+b^{\ast}a\in N^{\ast}+N.
\]
Consequently \(M=M^{\ast}\subset N^{\ast}+N\).

% Source PDF physical page 056
\ResultHeading{2.9.2}{2.9.2. Lemma.}
\begin{itshape}
Let \(A\) be a unital \(C^{*}\)-algebra, let \(L\) be a closed left ideal
of \(A\), and let \(x\in A^{+}\). If, for every \(\varepsilon>0\), there
exists a positive element \(u_{\varepsilon}\in L\) such that
\[
 x\leq u_{\varepsilon}+\varepsilon1,
\]
then \(x\in L\).
\end{itshape}

Put \(t_{\varepsilon}=u_{\varepsilon}^{1/2}\). This is a limit of
polynomials in \(u_{\varepsilon}\) without constant terms, and hence
belongs to \(L\). We have
\begin{align*}
 &\left\|x^{1/2}(t_{\varepsilon}+\sqrt{\varepsilon})^{-1}
 t_{\varepsilon}-x^{1/2}\right\|^{2}\\
 &\quad=
 \left\|x^{1/2}\sqrt{\varepsilon}
 (t_{\varepsilon}+\sqrt{\varepsilon})^{-1}\right\|^{2}\\
 &\quad=
 \varepsilon\left\|
 (t_{\varepsilon}+\sqrt{\varepsilon})^{-1}
 x(t_{\varepsilon}+\sqrt{\varepsilon})^{-1}\right\|.
\end{align*}
Since \(0\leq x\leq t_{\varepsilon}^{2}+\varepsilon\), we have, by
\XRef{1.6.8}{1.6.8},
\begin{align*}
 0&\leq
 (t_{\varepsilon}+\sqrt{\varepsilon})^{-1}
 x(t_{\varepsilon}+\sqrt{\varepsilon})^{-1}\\
 &\leq
 (t_{\varepsilon}+\sqrt{\varepsilon})^{-1}
 (t_{\varepsilon}^{2}+\varepsilon)
 (t_{\varepsilon}+\sqrt{\varepsilon})^{-1}\\
 &=(t_{\varepsilon}^{2}+\varepsilon)
 (t_{\varepsilon}^{2}+2\sqrt{\varepsilon}\,t_{\varepsilon}
   +\varepsilon)^{-1}\leq1.
\end{align*}
Therefore
\[
 \left\|x^{1/2}(t_{\varepsilon}+\sqrt{\varepsilon})^{-1}
 t_{\varepsilon}-x^{1/2}\right\|^{2}\leq\varepsilon.
\]
But
\[
 x^{1/2}(t_{\varepsilon}+\sqrt{\varepsilon})^{-1}
 t_{\varepsilon}\in L.
\]
Since \(L\) is closed, \(x^{1/2}\in L\), and hence \(x\in L\).

\ResultHeading{2.9.3}{2.9.3. Lemma.}
\begin{itshape}
Let \(A\) be a \(C^{*}\)-algebra and \(L\) a closed left ideal of \(A\).
Then \(L\) is the closed left ideal generated by \(L\cap A^{+}\).
\end{itshape}

This follows at once from \XRef{1.7.3}{1.7.3}.

\ResultHeading{2.9.4}{2.9.4. Lemma.}
\begin{itshape}
Let \(A\) be a \(C^{*}\)-algebra, and let \(L,L'\) be two closed left
ideals of \(A\) such that \(L\subset L'\). Suppose that every positive
functional on \(A\) that vanishes on \(L\) also vanishes on \(L'\). Then
\(L=L'\).
\end{itshape}

We may suppose that \(A\) has an identity element. Let
\(a\in L'\cap A^{+}\), and let \(\varepsilon>0\). The set
\(S_{\varepsilon}\) of positive functionals \(f\) on \(A\) such that
\[
 f(1)=\lVert f\rVert=1,
 \qquad
 f(a)\geq\varepsilon
\]
is weakly compact. If \(f\in S_{\varepsilon}\), then \(f\) is not
identically zero on \(L'\), and hence not on \(L\); thus there exists
\(x_f\in L\) such that \(f(x_f)\neq0\). Consequently there exists a weak
neighborhood \(U_f\) of \(f\) in \(S_{\varepsilon}\) such that
\(g(x_f)\neq0\) for \(g\in U_f\). By compactness of
\(S_{\varepsilon}\), there is a finite open cover
\((U_i)_{1\leq i\leq m}\) of \(S_{\varepsilon}\) and elements
\(a_1,\ldots,a_m\in L\) such that
\[
 0<|f(a_i)|^{2}\leq f(a_i^{\ast}a_i)
 \qquad(f\in U_i).
\]
Thus
\[
 f(a_1^{\ast}a_1+\cdots+a_m^{\ast}a_m)>0
 \qquad(f\in S_{\varepsilon}).
\]
Multiplying the \(a_i\) by sufficiently large scalars, we obtain
\[
 f(a_1^{\ast}a_1+\cdots+a_m^{\ast}a_m)\geq f(a)
 \qquad(f\in S_{\varepsilon}).
\]
Hence
\[
 f(a_1^{\ast}a_1+\cdots+a_m^{\ast}a_m+\varepsilon1-a)\geq0
\]
for every positive functional \(f\). It follows from
\XRef{2.6.2}{2.6.2} that
\[
 a\leq a_1^{\ast}a_1+\cdots+a_m^{\ast}a_m+\varepsilon1.
\]
Therefore \(a\in L\) by \XRef{2.9.2}{2.9.2}. Thus \(L\) and \(L'\) have
the same positive elements, and hence are equal by \XRef{2.9.3}{2.9.3}.

% Source PDF physical page 057
\ResultHeading{2.9.5}{2.9.5. Theorem.}
\begin{itshape}
Let \(A\) be a \(C^{*}\)-algebra.
\begin{enumerate}[label=\textup{(\roman*)}]
 \item Let \(f\) be a positive functional on \(A\), and let \(N_f\) be
 the left ideal formed by the \(x\in A\) such that \(f(x^{\ast}x)=0\).
 Then \(N_f\) is regular maximal if and only if \(f\) is pure.
 \item The map \(f\mapsto N_f\) defines a bijection from the set of pure
 states onto the set of regular maximal left ideals.
 \item Every closed left ideal of \(A\) is the intersection of the regular
 maximal left ideals that contain it.
\end{enumerate}
\end{itshape}

\emph{a.} Let \(f\) be a pure positive functional on \(A\). By
\XRef{2.8.4}{2.8.4} and \XRef{2.8.5}{2.8.5}, the canonical representation
\(\pi\) of \(A\) in \(A/N_f\) is nonzero and algebraically irreducible.
There exists \(\xi\in H_{\pi}=A/N_f\) such that \(N_f\) is the set of all
\(x\in A\) for which \(\pi(x)\xi=0\). Hence \(N_f\) is a regular maximal
left ideal of \(A\).

\emph{b.} Let \(L\) be a closed left ideal of \(A\). Let \(S\) be the set
of positive functionals on \(A\), of norm at most \(1\), that vanish on
\(L\). If \(f\in S\), then \(N_f\supset L\), because
\(x^{\ast}x\in L\) for every \(x\in L\), and hence
\(f(x^{\ast}x)=0\). By \XRef{2.9.4}{2.9.4},
\[
 L=\bigcap_{f\in S}N_f.
\]
Moreover, \(S\) is convex and weakly compact, and therefore is the weakly
closed convex hull of the set \(S'\) of its extreme points. Consequently,
if \(x\in A\) is such that \(f(x^{\ast}x)=0\) for every \(f\in S'\), then
the same holds for every \(f\in S\). This proves that
\[
 \bigcap_{f\in S}N_f=\bigcap_{f\in S'}N_f.
\]
Finally, if \(f\in S\) can be written \(f=f_1+f_2\), with \(f_1,f_2\)
positive, then, for every \(x\in L\),
\[
 0\leq f_1(x^{\ast}x)+f_2(x^{\ast}x)=f(x^{\ast}x)=0.
\]
Thus \(f_1(x^{\ast}x)=f_2(x^{\ast}x)=0\), and consequently
\(f_1(x)=f_2(x)=0\); hence \(f_1,f_2\in S\). This proves that the nonzero
elements of \(S'\) are pure functionals. In view of part \emph{a}, the
equality
\[
 L=\bigcap_{f\in S'}N_f
\]
proves (iii).

\emph{c.} At the same time, we see that if \(L\) is regular maximal, and
hence closed by \XRef{B.1}{B 1}, then \(L=N_f\) for some pure \(f\).
Thus the map considered in (ii) is surjective. If \(f,f'\) are pure states
such that \(N_f=N_{f'}\), Proposition \XRef{2.9.1}{2.9.1} shows that \(f\)
and \(f'\) have the same kernel. Hence \(f=\lambda f'\), with
\(\lambda\geq0\). Since \(\lVert f\rVert=\lVert f'\rVert=1\), we have
\(f=f'\), which completes the proof of (ii).

\emph{d.} Finally, let \(f\) be a positive functional such that \(N_f\) is
regular maximal. Then \(N_f=N_{f'}\), where \(f'\) is pure. The kernel of
\(f'\), which is
\[
 N_{f'}+N_{f'}^{\ast}=N_f+N_f^{\ast},
\]
is contained in the kernel of \(f\), by \XRef{2.9.1}{2.9.1}. Therefore
\(f\) and \(f'\) are proportional, so \(f\) is pure. Together with part
\emph{a}, this proves (i).

\ResultHeading{2.9.6}{2.9.6.}
Recall that two representations \(\pi_1,\pi_2\) of an algebra on vector
spaces \(E_1,E_2\) are said to be \emph{algebraically equivalent} if there
exists an isomorphism from \(E_1\) onto \(E_2\) that carries \(\pi_1\) into
\(\pi_2\).

% Source PDF physical page 058
\ResultLabel{Corollary.}
\begin{itshape}
Let \(A\) be a \(C^{*}\)-algebra.
\begin{enumerate}[label=\textup{(\roman*)}]
 \item Every algebraically irreducible representation of the
 noninvolutive algebra \(A\) on a complex vector space is algebraically
 equivalent to a representation of the \(C^{*}\)-algebra \(A\) on a
 Hilbert space.
 \item Let \(\pi,\pi'\) be two irreducible representations of the
 \(C^{*}\)-algebra \(A\) on Hilbert spaces \(H,H'\). If \(\pi,\pi'\) are
 algebraically equivalent, then they are equivalent in the sense of
 \XRef{2.2.1}{2.2.1}.
\end{enumerate}
\end{itshape}

Let \(\pi\) be a nonzero algebraically irreducible representation of \(A\)
on a complex vector space. There exists a regular maximal left ideal \(L\)
of \(A\) such that \(\pi\) identifies with the regular representation of
\(A\) on \(A/L\). By \XRef{2.9.5}{2.9.5}, there is then a pure positive
functional \(f\) on \(A\) such that \(L=N_f\). The space \(A/L\) can
therefore be given a Hilbert-space structure for which \(\pi\) is a
representation of the \(C^{*}\)-algebra \(A\). This proves (i).

Let \(\pi,\pi'\) be two nonzero irreducible representations of the
\(C^{*}\)-algebra \(A\) on Hilbert spaces \(H,H'\). Let \(\xi\)
(respectively \(\xi'\)) be a unit vector of \(H\) (respectively \(H'\)),
and let \(L\) (respectively \(L'\)) be the set of \(x\in A\) such that
\(\pi(x)\xi=0\) (respectively \(\pi'(x)\xi'=0\)). Put
\[
 f(x)=(\pi(x)\xi\mid\xi),
 \qquad
 f'(x)=(\pi'(x)\xi'\mid\xi').
\]
Then \(L\) (respectively \(L'\)) is the set of \(x\in A\) such that
\(f(x^{\ast}x)=0\) (respectively \(f'(x^{\ast}x)=0\)). If \(\pi,\pi'\)
are algebraically equivalent, \(\xi,\xi'\) may be chosen so that
\(L=L'\). Then \(f=f'\) by \XRef{2.9.5}{2.9.5}, and hence
\(\pi\sim\pi'\) by \XRef{2.4.1}{2.4.1}. This proves (ii).

We may therefore henceforth speak of classes of irreducible
representations of a \(C^{*}\)-algebra without specifying whether the
equivalence is purely algebraic or unitary.

\ResultHeading{2.9.7}{2.9.7.}
Recall that a \emph{primitive two-sided ideal} of an algebra \(A\) is the
kernel of a nonzero algebraically irreducible representation of \(A\) on
a vector space. We denote by \(\operatorname{Prim}(A)\) the set of
primitive two-sided ideals of \(A\).

\ResultLabel{Theorem.}
\begin{itshape}
Let \(A\) be a \(C^{*}\)-algebra.
\begin{enumerate}[label=\textup{(\roman*)}]
 \item The primitive two-sided ideals of \(A\) are the kernels of the
 nonzero topologically irreducible representations of \(A\) on a Hilbert
 space.
 \item Every closed two-sided ideal of \(A\) is the intersection of the
 primitive two-sided ideals that contain it.
\end{enumerate}
\end{itshape}

Assertion (i) follows from \XRef{2.8.4}{2.8.4} and
\XRef{2.9.6}{2.9.6(i)}.

Let \(I\) be a closed two-sided ideal of \(A\). Then \(A/I\) is a
\(C^{*}\)-algebra by \XRef{1.8.2}{1.8.2}. Hence there exists a family
\((\pi_i)\) of irreducible representations of \(A\), vanishing on \(I\),
whose kernels have intersection \(I\), by \XRef{2.7.3}{2.7.3}. But
\(\ker\pi_i\) is primitive for every \(i\).

% Source PDF physical page 059
Theorem \XRef{2.9.7}{2.9.7} defines a canonical map
\[
 \widehat A\longrightarrow\operatorname{Prim}(A)
\]
which is surjective. We shall see in \XRef{4.3.7}{4.3.7} very broad
classes of cases in which this map is bijective.

Bibliography: \cite{ref36,ref69,ref74,ref126,ref135}.

\BookSubsection[Extension of Representations of C*-Algebras]{2.10}{Extension of Representations of \(C^{*}\)-Algebras}

\ResultHeading{2.10.1}{2.10.1. Lemma.}
\begin{itshape}
Let \(A\) be a unital \(C^{*}\)-algebra, and let \(B\) be a self-adjoint
vector subspace of \(A\) such that \(1\in B\). Let \(F\) be the set of
linear functionals \(g\) on \(B\) such that
\[
 g(x^{\ast})=\overline{g(x)}\quad(x\in B),\qquad
 g(x)\geq0\quad(x\in B\cap A^{+}),\qquad
 g(1)=1.
\]
\begin{enumerate}[label=\textup{(\roman*)}]
 \item Every element of \(F\) extends to a state of \(A\).
 \item Every extreme element of \(F\) extends to a pure state of \(A\).
 \item Let \(g\) be an extreme element of \(F\). If there exists only one
 pure state \(f\) of \(A\) extending \(g\), then \(f\) is the only state
 of \(A\) extending \(g\), and, for every hermitian element \(x\in A\),
 \[
  f(x)=
  \sup_{\substack{y=y^{\ast}\in B\\y\leq x}}g(y).
 \]
\end{enumerate}
\end{itshape}

Let \(A_h\) (respectively \(B_h\)) be the set of hermitian elements of
\(A\) (respectively \(B\)). Let \(g\in F\), and let \(g'\) be its
restriction to \(B_h\); it is real and nonnegative on \(B\cap A^{+}\).
The element \(1\) is an interior point, relative to \(B_h\), of the
positive cone. Hence, by \XRef{B.6}{B 6}, \(g'\) extends to a real linear
functional \(f'\) on \(A_h\) that is nonnegative on \(A^{+}\). Next
\(f'\) extends to a hermitian linear functional \(f\) on \(A\), and \(f\)
is a state of \(A\). The functionals \(g\) and \(f|_B\) coincide on
\(B_h\), and hence on \(B_h+iB_h=B\). This proves (i).

Let \(g\) be an extreme element of \(F\), and let \(K\) be the set of
states of \(A\) that extend \(g\). By (i), \(K\) is nonempty, and it is
clearly convex and weakly compact. Choose an extreme point \(f\) of \(K\).
We shall show that \(f\) is a pure state. Suppose \(f_1,f_2\) are states
of \(A\) such that
\[
 f=\frac12(f_1+f_2).
\]
Then
\[
 g=\frac12\bigl((f_1|_B)+(f_2|_B)\bigr).
\]
Since \(g\) is extreme in \(F\), \(f_1|_B=f_2|_B=g\), so
\(f_1,f_2\in K\). Since \(f\) is extreme in \(K\), \(f=f_1=f_2\). This
proves (ii).

Let \(g\) be an extreme element of \(F\), and suppose there exists only one
pure state \(f\) of \(A\) extending \(g\). By what precedes, \(K\) has
only the one extreme point \(f\), and hence \(K=\{f\}\). Let
\(x\in A_h\), \(x\notin B\), and put
\[
 \alpha=\sup_{\substack{y\in B_h\\y\leq x}}g(y).
\]
There exists one and only one linear functional \(g_1\) on
\(B+\mathbb Cx\) extending \(g\) and satisfying \(g_1(x)=\alpha\). It is
clear that
\[
 g_1(y^{\ast})=\overline{g_1(y)}
 \qquad(y\in B+\mathbb Cx).
\]
If \(y+\lambda x\geq0\), where \(y\in B_h\) and
\(\lambda\in\mathbb R\), then \(g_1(y+\lambda x)\geq0\) (this is evident
if \(\lambda=0\);
% ===== ASSEMBLED TRANSLATION CHUNK 03: chunk_060_089.tex =====
% Source PDF physical page 060
if $\lambda>0$, then $-\lambda^{-1}y\leq x$, hence
$g(-\lambda^{-1}y)\leq\alpha$, that is, $\lambda\alpha+g(y)\geq0$;
if $\lambda<0$, then $-\lambda^{-1}y\geq x$, hence
$g(-\lambda^{-1}y)\geq\alpha$, that is, $\lambda\alpha+g(y)\geq0$.
By what precedes, $g_1$ extends to a state of $A$, which necessarily
coincides with $f$. Thus $f(x)=g_1(x)=\alpha$.

\ResultHeading{2.10.2}{2.10.2. Proposition.}
\begin{itshape}
Let $A$ be a $C^{*}$-algebra, $B$ a $C^{*}$-subalgebra of $A$, and
$\rho$ a representation of $B$ on a Hilbert space $K$. There exist a Hilbert
space $H$ containing $K$ as a subspace (with the norm of $H$ inducing that of
$K$), and a representation $\pi$ of $A$ on $H$, such that
$\rho(x)=\pi(x)|_K$ for every $x\in B$. If $\rho$ is irreducible, $\pi$ may be
chosen irreducible.
\end{itshape}

Let $\widetilde A$ be the $C^{*}$-algebra obtained from $A$ by adjoining an
identity element. Let $\widetilde B=B+\mathbb C1$, which is a
$C^{*}$-subalgebra of $\widetilde A$. The representation $\rho$ extends uniquely
to a representation $\widetilde\rho$ of $\widetilde B$ such that
$\widetilde\rho(1)=1$, and $\widetilde\rho$ is irreducible if and only if
$\rho$ is irreducible. We are therefore reduced to the case in which $A$ has an
identity element and $1\in B$. Moreover, by \XRef{2.2.6}{2.2.6} and
\XRef{2.2.7}{2.2.7}, we may suppose that $\rho$ has a cyclic unit vector $\xi$.
Set
\[
g(x)=(\rho(x)\xi\mid\xi)\qquad(x\in B).
\]
Then $g$ is a state of $B$, pure if $\rho$ is irreducible. Let $f$ be a state
of $A$ extending $g$, pure if $\rho$ is irreducible
(\XRef{2.10.1}{2.10.1}). Let $\pi=\pi_f$ and $\eta=\xi_f$, and put
$H_0=\overline{\pi(B)\eta}\subset H_\pi$. For each $x\in B$, let $\pi'(x)$ be
the restriction of $\pi(x)$ to $H_0$. Then $\pi'$ is a representation of $B$
on $H_0$, $\eta$ is cyclic for $\pi'$, and, for every $x\in B$,
\[
(\pi'(x)\eta\mid\eta)=f(x)=g(x)=(\rho(x)\xi\mid\xi).
\]
Consequently there is an isomorphism from $K$ onto $H_0$ which carries $\xi$
to $\eta$ and $\rho$ to $\pi'$ (\XRef{2.4.1}{2.4.1}). Thus we may identify
$K$ with $H_0$, and $\rho$ with $\pi'$. Finally, if $\rho$ is irreducible,
$f$ is a pure state by construction, and hence $\pi$ is irreducible.

\ResultHeading{2.10.3}{2.10.3. Lemma.}
\begin{itshape}
Let $A$ be a $C^{*}$-algebra, $I$ a closed two-sided ideal of $A$, $\pi$ a
representation of $A$ on $H$, and $H'$ a closed vector subspace of $H$
invariant under $\pi(I)$; for each $x\in I$, let
$\rho(x)=\pi(x)|_{H'}$. Suppose that the representation $\rho$ of $I$ on
$H'$ is nondegenerate.
\begin{enumerate}[label=\textup{(\roman*)}]
\item $H'$ is invariant under $\pi(A)$; let $\pi'$ be the subrepresentation of
$\pi$ defined by $H'$.
\item $\pi'(I)$ is strongly dense in $\pi'(A)$.
\end{enumerate}
\end{itshape}

Let $(u_\lambda)$ be an approximate identity of the $C^{*}$-algebra $I$.
Since $\rho$ is nondegenerate, $\rho(u_\lambda)$ converges strongly to $1$
(\XRef{2.2.10}{2.2.10}). If $x\in A$, then $xu_\lambda\in I$; hence, for
every $\xi\in H'$,
\[
\pi(x)\rho(u_\lambda)\xi
=\pi(x)\pi(u_\lambda)\xi
=\pi(xu_\lambda)\xi
=\rho(xu_\lambda)\xi\in H'.
\]
Passing to the limit gives $\pi(x)\xi\in H'$. This proves (i). Moreover,
$\pi'(x)$ is the strong limit of $\pi'(xu_\lambda)$, which proves (ii).

% Source PDF physical page 061
\ResultHeading{2.10.4}{2.10.4. Proposition.}
\begin{itshape}
Let $A$ be a $C^{*}$-algebra, $I$ a closed two-sided ideal of $A$, and $\rho$
a nondegenerate representation of $I$ on $H$.
\begin{enumerate}[label=\textup{(\roman*)}]
\item There exists a unique representation $\pi$ of $A$ on $H$ extending
$\rho$.
\item $\rho(I)$ is strongly dense in $\pi(A)$.
\end{enumerate}
\end{itshape}

By \XRef{2.10.2}{2.10.2}, there exist a Hilbert space $H'\supset H$ and a
representation $\pi_1$ of $A$ on $H'$ such that
$\rho(x)=\pi_1(x)|_H$ for every $x\in I$. By \XRef{2.10.3}{2.10.3}, $H$ is
invariant under $\pi_1(A)$. Taking for $\pi$ the subrepresentation of $\pi_1$
defined by $H$ gives the existence assertion in (i), and (ii) follows from
\XRef{2.10.3}{2.10.3}(ii). Finally, let $\nu$ be another representation of
$A$ on $H$ extending $\rho$. If $(u_\lambda)$ is an approximate identity of
$I$, then for every $x\in A$, $\rho(xu_\lambda)$ converges strongly both to
$\pi(x)$ and to $\nu(x)$; hence $\pi(x)=\nu(x)$.

Bibliography: \cite{ref35,ref45,ref108,ref134}.

\BookSubsection{2.11}{Passage to an Ideal and to a Quotient Algebra}

Let $A$ be a $C^{*}$-algebra. As we shall see, the choice of a closed
two-sided ideal of $A$ determines a decomposition of the sets
$P(A)$, $\widehat A$, and $\operatorname{Prim}(A)$.

\ResultHeading{2.11.1}{2.11.1. Lemma.}
\begin{itshape}
Let $A$ be a $C^{*}$-algebra, $I$ a closed two-sided ideal of $A$, $\pi$ a
representation of $A$ on a Hilbert space $H$, $K$ the essential subspace of
$\pi|_I$, and $B$ the strong closure of $\pi(A)$.
\begin{enumerate}[label=\textup{(\roman*)}]
\item $P_K\in B\cap B'$.
\item The subrepresentation of $\pi$ defined by $H\ominus K$ vanishes on $I$.
\item If $\pi'$ denotes the subrepresentation of $\pi$ defined by $K$, then
$\pi'(I)$ is strongly dense in $\pi'(A)$.
\end{enumerate}
\end{itshape}

The operator $P_K$ is in the strong closure of $\pi(I)$, and hence
$P_K\in B$. By \XRef{2.10.3}{2.10.3}, $K$ is invariant under $\pi(A)$, and
hence $P_K\in B'$. Assertion (ii) is evident, and (iii) follows from
\XRef{2.10.3}{2.10.3}.

\ResultHeading{2.11.2}{2.11.2. Proposition.}
\begin{itshape}
Let $A$ be a $C^{*}$-algebra, $I$ a closed two-sided ideal of $A$, and let
$\widehat A_I$ (respectively $\widehat A^{I}$) be the set of
$\pi\in\widehat A$ such that $\pi(I)=0$ (respectively $\pi(I)\ne0$).
\begin{enumerate}[label=\textup{(\roman*)}]
\item For every $\pi\in\widehat A_I$, let $\pi'$ be the representation of
$A/I$ induced from $\pi$ by passage to the quotient. Then
$\pi\mapsto\pi'$ is a bijection from $\widehat A_I$ onto $\widehat{A/I}$.
\item The map $\pi\mapsto\pi|_I$ is a bijection from $\widehat A^{I}$ onto
$\widehat I$.
\end{enumerate}
\end{itshape}

Assertion (i) is evident. We prove (ii). Let $\pi\in\widehat A^{I}$ and adopt
the notation of \XRef{2.11.1}{2.11.1}. Since $K\ne\{0\}$, irreducibility of
$\pi$ gives $K=H$. Then \XRef{2.11.1}{2.11.1}(iii) proves that $\pi|_I$ is
irreducible (see also \XRef{2.11.3}{2.11.3}). The map
$\pi\mapsto\pi|_I$ from $\widehat A^{I}$ to $\widehat I$ is bijective by
\XRef{2.10.4}{2.10.4}.

% Source PDF physical page 062
\ResultHeading{2.11.3}{2.11.3. Lemma.}
\begin{itshape}
Let $A$ be an algebra over a commutative field, and let $\pi$ be an
irreducible representation of $A$ on a vector space $E$.
\begin{enumerate}[label=\textup{(\roman*)}]
\item Let $I$ be a two-sided ideal of $A$. If $\pi(I)\ne0$, then $\pi|_I$ is
irreducible.
\item Let $I_1,I_2$ be two two-sided ideals of $A$ such that
$\pi(I_1)\ne0$ and $\pi(I_2)\ne0$. Then $\pi(I_1I_2)\ne0$.
\end{enumerate}
\end{itshape}

The set of elements of $E$ annihilated by $\pi(I)$ is invariant under
$\pi(A)$ and is distinct from $E$, and is therefore zero. Thus, if $\xi$ is
a nonzero element of $E$, then $\pi(I)\xi\ne0$. Since $\pi(I)\xi$ is invariant
under $\pi(A)$, we have $\pi(I)\xi=E$, which proves (i). Moreover, what
precedes shows that $\pi(I_2)E=E$, and hence
$\pi(I_1)\pi(I_2)E=E$; therefore $\pi(I_1I_2)\ne0$.

\ResultHeading{2.11.4}{2.11.4. Lemma.}
\begin{itshape}
Let $I_1,I_2$ be two two-sided ideals of the algebra $A$, and let $I$
be a primitive ideal of $A$. If $I\supset I_1I_2$ (in particular, if
$I\supset I_1\cap I_2$), then $I\supset I_1$ or $I\supset I_2$.
\end{itshape}

Suppose that $I\not\supset I_1$ and $I\not\supset I_2$. Applying
\XRef{2.11.3}{2.11.3}(ii) to an irreducible representation $\pi$ of $A$ with
kernel $I$, we obtain $\pi(I_1I_2)\ne0$, and hence
$I\not\supset I_1I_2$.

\ResultHeading{2.11.5}{2.11.5. Proposition.}
\begin{itshape}
Let $A$ be a $C^{*}$-algebra, $I$ a closed two-sided ideal of $A$, and let
$\operatorname{Prim}_I(A)$ (respectively $\operatorname{Prim}^{I}(A)$) be the
set of primitive two-sided ideals of $A$ that contain $I$ (respectively do not
contain $I$).
\begin{enumerate}[label=\textup{(\roman*)}]
\item The map $J\mapsto J/I$ is a bijection from $\operatorname{Prim}_I(A)$
onto $\operatorname{Prim}(A/I)$.
\item The map $J\mapsto J\cap I$ is a bijection from
$\operatorname{Prim}^{I}(A)$ onto $\operatorname{Prim}(I)$.
\end{enumerate}
\end{itshape}

Assertion (i) is evident. We prove (ii). If
$J\in\operatorname{Prim}^{I}(A)$, then $J$ is the kernel of an irreducible
representation $\pi$ of $A$ such that $\pi(I)\ne0$; hence
$\pi|_I\in\widehat I$ (\XRef{2.11.2}{2.11.2}) and
$J\cap I=\ker(\pi|_I)\in\operatorname{Prim}(I)$. If
$J'\in\operatorname{Prim}(I)$, there exists $\pi'\in\widehat I$ such that
$J'=\ker\pi'$, and then there exists $\pi\in\widehat A$ such that
$\pi'=\pi|_I$ (\XRef{2.10.4}{2.10.4}); thus
$J'=(\ker\pi)\cap I$, and the map in (ii) is surjective. Let
$J_1,J_2\in\operatorname{Prim}^{I}(A)$ and suppose that
$J_1\cap I=J_2\cap I$. We have $J_2\supset J_1\cap I$ and
$J_2\not\supset I$, and hence $J_2\supset J_1$ by
\XRef{2.11.4}{2.11.4}. Similarly $J_1\supset J_2$, so $J_1=J_2$: the map
in (ii) is injective.

\ResultHeading{2.11.6}{2.11.6.}
Before establishing for pure states a result analogous to Propositions
\XRef{2.11.2}{2.11.2} and \XRef{2.11.5}{2.11.5}, we first establish some
properties of arbitrary positive forms.

Let $A$ be a $C^{*}$-algebra, $I$ a closed two-sided ideal of $A$, and
$\omega\colon A\to A/I$ the canonical homomorphism. Since
$\omega(A^+)=(A/I)^+$, the map $k\mapsto k\circ\omega$ is a bijection from
the set of positive forms on $A/I$ onto the set of positive forms on $A$
that vanish on $I$. By the general properties of Banach spaces, this
bijection preserves norms and is bicontinuous for the weak topologies.

% Source PDF physical page 063
\ResultHeading{2.11.7}{2.11.7. Proposition.}
\begin{itshape}
Let $A$ be a $C^{*}$-algebra, $I$ a closed two-sided ideal of $A$, and $f$ a
positive form on $A$.
\begin{enumerate}[label=\textup{(\roman*)}]
\item There is a unique decomposition $f=f_1+f_2$, where $f_1,f_2$ are
positive forms on $A$ such that $\lVert f_1\rVert=\lVert f_1|_I\rVert$ and
$f_2(I)=0$.
\item The pair $(\pi_f,\xi_f)$ is identified with
$(\pi_{f_1}\oplus\pi_{f_2},\xi_{f_1}+\xi_{f_2})$; $\pi_{f_2}$ vanishes on
$I$, and $\pi_{f_1}|_I$ is nondegenerate.
\end{enumerate}
\end{itshape}

Put $\pi_f=\pi$. Let $K_1$ be the essential subspace of $\pi|_I$, let
$K_2=H_\pi\ominus K_1$, and set
\[
\xi_1=P_{K_1}\xi_f,\qquad \xi_2=P_{K_2}\xi_f,
\]
with $\pi_1$ and $\pi_2$ the subrepresentations of $\pi$ defined by $K_1$
and $K_2$ (\XRef{2.11.1}{2.11.1}(ii)). Then $\xi_i$ is cyclic for $\pi_i$;
hence, if $f_i=\omega_{\xi_i}\circ\pi_i$, the pair $(\pi_i,\xi_i)$ is
identified with $(\pi_{f_i},\xi_{f_i})$ (\XRef{2.4.6}{2.4.6}). We have
$\pi_2(I)=0$, and consequently $f_2(I)=0$. The representations $\pi_1$ and
$\pi_1|_I$ are nondegenerate, and therefore
\[
\lVert f_1|_I\rVert=\lVert\xi_1\rVert^2=\lVert f_1\rVert
\]
by \XRef{2.4.3}{2.4.3}. The equality $\xi_f=\xi_1+\xi_2$ gives
$f=f_1+f_2$. In particular, $f_2=0$ if
$\lVert f\rVert=\lVert f|_I\rVert$.

It remains to prove the uniqueness assertion in (i). Let $f'_1,f'_2$ be
positive forms on $A$ such that
\[
f=f'_1+f'_2,\qquad \lVert f'_1\rVert=\lVert f'_1|_I\rVert,\qquad
f'_2(I)=0.
\]
The preceding argument, applied to $f'_1$ in place of $f$, proves that
$\pi_{f'_1}|_I$ is nondegenerate. Thus, if $(u_\lambda)$ is an approximate
identity of $I$, then for every $x\in A$,
$f'_1(x)=\lim_\lambda f'_1(xu_\lambda)$; likewise
$f_1(x)=\lim_\lambda f_1(xu_\lambda)$. But $f_1$ and $f'_1$ coincide on
$I$, since $f_2(I)=f'_2(I)=0$. Hence $f_1=f'_1$ and $f_2=f'_2$.

\ResultHeading{2.11.8}{2.11.8. Proposition.}
\begin{itshape}
Let $A$ be a $C^{*}$-algebra, $I$ a closed two-sided ideal of $A$, and let
$P_I(A)$ (respectively $P^I(A)$) be the set of pure states of $A$ that vanish
on $I$ (respectively do not vanish on $I$).
\begin{enumerate}[label=\textup{(\roman*)}]
\item For every $f\in P_I(A)$, let $f'$ be the positive form on $A/I$ induced
from $f$ by passage to the quotient. Then $f\mapsto f'$ is a bijection from
$P_I(A)$ onto $P(A/I)$.
\item The map $f\mapsto f|_I$ is a bijection from $P^I(A)$ onto $P(I)$.
\end{enumerate}
\end{itshape}

Let $f\in P_I(A)$, and denote by $\omega$ the canonical homomorphism from
$A$ onto $A/I$. If $f'$ majorizes a positive form $g'$ on $A/I$, then $f$
majorizes $g=g'\circ\omega$; hence $g=\lambda f$, with
$0\leq\lambda\leq1$, so $g'=\lambda f'$ and $f'\in P(A/I)$. It is clear
that the map $f\mapsto f'$ is injective. We show that it is surjective. Let
$h\in P(A/I)$ and $f=h\circ\omega$. If $f$ majorizes a positive form $g$,
then $g$ vanishes on $I^+$ and hence on $I$, so $g=k\circ\omega$, where $k$
is a positive form majorized by $h$. Thus $k=\lambda h$, with
$0\leq\lambda\leq1$, whence $g=\lambda f$ and $f\in P_I(A)$. Since
$h=f'$, surjectivity has been proved.

Let $f\in P^I(A)$. In the decomposition $f=f_1+f_2$ of
\XRef{2.11.7}{2.11.7}, $f_1$ and $f_2$ are proportional to $f$. But $f_2$
vanishes on $I$ whereas $f$ does not, so $f_2=0$ and $f=f_1$. Hence
$\lVert f|_I\rVert=\lVert f\rVert=1$, and $f|_I$ is a state of $I$. By
\XRef{2.10.3}{2.10.3} and \XRef{2.11.7}{2.11.7}(ii), $\pi_f|_I$ is
irreducible, and therefore $f|_I\in P(I)$. If $f,f'\in P^I(A)$ and
$f|_I=f'|_I$, there is an isomorphism from $H_{\pi_f}$ onto $H_{\pi_{f'}}$
which carries $\xi_{f|_I}$ (that is, $\xi_f$) to $\xi_{f'|_I}$ (that is,
$\xi_{f'}$),
% Source PDF physical page 064
and carries $\pi_{f|_I}$ to $\pi_{f'|_I}$, and hence $\pi_f$ to $\pi_{f'}$
(\XRef{2.10.4}{2.10.4}); consequently $f=f'$, and the map in (ii) is
injective. Finally, if $g\in P(I)$, the irreducible representation $\pi_g$
of $I$ extends to an irreducible representation $\pi$ of $A$
(\XRef{2.10.4}{2.10.4}); since $\lVert\xi_g\rVert=1$,
$f=\omega_{\xi_g}\circ\pi$ is a pure state of $A$ extending $g$.
Therefore $f\in P^I(A)$ and $f|_I=g$, which proves that the map in (ii) is
surjective.

\ResultHeading{2.11.9}{2.11.9.}
The six bijections in Propositions \XRef{2.11.2}{2.11.2},
\XRef{2.11.5}{2.11.5}, and \XRef{2.11.8}{2.11.8} are called canonical.

In summary, we have the following diagram of canonical maps:
\[
\begin{array}{ccccc}
P(A/I)&\longrightarrow&P(A)&\longleftarrow&P(I)\\
\downarrow&&\downarrow&&\downarrow\\
\widehat{A/I}&\longrightarrow&\widehat A&\longleftarrow&\widehat I\\
\downarrow&&\downarrow&&\downarrow\\
\operatorname{Prim}(A/I)&\longrightarrow&\operatorname{Prim}(A)&\longleftarrow&
\operatorname{Prim}(I).
\end{array}
\]
The vertical arrows are surjective maps. On each line, the two horizontal
arrows are injections onto complementary subsets of the middle set. It is
easy to verify that this diagram is commutative.

Bibliography: \cite{ref21,ref64,ref77,ref135}.

\BookSubsection{2.12}{Complements}

\ResultHeading{2.12.1}{2.12.1.}
Let $A$ be a unital $C^{*}$-algebra, and let $f,g$ be two pure states of $A$
such that $\lVert f-g\rVert<2$. Then $\pi_f\simeq\pi_g$.
\cite{ref49}

\ResultHeading{2.12.2}{2.12.2.}
Let $A$ be a unital $C^{*}$-algebra, let $E(A)$ be the set of states of $A$,
and let $f\in E(A)$. The following conditions are equivalent:
\begin{enumerate}[label=\textup{(\roman*)}]
\item for every neighborhood $U$ of $f$ in $E(A)$, there exist $\delta>0$
and $a\in A$ such that $0\leq a\leq1$, $f(a)=1$, and
$g(a)<1-\delta$ for $g\in E(A)\setminus U$;
\item for every closed $G_\delta$ subset $S$ of $E(A)$ containing $f$, there
exists $a\in A$ such that $\lVert a\rVert=|f(a)|$ and $S$ contains the set of
$g\in E(A)$ for which $\lVert a\rVert=|g(a)|$;
\item $f$ is a pure state.
\end{enumerate}
\cite{ref46}

\ResultHeading{2.12.3}{2.12.3.}
Let $A$ be a $C^{*}$-algebra, and let $K_n$ be the set of states $f$ of $A$
such that $\dim\pi_f=n<+\infty$. Then $P(A)\cap K_n$ is an open subset of
$K_n$. \cite{ref52}

\ResultHeading{2.12.4}{2.12.4.}
Let $A$ be a von Neumann algebra, $\mathfrak Z$ its center, and $\Omega$ the
spectrum of $\mathfrak Z$. For every $\omega\in\Omega$, let $I_\omega$ be the
norm-closed two-sided ideal of $A$ generated by $\ker\omega$. For every
$T\in A$, let $T_\omega$ be its canonical image in $A/I_\omega$.
\begin{enumerate}[label=\textit{\alph*.}]
\item If $f$ is a pure state of $A$, then $f|_{\mathfrak Z}$ is an element
$\omega$ of $\Omega$ and $f$ vanishes on $I_\omega$.
\item For every $T\in A$, the function $\omega\mapsto\lVert T_\omega\rVert$
is continuous on $\Omega$.
\item In $A/I_\omega$, the product of two nonzero two-sided ideals is
nonzero.
\end{enumerate}
\cite{ref45}

% Source PDF physical page 065
\ResultHeading{2.12.5}{2.12.5.}
Let $A$ be a unital involutive Banach algebra. Suppose that, for every
continuous Hermitian linear form $f$ on $A$, there exist two positive forms
$g,h$ such that
\[
f=g-h,\qquad \lVert f\rVert=\lVert g\rVert+\lVert h\rVert.
\]
Then $A$ is a $C^{*}$-algebra. \cite{ref55}

\ResultHeading{2.12.6}{*2.12.6.}
Let $A_1$ and $A_2$ be two unital $C^{*}$-algebras, and let
$\overline{P(A_i)}$ be the weak closure of $P(A_i)$ in the dual of $A_i$.
Let $\mathcal E_i$ be the set of functions $f\mapsto f(x)$ on
$\overline{P(A_i)}$, with $x\in A_i$. If there exists a homeomorphism from
$\overline{P(A_1)}$ onto $\overline{P(A_2)}$ which carries $\mathcal E_1$
onto $\mathcal E_2$, then there exists a linear isometric bijection
$\varphi\colon A_1\to A_2$ such that
\[
\varphi(x^{\ast})=\varphi(x)^{\ast}\quad(x\in A_1),
\qquad
\varphi(x^2)=\varphi(x)^2\quad(x=x^{\ast}\in A_1).
\]
(See \XRef{1.9.10}{1.9.10}.) \cite{ref66}

\ResultHeading{2.12.7}{2.12.7.}
Let $A$ be a unital $C^{*}$-algebra, let $\overline{P(A)}$ be the weak closure
of $P(A)$ in the dual of $A$, and let $\mathcal E$ be the set of functions
$f\mapsto f(x)$ on $\overline{P(A)}$, with $x\in A$. If the product of any
two elements of $\mathcal E$ belongs to $\mathcal E$, then $A$ is
commutative. \cite{ref66}

\ResultHeading{2.12.8}{2.12.8.}
Let $A$ be a unital $C^{*}$-algebra, let $f$ be a state of $A$, and let $N$
be the set of $x\in A$ such that $f(x^{\ast}x)=0$. If
$\ker f=N+N^{\ast}$, then $f$ is pure. \cite{ref69}

\ResultHeading{2.12.9}{*2.12.9.}
A unital $C^{*}$-algebra $A$ is isomorphic to a von Neumann algebra if and
only if it satisfies the following conditions:
\begin{enumerate}[label=\textup{(\roman*)}]
\item every increasing, bounded, directed family of Hermitian elements of
$A$ has a supremum;
\item for every nonzero $x\in A^+$, there exists a state $f$ of $A$ with
$f(x)\ne0$ such that, for every increasing directed family $(y_i)$ in $A^+$
having supremum $y\in A^+$,
\[
f(y)=\sup_i f(y_i).
\]
\end{enumerate}
\cite{ref68,ref151,ref152}

\ResultHeading{2.12.10}{*2.12.10.}
Let $A$ be a unital $C^{*}$-algebra, and let $X$ be the compact space
$\overline{P(A)}$. For each $x\in A$, let $\varphi_x$ be the function
$f\mapsto f(x)$ on $\overline{P(A)}$. Let $\mathcal E$ be the set of all
$\varphi_x$, $x\in A$. Under the map $x\mapsto\varphi_x$, every state $g$ of
$A$ becomes a positive linear form on $\mathcal E$; in general this form
extends in infinitely many ways to a positive measure on $X$. Denote by
$N_g$ the collection of subsets of $X$ that are negligible for all these
measures.

If $\pi$ is a representation of $A$, denote by $N_\pi$ the intersection of
the $N_{h\circ\pi}$, where $h$ ranges over the normal states of the weak
closure of $\pi(A)$.

Let $H$ be a separable Hilbert space, $A$ a $C^{*}$-subalgebra of
$\mathcal L(H)$ containing $1$, $B$ the weak closure of $A$, and $\pi$ a
representation of $A$ on a separable Hilbert space. The representation
$\pi$ extends to an ultraweakly continuous representation of $B$ on $H_\pi$
if and only if $N_\pi\supset N_\iota$, where $\iota$ denotes the identity
representation of $A$. \cite{ref70}

\ResultHeading{2.12.11}{2.12.11.}
Let $A$ be a $C^{*}$-algebra and $G$ a topological group. For every $g\in G$,
let $\zeta_g$ be an automorphism of $A$. Suppose that $g\mapsto\zeta_g$ is a
representation of $G$, and that $g\mapsto\zeta_g(x)$ is continuous for every
$x\in A$. A state $f$ of $A$ is said to be \emph{stationary} for $\zeta$ if
$f(\zeta_g(x))=f(x)$ for every $x\in A$ and every $g\in G$. Suppose this is
the case. There exists a unique continuous unitary representation $\rho$ of
$G$ on $H_{\pi_f}$ such that:
\begin{enumerate}[label=\textup{\arabic*\textdegree}]
\item $\pi_f(\zeta_g(x))=\rho(g)\pi_f(x)\rho(g)^{-1}$ for every $x\in A$ and
every $g\in G$;
\item if $\eta_f$ denotes the canonical map from $A$ into $H_{\pi_f}$, then
$\eta_f(\zeta_g(x))=\rho(g)\eta_f(x)$ for every $x\in A$ and every $g\in G$.
\end{enumerate}
\cite{ref138}

% Source PDF physical page 066
\ResultHeading{2.12.12}{2.12.12.}
Let $A$ be a unital $C^{*}$-algebra, let $x$ be a normal element of $A$, and
let $\lambda\in\operatorname{Sp}_A x$. There exist an irreducible
representation $\pi$ of $A$ and a nonzero $\xi\in H_\pi$ such that
$\pi(x)\xi=\lambda\xi$. (Apply \XRef{2.10.2}{2.10.2} to the
$C^{*}$-subalgebra $B$ of $A$ generated by $1$ and $x$, and to the character
of $B$ whose value at $x$ is $\lambda$.) \cite{ref134}

\ResultHeading{2.12.13}{2.12.13.}
Let $A$ be a $C^{*}$-algebra without an identity element. Then $0$ belongs to
the weak closure of $P(A)$. [Indeed, let $x_1,\ldots,x_n\in A^+$ and
$\varepsilon>0$; we must construct $f\in P(A)$ such that
$f(x_i)\leq\varepsilon$ for $i=1,\ldots,n$. Replacing
$x_1,\ldots,x_n$ by $x_1+\cdots+x_n$, we may suppose
$n=1$; using \XRef{2.5.5}{2.5.5}, it is enough to construct a state $g$ of
$A$ such that $g(x_1)\leq\varepsilon$. Realize $A$ on a Hilbert space $H$
by a nondegenerate representation. Then $x_1$ is not invertible in
$A+\mathbb C1$, and hence is not invertible in $\mathcal L(H)$; there is a
unit vector $\xi\in H$ such that $(x_1\xi\mid\xi)\leq\varepsilon$, and we may
take $g=\omega_\xi$.] This was pointed out to me by J.~Glimm.

\ResultHeading{2.12.14}{2.12.14.}
Let $A$ be a $C^{*}$-algebra, let $I,J$ be two closed two-sided ideals of
$A$, and let $h$ be a state of $A$ that vanishes on $I\cap J$. Then
$h=\lambda f+\mu g$, where $\lambda\geq0$, $\mu\geq0$,
$\lambda+\mu=1$, and $f$ (respectively $g$) is a state of $A$ that vanishes
on $I$ (respectively $J$). (Reduce to the case $I\cap J=0$. The essential
subspaces of $\pi_h|_I$ and $\pi_h|_J$ are then orthogonal. Apply
\XRef{2.11.7}{2.11.7}.) \cite{ref135}

\ResultHeading{2.12.15}{2.12.15.}
Let $A$ and $B$ be $C^{*}$-algebras, let $\pi,\pi'$ be faithful
representations of $A$, and let $\rho,\rho'$ be faithful representations of
$B$. Let $D$ (respectively $D'$) be the $C^{*}$-algebra of operators on
$H_\pi\otimes H_\rho$ (respectively $H_{\pi'}\otimes H_{\rho'}$) generated
by the $\pi(x)\otimes\rho(y)$ (respectively
$\pi'(x)\otimes\rho'(y)$), where $x\in A$, $y\in B$. There is a unique
isomorphism from $D$ onto $D'$ which carries
$\pi(x)\otimes\rho(y)$ to $\pi'(x)\otimes\rho'(y)$ for all $x\in A$ and
$y\in B$. Thus $D$, considered as an abstract $C^{*}$-algebra, depends only
on $A$ and $B$ and is called the $C^{*}$-tensor product of $A$ and $B$.
\cite{ref172,ref173,ref174,ref184}

\ResultHeading{2.12.16}{*2.12.16.}
\begin{enumerate}[label=\textit{\alph*.}]
\item Let $A$ and $B$ be $C^{*}$-algebras, and let $\mu$ be a linear map from
$A$ into $B$. We say that $\mu$ is positive if $\mu(A^+)\subset B^+$. Let
$A^{(n)}$ be the $C^{*}$-algebra of $n$-by-$n$ matrices with entries in $A$.
Applying $\mu$ to each entry of such a matrix gives a map
$\mu^{(n)}\colon A^{(n)}\to B^{(n)}$. We say that $\mu$ is completely
positive if $\mu^{(n)}$ is positive for every $n$. In general complete
positivity is strictly stronger than positivity. The two notions are
equivalent when $A$ is commutative.
\item Let $A$ be a unital $C^{*}$-algebra, $H$ a Hilbert space, and
$\mu\colon A\to\mathcal L(H)$ a linear map such that $\mu(1)=1$. There exist
a Hilbert space $K$ containing $H$ as a subspace, and a homomorphism
$\rho$ from $A$ into $\mathcal L(K)$, such that
\[
\mu(x)=P_H\rho(x)|_H\qquad(x\in A),
\]
if and only if $\mu$ is completely positive.
\end{enumerate}
\cite{ref114}

\ResultHeading{2.12.17}{2.12.17.}
Let $A$ be a $C^{*}$-algebra, let $A'$ be the Banach dual of $A$, and let
$A_h$ and $A'_h$ be the spaces of Hermitian elements of $A$ and $A'$,
which are ordered vector spaces. The following conditions are equivalent:
$1^\circ$ $A_h$ is a lattice; $2^\circ$ $A'_h$ is a lattice; $3^\circ$ $A$
is commutative. \cite{ref38}

\ResultHeading{2.12.18}{2.12.18.}
Let $A$ be a $C^{*}$-algebra, let $f$ be a pure state of $A$, and let $N_f$
be the set of $x\in A$ such that $f(x^{\ast}x)=0$, so that $A/N_f$, equipped
with the inner product
% Source PDF physical page 067
induced by $f(y^{\ast}x)$, is a Hilbert space. The norm of this space is also
the norm of the quotient Banach space $A/N_f$. \cite{ref156}

\ResultHeading{2.12.19}{*2.12.19.}
Let $A$ be a unital $C^{*}$-algebra, and let $f$ be a positive form on $A$.
\begin{enumerate}[label=\textit{\alph*.}]
\item Let $E_f$ be the set of linear forms on $A$ of the form
\[
x\longmapsto f(x_0xx_0^{\ast}),\qquad x_0\in A.
\]
Let $F_f$ be the norm closure of $E_f$ in the dual $A'$ of $A$. If $g$ is a
positive form on $A$, the following conditions are equivalent:
\begin{enumerate}[label=\textup{(\roman*)}]
\item $g\in F_f$;
\item $\pi_g\leq\pi_f$;
\item there exists $\xi\in H_{\pi_f}$ such that
$g(x)=(\pi_f(x)\xi\mid\xi)$ for every $x\in A$.
\end{enumerate}
\item If $f$ is pure, then $F_f=E_f$.
\item Let $E'_f$ be the set of positive forms on $A$ majorized by a multiple
$\lambda f$ of $f$, where $\lambda\geq0$. Let $F'_f$ be the norm closure of
$E'_f$ in $A'$. Then $F'_f$ is the set of forms
\[
x\longmapsto(\pi_f(x)\eta\mid\eta)\quad\text{on }A,
\qquad
\eta\in\overline{\pi_f(A)'\xi_f}.
\]
\end{enumerate}
\cite{ref71}

\ResultHeading{2.12.20}{2.12.20.}
Let $A$ be an involutive Banach algebra possessing an approximate identity,
and let $f$ be a continuous positive form on $A$. Denote by $L^2(f)$ the set
of $g\in A'$ such that
\[
\sup_{\substack{x\in A\\ f(x^{\ast}x)\ne0}}
|g(x)|\,f(x^{\ast}x)^{-1/2}<+\infty.
\]
Let $H_f$ be the space of $\pi_f$, and let $\eta_f$ be the canonical map from
$A$ into $H_f$. By transposition, $\eta_f$ defines a map
$\zeta_f\colon H_f\to A'$ such that
\[
\langle\zeta_f(\xi),x\rangle=(\eta_f(x)\mid\xi)
\qquad(\xi\in H_f,\ x\in A).
\]
The map $\zeta_f$ is a semilinear bijection from $H_f$ onto $L^2(f)$ which
carries $\xi_f$ to $f$. For every $x\in A$, let $\pi'(x)$ be the transpose
of left multiplication by $x$ on $A$, which is a continuous linear operator
on $A'$. Then $\zeta_f$ carries $\pi_f(x)$ to
$\pi'(x^{\ast})|_{L^2(f)}$. Finally, $\zeta_f$ is isometric if $L^2(f)$ is
equipped with the norm
\[
\lVert g\rVert_{L^2(f)}=
\sup_{\substack{x\in A\\ f(x^{\ast}x)\ne0}}
|g(x)|\,f(x^{\ast}x)^{-1/2}.
\]
\cite{ref163}

\ResultHeading{2.12.21}{2.12.21.}
If $A$ is a $C^{*}$-algebra whose only nilpotent element is $0$, then $A$ is
commutative. [Let $x,y\in A$, with $x=x^{\ast}$, and let $f,g$ be continuous
functions on $\operatorname{Sp}'x$ such that $fg=0$. Then
$(f(x)yg(x))^2=0$, and hence $f(x)yg(x)=0$. Therefore
$f(\pi(x))Tg(\pi(x))=0$ for every irreducible representation $\pi$ of $A$
and every $T\in\mathcal L(H_\pi)$; consequently $\pi(A)$ consists of
scalars.] (I.~Kaplansky, unpublished; communicated by R.~V.~Kadison.)

\BookSection[Spectrum of a C*-Algebra]{3}{Spectrum of a $C^{*}$-Algebra}

Our two principal aims will now be the following: $1^\circ$ the decomposition
of a representation into irreducible representations; $2^\circ$ the
structure of $C^{*}$-algebras. For both problems, it is important to
generalize to arbitrary $C^{*}$-algebras the spectrum of a commutative
$C^{*}$-algebra.

% Source PDF physical page 068
We already know three candidates: the sets $P(A)$, $\widehat A$, and
$\operatorname{Prim}(A)$. When $A$ is commutative, these sets reduce to the
set of characters of $A$: this is clear for $\widehat A$, it was noted for
$P(A)$ (\XRef{2.5.2}{2.5.2}), and for $\operatorname{Prim}(A)$ it follows
from the fact that the characters of a commutative $C^{*}$-algebra are in
bijective correspondence with their kernels. So far only $P(A)$ has been
equipped with a topology, namely the weak topology (the topology defined by
the norm of the dual of $A$ is almost never useful). Unfortunately, this
space is too large (it has no tendency to be finite-dimensional when $A$ is
the $C^{*}$-algebra of a Lie group), and it is not locally compact. In this
section we shall define topologies on $\widehat A$ and on
$\operatorname{Prim}(A)$. In favorable cases these two spaces are
homeomorphic. [In unfavorable cases, $\operatorname{Prim}(A)$ may reduce to
a single point while $\widehat A$ is very complicated, thereby reflecting
the complexity of $A$; it is for this admittedly not very convincing reason
that preference is given to $\widehat A$ over $\operatorname{Prim}(A)$.]
These spaces are locally quasi-compact. In general they are not Hausdorff,
but we shall see that this defect is rather limited. Moreover, in the
examples this failure of separation corresponds well to certain structural
peculiarities of the algebras under consideration.

\BookSubsection{3.1}{Jacobson Topology}

\ResultHeading{3.1.1}{3.1.1.}
Let $A$ be an algebra over a commutative field. For every subset $T$ of
$\operatorname{Prim}(A)$, let $I(T)$ be the intersection of the elements of
$T$. This set is a two-sided ideal of $A$. Let $\overline T$ be the set of
primitive ideals of $A$ containing $I(T)$.

\ResultLabel{Lemma.}
\begin{itshape}
\begin{enumerate}[label=\textup{(\roman*)}]
\item $\overline\varnothing=\varnothing$;
\item $T\subset\overline T$ for $T\subset\operatorname{Prim}(A)$;
\item $\overline{\overline T}=\overline T$ for
$T\subset\operatorname{Prim}(A)$;
\item $\overline{T_1\cup T_2}=\overline{T_1}\cup\overline{T_2}$ for
$T_1,T_2\subset\operatorname{Prim}(A)$.
\end{enumerate}
\end{itshape}

Assertions (i) and (ii) are evident. Clearly
$I(\overline T)=I(T)$, whence $\overline{\overline T}=\overline T$. Let
$I_1=I(T_1)$ and $I_2=I(T_2)$. We have
$I(T_1\cup T_2)=I_1\cap I_2$. Hence, by \XRef{2.11.4}{2.11.4},
$\overline{T_1\cup T_2}$ is the set of primitive ideals of $A$ containing
$I_1$ or $I_2$, which proves (iv).

There exists a unique topology on $\operatorname{Prim}(A)$ such that, for
every $T\subset\operatorname{Prim}(A)$, $\overline T$ is the closure of $T$
for this topology.

\ResultLabel{Definition.}
\begin{itshape}
This topology is called the Jacobson topology on
$\operatorname{Prim}(A)$.
\end{itshape}

% Source PDF physical page 069
\ResultHeading{3.1.2}{3.1.2. Proposition.}
\begin{itshape}
Let $A$ be an algebra, and let $T$ be a subset of
$\operatorname{Prim}(A)$. The set $T$ is closed if and only if it is the set
of primitive ideals containing some given subset of $A$.
\end{itshape}

If $T$ is closed, then $T=\overline T$, so $T$ is the set of primitive
ideals of $A$ containing $I(T)$. Let $M\subset A$. If $T$ is the set of
primitive ideals of $A$ containing $M$, then $I(T)\supset M$, and hence
$\overline T\subset T$ and $\overline T=T$.

\ResultHeading{3.1.3}{3.1.3.}
Recall that a topological space is called a $T_0$-space if, for any two
distinct points of the space, one of the two points has a neighborhood that
does not contain the other.

\ResultLabel{Proposition.}
\begin{itshape}
The space $\operatorname{Prim}(A)$ is a $T_0$-space.
\end{itshape}

Let $I_1,I_2$ be two distinct points of $\operatorname{Prim}(A)$. For
example, suppose $I_1\not\subset I_2$. Then the set of
$I\in\operatorname{Prim}(A)$ which contain $I_1$ is a closed set $T$
(\XRef{3.1.2}{3.1.2}) such that $I_1\in T$ and $I_2\notin T$.

\ResultHeading{3.1.4}{3.1.4. Proposition.}
\begin{itshape}
Let $I\in\operatorname{Prim}(A)$. The singleton $\{I\}$ is closed in
$\operatorname{Prim}(A)$ if and only if $I$ is a maximal primitive ideal.
\end{itshape}

Indeed, the closure of $\{I\}$ consists of the primitive ideals of $A$
containing $I$.

\ResultHeading{3.1.5}{3.1.5.}
Let $A$ be an algebra, and let $\widehat A$ be the set of equivalence classes
of nonzero irreducible representations of $A$. The map
$\pi\mapsto\ker\pi$ is a canonical surjection from $\widehat A$ onto
$\operatorname{Prim}(A)$.

\ResultLabel{Definition.}
\begin{itshape}
The spectrum of $A$ is the set $\widehat A$ equipped with the inverse
image topology of the Jacobson topology under the canonical map
$\widehat A\to\operatorname{Prim}(A)$.
\end{itshape}

\ResultHeading{3.1.6}{3.1.6. Proposition.}
\begin{itshape}
The following three conditions are equivalent:
\begin{enumerate}[label=\textup{(\roman*)}]
\item $\widehat A$ is a $T_0$-space;
\item two irreducible representations of $A$ having the same kernel are
equivalent;
\item the canonical map $\widehat A\to\operatorname{Prim}(A)$ is a
homeomorphism.
\end{enumerate}
\end{itshape}

(ii)$\Rightarrow$(iii) is evident; (iii)$\Rightarrow$(i) follows from
\XRef{3.1.3}{3.1.3}; for (i)$\Rightarrow$(ii), let $\pi_1,\pi_2$ be two
representations with the same kernel and belonging to $\widehat A$. Every
open subset of $\widehat A$ containing $\pi_1$ also contains $\pi_2$; hence
$\pi_1=\pi_2$ if $\widehat A$ is a $T_0$-space.

\ResultHeading{3.1.7}{3.1.7. Lemma.}
\begin{itshape}
Suppose that $A$ has an identity element. Let $I$ be a maximal
two-sided ideal of $A$. Then $I$ is a primitive ideal.
\end{itshape}

% Source PDF physical page 070
Let $M$ be a maximal left ideal containing $I$. Let $\pi$ be the canonical
representation of $A$ on $A/M$, which is irreducible. Since
$IA\subset I\subset M$, the kernel $I'$ of $\pi$ contains $I$ and does not
contain $1$. Therefore $I'=I$, and $I$ is primitive.

\ResultHeading{3.1.8}{3.1.8. Proposition.}
\begin{itshape}
If $A$ has an identity element, then $\operatorname{Prim}(A)$ and
$\widehat A$ are quasi-compact (that is, they satisfy the Borel--Lebesgue
axiom without necessarily being Hausdorff).
\end{itshape}

It is enough to prove this for $\operatorname{Prim}(A)$. Let $(T_i)$ be a
family of closed subsets of $\operatorname{Prim}(A)$ with empty intersection.
If $\sum I(T_i)\ne A$, the ideal $\sum I(T_i)$ is contained in a maximal
two-sided ideal $I$; by \XRef{3.1.7}{3.1.7}, $I$ is primitive. But
$I\in T_i$ for every $i$, because every $T_i$ is closed, which is a
contradiction. Therefore $\sum I(T_i)=A$. Hence there exist
$x_1\in I(T_{i_1}),\ldots,x_n\in I(T_{i_n})$ such that
$1=x_1+\cdots+x_n$. This implies
$I(T_{i_1})+\cdots+I(T_{i_n})=A$, and consequently
$T_{i_1}\cap\cdots\cap T_{i_n}=\varnothing$.

\ResultHeading{3.1.9}{3.1.9.}
Let $A$ be an algebra and $\varphi$ an automorphism of $A$. For every
irreducible representation $\pi$ of $A$, $\pi\circ\varphi^{-1}$ is an
irreducible representation of $A$ whose kernel is the image under $\varphi$
of the kernel of $\pi$. We thus obtain bijections
\[
\operatorname{Prim}(\varphi)\colon\operatorname{Prim}(A)\longrightarrow
\operatorname{Prim}(A),
\qquad
\widehat\varphi\colon\widehat A\longrightarrow\widehat A,
\]
which are plainly homeomorphisms. If $\varphi_1,\varphi_2$ are two
automorphisms of $A$, then
\[
\operatorname{Prim}(\varphi_1\varphi_2)
=\operatorname{Prim}(\varphi_1)\operatorname{Prim}(\varphi_2),
\qquad
\widehat{\varphi_1\varphi_2}=\widehat\varphi_1\widehat\varphi_2.
\]

Many authors say ``dual'' instead of ``spectrum,'' but this causes confusion
with the dual of the vector space $A$ or of the Banach space $A$.

Bibliography: \cite{ref64}.

\BookSubsection{3.2}{Spectrum of an Ideal and of a Quotient Algebra}

We shall restrict ourselves to the case of $C^{*}$-algebras, although the
following proposition can be extended to the general case.

\ResultHeading{3.2.1}{3.2.1. Proposition.}
\begin{itshape}
Let $A$ be a $C^{*}$-algebra and $I$ a closed two-sided ideal of $A$. The
canonical bijections (see \XRef{2.11.2}{2.11.2} and
\XRef{2.11.5}{2.11.5})
\[
\begin{gathered}
\operatorname{Prim}_I(A)\longrightarrow\operatorname{Prim}(A/I),
\qquad
\widehat A_I\longrightarrow\widehat{A/I},\\
\operatorname{Prim}^{I}(A)\longrightarrow\operatorname{Prim}(I),
\qquad
\widehat A^{I}\longrightarrow\widehat I
\end{gathered}
\]
are homeomorphisms. The sets $\operatorname{Prim}_I(A)$ and $\widehat A_I$
are closed in $\operatorname{Prim}(A)$ and $\widehat A$, respectively. The
sets $\operatorname{Prim}^{I}(A)$ and $\widehat A^{I}$ are open in
$\operatorname{Prim}(A)$ and $\widehat A$, respectively.
\end{itshape}

By \XRef{3.1.2}{3.1.2}, $\operatorname{Prim}_I(A)$ is closed in
$\operatorname{Prim}(A)$, and therefore $\operatorname{Prim}^{I}(A)$ is
open; consequently, in $\widehat A$, $\widehat A_I$ is closed and
$\widehat A^{I}$ is open. It follows immediately from the definition of the
Jacobson topology that the bijection
$\operatorname{Prim}_I(A)\to\operatorname{Prim}(A/I)$ is a homeomorphism,
and hence so is the bijection $\widehat A_I\to\widehat{A/I}$. To prove that
the bijection $\operatorname{Prim}^{I}(A)\to\operatorname{Prim}(I)$ is a
homeomorphism, it is enough to prove the following: if $(J_\lambda)$ is a
family of elements of $\operatorname{Prim}^{I}(A)$ and $J$ is an element of
$\operatorname{Prim}^{I}(A)$, then
\[
J\supset\bigcap_\lambda J_\lambda
\quad\Longleftrightarrow\quad
J\cap I\supset\bigcap_\lambda(J_\lambda\cap I).
\]
% Source PDF physical page 071
The forward implication is evident. Suppose that
\[
J\cap I\supset\bigcap_\lambda(J_\lambda\cap I).
\]
Then $J\supset(\bigcap_\lambda J_\lambda)\cap I$, while
$J\not\supset I$; hence, by \XRef{2.11.4}{2.11.4},
$J\supset\bigcap_\lambda J_\lambda$, which proves the asserted equivalence.
The fact that the bijection $\widehat A^{I}\to\widehat I$ is a homeomorphism
follows at once.

\ResultHeading{3.2.2}{3.2.2.}
Let $A$ be a $C^{*}$-algebra. If $I$ is a closed two-sided ideal of $A$, we
shall identify $\widehat{A/I}$ with a closed subset of $\widehat A$, and
$\widehat I$ with the complementary open subset, by means of Proposition
\XRef{3.2.1}{3.2.1}.

\ResultLabel{Proposition.}
\begin{itshape}
Let $A$ be a $C^{*}$-algebra. Then $I\mapsto\widehat I$ is a
bijection from the set of closed two-sided ideals of $A$ onto the set of open
subsets of $\widehat A$. Moreover,
$I_1\subset I_2\Longleftrightarrow\widehat I_1\subset\widehat I_2$.
\end{itshape}

Let $U$ be an open subset of $\widehat A$, and put
$F=\widehat A\setminus U$. Since every primitive ideal of $A$ is closed
(\XRef{2.9.7}{2.9.7}), the intersection of the kernels of the
$\pi\in F$ is a closed two-sided ideal $I$ of $A$. Since $F$ is closed in
$\widehat A$, every irreducible representation of $A$ whose kernel contains
$I$ belongs to $F$. Therefore $\widehat I=\widehat A\setminus F=U$, and the
map in the proposition is surjective. Let $I_1,I_2$ be two closed two-sided
ideals of $A$. If $\widehat I_1=\widehat I_2$, the primitive ideals of $A$
containing $I_1$ are precisely those containing $I_2$. But every closed
two-sided ideal of a $C^{*}$-algebra is the intersection of the primitive
ideals containing it (\XRef{2.9.7}{2.9.7}). Hence $I_1=I_2$, and the map
$I\mapsto\widehat I$ is injective. The assertion
$I_1\subset I_2\Longleftrightarrow\widehat I_1\subset\widehat I_2$ is
evident.

\ResultHeading{3.2.3}{3.2.3.}
Recall that, in a $C^{*}$-algebra, the sum of two closed two-sided ideals is a
closed two-sided ideal (\XRef{1.8.4}{1.8.4}).

\ResultLabel{Proposition.}
\begin{itshape}
Let $A$ be a $C^{*}$-algebra, and let $I,J$ be two closed two-sided ideals of
$A$. Then
\[
\widehat{I+J}=\widehat I\cup\widehat J,
\qquad
\widehat{I\cap J}=\widehat I\cap\widehat J.
\]
In particular, $I\cap J=0$ if and only if
$\widehat I\cap\widehat J=\varnothing$.
\end{itshape}

This follows from the preceding proposition, since $I+J$ and $I\cap J$ are
the supremum and the infimum of $\{I,J\}$ in the ordered set of closed
two-sided ideals of $A$.

\ResultHeading{3.2.4}{3.2.4.}
Let $A$ be a $C^{*}$-algebra and $B$ the $C^{*}$-algebra obtained from $A$ by
adjoining an identity element. Then $A$ is a closed two-sided ideal of
% Source PDF physical page 072
codimension $1$ in $B$, and is the kernel of the one-dimensional
representation $\omega$ of $B$ defined by
\[
\omega(\lambda1+x)=\lambda
\qquad(\lambda\in\mathbb C,\ x\in A).
\]
The only primitive ideal of $B$ containing $A$ is $A$ itself. The only
irreducible representation of $B$ with kernel $A$ is $\omega$. Thus $\{A\}$
is closed in $\operatorname{Prim}(B)$, and $\{\omega\}$ is closed in
$\widehat B$. By \XRef{3.2.1}{3.2.1}, the spaces
$\operatorname{Prim}(A)$ and $\widehat A$ are identified with
$\operatorname{Prim}(B)\setminus\{A\}$ and
$\widehat B\setminus\{\omega\}$, respectively.

Suppose that $A$ has an identity element $e$. For every
$\pi\in\widehat B\setminus\{\omega\}$, $\pi(e)$ is a nonzero idempotent
commuting with $\pi(A)$, and hence $\pi(e)=1$; thus $\pi(e-1)=0$. It follows
that, for every $I\in\operatorname{Prim}(B)\setminus\{A\}$, one has
$e-1\in I$, whereas $e-1\notin A$. Thus $A$ is not in the closure of
$\operatorname{Prim}(B)\setminus\{A\}$ in $\operatorname{Prim}(B)$. In other
words, $\{A\}$ is not only closed but also open in $\operatorname{Prim}(B)$;
hence $\{\omega\}$ is not only closed but also open in $\widehat B$.

Conversely, suppose that $\{\omega\}$ is open in $\widehat B$. By
\XRef{3.2.2}{3.2.2} and \XRef{3.2.3}{3.2.3}, $B$ is the direct sum of two
closed two-sided ideals, one of which is $A$. Since $B$ has an identity
element, so does $A$.

Bibliography: \cite{ref77}.

\BookSubsection{3.3}{Norm and Topology}

\ResultHeading{3.3.1}{3.3.1. Lemma.}
\begin{itshape}
Let $A$ be a $C^{*}$-algebra, let $a$ be a Hermitian element of $A$, and let
$L$ be a closed set of real numbers.
\begin{enumerate}[label=\textup{(\roman*)}]
\item The set $Z'$ of $\pi\in\widehat A$ such that
$\operatorname{Sp}'\pi(a)\subset\{0\}\cup L$ is closed in $\widehat A$.
\item If $A$ has an identity element, the set $Z$ of $\pi\in\widehat A$ such
that $\operatorname{Sp}\pi(a)\subset L$ is closed in $\widehat A$.
\end{enumerate}
\end{itshape}

Let $\rho\in\overline{Z'}$. Suppose that $\operatorname{Sp}'\rho(a)$
contains a point $\alpha\notin\{0\}\cup L$. Let $f\colon\mathbb R\to\mathbb R$
be a continuous function which is nonzero at $\alpha$ and vanishes on
$\{0\}\cup L$. Then
\[
\pi(f(a))=f(\pi(a))
\]
vanishes for $\pi\in Z'$ and is nonzero for $\pi=\rho$, contradicting the
assumption $\rho\in\overline{Z'}$. This proves (i). Assertion (ii) is proved
in the same way.

\ResultHeading{3.3.2}{3.3.2. Proposition.}
\begin{itshape}
Let $A$ be a $C^{*}$-algebra and $x\in A$. The function
$\pi\mapsto\lVert\pi(x)\rVert$ is lower semicontinuous on $\widehat A$.
\end{itshape}

We are immediately reduced to the case where $A$ has an identity element.
Moreover, since
$\lVert\pi(x)\rVert^2=\lVert\pi(x^{\ast}x)\rVert$, we may suppose
$x\in A^+$. Let $k\geq0$, and let $Z$ be the set of $\pi\in\widehat A$ such
that $\lVert\pi(x)\rVert\leq k$. The set $Z$ is also the set of
$\pi\in\widehat A$ such that $\operatorname{Sp}\pi(x)\subset[0,k]$. By
\XRef{3.3.1}{3.3.1}, $Z$ is closed, which proves the proposition.

\ResultHeading{3.3.3}{3.3.3. Lemma.}
\begin{itshape}
Let $A$ be a $C^{*}$-algebra, let $(x_i)$ be a dense family of
elements of $A$, and let $Z_i$ be the set of $\pi\in\widehat A$ such that
$\lVert\pi(x_i)\rVert>1$. Then the $Z_i$ form a base for the topology of
$\widehat A$.
\end{itshape}

% Source PDF physical page 073
Let $U$ be an open subset of $\widehat A$ and let $\pi\in U$. Let $I\subset A$
be the intersection of the kernels of the representations belonging to
$\widehat A\setminus U$. There exists $x\in I$ such that
$\lVert\pi(x)\rVert=2$. Choose an index $i$ such that
$\lVert x-x_i\rVert<1$. Then $\lVert\pi(x_i)\rVert>1$, whereas
$\lVert\rho(x_i)\rVert<1$ for $\rho\in\widehat A\setminus U$. Thus
$\pi\in Z_i\subset U$. Since $Z_i$ is open by \XRef{3.3.2}{3.3.2}, this
proves the lemma.

\ResultHeading{3.3.4}{3.3.4. Proposition.}
\begin{itshape}
If $A$ is separable, the topology of $\widehat A$ has a countable
base.
\end{itshape}

This follows at once from \XRef{3.3.3}{3.3.3}.

\ResultHeading{3.3.5}{3.3.5. Lemma.}
\begin{itshape}
Let $A$ be a unital involutive algebra, let $x$ be a Hermitian
element of $A$, and let $\lambda\in\mathbb R$. For
$\lambda\in\operatorname{Sp}_A x$, it is necessary and sufficient that there
exist a primitive ideal $I$ of $A$ such that
$\lambda\in\operatorname{Sp}_{A/I}\dot x$, where $\dot x$ denotes the
canonical image of $x$ in $A/I$.
\end{itshape}

If $x-\lambda$ is invertible in $A$, its image is invertible in $A/I$ for
every primitive ideal $I$. Suppose that $x-\lambda$ is not invertible. Then
$x-\lambda$ is not left invertible (for a Hermitian element, left
invertibility is equivalent to right invertibility), and is therefore
contained in a maximal left ideal $L$. The kernel $I$ of the canonical
representation of $A$ on $A/L$ is a primitive ideal of $A$, and the image of
$x-\lambda$ in $A/I$ is not invertible.

\ResultHeading{3.3.6}{3.3.6. Lemma.}
\begin{itshape}
Let $A$ be a $C^{*}$-algebra. For every $x\in A$, the function
$\pi\mapsto\lVert\pi(x)\rVert$ on $\widehat A$ attains its upper bound
$\lVert x\rVert$.
\end{itshape}

Replacing $x$ by $x^{\ast}x$, we are reduced to the case $x\in A^+$. We may
also suppose that $A$ has an identity element. Then
$\lVert x\rVert\in\operatorname{Sp}_A x$. By \XRef{3.3.5}{3.3.5}, there
exists $\pi\in\widehat A$ such that $\lVert x\rVert$ belongs to the spectrum
of the operator $\pi(x)$; hence
\[
\lVert x\rVert\leq\lVert\pi(x)\rVert\leq\lVert x\rVert.
\]

\ResultHeading{3.3.7}{3.3.7. Proposition.}
\begin{itshape}
Let $A$ be a $C^{*}$-algebra, $x\in A$, and $\alpha>0$. The set $Z$
of $\pi\in\widehat A$ such that $\lVert\pi(x)\rVert\geq\alpha$ is
quasi-compact.
\end{itshape}

Let $(Z_i)$ be a decreasing directed family of nonempty subsets of $Z$ that
are relatively closed in $Z$. For every $i$, let $J_i$ be the intersection
of the kernels of the elements of $Z_i$. The canonical image of $x$ modulo
$J_i$ has norm at least $\alpha$. The $J_i$ form an increasing directed
family; let $J$ be the closure of their union, which is a closed two-sided
ideal of $A$. The definition of the norm in a quotient normed space shows
that the canonical image of $x$ modulo $J$ has norm at least $\alpha$. By
\XRef{3.3.6}{3.3.6}, there exists $\pi\in\widehat A$ whose kernel contains
$J$ and such that $\lVert\pi(x)\rVert\geq\alpha$. We have $\pi\in Z$, and
$\pi$ belongs to the closure of $Z_i$ for every $i$; hence
$\pi\in\bigcap_i Z_i$, which proves the proposition.

\ResultHeading{3.3.8}{3.3.8.}
Recall that a topological space is said to be locally quasi-compact if every
point has a neighborhood base consisting of quasi-compact neighborhoods.

% Source PDF physical page 074
\ResultLabel{Corollary.}
\begin{itshape}
Let $A$ be a $C^{*}$-algebra. Its spectrum $\widehat A$ is locally
quasi-compact.
\end{itshape}

Let $\pi\in\widehat A$ and let $U$ be an open neighborhood of $\pi$ in
$\widehat A$. Since $\widehat A\setminus U$ is closed, there exists $x\in A$
such that $\pi(x)\ne0$ and $\rho(x)=0$ for
$\rho\in\widehat A\setminus U$. Let $V$ (respectively $W$) be the set of
$\rho\in\widehat A$ such that
\[
\lVert\rho(x)\rVert>\tfrac12\lVert\pi(x)\rVert
\quad\text{(respectively }\lVert\rho(x)\rVert\geq
\tfrac12\lVert\pi(x)\rVert\text{)}.
\]
By \XRef{3.3.2}{3.3.2}, $V$ is an open neighborhood of $\pi$. Thus $W$ is a
neighborhood of $\pi$ contained in $U$, and $W$ is quasi-compact by
\XRef{3.3.7}{3.3.7}.

\ResultHeading{3.3.9}{3.3.9. Corollary.}
\begin{itshape}
Let $A$ be a $C^{*}$-algebra such that $\widehat A$ is Hausdorff. For
every $x\in A$, the function $\pi\mapsto\lVert\pi(x)\rVert$ is continuous on
$\widehat A$.
\end{itshape}

The set $Z$ of \XRef{3.3.7}{3.3.7} is then closed, and hence the function
$\pi\mapsto\lVert\pi(x)\rVert$ is upper semicontinuous. Now apply
\XRef{3.3.2}{3.3.2}.

Bibliography: \cite{ref27,ref28,ref29,ref77,ref126}.

\BookSubsection{3.4}{Second Definition of the Topology on the Spectrum}

\ResultHeading{3.4.1}{3.4.1. Lemma.}
\begin{itshape}
Let $A$ be a unital $C^{*}$-algebra, let $A_h$ be the set of Hermitian
elements of $A$, let $E(A)$ be the set of states of $A$, and let $Q$ be a
subset of $E(A)$. Suppose that, if $x\in A_h$ satisfies $f(x)\geq0$ for every
$f\in Q$, then $x\in A^+$. Under these conditions:
\begin{enumerate}[label=\textup{(\roman*)}]
\item the weak closure of $Q$ contains $P(A)$;
\item the weakly closed convex hull of $Q$ is $E(A)$.
\end{enumerate}
\end{itshape}

Let $Q^\circ$ be the polar of $Q$ in $A_h$. If $x\in A_h$, then
\begin{align*}
x\in Q^\circ
&\Longleftrightarrow f(x)\leq1 &&\text{for every }f\in Q,\\
&\Longleftrightarrow f(1-x)\geq0 &&\text{for every }f\in Q,\\
&\Longleftrightarrow 1-x\in A^+,\\
&\Longleftrightarrow f(x)\leq1 &&\text{for every }f\in E(A).
\end{align*}
The weakly closed convex hull of $Q$ is plainly contained in $E(A)$; it is
equal to the bipolar of $Q$, and hence, by what precedes, contains $E(A)$.
This proves (ii), and (ii), together with \XRef{B.14}{B 14}, implies that
$\overline Q$ contains the set of extreme points of $E(A)$, namely $P(A)$
(\XRef{2.5.6}{2.5.6}).

\ResultHeading{3.4.2}{3.4.2. Proposition.}
\begin{itshape}
Let $A$ be a $C^{*}$-algebra, and let $(\pi_i)_{i\in I}$ be a family of
representations of $A$ on spaces $H_i$.
\begin{enumerate}[label=\textup{(\roman*)}]
\item Every state of $A$ that vanishes on $\bigcap_i\ker\pi_i$ is a weak
limit of states of the form
\[
\omega_{\xi_1}\circ\pi_{i_1}+\cdots+
\omega_{\xi_n}\circ\pi_{i_n},
\]
where $i_1,\ldots,i_n\in I$ and
$\xi_1\in H_{i_1},\ldots,\xi_n\in H_{i_n}$.
% Source PDF physical page 075
\item Every pure state of $A$ that vanishes on $\bigcap_i\ker\pi_i$ is a weak
limit of states of the form $\omega_\xi\circ\pi_i$, where $i\in I$ and
$\xi\in H_i$.
\end{enumerate}
\end{itshape}

By \XRef{2.2.6}{2.2.6}, we may suppose that every $\pi_i$ is nondegenerate.
Let $H=\bigoplus_i H_i$. Passing to the quotient by
$\bigcap_i\ker\pi_i$, we may suppose that $\bigcap_i\ker\pi_i=0$. The
representation $\rho=\bigoplus_i\pi_i$ is faithful, and we identify $A$ with
the $C^{*}$-subalgebra $\rho(A)$ of $\mathcal L(H)$; thus the identity
representation of $A$ is nondegenerate. By \XRef{2.4.3}{2.4.3}, we may
replace $A$ by $A+\mathbb C1$. Let $Q$ be the set of states of $A$ of the
form $\omega_\xi\circ\pi_i$, where $i\in I$, $\xi\in H_i$, and
$\lVert\xi\rVert=1$. If $x$ is a Hermitian element of $A$ such that
$(\omega_\xi\circ\pi_i)(x)\geq0$ for every $i\in I$ and every unit vector
$\xi\in H_i$, then $\pi_i(x)\geq0$ for every $i$, and hence $x\geq0$. It is
therefore enough to apply \XRef{3.4.1}{3.4.1}.

\ResultHeading{3.4.3}{3.4.3. Corollary.}
\begin{itshape}
Let $A$ be a $C^{*}$-algebra and $f$ a state of $A$. The states of
$A$ that vanish on $\ker\pi_f$ are weak limits of states of the form
\[
x\longmapsto\sum_i f(y_i^{\ast}xy_i),
\qquad y_1,\ldots,y_n\in A.
\]
The pure states of $A$ that vanish on $\ker\pi_f$ are weak limits of states
of the form $x\mapsto f(y^{\ast}xy)$, where $y\in A$.
\end{itshape}

This follows from \XRef{3.4.2}{3.4.2} and \XRef{2.4.8}{2.4.8}(ii).

\ResultHeading{3.4.4}{3.4.4. Theorem.}
\begin{itshape}
Let $A$ be a $C^{*}$-algebra, $\pi$ a representation of $A$, and $S$ a set
of representations of $A$. The following conditions are equivalent:
\begin{enumerate}[label=\textup{(\roman*)}]
\item the kernel of $\pi$ contains the intersection of the kernels of the
elements of $S$;
\item every positive form on $A$ associated with $\pi$ is a weak limit of
linear combinations of positive forms associated with $S$;
\item every state of $A$ associated with $\pi$ is a weak limit of states
that are sums of positive forms associated with $S$.
\end{enumerate}
\end{itshape}

(iii)$\Rightarrow$(ii) is evident. For (ii)$\Rightarrow$(i), suppose that
(ii) holds. If $a\in\bigcap_{\rho\in S}\ker\rho$, every positive form
associated with $S$ vanishes at $a^{\ast}a$, and hence every positive form
associated with $\pi$ vanishes at $a^{\ast}a$; thus $\pi(a)=0$. For
(i)$\Rightarrow$(iii), if
$\ker\pi\supset\bigcap_{\rho\in S}\ker\rho$, every state of $A$ associated
with $\pi$ vanishes on $\bigcap_{\rho\in S}\ker\rho$, and we apply
\XRef{3.4.2}{3.4.2}(i).

\ResultHeading{3.4.5}{3.4.5.}
It is clear that, if $\pi$ is contained in one of the elements of $S$, the
conditions of Theorem \XRef{3.4.4}{3.4.4} are satisfied. The converse is
% Source PDF physical page 076
false, even if $S$ consists of a single element. We therefore make the
following definition.

\ResultLabel{Definition.}
\begin{itshape}
If $\pi$ and $S$ satisfy the equivalent conditions (i), (ii), and
(iii) of Theorem \XRef{3.4.4}{3.4.4}, we say that $\pi$ is weakly contained
in $S$.
\end{itshape}

If $S$ and $T$ are two sets of representations of $A$, we say that $T$ is
weakly contained in $S$ if every element of $T$ is weakly contained in $S$.

\ResultHeading{3.4.6}{3.4.6. Definition.}
\begin{itshape}
Let $A$ be a $C^{*}$-algebra and $\pi$ a representation of $A$. The
support of $\pi$ is the set of $\rho\in\widehat A$ that are weakly contained
in $\pi$.
\end{itshape}

Using condition (i) of \XRef{3.4.4}{3.4.4}, we see that this support $S$ is a
closed subset of $\widehat A$. Since $\ker\pi$ is the intersection of the
primitive ideals containing it, $\pi$ is weakly contained in $S$.

\ResultHeading{3.4.7}{3.4.7.}
Let $A$ be a $C^{*}$-algebra, $\pi$ a representation of $A$, $S$ a set of
representations of $A$, $\widetilde A$ the $C^{*}$-algebra obtained from $A$
by adjoining an identity element, $\widetilde\pi$ the canonical extension of
$\pi$ to $\widetilde A$, and $\widetilde S$ the set of canonical extensions
to $\widetilde A$ of the elements of $S$. If $\widetilde\pi$ is weakly
contained in $\widetilde S$, it is clear that $\pi$ is weakly contained in
$S$ [condition (i) of Theorem \XRef{3.4.4}{3.4.4}]. The converse is true if
$\pi$ is nondegenerate. Indeed, let $x\in A$ and $\lambda\in\mathbb C$, and
suppose that $\rho(x)=\lambda1$ for every $\rho\in S$. Then
$\rho(xy-\lambda y)=0$ for every $y\in A$ and every $\rho\in S$, and hence
$\pi(xy-\lambda y)=0$ for every $y\in A$. Thus
$(\pi(x)-\lambda1)\pi(y)=0$ for every $y\in A$, and consequently
$\pi(x)=\lambda1$ because $\pi$ is nondegenerate.

\ResultHeading{3.4.8}{3.4.8.}
Let $A$ be a $C^{*}$-algebra. To every set $S$ of representations of $A$ we
associate the set $Q(S)$ of positive forms on $A$ that are weak limits of
sums of positive forms associated with $S$. It is clear that $Q(S)$ is a
weakly closed convex cone in the dual of $A$. For $T$ to be weakly contained
in $S$, it is necessary and sufficient that $Q(T)\subset Q(S)$.

\ResultHeading{3.4.9}{3.4.9. Proposition.}
\begin{itshape}
Retain the notation of Theorem \XRef{3.4.4}{3.4.4}. If $\pi$ has a cyclic
vector $\xi$, the conditions of that theorem are also equivalent to the
following condition:

\textup{(ii$'$)} The positive form
$x\mapsto f(x)=(\pi(x)\xi\mid\xi)$ on $A$ is a weak limit of linear
combinations of positive forms associated with $S$.
\end{itshape}

It is clear that (ii) implies (ii$'$). Suppose that (ii$'$) holds, and let us
show that (ii) holds. By \XRef{2.4.8}{2.4.8}(ii), it is enough to show that,
for every $y\in A$, the form
\[
x\longmapsto g(x)=f(y^{\ast}xy)
\]
is a weak limit of linear combinations of positive forms associated with
$S$. But $f$ is a
% Source PDF physical page 077
weak limit of forms $f_i$ that are linear combinations of positive forms
associated with $S$. Thus $g$ is a weak limit of the forms
\[
x\longmapsto g_i(x)=f_i(y^{\ast}xy).
\]
By \XRef{2.4.8}{2.4.8}(i), each $g_i$ is a linear combination of positive
forms associated with $S$.

\ResultHeading{3.4.10}{3.4.10. Theorem.}
\begin{itshape}
Let $A$ be a $C^{*}$-algebra, $\pi\in\widehat A$, and
$S\subset\widehat A$. The following conditions are equivalent:
\begin{enumerate}[label=\textup{(\roman*)}]
\item $\pi\in\overline S$;
\item $\pi$ is weakly contained in $S$;
\item one of the nonzero positive forms associated with $\pi$ is a weak limit
of positive forms associated with $S$;
\item every state associated with $\pi$ is a weak limit of states associated
with $S$.
\end{enumerate}
\end{itshape}

(i)$\Longleftrightarrow$(ii) follows immediately from condition (i) of
Theorem \XRef{3.4.4}{3.4.4}; (iv)$\Rightarrow$(iii) is evident;
(iii)$\Rightarrow$(ii) follows from \XRef{3.4.9}{3.4.9}; and
(ii)$\Rightarrow$(iv) follows from \XRef{3.4.2}{3.4.2}(ii).

\ResultHeading{3.4.11}{3.4.11. Theorem.}
\begin{itshape}
Let $A$ be a $C^{*}$-algebra. The canonical map
$P(A)\to\widehat A$ is continuous and open.
\end{itshape}

Let $S$ be a subset of $\widehat A$, let $Q$ be its inverse image in $P(A)$,
let $f$ be a point of $P(A)$, and let $\pi$ be its image in $\widehat A$.
For $\pi$ to belong to $\overline S$, it is necessary and sufficient that
$f$ belong to $\overline Q$ (\XRef{3.4.10}{3.4.10}). Hence, for $S$ to be
closed in $\widehat A$, it is necessary and sufficient that $Q$ be closed in
$P(A)$. This first proves that the canonical map $P(A)\to\widehat A$ is
continuous. Let $U$ be an open subset of $P(A)$ and let $V$ be its image in
$\widehat A$. We show that $V$ is open.

Let $\rho\in V$, and let $g$ be a point of $U$ whose image in $\widehat A$ is
$\rho$. Then $g$ does not belong to the closure of $P(A)\setminus U$; a
fortiori, $g$ does not belong to the closure of the inverse image in $P(A)$
of $\widehat A\setminus V$. Thus $\rho$ does not belong to the closure of
$\widehat A\setminus V$, and consequently $V$ is open. This proves that the
canonical map $P(A)\to\widehat A$ is open.

\ResultHeading{3.4.12}{3.4.12.}
The second announced definition of the topology of $\widehat A$ is therefore
the following: it is the quotient topology of the topology of $P(A)$ under
the canonical map $P(A)\to\widehat A$. We see that, when $A$ is commutative,
the topological space $\widehat A$ is indeed the usual spectrum of $A$.

\ResultHeading{3.4.13}{3.4.13. Corollary.}
\begin{itshape}
Let $A$ be a $C^{*}$-algebra. The space $\widehat A$ is a Baire
space.
\end{itshape}

Let $(V_1,V_2,\ldots)$ be a decreasing sequence of open dense subsets of
$\widehat A$. Let $U_n$ be the inverse image of $V_n$ in $P(A)$. By
\XRef{3.4.11}{3.4.11}, $U_n$ is an open dense subset of $P(A)$. Since $P(A)$
is a Baire space (\XRef{B.14}{B 14}), $\bigcap_n U_n$ is dense in $P(A)$.
% Source PDF physical page 078
By \XRef{3.4.11}{3.4.11}, $\bigcap_n V_n$, which is the image of
$\bigcap_n U_n$ in $\widehat A$, is dense in $\widehat A$.

Bibliography: \cite{ref3,ref11,ref27,ref48,ref77,ref154}.

\BookSubsection{3.5}{Third Definition of the Topology on the Spectrum}

\ResultHeading{3.5.1}{3.5.1.}
For every cardinal $n$, we shall call the Hilbert space of families
$(\xi_i)$ of complex numbers such that
\[
\sum_i|\xi_i|^2<+\infty,
\]
indexed by a set of cardinality $n$ (for example, the least ordinal segment
of cardinality $n$), the \emph{model Hilbert space of dimension $n$}. This
choice will be immaterial in what follows; it is intended only to fix ideas.

Let $n$ be a cardinal, let $H_n$ be the model Hilbert space of dimension
$n$, let $A$ be a $C^{*}$-algebra, let $\operatorname{Rep}_n(A)$ be the set
of representations of $A$ on $H_n$, and let $\operatorname{Irr}_n(A)$ be the
set of nonzero irreducible representations of $A$ on $H_n$. There is an
evident canonical map from $\operatorname{Irr}_n(A)$ into $\widehat A$,
which assigns to each element of $\operatorname{Irr}_n(A)$ its class. The
canonical image of $\operatorname{Irr}_n(A)$ in $\widehat A$ is denoted
$\widehat A_n$: it is the set of classes of nonzero irreducible
representations of dimension $n$.

\ResultHeading{3.5.2}{3.5.2.}
We equip $\operatorname{Rep}_n(A)$ with the topology of pointwise weak
convergence on $A$. Thus $\pi_\lambda\to\pi$, where
$\pi_\lambda,\pi\in\operatorname{Rep}_n(A)$, means that
\[
(\pi_\lambda(a)\xi\mid\eta)\longrightarrow(\pi(a)\xi\mid\eta)
\]
for all $a\in A$ and $\xi,\eta\in H_n$. This topology is identical with the
topology of pointwise strong convergence. Indeed, if $\pi_\lambda\to\pi$ in
the preceding sense, then, for every $a\in A$ and every $\xi\in H_n$,
\begin{align*}
\lVert\pi_\lambda(a)\xi-\pi(a)\xi\rVert^2
={}&(\pi_\lambda(a)\xi\mid\pi_\lambda(a)\xi)
-2\operatorname{Re}(\pi_\lambda(a)\xi\mid\pi(a)\xi)
+(\pi(a)\xi\mid\pi(a)\xi)\\
={}&(\pi_\lambda(a^{\ast}a)\xi\mid\xi)
-2\operatorname{Re}(\pi_\lambda(a)\xi\mid\pi(a)\xi)
+(\pi(a^{\ast}a)\xi\mid\xi)\\
&\longrightarrow
(\pi(a^{\ast}a)\xi\mid\xi)
-2\operatorname{Re}(\pi(a)\xi\mid\pi(a)\xi)
+(\pi(a^{\ast}a)\xi\mid\xi)=0.
\end{align*}

\ResultHeading{3.5.3}{3.5.3.}
Let $\pi_\lambda,\pi\in\operatorname{Rep}_n(A)$, and suppose that $\pi$ has
a cyclic vector $\xi$. In order that $\pi_\lambda\to\pi$, it is enough that
$\pi_\lambda(a)\xi\to\pi(a)\xi$ for every $a\in A$. Indeed, in that case,
for all $a,b\in A$,
\begin{align*}
\lVert\pi_\lambda(a)\pi(b)\xi-\pi(a)\pi(b)\xi\rVert
\leq{}&\lVert\pi_\lambda(a)\pi(b)\xi
-\pi_\lambda(a)\pi_\lambda(b)\xi\rVert\\
&+\lVert\pi_\lambda(a)\pi_\lambda(b)\xi
-\pi(a)\pi(b)\xi\rVert\\
\leq{}&\lVert a\rVert\,
\lVert\pi(b)\xi-\pi_\lambda(b)\xi\rVert
+\lVert\pi_\lambda(ab)\xi-\pi(ab)\xi\rVert
\longrightarrow0.
\end{align*}
Since the $\pi(b)\xi$ are dense in $H_n$ and
$\lVert\pi_\lambda(a)\rVert\leq\lVert a\rVert$, it follows that
$\pi_\lambda(a)\eta\to\pi(a)\eta$ for every $\eta\in H_n$.

\ResultHeading{3.5.4}{3.5.4.}
Let $\pi_\lambda,\pi\in\operatorname{Rep}_n(A)$, and let $A'$ be a dense
subset of $A$. In order that $\pi_\lambda\to\pi$, it is enough that
$(\pi_\lambda(a)\xi\mid\eta)\to(\pi(a)\xi\mid\eta)$
% Source PDF physical page 079
for all $a\in A'$ and $\xi,\eta\in H_n$; this follows immediately from
\[
\lVert\pi_\lambda(a)-\pi_\lambda(a')\rVert\leq\lVert a-a'\rVert.
\]

\ResultHeading{3.5.5}{3.5.5.}
In $\operatorname{Rep}_n(A)$, equivalence of representations is an open
equivalence relation. Indeed, let $G$ be the group of unitary operators on
$H_n$. For $U\in G$ and $\pi\in\operatorname{Rep}_n(A)$, let
$U\pi\in\operatorname{Rep}_n(A)$ be the representation defined by
\[
(U\pi)(x)=U\pi(x)U^{-1}\qquad(x\in A).
\]
It is immediately seen that, for each $U$, the map $\pi\mapsto U\pi$ is a
homeomorphism of $\operatorname{Rep}_n(A)$ onto itself. Hence the saturation
of an open subset $\Omega$ of $\operatorname{Rep}_n(A)$, being the union of
the $U\Omega$ for $U\in G$, is open in $\operatorname{Rep}_n(A)$.

\ResultHeading{3.5.6}{3.5.6. Lemma.}
\begin{itshape}
Let $H$ be a Hilbert space, let $\xi_1,\ldots,\xi_n$ be vectors of $H$, and
let $\varepsilon>0$. There exists $\varepsilon'>0$ with the following
property: if $\eta_1,\ldots,\eta_n\in H$ satisfy
\[
\bigl|(\eta_i\mid\eta_j)-(\xi_i\mid\xi_j)\bigr|\leq\varepsilon'
\qquad\text{for all }i,j,
\]
then there exists a unitary operator $U$ on $H$ such that
$\lVert U\eta_i-\xi_i\rVert\leq\varepsilon$ for every $i$.
\end{itshape}

It is enough to prove the following. Let
$(\eta_1^p,\eta_2^p,\ldots,\eta_n^p)$, $p=1,2,\ldots$, be a sequence of
systems of vectors such that
\[
(\eta_i^p\mid\eta_j^p)\longrightarrow(\xi_i\mid\xi_j)
\qquad\text{for all }i,j.
\]
Then, denoting by $G$ the unitary group of $H$, we have
\[
\inf_{U\in G}
\bigl(\lVert U\eta_1^p-\xi_1\rVert+\cdots+
\lVert U\eta_n^p-\xi_n\rVert\bigr)
\longrightarrow0\qquad(p\to+\infty).
\]
This is evident for $n=0$. Suppose it has been proved when one starts with
$n-1$ vectors.

Let $H'$ be the vector subspace of $H$ spanned by
$\xi_1,\ldots,\xi_{n-1}$, and let $H''=H\ominus H'$. By the induction
hypothesis, there exist $U_1,U_2,\ldots\in G$ such that
\[
\lVert U_p\eta_1^p-\xi_1\rVert+\cdots+
\lVert U_p\eta_{n-1}^p-\xi_{n-1}\rVert\longrightarrow0.
\]
Put
\[
\xi_n'=P_{H'}\xi_n,\quad \xi_n''=P_{H''}\xi_n,\quad
\eta_n^{\prime p}=P_{H'}U_p\eta_n^p,\quad
\eta_n^{\prime\prime p}=P_{H''}U_p\eta_n^p.
\]
For $1\leq i\leq n-1$, we have
\begin{align*}
\bigl|((\xi_n'-\eta_n^{\prime p})\mid\xi_i)\bigr|
&=\bigl|((\xi_n-U_p\eta_n^p)\mid\xi_i)\bigr|\\
&\leq\bigl|(\xi_n\mid\xi_i)-(\eta_n^p\mid\eta_i^p)\bigr|
+\bigl|(U_p\eta_n^p\mid U_p\eta_i^p)
-(U_p\eta_n^p\mid\xi_i)\bigr|\\
&\leq\bigl|(\xi_n\mid\xi_i)-(\eta_n^p\mid\eta_i^p)\bigr|
+\lVert\eta_n^p\rVert\,
\lVert U_p\eta_i^p-\xi_i\rVert\longrightarrow0.
\end{align*}
Therefore $\eta_n^{\prime p}\to\xi_n'$ in the finite-dimensional space
$H'$. At the same time,
\[
\lVert U_p\eta_n^p\rVert^2=\lVert\eta_n^p\rVert^2
\longrightarrow\lVert\xi_n\rVert^2,
\qquad
\lVert\eta_n^{\prime\prime p}\rVert^2
\longrightarrow\lVert\xi_n''\rVert^2.
\]
There is therefore a unitary operator $V_p$ on $H$, acting as the identity
on $H'$, such that $V_p\eta_n^{\prime\prime p}\to\xi_n''$. For
$1\leq i\leq n-1$,
\[
\lVert P_{H''}U_p\eta_i^p\rVert\longrightarrow
\lVert P_{H''}\xi_i\rVert=0,
\qquad
\lVert V_pU_p\eta_i^p-U_p\eta_i^p\rVert\longrightarrow0.
\]
% Source PDF physical page 080
Consequently
\[
\lVert V_pU_p\eta_1^p-\xi_1\rVert+\cdots+
\lVert V_pU_p\eta_{n-1}^p-\xi_{n-1}\rVert\longrightarrow0.
\]
Moreover,
\[
\lVert V_pU_p\eta_n^p-\xi_n\rVert^2
=\lVert\eta_n^{\prime p}-\xi_n'\rVert^2
+\lVert V_p\eta_n^{\prime\prime p}-\xi_n''\rVert^2
\longrightarrow0.
\]
Thus
\[
\inf_{U\in G}\sum_{i=1}^n\lVert U\eta_i^p-\xi_i\rVert
\longrightarrow0\qquad(p\to+\infty).
\]

\ResultHeading{3.5.7}{3.5.7. Lemma.}
\begin{itshape}
Let $H$ be a Hilbert space, $A$ a $C^{*}$-algebra, $\pi_0$ a nonzero
irreducible representation of $A$ on a closed vector subspace $H_0$ of $H$,
and $S$ the set of elements of $\widehat A$ whose dimension is at most
$\dim H$. Let $a_1,\ldots,a_p\in A$, let
$\xi_1,\ldots,\xi_n\in H_0$, let $\varepsilon>0$, and let
$\tau_0\in S$ be the class of $\pi_0$. There exists a neighborhood $V$ of
$\tau_0$ in $S$ with the following property: if $\tau\in V$, there exists a
representation $\pi$ of $A$ on a closed vector subspace $H_\pi$ of $H$ such
that:
\begin{enumerate}[label=\textup{\arabic*\textdegree}]
\item the class of $\pi$ is $\tau$;
\item
$\lVert\pi(a_i)P_{H_\pi}\xi_j-\pi_0(a_i)\xi_j\rVert\leq\varepsilon$ for all
$i,j$;
\item if $\dim\tau=\dim H$, then $H_\pi=H$.
\end{enumerate}
\end{itshape}

We are immediately reduced to the case in which $A$ has an identity element,
$a_1=1$, $\lVert a_i\rVert\leq1$, and $\xi_1\ne0$.

$1^\circ$ Suppose first that $n=1$. Let $\varepsilon_1>0$. There exists a
neighborhood $W$ of $\tau_0$ in $\widehat A$ with the following property: if
$\tau\in W$, there exist a representation $\rho$ of $A$ of class $\tau$ and
a vector $\xi\in H_\rho$ such that
\begin{equation}
\bigl|(\rho(a_j^{\ast}a_i)\xi\mid\xi)
-(\pi_0(a_j^{\ast}a_i)\xi_1\mid\xi_1)\bigr|
\leq\varepsilon_1
\qquad\text{for all }i,j.
\tag{1}\label{eq:3.5.7-1}
\end{equation}
This follows from Theorem \XRef{3.4.11}{3.4.11}. If $\tau\in V=W\cap S$,
we may suppose that $H_\rho$ is a closed vector subspace of $H$, equal to
$H$ if $\dim\tau=\dim H$. For $x\in A$, extend $\rho(x)$ to an operator
$\rho'(x)\in\mathcal L(H)$ which is zero on $H\ominus H_\rho$. Inequality
\eqref{eq:3.5.7-1} may be written
\[
\bigl|(\rho(a_i)\xi\mid\rho(a_j)\xi)
-(\pi_0(a_i)\xi_1\mid\pi_0(a_j)\xi_1)\bigr|
\leq\varepsilon_1.
\]
If $\varepsilon_1$ has been suitably chosen, \XRef{3.5.6}{3.5.6} gives a
unitary operator $U$ on $H$ such that
\[
\lVert U\rho(a_i)\xi-\pi_0(a_i)\xi_1\rVert\leq\frac{\varepsilon}{2}
\qquad\text{for every }i.
\]
In particular, taking $i=1$, we obtain
$\lVert U\xi-\xi_1\rVert\leq\varepsilon/2$. Therefore
\begin{align*}
\lVert U\rho'(a_i)U^{-1}\xi_1-\pi_0(a_i)\xi_1\rVert
&\leq\lVert U\rho'(a_i)U^{-1}\xi_1-U\rho(a_i)\xi\rVert\\
&\quad+\lVert U\rho(a_i)\xi-\pi_0(a_i)\xi_1\rVert\\
&\leq\lVert a_i\rVert\,\lVert\xi_1-U\xi\rVert
+\frac{\varepsilon}{2}\leq\varepsilon.
\end{align*}
We may therefore take for $\pi$ the representation
$x\mapsto U\rho'(x)U^{-1}$, acting on $U(H_\rho)$.

% Source PDF physical page 081
$2^\circ$ Now let $n$ be arbitrary. There exist $b_1,\ldots,b_n\in A$ such
that
\[
\lVert\xi_j-\pi_0(b_j)\xi_1\rVert\leq\frac{\varepsilon}{4}.
\]
By $1^\circ$, there exists a neighborhood $V$ of $\tau_0$ in $S$ with the
following property: if $\tau\in V$, there exists a representation $\pi$ of
$A$, of class $\tau$, on a closed vector subspace of $H$ (equal to $H$ if
$\dim\tau=\dim H$), such that
\[
\begin{gathered}
\lVert\pi(a_ib_j)P_{H_\pi}\xi_1-
\pi_0(a_ib_j)\xi_1\rVert\leq\frac{\varepsilon}{4},\\
\lVert\pi(b_j)P_{H_\pi}\xi_1-
\pi_0(b_j)\xi_1\rVert\leq\frac{\varepsilon}{4}
\end{gathered}
\qquad\text{for all }i,j.
\]
Then
\begin{align*}
&\lVert\pi(a_i)P_{H_\pi}\xi_j-\pi_0(a_i)\xi_j\rVert\\
&\quad\leq
\lVert\pi(a_i)P_{H_\pi}(\xi_j-\pi_0(b_j)\xi_1)\rVert\\
&\qquad+
\lVert\pi(a_i)P_{H_\pi}\pi_0(b_j)\xi_1
-\pi(a_i)\pi(b_j)P_{H_\pi}\xi_1\rVert\\
&\qquad+
\lVert\pi(a_ib_j)P_{H_\pi}\xi_1-\pi_0(a_ib_j)\xi_1\rVert\\
&\qquad+
\lVert\pi_0(a_i)\pi_0(b_j)\xi_1-\pi_0(a_i)\xi_j\rVert\\
&\quad\leq
\frac{\varepsilon}{4}\lVert a_i\rVert
+\frac{\varepsilon}{4}\lVert a_i\rVert
+\frac{\varepsilon}{4}
+\frac{\varepsilon}{4}\lVert a_i\rVert
\leq\varepsilon.
\end{align*}

\ResultHeading{3.5.8}{3.5.8. Theorem.}
\begin{itshape}
Let $A$ be a $C^{*}$-algebra and $n$ a cardinal. The canonical map
from $\operatorname{Irr}_n(A)$ onto $\widehat A_n$ is continuous and open.
\end{itshape}

Let $Q$ be a subset of $\operatorname{Irr}_n(A)$, let
$\pi_0\in\operatorname{Irr}_n(A)$, and let $S$ and $\tau_0$ be their images
in $\widehat A_n$. If $\pi_0$ belongs to the closure of $Q$, every positive
form associated with $\pi_0$ is a weak limit of positive forms associated
with $Q$, and hence $\tau_0$ belongs to the closure of $S$. This proves that
the map $\operatorname{Irr}_n(A)\to\widehat A_n$ is continuous. Moreover,
every neighborhood of $\pi_0$ in $\operatorname{Irr}_n(A)$ contains a
neighborhood constructed as follows: choose
$a_1,\ldots,a_p\in A$, $\xi_1,\ldots,\xi_m\in H_n$, and $\varepsilon>0$,
and take the set of $\pi\in\operatorname{Irr}_n(A)$ such that
\[
\lVert\pi(a_i)\xi_j-\pi_0(a_i)\xi_j\rVert\leq\varepsilon
\qquad\text{for all }i,j.
\]
By \XRef{3.5.7}{3.5.7}, the image of such a neighborhood in $\widehat A_n$
is a neighborhood of $\tau_0$ in $\widehat A_n$. This proves that the map
$\operatorname{Irr}_n(A)\to\widehat A_n$ is open.

Thus the topology of $\widehat A_n$ is the quotient topology of the topology
of $\operatorname{Irr}_n(A)$ under the canonical map
$\operatorname{Irr}_n(A)\to\widehat A_n$.

\ResultHeading{3.5.9}{3.5.9. Proposition.}
\begin{itshape}
Let $A$ be a $C^{*}$-algebra and $x\in A^+$. For every
$\pi\in\widehat A$, $\operatorname{Tr}\pi(x)$ is a nonnegative number,
finite or not. The function
$\pi\mapsto\operatorname{Tr}\pi(x)$ is lower semicontinuous on
$\widehat A$.
\end{itshape}

Let $H$ be a Hilbert space such that every irreducible representation of $A$
can be realized on a closed vector subspace of $H$
(\XRef{2.3.3}{2.3.3}). Let $\pi_0\in\widehat A$ be realized on a closed
vector subspace $H_0$ of $H$. Let
$\alpha<\operatorname{Tr}\pi_0(x)$, and choose an orthonormal system
$(\xi_1,\ldots,\xi_n)$ in $H_0$ such that
\[
\sum_{j=1}^n(\pi_0(x)\xi_j\mid\xi_j)-\alpha=\beta>0.
\]
% Source PDF physical page 082
By \XRef{3.5.7}{3.5.7}, there exists a neighborhood $V$ of $\pi_0$ in
$\widehat A$ such that every $\pi\in V$ can be realized on a closed vector
subspace of $H$ in such a way that
\[
\lVert\pi(x)P_{H_\pi}\xi_j-\pi_0(x)\xi_j\rVert
\leq\frac{\beta}{n}
\qquad\text{for every }j.
\]
Then
\[
\operatorname{Tr}\pi(x)
\geq\sum_{j=1}^n(\pi(x)P_{H_\pi}\xi_j\mid\xi_j)
\geq\sum_{j=1}^n(\pi_0(x)\xi_j\mid\xi_j)-n\frac{\beta}{n}
=\alpha.
\]

Bibliography: \cite{ref14,ref27,ref28}.

\BookSubsection{3.6}{Finite-Dimensional Representations}

\ResultHeading{3.6.1}{3.6.1.}
Let $k$ be a commutative field, let $n$ be an integer $>0$, and let $M_n(k)$
be the algebra of $n$-by-$n$ matrices over $k$. This algebra has dimension
$n^2$. Hence, if $r>n^2$, then
\begin{equation}
\sum_{\sigma\in\mathfrak S_r}\varepsilon_\sigma
X_{\sigma(1)}X_{\sigma(2)}\cdots X_{\sigma(r)}=0
\tag{1}\label{eq:3.6.1-1}
\end{equation}
for all $X_1,\ldots,X_r\in M_n(k)$, where $\mathfrak S_r$ denotes the
symmetric group of $\{1,2,\ldots,r\}$ and $\varepsilon_\sigma$ is the sign of
the permutation $\sigma$. Let $r(n)$ be the least integer $r$ for which
identity \eqref{eq:3.6.1-1} holds in $M_n(k)$.

\ResultHeading{3.6.2}{3.6.2. Lemma.}
\begin{itshape}
One has $r(n)\geq r(n-1)+2$.
\end{itshape}

Let $t=r(n-1)-1$. Choose $X_1,\ldots,X_t\in M_{n-1}(k)$ such that
\[
Y=\sum_{\sigma\in\mathfrak S_t}\varepsilon_\sigma
X_{\sigma(1)}\cdots X_{\sigma(t)}\ne0.
\]
Write $Y=(y_{ij})$; there exist $h,l$ such that $y_{hl}\ne0$. Let $X_i'$
(respectively $Y'$) be the matrix obtained from $X_i$ (respectively $Y$) by
adjoining an $n$th row and an $n$th column of zeros. Let $X'_{t+1}$
(respectively $X'_{t+2}$) be the matrix in $M_n(k)$ all of whose entries are
zero except the entry in the $l$th row and $n$th column (respectively the
entry in the $n$th row and $n$th column), which is equal to $1$. In the sum
\begin{equation}
\sum_{\tau\in\mathfrak S_{t+2}}\varepsilon_\tau
X'_{\tau(1)}X'_{\tau(2)}\cdots X'_{\tau(t+2)},
\tag{1}\label{eq:3.6.2-1}
\end{equation}
all terms in which $X'_{t+2}$ is not in the last position vanish (because
$X'_{t+2}X'_j=0$ for $j<t+2$); among the terms ending in $X'_{t+2}$, those
in which $X'_{t+2}$ is not preceded by $X'_{t+1}$ vanish (because
$X'_jX'_{t+2}=0$ for $j<t+1$). The sum \eqref{eq:3.6.2-1} is therefore equal
to
\[
Y'X'_{t+1}X'_{t+2}=Y'X'_{t+1},
\]
and this matrix is nonzero because $y_{hl}\ne0$. Hence
$r(n)>t+2=r(n-1)+1$.

% Source PDF physical page 083
\ResultHeading{3.6.3}{3.6.3. Proposition.}
\begin{itshape}
Let $A$ be a $C^{*}$-algebra and $n$ an integer $>0$. Let
${}_{n}\widehat A$ be the set of $\pi\in\widehat A$ such that
$\dim\pi\leq n$, and let $\widehat A_n$ be the set of
$\pi\in\widehat A$ such that $\dim\pi=n$.
\begin{enumerate}[label=\textup{(\roman*)}]
\item ${}_{n}\widehat A$ is closed in $\widehat A$, and $\widehat A_n$ is
open in ${}_{n}\widehat A$.
\item Let $I_n$ be the intersection of the kernels of the
$\pi\in{}_{n}\widehat A$. Then ${}_{n}\widehat A$ is canonically
homeomorphic to the spectrum of $A/I_n$, and $\widehat A_n$ is canonically
homeomorphic to the spectrum of $I_{n-1}/I_n$.
\end{enumerate}
\end{itshape}

Let $l=r(n)$ and $x_1,\ldots,x_l\in A$. The element
\[
\sum_{\sigma\in\mathfrak S_l}\varepsilon_\sigma
x_{\sigma(1)}\cdots x_{\sigma(l)}
\]
belongs to the kernel of every representation in ${}_{n}\widehat A$, and
hence to $I_n$. Let $\pi\in\widehat A$ satisfy $\ker\pi\supset I_n$. Then
\[
\sum_{\sigma\in\mathfrak S_l}\varepsilon_\sigma
\pi(x_{\sigma(1)})\cdots\pi(x_{\sigma(l)})
=\pi\left(\sum_{\sigma\in\mathfrak S_l}\varepsilon_\sigma
x_{\sigma(1)}\cdots x_{\sigma(l)}\right)=0
\]
for all $x_1,\ldots,x_l\in A$. Since $\pi(A)$ is strongly dense in
$\mathcal L(H_\pi)$, it follows that
\[
\sum_{\sigma\in\mathfrak S_l}\varepsilon_\sigma
u_{\sigma(1)}\cdots u_{\sigma(l)}=0
\]
for all $u_1,\ldots,u_l\in\mathcal L(H_\pi)$. By \XRef{3.6.2}{3.6.2},
$H_\pi$ contains no finite-dimensional vector subspace of dimension $>n$;
hence $\dim H_\pi\leq n$, and $\pi\in{}_{n}\widehat A$. Thus
${}_{n}\widehat A$ is closed in $\widehat A$. Consequently
$\widehat A_n={}_{n}\widehat A\setminus{}_{n-1}\widehat A$ is open in
${}_{n}\widehat A$. Assertion (ii) follows from (i) and
\XRef{3.2.1}{3.2.1}.

\ResultHeading{3.6.4}{3.6.4. Proposition.}
\begin{itshape}
\begin{enumerate}[label=\textup{(\roman*)}]
\item The topology of $\widehat A_n$ is locally compact.
\item This topology is the coarsest topology making all the functions
$\pi\mapsto\operatorname{Tr}\pi(x)$ on $\widehat A_n$, $x\in A$,
continuous.
\end{enumerate}
\end{itshape}

Let $H_n$ be the model Hilbert space of dimension $n$. For every $x\in A$,
the map $\pi\mapsto\pi(x)$ is continuous on $\operatorname{Irr}_n(A)$; since
$n<+\infty$, it follows that $\pi\mapsto\operatorname{Tr}\pi(x)$ is
continuous on $\operatorname{Irr}_n(A)$. Since it takes the same value on
equivalent representations, it induces a continuous function $\varphi_x$ on
$\widehat A_n$ (\XRef{3.5.8}{3.5.8}). Moreover, if $\pi$ and $\pi'$ are two
inequivalent irreducible representations of $A$ on $H_n$, there exists
$x\in A$ such that
$\operatorname{Tr}\pi(x)\ne\operatorname{Tr}\pi'(x)$
(\XRef{2.8.3}{2.8.3} or \XRef{B.27}{B 27}). Thus the functions $\varphi_x$
separate the points of $\widehat A_n$. This first proves that the topology of
$\widehat A_n$ is Hausdorff. On the other hand, $\widehat A_n$ is identified
with the spectrum of a $C^{*}$-algebra [\XRef{3.6.3}{3.6.3}(ii)], and is
therefore locally compact (\XRef{3.3.8}{3.3.8}). Let
$Z=\widehat A_n\cup\{0\}$ be the Alexandroff compactification of
$\widehat A_n$. By \XRef{3.3.7}{3.3.7}, the functions $\varphi_x$ extend to
continuous functions $\psi_x$ on $Z$ that take the value $0$ at the point
$0$. The functions $\psi_x$ separate the points of $Z$; indeed,
% Source PDF physical page 084
by what precedes, it is enough to prove that, for every
$\pi\in\widehat A_n$, there exists $x\in A$ such that $\psi_x(\pi)\ne0$.
But if $y\in A$ satisfies $\pi(y)\ne0$, then
\[
\psi_{y^{\ast}y}(\pi)
=\operatorname{Tr}\pi(y^{\ast}y)
=\operatorname{Tr}\bigl(\pi(y)^{\ast}\pi(y)\bigr)>0.
\]
Thus the topology on $Z$ defined by the functions $\psi_x$ is Hausdorff and
coarser than the topology of $Z$, and hence is identical with it. This proves
(ii).

Bibliography: \cite{ref27,ref75,ref77}.

\BookSubsection[Complements on the Spaces Rep-n(A)]{3.7}{Complements on the Spaces $\operatorname{Rep}_n(A)$}

\ResultHeading{3.7.1}{3.7.1. Proposition.}
\begin{itshape}
Let $A$ be a separable $C^{*}$-algebra and $n$ a cardinal
$\leq\aleph_0$. The topological space $\operatorname{Rep}_n(A)$ is Polish.
\end{itshape}

The model Hilbert space $H_n$ of dimension $n$ is separable. Let
$\mathcal L_s(H_n)$ be $\mathcal L(H_n)$ equipped with the strong topology.
It is quasi-complete (\XRef{B.11}{B 11}). Let
$\mathcal L_s(A,\mathcal L_s(H_n))$ be the space of continuous linear maps
from $A$ into $\mathcal L_s(H_n)$, equipped with the topology of pointwise
convergence. Every closed equicontinuous subset of
$\mathcal L_s(A,\mathcal L_s(H_n))$ is a complete uniform subspace of
$\mathcal L_s(A,\mathcal L_s(H_n))$ (\XRef{B.12}{B 12}). In particular, the
set of linear maps $\pi$ from $A$ into $\mathcal L_s(H_n)$ such that
$\lVert\pi(x)\rVert\leq\lVert x\rVert$ for every $x\in A$ is a complete
uniform subspace $B$ of $\mathcal L_s(A,\mathcal L_s(H_n))$.

Let $(x_i)_{i\in I}$ be a dense sequence in $A$ such that
$\lVert x_i\rVert\leq1$ for every $i$. Let $(\xi_j)_{j\in J}$ be a dense
sequence in $H_n$ such that $\lVert\xi_j\rVert\leq1$ for every $j$. The
uniform structure of $B$ has as a fundamental system of entourages the
entourages defined by the inequalities
\[
\lVert\pi(x_i)\xi_j-\pi'(x_i)\xi_j\rVert\leq\frac1k
\qquad(i\in I,\ j\in J,\ k=1,2,\ldots).
\]
This uniform structure is therefore metrizable. The map
\[
\pi\longmapsto(\pi(x_i)\xi_j)_{i\in I,j\in J}
\]
from $B$ onto a subspace of $H_n^{I\times J}$ is bicontinuous, and hence $B$
is separable. In short, $B$, equipped with the topology of pointwise strong
convergence, is a Polish space. Finally, $\operatorname{Rep}_n(A)$ is the
set of $\pi\in B$ such that
\[
\pi(xy)=\pi(x)\pi(y),\qquad \pi(x^{\ast})=\pi(x)^{\ast}
\qquad(x,y\in A).
\]
The second condition is equivalent to
\[
(\pi(x^{\ast})\xi\mid\eta)=(\xi\mid\pi(x)\eta)
\qquad(x\in A,\ \xi,\eta\in H_n).
\]
Thus $\operatorname{Rep}_n(A)$ is a closed subspace of $B$, and hence is a
Polish space.

\ResultHeading{3.7.2}{3.7.2. Lemma.}
\begin{itshape}
Let $E,F$ be two complex Banach spaces with $F$ separable, let $M$ be a
topological space, and let $m\mapsto U_m$ be a continuous map from $M$ into
$\mathcal L_s(E,F)$ (the space of continuous linear maps from $E$ into $F$,
equipped with the topology of pointwise convergence). Let $n$ be an integer.
The set of $m\in M$ such that
$\operatorname{codim}\overline{U_m(E)}\leq n$ is a $G_\delta$ subset of $M$.
\end{itshape}

% Source PDF physical page 085
Let $f_1,f_2,\ldots,f_k$ be elements of $F$. Let $R$ be the set of sequences
$r=(r_1,r_2,\ldots,r_k)$ of complex numbers such that
\[
\frac12\leq |r_1|^2+\cdots+|r_k|^2\leq1.
\]
For every $r\in R$, every $e\in E$, and every integer $p>0$, the set
$M(r,e,p)$ of $m\in M$ such that
\[
\lVert r_1f_1+\cdots+r_kf_k-U_m(e)\rVert<\frac1p
\]
is open in $M$. Now $f_1,\ldots,f_k$ are linearly dependent modulo
$\overline{U_m(E)}$ if and only if, for every $p$, there exist $r\in R$ and
$e\in E$ such that
\[
\lVert r_1f_1+\cdots+r_kf_k-U_m(e)\rVert<\frac1p;
\]
in other words, if and only if
\[
m\in\bigcap_{p=1}^{\infty}\ \bigcup_{\substack{r\in R\\e\in E}}M(r,e,p).
\]
[Indeed, if this is so, we may suppose that, as $p\to+\infty$,
$(r_1,\ldots,r_k)$ tends to a limit. Then $U_m(e)$ has a limit which belongs
to $\overline{U_m(E)}$ and is a linear combination, with coefficients not all
zero, of $f_1,\ldots,f_k$. Conversely, if $f_1,\ldots,f_k$ are linearly
dependent modulo $\overline{U_m(E)}$, then for every $p$ there exist
$r=(r_1,\ldots,r_k)\in R$ and $e\in E$ such that
\[
\lVert r_1f_1+\cdots+r_kf_k-U_m(e)\rVert<\frac1p,
\]
which proves the assertion.]

Thus the set of points $m\in M$ for which $f_1,\ldots,f_k$ are linearly
dependent modulo $\overline{U_m(E)}$ is a $G_\delta$ subset of $M$. Let
$(g_1,g_2,\ldots)$ be a dense sequence in $F$. We have
$\operatorname{codim}\overline{U_m(E)}\leq n$ if and only if
$g_{i_1},g_{i_2},\ldots,g_{i_{n+1}}$ are linearly dependent modulo
$\overline{U_m(E)}$ for all indices $i_1,\ldots,i_{n+1}$. Since the family of
systems $(i_1,\ldots,i_{n+1})$ is countable, the set in Lemma
\XRef{3.7.2}{3.7.2} appears as a countable intersection of $G_\delta$ sets,
and hence is a $G_\delta$ set.

\ResultHeading{3.7.3}{3.7.3. Lemma.}
\begin{itshape}
Let $A$ be a $C^{*}$-algebra, let $n$ be a cardinal $\leq\aleph_0$, let $M$
be a topological space, let $m\mapsto\pi_m$ and $m\mapsto\pi'_m$ be two
continuous maps from $M$ into $\operatorname{Rep}_n(A)$, and let $p$ be an
integer $\geq0$. The set of $m\in M$ for which the intertwining number of
$\pi_m$ and $\pi'_m$ is at most $p$ is a $G_\delta$ subset of $M$.
\end{itshape}

Let $H_n$ be the model Hilbert space of dimension $n$, let $K$ be the unit
ball of $H_n$, and let $B$ be the unit ball of $A$. Denote by $Y$ the Banach
space of families $(\mu_{ijk})_{i\in B,j\in K,k\in K}$ of complex numbers
such that
% Source PDF physical page 086
\[
\lVert(\mu_{ijk})\rVert
=\sum_{\substack{i\in B\\j\in K,\ k\in K}}|\mu_{ijk}|<+\infty.
\]
Its dual is the Banach space $X$ of families
$(\lambda_{ijk})_{i\in B,j\in K,k\in K}$ of complex numbers such that
\[
\lVert(\lambda_{ijk})\rVert=\sup_{i,j,k}|\lambda_{ijk}|<+\infty.
\]
We denote by $P$ the predual of the von Neumann algebra $\mathcal L(H_n)$
(\XRef{A.23}{A 23}). It is separable.

For every $m\in M$, let $V_m$ be the continuous linear map from
$\mathcal L(H_n)$ into $X$ defined by
\[
V_m(T)=
\bigl(((\pi_m(i)T-T\pi'_m(i))j\mid k)\bigr)_{i\in B,j\in K,k\in K}.
\]
The intertwining number of $\pi_m$ and $\pi'_m$ is then the dimension of
$\ker V_m$. For $T\in\mathcal L(H_n)$ and $(\mu_{ijk})\in Y$, writing
$x_i,\xi_j$ in place of $i,j$, we have
\begin{align*}
\langle(\mu_{ijk}),V_m(T)\rangle
&=\sum_{ijk}\mu_{ijk}
\bigl((\pi_m(x_i)T-T\pi'_m(x_i))\xi_j\mid\xi_k\bigr)\\
&=\sum_{ijk}\mu_{ijk}
\left((T\xi_j\mid\pi_m(x_i^{\ast})\xi_k)
-(T\pi'_m(x_i)\xi_j\mid\xi_k)\right).
\end{align*}
Given two vectors $\eta,\zeta\in H_n$, denote by $\varphi(\eta,\zeta)$ the
element of $P$ defined by
\[
\langle\varphi(\eta,\zeta),T\rangle=(T\eta\mid\zeta);
\qquad
\lVert\varphi(\eta,\zeta)\rVert=\lVert\eta\rVert\,\lVert\zeta\rVert.
\]
Then
\[
\langle(\mu_{ijk}),V_m(T)\rangle
=\left\langle
\sum_{ijk}\mu_{ijk}
\bigl(\varphi(\xi_j,\pi_m(x_i^{\ast})\xi_k)
-\varphi(\pi'_m(x_i)\xi_j,\xi_k)\bigr),T
\right\rangle.
\]
Thus $V_m$ is the transpose of a continuous linear map $U_m$ from $Y$ into
$P$, defined by
\begin{equation}
U_m((\mu_{ijk}))=
\sum_{ijk}\mu_{ijk}
\bigl(\varphi(\xi_j,\pi_m(x_i^{\ast})\xi_k)
-\varphi(\pi'_m(x_i)\xi_j,\xi_k)\bigr),
\tag{1}\label{eq:3.7.3-1}
\end{equation}
and the intertwining number of $\pi_m$ and $\pi'_m$ is
$\operatorname{codim}\overline{U_m(Y)}$. By \XRef{3.7.2}{3.7.2}, to prove
\XRef{3.7.3}{3.7.3} it is enough to prove that the map
$m\mapsto U_m$ from $M$ into $\mathcal L_s(Y,P)$ is continuous. For fixed
$(\mu_{ijk})\in Y$, we therefore show that the element
\eqref{eq:3.7.3-1} of $P$ depends continuously on $m$ in the norm of $P$.
Since
\[
\left\lVert\mu_{ijk}
\bigl(\varphi(\xi_j,\pi_m(x_i^{\ast})\xi_k)
-\varphi(\pi'_m(x_i)\xi_j,\xi_k)\bigr)\right\rVert
\leq2|\mu_{ijk}|,
\]
it is enough to prove that each term
$\varphi(\xi_j,\pi_m(x_i^{\ast})\xi_k)$ and
$\varphi(\pi'_m(x_i)\xi_j,\xi_k)$ depends continuously on $m$. But the maps
$m\mapsto\pi_m(x_i^{\ast})\xi_k$ and
$m\mapsto\pi'_m(x_i)\xi_j$ from $M$ into $H_n$ are continuous by the
definition of the topology
% Source PDF physical page 087
of $\operatorname{Rep}_n(A)$, and the map
$(\eta,\zeta)\mapsto\varphi(\eta,\zeta)$ from $H_n\times H_n$ into $P$ is
continuous by the equality
$\lVert\varphi(\eta,\zeta)\rVert=\lVert\eta\rVert\,\lVert\zeta\rVert$.

\ResultHeading{3.7.4}{3.7.4. Proposition.}
\begin{itshape}
Let $A$ be a separable $C^{*}$-algebra and $n$ a cardinal
$\leq\aleph_0$. The topological space $\operatorname{Irr}_n(A)$ is Polish.
\end{itshape}

Apply \XRef{3.7.3}{3.7.3} with $M=\operatorname{Rep}_n(A)$,
$\pi_\rho=\pi'_\rho=\rho$ for every
$\rho\in\operatorname{Rep}_n(A)$, and $p=1$. If $n>0$, the set thus obtained
in $M=\operatorname{Rep}_n(A)$ is $\operatorname{Irr}_n(A)$, possibly with
the zero representation added when $n=1$. Thus $\operatorname{Irr}_n(A)$ is
a $G_\delta$ subset of $\operatorname{Rep}_n(A)$, which is Polish
(\XRef{3.7.1}{3.7.1}). Hence $\operatorname{Irr}_n(A)$ is Polish
(\XRef{B.15}{B 15}).

\ResultHeading{3.7.5}{3.7.5.}
Apply \XRef{3.7.3}{3.7.3} with $M=\operatorname{Rep}_n(A)$,
$\pi_\rho=\rho$, $\pi'_\rho$ the zero representation of $A$ on $H_n$, and
$p=0$. We obtain the set of nondegenerate representations of $A$ on $H_n$.
If $A$ is a separable $C^{*}$-algebra, this set is a $G_\delta$ subset of
$\operatorname{Rep}_n(A)$, and hence is a Polish space.

Bibliography: \cite{ref12,ref89}.

\BookSubsection{3.8}{Mackey Borel Structure}

\ResultHeading{3.8.1}{3.8.1.}
Let $A$ be a separable $C^{*}$-algebra. For every cardinal $n\leq\aleph_0$,
we equip $\operatorname{Rep}_n(A)$ and $\operatorname{Irr}_n(A)$ with the
Borel structures (\XRef{B.19}{B 19}) underlying their topologies. These are
the coarsest Borel structures for which the functions
\[
\pi\longmapsto(\pi(x)\xi\mid\eta)
\qquad(x\in A,\ \xi,\eta\in H_n)
\]
are Borel. Let $\operatorname{Rep}(A)$ (respectively
$\operatorname{Irr}(A)$) be the union of the
$\operatorname{Rep}_n(A)$ (respectively $\operatorname{Irr}_n(A)$), for
$n=1,2,\ldots,\aleph_0$, equipped with the Borel sum structure of the
structures on the $\operatorname{Rep}_n(A)$ (respectively
$\operatorname{Irr}_n(A)$). The Borel space $\operatorname{Irr}(A)$ is a
Borel subset of $\operatorname{Rep}(A)$ by the proof of
\XRef{3.7.4}{3.7.4}; and the structure induced on $\operatorname{Irr}(A)$ by
the Borel structure of $\operatorname{Rep}(A)$ is the Borel structure of
$\operatorname{Irr}(A)$. By \XRef{3.7.1}{3.7.1} and
\XRef{3.7.4}{3.7.4}, the Borel spaces $\operatorname{Rep}(A)$ and
$\operatorname{Irr}(A)$ are standard (\XRef{B.20}{B 20}).

\ResultHeading{3.8.2}{3.8.2.}
Since $A$ is separable, the canonical map
$\operatorname{Irr}(A)\to\widehat A$ that assigns to every element of
$\operatorname{Irr}(A)$ its class is surjective (\XRef{2.3.3}{2.3.3}).

\ResultLabel{Definition.}
\begin{itshape}
Let $A$ be a separable $C^{*}$-algebra. The Mackey Borel structure on
$\widehat A$ is the quotient structure of the Borel structure of
$\operatorname{Irr}(A)$ under the canonical map
$\operatorname{Irr}(A)\to\widehat A$.
\end{itshape}

This Borel structure is not standard in general (see, however,
\XRef{4.6.1}{4.6.1}).

% Source PDF physical page 088
\ResultHeading{3.8.3}{3.8.3. Proposition.}
\begin{itshape}
Let $A$ be a separable $C^{*}$-algebra. The Mackey Borel structure
on $\widehat A$ is finer than the Borel structure underlying the topology of
$\widehat A$.
\end{itshape}

Let $\Omega$ be an open subset of $\widehat A$. For
$n=1,2,\ldots,\aleph_0$, the inverse image of $\Omega$ in
$\operatorname{Irr}_n(A)$ is open (\XRef{3.5.8}{3.5.8}), and hence Borel.
Thus the inverse image of $\Omega$ in $\operatorname{Irr}(A)$ is Borel, so
$\Omega$ is Borel for the Mackey structure. It follows that every subset of
$\widehat A$ which is Borel for the topology of $\widehat A$ is Borel for
the Mackey structure.

\ResultHeading{3.8.4}{3.8.4. Proposition.}
\begin{itshape}
Let $A$ be a separable $C^{*}$-algebra. Every point of $\widehat A$
is Borel for the Mackey structure.
\end{itshape}

Let $\pi_0\in\operatorname{Irr}_n(A)$ and let $\tau_0$ be its image in
$\widehat A$. The set of $\pi\in\operatorname{Irr}_n(A)$ for which the
intertwining number of $\pi_0$ and $\pi$ is at most $0$ is a $G_\delta$
subset of $\operatorname{Irr}_n(A)$ (\XRef{3.7.3}{3.7.3}). Hence, by
\XRef{2.3.4}{2.3.4}, the set of $\pi\in\operatorname{Irr}_n(A)$ inequivalent
to $\pi_0$ is a $G_\delta$ subset of $\operatorname{Irr}_n(A)$. Therefore the
equivalence class of $\pi_0$ in $\operatorname{Irr}_n(A)$ is Borel, and so
$\{\tau_0\}$ is Borel in $\widehat A$ for the Mackey structure.

Bibliography: \cite{ref28,ref91}.

\BookSubsection{3.9}{Complements}

\ResultHeading{3.9.1}{3.9.1.}
Let $A$ be a $C^{*}$-algebra.
\begin{enumerate}[label=\textit{\alph*.}]
\item If the zero ideal is primitive, there is a dense point in $\widehat A$.
In particular, any two nonempty open subsets of $\widehat A$ intersect.
\item If the zero ideal is not primitive and $A$ is separable, there exist
two disjoint nonempty open subsets of $\widehat A$. (Use
\XRef{3.3.2}{3.3.2} and \XRef{3.4.13}{3.4.13}.)
\item If $A$ is separable and $I$ is a closed two-sided ideal of $A$, it is
equivalent to say that $I$ is primitive or that $I$ is prime
(\XRef{1.9.13}{1.9.13}). Problem: can the separability hypothesis be removed
in \emph{b} and \emph{c}?
\item[*\textit{d}.] There exist separable $C^{*}$-algebras $A$ such that the
zero ideal is primitive and such that, in $\widehat A$, every nonempty open
subset is non-Hausdorff. \cite{ref11}
\end{enumerate}

\ResultHeading{3.9.2}{3.9.2.}
Let $A$ be a $C^{*}$-algebra.
\begin{enumerate}[label=\textit{\alph*.}]
\item Let $\mathfrak N(A)$ be the compact space of $C^{*}$-seminorms
(\XRef{1.9.13}{1.9.13}). Let $\Phi$ be the canonical map
$I\mapsto N_I$ (\XRef{1.9.13}{1.9.13}). Equip $\operatorname{Prim}(A)$ with
the Borel structure underlying its topology. Then $\Phi$ is a Borel
isomorphism from $\operatorname{Prim}(A)$ onto a subset of $\mathfrak N(A)$.
\item If $A$ is separable, $\operatorname{Prim}(A)$ is a standard Borel
space. (Use \XRef{1.9.13}{1.9.13} and \XRef{3.9.1}{3.9.1}.)
\item If $A$ is separable, the Borel structure of $\operatorname{Prim}(A)$ is
the quotient structure of that of $\operatorname{Rep}(A)$ under the
canonical map
% Source PDF physical page 089
$\operatorname{Rep}(A)\to\operatorname{Prim}(A)$. (Use \emph{b} and
\XRef{B.22}{B 22}.) We may therefore speak of the Borel structure of
$\operatorname{Prim}(A)$ without further qualification. \cite{ref20}
\end{enumerate}

\ResultHeading{3.9.3}{3.9.3.}
Let $X$ be a locally quasi-compact space, and let $\mathcal F(X)$ be the set
of closed subsets of $X$. For every quasi-compact subset $C$ of $X$ and
every finite set $\mathcal O$ of nonempty open subsets of $X$, let
$U(C,\mathcal O)$ be the set of $Y\in\mathcal F(X)$ such that
\[
Y\cap C=\varnothing,
\qquad
Y\cap B\ne\varnothing\quad\text{for every }B\in\mathcal O.
\]
Then the sets $U(C,\mathcal O)$ form a base for a topology on
$\mathcal F(X)$, and $\mathcal F(X)$ thereby becomes a compact space.

Now let $A$ be a $C^{*}$-algebra. For every
$F\in\mathcal F(\widehat A)$, let $N_F$ be the function on $A$ defined by
\[
N_F(x)=\sup_{\pi\in F}\lVert\pi(x)\rVert.
\]
This is a $C^{*}$-seminorm on $A$ (\XRef{1.9.13}{1.9.13}). The map
$F\mapsto N_F$ is a homeomorphism from $\mathcal F(\widehat A)$ onto the
compact space $\mathfrak N(A)$ of \XRef{1.9.13}{1.9.13}. \cite{ref29}

\ResultHeading{3.9.4}{3.9.4.}
Let $X$ be a topological space. A point $b$ of $X$ is said to be
\emph{separated in $X$} if, for every point $b'$ of $X$ that is not adherent
to $b$, the points $b$ and $b'$ have disjoint neighborhoods. Let $A$ be a
$C^{*}$-algebra.
\begin{enumerate}[label=\textit{\alph*.}]
\item For $\pi_0\in\widehat A$ to be separated in $\widehat A$, it is
necessary and sufficient that, for every $x\in A$, the function
$\pi\mapsto\lVert\pi(x)\rVert$ on $\widehat A$ be continuous at $\pi_0$.
\item If $A$ is separable, the set of separated points of $\widehat A$ is a
dense $G_\delta$ subset of $\widehat A$. (Use \emph{a} and
\XRef{B.18}{B 18}.) \cite{ref14}
\end{enumerate}

\ResultHeading{3.9.5}{*3.9.5.}
Let $A$ be a $C^{*}$-algebra, and let $\widetilde A$ be the $C^{*}$-algebra
obtained from $A$ by adjoining an identity element $1$. For every
representation $\pi$ of $A$, let $\widetilde\pi$ be the canonical extension
of $\pi$ to $\widetilde A$. If $I$ is a closed two-sided ideal of $A$, let
$I_1$ be the set of $x+\lambda1$, with $x\in A$ and $\lambda\in\mathbb C$,
such that
\[
xy+\lambda y\in I\qquad\text{for every }y\in A.
\]
Then $I_1$ is a closed two-sided ideal of $\widetilde A$ such that
$A\cap I_1=I$. If $I=\ker\pi$ and $\pi$ is nondegenerate, then
$I_1=\ker\widetilde\pi$. Moreover, let $(\pi_\alpha)$ be a family of
nondegenerate representations of $A$. Then $\pi$ is weakly contained in the
set of the $\pi_\alpha$ if and only if $\widetilde\pi$ is weakly contained in
the set of the $\widetilde\pi_\alpha$. \cite{ref27}

\ResultHeading{3.9.6}{*3.9.6.}
Let $A$ be a $C^{*}$-algebra, let $n$ be an integer $>0$, and let
$(\pi_i)_{i\in I}$ be a family of elements of $\widehat A$ of dimension $n$,
indexed by a directed set. Let $\beta_1,\ldots,\beta_r$ be points of
$\widehat A$. Suppose that $\{\beta_1,\ldots,\beta_r\}$ is the set of limits
of all subnets of $(\pi_i)$. Then there exist a directed subset $J$ of $I$
and integers $m_1,\ldots,m_r>0$ such that
\[
\sum_{k=1}^r m_k\dim\beta_k\leq n,
\qquad
\lim_{j\in J}\operatorname{Tr}\pi_j(x)
=\sum_{k=1}^r m_k\operatorname{Tr}\beta_k(x)
\quad\text{for every }x\in A.
\]
\cite{ref27}

\ResultHeading{3.9.7}{3.9.7.}
Let $B$ be a $C^{*}$-algebra and $A$ a $C^{*}$-subalgebra of $B$ such that,
for every $\pi\in\widehat B$, $\pi|_A$ is irreducible. The relation
$\pi_1|_A\simeq\pi_2|_A$ on $\widehat B$ is an equivalence relation $R$.
Then $\widehat A$ is canonically identified with the space
$\widehat B/R$. \cite{ref29}
% ===== ASSEMBLED TRANSLATION CHUNK 04: chunk_090_119.tex =====
% Source PDF physical page 090
\ResultHeading{3.9.8}{3.9.8.}
Let $A$ be a $C^{*}$-algebra, let $H$ be a Hilbert space, and let
$\operatorname{Rep}(A,H)$ be the set of representations of $A$ on $H$.
For $\pi\in\operatorname{Rep}(A,H)$, $\varepsilon>0$,
$a_1,\ldots,a_p\in A$, and $\xi_1,\ldots,\xi_n$ in the essential subspace
of $\pi$, let
$V(\varepsilon,a_1,\ldots,a_p,\xi_1,\ldots,\xi_n,\pi)$ be the set of
$\pi'\in\operatorname{Rep}(A,H)$ such that
\[
 \lVert\pi'(a_i)\xi_j-\pi(a_i)\xi_j\rVert\leq\varepsilon
 \qquad\text{for every $i$ and every $j$.}
\]
Let $\mathcal T$ be the topology on $\operatorname{Rep}(A,H)$ for which the
sets just defined form a fundamental system of neighborhoods of $\pi$.
This topology is not
Hausdorff, but it induces a Hausdorff topology on the set of nondegenerate
representations of $A$ on $H$. We say that
$\pi\in\operatorname{Rep}(A,H)$ and
$\pi'\in\operatorname{Rep}(A,H)$ are almost equivalent if there exists a
partially isometric operator $U$, whose initial and final subspaces are the
essential subspaces of $\pi$ and $\pi'$, respectively, and such that
\[
 U\pi(x)=\pi'(x)U\qquad\text{for every }x\in A.
\]
Let $Q$ be the quotient space of $\operatorname{Rep}(A,H)$ by this
equivalence relation. The set $Q$ is canonically identified with the set of
classes of nondegenerate representations of $A$ whose dimension is at most
$\dim H$. Let $\pi\in Q$ and $S\subset Q$. For $\pi$ to be weakly contained
in $S$, it is necessary and sufficient that $\pi$ belong to the closure of
the set of Hilbert sums of finitely many elements of $S$. If $\dim H$ has
been chosen sufficiently large, then $\widehat A\subset Q$; the topology
induced on $\widehat A$ by that of $Q$ is the topology studied in this
section. \cite{ref31}

\ResultHeading{3.9.9}{3.9.9.}
Let $A$ be a $C^{*}$-algebra and $G$ a topological group. Let $\omega$ be a
homomorphism of $G$ into the group of automorphisms of $A$, continuous for
the topology of pointwise convergence. For every $\pi\in\widehat A$ and
every $s\in G$, let $s\cdot\pi\in\widehat A$ be the representation
\[
 x\longmapsto\pi\bigl(\omega(s)^{-1}x\bigr)
\]
of $A$. Then $\pi\mapsto s\cdot\pi$ is a homeomorphism of $\widehat A$ onto
$\widehat A$, and the map $(s,\pi)\mapsto s\cdot\pi$ from
$G\times\widehat A$ into $\widehat A$ is continuous. \cite{ref47}

\ResultHeading{3.9.10}{3.9.10.}
Let $H=L^2([0,1])$, and let $A$ be the $C^{*}$-algebra of continuous complex
functions on $[0,1]$. For every $f\in A$, let $\pi(f)$ be the operator of
multiplication by $f$ on $H$. Then $\pi$ is a representation of $A$ on
$H$; every character of $A$ is weakly contained in $\pi$, but no character
of $A$ is contained in $\pi$.

\ResultHeading{3.9.11}{3.9.11.}
Let $A$ be a $C^{*}$-algebra such that $\operatorname{Prim}(A)$ is
Hausdorff. For every $x\in A$, the function
$\pi\mapsto\lVert\pi(x)\rVert$ is continuous on $\widehat A$.
\cite{ref77}

\ResultHeading{3.9.12}{3.9.12.}
Let $A$ be a $C^{*}$-algebra and $Z$ its center. Suppose that, whenever two
primitive ideals $I$ and $J$ satisfy $I\cap Z=J\cap Z$, one has $I=J$.
Under these conditions, $I\mapsto I\cap Z$ is a homeomorphism of
$\operatorname{Prim}(A)$ onto $\operatorname{Prim}(Z)$. \cite{ref74}

\ResultHeading{3.9.13}{3.9.13.}
Let $A$ be a separable $C^{*}$-algebra and $B$ a $C^{*}$-subalgebra of
$A$. Then
\[
 \aleph_0\operatorname{Card}\widehat A\geq\operatorname{Card}\widehat B.
\]
(For every $\pi\in\widehat A$, let $E_\pi$ be the set of
$\rho\in\widehat B$ which are contained in $\pi|_B$. Then $E_\pi$ is
countable, and $\bigcup_{\pi\in\widehat A}E_\pi=\widehat B$.)

% Source PDF physical page 091
\ResultHeading{3.9.14}{3.9.14.}
Let $A$ be an algebra over a commutative field, let
$\operatorname{Prim}(A)$ be the set of primitive ideals of $A$, equipped
with the Jacobson topology, and let $\widehat A$ be the set of classes of
irreducible representations of $A$, equipped with the inverse-image
topology of that of $\operatorname{Prim}(A)$ under the canonical map
$\widehat A\to\operatorname{Prim}(A)$. Let $I$ be a two-sided ideal of $A$.
Then, using the same notation as in \XRef{3.2.1}{3.2.1}, there are canonical
homeomorphisms
\[
 \operatorname{Prim}_I(A)\longrightarrow\operatorname{Prim}(A/I),
 \qquad
 \widehat A_I\longrightarrow\widehat{A/I},
\]
\[
 \operatorname{Prim}^{I}(A)\longrightarrow\operatorname{Prim}(I),
 \qquad
 \widehat A^{I}\longrightarrow\widehat I.
\]
\cite{ref77}

\ResultHeading{3.9.15}{3.9.15. Problem.}
Given two $C^{*}$-algebras $A,A'$, classify the $C^{*}$-algebras $B$ having
the following properties:
\begin{enumerate}[label=\arabic*\textsuperscript{o}]
\item $B$ has a closed two-sided ideal $I$ isomorphic to $A$;
\item $B/I$ is isomorphic to $A'$;
\item $\widehat I$ is dense in $\widehat B$.
\end{enumerate}
\cite{ref262,ref307}

\BookSection[Liminal C*-Algebras]{4}{Liminal $C^{*}$-Algebras}

The two problems mentioned at the beginning of Section~3 are far from being
solved in general.

For certain special classes of $C^{*}$-algebras, on the other hand, one
comes close to a satisfactory answer. These are the classes of
$C^{*}$-algebras that we shall now introduce.

\BookSubsection{4.1}{The Algebra of Compact Operators}

\ResultHeading{4.1.1}{4.1.1.}
Let $H$ be a Hilbert space. The compact operators on $H$ form a self-adjoint
two-sided ideal of $\mathcal L(H)$, and a norm limit of compact operators is
compact. The set of compact operators on $H$ is therefore a
$C^{*}$-algebra, which we shall denote by $\mathcal{LC}(H)$.

In the study of finite-dimensional algebras over the complex field, it is
known that the algebras $\mathcal L(H)$, with $H$ a finite-dimensional
vector space, play a fundamental role. In the study of $C^{*}$-algebras,
this role is played not by the algebras $\mathcal L(H)$, with $H$ a Hilbert
space, but by the algebras $\mathcal{LC}(H)$. [Recall that
$\mathcal L(H)=\mathcal{LC}(H)$ if $H$ is finite-dimensional.]

A $C^{*}$-algebra $A$ is called \emph{elementary} if there exists a Hilbert
space $H$ such that the $C^{*}$-algebras $A$ and $\mathcal{LC}(H)$ are
isomorphic.

\ResultHeading{4.1.2}{4.1.2. Proposition.}
\begin{itshape}
Let $H$ be a Hilbert space, $A=\mathcal{LC}(H)$, and $T$ the vector space
of trace-class operators on $H$. For every $t\in T$, let $f_t$ be the linear
form $x\mapsto\operatorname{Tr}(tx)$ on $A$. Then $t\mapsto f_t$ is a
linear bijection of $T$ onto the dual of the Banach space $A$. The form
$f_t$ is hermitian (respectively positive) if and only if $t$ is hermitian
(respectively positive).
\end{itshape}

We have
\[
 |\operatorname{Tr}(tx)|\leq(\operatorname{Tr}|t|)\lVert x\rVert
 \qquad(\XRef{A.32}{A 32}),
\]
so $f_t$ is a continuous linear form on $A$. The map $t\mapsto f_t$ is
clearly linear. If $f_t=0$, then
$0=f_t(t^{*})=\operatorname{Tr}(tt^{*})$, hence $t=0$, and the map
$t\mapsto f_t$ is injective. Let $f$ be a continuous linear form on $A$;
we show that there exists $t\in T$ such that $f=f_t$. It suffices to do so
under the assumption that $f$ is hermitian. Let $A^0\subset A$ be
% Source PDF physical page 092
the set of Hilbert--Schmidt operators; it carries a natural Hilbert-space
structure, and the canonical injection of $A^0$ into $A$ is continuous.
Thus the restriction $f'$ of $f$ to $A^0$ is a continuous linear form on
the Hilbert space $A^0$. Consequently, there exists $t_0\in A^0$ such that
\[
 f'(x)=(x\mid t_0)=\operatorname{Tr}(xt_0^{*})
 \qquad\text{for every }x\in A^0.
\]
For every $x\in A^0$, we have
\begin{align*}
 \operatorname{Tr}(xt_0)
 &=\overline{\operatorname{Tr}(t_0^{*}x^{*})}
  =\overline{\operatorname{Tr}(x^{*}t_0^{*})}
  =\overline{f'(x^{*})} \\
 &=f'(x)=\operatorname{Tr}(xt_0^{*}),
\end{align*}
hence $t_0=t_0^{*}$. Take the spectral decomposition of $t_0$: there are
pairwise orthogonal rank-one projections $e_i$ on $H$ and real numbers
$\lambda_i$ such that
\[
 \sum_i\lambda_i^2<+\infty,
 \qquad
 t_0=\lambda_1e_1+\lambda_2e_2+\cdots.
\]
Put $\varepsilon_i=1$ if $\lambda_i\geq0$, and
$\varepsilon_i=-1$ if $\lambda_i<0$. Then
\begin{align*}
 |\lambda_1|+\cdots+|\lambda_n|
 &=\operatorname{Tr}\bigl((\varepsilon_1e_1+\cdots+
    \varepsilon_ne_n)t_0\bigr) \\
 &=f(\varepsilon_1e_1+\cdots+\varepsilon_ne_n) \\
 &\leq\lVert f\rVert\,
   \lVert\varepsilon_1e_1+\cdots+\varepsilon_ne_n\rVert
  =\lVert f\rVert.
\end{align*}
Thus $\sum_i|\lambda_i|<+\infty$, and hence $t_0\in T$. The forms $f$ and
$f_{t_0}$ agree on $A^0$ by construction of $t_0$; therefore $f=f_{t_0}$,
since $A^0$ is dense in $A$. This proves that $t\mapsto f_t$ is a bijection
of $T$ onto the dual of the Banach space $A$.

For $t\in T$ and $x\in A$, we have
\[
 f_{t^{*}}(x)=\operatorname{Tr}(t^{*}x)
 =\overline{\operatorname{Tr}(tx^{*})}
 =\overline{f_t(x^{*})},
\]
and hence
\begin{equation}
 f_{t^{*}}=(f_t)^{*}.
 \tag{1}\label{eq:4.1.2-1}
\end{equation}
This proves that $f_t$ is hermitian if and only if $t$ is hermitian.
If $t\geq0$, then $t^{1/2}$ is a Hilbert--Schmidt operator, and for every
$x\in A$ we have
\[
 f_t(x^{*}x)
 =\operatorname{Tr}(t^{1/2}x^{*}xt^{1/2})
 =\operatorname{Tr}\bigl((xt^{1/2})^{*}(xt^{1/2})\bigr)\geq0;
\]
thus $f_t\geq0$. Finally, let $t\in T$ be such that $f_t\geq0$. If $\xi$
is a unit vector and $e$ is the orthogonal projection onto $\mathbb C\xi$,
then
\[
 0\leq f_t(e)=\operatorname{Tr}(te)=\operatorname{Tr}(ete)
 =(t\xi\mid\xi).
\]
This proves that $t\geq0$.

\ResultHeading{4.1.3}{4.1.3. Corollary.}
\begin{itshape}
The positive forms on $A$ are the forms
\[
 x\longmapsto\lambda_1(x\xi_1\mid\xi_1)
 +\lambda_2(x\xi_2\mid\xi_2)+\cdots,
\]
% Source PDF physical page 093
where $(\xi_1,\xi_2,\ldots)$ is an orthonormal system in $H$,
$\lambda_i\geq0$, and $\sum_i\lambda_i<+\infty$.
\end{itshape}

The positive trace-class operators on $H$ are the operators of the form
\[
 \lambda_1P_{\mathbb C\xi_1}+\lambda_2P_{\mathbb C\xi_2}+\cdots,
\]
where $(\xi_1,\xi_2,\ldots)$ is an orthonormal system in $H$,
$\lambda_i\geq0$, and $\sum_i\lambda_i<+\infty$. Since
\[
 \operatorname{Tr}(P_{\mathbb C\xi_i}x)
 =\operatorname{Tr}(P_{\mathbb C\xi_i}xP_{\mathbb C\xi_i})
 =(x\xi_i\mid\xi_i),
\]
the corollary follows at once from \XRef{4.1.2}{4.1.2}.

\ResultHeading{4.1.4}{4.1.4. Corollary.}
\begin{itshape}
The pure positive forms on $A$ are the forms
$x\mapsto(x\xi\mid\xi)$, where $\xi$ is a nonzero vector of $H$.
\end{itshape}

Since the identity representation of $A$ is irreducible when $H\ne0$, the
forms $x\mapsto(x\xi\mid\xi)$, with $\xi\ne0$, are pure
(\XRef{2.5.4}{2.5.4}). Conversely, if $f$ is a pure positive form on $A$,
all positive forms dominated by $f$ are proportional to $f$; hence, by
\XRef{4.1.3}{4.1.3}, $f$ is of the form $x\mapsto(x\xi\mid\xi)$.

\ResultHeading{4.1.5}{4.1.5. Corollary.}
\begin{itshape}
Every nonzero irreducible representation of $A$ is equivalent to the
identity representation.
\end{itshape}

This follows from \XRef{2.5.4}{2.5.4}, \XRef{4.1.4}{4.1.4}, and
\XRef{2.4.1}{2.4.1 (ii)}.

\ResultHeading{4.1.6}{4.1.6. Corollary.}
\begin{itshape}
Let $B\ne0$ be a $C^{*}$-subalgebra of $A$, and suppose that the identity
representation of $B$ is irreducible. Then $B=A$.
\end{itshape}

The commutant of $B$ in $\mathcal L(H)$ consists only of scalars. Since
$B\ne0$, it follows from \XRef{A.12}{A 12} that $B$ is ultra-strongly dense
in $\mathcal L(H)$, and therefore in $A$. With the notation of
\XRef{4.1.3}{4.1.3}, the form
\[
 x\longmapsto\sum_i\lambda_i(x\xi_i\mid\xi_i)
 =\sum_i\omega_{\lambda_i^{1/2}\xi_i}(x)
\]
is ultra-strongly continuous on $A$; thus every continuous linear form on
$A$ is ultra-strongly continuous on $A$. What precedes proves that every
continuous linear form on $A$ which vanishes on $B$ vanishes identically.
The Hahn--Banach theorem then shows that $B=A$.

\ResultHeading{4.1.7}{4.1.7. Corollary.}
\begin{itshape}
The only closed two-sided ideals of $A$ are $0$ and $A$.
\end{itshape}

By \XRef{4.1.5}{4.1.5}, the zero ideal is the only primitive ideal of $A$.
But every closed two-sided ideal of $A$ is an intersection of primitive
ideals [\XRef{2.9.7}{2.9.7 (ii)}].

\ResultHeading{4.1.8}{4.1.8. Corollary.}
\begin{itshape}
Let $H$ and $H'$ be two Hilbert spaces, and let $\varphi$ be an isomorphism
of the $C^{*}$-algebra $\mathcal{LC}(H)$ onto the $C^{*}$-algebra
$\mathcal{LC}(H')$.
\begin{enumerate}[label=(\roman*)]
\item There exists an isomorphism $U$ of $H$ onto $H'$ which implements
$\varphi$.
% Source PDF physical page 094
\item Let $U'$ be an isomorphism of $H$ onto $H'$. Then $U'$ implements
$\varphi$ if and only if $U'=\lambda U$, where $\lambda$ is a complex
number of modulus $1$.
\end{enumerate}
\end{itshape}

Since $\varphi$ is a nonzero irreducible representation of
$\mathcal{LC}(H)$ when $H,H'\ne0$, assertion (i) follows from
\XRef{4.1.5}{4.1.5}. For $U'$ to implement $\varphi$, it is necessary and
sufficient that $U'^{-1}U$ implement the identity automorphism of
$\mathcal{LC}(H)$; in other words, that it commute with
$\mathcal{LC}(H)$; in other words, that it be scalar. This proves (ii).

\ResultHeading{4.1.9}{4.1.9.}
Retain the notation of \XRef{4.1.8}{4.1.8}. We see that $\varphi$ maps the
set of Hilbert--Schmidt elements (respectively the set of elements of
finite rank, of rank one, and so on) of $\mathcal{LC}(H)$ onto the set of
Hilbert--Schmidt elements (respectively the set of elements of finite rank,
of rank one, and so on) of $\mathcal{LC}(H')$. Given an elementary
$C^{*}$-algebra $A$, we may therefore speak of the Hilbert--Schmidt
elements, and so on, of $A$, independently of the realization of $A$ as an
algebra of compact operators.

Similarly, the Hilbert dimension $a$ of $H$ is intrinsically attached to
the $C^{*}$-algebra $\mathcal{LC}(H)$; the latter is said to have rank
$a^2$. For finite $a$, this recovers the usual notion of the rank of an
algebra. Notice that, for infinite $a$, one has $a^2=a$.

\ResultHeading{4.1.10}{4.1.10. Corollary.}
\begin{itshape}
Let $A$ be a $C^{*}$-algebra and $\pi$ an irreducible representation of
$A$. If $\pi(A)\cap\mathcal{LC}(H_\pi)\ne0$, then
$\pi(A)\supset\mathcal{LC}(H_\pi)$, and every irreducible representation
of $A$ having the same kernel as $\pi$ is equivalent to $\pi$.
\end{itshape}

Let $A'=\pi^{-1}(\mathcal{LC}(H_\pi))$. Then $A'$ is a closed two-sided
ideal of $A$. We have $\pi(A')\ne0$, and hence $\pi(A')$ is irreducible on
$H_\pi$ [\XRef{2.11.2}{2.11.2 (ii)}]. Therefore
$\pi(A')=\mathcal{LC}(H_\pi)$ by \XRef{4.1.6}{4.1.6}. Every representation
of $A$ having the same kernel as $\pi$ is the composite of $\pi$ and an
injective representation of $\pi(A)$. To prove the second assertion of
\XRef{4.1.10}{4.1.10}, it therefore suffices to establish the following
point. Let $H$ be a Hilbert space and $B$ a $C^{*}$-subalgebra of
$\mathcal L(H)$ containing $\mathcal{LC}(H)$; then every injective
irreducible representation $\rho$ of $B$ is equivalent to the identity
representation. Since $\rho|_{\mathcal{LC}(H)}$ is equivalent to the
identity representation of $\mathcal{LC}(H)$ by \XRef{4.1.5}{4.1.5}, we
may suppose that $H_\rho=H$ and that $\rho|_{\mathcal{LC}(H)}$ is the
identity map of $\mathcal{LC}(H)$. Then \XRef{2.10.4}{2.10.4 (i)} proves
that $\rho$ is the identity map of $B$.

\ResultHeading{4.1.11}{4.1.11. Corollary.}
\begin{itshape}
Let $A$ be a $C^{*}$-algebra and $\pi$ a nonzero irreducible representation
of $A$. Suppose that $\pi(A)\subset\mathcal{LC}(H_\pi)$.
\begin{enumerate}[label=(\roman*)]
\item $\pi(A)=\mathcal{LC}(H_\pi)$.
\item The kernel of $\pi$ is a maximal closed two-sided ideal of $A$.
\end{enumerate}
\end{itshape}

Assertion (i) follows from \XRef{4.1.10}{4.1.10}; assertion (ii) follows
from (i) and \XRef{4.1.7}{4.1.7}.

Bibliography: \cite{ref74,ref77,ref105,ref106,ref108,ref126}.

% Source PDF physical page 095
\BookSubsection[Liminal C*-Algebras]{4.2}{Liminal $C^{*}$-Algebras}

\ResultHeading{4.2.1}{4.2.1. Definition.}
\begin{itshape}
A $C^{*}$-algebra $A$ is called liminal if, for every irreducible
representation $\pi$ of $A$ and every $x\in A$, the operator $\pi(x)$ is
compact.
\end{itshape}

\ResultHeading{4.2.2}{4.2.2.}
Let $H$ be a Hilbert space. The $C^{*}$-algebra $\mathcal{LC}(H)$ is liminal
(\XRef{4.1.5}{4.1.5}), and the zero ideal of this algebra is primitive if
$H\ne0$. Conversely, if $A$ is a liminal $C^{*}$-algebra in which the zero
ideal is primitive, then $A$ is elementary by
\XRef{4.1.11}{4.1.11 (i)}.

Every commutative $C^{*}$-algebra is liminal.

Later we shall see (\XRef{10.4.5}{10.4.5},
\XRef{13.11.12}{13.11.12}, \XRef{15.5.6}{15.5.6}) other examples of
liminal $C^{*}$-algebras.

\ResultHeading{4.2.3}{4.2.3.}
Let $A$ be a liminal $C^{*}$-algebra. Every primitive ideal of $A$ is a
maximal closed two-sided ideal [\XRef{4.1.11}{4.1.11 (ii)}]. If $\pi$ is a
nonzero irreducible representation of $A$, then
$\pi(A)=\mathcal{LC}(H_\pi)$ [\XRef{4.1.11}{4.1.11 (i)}].

\ResultHeading{4.2.4}{4.2.4. Proposition.}
\begin{itshape}
Let $A$ be a liminal $C^{*}$-algebra. Every $C^{*}$-subalgebra of $A$ and
every quotient $C^{*}$-algebra of $A$ are liminal.
\end{itshape}

This is clear for a quotient $C^{*}$-algebra $A/I$ of $A$, since every
irreducible representation of $A/I$ is identified with an irreducible
representation of $A$. Now let $B$ be a $C^{*}$-subalgebra of $A$, and let
$\pi$ be an irreducible representation of $B$. There exist a Hilbert space
$K\supset H_\pi$ and an irreducible representation $\rho$ of $A$ on $K$
such that
\[
 \pi(x)=\rho(x)|_{H_\pi}\qquad\text{for every }x\in B
\]
(\XRef{2.10.2}{2.10.2}). Since $A$ is liminal, the operators $\rho(x)$ are
compact, and therefore the operators $\pi(x)$ are compact. Thus $B$ is
liminal.

\ResultHeading{4.2.5}{4.2.5. Proposition.}
\begin{itshape}
Let $A$ be a liminal $C^{*}$-algebra, let
$\pi_1,\ldots,\pi_n$ be pairwise inequivalent elements of $\widehat A$, and
let $T_1,\ldots,T_n$ be compact operators on
$H_{\pi_1},\ldots,H_{\pi_n}$, respectively. There exists $z\in A$ such that
\[
 \pi_1(z)=T_1,\ldots,\pi_n(z)=T_n.
\]
\end{itshape}

This is clear for $n=1$; suppose it has been proved for $n-1$. There exist
$x,y\in A$ such that
\[
 \pi_1(x)=T_1,\ldots,\pi_{n-1}(x)=T_{n-1},
 \qquad \pi_n(y)=T_n.
\]
Put $I_k=\ker\pi_k$. By \XRef{4.1.10}{4.1.10} and
\XRef{4.2.3}{4.2.3}, $I_n\not\supset I_k$ for $k=1,\ldots,n-1$; hence
\[
 I_n\not\supset I=I_1\cap\cdots\cap I_{n-1}
 \qquad(\XRef{2.11.4}{2.11.4}).
\]
There exist $x'\in I$ and $y'\in I_n$ such that
$x+x'=y+y'=z$. Then $\pi_k(z)=T_k$ for $k=1,\ldots,n$.

\ResultHeading{4.2.6}{4.2.6. Proposition.}
\begin{itshape}
Let $A$ be a $C^{*}$-algebra. Let $I$ be the set of $x\in A$ such that
$\pi(x)$ is compact for every irreducible representation $\pi$ of $A$.
Then $I$ is the largest closed two-sided liminal ideal of $A$.
\end{itshape}

It is clear that $I$ is a closed two-sided ideal of $A$. Let $\rho$ be an
irreducible representation of $I$. There exists an irreducible
representation $\pi$ of $A$ on $H_\rho$ which extends $\rho$
(\XRef{2.10.4}{2.10.4}). Then $\rho(x)=\pi(x)$ is compact for every
$x\in I$, and hence $I$ is liminal. Finally, let $J$ be a closed two-sided
liminal ideal
% Source PDF physical page 096
of $A$. We show that $J\subset I$, that is, that $\sigma(x)$ is compact for
every irreducible representation $\sigma$ of $A$ and every $x\in J$. This
is clear if $\sigma(J)=0$. Otherwise, $\sigma|_J$ is irreducible
[\XRef{2.11.2}{2.11.2 (ii)}], and $\sigma(x)$ is compact because $J$ is
liminal.

Kaplansky introduced liminal $C^{*}$-algebras in 1951 under the name
CCR-algebras (here CCR stands for ``completely continuous representations'';
``completely continuous'' is synonymous with ``compact''). Kaplansky and
Glimm also introduced the notations GCR and NGCR for the $C^{*}$-algebras
called below postliminal and antiliminal. These notions are fundamental,
but the accepted terminology seemed to me somewhat unaesthetic.

Bibliography: \cite{ref77}. The proof of \XRef{4.2.5}{4.2.5} was pointed
out to me by J. Tomiyama.

\BookSubsection[Postliminal C*-Algebras]{4.3}{Postliminal $C^{*}$-Algebras}

\ResultHeading{4.3.1}{4.3.1. Definition.}
\begin{itshape}
A $C^{*}$-algebra $A$ is called postliminal if every nonzero quotient
$C^{*}$-algebra of $A$ has a nonzero closed two-sided liminal ideal.
\end{itshape}

These are the most important $C^{*}$-algebras for what follows. Every
liminal $C^{*}$-algebra is postliminal (\XRef{4.2.4}{4.2.4}), but the
converse is not true.

\ResultLabel{Definition.}
\begin{itshape}
A $C^{*}$-algebra is called antiliminal if it has no nonzero closed
two-sided liminal ideal.
\end{itshape}

By \XRef{1.8.5}{1.8.5}, such a $C^{*}$-algebra has no nonzero closed
two-sided postliminal ideal.

For examples of postliminal or antiliminal $C^{*}$-algebras, see the
Complements (Sections~4, 5, 9, 10, 13, and 17).

\ResultHeading{4.3.2}{4.3.2. Definition.}
\begin{itshape}
Let $A$ be a $C^{*}$-algebra. A \emph{composition series} of $A$ is an
increasing family $(I_\rho)_{0\leq\rho\leq\alpha}$ of closed two-sided
ideals of $A$, indexed by the ordinals between $0$ and some ordinal
$\alpha$, and having the following properties:
\begin{enumerate}[label=(\roman*)]
\item $I_0=0$ and $I_\alpha=A$;
\item if $\rho\leq\alpha$ is a limit ordinal, then $I_\rho$ is the closure
of $\bigcup_{\rho'<\rho}I_{\rho'}$.
\end{enumerate}
\end{itshape}

\ResultHeading{4.3.3}{4.3.3. Proposition.}
\begin{itshape}
Let $A$ be a $C^{*}$-algebra. There exist an ordinal $\alpha$ and an
increasing family $(I_\rho)_{0\leq\rho\leq\alpha}$ of closed two-sided
ideals of $A$ having the following properties:
\begin{enumerate}[label=(\roman*)]
\item $I_0=0$, and $A/I_\alpha$ is antiliminal;
\item if $\rho\leq\alpha$ is a limit ordinal, then $I_\rho$ is the closure
of $\bigcup_{\rho'<\rho}I_{\rho'}$;
\item if $\rho<\alpha$, then $I_{\rho+1}/I_\rho$ is the largest liminal
two-sided ideal of $A/I_\rho$, and is nonzero.
\end{enumerate}
Moreover, $\alpha$ and the family $(I_\rho)$ are uniquely determined by
the preceding conditions.
\end{itshape}

% Source PDF physical page 097
The $I_\rho$ are constructed by transfinite induction. If $\rho$ is a
limit ordinal, condition (ii) prescribes the construction of $I_\rho$.
Otherwise $\rho=\rho'+1$, and condition (iii) prescribes the construction
of $I_\rho$, unless $A/I_{\rho'}$ is antiliminal, in which case the
construction stops with $\alpha=\rho'$.

\ResultHeading{4.3.4}{4.3.4. Proposition.}
\begin{itshape}
Let $A$ be a $C^{*}$-algebra. The following conditions are equivalent:
\begin{enumerate}[label=(\roman*)]
\item $A$ is postliminal;
\item $A$ has a composition series
$(J_\rho)_{0\leq\rho\leq\beta}$ such that all the
$J_{\rho+1}/J_\rho$ are postliminal;
\item $A$ has a composition series
$(I_\rho)_{0\leq\rho\leq\alpha}$ such that all the
$I_{\rho+1}/I_\rho$ are liminal.
\end{enumerate}
\end{itshape}

(i) $\Rightarrow$ (iii): suppose that $A$ is postliminal. With the notation
of \XRef{4.3.3}{4.3.3}, we cannot have $A/I_\alpha\ne0$; hence
$I_\alpha=A$, and $(I_\rho)_{0\leq\rho\leq\alpha}$ is a composition series
of $A$ with the required properties.

(iii) $\Rightarrow$ (ii) is clear.

(ii) $\Rightarrow$ (i): suppose there exists a composition series
$(J_\rho)_{0\leq\rho\leq\beta}$ such that the
$J_{\rho+1}/J_\rho$ are postliminal. Let $J\ne A$ be a closed two-sided
ideal of $A$. We must show that $A/J$ has a nonzero liminal two-sided ideal.
There is a least ordinal $\rho\leq\beta$ such that $J_\rho\not\subset J$.
If $\rho'<\rho$, then $J_{\rho'}\subset J$; hence the closure of
$\bigcup_{\rho'<\rho}J_{\rho'}$ is contained in $J$, so $\rho$ is not a
limit ordinal. Write $\rho=\rho'+1$ and put $I=J_\rho\cap J$. Then
$J_\rho/I$ is identified with a closed two-sided ideal of $A/J$
(\XRef{1.8.4}{1.8.4}). Moreover, $J_\rho/I$ is a nonzero quotient of the
postliminal $C^{*}$-algebra $J_{\rho'+1}/J_{\rho'}$, and hence has a
nonzero liminal two-sided ideal. By \XRef{1.8.5}{1.8.5}, this ideal is also
a two-sided ideal of $A/J$.

\ResultHeading{4.3.5}{4.3.5. Proposition.}
\begin{itshape}
Let $A$ be a postliminal $C^{*}$-algebra. Every $C^{*}$-subalgebra of $A$
and every quotient $C^{*}$-algebra of $A$ are postliminal.
\end{itshape}

This is clear for quotient $C^{*}$-algebras. Now let $B$ be a
$C^{*}$-subalgebra of $A$. Let $(I_\rho)_{0\leq\rho\leq\alpha}$ be a
composition series of $A$ such that the $I_{\rho+1}/I_\rho$ are liminal
(\XRef{4.3.4}{4.3.4}). Then
\[
 (I_{\rho+1}\cap B)/(I_\rho\cap B)
\]
is identified with a $C^{*}$-subalgebra of $I_{\rho+1}/I_\rho$
(\XRef{1.8.4}{1.8.4}), and hence is liminal (\XRef{4.2.4}{4.2.4}). To
complete the proof by \XRef{4.3.4}{4.3.4}, we show that
$(I_\rho\cap B)_{0\leq\rho\leq\alpha}$ is a composition series of $B$.
It suffices to show that, if $\rho$ is a limit ordinal and
$x\in I_\rho\cap B$, then $x$ is a limit of points of
$\bigcup_{\rho'<\rho}(I_{\rho'}\cap B)$. Let $\varepsilon>0$. There is a
$\rho'<\rho$ such that the distance from $x$ to $I_{\rho'}$ is less than
$\varepsilon$. Since $B/(I_{\rho'}\cap B)$ is identified with
$(B+I_{\rho'})/I_{\rho'}$, there is a point of $I_{\rho'}\cap B$ whose
distance from $x$ is less than $\varepsilon$, proving our assertion.

% Source PDF physical page 098
\ResultHeading{4.3.6}{4.3.6. Proposition.}
\begin{itshape}
Let $A$ be a $C^{*}$-algebra. The ideal $I_\alpha$ of
\XRef{4.3.3}{4.3.3} is the largest postliminal two-sided ideal of $A$, and
the smallest closed two-sided ideal for which the corresponding quotient
of $A$ is antiliminal.
\end{itshape}

The ideal $I_\alpha$ is postliminal by \XRef{4.3.4}{4.3.4}. Let $I$ be a
postliminal two-sided ideal of $A$. Then $I/(I\cap I_\alpha)$ is identified
with a postliminal two-sided ideal of $A/I_\alpha$. Since $A/I_\alpha$ is
antiliminal, $I/(I\cap I_\alpha)=0$, and hence $I\subset I_\alpha$.
Finally, let $J$ be a closed two-sided ideal of $A$ such that $A/J$ is
antiliminal. Then $I_\alpha/(J\cap I_\alpha)$ is identified with a
postliminal two-sided ideal of $A/J$, so
$I_\alpha/(J\cap I_\alpha)=0$ and $J\supset I_\alpha$.

Proposition~\XRef{4.3.6}{4.3.6} reduces, to some extent, the study of
arbitrary $C^{*}$-algebras to the study of postliminal $C^{*}$-algebras on
the one hand and antiliminal $C^{*}$-algebras on the other. Very little is
known about antiliminal $C^{*}$-algebras. By contrast, fairly detailed
information is available about postliminal $C^{*}$-algebras.

\ResultHeading{4.3.7}{4.3.7. Theorem.}
\begin{itshape}
Let $A$ be a postliminal $C^{*}$-algebra.
\begin{enumerate}[label=(\roman*)]
\item For every nonzero irreducible representation $\pi$ of $A$,
$\pi(A)$ contains all the compact operators on $H_\pi$.
\item If $\pi_1$ and $\pi_2$ are two irreducible representations of $A$
with the same kernel, then $\pi_1$ and $\pi_2$ are equivalent.
\end{enumerate}
\end{itshape}

Let $\pi$ be a nonzero irreducible representation of $A$. Then
$A/\ker\pi$ contains a nonzero liminal two-sided ideal $I/\ker\pi$. The
restriction of $\pi$ to $I$ is nonzero and irreducible
[\XRef{2.11.2}{2.11.2 (ii)}], and therefore
$\pi(I)=\mathcal{LC}(H_\pi)$ because $I/\ker\pi$ is liminal
(\XRef{4.2.3}{4.2.3}). This proves (i); assertion (ii) follows from (i) and
\XRef{4.1.10}{4.1.10}.

Theorem~\XRef{4.3.7}{4.3.7} has a converse (\XRef{9.1}{9.1}).

\ResultHeading{4.3.8}{4.3.8.}
Let $A$ be a separable $C^{*}$-algebra and
$(I_\rho)_{0\leq\rho\leq\alpha}$ a composition series of $A$. If the
$I_{\rho+1}/I_\rho$ are nonzero, the family $(I_\rho)$ is countable. Indeed,
for every ordinal $\rho<\alpha$, choose $x_\rho\in I_{\rho+1}$ whose
distance from $I_\rho$ is at least $1$. Then the mutual distances of the
$x_\rho$ are at least $1$, and hence the family of the $x_\rho$ is
countable.

\ResultHeading{4.3.9}{4.3.9.}
Let $A$ be a $C^{*}$-algebra and $\widetilde A$ the $C^{*}$-algebra
obtained from $A$ by adjoining an identity element. If $A$ is postliminal,
then $\widetilde A$ is postliminal (\XRef{4.3.4}{4.3.4}). If
$\widetilde A$ is postliminal, then $A$ is postliminal
(\XRef{4.3.5}{4.3.5}). If $\widetilde A$ is antiliminal, then $A$ is
obviously antiliminal. Suppose that $A$ is antiliminal. If $A$ has an
identity element, then $\widetilde A$ has a one-dimensional ideal, and so
is not antiliminal. But if $A$ has no identity element, then
$\widetilde A$ is antiliminal; for if $\widetilde A$ had a nonzero liminal
two-sided ideal $I$, then $I\cap A=0$, so $I$ would be a complement of $A$
in $\widetilde A$; consequently $A$ would be a quotient of $\widetilde A$
and would have an identity element.

Bibliography: \cite{ref11,ref28,ref46,ref77}.

% Source PDF physical page 099
\BookSubsection{4.4}{Spectrum of a Postliminal Algebra}

\ResultHeading{4.4.1}{4.4.1.}
If $A$ is a postliminal $C^{*}$-algebra, Theorem~\XRef{4.3.7}{4.3.7}
shows that the canonical map $\widehat A\to\operatorname{Prim}(A)$ is a
homeomorphism. Thus $\widehat A$ is a $T_0$-space
(\XRef{3.1.3}{3.1.3}). If, moreover, $A$ is liminal, all the primitive
ideals of $A$ are maximal (\XRef{4.2.3}{4.2.3}), and hence the points of
$\widehat A$ are closed (\XRef{3.1.4}{3.1.4}).

\ResultHeading{4.4.2}{4.4.2. Lemma.}
\begin{itshape}
Let $A$ be a $C^{*}$-algebra, let $\pi_0\in\widehat A$, and let $x\in A^+$
be such that the function $\pi\mapsto\operatorname{Tr}\pi(x)$ on
$\widehat A$ is finite and continuous at $\pi_0$.
\begin{enumerate}[label=(\roman*)]
\item If $y\in A^+$ is dominated by $x$, the function
$\pi\mapsto\operatorname{Tr}\pi(y)$ on $\widehat A$ is continuous at
$\pi_0$.
\item If $\pi_0(x)\ne0$, there exists $z\in A^+$ such that $\pi(z)$ is a
rank-one projection for every $\pi$ in a neighborhood of $\pi_0$.
\end{enumerate}
\end{itshape}

Put $y'=x-y\in A^+$. The functions
$\pi\mapsto\operatorname{Tr}\pi(y)$ and
$\pi\mapsto\operatorname{Tr}\pi(y')$ are lower semicontinuous
(\XRef{3.5.9}{3.5.9}); their finite sum is continuous at $\pi_0$, and hence
both functions are continuous at $\pi_0$.

To prove (ii), we may suppose that $\lVert\pi_0(x)\rVert=1$. Let $K$ be a
one-dimensional eigenspace of $\pi_0(x)$ corresponding to the eigenvalue
$1$. The irreducible algebra $\pi_0(A)$ contains the nonzero compact
operator $\pi_0(x)$, and hence contains $\mathcal{LC}(H_{\pi_0})$
(\XRef{4.1.10}{4.1.10}); consequently, there exists $z_1\in A^+$ such that
$\pi_0(z_1)=P_K$. Put $z_2=f(z_1)$, where $f(t)=t$ for $0\leq t\leq1$ and
$f(t)=1$ for $t\geq1$. Then $\lVert z_2\rVert\leq1$ and
$\pi_0(z_2)=f(P_K)=P_K$. Put
\[
 z_3=x^{1/2}z_2x^{1/2}\leq x;
\]
then
\[
 \pi_0(z_3)=\pi_0(x)^{1/2}P_K\pi_0(x)^{1/2}=P_K.
\]
By (i), $\operatorname{Tr}\pi(z_3)\leq5/4$ in a neighborhood of $\pi_0$;
on the other hand, $\lVert\pi(z_3)\rVert\geq3/4$ in a neighborhood of
$\pi_0$ (\XRef{3.3.2}{3.3.2}). Thus, in a neighborhood $V$ of $\pi_0$,
the largest eigenvalue of $\pi(z_3)$ is at least $3/4$ and has multiplicity
one, while all the following eigenvalues are at most
$5/4-3/4=1/2$. Let $g\colon\mathbb R\to\mathbb R$ be a continuous
nonnegative function equal to $0$ for $t\leq1/2$ and to $1$ for
$t\geq3/4$, and put $z=g(z_3)\in A^+$. Then
$\pi(z)=g(\pi(z_3))$ is a rank-one projection for every $\pi\in V$.

\ResultHeading{4.4.3}{4.4.3. Lemma.}
\begin{itshape}
Let $A$ be a liminal $C^{*}$-algebra. Let $I$ be the set of $x\in A$ such
that $\pi(x)$ has finite rank for every $\pi\in\widehat A$. Then $I$ is a
dense self-adjoint two-sided ideal of $A$.
\end{itshape}

It is clear that $I$ is a self-adjoint two-sided ideal of $A$. For
$n=1,2,\ldots$, let $f_n\colon\mathbb R\to\mathbb R$ be a continuous
function which vanishes in a neighborhood of $0$ (depending on $n$), and
such that $f_n(t)$ tends uniformly
% Source PDF physical page 100
to $t$ as $n\to+\infty$. For every hermitian element $x$ of $A$, one has
$f_n(x)\to x$ as $n\to+\infty$. On the other hand, for every
$\pi\in\widehat A$, the operator $\pi(x)$ is compact, and therefore its
eigenvalues, each counted with its multiplicity, tend to zero. Thus
\[
 \pi(f_n(x))=f_n(\pi(x))
\]
has only finitely many nonzero eigenvalues, each of finite multiplicity;
hence $\pi(f_n(x))$ has finite rank. Therefore $I$ is dense in $A$.

\ResultHeading{4.4.4}{4.4.4. Lemma.}
\begin{itshape}
Let $A$ be a nonzero postliminal $C^{*}$-algebra. Then $A$ has a nonzero
closed two-sided ideal $I$ with the following property: for any two
distinct points $\pi_1,\pi_2$ of the open subset $\widehat I$ of
$\widehat A$, there exists $x\in I^+$ such that
\begin{enumerate}[label=\alph*.]
\item $\pi(x)$ has rank at most one for every $\pi\in\widehat I$;
\item the function $\pi\mapsto\operatorname{Tr}\pi(x)$ is continuous on
$\widehat I$;
\item $\operatorname{Tr}\pi_1(x)=1$ and
$\operatorname{Tr}\pi_2(x)=0$.
\end{enumerate}
\end{itshape}

We may suppose that $A$ is liminal. By \XRef{4.4.3}{4.4.3}, there exists a
nonzero $x_1\in A^+$ such that $\pi(x_1)$ has finite rank for every
$\pi\in\widehat A$; thus the function
\[
 \pi\longmapsto f(\pi)=\operatorname{Tr}\pi(x_1)
\]
on $\widehat A$ is finite everywhere. This function is lower
semicontinuous and not identically zero, and hence remains at least some
$\varepsilon>0$ on a nonempty open subset $U$ of $\widehat A$. Since $U$
is a Baire space (\XRef{3.2.2}{3.2.2} and
\XRef{3.4.13}{3.4.13}), $f$ has a point of continuity $\pi_0$ in $U$
(\XRef{B.18}{B 18}). By \XRef{4.4.2}{4.4.2 (ii)}, there exist
$x_2\in A^+$ and an open neighborhood $V$ of $\pi_0$ such that
$\pi(x_2)$ is a rank-one projection for every $\pi\in V$. Let $I$ be the
closed two-sided ideal of $A$ such that $\widehat I=V$. Let
$\pi_1,\pi_2$ be two distinct points of $V$. There exists $x_3\in I^+$ such
that
\[
 \pi_1(x_3)=\pi_1(x_2),\qquad \pi_2(x_3)=0
\]
(\XRef{4.2.4}{4.2.4} and \XRef{4.2.5}{4.2.5}). Put
\[
 x=x_2^{1/2}x_3x_2^{1/2}\in I^+.
\]
Since $x\leq\lVert x_3\rVert x_2$, the function
$\pi\mapsto\operatorname{Tr}\pi(x)$ is continuous on $V$ by
\XRef{4.4.2}{4.4.2 (i)}. Moreover,
\[
 \operatorname{Tr}\pi_1(x)=\operatorname{Tr}\pi_1(x_2)^2=1,
 \qquad
 \operatorname{Tr}\pi_2(x)=0.
\]
Finally, the rank of $\pi(x)$ is at most the rank of $\pi(x_2)$, and hence
at most one, for every $\pi\in V$.

\ResultHeading{4.4.5}{4.4.5. Theorem.}
\begin{itshape}
Let $A$ be a postliminal $C^{*}$-algebra. There exists in $\widehat A$ a
dense open subset which is locally compact.
\end{itshape}

With the notation of \XRef{4.4.4}{4.4.4}, the points of $\widehat I$ are
separated by continuous real-valued functions, and hence $\widehat I$ is
Hausdorff. The set of Hausdorff subsets of $\widehat A$ is clearly
inductive under inclusion. Let $S\subset\widehat A$ be a maximal Hausdorff
open subset. Suppose that there is a nonempty open subset $U$ of
$\widehat A$ disjoint from $S$. Then $U$ is identified with the spectrum
of a nonzero closed two-sided ideal of $A$, which is a postliminal
$C^{*}$-algebra. By the beginning of the proof, $U$ contains a nonempty
Hausdorff open subset $S'$. It is immediate that $S\cup S'$ is open and
Hausdorff in $\widehat A$, which
% Source PDF physical page 101
contradicts the maximality of $S$. Thus $\overline S=\widehat A$. By
\XRef{3.3.8}{3.3.8}, $S$ is locally compact.

Bibliography: \cite{ref16,ref77}.

\BookSubsection[Continuous-Trace C*-Algebras]{4.5}{Continuous-Trace
$C^{*}$-Algebras}

\ResultHeading{4.5.1}{4.5.1. Lemma.}
\begin{itshape}
Let $A$ be a $C^{*}$-algebra and let $\mathfrak p$ be a subset of $A^+$
such that:
\begin{enumerate}[label=\alph*.]
\item $\mathfrak p+\mathfrak p\subset\mathfrak p$;
\item if $x\in A$ and $xx^{*}\in\mathfrak p$, then
$x^{*}x\in\mathfrak p$;
\item every element of $A^+$ dominated by an element of $\mathfrak p$
belongs to $\mathfrak p$.
\end{enumerate}
Then:
\begin{enumerate}[label=(\roman*)]
\item the set
$\mathfrak n=\{x\in A:xx^{*}\in\mathfrak p\}$ is a self-adjoint two-sided
ideal of $A$;
\item the two-sided ideal $\mathfrak m=\mathfrak n^2$ is the set of linear
combinations of elements of
$\mathfrak m^+=\mathfrak m\cap A^+$, and it has the same closure as
$\mathfrak n$ in $A$;
\item $\mathfrak m^+=\mathfrak p$.
\end{enumerate}
\end{itshape}

It is clear that $x\in\mathfrak n$ implies $x^{*}\in\mathfrak n$. Suppose
$x,y\in\mathfrak n$. We have
\begin{align*}
 2xx^{*}+2yy^{*}-(x+y)(x+y)^{*}
 &=xx^{*}+yy^{*}-xy^{*}-yx^{*} \\
 &=(x-y)(x-y)^{*}\geq0,
\end{align*}
and hence
\[
 (x+y)(x+y)^{*}\leq2xx^{*}+2yy^{*}\in\mathfrak p.
\]
Thus $(x+y)(x+y)^{*}\in\mathfrak p$, so $x+y\in\mathfrak n$. If
$x\in\mathfrak n$ and $z\in A$, then
\[
 xzz^{*}x^{*}\leq\lVert zz^{*}\rVert xx^{*}\in\mathfrak p;
\]
therefore $(xz)(xz)^{*}\in\mathfrak p$, so $xz\in\mathfrak n$. Thus
$\mathfrak n$ is a right ideal of $A$, and consequently a two-sided ideal
because $\mathfrak n=\mathfrak n^{*}$. This proves (i).

Every element of $\overline{\mathfrak n}$ is a linear combination of
products of two elements of $\overline{\mathfrak n}$
(\XRef{1.5.8}{1.5.8}), and hence
is a limit of elements of $\mathfrak m$. Thus
$\overline{\mathfrak n}=\overline{\mathfrak m}$. Let
$x\in\mathfrak m$. Write $x=\sum_j a_jb_j^{*}$, with
$a_j,b_j\in\mathfrak n$. Then
\begin{align*}
 4x={}&\sum_j(a_j+b_j)(a_j+b_j)^{*}
       -\sum_j(a_j-b_j)(a_j-b_j)^{*} \\
 &+i\sum_j(a_j+ib_j)(a_j+ib_j)^{*}
       -i\sum_j(a_j-ib_j)(a_j-ib_j)^{*},
\end{align*}
and
\[
 (a_j+b_j)(a_j+b_j)^{*},\ldots,
 (a_j-ib_j)(a_j-ib_j)^{*}\in\mathfrak m\cap A^+.
\]
This proves (ii).

If the preceding element $x$ is hermitian, then
\[
 4x=\sum_j(a_j+b_j)(a_j+b_j)^{*}
     -\sum_j(a_j-b_j)(a_j-b_j)^{*}
 \leq\sum_j(a_j+b_j)(a_j+b_j)^{*}\in\mathfrak p.
\]
% Source PDF physical page 102
If $x\geq0$, it follows that $x\in\mathfrak p$. Hence
$\mathfrak m^+\subset\mathfrak p$. Conversely, let $x\in\mathfrak p$.
Then $x^{1/2}\in\mathfrak n$, and hence
$x=x^{1/2}x^{1/2}\in\mathfrak m^+$. This proves (iii).

\ResultHeading{4.5.2}{4.5.2.}
Let $A$ be a $C^{*}$-algebra. Let $\mathfrak p$ be the set of $x\in A^+$
such that the function $\pi\mapsto\operatorname{Tr}\pi(x)$ is finite and
continuous on $\widehat A$. It is clear that
$\mathfrak p+\mathfrak p\subset\mathfrak p$. If $x\in A$ is such that
$xx^{*}\in\mathfrak p$, then $x^{*}x\in\mathfrak p$, since
\[
 \operatorname{Tr}\pi(x^{*}x)=\operatorname{Tr}\pi(xx^{*})
 \qquad\text{for every }\pi\in\widehat A.
\]
If $x\in\mathfrak p$ and $y\in A^+$ is dominated by $x$, then
$y\in\mathfrak p$ by \XRef{4.4.2}{4.4.2 (i)}. Lemma~\XRef{4.5.1}{4.5.1}
therefore proves that $\mathfrak p$ is the positive part of a self-adjoint
two-sided ideal $\mathfrak m$ of $A$, namely the set of linear combinations
of elements of $\mathfrak p$. If $x\in\mathfrak m$, then $\pi(x)$ is a
trace-class operator for every $\pi\in\widehat A$, and the function
$\pi\mapsto\operatorname{Tr}\pi(x)$ is continuous on $\widehat A$. We
write $\mathfrak m=\mathfrak m(A)$.

\ResultLabel{Definition.}
\begin{itshape}
A $C^{*}$-algebra $A$ is called a continuous-trace algebra if the ideal
$\mathfrak m(A)$ is dense in $A$.
\end{itshape}

Since $\overline{\mathfrak m(A)}$ is the intersection of the primitive
ideals of $A$ which contain it, this is equivalent to saying that, for
every $\pi\in\widehat A$, there exists $x\in\mathfrak m(A)$ such that
$\pi(x)\ne0$.

By \XRef{4.5.1}{4.5.1 (ii)}, it is also equivalent to say that the set of
$x\in A$ such that the function
\[
 \pi\longmapsto\operatorname{Tr}\bigl(\pi(x)\pi(x)^{*}\bigr)
\]
is finite and continuous on $\widehat A$ is dense in $A$. This set is a
self-adjoint two-sided ideal $\mathfrak n$ such that
$\mathfrak n^2=\mathfrak m(A)$.

Examples of continuous-trace $C^{*}$-algebras will be given later
(\XRef{10.5.8}{10.5.8} and \XRef{10.9.4}{10.9.4}).

\ResultHeading{4.5.3}{4.5.3. Proposition.}
\begin{itshape}
Let $A$ be a continuous-trace $C^{*}$-algebra.
\begin{enumerate}[label=(\roman*)]
\item $A$ is liminal.
\item $\widehat A$ is Hausdorff.
\item For every $\pi_0\in\widehat A$, there exists $y\in A^+$ such that
$\pi(y)$ is a rank-one projection for every $\pi$ in a neighborhood of
$\pi_0$.
\end{enumerate}
\end{itshape}

Let $\pi\in\widehat A$. For every $x\in\mathfrak m(A)$, $\pi(x)$ is a
trace-class operator. Hence, for every
$x\in\overline{\mathfrak m(A)}=A$, $\pi(x)$ is a norm limit of trace-class
operators and is therefore compact. Thus $A$ is liminal. Assertion (iii)
follows from \XRef{4.4.2}{4.4.2 (ii)}. Let $\pi_1$ and $\pi_2$ be two
distinct points of $\widehat A$, let $I=\ker\pi_1$, and choose
$x\in\mathfrak m(A)$ such that $\pi_2(x)\ne0$. By
\XRef{4.2.5}{4.2.5}, there exists $x_1\in I$ such that
$\pi_2(x_1)=\pi_2(x)^{*}$. Then $x_1x\in\mathfrak m(A)$, so the function
$\pi\mapsto\operatorname{Tr}\pi(x_1x)$ is finite and continuous on
$\widehat A$; moreover,
\[
 \operatorname{Tr}\pi_1(x_1x)=\operatorname{Tr}(0)=0,
 \qquad
 \operatorname{Tr}\pi_2(x_1x)
 =\operatorname{Tr}\bigl(\pi_2(x)^{*}\pi_2(x)\bigr)\ne0.
\]
Thus $\widehat A$ is Hausdorff.

% Source PDF physical page 103
\ResultHeading{4.5.4}{4.5.4. Proposition.}
\begin{itshape}
Let $A$ be a $C^{*}$-algebra satisfying conditions (ii) and (iii) of
\XRef{4.5.3}{4.5.3}. Then $A$ is a continuous-trace $C^{*}$-algebra.
\end{itshape}

Let $\pi_0\in\widehat A$. There exist an open neighborhood $V$ of $\pi_0$
and $y\in A^+$ such that $\pi(y)$ is a rank-one projection for every
$\pi\in V$. Since $\widehat A$ is Hausdorff, it is locally compact
(\XRef{3.3.8}{3.3.8}), and there is an open neighborhood $W$ of $\pi_0$
such that $\overline W\subset V$. Let $I$ be the closed two-sided ideal of
$A$ such that $\widehat I=W$. Since $\pi_0\in W$, the restriction
$\pi_0|_I$ is irreducible. Hence there is a hermitian element $x$ of $I$
such that $\pi_0(x)$ acts identically on the one-dimensional subspace
$\pi_0(y)(H_{\pi_0})$ (\XRef{2.8.3}{2.8.3}); replacing $x$ by $x^+$, we
may suppose that $x\in I^+$. Put
\[
 z=y^{1/2}xy^{1/2}\in I^+.
\]
The function
\[
 \pi\longmapsto f(\pi)=\operatorname{Tr}\pi(z)
\]
on $\widehat A$ vanishes on $\widehat A-W$, and hence is continuous at
every point of $\widehat A-V$. On the other hand,
$z\leq\lVert x\rVert y$, so $f$ is continuous at every point of $V$ by
\XRef{4.4.2}{4.4.2 (i)}. Thus $z\in\mathfrak m(A)^+$. Since
$\pi_0(z)=\pi_0(y)\ne0$, this proves that $A$ is a continuous-trace
$C^{*}$-algebra.

\ResultHeading{4.5.5}{4.5.5. Theorem.}
\begin{itshape}
Let $A$ be a postliminal $C^{*}$-algebra. Then $A$ has a composition series
$(I_\rho)_{0\leq\rho\leq\alpha}$ such that the quotients
$I_{\rho+1}/I_\rho$ are continuous-trace $C^{*}$-algebras.
\end{itshape}

If $A\ne0$, then $A$ has a nonzero closed two-sided ideal $I_1$ which is a
continuous-trace $C^{*}$-algebra (\XRef{4.4.4}{4.4.4}). If $I_1\ne A$,
then $A/I_1$ has a nonzero closed two-sided ideal $I_2/I_1$ which is a
continuous-trace $C^{*}$-algebra. Continuing transfinitely, one obtains the
composition series $(I_\rho)$.

This theorem reduces, to some extent, the study of the structure of
postliminal $C^{*}$-algebras to the study of the structure of
continuous-trace $C^{*}$-algebras. The latter will be investigated more
deeply in Section~10.

\ResultHeading{4.5.6}{4.5.6.}
If a $C^{*}$-algebra $A$ has a composition series
$(I_\rho)_{0\leq\rho\leq\alpha}$, the $\widehat I_\rho$ are increasing open
subsets of $\widehat A$, and, for a limit ordinal $\rho$,
\[
 \widehat I_\rho=\bigcup_{\rho'<\rho}\widehat I_{\rho'};
\]
this follows at once from \XRef{3.2.2}{3.2.2}. Thus $\widehat A$ is the
union of the $\widehat I_{\rho+1}-\widehat I_\rho$. If $A$ is postliminal
and the $I_\rho$ are chosen so that the $I_{\rho+1}/I_\rho$ have continuous
trace, then the spaces $\widehat I_{\rho+1}-\widehat I_\rho$ are locally
compact [\XRef{4.5.3}{4.5.3 (ii)}]. The existence of such a family of open
subsets of $\widehat A$ also follows easily from
Theorem~\XRef{4.4.5}{4.4.5}.

\ResultHeading{4.5.7}{4.5.7.}
Let $A$ be a separable postliminal $C^{*}$-algebra. Let
$(I_\rho)_{0\leq\rho\leq\alpha}$ be a composition series of $A$ such that
the $I_{\rho+1}/I_\rho$ are nonzero continuous-trace $C^{*}$-algebras.
Then the family $(I_\rho)$ is countable (\XRef{4.3.8}{4.3.8}). Thus
$\widehat A$ is the union of the countable family of subsets
$\widehat I_{\rho+1}-\widehat I_\rho$. Each of these subsets is the
intersection of an open subset and a closed subset of $\widehat A$.
Moreover, each space $\widehat I_{\rho+1}-\widehat I_\rho$ is locally
% Source PDF physical page 104
compact and has a countable base, since
\[
 \widehat I_{\rho+1}-\widehat I_\rho
 =\widehat{I_{\rho+1}/I_\rho}
\]
and $I_{\rho+1}/I_\rho$ is a separable $C^{*}$-algebra
(\XRef{3.3.4}{3.3.4}).

Bibliography: \cite{ref16,ref29}.

\BookSubsection[Borel Structure on the Spectrum of a Postliminal C*-Algebra]
{4.6}{Borel Structure on the Spectrum of a Postliminal $C^{*}$-Algebra}

\ResultHeading{4.6.1}{4.6.1. Proposition.}
\begin{itshape}
Let $A$ be a separable postliminal $C^{*}$-algebra. The Mackey Borel
structure on $\widehat A$ is the Borel structure underlying its topology,
and makes $\widehat A$ a standard Borel space.
\end{itshape}

Let $S_1$ be the Mackey Borel structure on $\widehat A$, and let $S_2$ be
the Borel structure underlying its topology. Let $\widehat A_1$ and
$\widehat A_2$ be the corresponding Borel spaces. By
\XRef{4.5.7}{4.5.7}, $\widehat A_2$ is the union of a sequence $(X_n)$ of
disjoint Borel subsets, each of which, equipped with the structure induced
by $S_2$, is a standard Borel space. Therefore $\widehat A_2$, which is the
sum of the Borel spaces $X_n$, is standard. On the other hand, $S_1$ is
finer than $S_2$ (\XRef{3.8.3}{3.8.3}). Since $\widehat A_1$ is a quotient
of the standard Borel space $\operatorname{Irr}(A)$
(\XRef{3.8.1}{3.8.1}), it follows that $S_1=S_2$
(\XRef{B.22}{B 22}).

Thus, when $A$ is a separable postliminal $C^{*}$-algebra, we may speak of
the Borel space $\widehat A$ without risk of ambiguity.

\ResultHeading{4.6.2}{4.6.2. Proposition.}
\begin{itshape}
Let $A$ be a separable postliminal $C^{*}$-algebra. There exists a Borel map
$\xi\mapsto\pi(\xi)$ from $\widehat A$ into $\operatorname{Irr}(A)$ such
that, for every $\xi\in\widehat A$, the representation $\pi(\xi)$ belongs
to the class $\xi$.
\end{itshape}

Let $\widehat A_p$ be the canonical image of
$\operatorname{Irr}_p(A)$ in $\widehat A$. The topological space
$\operatorname{Irr}_p(A)$ is Polish (\XRef{3.7.4}{3.7.4}), and the
canonical map $\operatorname{Irr}_p(A)\to\widehat A_p$ is continuous and
open (\XRef{3.5.8}{3.5.8}). On the other hand, $\widehat A$ is the union of
a sequence $(X_n)$ of disjoint subsets, each of which is the intersection
of an open and a closed subset of $\widehat A$, and is Hausdorff
(\XRef{4.5.7}{4.5.7}). Let $\operatorname{Irr}_{p,n}(A)$ be the inverse
image of $X_n$ in $\operatorname{Irr}_p(A)$. Then
$\operatorname{Irr}_{p,n}(A)$ is the intersection of an open and a closed
subset of $\operatorname{Irr}_p(A)$, and hence is a $G_\delta$;
consequently $\operatorname{Irr}_{p,n}(A)$ is a Polish space
(\XRef{B.15}{B 15}). The saturation in $\operatorname{Irr}_{p,n}(A)$ of an
open subset is open, since equivalence in $\operatorname{Rep}_p(A)$ is
defined by a group of homeomorphisms (\XRef{3.5.5}{3.5.5}). Finally, since
$X_n$ is Hausdorff, the equivalence classes are closed in
$\operatorname{Irr}_{p,n}(A)$. Hence, by \XRef{B.17}{B 17}, there is a
Borel subset $B_{p,n}$ of $\operatorname{Irr}_{p,n}(A)$ which meets each
equivalence class of $\operatorname{Irr}_{p,n}(A)$ in exactly one point.
The union $B$ of the $B_{p,n}$ is a Borel subset of
$\operatorname{Irr}(A)$ which meets each equivalence class of
$\operatorname{Irr}(A)$ in exactly one point. Let $\pi$ be the map from
$\widehat A$ into $\operatorname{Irr}(A)$ such that $\pi(\xi)$ is the
unique representative of $\xi$ belonging
% Source PDF physical page 105
to $B$. Let $\psi$ be the canonical map
$\operatorname{Irr}(A)\to\widehat A$. Since $\psi|_B$ is a Borel bijection
of the standard Borel space $B$ onto the standard Borel space $\widehat A$
(\XRef{4.6.1}{4.6.1}), $\psi|_B$ is a Borel isomorphism of $B$ onto
$\widehat A$ (\XRef{B.21}{B 21}). Therefore $\pi$ is a Borel map from
$\widehat A$ into $\operatorname{Irr}(A)$.

Bibliography: \cite{ref12,ref15,ref28}.

\BookSubsection{4.7}{Complements}

\ResultHeading{4.7.1}{4.7.1.}
Let $H$ be a Hilbert space, $t$ a trace-class operator on $H$, and $f_t$ the
linear form $x\mapsto\operatorname{Tr}(tx)$ on $\mathcal{LC}(H)$. Then
\[
 \lVert f_t\rVert=\operatorname{Tr}|t|.
\]

\ResultHeading{4.7.2}{*4.7.2.}
Let $A$ be a separable $C^{*}$-algebra and $\pi$ an irreducible
representation of $A$. If
$\pi(A)\cap\mathcal{LC}(H_\pi)=0$, there exists a family, of the cardinality
of the continuum, of pairwise inequivalent irreducible representations of
$A$ having the same kernel as $\pi$. \cite{ref199}

\ResultHeading{4.7.3}{4.7.3.}
If a separable $C^{*}$-algebra $A$ is such that $\widehat A$ consists of one
point, then $A$ is isomorphic to $\mathcal{LC}(H)$ for some Hilbert space
$H$. (Use \XRef{4.7.2}{4.7.2}.) \cite{ref106,ref127} Problem: can the
separability hypothesis be omitted?

\ResultHeading{4.7.4}{4.7.4.}
\begin{enumerate}[label=\textit{\alph*.}]
\item Let $A$ be a $C^{*}$-algebra and $I$ a closed two-sided ideal of $A$.
If $I$ and $A/I$ are antiliminal, then $A$ is antiliminal.
\item A $C^{*}$-subalgebra of an antiliminal $C^{*}$-algebra $A$ need not
be antiliminal (it may be commutative!). But every closed two-sided ideal
of $A$ is antiliminal.
\item A quotient $C^{*}$-algebra of an antiliminal $C^{*}$-algebra need not
be antiliminal. [By \XRef{4.3.9}{4.3.9}, it suffices to construct a
nonunital antiliminal $C^{*}$-algebra. For example, if $B$ is an
antiliminal $C^{*}$-algebra, one may form the restricted product of the
sequence $(B,B,B,\ldots)$.] \cite{ref46}
\end{enumerate}

\ResultHeading{4.7.5}{*4.7.5.}
Some antiliminal $C^{*}$-algebras $A$ have a family $(\pi_i)$ of
finite-dimensional irreducible representations for which
\begin{equation*}
 \bigcap_i\ker\pi_i=0.\tag*{\cite{ref57}}
\end{equation*}

\ResultHeading{4.7.6}{4.7.6.}
For $i=1,2,\ldots$, let $H_i$ be a Hilbert space of dimension $i$. Then the
product $C^{*}$-algebra of the $\mathcal L(H_i)$ is neither postliminal nor
antiliminal. \cite{ref77}

\ResultHeading{4.7.7}{4.7.7.}
\begin{enumerate}[label=\textit{\alph*.}]
\item Let $p_1<p_2<\cdots$ be a sequence of integers such that $p_i$
divides $p_{i+1}$, and let $H$ be an infinite-dimensional Hilbert space.
There exists a sequence of $C^{*}$-subalgebras
$A_1,A_2,\ldots$ of $\mathcal L(H)$ such that:
\begin{enumerate}[label=(\roman*)]
\item $A_i$ is a factor of type $\mathrm I_{p_i}$;
\item $A_i\subset A_{i+1}$.
\end{enumerate}
Let $A$ be the closure of $\bigcup_iA_i$. Then $A$ is a unital, separable,
antiliminal $C^{*}$-algebra whose only closed two-sided ideals are $0$ and
$A$. The data of the $p_i$ determine $A$ up to isomorphism.
% Source PDF physical page 106
\item[*\textit{b}.] There exist sequences $(p_i)$ and $(p'_i)$ which give
rise to nonisomorphic $C^{*}$-algebras. \cite{ref44}
\end{enumerate}

\ResultHeading{4.7.8}{4.7.8.}
Let $A$ be a postliminal $C^{*}$-algebra and $I$ its largest closed
two-sided liminal ideal. Then $\widehat I$ is dense in $\widehat A$. (Use
\XRef{1.9.12}{1.9.12 b}.) Every $\pi\in\widehat I$ has as its kernel a
minimal primitive ideal of $A$. Problem: is every minimal primitive ideal
of $A$ the kernel of some $\pi\in\widehat I$? \cite{ref11}

\ResultHeading{4.7.9}{*4.7.9.}
There exist separable liminal $C^{*}$-algebras $A$ such that the set of
nonseparated points of $\widehat A$ (\XRef{3.9.4}{3.9.4}) is dense in
$\widehat A$. There exist liminal $C^{*}$-algebras $A$ such that
$\widehat A$ has no separated point. \cite{ref14} Problem: if $A$ is
postliminal, is the set of separated points of $\widehat A$ contained in a
Hausdorff open subset of $\widehat A$?

\ResultHeading{4.7.10}{4.7.10.}
Let $A$ be a separable liminal $C^{*}$-algebra and $X$ a quasi-compact
subset of $\widehat A$. Then $X$ is a $G_\delta$ subset of $\widehat A$.
[Let $(U_1,U_2,\ldots)$ be a base for the topology of $\widehat A$. For
every $\pi\in\widehat A-X$, there is a finite covering of $X$ by finitely
many of the $U_i$ which do not contain $\pi$.] By contrast, if $A$ is
separable and postliminal, a quasi-compact subset of $\widehat A$ need not
be Borel in $\widehat A$.

\ResultHeading{4.7.11}{*4.7.11.}
Let $A$ be a $C^{*}$-algebra. The set $I$ of $x\in A$ such that the rank of
$\pi(x)$ is finite and bounded as $\pi$ ranges over $\widehat A$ is a
self-adjoint two-sided ideal of $A$. Suppose that $I$ is dense in $A$. Let
$(\pi_i)$ be a family of elements of $\widehat A$ indexed by a directed
set. Let $\rho_1,\ldots,\rho_r\in\widehat A$ be such that, for every
$x\in I$,
\[
 \lim_i\operatorname{Tr}\pi_i(x)
 =\sum_{k=1}^r\operatorname{Tr}\rho_k(x).
\]
Then $\{\rho_1,\ldots,\rho_r\}$ is the set of limits of $(\pi_i)$ in
$\widehat A$. \cite{ref27} The $C^{*}$-algebra $A$ is liminal, and
$\widehat A$ has an open covering $(U_i)$ such that, for every $i$, every
family of elements of $\widehat A$ indexed by a directed set has at most a
finite number of distinct limits in $U_i$. (Fell, unpublished.)

\ResultHeading{4.7.12}{4.7.12.}
\begin{enumerate}[label=\textit{\alph*.}]
\item For every $C^{*}$-algebra $A$, put
$J(A)=\overline{\mathfrak m(A)}$ (\XRef{4.5.2}{4.5.2}). There exist an
ordinal $\alpha$ and an increasing family
$(J_\rho)_{0\leq\rho\leq\alpha}$ of closed two-sided ideals of $A$ with the
following properties:
\begin{enumerate}[label=(\alph*)]
\item $J_0=0$ and $J(A/J_\alpha)=0$;
\item if $\rho\leq\alpha$ is a limit ordinal, then $J_\rho$ is the closure
of $\bigcup_{\rho'<\rho}J_{\rho'}$;
\item if $\rho<\alpha$, then
$J_{\rho+1}/J_\rho=J(A/J_\rho)\ne0$.
\end{enumerate}
Moreover, $\alpha$ and the family $(J_\rho)$ are uniquely determined by
the preceding conditions. Put $J_\alpha=K(A)$.
\item For $K(A)=A$, it is necessary and sufficient that
$J(A/I)\ne0$ for every closed two-sided ideal $I\ne A$ of $A$. In this
case $A$ is called a \emph{generalized continuous-trace $C^{*}$-algebra}.
% Source PDF physical page 107
\item Suppose that $A$ has generalized continuous trace. For every
$\rho<\alpha$, the set
$\widehat J_{\rho+1}-\widehat J_\rho$ is open and dense in
$\widehat A-\widehat J_\rho$, and each point of
$\widehat J_{\rho+1}-\widehat J_\rho$ has in
$\widehat A-\widehat J_\rho$ a fundamental system of closed
neighborhoods. The algebra $A$ is liminal (use \XRef{4.7.15}{4.7.15}), but
a liminal algebra need not have generalized continuous trace
(\XRef{4.7.9}{4.7.9}).
\item Let $A$ be a postliminal $C^{*}$-algebra, let $\alpha$ be an ordinal,
and let $(V_\rho)_{0\leq\rho\leq\alpha}$ be an increasing family of open
subsets of $\widehat A$ with the following properties:
\begin{enumerate}[label=\arabic*\textsuperscript{o}]
\item $V_0=\varnothing$ and $V_\alpha=\widehat A$;
\item if $\rho\leq\alpha$ is a limit ordinal, then
$V_\rho=\bigcup_{\rho'<\rho}V_{\rho'}$;
\item every point of $V_{\rho+1}-V_\rho$ has in
$\widehat A-V_\rho$ a fundamental system of closed neighborhoods.
\end{enumerate}
Then $A$ has generalized continuous trace. In particular, if $\widehat A$
is Hausdorff, then $A$ has generalized continuous trace. \cite{ref16}
\end{enumerate}

\ResultHeading{4.7.13}{4.7.13.}
Let $A$ be a continuous-trace $C^{*}$-algebra, and let $x\in A$ be such
that the rank of $\pi(x)$ is finite and bounded as $\pi$ ranges over
$\widehat A$. Then the function
$\pi\mapsto\operatorname{Tr}\pi(x)$ is continuous on $\widehat A$.
\cite{ref29}

\ResultHeading{4.7.14}{4.7.14.}
\begin{enumerate}[label=\textit{\alph*.}]
\item If $A$ is a liminal $C^{*}$-algebra, the $C^{*}$-algebra
$\widetilde A$ obtained from $A$ by adjoining an identity element need not
be liminal. [Example: $A=\mathcal{LC}(H)$, with $H$ an
infinite-dimensional Hilbert space. In this example the spectrum of
$\widetilde A$ consists of two points $\pi_1,\pi_2$, and its closed subsets
are $\varnothing$, $\{\pi_1\}$, and $\{\pi_1,\pi_2\}$.]
\item If a liminal $C^{*}$-algebra $A$ has an identity element, all
its irreducible representations are finite-dimensional. \cite{ref77}
\end{enumerate}

\ResultHeading{4.7.15}{4.7.15.}
Let $A$ be a postliminal $C^{*}$-algebra. For $A$ to be liminal, it is
necessary and sufficient that every point of $\widehat A$ be closed. (Use
\XRef{4.3.7}{4.3.7 (i)}.) \cite{ref46}

\ResultHeading{4.7.16}{*4.7.16.}
Let $A$ be a $C^{*}$-algebra.
\begin{enumerate}[label=\textit{\alph*.}]
\item For every weak limit of pure states to be proportional to a pure
state, it is necessary and sufficient that the following conditions hold:
\begin{enumerate}[label=(\roman*)]
\item $A$ is liminal;
\item $\widehat A$ is Hausdorff;
\item $A/J(A)$ (\XRef{4.7.12}{4.7.12}) is commutative.
\end{enumerate}
\cite{ref46}
\item Suppose that $A$ is unital and has no nonzero one-dimensional
irreducible representation. For $P(A)$ to be weakly closed in the dual of
$A$, it is necessary and sufficient that $A$ be a product of finitely many
$C^{*}$-algebras $A_1,\ldots,A_n$ having the following property: for every
$\pi\in\widehat A_i$, $\dim\pi$ is finite and independent of $\pi$.
\cite{ref169}
\end{enumerate}

\ResultHeading{4.7.17}{4.7.17.}
Let $\alpha$ be an ordinal. Let $E$ be the set of maps
$f\colon[0,\alpha]\to\mathbb N$ such that $f(\rho)=0$ except for finitely
many $\rho$. Let $H=\ell^2_E$, with its canonical orthonormal basis
$(\xi_f)_{f\in E}$. If $S\subset[0,\alpha]$, let $E(S)$ be the set of
$f\in E$ which vanish outside $S$. Let $\rho\leq\alpha$. We identify $E$
with $E([0,\rho[)\times E([\rho,\alpha])$. For
$g\in E([0,\rho[)$, let $H_g$ be the Hilbert subspace of $H$ with
orthonormal basis
$(\xi_{(g,h)})_{h\in E([\rho,\alpha])}$; for
$g,g'\in E([0,\rho[)$, let $U_{g,g'}$
% Source PDF physical page 108
be the isomorphism from $H_g$ onto $H_{g'}$ which maps
$\xi_{(g,h)}$ to $\xi_{(g',h)}$ for every $h\in E([\rho,\alpha])$. Let
$B_\rho$ be the set of $x\in\mathcal L(H)$ such that:
\begin{enumerate}[label=\arabic*\textsuperscript{o}]
\item $x(H_g)\subset H_g$ for every $g\in E([0,\rho[)$;
\item the operators $x|_{H_g}$ are compact and are carried into one
another by the $U_{g,g'}$.
\end{enumerate}
Let $C_\rho$ be the norm closure of $\sum_{\sigma<\rho}B_\sigma$. The
$C_\rho$ are postliminal. For $C_\alpha$, the series of ideals in
\XRef{4.3.3}{4.3.3} is
$(C_\rho)_{0\leq\rho\leq\alpha}$.

\ResultHeading{4.7.18}{4.7.18.}
Let $H$ be a Hilbert space and let $A$ be a commutative von Neumann algebra
on $H$ not contained in $\mathcal{LC}(H)$. Then
$B=\mathcal{LC}(H)+A$ is postliminal but not liminal.

\ResultHeading{4.7.19}{4.7.19.}
Let $A$ be the $C^{*}$-algebra of complex two-by-two matrices. Let $B$ be
the $C^{*}$-algebra of sequences $x=(x_1,x_2,\ldots)$ of elements of $A$
such that $x_n$ tends, as $n\to+\infty$, to a diagonal matrix
\[
 \begin{pmatrix}\lambda(x)&0\\0&\mu(x)\end{pmatrix}.
\]
Then $\widehat B$ consists of the maps $x\mapsto x_n$, $x\mapsto\lambda(x)$,
and $x\mapsto\mu(x)$. The topological space $\widehat B$ is not Hausdorff.
The algebra $B$ has generalized continuous trace
(\XRef{4.7.12}{4.7.12}). \cite{ref74}

\ResultHeading{4.7.20}{4.7.20.}
Let $A$ be a $C^{*}$-algebra. For every subset $M$ of $A$, let $R(M)$
[respectively $L(M)$] denote the set of $x\in A$ such that $Mx=0$
[respectively $xM=0$]. The following conditions are equivalent:
\begin{enumerate}[label=(\roman*)]
\item \textup{[respectively (i$'$)]}: For every closed left (respectively right)
ideal $I$ of $A$, one has
$L(R(I))=I$ [respectively $R(L(I))=I$].
\item \textup{[respectively (ii$'$)]}: The sum of the minimal left (respectively
right) ideals of $A$ is dense in $A$.
\item $A$ is isomorphic to a $C^{*}$-subalgebra of an algebra
$\mathcal{LC}(H)$.
\item There exists a family $(A_i)$ of elementary $C^{*}$-algebras such
that $A$ is isomorphic to the restricted product of the $A_i$.
\item \textup{[respectively (v$'$)]}: For every $x\in A$, left (respectively right)
multiplication by $x$ on $A$ is a weakly compact operator.
\item The spectrum of every maximal commutative $C^{*}$-subalgebra of $A$
is discrete.
\item For every hermitian $x\in A$, every nonzero point of
$\operatorname{Sp}x$ is isolated in $\operatorname{Sp}x$.
\end{enumerate}
Such a $C^{*}$-algebra is called \emph{dual}. A dual $C^{*}$-algebra is
liminal. \cite{ref77,ref118,ref119}

\ResultHeading{4.7.21}{4.7.21.}
Let $A$ be a $C^{*}$-algebra. For all the operators of left and right
multiplication on $A$ by elements of $A$ to be compact, it is necessary and
sufficient that there exist a family $(A_i)$ of finite-dimensional
$C^{*}$-algebras such that $A$ is isomorphic to the restricted product of
the $A_i$. \cite{ref74}

\ResultHeading{4.7.22}{4.7.22.}
Let $H$ be a separable infinite-dimensional Hilbert space.
\begin{enumerate}[label=\textit{\alph*.}]
\item The only closed two-sided ideals of $\mathcal L(H)$ are
$0$, $\mathcal{LC}(H)$, and $\mathcal L(H)$.
\item The $C^{*}$-algebra $\mathcal L(H)/\mathcal{LC}(H)$ is antiliminal,
nonseparable, not isomorphic to a von Neumann algebra, and has a family of
$\mathfrak c$ orthogonal projections, where $\mathfrak c$ denotes the
cardinality of the continuum.
% Source PDF physical page 109
\item Every representation of $\mathcal L(H)$ can be written
\[
 \pi=\pi_0\oplus\Bigl(\bigoplus_{\sigma\in I}\pi_\sigma\Bigr),
\]
where $\pi_0$ vanishes on $\mathcal{LC}(H)$ and each $\pi_\sigma$ is
equivalent to the identity representation of $\mathcal L(H)$. Moreover,
either $H_{\pi_0}=0$ or $H_{\pi_0}$ is nonseparable. Consequently, there is
no $C^{*}$-subalgebra of $\mathcal L(H)$ complementary to
$\mathcal{LC}(H)$. (In fact, there is no closed vector subspace of
$\mathcal L(H)$ complementary to $\mathcal{LC}(H)$.)
\item $\operatorname{Card}\widehat{\mathcal L(H)}=2^{\mathfrak c}$. (Use
\XRef{2.10.1}{2.10.1}, \XRef{2.8.6}{2.8.6}, and a commutative
$C^{*}$-subalgebra $A$ of $\mathcal L(H)$ with
$\operatorname{Card}\widehat A=2^{\mathfrak c}$.)
\cite{ref26,ref34,ref73,ref78,ref105,ref159}
\end{enumerate}

\ResultHeading{4.7.23}{4.7.23.}
Let $A$ be a unital $C^{*}$-algebra having minimal right ideals and a least
nonzero closed two-sided ideal. For every subset $P$ of $A$, let $P'$ be
the set of $x\in A$ such that $xP=0$. Let $\mathcal F$ be the set of all
$P'$ as $P$ ranges over the subsets of $A$. Suppose that, if
$J_1,J_2\in\mathcal F$ and $J_1J_2^{\ast}=0$, then
$J_1+J_2\in\mathcal F$. Under these conditions, $A$ is isomorphic to a
$C^{*}$-algebra $\mathcal L(H)$. \cite{ref183}

\ResultHeading{4.7.24}{4.7.24. Problem.}
Let $A$ be a continuous-trace $C^{*}$-algebra. Is the set of $x\in A$ such
that $\pi(x)$ has bounded rank on $\widehat A$ and vanishes outside a
compact subset of $\widehat A$ the smallest self-adjoint two-sided ideal
$\mathfrak m$ of $A$ such that $\overline{\mathfrak m}=A$?

\ResultHeading{4.7.25}{4.7.25. Problem.}
Construct a liminal $C^{*}$-algebra $A$ having no finite composition series
$(I_i)$ such that the $\widehat{I_{i+1}/I_i}$ are Hausdorff.

\BookSection{5}{Type of a Representation}

Throughout this section, $A$ denotes an involutive algebra.

\BookSubsection{5.1}{Comparison of Representations and Comparison of
Projections}

\ResultHeading{5.1.1}{5.1.1.}
Given a family $(\pi_i)$ of representations of $A$, we may always regard
the $\pi_i$ as subrepresentations of a single representation $\rho$ of $A$
(for example, $\rho=\bigoplus_i\pi_i$). Let $E_i$ be the projection in
$\rho(A)'$ corresponding to $\pi_i$. As we shall see, relations among the
$\pi_i$ can often be interpreted as relations among the projections $E_i$
in the von Neumann algebra $\rho(A)'$.

\ResultHeading{5.1.2}{5.1.2. Proposition.}
\begin{itshape}
Let $\rho$ be a representation of $A$, let $\mathcal B$ be the von Neumann
algebra generated by $\rho(A)$, and let $\pi,\pi'$ be two
subrepresentations of $\rho$ corresponding to two projections $E,E'$ of
$\mathcal B'$. Let $T$ be a continuous linear map from $E(H_\rho)$ into
$E'(H_\rho)$. Then $T$ is an intertwining operator for $\pi$ and $\pi'$ if
and only if $TE\in\mathcal B'$.
\end{itshape}

Suppose that $T$ is an intertwining operator for $\pi$ and $\pi'$. Let
$x\in A$ and $\xi\in H_\rho$. We show that
$TE\rho(x)\xi=\rho(x)TE\xi$; it will follow that $TE\in\mathcal B'$. If
$\xi\in E(H_\rho)$, then $\rho(x)\xi\in E(H_\rho)$, and hence
\[
 TE\rho(x)\xi=T\rho(x)\xi=T\pi(x)\xi
 =\pi'(x)T\xi=\rho(x)T\xi=\rho(x)TE\xi.
\]
% Source PDF physical page 110
If $\xi\in(1-E)(H_\rho)$, then
$\rho(x)\xi\in(1-E)(H_\rho)$, so
$E\rho(x)\xi=E\xi=0$; hence
\[
 TE\rho(x)\xi=0=\rho(x)TE\xi.
\]
Conversely, if $TE\in\mathcal B'$, then, for every $x\in A$ and every
$\xi\in E(H_\rho)$,
\[
 T\pi(x)\xi=TE\rho(x)\xi=\rho(x)TE\xi=\pi'(x)T\xi,
\]
so $T$ is an intertwining operator for $\pi$ and $\pi'$.

\ResultHeading{5.1.3}{5.1.3. Corollary.}
\begin{itshape}
Retain the notation $\rho,\mathcal B,\pi,\pi',E,E'$ of
\XRef{5.1.2}{5.1.2}. Then $\pi\simeq\pi'$ if and only if
$E\sim E'$ relative to the von Neumann algebra $\mathcal B'$
(\XRef{A.41}{A 41}).
\end{itshape}

Indeed, $\pi\simeq\pi'$ means that there exists an intertwining operator
for $\pi$ and $\pi'$ which is an isomorphism from $E(H_\rho)$ onto
$E'(H_\rho)$. By \XRef{5.1.2}{5.1.2}, this is equivalent to the existence
of a partially isometric operator in $\mathcal B'$ whose initial and final
projections are $E$ and $E'$.

\ResultHeading{5.1.4}{5.1.4. Corollary.}
\begin{itshape}
Retain the notation $\rho,\mathcal B,\pi,\pi',E,E'$ of
\XRef{5.1.2}{5.1.2}. Then $\pi\leq\pi'$ if and only if $E\prec E'$ relative
to the von Neumann algebra $\mathcal B'$ (\XRef{A.41}{A 41}).
\end{itshape}

Indeed, $\pi\leq\pi'$ means that there exists a projection
$E'_1\in\mathcal B'$ dominated by $E'$ such that $\pi$ is equivalent to
the subrepresentation of $\rho$ defined by $E'_1$; by
\XRef{5.1.3}{5.1.3}, this means that $E\sim E'_1$ relative to
$\mathcal B'$.

\ResultHeading{5.1.5}{5.1.5. Corollary.}
\begin{itshape}
Let $\pi$ and $\pi_1$ be two representations of $A$. If
$\pi\leq\pi_1$ and $\pi_1\leq\pi$, then $\pi\simeq\pi_1$.
\end{itshape}

This follows from \XRef{5.1.1}{5.1.1}, \XRef{5.1.3}{5.1.3},
\XRef{5.1.4}{5.1.4}, and \XRef{A.42}{A 42}.

Bibliography: \cite{ref89,ref90}.

\BookSubsection{5.2}{Disjointness}

\ResultHeading{5.2.1}{5.2.1. Proposition.}
\begin{itshape}
Let $\pi$ and $\pi_1$ be two representations of $A$. The following
conditions are equivalent:
\begin{enumerate}[label=(\roman*)]
\item the only intertwining operator for $\pi$ and $\pi_1$ is $0$;
\item if $\tau$ (respectively $\tau_1$) is a subrepresentation of $\pi$
(respectively $\pi_1$) on a nonzero space, then $\tau$ and $\tau_1$ are
inequivalent.
\end{enumerate}
Suppose, in addition, that $\pi$ and $\pi_1$ are subrepresentations of a
representation $\rho$ of $A$. Let $\mathcal B$ be the von Neumann algebra
generated by $\rho(A)$, and let $E,E_1$ be the projections in $\mathcal B'$
corresponding to $\pi,\pi_1$. Then conditions (i) and (ii) are also
equivalent to the following condition:
\begin{enumerate}[label=(\roman*),start=3]
\item the central supports of $E$ and $E_1$ in $\mathcal B'$ are
orthogonal.
\end{enumerate}
\end{itshape}

We may suppose from the outset that $\pi$ and $\pi_1$ are
subrepresentations of $\rho$ (\XRef{5.1.1}{5.1.1}). We prove
(i) $\Rightarrow$ (ii) $\Rightarrow$ (iii) $\Rightarrow$ (i).

% Source PDF physical page 111
(i) $\Rightarrow$ (ii): if $\pi$ and $\pi_1$ have equivalent
subrepresentations $\tau$ and $\tau_1$ on nonzero spaces, then there is a
nonzero intertwining operator for $\pi$ and $\pi_1$, namely a partially
isometric operator whose initial and final subspaces are $H_\tau$ and
$H_{\tau_1}$.

(ii) $\Rightarrow$ (iii): if the central supports of $E$ and $E_1$ are not
orthogonal, there exist nonzero projections $F,F_1\in\mathcal B'$ which
are equivalent relative to $\mathcal B'$ and satisfy $F\leq E$ and
$F_1\leq E_1$ (\XRef{A.44}{A 44}). The corresponding subrepresentations of
$\pi$ and $\pi_1$ are equivalent (\XRef{5.1.3}{5.1.3}).

(iii) $\Rightarrow$ (i): suppose that the central supports $G,G_1$ of
$E,E_1$ are orthogonal. Let $T$ be an intertwining operator for $\pi$ and
$\pi_1$, and put $S=TE\in\mathcal B'$ (\XRef{5.1.2}{5.1.2}). The
projections onto $H_\pi\ominus\ker S$ and onto $\overline{S(H_\pi)}$ belong
to $\mathcal B'$ and are equivalent relative to $\mathcal B'$
(\XRef{A.45}{A 45}), and therefore have the same central support
(\XRef{A.43}{A 43}). But the projection onto
$H_\pi\ominus\ker S$ is dominated by $E$, hence by $G$, while the
projection onto $\overline{S(H_\pi)}$ is dominated by $E_1$, hence by
$G_1$. Since $GG_1=0$, it follows that $S(H_\pi)=0$, and hence
$S=0$ and $T=0$.

\ResultHeading{5.2.2}{5.2.2. Definition.}
\begin{itshape}
If two representations $\pi,\pi_1$ of $A$ satisfy the equivalent
conditions of \XRef{5.2.1}{5.2.1}, they are said to be disjoint.
\end{itshape}

We then write $\pi\perp\pi_1$. This relation is symmetric. If $\pi$ and
$\pi_1$ are topologically irreducible, to say that they are disjoint means
simply that they are inequivalent (\XRef{2.3.4}{2.3.4}). If $\pi$ and
$\pi_1$ are finite-dimensional representations, to say that they are
disjoint means that no irreducible representation occurs in both of their
decompositions into irreducible representations; this follows from
condition (ii) of \XRef{5.2.1}{5.2.1}.

\ResultHeading{5.2.3}{5.2.3. Proposition.}
\begin{itshape}
Let $\pi$ be a representation of $A$ and $(\pi_i)_{i\in I}$ a family of
representations of $A$. If $\pi$ and $\pi_i$ are disjoint for every
$i\in I$, then $\pi$ and $\bigoplus_i\pi_i$ are disjoint.
\end{itshape}

Indeed, let $H_i$ be the space of $\pi_i$, let $H$ be the Hilbert sum of
the $H_i$, and let $P_i$ be the projection of $H$ onto $H_i$. Let $T$ be
an intertwining operator for $\pi$ and $\bigoplus_i\pi_i$. For every $i$,
$P_iT$ is an intertwining operator for $\pi$ and $\pi_i$, and therefore
vanishes. Thus $T=0$, so $\pi$ and $\bigoplus_i\pi_i$ are disjoint.

\ResultHeading{5.2.4}{5.2.4. Proposition.}
\begin{itshape}
Let $\rho$ be a representation of $A$, let $\mathcal B$ be the von Neumann
algebra generated by $\rho(A)$, and let $E$ be a projection in
$\mathcal B'$. Let $\pi,\pi_1$ be the subrepresentations of $\rho$
corresponding to $E,1-E$. Then $\pi$ and $\pi_1$ are disjoint if and only
if $E$ belongs to the center of $\mathcal B$.
\end{itshape}

Let $G,G_1$ be the central supports of $E,1-E$. We have
$G\geq E$ and $G_1\geq1-E$. Therefore
\begin{align*}
 \pi\perp\pi_1
 &\Longleftrightarrow GG_1=0 \\
 &\Longleftrightarrow G=E\ \text{and}\ G_1=1-E \\
 &\Longleftrightarrow E\ \text{belongs to the center of }\mathcal B.
\end{align*}

% Source PDF physical page 112
\ResultHeading{5.2.5}{5.2.5. Corollary.}
\begin{itshape}
Let $\rho$ be a representation of $A$, and let $\mathcal B$ be the von
Neumann algebra generated by $\rho(A)$. The following conditions are
equivalent:
\begin{enumerate}[label=(\roman*)]
\item it is impossible to write $\rho$ as a Hilbert sum of two disjoint
subrepresentations on nonzero spaces;
\item $\mathcal B$ is a factor;
\item $\mathcal B'=\rho(A)'$ is a factor.
\end{enumerate}
\end{itshape}

By \XRef{5.2.4}{5.2.4}, condition (i) says that the only projections in the
common center of $\mathcal B$ and $\mathcal B'$ are $0$ and $1$, that is,
that this center consists only of scalars.

\ResultHeading{5.2.6}{5.2.6. Definition.}
\begin{itshape}
A representation of $A$ satisfying the conditions of
\XRef{5.2.5}{5.2.5} is called a factor representation.
\end{itshape}

By \XRef{2.2.6}{2.2.6}, such a representation is either zero or
nondegenerate.

If $\rho$ is topologically irreducible, it is clear that $\rho$ is a factor
representation. The converse is false. For example, a Hilbert sum of
equivalent irreducible representations is a factor representation
(\XRef{5.3.4}{5.3.4}); for finite-dimensional representations, this is in
fact the most general example of a factor representation, by condition (i)
of \XRef{5.2.5}{5.2.5}.

\ResultHeading{5.2.7}{5.2.7. Proposition.}
\begin{itshape}
Let $\rho$ be a representation of $A$, let $\mathcal B$ be the von Neumann
algebra generated by $\rho(A)$, let $\mathfrak Z$ be the center of
$\mathcal B$, let $E\ne0$ be a projection in $\mathcal B'$, let $\pi$ be
the corresponding subrepresentation of $\rho$, and let $F$ be the central
support of $E$. Then $\pi$ is a factor representation if and only if $F$ is
a minimal projection of $\mathfrak Z$.
\end{itshape}

The von Neumann algebra generated by $\pi(A)$ is the induced von Neumann
algebra $\mathcal B_E$ (\XRef{A.15}{A 15}), which is isomorphic to
$\mathcal B_F$ (\XRef{A.20}{A 20}); the center of $\mathcal B_F$ is
$\mathfrak Z_F$ (\XRef{A.15}{A 15}). Thus $\pi$ is a factor representation
if and only if $\mathfrak Z_F$ consists only of scalars, that is, if and
only if $F$ is a minimal projection of $\mathfrak Z$
(\XRef{A.34}{A 34}).

\ResultHeading{5.2.8}{5.2.8.}
Let $\rho$ be a representation of $A$, let $\mathcal B$ be the von Neumann
algebra generated by $\rho(A)$, and let $\mathfrak Z$ be the center of
$\mathcal B$. To decompose $\rho$ as a Hilbert sum of topologically
irreducible representations is to write the identity operator on $H_\rho$
as a sum of pairwise orthogonal minimal projections of $\mathcal B'$. To
decompose $\rho$ as a Hilbert sum of pairwise disjoint factor
representations is to write $1$ as a sum of pairwise orthogonal minimal
projections of $\mathfrak Z$, by \XRef{5.2.4}{5.2.4} and
\XRef{5.2.7}{5.2.7}. The second problem, concerning $\mathfrak Z$, is in
principle easier to solve than the first, which concerns $\mathcal B'$.

% Source PDF physical page 113
Let $(F_j)$ be the family of minimal projections of $\mathfrak Z$; these
are pairwise distinct and hence pairwise orthogonal. If
$1=\sum_jF_j$, we obtain a decomposition of $\rho$ as a Hilbert sum of
pairwise disjoint factor representations. The problem of writing $\rho$ as
a Hilbert sum of topologically irreducible representations is thus reduced
to the case in which $\rho$ is a factor representation. We shall see in
\XRef{5.4.11}{5.4.11} that, in the important special case of type I
representations, this problem is then quite easy to solve. For
representations not of type I, the problem, even after this reduction, is
on the contrary very difficult.

In general, however, one has $F=\sum_jF_j\ne1$. The preceding decomposition
is effective only for a subrepresentation of $\rho$, and it may
unfortunately happen that $F=0$. Later we shall obtain a satisfactory
theory by replacing the decomposition of $\rho$ into a Hilbert sum of
pairwise disjoint factor representations by a decomposition into a direct
integral.

\ResultHeading{5.2.9}{5.2.9. Proposition.}
\begin{itshape}
Let $\pi$ and $\pi_1$ be factor representations of $A$. One of the
relations
\[
 \pi\perp\pi_1,\qquad \pi\leq\pi_1,\qquad \pi_1\leq\pi
\]
holds.
\end{itshape}

We may regard $\pi$ and $\pi_1$ as subrepresentations of a representation
$\rho$ of $A$. Let $\mathcal B$ be the von Neumann algebra generated by
$\rho(A)$, let $E,E_1$ be the projections in $\mathcal B'$ corresponding
to $\pi,\pi_1$, and let $F,F_1$ be their central supports. These are zero or
minimal projections of $\mathcal B\cap\mathcal B'$
(\XRef{5.2.7}{5.2.7}); hence $F,F_1$ are either orthogonal, in which case
$\pi\perp\pi_1$, or equal. In the latter case, relative to the factor
$\mathcal B'_F$, either $E\prec E_1$ or $E_1\prec E$
(\XRef{A.46}{A 46}); hence either $\pi\leq\pi_1$ or $\pi_1\leq\pi$
(\XRef{5.1.4}{5.1.4}).

Bibliography: \cite{ref89,ref90}.

\BookSubsection{5.3}{Quasi-Equivalence}

\ResultHeading{5.3.1}{5.3.1. Proposition.}
\begin{itshape}
Let $\pi$ and $\pi_1$ be two representations of $A$, and let
$\mathcal B,\mathcal B_1$ be the von Neumann algebras generated by
$\pi(A),\pi_1(A)$, respectively. The following conditions are equivalent:
\begin{enumerate}[label=(\roman*)]
\item no subrepresentation of $\pi$ on a nonzero space is disjoint from
$\pi_1$, and no subrepresentation of $\pi_1$ on a nonzero space is
disjoint from $\pi$;
\item there exists an isomorphism $\Phi$ of $\mathcal B$ onto
$\mathcal B_1$ such that
\[
 \pi_1(x)=\Phi(\pi(x))\qquad\text{for every }x\in A;
\]
\item there exists a multiple $\pi_2$ of $\pi$ and, in $\pi_2(A)'$, a
projection of central support $1$ such that the corresponding
subrepresentation of $\pi_2$ is equivalent to $\pi_1$;
\item there exist a multiple of $\pi$ and a multiple of $\pi_1$ which are
equivalent.
\end{enumerate}

% Source PDF physical page 114
Suppose, in addition, that $\pi$ and $\pi_1$ are subrepresentations of a
representation $\rho$ of $A$. Let $\mathcal C$ be the von Neumann algebra
generated by $\rho(A)$, and let $E,E_1$ be the projections in $\mathcal C'$
corresponding to $\pi,\pi_1$. Then conditions (i)--(iv) are also equivalent
to the following condition:
\begin{enumerate}[label=(\roman*),start=5]
\item the central supports $F,F_1$ of $E,E_1$ are equal.
\end{enumerate}
\end{itshape}

(ii) $\Rightarrow$ (iii): by \XRef{A.22}{A 22}, $\Phi$ is the composite of
an amplification $\Psi$ of $\mathcal B$ onto a von Neumann algebra
$\mathcal B_2$, an induction $\Psi'$ of $\mathcal B_2$ onto
$\mathcal B_1$ defined by a projection $E\in\mathcal B_2'$ of central
support $1$, and a spatial isomorphism $\Psi''$. Put
\[
 \pi_2(x)=\Psi(\pi(x))\qquad\text{for every }x\in A.
\]
Then $\pi_2$ is a multiple of $\pi$. Moreover,
\[
 \Psi''^{-1}(\pi_1(x))
 =\Psi''^{-1}(\Phi(\pi(x)))
 =\Psi'(\pi_2(x))
\]
is the operator induced by $\pi_2(x)$ on $E(H_{\pi_2})$. Thus $\pi_1$ is
equivalent to the subrepresentation of $\pi_2$ defined by $E$.

(iii) $\Rightarrow$ (ii): suppose condition (iii) holds, and let
$\mathcal B_2$ be the von Neumann algebra generated by $\pi_2(A)$. There is
clearly an isomorphism of $\mathcal B$ onto $\mathcal B_2$ which carries
$\pi(x)$ to $\pi_2(x)$ for every $x\in A$. Moreover, if
$E\in\mathcal B_2'$ is a projection of central support $1$, induction from
$\mathcal B_2$ to $(\mathcal B_2)_E$ is an isomorphism
(\XRef{A.20}{A 20}). Hence there is an isomorphism of $\mathcal B_2$ onto
$\mathcal B_1$ which carries $\pi_2(x)$ to $\pi_1(x)$ for every $x\in A$.

(iii) $\Rightarrow$ (iv): suppose condition (iii) holds. There exists a
cardinal $n$ such that $\pi_1\leq n\pi$, and we may suppose that $n$ is
infinite. Since (iii) $\Leftrightarrow$ (ii), condition (iii) is symmetric
in $\pi,\pi_1$; hence there also exists an infinite cardinal $n_1$ such
that $\pi\leq n_1\pi_1$. Put $p=\sup(n,n_1)$. Then
\[
 p\pi\leq pn_1\pi_1=p\pi_1\leq pn\pi=p\pi,
\]
and hence $p\pi_1\simeq p\pi$ by \XRef{5.1.5}{5.1.5}.

(iv) $\Rightarrow$ (i): suppose that $p\pi\simeq p\pi_1$. Let $\tau$ be a
subrepresentation of $\pi$ on a nonzero space. Then $\tau$ is not disjoint
from $p\pi\simeq p\pi_1$, and therefore is not disjoint from $\pi_1$
(\XRef{5.2.3}{5.2.3}). Similarly, every subrepresentation of $\pi_1$ on a
nonzero space is not disjoint from $\pi$.

Now suppose that $\pi$ and $\pi_1$ are subrepresentations of $\rho$.

(i) $\Rightarrow$ (v): if $F\ne F_1$, there is a nonzero projection $G$ in
the center of $\mathcal C$ dominated by $F$ and orthogonal to $F_1$, after
interchanging $F$ and $F_1$ if necessary. We have $GE\ne0$, for otherwise
the central support of $E$ would be orthogonal to $G$; and $GE$ is
orthogonal to $F_1$, so the central support of $GE$ is orthogonal to $F_1$.
Thus there exists a subrepresentation of $\pi$ on a nonzero space which is
disjoint from $\pi_1$.

(v) $\Rightarrow$ (ii): the map $S\mapsto S_E$ (respectively
$T\mapsto T_{E_1}$) from $\mathcal C_F$ onto $\mathcal C_E$ (respectively
from $\mathcal C_{F_1}$ onto $\mathcal C_{E_1}$) is an isomorphism
(\XRef{A.20}{A 20}). Hence, if $F=F_1$, there is an isomorphism from
$\mathcal C_E=\mathcal B$ onto $\mathcal C_{E_1}=\mathcal B_1$ which
carries $\rho(x)_E=\pi(x)$ to $\rho(x)_{E_1}=\pi_1(x)$ for every $x\in A$.

% Source PDF physical page 115
\ResultHeading{5.3.2}{5.3.2. Definition.}
\begin{itshape}
If the equivalent conditions of \XRef{5.3.1}{5.3.1} hold, $\pi$ and
$\pi_1$ are said to be quasi-equivalent. We then write
$\pi\approx\pi_1$.
\end{itshape}

It is clear from condition (ii) of \XRef{5.3.1}{5.3.1} that this is an
equivalence relation. If $\pi$ and $\pi_1$ are finite-dimensional, then
$\pi\approx\pi_1$ means that the same irreducible representations occur in
the decompositions of $\pi$ and $\pi_1$ into irreducible components, though
possibly with different multiplicities; this follows from condition (i) of
\XRef{5.3.1}{5.3.1}.

\ResultHeading{5.3.3}{5.3.3. Proposition.}
\begin{itshape}
\begin{enumerate}[label=(\roman*)]
\item If $\pi\simeq\pi_1$, then $\pi\approx\pi_1$.
\item If $\pi$ and $\pi_1$ are topologically irreducible and
$\pi\approx\pi_1$, then $\pi\simeq\pi_1$.
\end{enumerate}
\end{itshape}

Assertion (i) is clear. If $\pi$ and $\pi_1$ are topologically irreducible
and quasi-equivalent, they are not disjoint by condition (i) of
\XRef{5.3.1}{5.3.1}, and hence are equivalent by condition (ii) of
\XRef{5.2.1}{5.2.1}.

\ResultHeading{5.3.4}{5.3.4. Proposition.}
\begin{itshape}
Let $\pi$ and $\pi_1$ be two quasi-equivalent representations of $A$. If
$\pi$ is a factor representation, then $\pi_1$ is a factor representation.
\end{itshape}

This follows at once from condition (ii) of \XRef{5.3.1}{5.3.1} and
condition (ii) of \XRef{5.2.5}{5.2.5}.

On the set of irreducible representations, quasi-equivalence coincides with
equivalence (\XRef{5.3.3}{5.3.3}). On the set of factor representations,
quasi-equivalence appears to be the ``right'' equivalence relation.

\ResultHeading{5.3.5}{5.3.5. Proposition.}
\begin{itshape}
Let $\pi$ be a factor representation of $A$. Every subrepresentation of
$\pi$ on a nonzero space is quasi-equivalent to $\pi$.
\end{itshape}

Since the only nonzero projection in the center of $\pi(A)'$ is $1$, this
follows from condition (iii) of \XRef{5.3.1}{5.3.1}.

\ResultHeading{5.3.6}{5.3.6. Corollary.}
\begin{itshape}
Two factor representations of $A$ are either quasi-equivalent or disjoint.
\end{itshape}

This follows from \XRef{5.2.9}{5.2.9} and \XRef{5.3.5}{5.3.5}.

\ResultHeading{5.3.7}{5.3.7. Corollary.}
\begin{itshape}
Let $A$ be a separable involutive Banach algebra. Every factor
representation of $A$ is quasi-equivalent to a factor representation on a
separable space.
\end{itshape}

Let $\pi$ be a factor representation of $A$ and $\xi$ a nonzero vector of
$H_\pi$. Then the closure $K$ of $\pi(A)\xi$ is separable and invariant
under $\pi(A)$. The subrepresentation of $\pi$ defined by $K$ is
quasi-equivalent to $\pi$ by \XRef{5.3.5}{5.3.5}.

\ResultHeading{5.3.8}{5.3.8.}
Let $\pi,\pi_1$ be two quasi-equivalent representations of $A$ on separable
spaces. Then
\[
 \aleph_0\pi\simeq\aleph_0\pi_1.
\]
It suffices to repeat the proof of
\XRef{5.3.1}{5.3.1}, (ii) $\Rightarrow$ (iii) $\Rightarrow$ (iv), and use
\XRef{A.22}{A 22}.

% Source PDF physical page 116
The customary notation for quasi-equivalence is $\pi\sim\pi_1$. Since this
conflicts with the notation $E\sim E_1$ for projections, we have preferred
to modify it slightly.

Bibliography: \cite{ref89,ref90}.

\BookSubsection{5.4}{Representations of Type I}

\ResultHeading{5.4.1}{5.4.1. Proposition.}
\begin{itshape}
Let $\pi$ be a representation of $A$. The following conditions are
equivalent:
\begin{enumerate}[label=(\roman*)]
\item the von Neumann algebra generated by $\pi(A)$ is of type I;
\item the von Neumann algebra $\pi(A)'$ is of type I;
\item $\pi$ is quasi-equivalent to a representation $\pi_1$ such that
$\pi_1(A)'$ is commutative.
\end{enumerate}
\end{itshape}

(iii) $\Rightarrow$ (i): suppose condition (iii) holds. Since
$\pi_1(A)'$ is commutative, $\pi_1(A)$ generates a von Neumann algebra of
type I (\XRef{A.35}{A 35}). Hence $\pi(A)$ generates a von Neumann algebra
of type I, by condition (ii) of \XRef{5.3.1}{5.3.1}.

(i) $\Rightarrow$ (ii) follows from \XRef{A.52}{A 52}.

(ii) $\Rightarrow$ (iii): suppose condition (ii) holds. There exists a
projection $E\in\pi(A)'$ of central support $1$ such that $\pi(A)'_E$ is
commutative (\XRef{A.35}{A 35}). The subrepresentation $\pi_1$ of $\pi$
defined by $E$ is quasi-equivalent to $\pi$, by condition (iii) of
\XRef{5.3.1}{5.3.1}, and
$\pi_1(A)'=\pi(A)'_E$ is commutative.

\ResultHeading{5.4.2}{5.4.2. Definition.}
\begin{itshape}
If a representation $\pi$ satisfies the equivalent conditions of
\XRef{5.4.1}{5.4.1}, it is said to be of type I.
\end{itshape}

Every finite-dimensional representation is of type I. Indeed, on a
finite-dimensional Hilbert space every von Neumann algebra is of type I.

\ResultHeading{5.4.3}{5.4.3.}
Every representation quasi-equivalent to a representation of type I is of
type I. Every subrepresentation of a representation of type I is of type I,
since every von Neumann algebra induced by a von Neumann algebra of type I
is of type I (\XRef{A.53}{A 53}). Every Hilbert sum
$\rho=\bigoplus_i\pi_i$ of representations $\pi_i$ of type I is of type I.
Indeed, let $\mathcal B$ be the von Neumann algebra generated by $\rho(A)$,
let $E_i\in\mathcal B'$ be the projection corresponding to $\pi_i$, and let
$F_i$ be its central support. Then $\mathcal B_{E_i}$ is of type I, and so
$\mathcal B_{F_i}$, which is isomorphic to $\mathcal B_{E_i}$
(\XRef{A.20}{A 20}), is of type I. Therefore the largest projection $F$ of
$\mathcal B\cap\mathcal B'$ for which $\mathcal B_F$ is of type I
(\XRef{A.39}{A 39}) equals $1$, since $\sup_iF_i=1$.

\ResultHeading{5.4.4}{5.4.4.}
\begin{itshape}
Let $\pi$ be a representation of $A$. The following conditions are
equivalent:
\begin{enumerate}[label=(\roman*)]
\item $\pi(A)'$ is commutative;
\item in every decomposition $\pi=\pi_1\oplus\pi_2$, the representations
$\pi_1$ and $\pi_2$ are disjoint.
\end{enumerate}
\end{itshape}

% Source PDF physical page 117
By \XRef{5.2.4}{5.2.4}, condition (ii) says that every projection in
$\pi(A)'$ belongs to the center of $\pi(A)'$, and hence that $\pi(A)'$ is
commutative.

\ResultHeading{5.4.5}{5.4.5. Definition.}
\begin{itshape}
If $\pi$ satisfies the equivalent conditions of \XRef{5.4.4}{5.4.4}, it is
said to be multiplicity-free.
\end{itshape}

Let $\pi$ be a finite-dimensional representation, and let
\[
 \pi=p_1\pi_1\oplus\cdots\oplus p_n\pi_n
\]
be the decomposition of $\pi$ into multiples of inequivalent irreducible
representations $\pi_1,\ldots,\pi_n$. Then $\pi$ is multiplicity-free if
and only if $p_1=\cdots=p_n=1$, by condition (ii) of
\XRef{5.4.4}{5.4.4}.

\ResultHeading{5.4.6}{5.4.6.}
A multiplicity-free representation is of type I. Every representation
$\pi$ of type I is quasi-equivalent to a multiplicity-free representation
$\pi'$ [condition (iii) of \XRef{5.4.1}{5.4.1}]; moreover, $\pi'$ is
determined by $\pi$ up to equivalence. Indeed, we have the following result,
which generalizes \XRef{5.3.3}{5.3.3 (ii)}.

\ResultLabel{Proposition.}
\begin{itshape}
Let $\pi,\rho$ be two multiplicity-free representations of $A$. If
$\pi\approx\rho$, then $\pi\simeq\rho$.
\end{itshape}

Let $(E_i,F_i)_{i\in I}$ be a maximal family among the families having the
following properties:
\begin{enumerate}[label=\arabic*\textsuperscript{o}]
\item $E_i\in\pi(A)'$ and $F_i\in\rho(A)'$;
\item the $E_i$ are nonzero pairwise orthogonal projections, and the $F_i$
are nonzero pairwise orthogonal projections;
\item the subrepresentation of $\pi$ defined by $E_i$ is equivalent to the
subrepresentation of $\rho$ defined by $F_i$.
\end{enumerate}
Put $E=\sum_iE_i$ and $F=\sum_iF_i$. Then the subrepresentation $\pi_1$ of
$\pi$ defined by $E$ is equivalent to the subrepresentation $\rho_1$ of
$\rho$ defined by $F$. Suppose that $E\ne1$. Since $\pi\approx\rho$, there
exist a nonzero projection $E'\in\pi(A)'$ dominated by $1-E$ and a nonzero
projection $F'\in\rho(A)'$ such that the subrepresentations $\pi',\rho'$ of
$\pi,\rho$ defined by $E',F'$ are equivalent, by condition (i) of
\XRef{5.3.1}{5.3.1}. Since $\rho(A)'$ is commutative, we have
\[
 F'=F'_1+F'_2,
 \qquad F'_1\leq F,
 \qquad F'_2\leq1-F,
\]
where $F'_1,F'_2$ are projections in $\rho(A)'$. Hence there are
projections $E'_1,E'_2$ in $\pi(A)'$ such that $E'=E'_1+E'_2$ and such that
the subrepresentations of $\pi$ and $\rho$ defined by $E'_j,F'_j$ are
equivalent for $j=1,2$. The maximality of the family $(E_i,F_i)$ implies
that $F'_2=E'_2=0$. Thus $F'\leq F$, and consequently $\pi'$ is equivalent
to a subrepresentation of $\rho_1$, hence to a subrepresentation of
$\pi_1$, and therefore to a subrepresentation of $\pi$. But this
contradicts condition (ii) of \XRef{5.4.4}{5.4.4}. Hence $E=1$, and
similarly $F=1$. Thus $\pi\simeq\rho$.

\ResultHeading{5.4.7}{5.4.7. Proposition.}
\begin{itshape}
Let $\pi$ be a representation of $A$ and $n$ a cardinal. The following
conditions are equivalent:
\begin{enumerate}[label=(\roman*)]
\item $\pi(A)'$ is a von Neumann algebra of type $\mathrm I_n$
(\XRef{A.47}{A 47});
\item $\pi$ is a Hilbert sum of $n$ equivalent multiplicity-free
representations.
\end{enumerate}
\end{itshape}

Suppose condition (ii) holds. There exists a multiplicity-free
representation $\rho$ of $A$ such that $\pi=n\rho$. Then $H_\pi$ is
identified with the Hilbert space
% Source PDF physical page 118
$H_\rho\otimes K$, where $K$ is a Hilbert space of dimension $n$, in such a
way that $\pi(x)$ is identified with $\rho(x)\otimes1$ for every $x\in A$
(\XRef{A.17}{A 17}). Then
\[
 \pi(A)'=\rho(A)'\mathbin{\overline\otimes}\mathcal L(K)
 \qquad(\XRef{A.18}{A 18}).
\]
Since $\rho(A)'$ is commutative, $\pi(A)'$ is of type $\mathrm I_n$
(\XRef{A.47}{A 47}). The converse is proved similarly.

\ResultHeading{5.4.8}{5.4.8. Definition.}
\begin{itshape}
If $\pi$ satisfies the equivalent conditions of \XRef{5.4.7}{5.4.7}, it is
said to have multiplicity $n$.
\end{itshape}

A representation of multiplicity $n$ is of type I. A multiplicity-free
representation is simply a representation of multiplicity $1$. The number
$n$ in \XRef{5.4.7}{5.4.7} is determined by $\pi$
(\XRef{A.47}{A 47}). Moreover, if $\pi\simeq n\pi'$ with $\pi'$
multiplicity-free, the class of $\pi'$ is determined by $\pi$. Indeed, if
$\pi\simeq n\pi'_1$ with $\pi'_1$ multiplicity-free, then
$\pi'\approx\pi'_1$ by condition (iv) of \XRef{5.3.1}{5.3.1}, and hence
$\pi'\simeq\pi'_1$ by \XRef{5.4.6}{5.4.6}.

Let $\pi$ be a finite-dimensional representation, and let
\[
 \pi=p_1\pi_1\oplus\cdots\oplus p_r\pi_r
\]
be the decomposition of $\pi$ into multiples of inequivalent irreducible
representations $\pi_1,\ldots,\pi_r$. Then $\pi$ has multiplicity $n$ if
and only if $p_1=\cdots=p_r=n$, by condition (ii) of
\XRef{5.4.7}{5.4.7}.

\ResultHeading{5.4.9}{5.4.9. Proposition.}
\begin{itshape}
Let $\pi$ be a representation of type I of $A$. There exists a family
$(E_j)_{j\in J}$ of nonzero pairwise orthogonal projections in $\pi(A)'$
whose sum is $1$, such that:
\begin{enumerate}[label=\alph*.]
\item the subrepresentations $\pi_j$ of $\pi$ defined by the $E_j$ are
pairwise disjoint;
\item $\pi_j$ has multiplicity $n_j$;
\item the $n_j$ are pairwise distinct.
\end{enumerate}
The family $(E_j)_{j\in J}$ is unique up to a bijection of the index set.
The $E_j$ belong to the center of $\pi(A)'$.
\end{itshape}

Let $\mathfrak Z$ be the center of $\pi(A)'$. In view of
\XRef{5.2.4}{5.2.4}, the proposition says that there is a family
$(E_j)_{j\in J}$ of nonzero pairwise orthogonal projections in
$\mathfrak Z$, with sum $1$, such that the von Neumann algebras
$\pi(A)'_{E_j}$ are of type $\mathrm I_{n_j}$, where the $n_j$ are pairwise
distinct; and that this family $(E_j)$ is unique up to a bijection of the
index set. Since $\pi(A)'$ is of type I by condition (ii) of
\XRef{5.4.1}{5.4.1}, this follows from \XRef{A.50}{A 50}.

What precedes completely reduces the study of representations of type I to
that of multiplicity-free representations.

\ResultHeading{5.4.10}{5.4.10.}
Retain the notation of \XRef{5.4.9}{5.4.9}. Let $U$ be a unitary operator
on $H_\pi$ such that
\[
 U\pi(A)U^{-1}=\pi(A).
\]
Then $U\mathfrak ZU^{-1}=\mathfrak Z$, so the
$E'_j=UE_jU^{-1}$ are nonzero pairwise orthogonal projections in
$\mathfrak Z$, with sum $1$, such that $\pi(A)'_{E'_j}$ is of type
$\mathrm I_{n_j}$. Hence $E'_j=E_j$. In other words, $U$ commutes with the
$E_j$.

\ResultHeading{5.4.11}{5.4.11. Proposition.}
\begin{itshape}
Let $\pi$ be a representation of $A$. The following conditions are
equivalent:
\begin{enumerate}[label=(\roman*)]
\item $\pi$ is a factor representation of type I;
\item $\pi$ is quasi-equivalent to a topologically irreducible
representation;
% Source PDF physical page 119
\item $\pi$ is of the form $n\pi'$, with $\pi'$ topologically irreducible.
\end{enumerate}
Then $\pi$ has multiplicity $n$.
\end{itshape}

(iii) $\Rightarrow$ (ii) follows from condition (iv) of
\XRef{5.3.1}{5.3.1}.

(ii) $\Rightarrow$ (i) follows from \XRef{5.3.4}{5.3.4} and condition
(iii) of \XRef{5.4.1}{5.4.1}.

(i) $\Rightarrow$ (iii): if $\pi$ is a factor representation of type I,
\XRef{5.4.9}{5.4.9} shows that $\pi$ has multiplicity $n$ for some cardinal
$n$. Hence $\pi=n\rho$, with $\rho$ multiplicity-free by condition (ii) of
\XRef{5.4.7}{5.4.7}, and $\rho$ is a factor representation because it is
quasi-equivalent to $\pi$. Thus $\rho(A)'$ is a commutative factor, and
hence consists only of scalars; therefore $\rho$ is topologically
irreducible.

\ResultHeading{5.4.12}{5.4.12.}
Let $\pi$ be a representation of $A$, and let $\mathfrak Z$ be the center
of $\pi(A)'$. We have seen that $\pi$ is a Hilbert sum of pairwise disjoint
factor representations if and only if the identity operator is a sum of
minimal projections of $\mathfrak Z$. Although this situation is far from
being the general one, it is useful to interpret in this case the
definitions in \XRef{5.4.2}{5.4.2}, \XRef{5.4.5}{5.4.5}, and
\XRef{5.4.8}{5.4.8}.

\ResultLabel{Proposition.}
\begin{itshape}
Let $(\pi_i)$ be a family of pairwise disjoint factor representations of
$A$, and let $\pi=\bigoplus_i\pi_i$.
\begin{enumerate}[label=(\roman*)]
\item $\pi$ is of type I if and only if every $\pi_i$ is of type I.
\item $\pi$ has multiplicity $n$ if and only if every $\pi_i$ has
multiplicity $n$.
\item $\pi$ is multiplicity-free if and only if every $\pi_i$ is
topologically irreducible.
\end{enumerate}
\end{itshape}

Assertion (i) follows from \XRef{5.4.3}{5.4.3}; assertion (ii) follows from
the fact that a product of von Neumann algebras $\mathcal B_i$ is of type
$\mathrm I_n$ if and only if each $\mathcal B_i$ is of type
$\mathrm I_n$ (\XRef{A.50}{A 50}); assertion (iii) is a special case of
(ii).

\ResultHeading{5.4.13}{5.4.13.}
We saw in \XRef{2.3.5}{2.3.5} that a finite-dimensional representation is
a Hilbert sum of irreducible representations. More generally:

\ResultLabel{Proposition.}
\begin{itshape}
Let $\pi$ be a representation of $A$ such that $\pi(x)$ is compact for
every $x\in A$.
\begin{enumerate}[label=(\roman*)]
\item The representation $\pi$ is a Hilbert sum of a family
$(\pi_i)_{i\in I}$ of irreducible representations.
\item If $i\in I$ and $\pi_i(A)\ne0$, there are only finitely many indices
$j\in I$ such that $\pi_j\simeq\pi_i$.
\end{enumerate}
\end{itshape}

We show that, if $H_\pi\ne0$, then $\pi$ has an irreducible
subrepresentation. We may suppose that $\pi\ne0$. Let $\mathfrak B$ be the
von Neumann algebra generated by $\pi(A)$, and let $\mathfrak Z$ be its
center. There exists a hermitian $x\in A$ such that $\pi(x)\ne0$. Since
$\pi(x)$ is compact, it has a nonzero eigenvalue, and the corresponding
spectral projection $E$ is nonzero and of finite rank. We have
$E\in\mathfrak B$. Let $F$ be a projection of minimum rank among the
nonzero projections of $\mathfrak B$ which are dominated by $E$. Then $F$ is a
minimal projection
% ===== ASSEMBLED TRANSLATION CHUNK 05: chunk_120_149.tex =====
% Source PDF physical page 120
of $\mathfrak B$. Its central support $F'$ is a minimal projection of
$\mathfrak Z$. The induced von Neumann algebra $\mathfrak B_{F'}$ is therefore
a factor possessing a minimal projection, hence a factor of type I
(\XRef{A.36}{A 36}). Thus $\pi$ has a factorial subrepresentation of type I
on a nonzero space, and consequently, by \XRef{5.4.11}{5.4.11}, an
irreducible subrepresentation on a nonzero space.

This being established, let $(E_i)_{i\in I}$ be a maximal family of nonzero,
pairwise orthogonal projections of $\pi(A)'$ such that the corresponding
subrepresentations of $\pi$ are irreducible. If $\sum_iE_i\ne1$, let $\pi'$ be
the subrepresentation of $\pi$ corresponding to $1-\sum_iE_i$. For every
$x\in A$, $\pi'(x)$ is compact. The preceding result, applied to $\pi'$,
then contradicts the maximality of $(E_i)$. Hence $\sum_iE_i=1$, which proves
(i). If $x\in A$ and $i\in I$ are such that $\pi_i(x)\ne0$, and if there are
infinitely many indices $j\in I$ such that $\pi_j\simeq\pi_i$, it is clear
that $\pi(x)$ is not compact. This proves (ii).

Bibliography: \cite{ref77,ref89,ref90}.

\BookSubsection{5.5}{Involutive Algebras of Type I}

\ResultHeading{5.5.1}{5.5.1.}
We shall see that the decomposition of representations is simpler in the case
of representations of type I. Hence the interest of the following notion.

\ResultLabel{Definition.}
\begin{itshape}
An involutive algebra is said to be \emph{of type I} if all its representations
are of type I.
\end{itshape}

For algebras of type I, the notion of a factor representation may therefore be
set aside in favor of the more natural notion of a topologically irreducible
representation (\XRef{5.4.11}{5.4.11}). Unfortunately, there are
$C^{*}$-algebras that are not of type I, and for these the consideration of
factor representations sometimes leads to simpler theorems than does the
consideration of topologically irreducible representations.

\ResultHeading{5.5.2}{5.5.2. Theorem.}
\begin{itshape}
Every postliminal $C^{*}$-algebra is of type I.
\end{itshape}

Let $A$ be a nonzero postliminal $C^{*}$-algebra. By \XRef{4.4.4}{4.4.4},
there is a nonzero $x\in A^{+}$ such that $\pi(x)$ is zero or of rank $1$ for
every $\pi\in\widehat A$. Then
$\pi(xAx)=\pi(x)\pi(A)\pi(x)$ is commutative for every $\pi\in\widehat A$;
hence $xAx$ is commutative.

Now let $\pi$ be a representation of $A$. Let $(E_i)$ be a maximal family of
nonzero, pairwise orthogonal projections in the center of $\pi(A)'$ such that
the corresponding subrepresentations of $\pi$ are of type I. By
\XRef{5.4.3}{5.4.3}, it is enough to prove that $\sum_iE_i=1$. Suppose
% Source PDF physical page 121
$E=1-\sum_iE_i\ne0$, and let $\rho$ be the subrepresentation of $\pi$ defined
by $E$. If $\rho(A)=0$, we could add $E$ to the family $(E_i)$, contrary to
its maximality. Thus $\rho(A)\ne0$. On the other hand, $\rho(A)$ is
postliminal (\XRef{4.3.5}{4.3.5}). By the beginning of the proof, there exists
$x\in A^{+}$ such that $\rho(x)$ is nonzero and
$\rho(x)\rho(A)\rho(x)$ is commutative. If $\mathfrak C$ denotes the von
Neumann algebra generated by $\rho(A)$, then
$\rho(x)\mathfrak C\rho(x)$ is commutative. Let
\[
 \rho(x)=\int_{0}^{+\infty}\lambda\,de_{\lambda}
\]
be the spectral decomposition of $\rho(x)$. We have $e_{\lambda}\in\mathfrak C$
for every $\lambda$; hence
\[
 y_{\varepsilon}=\int_{\varepsilon}^{+\infty}\lambda^{-1}\,de_{\lambda}
 \in\mathfrak C \qquad(\varepsilon>0).
\]
If $\varepsilon$ is sufficiently small,
$\rho(x)y_{\varepsilon}=y_{\varepsilon}\rho(x)$ is a nonzero projection $F$
of $\mathfrak C$. We have
$F\mathfrak C F\subset\rho(x)\mathfrak C\rho(x)$, so that
$F\mathfrak C F$ is commutative. Let $G\in\mathfrak C\cap\mathfrak C'$ be
the central support of $F$. The von Neumann algebra $\mathfrak C_G$ is of
type I (\XRef{A.35}{A 35}). The projection on $H_\pi$ which is zero on
$(1-E)(H_\pi)$ and agrees with $G$ on $E(H_\pi)$ is therefore a nonzero
element of the center of $\pi(A)'$, orthogonal to the $E_i$, and the
corresponding subrepresentation of $\pi$ is of type I. This contradicts the
maximality of $(E_i)$. Hence $\sum_iE_i=1$, and $\pi$ is of type I.

Later, in \SectionRef{9.1}{9.1}, we shall prove a converse to Theorem
\XRef{5.5.2}{5.5.2}.

\ResultHeading{5.5.3}{5.5.3. Corollary.}
\begin{itshape}
Let $A$ be a postliminal $C^{*}$-algebra. Every factor representation of $A$
is quasi-equivalent to an irreducible representation.
\end{itshape}

This follows from \XRef{5.5.2}{5.5.2} and \XRef{5.4.11}{5.4.11}.

Bibliography: \cite{ref76,ref77}.

\BookSubsection{5.6}{Representations of Types II and III}

\ResultHeading{5.6.1}{5.6.1.}
Besides von Neumann algebras of type I, there are the more complicated von
Neumann algebras of type II or type III (\XRef{A.38}{A 38}). A representation
$\pi$ of $A$ is said to be of type II (respectively, of type III) if the von
Neumann algebra generated by $\pi(A)$ is of type II (respectively, III). This
is equivalent to saying that $\pi(A)'$ is a von Neumann algebra of type II
(respectively, III) (\XRef{A.52}{A 52}).

\ResultHeading{5.6.2}{5.6.2.}
Every representation quasi-equivalent to a representation of type II
(respectively, III) is of type II (respectively, III). Every subrepresentation
of a representation of type II (respectively, III) is of type II
(respectively, III). Every Hilbert direct sum of representations of type II
(respectively, III) is of type II (respectively, III). This follows by the
same arguments as for representations of type I.

% Source PDF physical page 122
\ResultHeading{5.6.3}{5.6.3.}
A representation on a nonzero space cannot be of two different types at once
(\XRef{A.38}{A 38}). It follows that two representations of distinct types
are disjoint.

\ResultHeading{5.6.4}{5.6.4. Proposition.}
\begin{itshape}
Let $\pi$ be a representation of $A$. There exist projections
$E_{\mathrm I},E_{\mathrm {II}},E_{\mathrm {III}}\in\pi(A)'$, uniquely
determined by the following conditions:
\begin{enumerate}[label=\textup{\arabic*\textsuperscript{o}},leftmargin=3em]
 \item $E_{\mathrm I},E_{\mathrm {II}},E_{\mathrm {III}}$ are pairwise
 orthogonal and have sum $1$;
 \item the subrepresentations of $\pi$ defined by
 $E_{\mathrm I},E_{\mathrm {II}},E_{\mathrm {III}}$ are respectively of
 types I, II, and III.
\end{enumerate}
These projections belong to the center of $\pi(A)'$.
\end{itshape}

This follows from the corresponding property of von Neumann algebras
(\XRef{A.39}{A 39}).

Thus the study of arbitrary representations of $A$ is reduced to that of
representations of type I, representations of type II, and representations of
type III. This, however, does not simplify matters very much.

\ResultHeading{5.6.5}{5.6.5. Proposition.}
\begin{itshape}
Let $\pi$ be a representation of $A$. In order that $\pi$ be a Hilbert direct
sum of a representation of type II and a representation of type III, it is
necessary and sufficient that every subrepresentation of $\pi$ be a Hilbert
direct sum of two equivalent representations.
\end{itshape}

For $\pi$ to be a Hilbert direct sum of a representation of type II and a
representation of type III, it is necessary and sufficient that the projection
$E_{\mathrm I}$ of \XRef{5.6.4}{5.6.4} be zero, in other words, that the von
Neumann algebra $\pi(A)'$ be continuous (\XRef{A.39}{A 39}). For this it is
necessary and sufficient that every projection of $\pi(A)'$ be a sum of two
orthogonal equivalent projections of $\pi(A)'$ (\XRef{A.48}{A 48}). It then
suffices to apply \XRef{5.1.3}{5.1.3}.

\ResultHeading{5.6.6}{5.6.6.}
Let $\pi$ be a representation of type I of $A$, and let $C$ be the
quasi-equivalence class of $\pi$. The different equivalence classes contained
in $C$ are distinguished by multiplicities (\XRef{5.4.9}{5.4.9}). For
representations not of type I there is an analogous, but more complicated,
theory, which we shall not present. We do, however, have the following very
simple result.

\ResultLabel{Proposition.}
\begin{itshape}
Let $\pi_1,\pi_2$ be two representations of type III of $A$ on separable
spaces. Then $\pi_1\approx\pi_2$ implies $\pi_1\simeq\pi_2$.
\end{itshape}

Indeed, let $\mathfrak B_1,\mathfrak B_2$ be the von Neumann algebras
generated by $\pi_1(A)$ and $\pi_2(A)$. If $\pi_1\approx\pi_2$, there is an
isomorphism $\Phi$ of $\mathfrak B_1$ onto $\mathfrak B_2$ which transforms
$\pi_1(x)$ into $\pi_2(x)$ for every $x\in A$. But
$\mathfrak B_1,\mathfrak B_2$ are of type III and act on separable spaces.
Thus $\Phi$ is implemented by an isomorphism of $H_{\pi_1}$ onto
$H_{\pi_2}$ (\XRef{A.51}{A 51}).

\ResultHeading{5.6.7}{5.6.7.}
The classification of von Neumann algebras also includes the definition of
von Neumann algebras of type $\mathrm{II}_1$, of type
$\mathrm{II}_{\infty}$,
% Source PDF physical page 123
semifinite, and so forth. A representation $\pi$ of $A$ is said to be of type
$\mathrm{II}_1$, of type $\mathrm{II}_{\infty}$, finite, and so forth, if
$\pi(A)$ generates a von Neumann algebra of type $\mathrm{II}_1$, of type
$\mathrm{II}_{\infty}$, finite, and so forth.

Bibliography: \cite{ref89,ref90}.

\BookSubsection{5.7}{Complements}

\ResultHeading{5.7.1}{5.7.1.}
A finite-dimensional $C^{*}$-algebra $A$ is a finite product of
$C^{*}$-algebras $\mathcal L(H_i)$, where each $H_i$ is a finite-dimensional
Hilbert space. (The irreducible representations of $A$ are finite-dimensional,
and $A$ admits an injective finite-dimensional representation which is a
Hilbert direct sum of pairwise disjoint irreducible representations.)

\ResultHeading{5.7.2}{5.7.2.}
Let $A$ be an involutive algebra, and let $\pi$ and $\rho$ be representations
of $A$. There exist a projection $E$ in the center of $\pi(A)'$ and a
projection $F$ in the center of $\rho(A)'$ such that
\[
 \pi_E\approx\rho_F,\qquad \pi_{1-E}\perp\rho,
 \qquad \pi\perp\rho_{1-F}.
\]
\cite{ref90}

\ResultHeading{5.7.3}{5.7.3.}
Let $A$ be an involutive algebra and let $\pi$ be a factor representation of
$A$ on a separable space. There exists a function $m$ which assigns to every
representation $\rho$ on a separable space quasi-equivalent to $\pi$ a number
$m(\rho)\in[0,+\infty]$, with the following properties:
\begin{enumerate}[label=\textup{(\roman*)}]
 \item $m(\rho_1)=m(\rho_2)$ is equivalent to $\rho_1\simeq\rho_2$;
 \item for every sequence $\rho_1,\rho_2,\ldots$ of representations on
 separable spaces quasi-equivalent to $\pi$,
 \[
   m\!\left(\bigoplus_i\rho_i\right)=\sum_i m(\rho_i).
 \]
\end{enumerate}
Moreover, the function $m$ is unique up to multiplication by a constant
$>0$. \cite{ref90}

\ResultHeading{5.7.4}{5.7.4.}
Let $A$ be a $C^{*}$-algebra and $J$ a closed two-sided ideal of $A$. In order
that $A$ be of type I, it is necessary and sufficient that $J$ and $A/J$ be
of type I. (Use \XRef{2.11.2}{2.11.2}.) \cite{ref184}

\ResultHeading{5.7.5}{5.7.5.}
Let $\mathfrak C$ be a continuous von Neumann algebra, and let $B$ be a
$C^{*}$-algebra weakly dense in $\mathfrak C$. Then $B$ is antiliminal. (Use
\XRef{2.11.1}{2.11.1} and \XRef{5.5.2}{5.5.2}.) A von Neumann algebra of
type I can have weakly dense $C^{*}$-subalgebras which are antiliminal
(consider an irreducible representation of a factor of type $\mathrm{II}_1$).
\cite{ref45}

\ResultHeading{5.7.6}{5.7.6.}
Let $A$ be a $C^{*}$-algebra and $\pi$ a representation of $A$.

\begin{enumerate}[label=\textit{\alph*.},leftmargin=2em]
 \item A representation $\pi$ is called \emph{homogeneous} if, for every
 projection $E\in\pi(A)'$, $\pi_E$ has the same kernel as $\pi$. One then has
 $\lVert\pi(x)E\rVert=\lVert\pi(x)\rVert$ for every $x\in A$.
 \item A factor representation is homogeneous. The converse is false.
 \item For every closed two-sided ideal $I$ of $A$, let $Q(I)$ be the
 projection of $H_\pi$ onto the essential subspace of $\pi|_I$. As $I$ varies,
 the $Q(I)$ generate a von Neumann subalgebra $\mathfrak B$ of the center of
 $\pi(A)'$, called the \emph{ideal center associated with $\pi$}. In order that
 $\pi$ be homogeneous, it is necessary and sufficient that $\mathfrak B$
 reduce to the scalars.
 \item The kernel of a homogeneous representation is prime. \cite{ref20}
 \item[\textit{*e.}] For every open subset $U$ of $\widehat A$, let $I(U)$ be the closed
 two-sided ideal
% Source PDF physical page 124
 of $A$ with spectrum $U$, and put $Q'(U)=Q(I(U))$. Then
 $U\mapsto Q'(U)$ extends to a map $Q''$ from the set of Borel subsets of
 $\widehat A$ into $\mathfrak B$ such that
 \[
   Q''(X_1\cup X_2\cup\cdots)=\sum_iQ''(X_i)
 \]
 whenever $X_1,X_2,\ldots$ are pairwise disjoint Borel subsets of
 $\widehat A$. \cite{ref47}
\end{enumerate}

\ResultHeading{5.7.7}{*5.7.7.}
Let $A$ be a unital $C^{*}$-algebra. We use the notation of
\XRef{2.12.19}{2.12.19}. Let $f,g$ be positive functionals on $A$; we write
$g\prec f$ if $F_g\subset F_f$, or, equivalently, if $g\in F_f$. The
relation $F_g=F_f$ is an equivalence relation, and hence gives a notion of a
class of positive functionals.

\begin{enumerate}[label=\textit{\alph*.},leftmargin=2em]
 \item $\pi_f$ is factorial if and only if the relation $\prec$ totally
 preorders $F_f$. In that case $f$ is called \emph{primary}.
 \item Suppose that $f$ is primary. In order that $\pi_f$ be of type I, it is
 necessary and sufficient that there exist a pure state $g$ such that
 $g\prec f$. In order that $\pi_f$ be of type III, it is necessary and
 sufficient that the class of $f$ be minimal and maximal in the set of classes
 of primary positive functionals and that $f$ not be a homomorphism of $A$
 into $\mathbb C$. \cite{ref71}
\end{enumerate}

\BookSectionHere{6}{Traces and Representations}

Let $A$ be a $C^{*}$-algebra. In \SectionRef{2}{Section 2}, we associated
with every positive functional $f$ on $A$ a representation $\pi_f$. All
representations are obtained by this procedure, or at least all those that
admit a cyclic vector (recall that every nondegenerate representation is a
sum of representations admitting cyclic vectors). But the map
$f\mapsto\pi_f$ is by no means injective. We now wish to parametrize
representations bijectively. Let $\pi$ and $\pi'$ be two finite-dimensional
representations of $A$. If
\[
 \operatorname{Tr}\pi(x)=\operatorname{Tr}\pi'(x)\qquad(x\in A),
\]
then $\pi$ and $\pi'$ are equivalent. It is this fact that we shall partially
generalize. In fact, only traceable representations, or rather their
quasi-equivalence classes, can be parametrized.

\BookSubsection{6.1}{Traces}

\ResultHeading{6.1.1}{6.1.1. Definition.}
\begin{itshape}
Let $A$ be a $C^{*}$-algebra. A \emph{trace} on $A^{+}$ is a function
$f\colon A^{+}\to[0,+\infty]$ satisfying the following axioms:
\begin{enumerate}[label=\textup{(\roman*)}]
 \item if $x,y\in A^{+}$, then $f(x+y)=f(x)+f(y)$;
 \item if $x\in A^{+}$ and $\lambda\geq0$, then
 $f(\lambda x)=\lambda f(x)$ (with the convention $0\cdot+\infty=0$);
 \item if $z\in A$, then $f(zz^{\ast})=f(z^{\ast}z)$.
\end{enumerate}
A trace $f$ is said to be \emph{finite} if $f(x)<+\infty$ for every
$x\in A^{+}$.

A trace $f$ is said to be \emph{semifinite} if, for every $x\in A^{+}$,
$f(x)$ is the supremum of the numbers $f(y)$ for $y\in A^{+}$ such that
$y\leq x$ and $f(y)<+\infty$.

[In the definition of semifinite traces, it is plainly enough to consider
those $x$ for which $f(x)=+\infty$.]

% Source PDF physical page 125
We shall be concerned almost exclusively with traces on $A^{+}$ which are
lower semicontinuous.

If $A$ is a von Neumann algebra, these definitions agree with the familiar
ones (see, however, \XRef{A.28}{A 28}).
\end{itshape}

\ResultHeading{6.1.2}{6.1.2. Proposition.}
\begin{itshape}
Let $A$ be a $C^{*}$-algebra and let $f$ be a trace on $A^{+}$.
\begin{enumerate}[label=\textup{(\roman*)}]
 \item Let $\mathfrak n$ be the set of $x\in A$ such that
 $f(xx^{\ast})<+\infty$. Then $\mathfrak n$ is a self-adjoint two-sided ideal
 of $A$; the two-sided ideal $\mathfrak m=\mathfrak n^{2}$ is the set of
 linear combinations of elements of
 $\mathfrak m^{+}=\mathfrak m\cap A^{+}$, and $\mathfrak m^{+}$ is the set
 of $x\in A^{+}$ such that $f(x)<+\infty$. The ideals $\mathfrak m$ and
 $\mathfrak n$ have the same closure in $A$.
 \item There is a unique linear functional $f'$ on $\mathfrak m$ which
 agrees with $f$ on $\mathfrak m^{+}$.
 \item One has
 \begin{align*}
  f'(x^{\ast})&=\overline{f'(x)} &&\text{for every }x\in\mathfrak m,\\
  f'(zx)&=f'(xz) &&\text{for all }x\in\mathfrak m,\ z\in A,\\
  f'(uv)&=f'(vu) &&\text{for all }u,v\in\mathfrak n.
 \end{align*}
\end{enumerate}
\end{itshape}

Let $\mathfrak p$ be the set of $x\in A^{+}$ such that $f(x)<+\infty$.
Applying Lemma \XRef{4.5.1}{4.5.1} to $\mathfrak p$ gives (i). Assertion (ii)
then follows immediately from (i). The functional $f'$ is real on the
self-adjoint elements of $\mathfrak m$, so
$f'(x^{\ast})=\overline{f'(x)}$ for every $x\in\mathfrak m$. If
$u\in\mathfrak n$, then $uu^{\ast}\in\mathfrak m$ and
$u^{\ast}u\in\mathfrak m$, and hence
\[
 f'(uu^{\ast})=f(uu^{\ast})=f(u^{\ast}u)=f'(u^{\ast}u).
\]
Applying this to $u+v,u-v,u+iv,u-iv$, we deduce that
$f'(uv^{\ast})=f'(v^{\ast}u)$ for $u,v\in\mathfrak n$. Finally, let
$z\in A$ and $x\in\mathfrak m$. We have
$x=\sum_j u_jv_j$, with $u_j,v_j\in\mathfrak n$, and therefore
\begin{align*}
 f'\!\left(z\sum_j u_jv_j\right)
 &=\sum_j f'((zu_j)v_j)
  =\sum_j f'(v_j(zu_j))\\
 &=\sum_j f'((v_jz)u_j)
  =\sum_j f'(u_j(v_jz))
  =f'\!\left(\left(\sum_j u_jv_j\right)z\right).
\end{align*}
This proves the proposition.

By an abuse of language, the ideal $\mathfrak m$ will be called the
\emph{ideal of definition} of $f$, and will be denoted by $\mathfrak m_f$.

\ResultHeading{6.1.3}{6.1.3.}
If $\overline{\mathfrak m_f}=A$ and $f$ is lower semicontinuous, then $f$ is
semifinite. Indeed, let $(u_i)$ be an increasing directed approximate identity
of $A$, consisting of elements of $\mathfrak m_f$ (\XRef{1.7.2}{1.7.2}).
Let $x\in A^{+}$. The elements $x^{1/2}u_i x^{1/2}$ form an increasing
directed family of elements of $A^{+}$, bounded above by $x$
(\XRef{1.6.8}{1.6.8}), and tending to $x$; hence
\[
 \lim_i f(x^{1/2}u_i x^{1/2})\geq f(x),
\]
% Source PDF physical page 126
and consequently $f(x^{1/2}u_i x^{1/2})\to f(x)$, since
$f(x^{1/2}u_i x^{1/2})\leq f(x)$. But
\[
 x^{1/2}u_i x^{1/2}\in\mathfrak m_f,
 \qquad f(x^{1/2}u_i x^{1/2})<+\infty,
\]
which proves the assertion.

\ResultHeading{6.1.4}{6.1.4.}
Let $A$ and $B$ be two $C^{*}$-algebras, let $\rho$ be a homomorphism from
$A$ into $B$, and let $f$ be a trace on $B^{+}$. Then
$(f\circ\rho)|_{A^{+}}$ is a trace on $A^{+}$. If $f$ is lower
semicontinuous, $(f\circ\rho)|_{A^{+}}$ is lower semicontinuous. If $f$ is
finite, $(f\circ\rho)|_{A^{+}}$ is finite. If $f$ is semifinite,
$(f\circ\rho)|_{A^{+}}$ need not be semifinite in general. Nevertheless, if
$A$ is a closed two-sided ideal of $B$ and $\rho$ is the canonical injection
of $A$ into $B$, then $(f\circ\rho)|_{A^{+}}=f|_{A^{+}}$ is semifinite; for
if $x\in A^{+}$ and if $y\in B^{+}$ is bounded above by $x$, then $y\in A^{+}$,
as is seen by considering the canonical map of $B$ onto $B/A$.

\ResultHeading{6.1.5}{6.1.5. Lemma.}
\begin{itshape}
Let $A$ be a $C^{*}$-algebra, $\pi$ a representation of $A$, $\mathfrak U$
the von Neumann algebra generated by $\pi(A)$, $t$ a normal trace on
$\mathfrak U^{+}$, and $\mathfrak m$ a self-adjoint two-sided ideal of $A$
such that $\pi(\mathfrak m)$ is contained in $\mathfrak m_t$ and is strongly
dense in $\mathfrak U$. Then
$f=(t\circ\pi)|_{A^{+}}$ is semifinite and lower semicontinuous; moreover, if
$(u_i)$ is an increasing directed approximate identity of
$\overline{\mathfrak m}$, then
\[
 f(x)=\lim_i f(x^{1/2}u_i x^{1/2})
 \qquad(x\in A^{+}).
\]
\end{itshape}

Since $\pi(\mathfrak m)$ is strongly dense in $\mathfrak U$,
$\pi|_{\overline{\mathfrak m}}$ is nondegenerate, and hence $\pi(u_i)$
converges strongly to $1$ (\XRef{2.2.10}{2.2.10}). Let $x\in A^{+}$. The
operators
$\pi(x)^{1/2}\pi(u_i)\pi(x)^{1/2}$ form an increasing directed family in
$\mathfrak U^{+}$ which converges strongly to $\pi(x)$. Since $t$ is normal,
\[
 t\bigl(\pi(x)^{1/2}\pi(u_i)\pi(x)^{1/2}\bigr)\longrightarrow t(\pi(x)),
\]
that is,
$f(x^{1/2}u_i x^{1/2})\to f(x)$. We may take the $u_i$ in
$\mathfrak m$ (\XRef{1.7.2}{1.7.2}); then
\[
 \pi(x^{1/2}u_i x^{1/2})\in\pi(\mathfrak m)\subset\mathfrak m_t,
 \qquad f(x^{1/2}u_i x^{1/2})<+\infty,
\]
so $f$ is semifinite. Since $t$ is lower semicontinuous for the weak topology,
and \emph{a fortiori} for the topology induced by the norm, $f$ is lower
semicontinuous.

Bibliography: \cite{ref16}.

\BookSubsection{6.2}{Bitraces}

\ResultHeading{6.2.1}{6.2.1. Definition.}
\begin{itshape}
Let $A$ be a $C^{*}$-algebra. A \emph{bitrace} of $A$ is a function
$s\colon\mathfrak n\times\mathfrak n\to\mathbb C$, where $\mathfrak n$ is a
self-adjoint two-sided ideal of $A$, satisfying the following axioms:
\begin{enumerate}[label=\textup{(\roman*)}]
 \item $s(x,y)$ is linear in $x$ and conjugate-linear in $y$,
 $s(y,x)=\overline{s(x,y)}$, and $s(x,x)\geq0$ (in other words, $s$ is a
 positive Hermitian sesquilinear form);
% Source PDF physical page 127
 \item $s(y,x)=s(x^{\ast},y^{\ast})$ for $x,y\in\mathfrak n$;
 \item $s(zx,y)=s(x,z^{\ast}y)$ for $x,y\in\mathfrak n$, $z\in A$;
 \item for every $z\in A$, the map $x\mapsto zx$ from $\mathfrak n$ into
 $\mathfrak n$ is continuous for the pre-Hilbert structure defined by $s$;
 \item the elements $xy$, where $x,y\in\mathfrak n$, are dense in
 $\mathfrak n$ for the pre-Hilbert structure defined by $s$.
\end{enumerate}
\end{itshape}

The ideal $\mathfrak n$ is called the \emph{ideal of definition} of $s$ and
is denoted by $\mathfrak n_s$.

It follows readily from the axioms that
$s(xz,y)=s(x,yz^{\ast})$ for $x,y\in\mathfrak n$, $z\in A$, and that, for
every $z\in A$, the map $x\mapsto xz$ from $\mathfrak n$ into $\mathfrak n$
is continuous for the pre-Hilbert structure defined by $s$.

\ResultHeading{6.2.2}{6.2.2.}
Let $N_s$ be the set of $x\in\mathfrak n_s$ such that $s(x,x)=0$. By (i), it
is a vector subspace of $\mathfrak n_s$; by (ii), it is self-adjoint; and by
(iii), it is a left ideal of $A$, hence a two-sided ideal of $A$. Let
$\Lambda_s$ be the canonical map from $\mathfrak n_s$ onto
$\mathfrak n_s/N_s$. The form $s$ passes to the quotient and defines on
$\mathfrak n_s/N_s$ a separated pre-Hilbert-space structure. We denote its
Hilbert-space completion by $H_s$ and its inner product by $(\,\cdot\mid\cdot\,)$.
By (ii), the map $x\mapsto x^{\ast}$ passes to the quotient and defines an
isometric map of $\mathfrak n_s/N_s$ onto itself. Extending it by continuity,
we obtain a conjugate-linear map $J_s\colon H_s\to H_s$ such that $J_s^2=1$
and
\[
 (J_sa\mid J_sb)=(b\mid a)\qquad(a,b\in H_s).
\]
Let $z\in A$. The maps $x\mapsto zx$ and $x\mapsto xz$ on
$\mathfrak n_s$ pass to the quotient $\mathfrak n_s/N_s$ and extend by
continuity to bounded linear operators $\lambda_s(z)$ and $\rho_s(z)$ on
$H_s$. It is easy to see that $\lambda_s$ is a representation of $A$ on
$H_s$; for example, by (iii), for $x,y\in\mathfrak n_s$,
\begin{align*}
 (\lambda_s(z^{\ast})\Lambda_sx\mid\Lambda_sy)
 &= (\Lambda_s(z^{\ast}x)\mid\Lambda_sy)
  =s(z^{\ast}x,y)=s(x,zy)\\
 &= (\Lambda_sx\mid\Lambda_s(zy))
  =(\Lambda_sx\mid\lambda_s(z)\Lambda_sy)\\
 &= (\lambda_s(z)^{\ast}\Lambda_sx\mid\Lambda_sy),
\end{align*}
whence $\lambda_s(z^{\ast})=\lambda_s(z)^{\ast}$. Similarly, $\rho_s$ is a
representation on $H_s$ of the $C^{*}$-algebra $A^{\mathrm o}$ opposite to
$A$. The representations $\lambda_s$ and $\rho_s$ are nondegenerate by (v).
For $x\in\mathfrak n_s$ and $z\in A$, we have
\[
 \rho_s(z^{\ast})\Lambda_sx^{\ast}
 =\Lambda_s(x^{\ast}z^{\ast})
 =\Lambda_s((zx)^{\ast})
 =J_s\lambda_s(z)J_s\Lambda_sx^{\ast},
\]
and hence
\[
 \rho_s(z^{\ast})=J_s\lambda_s(z)J_s.
\]
Since $J_s$ is an isomorphism of $H_s$ onto the conjugate Hilbert space, the
representations $\bar\lambda_s^{\mathrm o}$ (\XRef{2.2.8}{2.2.8}) and $\rho_s$
are equivalent.

\ResultHeading{6.2.3}{6.2.3.}
The axioms (i)--(v) imply that the involutive algebra
$\mathfrak n_s/N_s$, equipped with the inner product, is a Hilbert algebra
(\XRef{A.54}{A 54}). The left and right von Neumann algebras
$\mathfrak U_s,\mathfrak V_s$ associated with $\mathfrak n_s/N_s$ are the
von Neumann algebras generated by $\lambda_s(\mathfrak n_s)$
% Source PDF physical page 128
and $\rho_s(\mathfrak n_s)$; each is the commutant of the other
(\XRef{A.54}{A 54}). The von Neumann algebras generated by $\lambda_s(A)$
and $\rho_s(A)$ are \emph{a priori} larger than the preceding ones; but since
they plainly commute, they are in fact equal to $\mathfrak U_s$ and
$\mathfrak V_s$, respectively.

\ResultHeading{6.2.4}{6.2.4.}
The Hilbert algebra $\mathfrak n_s/N_s$ defines on $\mathfrak U_s^{+}$ a
natural trace, denoted by $t_s$ (\XRef{A.60}{A 60}). This trace is normal.
For every $x\in\mathfrak n_s$, we have
\[
 t_s\bigl(\lambda_s(x)\lambda_s(x)^{\ast}\bigr)
   =(\Lambda_sx\mid\Lambda_sx)
 \qquad\text{(\XRef{A.60}{A 60})},
\]
and hence
\begin{equation}
 t_s(\lambda_s(xx^{\ast}))=s(x,x)<+\infty.
 \tag{1}\label{eq:6.2.4-1}
\end{equation}
If $t_s'$ denotes the linear extension of
$t_s|_{\mathfrak m_{t_s}^{+}}$ to $\mathfrak m_{t_s}$, then, for every
$x,y\in\mathfrak n_s$,
\[
 t_s'\bigl(\lambda_s(x)\lambda_s(y)^{\ast}\bigr)
   =(\Lambda_sx\mid\Lambda_sy)
 \qquad\text{(\XRef{A.60}{A 60})},
\]
and hence
\begin{equation}
 t_s'(\lambda_s(xy^{\ast}))=s(x,y).
 \tag{2}\label{eq:6.2.4-2}
\end{equation}
In the same way one can define a trace on $\mathfrak V_s^{+}$, which can also
be obtained from $t_s$ with the aid of $J_s$, but we shall not need it.

Bibliography: \cite{ref54,ref57}.

\BookSubsection{6.3}{Maximal Bitraces}

\ResultHeading{6.3.1}{6.3.1. Lemma.}
\begin{itshape}
Retain the preceding notation. Let
$q=\mathfrak m_{t_s}$ and $q'=q^{1/2}$ (\XRef{A.9}{A 9}). If, for
$x,y\in\lambda_s^{-1}(q')$, one puts
\[
 \widetilde s(x,y)=t_s'(\lambda_s(xy^{\ast})),
\]
then $\widetilde s$ is a bitrace extending $s$, and the pre-Hilbert space
$\mathfrak n_s$ is dense in the pre-Hilbert space
$\mathfrak n_{\widetilde s}$.
\end{itshape}

If $x\in\mathfrak n_s$, then
$\lambda_s(xx^{\ast})\in\mathfrak m_{t_s}$, and therefore
$\lambda_s(x)\in q'$, whence $x\in\lambda_s^{-1}(q')$. This and formula
\eqref{eq:6.2.4-2} prove that $\widetilde s$ extends $s$. Since $q'$ is the
completed Hilbert algebra corresponding to the Hilbert algebra
$\lambda_s(\mathfrak n_s)$ (\XRef{A.60}{A 60} and \XRef{A.61}{A 61}), the
image under $\lambda_s$ of $\lambda_s^{-1}(q')$ is a Hilbert algebra in which
$\lambda_s(\mathfrak n_s)$ is dense. It follows that $\widetilde s$ is a
bitrace and that the pre-Hilbert space $\mathfrak n_s$ is dense in the
pre-Hilbert space $\mathfrak n_{\widetilde s}$.

\ResultHeading{6.3.2}{6.3.2. Lemma.}
\begin{itshape}
Let $s'$ be a bitrace extending $s$. Then
\[
 \mathfrak n_s\subset\mathfrak n_{s'}
 \subset\mathfrak n_{\widetilde{s'}}
 \subset\mathfrak n_{\widetilde s},
 \qquad
 \widetilde{s'}(x,x)\geq\widetilde s(x,x)
 \quad(x\in\mathfrak n_{\widetilde{s'}}).
\]
\end{itshape}

Plainly $N_s=N_{s'}\cap\mathfrak n_s$, and hence there is an isometric linear
map $\mathfrak n_s/N_s\to\mathfrak n_{s'}/N_{s'}$. This map extends to a
linear isometry $T\colon H_s\to H_{s'}$. Put $K=T(H_s)$. Since
$\mathfrak n_s$ is a two-sided ideal of $A$, $K$ is invariant under the
operators $\lambda_{s'}(z)$ and $\rho_{s'}(z)$ for every $z\in A$; hence the
orthogonal projection of $H_{s'}$
% Source PDF physical page 129
onto $K$ belongs to the center of $\mathfrak U_{s'}$. On identifying $H_s$
with $K$ by means of $T$, $\lambda_{s'}(z)$ induces $\lambda_s(z)$ on $H_s$.
Let $B'\subset H_{s'}$ be the Hilbert algebra of elements bounded relative to
the Hilbert algebra $\mathfrak n_{s'}/N_{s'}$, and let $B$ be the orthogonal
projection of $B'$ onto $H_s$. Then the left von Neumann algebra associated
with $B$ is the von Neumann algebra induced by $\mathfrak U_{s'}$ on $H_s$,
and the natural trace defined by $B$ is the trace induced by $t_{s'}$
(\XRef{A.65}{A 65}). Now $\mathfrak n_s/N_s$ is contained in $B'$ and is
dense in $H_s$, hence is contained in $B$ and is dense in $B$. Thus
$\mathfrak n_s/N_s$ defines the same von Neumann algebras as $B$ and the same
natural trace. In other words, $\mathfrak U_s$ and $t_s$ are respectively the
von Neumann algebra induced by $\mathfrak U_{s'}$ on $H_s$ and the trace
induced by $t_{s'}$ on $\mathfrak U_s^{+}$. Now let
$x\in\mathfrak n_{\widetilde{s'}}$. By what precedes,
\[
 t_s(\lambda_s(xx^{\ast}))
 \leq t_{s'}(\lambda_{s'}(xx^{\ast}))<+\infty;
\]
hence $x\in\mathfrak n_{\widetilde s}$ and
$\widetilde s(x,x)\leq\widetilde{s'}(x,x)$.

\ResultHeading{6.3.3}{6.3.3. Definition.}
\begin{itshape}
A bitrace is said to be \emph{maximal} if there is no bitrace properly
extending it.
\end{itshape}

\ResultHeading{6.3.4}{6.3.4. Proposition.}
\begin{itshape}
Let $s$ be a bitrace.
\begin{enumerate}[label=\textup{(\roman*)}]
 \item The bitrace $s$ is maximal if and only if $s=\widetilde s$.
 \item The bitrace $\widetilde s$ is maximal.
\end{enumerate}
\end{itshape}

If $s'$ is a bitrace extending $\widetilde s$, hence extending $s$, then
$\mathfrak n_{s'}\subset\mathfrak n_{\widetilde s}$ by
\XRef{6.3.2}{6.3.2}; hence $s'=\widetilde s$, so $\widetilde s$ is maximal.
Consequently, if $s=\widetilde s$, then $s$ is maximal. The converse is clear.

\ResultHeading{6.3.5}{6.3.5. Definition.}
\begin{itshape}
Let $s$ be a bitrace. The bitrace $\widetilde s$ is called the
\emph{canonical maximal extension} of $s$.
\end{itshape}

In general, $s$ has maximal extensions other than $\widetilde s$: it is enough
to consider the case in which $A$ is the product of $\mathfrak n_s$ and
another $C^{*}$-algebra.

\ResultHeading{6.3.6}{6.3.6.}
Equip $\mathfrak n_s$ with the norm induced by that of $A$. In general the
map $x\mapsto\Lambda_sx$ from $\mathfrak n_s$ into $H_s$ is not continuous.
But the map
\[
 (x,y)\longmapsto\Lambda_s(xy)
 =\lambda_s(x)\Lambda_sy=\rho_s(y)\Lambda_sx
\]
from $\mathfrak n_s\times\mathfrak n_s$ into $H_s$ is separately continuous,
since the representations $\lambda_s$ and $\rho_s$ of $A$ are continuous.
It follows from this and from axiom (v) for bitraces that, if
$E\subset\mathfrak n_s$ is norm-dense in $\mathfrak n_s$, the vectors
$\Lambda_s(xy)$, $x,y\in E$, are dense in $H_s$. In particular, if $A$ is
separable, then $H_s$ is separable.

Bibliography: \cite{ref54,ref57}.

% Source PDF physical page 130
\BookSubsection{6.4}{Relations Between Traces and Bitraces}

\ResultHeading{6.4.1}{6.4.1. Lemma.}
\begin{itshape}
Let $A$ be a $C^{*}$-algebra and let $f$ be a lower semicontinuous trace on
$A^{+}$. Use the notation of \XRef{6.1.2}{6.1.2}. For
$x,y\in\mathfrak n$, put
\[
 s(x,y)=f'(xy^{\ast})=f'(y^{\ast}x).
\]
Then $s$ is a bitrace of $A$.
\end{itshape}

It is clear that $s$ is sesquilinear. For $x,y\in\mathfrak n$ and $z\in A$,
we have, by \XRef{6.1.2}{6.1.2}(iii),
\begin{align*}
 s(y,x)&=f'(yx^{\ast})=\overline{f'(xy^{\ast})}
       =\overline{s(x,y)},\\
 s(x,x)&=f'(xx^{\ast})\geq0,\\
 s(x^{\ast},y^{\ast})&=f'(x^{\ast}y)=s(y,x),\\
 s(zx,y)&=f'(yx^{\ast}z^{\ast})
         =f'((z^{\ast}y)^{\ast}x)=s(x,z^{\ast}y),\\
 s(zx,zx)&=f'(x^{\ast}z^{\ast}zx)
 \leq\lVert z^{\ast}z\rVert f'(x^{\ast}x)
 =\lVert z^{\ast}z\rVert s(x,x).
\end{align*}
Thus axioms (i)--(iv) for bitraces are satisfied. Now let $(u_i)$ be an
increasing directed approximate identity of the $C^{*}$-algebra
$\overline{\mathfrak n}$. For every $x\in\mathfrak n$,
\begin{align*}
 s(u_ix-x,u_ix-x)
 &=s(u_ix,u_ix)-s(u_ix,x)-s(x,u_ix)+s(x,x)\\
 &\leq\lVert u_i^{\ast}u_i\rVert s(x,x)
       -f'(x^{\ast}u_ix)-f'(x^{\ast}u_ix)+s(x,x)\\
 &\leq2\bigl[f'(x^{\ast}x)-f'(x^{\ast}u_ix)\bigr].
\end{align*}
But the $x^{\ast}u_ix$ form an increasing directed family of elements of
$A^{+}$ tending to $x^{\ast}x$, and hence
$f'(x^{\ast}u_ix)\to f'(x^{\ast}x)$ because $f$ is lower semicontinuous.
Thus $x$ is the limit of $u_ix$ for the pre-Hilbert structure defined by $s$.
By \XRef{1.7.2}{1.7.2}, the $u_i$ may be chosen in $\mathfrak n$.

The bitrace $s$ is called the \emph{bitrace associated with $f$}.

\ResultHeading{6.4.2}{6.4.2.}
Conversely, let $A$ be a $C^{*}$-algebra and let $s$ be a bitrace of $A$.
For every $x\in A^{+}$, put
\[
 f(x)=t_s(\lambda_s(x)).
\]
We have $\lambda_s(\mathfrak n_s^2)\subset\mathfrak m_{t_s}$, and
$\lambda_s(\mathfrak n_s^2)$ is strongly dense in $\mathfrak U_s$. Hence
$f$ is a semifinite lower semicontinuous trace on $A^{+}$
(\XRef{6.1.5}{6.1.5}), which will be called the trace associated with $s$.

\ResultHeading{6.4.3}{6.4.3. Lemma.}
\begin{itshape}
Let $A$ be a $C^{*}$-algebra, $s$ a bitrace of $A$, and $f$ its associated
trace. The bitrace $s'$ associated with $f$ is the canonical maximal extension
of $s$.
\end{itshape}

Let $q=\mathfrak m_{t_s}$ and $q'=q^{1/2}$. For every $x\in A$, we have
\[
 \begin{aligned}
 x\in\mathfrak n_{s'}
 &\Longleftrightarrow f(xx^{\ast})<+\infty
 \Longleftrightarrow t_s(\lambda_s(xx^{\ast}))<+\infty\\
 &\Longleftrightarrow \lambda_s(x)\in q'
 \Longleftrightarrow x\in\lambda_s^{-1}(q').
 \end{aligned}
\]
Therefore $\mathfrak n_{s'}=\mathfrak n_{\widetilde s}$, and for
$x\in\mathfrak n_{s'}$,
\[
 s'(x,x)=f(xx^{\ast})=t_s(\lambda_s(xx^{\ast}))=\widetilde s(x,x).
\]

\ResultHeading{6.4.4}{6.4.4. Lemma.}
\begin{itshape}
Let $A$ be a $C^{*}$-algebra, let $f$ be a semifinite lower semicontinuous
trace on $A^{+}$, and let $s$ be the bitrace associated with $f$. Then $s$ is
maximal, and the trace associated with $s$ is $f$.
\end{itshape}

% Source PDF physical page 131
Let $f'$ be the trace associated with $s$, and let $s'$ be the bitrace
associated with $f'$. By \XRef{6.4.3}{6.4.3}, $s'$ is the canonical maximal
extension of $s$. For $y\in\mathfrak n_s$, we have
\[
 f'(yy^{\ast})=s'(y,y)=s(y,y)=f(yy^{\ast}).
\]
Consequently, if $x\in A^{+}$ and $f(x)<+\infty$, then $f'(x)=f(x)$. Let
$x\in A^{+}$ satisfy $f(x)=+\infty$. For every finite real number $\alpha$,
there exists $x_1\in A^{+}$ such that
$x_1\leq x$ and $\alpha\leq f(x_1)<+\infty$, since $f$ is semifinite. Hence
\[
 f'(x)\geq f'(x_1)=f(x_1)\geq\alpha,
 \qquad\text{so that}\qquad f'(x)=+\infty.
\]
This proves that $f'=f$. Therefore $s'=s$, and $s$ is maximal.

\ResultHeading{6.4.5}{6.4.5. Proposition.}
\begin{itshape}
Let $A$ be a $C^{*}$-algebra. Let $E$ be the set of semifinite lower
semicontinuous traces on $A^{+}$, and let $E'$ be the set of maximal bitraces
of $A$. The maps
\[
 \text{trace}\longmapsto\text{associated bitrace},\qquad
 \text{bitrace}\longmapsto\text{associated trace}
\]
define mutually inverse bijections $E\to E'$ and $E'\to E$.
\end{itshape}

This follows from \XRef{6.4.1}{6.4.1}, \XRef{6.4.3}{6.4.3}, and
\XRef{6.4.4}{6.4.4}.

Let $f$ be a semifinite lower semicontinuous trace on $A^{+}$, and let $s$ be
the associated maximal bitrace. We shall write
\[
 N_f=N_s,\qquad \mathfrak n_f=\mathfrak n_s,\qquad \ldots,
 \qquad \mathfrak U_f=\mathfrak U_s,\qquad t_f=t_s;
\]
and, conversely, $\mathfrak m_s=\mathfrak m_f$. Recall that
$\overline{\mathfrak m_f}=\overline{\mathfrak n_f}$
[\XRef{6.1.2}{6.1.2}(i)].

Bibliography: \cite{ref16}.

\BookSubsection{6.5}{Sum of Two Traces}

\ResultHeading{6.5.1}{6.5.1.}
Let $A$ be a $C^{*}$-algebra, and let $f,f'$ be two traces on $A^{+}$. It is
clear that $f+f'$ is a trace. If $f$ and $f'$ are lower semicontinuous, then
$f+f'$ is lower semicontinuous. Suppose that $f$ and $f'$ are semifinite,
and let us prove that $f+f'$ is semifinite. Let $x\in A^{+}$ satisfy
$(f+f')(x)=+\infty$, and let $\alpha$ be a finite real number. For example,
suppose $f(x)=+\infty$. There is $y\in A^{+}$ such that
$y\leq x$ and $\alpha\leq f(y)<+\infty$. If $f'(y)<+\infty$, then
$\alpha\leq(f+f')(y)<+\infty$. If $f'(y)=+\infty$, there exists
$z\in A^{+}$ such that $z\leq y$ and
$\alpha\leq f'(z)<+\infty$; hence $z\leq x$ and
$\alpha\leq(f+f')(z)<+\infty$. In both cases we see that $f+f'$ is
semifinite.

\ResultHeading{6.5.2}{6.5.2.}
The set of traces on $A^{+}$ is plainly ordered by the relation $f\leq f'$,
that is, by the relation
\[
 f(x)\leq f'(x)\qquad\text{for every }x\in A^{+}.
\]
Suppose $f,f'$ are semifinite and lower semicontinuous, and let $s,s'$ be the
associated bitraces. Then $f\leq f'$ is equivalent to
\[
 \mathfrak n_{s'}\subset\mathfrak n_s,
 \qquad s(x,x)\leq s'(x,x)\quad\text{for every }x\in\mathfrak n_{s'}.
\]

\ResultHeading{6.5.3}{6.5.3. Lemma.}
\begin{itshape}
Let $A$ be a $C^{*}$-algebra and let $s,s'$ be two maximal bitraces of $A$.
Suppose there exists an involutive subalgebra $B$ of $A$, dense in $A$ and
% Source PDF physical page 132
contained in $\mathfrak n_s$ and $\mathfrak n_{s'}$, such that $s$ and $s'$
agree on $B\times B$. Then $s=s'$.
\end{itshape}

By \XRef{6.3.6}{6.3.6}, $B$ is dense in $\mathfrak n_s$ (respectively,
$\mathfrak n_{s'}$) for the pre-Hilbert structure of $\mathfrak n_s$
(respectively, $\mathfrak n_{s'}$). Since $s$ and $s'$ agree on
$B\times B$, there is an isomorphism $\Theta$ from $H_s$ onto $H_{s'}$ which,
for every $x\in B$, transforms $\Lambda_sx$ into $\Lambda_{s'}x$. For
$x,y\in B$, we have
\[
 \Theta\lambda_s(x)\Lambda_sy
 =\Theta\Lambda_s(xy)=\Lambda_{s'}(xy)
 =\lambda_{s'}(x)\Lambda_{s'}y
 =\lambda_{s'}(x)\Theta\Lambda_sy;
\]
hence $\Theta$ transforms $\lambda_s(x)$ into $\lambda_{s'}(x)$ for every
$x\in B$, and consequently for every $x\in A$. Moreover, since $\Theta$
defines an isomorphism from the Hilbert algebra $\lambda_s(B)$, dense in
$\mathfrak n_s/N_s$, onto the Hilbert algebra $\lambda_{s'}(B)$, dense in
$\mathfrak n_{s'}/N_{s'}$, $\Theta$ transforms $\mathfrak U_s$ into
$\mathfrak U_{s'}$ and $t_s$ into $t_{s'}$. Let $f,f'$ be the traces
associated with $s,s'$. For every $x\in A^{+}$,
\[
 f(x)=t_s(\lambda_s(x))=t_{s'}(\lambda_{s'}(x))=f'(x).
\]
Since $s$ and $s'$ are maximal, $s=s'$ by \XRef{6.4.3}{6.4.3}.

\ResultHeading{6.5.4}{6.5.4. Lemma.}
\begin{itshape}
Let $A$ be a $C^{*}$-algebra, let $f,f'$ be two semifinite lower
semicontinuous traces on $A^{+}$, and let $s,s'$ be the associated bitraces.
Suppose that $f'\leq f$ on $(\overline{\mathfrak n_s})^{+}$.
\begin{enumerate}[label=\textup{(\roman*)}]
 \item There exists a unique $T\in\mathcal L(H_s)$ such that
 \[
   s'(x,y)=(T\Lambda_sx\mid\Lambda_sy)
   \qquad(x,y\in\mathfrak n_s).
 \]
 \item One has $0\leq T\leq1$ and
 $T\in\mathfrak U_s\cap\mathfrak V_s$.
 \item If
 \[
   f_1(x)=t_s(T\lambda_s(x))\qquad(x\in A^{+}),
 \]
 then $f_1$ is a semifinite lower semicontinuous trace on $A^{+}$, is
 majorized by $f'$, and agrees with $f'$ on
 $(\overline{\mathfrak n_s})^{+}$.
\end{enumerate}
\end{itshape}

For $x\in\mathfrak n_s$, we have $s'(x,x)\leq s(x,x)$. Hence, for
$x,y\in\mathfrak n_s$, $s'(x,y)$ depends only on $\Lambda_sx$ and
$\Lambda_sy$, and is continuous in these two variables for the Hilbert-space
structure of $H_s$. This proves the existence of a $T\in\mathcal L(H_s)$
satisfying (i). Its uniqueness follows because
$\Lambda_s(\mathfrak n_s)$ is dense in $H_s$. We have
$0\leq s'(x,x)\leq s(x,x)$ for every $x\in\mathfrak n_s$, and therefore
\[
 0\leq(T\Lambda_sx\mid\Lambda_sx)
 \leq(\Lambda_sx\mid\Lambda_sx),
 \qquad\text{so that}\qquad 0\leq T\leq1.
\]
If $z\in A$ and $x,y\in\mathfrak n_s$, then
\begin{align*}
 (T\lambda_s(z)\Lambda_sx\mid\Lambda_sy)
 &= (T\Lambda_s(zx)\mid\Lambda_sy)=s'(zx,y)=s'(x,z^{\ast}y)\\
 &= (T\Lambda_sx\mid\Lambda_s(z^{\ast}y))
  =(T\Lambda_sx\mid\lambda_s(z^{\ast})\Lambda_sy)\\
 &= (\lambda_s(z)T\Lambda_sx\mid\Lambda_sy),
\end{align*}
and hence $T\lambda_s(z)=\lambda_s(z)T$. In the same way one sees that
$T\rho_s(z)=\rho_s(z)T$, whence
$T\in\mathfrak U_s\cap\mathfrak V_s$.

The function $S\mapsto t_s(TS)$ on $\mathfrak U_s^{+}$ is a normal trace
$t_1$. We have $\lambda_s(\mathfrak m_s)\subset\mathfrak m_{t_1}$; indeed,
if $x\in\mathfrak m_s^{+}$, then $t_s(\lambda_s(x))<+\infty$, and therefore
\[
 t_1(\lambda_s(x))\leq\lVert T\rVert t_s(\lambda_s(x))<+\infty.
\]
% Source PDF physical page 133
Moreover, $\lambda_s(\mathfrak m_s)$ is strongly dense in
$\mathfrak U_s$; consequently $f_1$ is a semifinite lower semicontinuous trace
on $A^{+}$ (\XRef{6.1.5}{6.1.5}). For $x\in\mathfrak n_s$,
\[
 f_1(xx^{\ast})=t_s(T\lambda_s(x)\lambda_s(x^{\ast})).
\]
Since $T\in\mathfrak U_s$, the vector $T\Lambda_sx$ is a bounded element
relative to the Hilbert algebra $\lambda_s(\mathfrak n_s)$
(\XRef{A.59}{A 59}), and
\[
 t_s(T\lambda_s(x)\lambda_s(x^{\ast}))
 =(T\Lambda_sx\mid\Lambda_sx)=s'(x,x)=f'(xx^{\ast}).
\]
Thus $f_1|_{(\overline{\mathfrak n_s})^{+}}$ and
$f'|_{(\overline{\mathfrak n_s})^{+}}$ are semifinite lower semicontinuous
traces on $(\overline{\mathfrak n_s})^{+}$ (\XRef{6.1.4}{6.1.4}), finite and
equal on $(\mathfrak n_s^2)^{+}$ by what precedes, and hence equal by
\XRef{6.5.3}{6.5.3}. On the other hand, let $(u_i)$ be an increasing directed
approximate identity of $\overline{\mathfrak n_s}$. By
\XRef{6.1.5}{6.1.5}, for every $x\in A^{+}$,
\[
 f_1(x)=\lim_i f_1(x^{1/2}u_ix^{1/2})
       =\lim_i f'(x^{1/2}u_ix^{1/2})\leq f'(x),
\]
because $x^{1/2}u_ix^{1/2}\in(\overline{\mathfrak n_s})^{+}$ and
$x^{1/2}u_ix^{1/2}\leq x$.

\ResultHeading{6.5.5}{6.5.5. Proposition.}
\begin{itshape}
Let $A$ be a $C^{*}$-algebra, and let $f,f'$ be two semifinite lower
semicontinuous traces on $A^{+}$. In order that $f'\leq f$, it is necessary
and sufficient that there exist a semifinite lower semicontinuous trace $f''$
on $A^{+}$ such that $f=f'+f''$.
\end{itshape}

The condition is plainly sufficient. Suppose $f'\leq f$, and use Lemma
\XRef{6.5.4}{6.5.4} and its notation. Put
\[
 f''(x)=t_s((1-T)\lambda_s(x))\qquad(x\in A^{+}).
\]
By \XRef{6.1.5}{6.1.5}, $f''$ is a semifinite lower semicontinuous trace on
$A^{+}$. For every $x\in A^{+}$,
\[
 f_1(x)+f''(x)
 =t_s(T\lambda_s(x))+t_s((1-T)\lambda_s(x))
 =t_s(\lambda_s(x))=f(x).
\]
By \XRef{6.5.4}{6.5.4}, $f'+f''$ agrees with $f$ on
$(\overline{\mathfrak n_s})^{+}$. If
$x\in A^{+}$ and $x\notin(\overline{\mathfrak n_s})^{+}$, then
\[
 f'(x)+f''(x)\geq f_1(x)+f''(x)=f(x)=+\infty,
\]
so again $f'(x)+f''(x)=f(x)$.

\ResultHeading{6.5.6}{6.5.6. Proposition.}
\begin{itshape}
Let $A$ be a $C^{*}$-algebra, and let $f,f'$ be two lower semicontinuous
traces on $A^{+}$ such that $f'\leq f$. Suppose $\mathfrak m_f$ is dense in
$A$, so that $f$ and $f'$ are semifinite (\XRef{6.1.3}{6.1.3}). There exists
a unique lower semicontinuous trace $f''$ on $A^{+}$ such that
$f=f'+f''$.
\end{itshape}

The existence of $f''$ follows from \XRef{6.5.5}{6.5.5}. Suppose that we have
two decompositions
\[
 f=f'+f'',\qquad f=f'+f_1''.
\]
For
$x\in\mathfrak m_f^{+}$, the numbers
$f(x),f'(x),f''(x),f_1''(x)$ are finite, and hence
$f''(x)=f_1''(x)$. Therefore $f''=f_1''$ by \XRef{6.5.3}{6.5.3}.

\ResultHeading{6.5.7}{6.5.7.}
The uniqueness of $f''$ fails when $\overline{\mathfrak m_f}\ne A$, even if
$f$ is semifinite (\XRef{6.9.3}{6.9.3}).

Bibliography: \cite{ref16,ref57}.

% Source PDF physical page 134
\BookSubsection{6.6}{Traces and Representations}

\ResultHeading{6.6.1}{6.6.1. Definition.}
\begin{itshape}
Let $A$ be a $C^{*}$-algebra. A \emph{traced representation} of $A$ is a
pair $(\pi,t)$ having the following properties:
\begin{enumerate}[label=\textup{(\roman*)}]
 \item $\pi$ is a nondegenerate representation of $A$ on a Hilbert space;
 \item $t$ is a faithful normal trace on $\mathfrak U^{+}$, where
 $\mathfrak U$ denotes the von Neumann algebra generated by $\pi(A)$;
 \item $\pi(A)\cap\mathfrak n_t$ is weakly dense in the von Neumann algebra
 $\mathfrak U$.
\end{enumerate}
\end{itshape}

These conditions imply that $t$ is semifinite, and hence that
$\mathfrak U$ is semifinite. Thus $\pi$ is a Hilbert direct sum of a
representation of type I and a representation of type II.

\ResultHeading{6.6.2}{6.6.2.}
Let $(\pi,t)$ and $(\pi',t')$ be two traced representations of $A$, and let
$\mathfrak U,\mathfrak U'$ be the von Neumann algebras generated by $\pi(A)$
and $\pi'(A)$. We say that $(\pi,t)$ and $(\pi',t')$ are quasi-equivalent
(respectively, equivalent) if there exists an isomorphism (respectively, a
spatial isomorphism) from $\mathfrak U$ onto $\mathfrak U'$ which transforms
$\pi$ into $\pi'$ and $t$ into $t'$. Then $\pi$ and $\pi'$ are
quasi-equivalent (respectively, equivalent).

\ResultHeading{6.6.3}{6.6.3.}
Let $A$ be a $C^{*}$-algebra and $(\pi,t)$ a traced representation of $A$.
Then
\[
 f=(t\circ\pi)|_{A^{+}}
\]
is a semifinite lower semicontinuous trace on $A^{+}$
(\XRef{6.1.5}{6.1.5}). We shall say that $f$, and the corresponding maximal
bitrace, are associated with $(\pi,t)$. If $(\pi,t)$ is replaced by a
quasi-equivalent traced representation, $f$ does not change.

\ResultHeading{6.6.4}{6.6.4.}
Conversely, let $f$ be a semifinite lower semicontinuous trace on $A^{+}$,
and let $s$ be the corresponding bitrace. Then $(\lambda_s,t_s)$ is a traced
representation of $A$, called the traced representation associated with $f$,
or with $s$. If $\lambda_s$ is of type I, of type II, and so forth, then $f$
and $s$ are said to be of type I, of type II, and so forth.

\ResultHeading{6.6.5}{6.6.5. Proposition.}
\begin{itshape}
\begin{enumerate}[label=\textup{(\roman*)}]
 \item Let $f$ be a semifinite lower semicontinuous trace on $A^{+}$, and let
 $(\pi,t)$ be the traced representation associated with $f$. Then the trace
 associated with $(\pi,t)$ is $f$.
 \item Let $(\pi,t)$ be a traced representation of $A$, and let $f$ be the
 trace associated with $(\pi,t)$. Then the traced representation associated
 with $f$ is quasi-equivalent to $(\pi,t)$.
\end{enumerate}
\end{itshape}

Let $f,\pi,t$ be as in (i), let $s$ be the bitrace associated with $f$, and
let $f'$ be the trace associated with $(\pi,t)$. Then
$f'=(t_s\circ\lambda_s)|_{A^{+}}$, and therefore $f'=f$ by
\XRef{6.4.4}{6.4.4}.

Let $(\pi,t)$ be a traced representation of $A$, and let $\mathfrak C$ be the
von Neumann algebra generated by $\pi(A)$. Recall the following facts, which
have nothing to do with $\pi$. Equipped with the inner product
\[
 (S,T)\longmapsto t'(ST^{\ast}),
\]
$\mathfrak n_t$ is a completed Hilbert algebra, where $t'$ denotes the linear
extension to $\mathfrak m_t$ of $t|_{\mathfrak m_t^{+}}$;
$\lambda_t$ is an isomorphism from the von Neumann algebra $\mathfrak C$ onto
the von Neumann algebra $\mathfrak U_t$ on $H_t$; and $\lambda_t$ transforms
$t$ into the natural trace $\theta$ defined
% Source PDF physical page 135
on $\mathfrak U_t^{+}$ by the Hilbert algebra $\mathfrak n_t$
(\XRef{A.61}{A 61}). Now let $s$ be the bitrace of $A$ associated with
$(\pi,t)$. We have $\mathfrak n_s=\pi^{-1}(\mathfrak n_t)$ and, for
$x,y\in\mathfrak n_s$,
\[
 s(x,y)=t'(\pi(xy^{\ast})).
\]
Thus $\pi|_{\mathfrak n_s}$ passes to the quotient and defines an isomorphism
$\Phi$ from the Hilbert algebra $\mathfrak n_s/N_s$ onto the Hilbert algebra
$\pi(\mathfrak n_s)\subset\mathfrak n_t$. Since $\pi(\mathfrak n_s)$
generates the von Neumann algebra $\mathfrak C$ by the definition of traced
representations, $\lambda_t(\pi(\mathfrak n_s))$ generates the von Neumann
algebra $\mathfrak U_t$; hence, by \XRef{A.55}{A 55},
$\pi(\mathfrak n_s)$ is dense in the completed Hilbert algebra
$\mathfrak n_t$. Therefore $\Phi$ extends to an isomorphism, still denoted by
$\Phi$, from $H_s$ onto $H_t$. It is easy to see that, for every $z\in A$,
$\Phi$ transforms $\lambda_s(z)$ into $\lambda_t(\pi(z))$. Moreover,
$\Phi$ transforms $t_s$ into $\theta$. Now $(\lambda_s,t_s)$ is the traced
representation associated with $s$, or with $f$. We see that it is equivalent
to $(\lambda_t\circ\pi,\theta)$. Since $\lambda_t$ is an isomorphism from
$\mathfrak C$ onto $\mathfrak U_t$ which transforms $t$ into $\theta$,
$(\lambda_s,t_s)$ is quasi-equivalent to $(\pi,t)$.

\ResultHeading{6.6.6}{6.6.6. Corollary.}
\begin{itshape}
There is a canonical bijection from the set of semifinite lower
semicontinuous traces on $A^{+}$ onto the set of quasi-equivalence classes of
traced representations of $A$.
\end{itshape}

\ResultHeading{6.6.7}{6.6.7.}
In fact, it is not the traced representations of $A$ that are of interest,
but the representations themselves. We therefore make the following
definition.

\ResultLabel{Definition.}
\begin{itshape}
Let $A$ be a $C^{*}$-algebra, let $\pi$ be a representation of $A$, and let
$\mathfrak U$ be the von Neumann algebra generated by $\pi(A)$. The
representation $\pi$ is said to be \emph{traceable} if there exists a trace
$t$ on $\mathfrak U^{+}$ such that $(\pi,t)$ is a traced representation.
\end{itshape}

A representation quasi-equivalent to a traceable representation is
traceable.

We shall study a case in which there is little difference between traced and
traceable representations.

Bibliography: \cite{ref16,ref54,ref57}.

\BookSubsection{6.7}{Characters and Traceable Factor Representations}

\ResultHeading{6.7.1}{6.7.1. Definition.}
\begin{itshape}
Let $A$ be a $C^{*}$-algebra. A \emph{character} of $A$ is a nonzero
semifinite lower semicontinuous trace $f$ on $A^{+}$ satisfying the following
condition: every semifinite lower semicontinuous trace on $A^{+}$ that is
majorized by $f$ is proportional to $f$ on $(\mathfrak m_f)^{+}$.
\end{itshape}

\ResultHeading{6.7.2}{6.7.2.}
Recall that, if $\mathfrak F$ is a semifinite factor, there exist faithful
normal semifinite traces $t$ on $\mathfrak F^{+}$, and all these traces are
proportional. Thus the ideal $\mathfrak n_t$ of $\mathfrak F$ is uniquely
determined; its elements are called the Hilbert--Schmidt operators relative
to $\mathfrak F$ (\XRef{A.32}{A 32}).

% Source PDF physical page 136
\ResultLabel{Proposition.}
\begin{itshape}
Let $A$ be a $C^{*}$-algebra, let $\pi$ be a nonzero factor representation of
$A$, and let $\mathfrak F$ be the factor generated by $\pi(A)$. The following
conditions are equivalent:
\begin{enumerate}[label=\textup{(\roman*)}]
 \item $\pi$ is traceable;
 \item $\mathfrak F$ is semifinite and $\pi(A)$ contains a nonzero
 Hilbert--Schmidt operator relative to $\mathfrak F$.
\end{enumerate}
\end{itshape}

(i)$\Rightarrow$(ii) is clear. Suppose (ii) holds, and let $t$ be a faithful
normal semifinite trace on $\mathfrak F^{+}$. By hypothesis,
$\mathfrak n=\pi(A)\cap\mathfrak n_t$ is a nonzero two-sided ideal of
$\pi(A)$. Hence the weak closure $B$ of $\mathfrak n$ is a nonzero weakly
closed two-sided ideal of the weak closure of $\pi(A)$, that is, of
$\mathfrak F$. Since $\mathfrak F$ is a factor, $B=\mathfrak F$
(\XRef{A.9}{A 9}).

\ResultHeading{6.7.3}{6.7.3. Theorem.}
\begin{itshape}
Let $A$ be a $C^{*}$-algebra, let $f$ be a semifinite lower semicontinuous
trace on $A^{+}$, and let $(\pi,t)$ be the traced representation associated
with $f$. In order that $\pi$ be nonzero and factorial, it is necessary and
sufficient that $f$ be a character.
\end{itshape}

Suppose that $\pi$ is nonzero and factorial. Then $f$ is nonzero. Let $f'$ be
a semifinite lower semicontinuous trace on $A^{+}$, majorized by $f$ on
$(\mathfrak m_f)^{+}$. Apply Lemma \XRef{6.5.4}{6.5.4}. With the notation of
that lemma, $T$ is scalar because $\pi$ is factorial. By
\XRef{6.5.4}{6.5.4}(iii), $f'$ is proportional to $f$ on
$(\mathfrak m_f)^{+}$.

Conversely, suppose that $f$ is a character. Then $\pi$ is nonzero. Let $T$
be an element of $\mathfrak U_f\cap\mathfrak V_f$ such that $0\leq T\leq1$.
If
\[
 f'(x)=t_f(T\lambda_f(x))\qquad(x\in A^{+}),
\]
then, by \XRef{6.1.5}{6.1.5}, $f'$ is a semifinite lower semicontinuous trace
on $A^{+}$, and it is clear that $f'$ is majorized by $f$. Thus there exists
$\lambda\in[0,1]$ such that $f'=\lambda f$ on $(\mathfrak m_f)^{+}$. For
every $x\in\mathfrak n_f$,
\begin{align*}
 (T\Lambda_fx\mid\Lambda_fx)
 &=t_f(T\lambda_f(x)\lambda_f(x^{\ast}))=f'(xx^{\ast})\\
 &=\lambda f(xx^{\ast})
  =\lambda(\Lambda_fx\mid\Lambda_fx),
\end{align*}
whence $T=\lambda$. This proves that $\mathfrak U_f$ is a factor.

\ResultHeading{6.7.4}{6.7.4. Corollary.}
\begin{itshape}
There is a canonical bijective correspondence between:
\begin{enumerate}[label=\textit{\alph*.},leftmargin=2em]
 \item the set of quasi-equivalence classes of nonzero traceable factor
 representations of $A$;
 \item the set of characters of $A$, defined up to a factor $>0$.
\end{enumerate}
\end{itshape}

For every character $f$ of $A$, $\lambda_f$ is a nonzero traceable factor
representation (\XRef{6.7.3}{6.7.3}). Every nonzero traceable factor
representation of $A$ is obtained in this way up to quasi-equivalence
(\XRef{6.6.5}{6.6.5} and \XRef{6.7.3}{6.7.3}). Finally, let $f,f'$ be two
characters of $A$ such that $\lambda_f$ and $\lambda_{f'}$ are
quasi-equivalent. There exists an isomorphism $\Phi$ from $\mathfrak U_f$ onto
$\mathfrak U_{f'}$ which transforms
% Source PDF physical page 137
$\lambda_f$ into $\lambda_{f'}$. Since $\mathfrak U_f$ and
$\mathfrak U_{f'}$ are factors, $\Phi$ transforms $t_f$ into
$\alpha t_{f'}$ for some $\alpha\in(0,+\infty)$. Hence, for every
$x\in A^{+}$,
\[
 f(x)=t_f(\lambda_f(x))
     =\alpha t_{f'}(\lambda_{f'}(x))=\alpha f'(x).
\]

In view of this result, one may speak, with a slight abuse of language, of the
character of a nonzero traceable factor representation. This character is
defined only up to a factor $>0$.

\ResultHeading{6.7.5}{6.7.5.}
Let $A$ be a $C^{*}$-algebra and let $\pi$ be a nonzero irreducible (hence
factorial) representation of $A$. In order that $\pi$ be traceable, it is
necessary and sufficient that
\[
\pi(A)\supset\mathcal{LC}(H_\pi).
\]
Indeed, this condition is plainly sufficient, and it is necessary by
\XRef{4.1.10}{4.1.10}. [The character of $\pi$ may be normalized by using the
usual trace on $\mathcal L(H_\pi)^{+}$.] In particular, every irreducible
representation of a postliminal $C^{*}$-algebra is traceable
(\XRef{4.3.7}{4.3.7}), and hence every factor representation of a postliminal
$C^{*}$-algebra is traceable (\XRef{5.5.3}{5.5.3}).

\ResultHeading{6.7.6}{6.7.6.}
Let $A$ be a postliminal $C^{*}$-algebra, and let $\pi$ be a nonzero
irreducible representation of $A$. For every $x\in A^{+}$, put
\[
 f_\pi(x)=\operatorname{Tr}\pi(x).
\]
Then $f_\pi$ is a character of $A$ (\XRef{6.7.3}{6.7.3} and
\XRef{6.7.5}{6.7.5}). If $\pi'$ is another irreducible representation of
$A$ and $f_\pi=f_{\pi'}$, then $\pi$ and $\pi'$ are quasi-equivalent
(\XRef{6.6.6}{6.6.6}), and hence equivalent (\XRef{5.3.3}{5.3.3}). Finally,
if $g$ is a character of $A$, then $\lambda_g$ is nonzero and factorial
(\XRef{6.7.3}{6.7.3}); it is therefore quasi-equivalent to a nonzero
irreducible representation $\pi$ (\XRef{5.5.3}{5.5.3}), and consequently
$g$ is proportional to $f_\pi$ (\XRef{6.7.4}{6.7.4}).

\ResultHeading{6.7.7}{6.7.7.}
Retain the notation of \XRef{6.7.6}{6.7.6}, and let $f=f_\pi$. Since
$\pi(A)\supset\mathcal{LC}(H_\pi)$, $\lambda_f$ passes to the quotient and
defines an isomorphism from $\mathfrak n_f/N_f$ onto the Hilbert algebra $B$
of Hilbert--Schmidt operators on $H_\pi$. Identify $\mathfrak n_f/N_f$ with
$B$ by this isomorphism. Then, for every $x\in A$, $\lambda_f(x)$ is the
operator of left multiplication by $\pi(x)$ on $B$, and
$\bar\rho_f^{\mathrm o}(x)$ is the operator of right multiplication by
$\pi(x)^{\ast}$ on the conjugate space of $B$. If $B$ is now canonically
identified with $H_\pi\otimes\overline{H_\pi}$, then $\lambda_f(x)$ becomes
$\pi(x)\otimes1$ and $\bar\rho_f^{\mathrm o}(x)$ becomes
$1\otimes\pi(x)$.
Thus $\mathfrak U_f$ is identified with
$\mathcal L(H_\pi)\otimes\mathbb C$, $\mathfrak V_f$ with
$\mathbb C\otimes\mathcal L(\overline{H_\pi})$, and $J_f$ with the canonical
involution of $H_\pi\otimes\overline{H_\pi}$. Finally,
\[
 t_f(T\otimes1)=\operatorname{Tr}T
 \qquad\text{for every }T\in\mathcal L(H_\pi)^{+}.
\]

\ResultHeading{6.7.8}{6.7.8.}
Let $A$ be a commutative $C^{*}$-algebra. By \XRef{6.7.6}{6.7.6}, the
characters of $A$ in the sense of \XRef{6.7.1}{6.7.1} are the functions
$x\mapsto\alpha g(x)$ on $A^{+}$, where $\alpha>0$ and $g$ is a character of
$A$ in the classical sense. There is therefore a slight inconsistency in the
terminology.

% Source PDF physical page 138
\ResultHeading{6.7.9}{6.7.9.}
The reader will have noticed an analogy between the developments of this
section and those of Section 2; the correspondence is as follows:
\[
\begin{array}{c@{\qquad}c}
 \text{Section 2} & \text{Section 6}\\[2pt]
 \text{Positive functionals} &
   \begin{array}{c}\text{Semifinite lower semicontinuous}\\[-2pt]
                     \text{traces}\end{array}\\[4pt]
 \text{Pure positive functionals} & \text{Characters}\\
 \text{Representations} & \text{Traceable representations}\\
 \text{Irreducible representations} &
   \text{Traceable factor representations}
\end{array}
\]
Unfortunately, the analogy breaks down at an essential point. Whereas in
Section 2 one can prove the existence of pure positive functionals by the
Krein--Milman theorem, characters do not always exist (see
\XRef{6.9.2}{6.9.2}, \XRef{6.9.10}{6.9.10}, and \XRef{9.5.7}{9.5.7}). The
question has been settled only for postliminal $C^{*}$-algebras, since in that
case irreducible representations, whose existence has been proved, are
automatically traceable.

On the other hand, the theories of Sections 2 and 6 have a common special
case, which we now study.

Bibliography: \cite{ref54,ref57}.

\BookSubsection{6.8}{Finite Traces}

\ResultHeading{6.8.1}{6.8.1.}
Let $A$ be a $C^{*}$-algebra and let $g$ be a linear functional on $A$. We say
that $g$ is \emph{central} if $g(xy)=g(yx)$ for all $x,y\in A$. Suppose $g$
is positive and central. Then the restriction $f$ of $g$ to $A^{+}$ is
plainly a finite trace. Conversely, let $f$ be a finite trace on $A^{+}$.
Apply Proposition \XRef{6.1.2}{6.1.2} and its notation. We have
$\mathfrak m=A$ (\XRef{1.5.8}{1.5.8}), and hence $f'$ is a linear functional
on $A$ extending $f$, and is therefore positive and central
[\XRef{6.1.2}{6.1.2}(iii)]. Thus we have established a bijective
correspondence between central positive functionals $g$ on $A$ and finite
traces $f$ on $A^{+}$; we shall sometimes identify $f$ and $g$. By
\XRef{2.1.8}{2.1.8}, finite traces are automatically continuous.

\ResultHeading{6.8.2}{6.8.2.}
Let $f$ be a lower semicontinuous trace on $A^{+}$, and let $s$ be the
associated bitrace. We have
$\mathfrak m_f=\mathfrak n_s^2\subset\mathfrak n_s$. By
\XRef{1.5.8}{1.5.8}, the conditions $\mathfrak m_f=A$ and
$\mathfrak n_s=A$ are equivalent. The condition $\mathfrak m_f=A$ plainly
means that $f$ is finite.

\ResultHeading{6.8.3}{6.8.3. Proposition.}
\begin{itshape}
Let $A$ be a $C^{*}$-algebra, let $f$ be a central positive functional on
$A$, and put $g=f|_{A^{+}}$.
\begin{enumerate}[label=\textup{(\roman*)}]
 \item One has $\pi_f=\lambda_g$.
 \item The trace $t_g$ on $\mathfrak U_g^{+}$ is finite. If $t_g'$ denotes
 its linear extension to $\mathfrak U_g$, then
 \[
   t_g'(\pi_f(z))=(\pi_f(z)\xi_f\mid\xi_f)\qquad(z\in A).
 \]
 \item $\mathfrak U_g$ and $\mathfrak V_g$ are finite von Neumann algebras.
 \item For every $z\in A$,
 \[
   \lambda_g(z)\xi_f=\rho_g(z)\xi_f=\Lambda_gz,
   \qquad J_g\xi_f=\xi_f.
 \]
 \item The vector $\xi_f$ is separating and cyclic for $\mathfrak U_g$ and
 $\mathfrak V_g$.
\end{enumerate}
\end{itshape}

% Source PDF physical page 139
Let $s$ be the bitrace associated with $g$. The space $H(\pi_f)$ is the
Hilbert space obtained by separation and completion of $A$ equipped with the
inner product
\[
 f(y^{\ast}x)=f(xy^{\ast})=s(x,y).
\]
Thus $H(\pi_f)=H_s=H_f$. The representations $\pi_f$ and $\lambda_g$ are
both obtained from the left regular representation of $A$, and hence
$\pi_f=\lambda_g$. For $x,y\in A$, writing $t_g'$ for the linear extension
of $t_g$ to $\mathfrak m_{t_g}$, we have
\begin{align}
 t_g'(\lambda_g(x)\lambda_g(y)^{\ast})
 &=s(x,y)=f(xy^{\ast})\notag\\
 &=(\pi_f(xy^{\ast})\xi_f\mid\xi_f)
  =(\pi_f(x)\pi_f(y)^{\ast}\xi_f\mid\xi_f).
 \tag{1}\label{eq:6.8.3-1}
\end{align}
Passing first to weak limits, we deduce that
\[
 (ST\xi_f\mid\xi_f)=(TS\xi_f\mid\xi_f)
 \qquad(S,T\in\mathfrak U_g).
\]
Thus $S\mapsto(S\xi_f\mid\xi_f)$ is a finite normal trace on
$\mathfrak U_g^{+}$. Moreover, \eqref{eq:6.8.3-1} proves that
$t_g'(\pi_f(z))=(\pi_f(z)\xi_f\mid\xi_f)$ for every $z\in A$.

Let $\widetilde A$ be the $C^{*}$-algebra obtained from $A$ by adjoining an
identity element $1$; let $\widetilde f$ be the canonical extension of $f$ to
$\widetilde A$, and let $\widetilde g=\widetilde f|_{\widetilde A^{+}}$. The
space $H(\pi_f)$ is also the Hilbert space obtained by separation and
completion of $\widetilde A$ equipped with the inner product
$\widetilde f(y^{\ast}x)=\widetilde f(xy^{\ast})$, and $\xi_f$ is the
canonical image of $1$. Let $z\in A$. Since the operators of left and right
multiplication by $z$ are continuous on the pre-Hilbert space
$\widetilde A$, as is the involution, we have
\[
 \lambda_g(z)\xi_f=\lambda_{\widetilde g}(z\cdot1)=\Lambda_gz,
 \qquad
 \rho_g(z)\xi_f=\rho_{\widetilde g}(1\cdot z)=\Lambda_gz,
\]
and
\[
 J_g\xi_f=\Lambda_{\widetilde g}1^{\ast}
          =\Lambda_{\widetilde g}1=\xi_f.
\]
It follows that $\xi_f$ is cyclic for $\mathfrak U_g$ and $\mathfrak V_g$,
and therefore separating for $\mathfrak U_g=\mathfrak V_g'$ and
$\mathfrak V_g=\mathfrak U_g'$. The normal trace
$S\mapsto(S\xi_f\mid\xi_f)$ on $\mathfrak U_g^{+}$ is consequently faithful,
and it agrees on $\pi_f(A^{+})$ with the faithful normal trace $t_g$; hence
the two traces are equal (\XRef{A.29}{A 29}). Thus $t_g$ is a finite trace,
and the von Neumann algebras $\mathfrak U_g,\mathfrak V_g$ are finite.

\ResultHeading{6.8.4}{6.8.4. Corollary.}
\begin{itshape}
Let $A$ be a $C^{*}$-algebra, and let $f$ be a semifinite lower
semicontinuous trace on $A^{+}$. In order that $t_f$ be finite on
$\mathfrak U_f^{+}$, it is necessary and sufficient that $f$ be finite.
\end{itshape}

The necessity is clear, since
$f(x)=t_f(\lambda_f(x))$ for every $x\in A^{+}$. The sufficiency follows
from \XRef{6.8.3}{6.8.3}.

\ResultHeading{6.8.5}{6.8.5. Corollary.}
\begin{itshape}
Let $A$ be a $C^{*}$-algebra and let $f$ be a character of $A$. In order that
$f$ be of finite type, it is necessary and sufficient that $f$ be finite.
\end{itshape}

This is a special case of \XRef{6.8.4}{6.8.4}.

\ResultHeading{6.8.6}{6.8.6. Corollary.}
\begin{itshape}
There is a canonical bijective correspondence between
% Source PDF physical page 140
\begin{enumerate}[label=\textit{\alph*.},leftmargin=2em]
 \item the set of quasi-equivalence classes of nonzero finite-type factor
 representations of $A$;
 \item the set of finite characters of norm $1$ of $A$.
\end{enumerate}
\end{itshape}

This follows from \XRef{6.7.4}{6.7.4} and \XRef{6.8.5}{6.8.5}.

Thus the character of a nonzero factor representation of finite type can be
normalized by requiring it to have norm $1$. This normalization is sometimes
incompatible with that of \XRef{6.7.5}{6.7.5}.

\ResultHeading{6.8.7}{6.8.7. Proposition.}
\begin{itshape}
Let $A$ be a $C^{*}$-algebra, and let $C$ be the set of central positive
functionals on $A$ of norm at most $1$.
\begin{enumerate}[label=\textup{(\roman*)}]
 \item $C$ is convex and compact for the weak topology of the dual of $A$.
 \item The extreme points of $C$ are $0$ and the finite characters of norm
 $1$.
 \item $C$ is the weakly closed convex hull of $0$ and of the set of finite
 characters of norm $1$.
\end{enumerate}
\end{itshape}

Assertion (i) is immediate; (iii) follows from (i), (ii), and the
Krein--Milman theorem. On the other hand, by \XRef{6.7.1}{6.7.1}, a finite
character of $A$ is a nonzero finite trace $f$ on $A^{+}$ such that every
finite trace majorized by $f$ is proportional to $f$. This being so, the
proof of (ii) is exactly the same as the corresponding proof of
\XRef{2.5.5}{2.5.5}.

Bibliography: \cite{ref54,ref57,ref111}.

\BookSubsection{6.9}{Complements}

\ResultHeading{6.9.1}{6.9.1.}
Let $A$ be a $C^{*}$-algebra, let $f$ be a semifinite lower semicontinuous
trace on $A^{+}$, and let $s$ be the bitrace associated with $f$. We shall see
that the choice of a closed two-sided ideal $J$ of $A$ defines a decomposition
of $f$.

\begin{enumerate}[label=\textit{\alph*.},leftmargin=2em]
 \item Let $f'=f|_{J^{+}}$, and let $s'$ be the bitrace of $J$ associated
 with $f'$. Then $s'$ is a bitrace of $A$. Let $s_1$ be its canonical maximal
 extension and let $f_1$ be the associated trace on $A^{+}$.
 \item One has
 \[
 H_{f_1}=H_{f'}\subset H_f,\qquad
 \mathfrak U_{f_1}=\mathfrak U_{f'},\qquad t_{f_1}=t_{f'}.
 \]
 For every $x\in A$, $\lambda_{f_1}(x)$ is the restriction of
 $\lambda_f(x)$ to $H_{f'}$. If $x\in A^{+}$, then
 \[
  f_1(x)=t_f(P\lambda_f(x)),
 \]
 where $P$ denotes the orthogonal projection of $H_f$ onto $H_{f'}$.
 \item Let $(u_i)$ be an increasing directed approximate identity of
 $\overline{\mathfrak m_{f'}}$. If $x\in A^{+}$, then
 \[
  f_1(x)=\lim_i f'(x^{1/2}u_ix^{1/2}).
 \]
 \item For every $x\in A^{+}$, put
 \[
  g(x)=t_f((1-P)\lambda_f(x)).
 \]
 Then $g$ is a semifinite lower semicontinuous trace on $A^{+}$ which
 vanishes on $J^{+}$, and $f=f_1+g$. \cite{ref16}
\end{enumerate}

\ResultHeading{6.9.2}{6.9.2.}
Let $H$ be a separable infinite-dimensional Hilbert space, and put
\[
 A=\mathcal L(H)/\mathcal{LC}(H).
\]
Let $t$ be a trace on $A^{+}$ such that
$\mathfrak m_t\ne0$. Then $t=0$. [Every two-sided ideal of $A$ is equal to
$A$, so $\mathfrak m_t=A$. Let $e_1,e_2$ be two infinite-rank orthogonal
projections with sum $1$ on $H$, and let $f_1,f_2$ be their canonical images
in $A$. Then $t(f_1)=t(1)$ and $t(f_2)=t(1)$, whence
$t(1)=2t(1)$ and $t(1)=0$.] \cite{ref46}

% Source PDF physical page 141
\ResultHeading{6.9.3}{6.9.3.}
Let $A$ be the $C^{*}$-algebra of continuous complex-valued functions on
$[0,1]$. For $x\in A^{+}$, put
\[
 f(x)=\int_0^1x(t)t^{-1}\,dt.
\]
Then $f$ is a semifinite lower semicontinuous trace, but
$\overline{\mathfrak m_f}\ne A$. If, for $x\in A^{+}$, one puts
$f'(x)=x(0)$, then $f'$ is a finite continuous trace and $f+f'=f$.
\cite{ref16}

\ResultHeading{6.9.4}{*6.9.4.}
Let $A$ be a $C^{*}$-algebra. Two traceable irreducible representations of
$A$ having the same kernel are equivalent (\XRef{6.7.5}{6.7.5}), but two
traceable factor representations of $A$ having the same kernel may be
disjoint. \cite{ref57}

\ResultHeading{6.9.5}{6.9.5.}
Return to the $C^{*}$-algebra $B$ of \XRef{4.7.18}{4.7.18}. The identity
representation of $B$ is irreducible; let $f$ be its character. Let $f'$ be a
character obtained by composing the canonical map
$B\to B/\mathcal{LC}(H)$ with a character of the commutative $C^{*}$-algebra
$B/\mathcal{LC}(H)$. Then $f$ majorizes $f'$ without being proportional to it.

\ResultHeading{6.9.6}{6.9.6.}
Let $A$ be a $C^{*}$-algebra and let $f$ be a character of $A$. In order that
there exist a family $(f_i)$ of pure positive functionals on $A$ such that
$f=\sum_i f_i$, it is necessary and sufficient that $f$ be of type I.
\cite{ref309}

\ResultHeading{6.9.7}{6.9.7.}
Let $A$ be a separable $C^{*}$-algebra.
\begin{enumerate}[label=\textit{\alph*.},leftmargin=2em]
 \item The set $E$ of $\pi\in\operatorname{Rep}(A)$ for which $\pi$ is
 irreducible and traceable is a Borel subset of $\operatorname{Rep}(A)$.
 \item Let $E',E''$ be the canonical images of $E$ in $\widehat A$ and
 $\operatorname{Prim}(A)$. Then $E'$ is Borel for the Mackey Borel structure,
 and $E''$ is Borel in $\operatorname{Prim}(A)$ (see
 \XRef{3.9.2}{3.9.2}). The set $E'$ is the largest subset of $\widehat A$ on
 which the restriction of $\widehat A\to\operatorname{Prim}(A)$ is injective
 (\XRef{4.7.2}{4.7.2}). In view of \XRef{3.9.2}{3.9.2}, $E'$ is a standard
 Borel subset of $\widehat A$ on which the Mackey Borel structure is the
 Borel structure underlying the topology. \cite{ref57}
\end{enumerate}

\ResultHeading{6.9.8}{6.9.8.}
Let $A$ be a separable $C^{*}$-algebra. The set of
$\pi\in\operatorname{Rep}(A)$ such that
$\pi(A)=\mathcal{LC}(H_\pi)$ is a Borel subset of $\operatorname{Rep}(A)$
(Guichardet, unpublished).

\ResultHeading{6.9.9}{6.9.9.}
Every representation of a postliminal $C^{*}$-algebra is traceable. (Use the
proof of \XRef{5.5.2}{5.5.2}.)

\ResultHeading{6.9.10}{6.9.10.}
Let $E_1$ be the class of separable $C^{*}$-algebras all of whose
representations are of type I. Let $E_1'$ be the class of separable
$C^{*}$-algebras admitting sufficiently many representations of type I [that
is, such that, for every nonzero $x$ in the algebra, there exists a
representation $\pi$ of type I of the algebra such that $\pi(x)\ne0$]. By
successively replacing ``of type I'' by ``semifinite,'' ``finite,'' and
``traceable,'' one obtains the classes
$E_2,E_2',E_3,E_3',E_4,E_4'$ of $C^{*}$-algebras. What can be said about
these classes?

% Source PDF physical page 142
The class $E_1$ is well known (see \SectionRef{9.1}{9.1}). We have
$E_1=E_2=E_4$ by \XRef{6.9.9}{6.9.9}, \SectionRef{9.1}{9.1}, and
\XRef{9.5.4}{9.5.4}. The class $E_1'$ (and \emph{a fortiori} $E_2'$) is the
class of all separable $C^{*}$-algebras, since an irreducible representation
is of type I. Next, $E_3'\subset E_2'=E_1'$ is the class of separable
$C^{*}$-algebras all of whose irreducible representations are
finite-dimensional. The classes $E_4,E_4'$ are distinct from the class of all
separable $C^{*}$-algebras, but are poorly understood. \cite{ref200}

\BookSectionHere{7}{Quasi-Spectrum}

Let $A$ be a separable $C^{*}$-algebra. In \SectionRef{3}{Section 3}, we
defined the Mackey Borel structure on the set of irreducible representations
of $A$, or rather on the set of their equivalence classes. We shall now define
a Borel structure on the set of factor representations of $A$, or rather on
the set of their quasi-equivalence classes. This will be useful in
\SectionRef{8}{Section 8}.

\BookSubsection{7.1}{The Space of Factor Representations}

\ResultHeading{7.1.1}{7.1.1.}
Let $n$ be a cardinal $\leq\aleph_0$, and let $H_n$ be the model Hilbert space
of dimension $n$ (\XRef{3.5.1}{3.5.1}). For every $r\geq0$, denote by
$\mathcal L_r(H_n)$ the closed ball of center $0$ and radius $r$ in
$\mathcal L(H_n)$.

\ResultLabel{Lemma.}
\begin{itshape}
Let $d$ be a metric on $\mathcal L_1(H_n)$ compatible with the weak topology
(\XRef{B.9}{B 9}), let $S\in\mathcal L_1(H_n)$, let $a,b\geq0$, and let
$x_1,\ldots,x_m\in A$. Let $Y$ be the set of $\pi\in\operatorname{Rep}_n(A)$
having the following property: there exist
$T_1,\ldots,T_m\in\mathcal L_b(H_n)\cap\pi(A)'$ such that
\[
 \pi(x_1)T_1+\cdots+\pi(x_m)T_m\in\mathcal L_1(H_n)
\]
and
\[
 d\bigl(\pi(x_1)T_1+\cdots+\pi(x_m)T_m,S\bigr)\leq a.
\]
Then $Y$ is closed in $\operatorname{Rep}_n(A)$.
\end{itshape}

Let $(\pi^i)$ be a net of elements of $Y$, indexed by a directed set, and let
$\pi\in\operatorname{Rep}_n(A)$, with
$\pi^i\to\pi$ in $\operatorname{Rep}_n(A)$. We prove that $\pi\in Y$. For
each $i$, there exist
\[
 T_1^i,\ldots,T_m^i\in\mathcal L_b(H_n)\cap\pi^i(A)'
\]
such that
\[
 \pi^i(x_1)T_1^i+\cdots+\pi^i(x_m)T_m^i\in\mathcal L_1(H_n)
\]
and
\[
 d\bigl(\pi^i(x_1)T_1^i+\cdots+\pi^i(x_m)T_m^i,S\bigr)\leq a.
\]
% Source PDF physical page 143
Since $\mathcal L_b(H_n)$ is weakly compact, we may, on passing to a subnet,
suppose that $T_1^i,\ldots,T_m^i$ converge weakly to
$T_1,\ldots,T_m\in\mathcal L_b(H_n)$. Let $x\in A$ and
$\varphi,\psi\in H_n$. We have
\begin{align*}
 (T_1^i\pi^i(x)\varphi\mid\psi)
 &=(\pi^i(x)\varphi\mid T_1^{i\ast}\psi)
 \longrightarrow(\pi(x)\varphi\mid T_1^{\ast}\psi)
 =(T_1\pi(x)\varphi\mid\psi),\\
 (\pi^i(x)T_1^i\varphi\mid\psi)
 &=(T_1^i\varphi\mid\pi^i(x^{\ast})\psi)
 \longrightarrow(T_1\varphi\mid\pi(x^{\ast})\psi)
 =(\pi(x)T_1\varphi\mid\psi)
 \quad\text{(\XRef{3.5.2}{3.5.2})}.
\end{align*}
Therefore
\[
 (T_1\pi(x)\varphi\mid\psi)=(\pi(x)T_1\varphi\mid\psi),
 \qquad\text{and hence }T_1\in\pi(A)'.
\]
Similarly, $T_2,\ldots,T_m\in\pi(A)'$. Moreover,
\[
 \bigl((\pi^i(x_1)T_1^i+\cdots+\pi^i(x_m)T_m^i)\varphi\mid\psi\bigr)
 \longrightarrow
 \bigl((\pi(x_1)T_1+\cdots+\pi(x_m)T_m)\varphi\mid\psi\bigr).
\]
Thus $\pi^i(x_1)T_1^i+\cdots+\pi^i(x_m)T_m^i$ converges weakly to
$\pi(x_1)T_1+\cdots+\pi(x_m)T_m$, whence
\[
 \pi(x_1)T_1+\cdots+\pi(x_m)T_m\in\mathcal L_1(H_n)
\]
and
\[
 d\bigl(\pi(x_1)T_1+\cdots+\pi(x_m)T_m,S\bigr)\leq a.
\]
Therefore $\pi\in Y$.

\ResultHeading{7.1.2}{7.1.2. Lemma.}
\begin{itshape}
Let $d$ be a metric on $\mathcal L_1(H_n)$ compatible with the weak topology,
let $(x_1,x_2,\ldots)$ be a dense sequence in $A$, and let
$(S_1,S_2,\ldots)$ be a weakly dense sequence in $\mathcal L_1(H_n)$
(\XRef{B.9}{B 9}). Let $k,m,p$ be integers. Let $Y_{k,m,p}$ be the set of
$\pi\in\operatorname{Rep}_n(A)$ having the following property: there exist
$T_1,T_2,\ldots,T_m\in\mathcal L_m(H_n)\cap\pi(A)'$ such that
\[
 \pi(x_1)T_1+\cdots+\pi(x_m)T_m\in\mathcal L_1(H_n)
\]
and
\[
 d\bigl(\pi(x_1)T_1+\cdots+\pi(x_m)T_m,S_k\bigr)\leq\frac1p.
\]
\begin{enumerate}[label=\textup{(\roman*)}]
 \item $Y_{k,m,p}$ is closed in $\operatorname{Rep}_n(A)$.
 \item The set
 \[
   Y=\bigcap_{k,p}\ \bigcup_m Y_{k,m,p}
 \]
 is the set of nonzero factor representations in
 $\operatorname{Rep}_n(A)$.
\end{enumerate}
\end{itshape}

Assertion (i) follows from \XRef{7.1.1}{7.1.1}.

For every $\pi\in\operatorname{Rep}_n(A)$, let $N(\pi)$ be the set of linear
combinations of operators of the form $RR'$, where
$R\in\pi(A)$ and $R'\in\pi(A)'$. Let $\mathfrak M(\pi)$ be the weak closure
of $N(\pi)$. Then $\pi$ is nonzero and factorial if and only if
$\mathfrak M(\pi)=\mathcal L(H_n)$.

% Source PDF physical page 144
First suppose $\pi\in Y$. We shall prove that $S_k\in\mathfrak M(\pi)$ for
every $k$, and this will imply that
$\mathfrak M(\pi)=\mathcal L(H_n)$. Let $p$ be a positive integer. There
exists $m$ such that $\pi\in Y_{k,m,p}$. Hence there exist
$T_1,\ldots,T_m\in\mathcal L_m(H_n)\cap\pi(A)'$ such that
\[
 \pi(x_1)T_1+\cdots+\pi(x_m)T_m\in\mathcal L_1(H_n)
\]
and
\[
 d\bigl(\pi(x_1)T_1+\cdots+\pi(x_m)T_m,S_k\bigr)\leq\frac1p.
\]
Since $\pi(x_1)T_1+\cdots+\pi(x_m)T_m\in N(\pi)$, we have
\[
 d\bigl(N(\pi)\cap\mathcal L_1(H_n),S_k\bigr)\leq\frac1p.
\]
Since $p$ is arbitrary, this proves that $S_k\in\mathfrak M(\pi)$. Thus
$\pi$ is nonzero and factorial.

Conversely, suppose $\pi\in\operatorname{Rep}_n(A)$ and
$\mathfrak M(\pi)=\mathcal L(H_n)$. Let $k,p$ be positive integers. We shall
prove that there exists an $m$ such that $\pi\in Y_{k,m,p}$; this will
complete the proof. By \XRef{A.13}{A 13}, there exists $P\in N(\pi)$ such
that $\lVert P\rVert<1$ and $d(P,S_k)<1/p$. We have
\[
 P=\pi(y_1)T_1+\cdots+\pi(y_q)T_q,
\]
where
\[
 y_1,\ldots,y_q\in A,
 \qquad T_1,\ldots,T_q\in\pi(A)'.
\]
There then exist $x_{i_1},\ldots,x_{i_q}$ sufficiently close to
$y_1,\ldots,y_q$ that
\[
 \lVert\pi(x_{i_1})T_1+\cdots+\pi(x_{i_q})T_q\rVert\leq1
\]
and
\[
 d\bigl(\pi(x_{i_1})T_1+\cdots+\pi(x_{i_q})T_q,S_k\bigr)\leq\frac1p.
\]
If
$m\geq i_1,\ldots,i_q,\lVert T_1\rVert,\ldots,\lVert T_q\rVert$, then
$\pi\in Y_{k,m,p}$, as required.

\ResultHeading{7.1.3}{7.1.3.}
Let $A$ be a separable $C^{*}$-algebra. We denote by
$\operatorname{Fac}_n(A)$ the set of nonzero factor representations of $A$
on $H_n$, and by $\operatorname{Fac}(A)$ the union of the
$\operatorname{Fac}_n(A)$ for $n=1,2,\ldots,\aleph_0$.

\ResultHeading{7.1.4}{7.1.4. Theorem.}
\begin{itshape}
Let $A$ be a separable $C^{*}$-algebra. The set
$\operatorname{Fac}(A)$ is Borel in $\operatorname{Rep}(A)$.
\end{itshape}

This follows immediately from \XRef{7.1.2}{7.1.2}.

Thus $\operatorname{Fac}(A)$, equipped with the Borel structure induced by
that of $\operatorname{Rep}(A)$, is a standard Borel space, since
$\operatorname{Rep}(A)$ is standard (\XRef{3.8.1}{3.8.1}).

Bibliography: \cite{ref25}.

% Source PDF physical page 145
\BookSubsection{7.2}{Definition of the Quasi-Spectrum}

\ResultHeading{7.2.1}{7.2.1.}
We denote by $\check A$ the set of quasi-equivalence classes of nonzero
factor representations of $A$. If $A$ is separable, every factor
representation is quasi-equivalent to a factor representation on a separable
space (\XRef{5.3.7}{5.3.7}); hence the evident canonical map
\[
 \operatorname{Fac}(A)\longrightarrow\check A
\]
is surjective.

\ResultHeading{7.2.2}{7.2.2. Definition.}
\begin{itshape}
Let $A$ be a separable $C^{*}$-algebra. The \emph{Mackey Borel structure} on
$\check A$ is the quotient structure of that on $\operatorname{Fac}(A)$
under the canonical map $\operatorname{Fac}(A)\to\check A$. The Borel
space $\check A$ is called the \emph{quasi-spectrum} of $A$.
\end{itshape}

\ResultHeading{7.2.3}{7.2.3. Proposition.}
\begin{itshape}
Let $A$ be a separable $C^{*}$-algebra. Let $B$ be a Borel subset of
$\operatorname{Fac}(A)$ which meets each quasi-equivalence class in at most
one point, and let
$\psi\colon\operatorname{Fac}(A)\to\check A$ be the canonical map. Then
$\psi(B)$ is Borel in $\check A$, and $\psi|_B$ is an isomorphism of the
Borel space $B$ onto the Borel space $\psi(B)$.
\end{itshape}

Let $\mathcal E(\pi,\pi')$ be the intertwining number of two representations
$\pi,\pi'$ of $A$. By \XRef{3.7.3}{3.7.3}, the set of pairs
$(\pi,\pi')\in\operatorname{Fac}(A)\times\operatorname{Fac}(A)$ such that
$\mathcal E(\pi,\pi')\leq0$ is Borel in
$\operatorname{Fac}(A)\times\operatorname{Fac}(A)$. Its complement $R$ is
the set of pairs $(\pi,\pi')$ in that product such that
$\pi\approx\pi'$ (\XRef{5.3.6}{5.3.6}). Thus
$S=R\cap(\operatorname{Fac}(A)\times B)$ is Borel in
$\operatorname{Fac}(A)\times\operatorname{Fac}(A)$. Let $p$ be the first
projection from this product onto $\operatorname{Fac}(A)$. If
$(\pi,\pi')\in S$ and $(\pi_1,\pi_1')\in S$ satisfy
$p(\pi,\pi')=p(\pi_1,\pi_1')$, then
$\pi'\approx\pi=\pi_1\approx\pi_1'$, and hence $\pi'=\pi_1'$ by the
hypothesis on $B$. Thus $p|_S$ is injective. Consequently $p(S)$ is a Borel
subset of $\operatorname{Fac}(A)$ (\XRef{B.21}{B 21}). But $p(S)$ is the
saturation of $B$ in $\operatorname{Fac}(A)$ for quasi-equivalence. Hence
$\psi(B)$ is Borel in $\check A$. The map $\psi|_B$ from $B$ onto
$\psi(B)$ is a Borel bijection. To prove that it is a Borel isomorphism, it
is enough to prove that the image $\psi(B')$ of every Borel subset $B'$ of
$B$ is Borel. This follows from the preceding argument applied to $B'$ in
place of $B$.

\ResultHeading{7.2.4}{7.2.4. Corollary.}
\begin{itshape}
Every point of $\check A$ is Borel in $\check A$.
\end{itshape}

The quasi-spectrum has hitherto been called the quasi-dual.

Bibliography: \cite{ref25}.

\BookSubsection{7.3}{Relations Between Spectrum and Quasi-Spectrum}

\ResultHeading{7.3.1}{7.3.1.}
Let $A$ be a separable $C^{*}$-algebra. For $p=1,2,\ldots,\infty$, let
$\operatorname{FacI}_p(A)$ be the set of
$\pi\in\operatorname{Fac}(A)$ for which $\pi(A)'$ is a factor of type
$\mathrm I_p$. We have
$\operatorname{FacI}_1(A)=\operatorname{Irr}(A)$. Let
$\operatorname{FacI}(A)$
% Source PDF physical page 146
be the union of the $\operatorname{FacI}_p(A)$. Then
$\operatorname{FacI}(A)$ is the set of $\pi\in\operatorname{Fac}(A)$ for
which $\pi(A)'$ is of type I, that is, the set of representations in
$\operatorname{Fac}(A)$ which are of type I. The sets
$\operatorname{FacI}_p(A)$ are saturated under equivalence of
representations, but not under quasi-equivalence. In contrast,
$\operatorname{FacI}(A)$ is saturated under quasi-equivalence. A
representation in $\operatorname{Rep}(A)$ belongs to
$\operatorname{FacI}_p(A)$ if and only if it is equivalent to a representation
of the form $p\pi$, where $\pi\in\operatorname{Irr}(A)$
(\XRef{5.4.11}{5.4.11}). Thus all the sets $\operatorname{FacI}_p(A)$ have
the same canonical image in $\check A$ as $\operatorname{Irr}(A)$. Denote
this image by $\check A_{\mathrm I}$; it is the set of quasi-equivalence
classes of factor representations of type I.

\ResultHeading{7.3.2}{7.3.2.}
Recall that, on the set of unitary operators on a Hilbert space, the weak and
strong topologies coincide and are compatible with the group structure. This
being so, we have the following lemma.

\ResultLabel{Lemma.}
\begin{itshape}
Let $H$ be a separable Hilbert space and let $\mathcal G$ be the group of
unitary operators on $H$. With the strong (or weak) topology,
$\mathcal G$ is a Polish group.
\end{itshape}

Indeed, the unit ball $\mathcal L_1(H)$ of $\mathcal L(H)$ is compact and
metrizable for the weak topology, and is therefore a Polish space. It is
enough to prove that $\mathcal G$ is a $G_\delta$ in $\mathcal L_1(H)$. If
$(\xi_i)$ is a dense sequence in the unit sphere of $H$, then
$\mathcal G$ is the set of $T\in\mathcal L_1(H)$ such that
\[
 \lVert T\xi_i\rVert>1-\frac1j,
 \qquad
 \lVert T^{\ast}\xi_i\rVert>1-\frac1j
\]
for all integers $i,j$. Each of these inequalities defines a weakly open
subset of $\mathcal L_1(H)$, since for every $\xi\in H$ the map
$T\mapsto\lVert T\xi\rVert$ is lower semicontinuous for the weak topology,
and $T\mapsto T^{\ast}$ is a homeomorphism for the weak topology.

\ResultHeading{7.3.3}{7.3.3. Lemma.}
\begin{itshape}
Let $A$ be a separable $C^{*}$-algebra, let $p$ be a cardinal
$\leq\aleph_0$, and let $B$ be a Borel subset of
$\operatorname{Irr}(A)$ saturated under equivalence. Let $C$ be the set of
elements of $\operatorname{FacI}_p(A)$ quasi-equivalent to an element of $B$.
Then $C$ is Borel in $\operatorname{Rep}(A)$.
\end{itshape}

Let $H_n$ be the model Hilbert space of dimension $n$. The sets
$B\cap\operatorname{Rep}_n(A)$ are saturated under equivalence and are Borel.
It is enough to argue on each of these sets. In other words, we shall henceforth
suppose that the elements of $B$ act on one and the same $H_n$. Then the
elements of $C$ act on $H_{np}$. We give the proof for
$n=p=\aleph_0$; the other cases are handled similarly and somewhat more
simply. Put $H=H_n=H_{np}$.

Let $(K_1,K_2,\ldots)$ be an infinite sequence of pairwise orthogonal closed
infinite-dimensional vector subspaces of $H$ whose Hilbert direct sum is
% Source PDF physical page 147
$H$. Let $J_n$ be an isomorphism from $H$ onto $K_n$. For every
$T\in\mathcal L(H)$, the operators $J_nTJ_n^{-1}$ define, by linearity and
continuity, an operator $T^{\sim}$ on $H$, and
$T\mapsto T^{\sim}$ is an isomorphism from $\mathcal L(H)$ onto a factor
$\mathfrak F$ of type $\mathrm I_\infty$: namely, the set of bounded linear
operators on $H$ which leave the $K_n$ invariant and whose restrictions to
the $K_n$ are carried into one another by the isomorphisms $J_nJ_m^{-1}$.
The commutant $\mathfrak F'$ of $\mathfrak F$ is also a factor of type
$\mathrm I_\infty$. Let $\pi$ be a representation of $A$ on $H$. Denote by
$\pi^{\sim}$ the representation $x\mapsto\pi(x)^{\sim}$ of $A$ on $H$. If
$\pi\in\operatorname{Irr}(A)$, then $\pi^{\sim}(A)$ generates the factor
$\mathfrak F$. The map $\pi\mapsto\pi^{\sim}$ is plainly injective. For every
$\xi$ belonging to one of the $K_i$ and every $x\in A$, the map
$\pi\mapsto\pi^{\sim}(x)\xi$ is continuous. Thus
$\pi\mapsto\pi^{\sim}$ is Borel. By \XRef{B.21}{B 21}, the image
$B^{\sim}$ of $B$ under this map is a Borel subset of
$\operatorname{Rep}(A)$ contained in $\operatorname{FacI}_\infty(A)$.

The set $C$ in the statement is the union of the
$UB^{\sim}U^{-1}$ as $U$ ranges over the group $\mathcal G$ of unitary
operators on $H$ (\XRef{5.3.3}{5.3.3} and \XRef{5.4.11}{5.4.11}). In other
words, $C$ is the image of $\mathcal G\times B^{\sim}$ under the map
\[
 \Theta\colon(U,\sigma)\longmapsto U\sigma U^{-1}
 \quad\text{from }\mathcal G\times B^{\sim}\text{ into }
 \operatorname{Rep}(A).
\]
For every $x\in A$ and every $\xi\in H$, the map
\[
 (U,\sigma)\longmapsto U\sigma(x)U^{-1}\xi
 \quad\text{from }\mathcal G\times\operatorname{Rep}_\infty(A)
 \text{ into }H
\]
is continuous, where $\mathcal G$ is equipped with the strong topology;
hence $\Theta$ is Borel. We shall prove that, for $U\in\mathcal G$, the
following conditions are equivalent, provided $B\ne\varnothing$, to which
case we may plainly restrict ourselves:
\begin{enumerate}[label=\textup{(\arabic*)}]
 \item $U\mathfrak F U^{-1}=\mathfrak F$;
 \item $U=VW$, where $V,W$ are unitary, $V\in\mathfrak F$ and
 $W\in\mathfrak F'$;
 \item $UB^{\sim}U^{-1}=B^{\sim}$;
 \item there exist $\pi_1,\pi_2\in B^{\sim}$ such that
 $U\pi_1U^{-1}=\pi_2$.
\end{enumerate}

(4)$\Rightarrow$(1): let $\pi_1,\pi_2\in B^{\sim}$ satisfy
$U\pi_1U^{-1}=\pi_2$. Then
$U\pi_1(A)U^{-1}=\pi_2(A)$, and, on taking weak closures,
$U\mathfrak F U^{-1}=\mathfrak F$.

(1)$\Rightarrow$(2): if $U\mathfrak F U^{-1}=\mathfrak F$, then $U$ defines
an automorphism of $\mathfrak F$. But every automorphism of a factor of type I
is implemented by a unitary operator belonging to that factor
(\XRef{A.37}{A 37}). There is therefore a unitary $V\in\mathfrak F$ such that
$W=V^{-1}U$ implements the identity automorphism of $\mathfrak F$, that is,
$W\in\mathfrak F'$, and $U=VW$.

(2)$\Rightarrow$(3): let $U=VW$, where $V,W$ are unitary,
$V\in\mathfrak F$, and $W\in\mathfrak F'$. Let $\pi\in B$. We have
$W\pi(x)^{\sim}W^{-1}=\pi(x)^{\sim}$ for every $x\in A$. Moreover, $V$ is
of the form $T^{\sim}$ for a unitary operator $T$ on $H$. Hence
\[
 V\pi(x)^{\sim}V^{-1}=(T\pi(x)T^{-1})^{\sim},
\]
and consequently
\[
 U\pi^{\sim}U^{-1}=(T\pi T^{-1})^{\sim}.
\]
% Source PDF physical page 148
But $T\pi T^{-1}\in B$, since $B$ is saturated under equivalence. Hence
$U\pi^{\sim}U^{-1}\in B^{\sim}$, and finally
$UB^{\sim}U^{-1}=B^{\sim}$.

(3)$\Rightarrow$(4) is evident.

Conditions (1)--(4) define a subgroup $\mathcal G'$ of $\mathcal G$. From
condition (1), it is clear that $\mathcal G'$ is closed in $\mathcal G$. By
\XRef{7.3.2}{7.3.2} and \XRef{B.17}{B 17}, there is a Borel subset
$\Sigma$ of $\mathcal G$ which meets each left coset of $\mathcal G'$ in
exactly one point. We prove that
\[
 \Theta(\mathcal G\times B^{\sim})
 =\Theta(\Sigma\times B^{\sim}).
\]
It is clear that the left-hand side contains the right-hand side. Conversely,
let $U\in\mathcal G$ and $\sigma\in B^{\sim}$. We have $U=U_1U_2$, with
$U_1\in\Sigma$ and $U_2\in\mathcal G'$. Hence
\[
 \Theta(U,\sigma)
 =U_1(U_2\sigma U_2^{-1})U_1^{-1}
 \in\Theta(\Sigma\times B^{\sim})
\]
by condition (3), proving the asserted equality.

We next prove that the restriction of $\Theta$ to
$\Sigma\times B^{\sim}$ is injective. Let
$\pi_1,\pi_2\in B^{\sim}$ and $V_1,V_2\in\Sigma$ satisfy
\[
 V_1\pi_1V_1^{-1}=V_2\pi_2V_2^{-1}.
\]
Then $V_2^{-1}V_1$ satisfies condition (4), so $V_1$ and $V_2$ belong to the
same left coset modulo $\mathcal G'$. Hence $V_1=V_2$, and then
$\pi_1=\pi_2$. It follows that
$C=\Theta(\Sigma\times B^{\sim})$ is Borel in $\operatorname{Rep}(A)$
(\XRef{B.21}{B 21}).

\ResultHeading{7.3.4}{7.3.4. Proposition.}
\begin{itshape}
Let $A$ be a separable $C^{*}$-algebra. Each set
$\operatorname{FacI}_p(A)$ is Borel in $\operatorname{Rep}(A)$.
\end{itshape}

This follows from \XRef{7.3.3}{7.3.3}, applied with
$B=\operatorname{Irr}(A)$.

\ResultHeading{7.3.5}{7.3.5.}
Equip the sets $\operatorname{FacI}_p(A)$ and
\[
 \operatorname{FacI}(A)=\bigcup_p\operatorname{FacI}_p(A)
\]
with the Borel structures induced by that of $\operatorname{Rep}(A)$, or of
$\operatorname{Fac}(A)$. The resulting Borel spaces are standard.

\ResultHeading{7.3.6}{7.3.6.}
The equivalence relation induced on $\operatorname{Irr}(A)$ by
quasi-equivalence is just equivalence of representations
(\XRef{5.3.3}{5.3.3}). Thus the canonical injection $i$ of
$\operatorname{Irr}(A)$ into $\operatorname{FacI}(A)$ passes to the quotient
and defines a bijection $\Phi$ from $\widehat A$ onto
$\check A_{\mathrm I}$ (\XRef{7.3.1}{7.3.1}).

\ResultLabel{Theorem.}
\begin{itshape}
Let $A$ be a separable $C^{*}$-algebra.
\begin{enumerate}[label=\textup{(\roman*)}]
 \item $\check A_{\mathrm I}$ is a Borel subset of $\check A$.
 \item The bijection $\Phi$ is an isomorphism of $\widehat A$ onto
 $\check A_{\mathrm I}$ for the Mackey Borel structures.
\end{enumerate}
\end{itshape}

Assertion (i) follows from \XRef{7.3.4}{7.3.4}. We prove (ii). Let
\[
 q\colon\operatorname{Irr}(A)\longrightarrow\widehat A,
 \qquad
 p\colon\operatorname{FacI}(A)\longrightarrow\check A_{\mathrm I}
\]
be the canonical maps. The diagram is
\[
\begin{array}{ccc}
 \operatorname{Irr}(A) & \xrightarrow{\ i\ } & \operatorname{FacI}(A)\\
 q\big\downarrow && \big\downarrow p\\
 \widehat A & \xrightarrow{\ \Phi\ } & \check A_{\mathrm I}.
\end{array}
\]
Let $S$ be a Borel subset of $\check A_{\mathrm I}$. Then $p^{-1}(S)$ is a
Borel subset of $\operatorname{FacI}(A)$ saturated under quasi-equivalence.
% Source PDF physical page 149
Thus
$i^{-1}(p^{-1}(S))=q^{-1}(\Phi^{-1}(S))$ is a Borel subset of
$\operatorname{Irr}(A)$ saturated under equivalence. Hence
$\Phi^{-1}(S)$ is Borel in $\widehat A$. Conversely, let $T$ be a Borel
subset of $\widehat A$. Then $q^{-1}(T)$ is a Borel subset of
$\operatorname{Irr}(A)$ saturated under equivalence. The set
$p^{-1}(\Phi(T))$ is the saturation of $i(q^{-1}(T))$ under
quasi-equivalence. Applying \XRef{7.3.3}{7.3.3} with
$B=q^{-1}(T)$, we see that $p^{-1}(\Phi(T))$ is Borel in
$\operatorname{FacI}(A)$. Therefore $\Phi(T)$ is Borel in
$\check A_{\mathrm I}$. This proves the theorem.

Henceforth we shall identify the Borel space $\widehat A$ with the Borel
subset $\check A_{\mathrm I}$ of $\check A$.

\ResultHeading{7.3.7}{7.3.7.}
In particular, if $A$ is a separable postliminal $C^{*}$-algebra, the Borel
spaces $\widehat A$ and $\check A$ are identified
(\XRef{5.5.3}{5.5.3}).

Bibliography: \cite{ref15}.

\BookSubsection{7.4}{Finite Part of the Quasi-Spectrum}

\ResultHeading{7.4.1}{7.4.1.}
It is not very clear what topology should be placed on $\check A$. We shall,
however, define a subset $\check A_f$ of $\check A$ on which there is a
fairly satisfactory topology.

Let $A$ be a separable $C^{*}$-algebra, let $H_n$ be the model Hilbert space
of dimension $n$, let $\operatorname{Facf}_n(A)$ be the set of nonzero
finite-type factor representations of $A$ on $H_n$, and let
$\operatorname{Facf}(A)$ be the union of the $\operatorname{Facf}_n(A)$ for
$n=1,2,\ldots,\aleph_0$. [In fact,
$\operatorname{Facf}_n(A)=\operatorname{Fac}_n(A)$ for $n<\aleph_0$.] Denote
by $\check A_f$ the canonical image of $\operatorname{Facf}(A)$ in
$\check A$, that is, the set of quasi-equivalence classes of nonzero
finite-type factor representations of $A$.

\ResultHeading{7.4.2}{7.4.2.}
On the other hand, let $C$ be the convex set of central positive functionals
on $A$ of norm at most $1$. Then $C$, equipped with the weak topology, is a
compact metrizable space. The set $D$ of nonzero extreme points of $C$ is a
$G_\delta$ in $C$ (\XRef{B.13}{B 13}), and is therefore a Polish space for
the weak topology. Moreover, $D$ is the set of finite characters of norm $1$
of $A$ (\XRef{6.8.7}{6.8.7}). By \XRef{6.8.6}{6.8.6}, the map
\[
 g\longmapsto\text{the quasi-equivalence class of }\lambda_g
\]
is a bijection from $D$ onto $\check A_f$. We equip $\check A_f$ with the
topology transported from the weak topology on $D$ by this bijection. Thus
$\check A_f$ becomes a Polish space.

\ResultHeading{7.4.3}{7.4.3. Proposition.}
\begin{itshape}
Let $A$ be a separable $C^{*}$-algebra.
\begin{enumerate}[label=\textup{(\roman*)},nosep]
 \item $\check A_f$ is a Borel subset of $\check A$.
\end{enumerate}
\end{itshape}
% The statement continues on source PDF physical page 150.
% ===== ASSEMBLED TRANSLATION CHUNK 06: chunk_150_179.tex =====
% Source PDF physical page 150
\begin{itshape}
\begin{enumerate}[label=\textup{(\roman*)},start=2,nosep]
\item The Borel structure induced on $\check A_f$ by that of $\check A$ is
the Borel structure underlying the topology of $\check A_f$.
\item The Borel space $\check A_f$ is standard.
\end{enumerate}
\end{itshape}

We shall construct a Borel map $\Omega$ from $D$ into
$\operatorname{Rep}(A)$ such that, for every $g\in D$, $\Omega(g)$ is a
finite-type factor representation with character $g$. If $g,g'\in D$ and
$g\ne g'$, the representations $\Omega(g)$ and $\Omega(g')$ are then not
quasi-equivalent. Since $D$ and $\operatorname{Rep}(A)$ are standard Borel
spaces, $\Omega(D)$ is a Borel subset of $\operatorname{Rep}(A)$ meeting
each quasi-equivalence class in at most one point, and $\Omega$ is a Borel
isomorphism from $D$ onto $\Omega(D)$ (\XRef{B.21}{B 21}). By
\XRef{7.2.3}{7.2.3}, if $\psi$ denotes the canonical map from
$\operatorname{Fac}(A)$ onto $\check A$, then $\psi(\Omega(D))$ is a Borel
subset of $\check A$, and
$\psi|_{\Omega(D)}$ is a Borel isomorphism from $\Omega(D)$ onto
$\psi(\Omega(D))$. But $\psi(\Omega(D))=\check A_f$ by
\XRef{6.8.6}{6.8.6}. Proposition \XRef{7.4.3}{7.4.3} will therefore be
proved.

Let $(x_1,x_2,\ldots)$ be a dense sequence in $A$. Retain the notation of
\XRef{6.2.2}{6.2.2}. For every $g\in D$, the vectors $\Lambda_gx_n$ are
dense in $H_g$, since $g$ is continuous. The functions
\[
g\longmapsto(\Lambda_gx_n\mid\Lambda_gx_p)=g(x_nx_p^{\ast})
\]
are continuous on $D$. Hence, by \XRef{A.98}{A 98}, there is on the field
$g\mapsto H_g$ one and only one Borel-field structure for which the vector
fields $g\mapsto\Lambda_gx_n$ are Borel. Let $D_p$ be the set of $g\in D$
such that
\[
\dim H_g=p\qquad(p=1,2,\ldots,\aleph_0);
\]
this is a Borel subset of $D$ (\XRef{A.96}{A 96}). It is enough to
construct $\Omega$ on each $D_p$. There exists an isomorphism
$g\mapsto V(g)$ from the field $g\mapsto H_g$ over $D_p$ onto the constant
field defined by the model Hilbert space $H_p$ (\XRef{A.96}{A 96}). It is
then enough to set
\[
\Omega(g)(x)=V(g)\lambda_g(x)V(g)^{-1}
\]
for every $x\in A$. Since $g\mapsto\lambda_g(x)$ is a Borel operator
field, the map $g\mapsto\Omega(g)(x)$ is Borel; thus $\Omega$ is a Borel
map from $D$ into $\operatorname{Rep}(A)$.

\ResultHeading{7.4.4}{7.4.4.}
It follows from \XRef{7.4.3}{7.4.3} that
$\operatorname{Facf}(A)$ and $\operatorname{Facf}_n(A)$ are Borel
subsets of $\operatorname{Fac}(A)$.

Bibliography: \cite{ref57}.

\BookSubsection{7.5}{Complements}

\ResultHeading{7.5.1}{7.5.1.}
Let $A$ be a $C^{*}$-algebra, $I$ a closed two-sided ideal of $A$, and
$\omega$ the canonical homomorphism from $A$ onto $A/I$. Let $E$ be the
set of $\rho\in\check A$ that vanish on $I$. Then
\[
\pi\longmapsto\pi\circ\omega
\]
defines a bijection from $\check{(A/I)}$ onto $E$, and
$\pi\mapsto\pi|_I$ defines a bijection from $\check A\setminus E$ onto
$\check I$. (Use \XRef{2.11.1}{2.11.1}.) If $A$ is
% Source PDF physical page 151
separable, $E$ is a Borel subset of $\check A$, and the preceding
bijections are Borel isomorphisms. (Use the existence of an approximate
identity in $I$.) In particular, $\check A$ is identified with the Borel
sum of a standard space and a space $\check B$, where $B$ is a separable
antiliminal $C^{*}$-algebra. \cite{ref17}

\ResultHeading{7.5.2}{7.5.2.}
Let $A$ be a separable $C^{*}$-algebra. Let $M$ be the set of
$\pi\in\operatorname{Rep}(A)$ that are homogeneous
(\XRef{5.7.6}{5.7.6}). Then $M$ is a Borel subset of
$\operatorname{Rep}(A)$. \cite{ref20}

\ResultHeading{7.5.3}{7.5.3.}
Let $A$ be a separable $C^{*}$-algebra. On
$\widehat A\cap\check A_f$, the topologies induced by the topology of
$\widehat A$ and by that of $\check A_f$ do not necessarily coincide
(take for $A$ the example of \XRef{4.7.19}{4.7.19}).

\ResultHeading{7.5.4}{7.5.4.}
Problems: let $A$ be a separable $C^{*}$-algebra. Is the set of
$\pi\in\operatorname{Fac}(A)$ for which $\pi$ is traceable (respectively,
of type II) a Borel subset of $\operatorname{Rep}(A)$? Is the set $E$ of
traceable elements of $\check A$ a standard Borel space (cf.
\XRef{6.9.7}{6.9.7})? In the absence of a reasonable topology on
$\check A$, is there a reasonable topology on $E$? \cite{ref57}

\BookSection{8}{Integration and Disintegration of Representations}

We now have enough tools to address directly one of the two problems
mentioned at the beginning of \S3: the decomposition of a representation
into irreducible representations. As has already been said, this will not
be a decomposition into a Hilbert-space direct sum of irreducible
representations, but a subtler decomposition (which we shall call
disintegration). This can be anticipated from the following very simple
example. Let $A$ be the $C^{*}$-algebra of continuous complex-valued
functions on $[0,1]$; for every $f\in A$, let $\pi(f)$ be the operator of
multiplication by $f$ in the Hilbert space $H=L^2([0,1])$. Then $\pi$ is a
representation of $A$ on $H$. The irreducible representations of $A$ are
its characters $f\mapsto f(t_0)$, where $t_0\in[0,1]$. To decompose $\pi$
as a Hilbert-space direct sum of irreducible representations would
therefore amount to finding an orthonormal basis of $H$ each of whose
elements is a simultaneous eigenvector of all the $\pi(f)$. But no vector
(except $0$) is a simultaneous eigenvector of all the $\pi(f)$. If we
cease to speak rigorously, however, the Dirac functions at the different
points of $[0,1]$ may be regarded as pairwise orthogonal simultaneous
eigenvectors of the $\pi(f)$, and every element of $H$ may be written as
an integral linear combination of these Dirac functions.

Throughout this section, $A$ denotes a separable $C^{*}$-algebra.

\BookSubsection{8.1}{Integration of Representations}

\ResultHeading{8.1.1}{8.1.1.}
Let $Z$ be a Borel space, $\mu$ a positive measure
(\XRef{B.30}{B 30}) on $Z$, and $\zeta\mapsto H(\zeta)$ a
$\mu$-measurable field of Hilbert spaces over $Z$ (\XRef{A.69}{A 69}).
% Source PDF physical page 152
For every $\zeta\in Z$, let $\pi(\zeta)$ be a representation of $A$ on
$H(\zeta)$; we say that $\zeta\mapsto\pi(\zeta)$ is a field of
representations of $A$.

\ResultLabel{Definition.}
\begin{itshape}
The field of representations $\zeta\mapsto\pi(\zeta)$ is called
measurable if, for every $x\in A$, the operator field
$\zeta\mapsto\pi(\zeta)(x)$ is measurable (\XRef{A.77}{A 77}).
\end{itshape}

\ResultHeading{8.1.2}{8.1.2.}
We have $\lVert\pi(\zeta)(x)\rVert\leq\lVert x\rVert$ for every $x\in A$
and every $\zeta\in Z$. If the field $\zeta\mapsto\pi(\zeta)$ is
measurable, then, for every $x\in A$, we may form the continuous operator
\[
\pi(x)=\int_Z^\oplus\pi(\zeta)(x)\,d\mu(\zeta)
\]
on the Hilbert space
\[
H=\int_Z^\oplus H(\zeta)\,d\mu(\zeta).
\]

\ResultLabel{Proposition.}
\begin{itshape}
The map $x\mapsto\pi(x)$ is a representation of $A$ on $H$.
\end{itshape}

For $x,y\in A$, we have
\begin{align*}
\pi(x+y)
 &=\int_Z^\oplus\pi(\zeta)(x+y)\,d\mu(\zeta)\\
 &=\int_Z^\oplus\bigl(\pi(\zeta)(x)+\pi(\zeta)(y)\bigr)\,d\mu(\zeta)\\
 &=\int_Z^\oplus\pi(\zeta)(x)\,d\mu(\zeta)
   +\int_Z^\oplus\pi(\zeta)(y)\,d\mu(\zeta)\\
 &=\pi(x)+\pi(y).
\end{align*}
Likewise,
\[
\pi(\lambda x)=\lambda\pi(x),\qquad
\pi(xy)=\pi(x)\pi(y),\qquad
\pi(x^{\ast})=\pi(x)^{\ast}.
\]

\ResultHeading{8.1.3}{8.1.3. Definition.}
\begin{itshape}
With the preceding notation, $\pi$ is called the direct integral of the
$\pi(\zeta)$, and we write
\[
\pi=\int_Z^\oplus\pi(\zeta)\,d\mu(\zeta).
\]
\end{itshape}

\ResultHeading{8.1.4}{8.1.4.}
Every operator $\pi(x)$ is decomposable, and hence commutes with the
algebra $\mathfrak Z$ of diagonalizable operators (\XRef{A.80}{A 80}).
In other words, $\mathfrak Z\subset\pi(A)'$.

\ResultHeading{8.1.5}{8.1.5.}
Let $K$ be the essential subspace of $\pi$ (\XRef{2.2.6}{2.2.6}). Then
$P_K$ commutes with $\pi(A)'$, hence commutes with $\mathfrak Z$; consequently
there is a measurable field $\zeta\mapsto E(\zeta)\in\mathcal L(H(\zeta))$
of projections such that
\[
P_K=\int_Z^\oplus E(\zeta)\,d\mu(\zeta).
\]
Let $(u_1,u_2,\ldots)$ be an approximate identity of $A$
(\XRef{1.7.2}{1.7.2}). We have $\pi(u_n)\to P_K$ strongly; hence, after
passing if necessary to a subsequence of $(u_n)$,
$\pi(\zeta)(u_n)\to E(\zeta)$ almost everywhere (\XRef{A.79}{A 79}).
Thus, almost everywhere, $E(\zeta)(H(\zeta))$ is the essential subspace of
$\pi(\zeta)$. In particular, $\pi$ is nondegenerate if and only if almost
all the $\pi(\zeta)$ are nondegenerate.

\ResultHeading{8.1.6}{8.1.6. Proposition.}
\begin{itshape}
Let $\zeta\mapsto\pi(\zeta)$ and $\zeta\mapsto\pi'(\zeta)$ be two
measurable fields of representations of $A$ relative to the same
measurable field $\zeta\mapsto H(\zeta)$. Let
\[
\pi=\int_Z^\oplus\pi(\zeta)\,d\mu(\zeta),
\qquad
\pi'=\int_Z^\oplus\pi'(\zeta)\,d\mu(\zeta).
\]
Then $\pi=\pi'$ if and only if
$\pi(\zeta)=\pi'(\zeta)$ almost everywhere.
\end{itshape}

% Source PDF physical page 153
The condition is plainly sufficient. Suppose that $\pi=\pi'$. Let
$(x_i)$ be a dense sequence in $A$. We have
$\pi(\zeta)(x_i)=\pi'(\zeta)(x_i)$ outside a null set $N_i$. Then
$N=\bigcup_iN_i$ is null, and, for $\zeta\notin N$, we have
$\pi(\zeta)(x_i)=\pi'(\zeta)(x_i)$ for every $i$, hence
$\pi(\zeta)=\pi'(\zeta)$.

It follows from \XRef{8.1.6}{8.1.6} that the $\pi(\zeta)$ are determined
by $\pi$ up to null sets.

\ResultHeading{8.1.7}{8.1.7. Proposition.}
\begin{itshape}
Let $\zeta\mapsto\pi(\zeta)$ be a measurable field of representations of
$A$ on the $H(\zeta)$. Let
\[
H=\int_Z^\oplus H(\zeta)\,d\mu(\zeta),
\qquad
\pi=\int_Z^\oplus\pi(\zeta)\,d\mu(\zeta).
\]
Suppose that there exists a representation $\pi_0$ on a Hilbert space
$H_0$ such that $\pi(\zeta)\simeq\pi_0$ for every $\zeta$, and suppose
that $\mu$ is standard. Then there exists an isomorphism from $H$ onto
$L^2(Z,\mu)\otimes H_0$ which carries $\pi(x)$ to
$1\otimes\pi_0(x)$ for every $x\in A$. Thus $\pi$ is a multiple of
$\pi_0$.
\end{itshape}

This is a special case of \XRef{A.83}{A 83}.

\ResultHeading{8.1.8}{8.1.8.}
Let $H_1,H_2,\ldots,H_{\aleph_0}$ be the model Hilbert spaces of
dimensions $1,2,\ldots,\aleph_0$.

\ResultLabel{Proposition.}
\begin{itshape}
Suppose that $Z$ is the union of pairwise disjoint Borel sets
$Z_1,\allowbreak Z_2,\allowbreak\ldots,\allowbreak Z_{\aleph_0}$, and that,
over $Z_n$, the field
$\zeta\mapsto H(\zeta)$ reduces to the constant field defined by $H_n$.
Let $\zeta\mapsto\pi(\zeta)$ be a field of representations of $A$ on the
$H(\zeta)$. Then $\zeta\mapsto\pi(\zeta)$ is measurable if and only if,
almost everywhere, it agrees with a Borel map from $Z$ into
$\operatorname{Rep}(A)$.
\end{itshape}

The condition is plainly sufficient. Conversely, let $(\xi_j)$ be an
orthonormal basis of $H_n$, and let $(x_i)$ be a dense sequence in $A$. If
the field $\zeta\mapsto\pi(\zeta)$ is measurable, the functions
\[
\zeta\longmapsto(\pi(\zeta)(x_i)\xi_j\mid\xi_k)
\]
are measurable on $Z_n$, and hence agree almost everywhere with Borel
functions. By changing $\zeta\mapsto\pi(\zeta)$ on a null subset of
$Z_n$, we may arrange that all the maps
\[
\zeta\longmapsto(\pi(\zeta)(x_i)\xi_j\mid\xi_k)
\]
are Borel. Then, for every $x\in A$ and all $\xi,\eta\in H_n$, the map
$\zeta\mapsto(\pi(\zeta)(x)\xi\mid\eta)$ is Borel; hence
$\zeta\mapsto\pi(\zeta)$ is a Borel map from $Z_n$ into
$\operatorname{Rep}_n(A)$.

Bibliography:
\cite{ref20,ref52,ref88,ref91,ref95,ref96,ref101,ref108,ref110,ref137,ref161}.

\BookSubsection{8.2}{Equivalence of Two Integrals of Representations}

\ResultHeading{8.2.1}{8.2.1. Proposition.}
\begin{itshape}
Let $Z$ be a Borel space, $\mu$ a positive measure on $Z$,
$\zeta\mapsto H(\zeta)$ a measurable field of Hilbert spaces over $Z$,
and
% Source PDF physical page 154
$\zeta\mapsto\pi(\zeta)$ a measurable field of representations of $A$ on
the $H(\zeta)$. Put
\[
H=\int_Z^\oplus H(\zeta)\,d\mu(\zeta),\qquad
\pi=\int_Z^\oplus\pi(\zeta)\,d\mu(\zeta),
\]
and let $\mathfrak Z$ be the algebra of diagonalizable operators on $H$.
Let $\widetilde\mu$ be a positive measure on $Z$ equivalent to $\mu$, and
put
\[
\widetilde H=\int_Z^\oplus H(\zeta)\,d\widetilde\mu(\zeta),\qquad
\widetilde\pi=\int_Z^\oplus\pi(\zeta)\,d\widetilde\mu(\zeta),
\]
with $\widetilde{\mathfrak Z}$ the algebra of diagonalizable operators on
$\widetilde H$. The canonical isomorphism from $H$ onto $\widetilde H$
(\XRef{A.75}{A 75}) carries $\pi$ to $\widetilde\pi$ and
$\mathfrak Z$ to $\widetilde{\mathfrak Z}$.
\end{itshape}

Let
\[
r(\zeta)=\frac{d\widetilde\mu(\zeta)}{d\mu(\zeta)}.
\]
The canonical isomorphism $U$ from $H$ onto $\widetilde H$ is the
isomorphism which carries
\[
x=\int_Z^\oplus x(\zeta)\,d\mu(\zeta)\in H
\]
to
\[
Ux=\int_Z^\oplus r(\zeta)^{-1/2}x(\zeta)\,
            d\widetilde\mu(\zeta).
\]
It is clear that $U$ carries $\mathfrak Z$ to
$\widetilde{\mathfrak Z}$. Let $a\in A$. Then
\begin{align*}
\widetilde\pi(a)Ux
 &=\int_Z^\oplus\pi(\zeta)(a)r(\zeta)^{-1/2}x(\zeta)\,
       d\widetilde\mu(\zeta)\\
 &=U\int_Z^\oplus\pi(\zeta)(a)x(\zeta)\,d\mu(\zeta)
  =U\pi(a)x,
\end{align*}
so $U$ carries $\pi$ to $\widetilde\pi$.

\ResultHeading{8.2.2}{8.2.2. Proposition.}
\begin{itshape}
Let $Z$ be a Borel space, $\mu$ a positive measure on $Z$,
$\zeta\mapsto H(\zeta)$ a $\mu$-measurable field of Hilbert spaces over
$Z$, and $\zeta\mapsto\pi(\zeta)$ a $\mu$-measurable field of
representations of $A$ on the $H(\zeta)$. Put
\[
H=\int_Z^\oplus H(\zeta)\,d\mu(\zeta),\qquad
\pi=\int_Z^\oplus\pi(\zeta)\,d\mu(\zeta),
\]
and let $\mathfrak Z$ be the algebra of diagonalizable operators on $H$.
Define analogously
\[
Z_1,\ \mu_1,\ \zeta_1\mapsto H_1(\zeta_1),\
\zeta_1\mapsto\pi_1(\zeta_1),\ H_1,\ \pi_1,\ \mathfrak Z_1.
\]
Suppose that there exist:
\begin{enumerate}[label=\textup{\arabic*\textdegree}]
\item a Borel $\mu$-null subset $N$ of $Z$ and a Borel
  $\mu_1$-null subset $N_1$ of $Z_1$;
\item a Borel isomorphism $\eta$ from $Z\setminus N$ onto
  $Z_1\setminus N_1$ carrying $\mu$ to $\mu_1$;
\item an $\eta$-isomorphism $\zeta\mapsto V(\zeta)$ from the field
  $\zeta\mapsto H(\zeta)$ over $Z\setminus N$ onto the field
  $\zeta_1\mapsto H_1(\zeta_1)$ over $Z_1\setminus N_1$
  (\XRef{A.70}{A 70}, \XRef{A.74}{A 74}), such that $V(\zeta)$ carries
  $\pi(\zeta)$ to $\pi_1(\eta(\zeta))$ for every $\zeta$.
\end{enumerate}
Then
\[
V=\int_Z^\oplus V(\zeta)\,d\mu(\zeta)
\]
carries $\mathfrak Z$ to $\mathfrak Z_1$ and $\pi$ to $\pi_1$.
\end{itshape}

% Source PDF physical page 155
Denote by $1_{H(\zeta)}$ and $1_{H_1(\zeta_1)}$ the identity operators on
$H(\zeta)$ and $H_1(\zeta_1)$. If $f\in L^\infty(Z,\mu)$ and $f_1$ is the
function on $Z_1\setminus N_1$ transported by $\eta$ from
$f|_{Z\setminus N}$, then $V$ carries
\[
\int_Z^\oplus f(\zeta)1_{H(\zeta)}\,d\mu(\zeta)
\quad\text{to}\quad
\int_{Z_1}^\oplus f_1(\zeta_1)1_{H_1(\zeta_1)}\,d\mu_1(\zeta_1);
\]
hence $V$ carries $\mathfrak Z$ to $\mathfrak Z_1$. On the other hand,
let $a\in A$ and
\[
x=\int_Z^\oplus x(\zeta)\,d\mu(\zeta)\in H.
\]
Then
\begin{align*}
V\pi(a)x
 &=V\left(\int_Z^\oplus\pi(\zeta)(a)x(\zeta)\,d\mu(\zeta)\right)\\
 &=\int_{Z_1}^\oplus
 V(\eta^{-1}(\zeta_1))\pi(\eta^{-1}(\zeta_1))(a)
 x(\eta^{-1}(\zeta_1))\,d\mu_1(\zeta_1)\\
 &=\int_{Z_1}^\oplus
 \pi_1(\zeta_1)(a)V(\eta^{-1}(\zeta_1))
 x(\eta^{-1}(\zeta_1))\,d\mu_1(\zeta_1)\\
 &=\pi_1(a)Vx,
\end{align*}
so $V$ carries $\pi$ to $\pi_1$.

\ResultHeading{8.2.3}{8.2.3. Proposition.}
\begin{itshape}
Retain the notation $Z,\mu,H(\zeta),\pi(\zeta),H,\pi,\mathfrak Z$ and
$Z_1,\ldots,\mathfrak Z_1$ of \XRef{8.2.2}{8.2.2}. Suppose that $N,N_1$
and $\eta$ exist with properties $1^\circ$ and $2^\circ$ of
\XRef{8.2.2}{8.2.2}. Suppose further that, for every
$\zeta\in Z\setminus N$, the representations $\pi(\zeta)$ and
$\pi_1(\eta(\zeta))$ are equivalent. If, in addition, $\mu$ is standard,
there exists an isomorphism $V$ from $H$ onto $H_1$ carrying
$\mathfrak Z$ to $\mathfrak Z_1$ and $\pi$ to $\pi_1$.
\end{itshape}

There exists (\XRef{A.82}{A 82}) an $\eta$-isomorphism
$(V(\zeta))_{\zeta\in Z\setminus N}$ from the field
$(H(\zeta))_{\zeta\in Z\setminus N}$ onto the field
$(H_1(\zeta_1))_{\zeta_1\in Z_1\setminus N_1}$ such that $V(\zeta)$
carries $\pi(\zeta)(x)$ to $\pi_1(\eta(\zeta))(x)$ for every
$\zeta\in Z$ and every $x\in D$, where $D$ is a countable dense subset of
$A$; hence it carries $\pi(\zeta)$ to $\pi_1(\eta(\zeta))$.
The proposition now follows from \XRef{8.2.2}{8.2.2}.

\ResultHeading{8.2.4}{8.2.4.}
Propositions \XRef{8.2.1}{8.2.1} and \XRef{8.2.2}{8.2.2} have a converse:

\ResultLabel{Proposition.}
\begin{itshape}
Retain the notation $Z,\mu,H(\zeta),\pi(\zeta),H,\pi,\mathfrak Z$ and
$Z_1,\ldots,\mathfrak Z_1$ of \XRef{8.2.2}{8.2.2}. Suppose that there is
an isomorphism $U$ from $H$ onto $H_1$ carrying $\mathfrak Z$ to
$\mathfrak Z_1$ and $\pi$ to $\pi_1$. If $\mu$ and $\mu_1$ are standard,
there exist:
\begin{enumerate}[label=\textup{\arabic*\textdegree}]
\item a Borel $\mu$-null set $N$ in $Z$ and a Borel $\mu_1$-null set
  $N_1$ in $Z_1$;
\item a Borel isomorphism $\eta$ from $Z\setminus N$ onto
  $Z_1\setminus N_1$ carrying $\mu$ to a measure
  $\widetilde\mu_1$ equivalent to $\mu_1$;
\item an $\eta$-isomorphism $\zeta\mapsto V(\zeta)$ from the field
  $\zeta\mapsto H(\zeta)$ over $Z\setminus N$ onto the field
  $\zeta_1\mapsto H_1(\zeta_1)$ over $Z_1\setminus N_1$, such that
  $V(\zeta)$ carries $\pi(\zeta)$ to $\pi_1(\eta(\zeta))$ for every
  $\zeta\in Z\setminus N$.
\end{enumerate}
Moreover, $U$ is the composition of
\[
\int_Z^\oplus V(\zeta)\,d\mu(\zeta)
\]
and the canonical isomorphism from
$\int_{Z_1}^\oplus H_1(\zeta_1)\,d\widetilde\mu_1(\zeta_1)$ onto $H_1$.
\end{itshape}

% Source PDF physical page 156
Since $U$ carries $\mathfrak Z$ to $\mathfrak Z_1$, there exist
$N,N_1,\eta,\widetilde\mu_1$ and $\zeta\mapsto V(\zeta)$ with all the
properties in \XRef{8.2.4}{8.2.4}, except for the assertion that
$V(\zeta)$ carries $\pi(\zeta)$ to $\pi_1(\eta(\zeta))$ for
$\zeta\in Z\setminus N$ (\XRef{A.85}{A 85}). Since $U$ is the composition
of $\int_Z^\oplus V(\zeta)\,d\mu(\zeta)$ and the canonical isomorphism
from
$\int_{Z_1}^\oplus H_1(\zeta_1)\,d\widetilde\mu_1(\zeta_1)$ onto $H_1$,
$U$ carries $\pi$ to
\[
\int_{Z_1}^\oplus\pi'_1(\zeta_1)\,d\mu_1(\zeta_1),
\]
where $\pi'_1(\zeta_1)$ is the transform of
$\pi(\eta^{-1}(\zeta_1))$ by $V(\eta^{-1}(\zeta_1))$
(\XRef{8.2.1}{8.2.1} and \XRef{8.2.2}{8.2.2}). Since $U$ carries $\pi$
to $\pi_1$, we have
$\pi_1(\zeta_1)=\pi'_1(\zeta_1)$ almost everywhere
(\XRef{8.1.6}{8.1.6}). By adjoining a suitable null set to $N$, we
therefore see that $V(\zeta)$ carries $\pi(\zeta)$ to
$\pi_1(\eta(\zeta))$ for $\zeta\in Z\setminus N$.

Bibliography:
\cite{ref52,ref88,ref91,ref95,ref96,ref101,ref108,ref110,ref137,ref161}.

\BookSubsection{8.3}{Disintegration of Representations}

\ResultHeading{8.3.1}{8.3.1. Lemma.}
\begin{itshape}
Let $Z$ be a Borel space, $\mu$ a positive measure on $Z$,
$\zeta\mapsto H(\zeta)$ a $\mu$-measurable field of Hilbert spaces over
$Z$,
\[
H=\int_Z^\oplus H(\zeta)\,d\mu(\zeta),
\]
and $\pi$ a representation of $A$ on $H$. Suppose that every operator
$\pi(x)$, $x\in A$, is decomposable. Then, for every $\zeta\in Z$, there
exists a representation $\pi(\zeta)$ of $A$ on $H(\zeta)$ such that
\[
\pi=\int_Z^\oplus\pi(\zeta)\,d\mu(\zeta).
\]
\end{itshape}

Let $(x_1,x_2,\ldots)$ be a sequence of distinct points of $A$, dense in
$A$. By enlarging this sequence if necessary, we may suppose that the set
of the $x_i$ is an involutive subalgebra $B$ of $A$ over the (countable)
field of rational complex numbers. For each $i$, choose arbitrarily a
measurable field $\zeta\mapsto T_i(\zeta)$ such that
\[
\pi(x_i)=\int_Z^\oplus T_i(\zeta)\,d\mu(\zeta).
\]
Let $N_i$ be the null set of $\zeta\in Z$ such that
$\lVert T_i(\zeta)\rVert>\lVert\pi(x_i)\rVert$. The union $N$ of the
$N_i$ is null. On the other hand, there are countably many relations
among the $x_i$ of the forms
\[
x_k=\lambda_1x_i+\lambda_2x_j,\qquad
x_k=x_ix_j,\qquad
x_j=x_i^{\ast},
\]
where $\lambda_1,\lambda_2$ are rational complex numbers. Almost
everywhere these relations imply, respectively,
\[
T_k(\zeta)=\lambda_1T_i(\zeta)+\lambda_2T_j(\zeta),\qquad
T_k(\zeta)=T_i(\zeta)T_j(\zeta),\qquad
T_j(\zeta)=T_i(\zeta)^{\ast}.
\]
There is therefore a null subset $N'$ of $Z$ such that, for
$\zeta\notin N'$, the map $x_i\mapsto T_i(\zeta)$ is a homomorphism from
$B$ into $\mathcal L(H(\zeta))$, regarded as an involutive algebra over
the field of rational complex numbers. Replace the $T_i(\zeta)$ by $0$
for $\zeta\in N\cup N'$. Then, for every $\zeta\in Z$,
\[
\lVert T_i(\zeta)\rVert\leq\lVert\pi(x_i)\rVert\leq\lVert x_i\rVert
\]
for every $i$, and the map $x_i\mapsto T_i(\zeta)$
% Source PDF physical page 157
is a homomorphism from $B$ into $\mathcal L(H(\zeta))$. This homomorphism
extends uniquely to a representation $\pi(\zeta)$ of $A$ on $H(\zeta)$.
Let $x\in A$, and write
\[
\pi(x)=\int_Z^\oplus T(\zeta)\,d\mu(\zeta).
\]
There exists a sequence $(x_{n_i})$ of elements of $B$ tending to $x$.
Almost everywhere,
\[
\lVert T(\zeta)-\pi(\zeta)(x_{n_i})\rVert
\leq\lVert x-x_{n_i}\rVert;
\]
hence $\pi(\zeta)(x_{n_i})$ tends to $T(\zeta)$ almost everywhere. Thus
$T(\zeta)=\pi(\zeta)(x)$ almost everywhere, and consequently
\[
\pi(x)=\int_Z^\oplus\pi(\zeta)(x)\,d\mu(\zeta).
\]

\ResultHeading{8.3.2}{8.3.2. Theorem.}
\begin{itshape}
Let $H$ be a separable Hilbert space, $\pi$ a representation of $A$ on
$H$, and $\mathfrak Z$ a commutative von Neumann subalgebra of
$\pi(A)'$. There exist a standard Borel space $Z$, a bounded positive
measure $\mu$ on $Z$, a measurable field $\zeta\mapsto H(\zeta)$ of
Hilbert spaces over $Z$, a measurable field $\zeta\mapsto\pi(\zeta)$ of
representations of $A$ on the $H(\zeta)$, and an isomorphism from $H$ onto
\[
\int_Z^\oplus H(\zeta)\,d\mu(\zeta)
\]
which carries $\mathfrak Z$ to the algebra of diagonalizable operators
and $\pi$ to
\[
\int_Z^\oplus\pi(\zeta)\,d\mu(\zeta).
\]
\end{itshape}

There exist (\XRef{A.84}{A 84}) a standard Borel space $Z$, a bounded
positive measure $\mu$ on $Z$, a measurable field
$\zeta\mapsto H(\zeta)$ of Hilbert spaces over $Z$, and an isomorphism
from $H$ onto $\int_Z^\oplus H(\zeta)\,d\mu(\zeta)$ which carries
$\mathfrak Z$ to the algebra of diagonalizable operators. Since
$\mathfrak Z\subset\pi(A)'$, every operator $\pi(x)$, $x\in A$, is
decomposable. It remains only to apply \XRef{8.3.1}{8.3.1}.

Bibliography:
\cite{ref52,ref88,ref91,ref95,ref96,ref101,ref108,ref110,ref137,ref161}.

\BookSubsection{8.4}{Central Disintegration}

\ResultHeading{8.4.1}{8.4.1. Lemma.}
\begin{itshape}
Let $Z$ be a Borel space, $\mu$ a positive measure on $Z$,
$\zeta\mapsto H(\zeta)$ a measurable field of Hilbert spaces over $Z$,
$\zeta\mapsto\pi(\zeta)$ a measurable field of representations of $A$ on
the $H(\zeta)$, $\mathcal A(\zeta)$ the strong closure of
$\pi(\zeta)(A)$,
\[
H=\int_Z^\oplus H(\zeta)\,d\mu(\zeta),\qquad
\pi=\int_Z^\oplus\pi(\zeta)\,d\mu(\zeta),
\]
$\mathfrak Z$ the algebra of diagonalizable operators, and $\mathcal A$
the strong closure of $\pi(A)$. Suppose that
$\mathfrak Z\subset\mathcal A$.
\begin{enumerate}[label=\textup{(\roman*)}]
\item We have
\[
\mathcal A=\int_Z^\oplus\mathcal A(\zeta)\,d\mu(\zeta).
\]
\item If $\mathfrak Z$ is the center of $\mathcal A$, then
$\pi(\zeta)$ is a factor representation almost everywhere.
\item If $\mu$ is standard, there exists a null subset $N$ of $Z$ such
that, whenever $\zeta,\zeta'\in Z\setminus N$ and $\zeta\ne\zeta'$, one
has $\pi(\zeta)\perp\pi(\zeta')$.
\end{enumerate}
\end{itshape}

% Source PDF physical page 158
Since $\mathfrak Z\subset\mathcal A$, $\pi$ is nondegenerate, and hence
$\pi(\zeta)$ is nondegenerate almost everywhere (\XRef{8.1.5}{8.1.5}).
Thus $\mathcal A$ (respectively, $\mathcal A(\zeta)$) is also the von
Neumann algebra generated by $\pi(A)$ (respectively,
$\pi(\zeta)(A)$). Observe that
$\mathfrak Z\subset\pi(A)'=\mathcal A'$, so $\mathfrak Z$ is contained
in the center of $\mathcal A$.

Let $(x_i)$ be a dense sequence in $A$. Then $\mathcal A(\zeta)$ is the
von Neumann algebra generated by the $\pi(\zeta)(x_i)$. On the other
hand, $\mathcal A$ may be regarded as the von Neumann algebra generated
by the $\pi(x_i)$ and by $\mathfrak Z$. But
\[
\pi(x_i)=\int_Z^\oplus\pi(\zeta)(x_i)\,d\mu(\zeta).
\]
Therefore, by \XRef{A.88}{A 88},
\[
\mathcal A=\int_Z^\oplus\mathcal A(\zeta)\,d\mu(\zeta).
\]
If $\mathfrak Z$ is the center of $\mathcal A$, then
$\mathcal A(\zeta)$ is a factor almost everywhere (\XRef{A.89}{A 89});
thus $\pi(\zeta)$ is a factor representation almost everywhere.

Let
\[
T=\int_Z^\oplus T(\zeta)\,d\mu(\zeta)\in\mathcal A.
\]
Since $T$ is a strong limit of elements of $\pi(A)$, there are a null
subset $N(T)$ of $Z$ and a sequence $(x_1,x_2,\ldots)$ in $A$ such that,
for $\zeta\in Z\setminus N(T)$, $\pi(\zeta)(x_i)$ tends strongly to
$T(\zeta)$ (\XRef{A.13}{A 13}, \XRef{A.79}{A 79}). Let
$\zeta,\zeta'\in Z\setminus N(T)$, and let $R$ be an intertwining operator
for $\pi(\zeta)$ and $\pi(\zeta')$. The equality
\[
R\pi(\zeta)(x_i)=\pi(\zeta')(x_i)R,
\]
valid for every $i$, implies
\[
RT(\zeta)=T(\zeta')R.
\]

This established, suppose that $Z$ is standard, and let
$(Y_1,Y_2,\ldots)$ be a sequence of Borel subsets of $Z$ separating the
points of $Z$. Let $\varphi_i$ be the characteristic function of $Y_i$;
it defines an element $E_i$ of $\mathfrak Z$. The union $N$ of the
$N(E_i)$ is null. Let $\zeta,\zeta'\in Z\setminus N$ with
$\zeta\ne\zeta'$, and let $R$ be an intertwining operator for
$\pi(\zeta)$ and $\pi(\zeta')$. By what was proved above, for every $i$,
\[
RE_i(\zeta)=E_i(\zeta')R,
\qquad\text{that is,}\qquad
\varphi_i(\zeta)R=\varphi_i(\zeta')R.
\]
There is an $i$ such that $\varphi_i(\zeta)=1$ and
$\varphi_i(\zeta')=0$, whence $R=0$. Thus
$\pi(\zeta)\perp\pi(\zeta')$. This proves (iii).

\ResultHeading{8.4.2}{8.4.2. Theorem.}
\begin{itshape}
Let $H$ be a separable Hilbert space, $\pi$ a nondegenerate
representation of $A$ on $H$, and $\mathfrak Z$ the center of the von
Neumann algebra generated by $\pi(A)$.
\begin{enumerate}[label=\textup{(\roman*)}]
\item There exist a standard measure $\mu$ on $\check A$, a
$\mu$-measurable field $\zeta\mapsto H(\zeta)$ of Hilbert spaces over
$\check A$, a $\mu$-measurable field
$\zeta\mapsto\pi(\zeta)$ of nonzero factor representations of $A$ on the
$H(\zeta)$ such that $\pi(\zeta)\in\zeta$, and an isomorphism $W$ from
$H$ onto
\[
\int_{\check A}^\oplus H(\zeta)\,d\mu(\zeta)
\]
which carries $\mathfrak Z$ to the algebra of diagonalizable operators
and $\pi$ to
\[
\int_{\check A}^\oplus\pi(\zeta)\,d\mu(\zeta).
\]
% Source PDF physical page 159
\item Let $\mu_1$, $\zeta\mapsto H_1(\zeta)$,
$\zeta\mapsto\pi_1(\zeta)$, and $W_1$ have the same properties. Then
$\mu$ and $\mu_1$ are equivalent. After changing the $H_1(\zeta)$ and
the $\pi_1(\zeta)$ on a null set, there exists an isomorphism
$\zeta\mapsto V(\zeta)$ from the field $\zeta\mapsto H(\zeta)$ onto the
field $\zeta\mapsto H_1(\zeta)$ such that $V(\zeta)$ carries
$\pi(\zeta)$ to $\pi_1(\zeta)$ for every $\zeta$.
\end{enumerate}
\end{itshape}

Apply Theorem \XRef{8.3.2}{8.3.2} to $\pi$ and $\mathfrak Z$, retaining
its notation $Z,\mu,\zeta\mapsto H(\zeta),\zeta\mapsto\pi(\zeta)$.
There is an isomorphism from $H$ onto
$\int_Z^\oplus H(\zeta)\,d\mu(\zeta)$ which carries $\mathfrak Z$ to the
algebra of diagonalizable operators and $\pi$ to
$\int_Z^\oplus\pi(\zeta)\,d\mu(\zeta)$. Since $\pi$ is nondegenerate,
$\mathfrak Z$ is contained in the strong closure of $\pi(A)$. By
\XRef{8.4.1}{8.4.1}, almost all the $\pi(\zeta)$ are factor
representations and are pairwise disjoint. By \XRef{8.1.5}{8.1.5},
almost all the $\pi(\zeta)$ are nonzero. After changing $H(\zeta)$ and
$\pi(\zeta)$ on a null set, we may therefore suppose that the
$\pi(\zeta)$ are nonzero, pairwise disjoint factor representations.
Next, by isomorphisms on the $H(\zeta)$, we may suppose that
$\zeta\mapsto\pi(\zeta)$ is a Borel map from $Z$ into
$\operatorname{Fac}(A)$ (\XRef{8.1.8}{8.1.8}). Since the
$\pi(\zeta)$ are pairwise disjoint, this map is injective; it is therefore
(\XRef{B.21}{B 21}) a Borel isomorphism $\theta$ from $Z$ onto a Borel
subset $B$ of $\operatorname{Fac}(A)$ which meets each quasi-equivalence
class in at most one point. Let $\psi$ be the restriction to $B$ of the
canonical map $\operatorname{Fac}(A)\to\check A$; it is a Borel
isomorphism from $B$ onto a Borel subset $C$ of $\check A$
(\XRef{7.2.3}{7.2.3}). Thus $\psi\circ\theta$ is a Borel isomorphism from
$Z$ onto $C$. We may therefore replace $Z$ by $C$ and suppose that the
quasi-equivalence class of $\pi(\zeta)$ is $\zeta$ for every
$\zeta\in C$. This proves (i), on extending $\mu$ to $\check A$ by
$\mu(\check A\setminus C)=0$.

Introduce the notation $\mu_1$, $\zeta\mapsto H_1(\zeta)$,
$\zeta\mapsto\pi_1(\zeta)$, and $W_1$ of (ii). Apply
\XRef{8.2.4}{8.2.4}, with $Z=Z_1=\check A$, obtaining
$N,N_1,\widetilde\mu_1,\eta$, and the $V(\zeta)$. For
$\zeta\in Z\setminus N$, $V(\zeta)$ carries $\pi(\zeta)$ to
$\pi_1(\eta(\zeta))$. Hence the quasi-equivalence classes $\zeta$ and
$\eta(\zeta)$ of $\pi(\zeta)$ and $\pi_1(\eta(\zeta))$ are equal. Thus
$\eta$ is the identity map on $Z\setminus N$. Consequently $\mu$ and
$\mu_1$ are equivalent, and (ii) is proved.

\ResultHeading{8.4.3}{8.4.3.}
The class of the measure $\mu$ on $\check A$ is therefore uniquely
determined by $\pi$.

\ResultLabel{Definition.}
\begin{itshape}
It is called the measure class associated with $\pi$.
\end{itshape}

The disintegration
\[
\pi\simeq\int_{\check A}^\oplus\pi(\zeta)\,d\mu(\zeta)
\]
of Theorem \XRef{8.4.2}{8.4.2} is called, with a slight abuse of language,
the central disintegration of $\pi$.

% Source PDF physical page 160
\ResultHeading{8.4.4}{8.4.4. Proposition.}
\begin{itshape}
Let $\pi,\rho$ be two nondegenerate representations of $A$ on separable
Hilbert spaces. Then $\pi$ and $\rho$ are quasi-equivalent if and only if
their associated measure classes on $\check A$ are equal.
\end{itshape}

Suppose that the measure classes on $\check A$ associated with $\pi$ and
$\rho$ are equal. Let $\mathfrak Y$ (respectively, $\mathfrak Z$) be the
center of the von Neumann algebra $\mathcal A$ (respectively,
$\mathcal B$) generated by $\pi(A)$ (respectively, $\rho(A)$). There are
a standard measure $\mu$ on $\check A$, a measurable field
$\zeta\mapsto H(\zeta)$ (respectively, $\zeta\mapsto K(\zeta)$) of Hilbert
spaces over $\check A$, and a measurable field
$\zeta\mapsto\pi(\zeta)$ (respectively, $\zeta\mapsto\rho(\zeta)$) of
factor representations of $A$ on the $H(\zeta)$ (respectively, the
$K(\zeta)$), such that
\[
\pi(\zeta)\in\zeta\quad
  \text{(respectively, }\rho(\zeta)\in\zeta\text{)},
\]
\[
\pi=\int_{\check A}^\oplus\pi(\zeta)\,d\mu(\zeta)
\quad
\left[
\text{respectively, }\rho=\int_{\check A}^\oplus
\rho(\zeta)\,d\mu(\zeta)
\right],
\]
and $\mathfrak Y$ (respectively, $\mathfrak Z$) is the algebra of
diagonalizable operators (all this up to isomorphisms). Since
$\pi(\zeta)\approx\rho(\zeta)$ for every $\zeta$, there is an isomorphism
$\Phi(\zeta)$ from the von Neumann algebra $\mathcal A(\zeta)$ generated
by $\pi(\zeta)(A)$ onto the von Neumann algebra $\mathcal B(\zeta)$
generated by $\rho(\zeta)(A)$ such that
\[
\Phi(\zeta)\bigl(\pi(\zeta)(x)\bigr)=\rho(\zeta)(x)
\qquad\text{for every }x\in A.
\]
We have
\[
\mathcal A=\int_{\check A}^\oplus\mathcal A(\zeta)\,d\mu(\zeta),
\qquad
\mathcal B=\int_{\check A}^\oplus\mathcal B(\zeta)\,d\mu(\zeta)
\]
by \XRef{8.4.1}{8.4.1}. Let us show that
$\zeta\mapsto\Phi(\zeta)$ is a measurable field of isomorphisms
(\XRef{A.91}{A 91}). Let
\[
T=\int_{\check A}^\oplus T(\zeta)\,d\mu(\zeta)\in\mathcal A.
\]
There is a sequence $(x_i)$ in $A$ such that $\pi(x_i)$ tends strongly to
$T$. On passing to a subsequence, we may suppose that
$\pi(\zeta)(x_i)$ tends strongly to $T(\zeta)$ almost everywhere. Then
\[
\Phi(\zeta)\bigl(\pi(\zeta)(x_i)\bigr)
\longrightarrow\Phi(\zeta)\bigl(T(\zeta)\bigr)
\]
strongly (\XRef{A.1}{A 1}, \XRef{A.27}{A 27}). Since the fields
\[
\zeta\longmapsto
\Phi(\zeta)\bigl(\pi(\zeta)(x_i)\bigr)=\rho(\zeta)(x_i)
\]
are measurable, the field
$\zeta\mapsto\Phi(\zeta)(T(\zeta))$ is measurable, which proves the
assertion. Hence
\[
\Phi=\int_{\check A}^\oplus\Phi(\zeta)\,d\mu(\zeta)
\]
is an isomorphism from $\mathcal A$ onto $\mathcal B$
(\XRef{A.91}{A 91}) which, for every $x\in A$, carries
\[
\pi(x)=\int_{\check A}^\oplus\pi(\zeta)(x)\,d\mu(\zeta)
\quad\text{to}\quad
\int_{\check A}^\oplus\rho(\zeta)(x)\,d\mu(\zeta)=\rho(x).
\]
Thus $\pi\approx\rho$.

Conversely, suppose that $\pi\approx\rho$. It is clear that if
\[
\pi\simeq\int_{\check A}^\oplus\pi(\zeta)\,d\mu(\zeta)
\]
and $n$ is a cardinal $\leq\aleph_0$, then
\[
n\pi\simeq\int_{\check A}^\oplus n\pi(\zeta)\,d\mu(\zeta).
\]
Thus the measure classes associated with $\pi$ and $n\pi$ are equal.
But $\aleph_0\pi\simeq\aleph_0\rho$ (\XRef{5.3.8}{5.3.8}), so the measure
classes associated with $\pi$ and $\rho$ are equal.

\ResultHeading{8.4.5}{8.4.5. Proposition.}
\begin{itshape}
Let $\pi,\rho$ be two nondegenerate representations of $A$ on separable
Hilbert spaces, and let $C,D$ be the measure classes on $\check A$
associated with $\pi,\rho$. Then $\rho$ is quasi-equivalent to a
subrepresentation of $\pi$ if and only if $D$ is based on $C$.
\end{itshape}

% Source PDF physical page 161
Suppose that $\rho$ is quasi-equivalent to a subrepresentation $\rho'$ of
$\pi$. Let $E\in\pi(A)'$ be the projection corresponding to $\rho'$,
let $\mathfrak Z$ be the center of $\pi(A)'$, let
$E_1\in\mathfrak Z$ be the central support of $E$, and let $\rho_1$ be
the subrepresentation of $\pi$ corresponding to $E_1$. We have
$\rho\approx\rho'\approx\rho_1$ (condition (v) of
\XRef{5.3.1}{5.3.1}), so $D$ is also the measure class associated with
$\rho_1$ (\XRef{8.4.4}{8.4.4}). Now let
\[
\pi=\int_{\check A}^\oplus\pi(\zeta)\,d\mu(\zeta)
\]
be the central disintegration of $\pi$; then $\mu\in C$. Let $Y$ be the
$\mu$-measurable subset of $\check A$ whose characteristic function
$\varphi_Y$ corresponds to $E_1$. We have
\[
\rho_1=\int_{\check A}^\oplus
\pi(\zeta)\,d(\varphi_Y\mu)(\zeta).
\]
Thus $D$ is the class of $\varphi_Y\mu$, and hence $D$ is based on $C$.

Conversely, suppose that $D$ is based on $C$. Choose $\mu\in C$. There is
a $\mu$-measurable subset $Y$ of $\check A$ such that
$\varphi_Y\mu\in D$, where $\varphi_Y$ denotes the characteristic
function of $Y$. Let
\[
\pi=\int_{\check A}^\oplus\pi(\zeta)\,d\mu(\zeta)
\]
be the central disintegration of $\pi$. Let $E_1$ be the diagonalizable
projection corresponding to $\varphi_Y$, and let $\rho_1$ be the
subrepresentation of $\pi$ corresponding to $E_1$. The measure class
associated with $\rho_1$ is the class of $\varphi_Y\mu$, hence is $D$.
Therefore $\rho\approx\rho_1$ (\XRef{8.4.4}{8.4.4}).

\ResultHeading{8.4.6}{8.4.6.}
The second part of the preceding proof shows that, if a measure class $C$
on $\check A$ is associated with a representation of $A$ on a separable
Hilbert space, then every measure class based on $C$ is also associated
with a representation of $A$.

\ResultHeading{8.4.7}{8.4.7. Proposition.}
\begin{itshape}
Let $\pi,\rho,C,D$ be as in \XRef{8.4.5}{8.4.5}. Then $\pi$ and $\rho$
are disjoint if and only if $C$ and $D$ are mutually singular.
\end{itshape}

Indeed, to say that $\pi$ and $\rho$ are not disjoint means that there is
a nonzero representation quasi-equivalent both to a subrepresentation of
$\pi$ and to a subrepresentation of $\rho$. By
\XRef{8.4.5}{8.4.5} and \XRef{8.4.6}{8.4.6}, this is equivalent to the
existence of a nonzero measure class based both on $C$ and on $D$, that
is, to $C$ and $D$ not being mutually singular.

\ResultHeading{8.4.8}{8.4.8. Proposition.}
\begin{itshape}
Let $\pi$ be a nondegenerate representation of $A$ on a separable Hilbert
space, and let
\[
\pi=\int_{\check A}^\oplus\pi(\zeta)\,d\mu(\zeta)
\]
be its central disintegration. Then $\pi$ is of type I if and only if
almost all the $\pi(\zeta)$ are of type I; equivalently, if and only if
$\mu$ is concentrated on $\widehat A$.
\end{itshape}

% Source PDF physical page 162
Let $\mathcal A$ and $\mathcal A(\zeta)$ be the von Neumann algebras
generated by $\pi(A)$ and $\pi(\zeta)(A)$. We have
\[
\mathcal A=\int_{\check A}^\oplus\mathcal A(\zeta)\,d\mu(\zeta)
\]
by \XRef{8.4.1}{8.4.1}(i); consequently $\mathcal A$ is of type I if and
only if almost all the $\mathcal A(\zeta)$ are of type I
(\XRef{A.90}{A 90}).

Bibliography: \cite{ref25,ref56,ref109}.

\BookSubsection{8.5}{Disintegration into Irreducible Representations}

\ResultHeading{8.5.1}{8.5.1. Lemma.}
\begin{itshape}
Let $Z$ be a Borel space, $\mu$ a positive measure on $Z$,
$\zeta\mapsto H(\zeta)$ a measurable field of Hilbert spaces over $Z$,
$\zeta\mapsto\pi(\zeta)$ a measurable field of representations of $A$ on
the $H(\zeta)$,
\[
H=\int_Z^\oplus H(\zeta)\,d\mu(\zeta),\qquad
\pi=\int_Z^\oplus\pi(\zeta)\,d\mu(\zeta),
\]
and $\mathfrak Z$ the algebra of diagonalizable operators. The following
two conditions are equivalent:
\begin{enumerate}[label=\textup{(\roman*)}]
\item $\mathfrak Z$ is a maximal commutative von Neumann subalgebra of
$\pi(A)'$;
\item almost all the $\pi(\zeta)$ are irreducible.
\end{enumerate}
\end{itshape}

Let $\mathcal A$ and $\mathcal A(\zeta)$ be the von Neumann algebras
generated by $\pi(A)$ and $\pi(\zeta)(A)$. Let $(x_i)$ be a dense sequence
in $A$. Then the $\pi(x_i)$ generate the von Neumann algebra
$\mathcal A$, the $\pi(\zeta)(x_i)$ generate the von Neumann algebra
$\mathcal A(\zeta)$, and
\[
\pi(x_i)=\int_Z^\oplus\pi(\zeta)(x_i)\,d\mu(\zeta).
\]
Hence, by \XRef{A.86}{A 86}, $\mathfrak Z$ is a maximal commutative
subalgebra of $\pi(A)'=\mathcal A'$ if and only if
$\mathcal A(\zeta)=\mathcal L(H(\zeta))$ almost everywhere, that is, if
and only if $\pi(\zeta)$ is irreducible almost everywhere.

\ResultHeading{8.5.2}{8.5.2. Theorem.}
\begin{itshape}
Let $H$ be a separable Hilbert space, $\pi$ a representation of $A$ on
$H$, and $\mathfrak Z$ a maximal commutative von Neumann subalgebra of
$\pi(A)'$. There exist a standard Borel space $Z$, a bounded positive
measure $\mu$ on $Z$, a measurable field $\zeta\mapsto H(\zeta)$ of
Hilbert spaces over $Z$, a measurable field
$\zeta\mapsto\pi(\zeta)$ of irreducible representations of $A$ on the
$H(\zeta)$, and an isomorphism from $H$ onto
\[
\int_Z^\oplus H(\zeta)\,d\mu(\zeta)
\]
which carries $\mathfrak Z$ to the algebra of diagonalizable operators
and $\pi$ to
\[
\int_Z^\oplus\pi(\zeta)\,d\mu(\zeta).
\]
\end{itshape}

This follows from \XRef{8.3.2}{8.3.2} and \XRef{8.5.1}{8.5.1}.

\ResultHeading{8.5.3}{8.5.3.}
Recall that, when $\pi$ and $\mathfrak Z$ are given, the disintegration
of $\pi$ provided by Theorem \XRef{8.5.2}{8.5.2} has a certain uniqueness
property (\XRef{8.2.4}{8.2.4}). But when $\pi$ is given and
$\mathfrak Z$ is changed, this disintegration may
% Source PDF physical page 163
change completely, to the point that no irreducible representation occurs
in both disintegrations (cf. \XRef{18.9.8}{18.9.8}). (In this respect,
central disintegration behaves much more satisfactorily.) This
pathological circumstance will, however, be excluded for postliminal
$C^{*}$-algebras.

Bibliography:
\cite{ref52,ref88,ref89,ref91,ref95,ref96,ref101,ref108,ref110,ref137,ref161}.

\BookSubsection[Case of Postliminal C*-Algebras]{8.6}{Case of Postliminal
$C^{*}$-Algebras}

\ResultHeading{8.6.1}{8.6.1.}
Let $A$ be a separable postliminal $C^{*}$-algebra. Then
$\check A=\widehat A$ is a standard Borel space
(\XRef{4.6.1}{4.6.1} and \XRef{7.3.7}{7.3.7}). Let
$\widehat A_n$, for $n=1,2,\ldots,\aleph_0$, be the set of elements of
$\widehat A$ of dimension $n$; this is a Borel subset of $\widehat A$
(\XRef{3.6.3}{3.6.3}(i)). Let $H_n$ be the model Hilbert space of
dimension $n$. There is one and only one Borel field
$\zeta\mapsto H(\zeta)$ of Hilbert spaces over $\widehat A$
(\XRef{A.96}{A 96}) which reduces over $\widehat A_n$ to the constant
field defined by $H_n$. We shall call it the canonical field. For every
positive measure $\mu$ on $\widehat A$, the canonical field
$\zeta\mapsto H(\zeta)$ automatically has the structure of a
$\mu$-measurable field (\XRef{A.97}{A 97}).

\ResultHeading{8.6.2}{8.6.2. Lemma.}
\begin{itshape}
Retain the preceding notation.
\begin{enumerate}[label=\textup{(\roman*)}]
\item There is a field $\zeta\mapsto\pi(\zeta)$,
$\zeta\in\widehat A$, of representations of $A$ on the $H(\zeta)$ such
that $\pi(\zeta)$ belongs to the equivalence class $\zeta$ for every
$\zeta\in\widehat A$, and which is measurable for every positive measure
on $\widehat A$.
\item Let $\zeta\mapsto\pi_1(\zeta)$ be another field of representations
with the properties in (i). For every positive measure $\mu$ on
$\widehat A$,
\[
\int_{\widehat A}^\oplus\pi(\zeta)\,d\mu(\zeta)
\simeq
\int_{\widehat A}^\oplus\pi_1(\zeta)\,d\mu(\zeta).
\]
\end{enumerate}
\end{itshape}

There is a Borel map $\zeta\mapsto\pi(\zeta)$ from $\widehat A$ into
$\operatorname{Irr}(A)$ such that, for every $\zeta\in\widehat A$,
$\pi(\zeta)$ belongs to the class $\zeta$ (\XRef{4.6.2}{4.6.2}). For
every positive measure on $\widehat A$, the field
$\zeta\mapsto\pi(\zeta)$ is measurable (\XRef{8.1.8}{8.1.8}). This proves
(i). Assertion (ii) follows from Proposition \XRef{8.2.3}{8.2.3}, applied
with
\[
Z=Z_1=\widehat A,\qquad
\mu=\mu_1,\qquad
H(\zeta)=H_1(\zeta),\qquad
N=N_1=\varnothing,
\]
and $\eta$ the identity map of $\widehat A$.

\ResultHeading{8.6.3}{8.6.3.}
Lemma \XRef{8.6.2}{8.6.2}(ii) shows that every positive measure $\mu$ on
$\widehat A$ determines a well-defined equivalence class of
representations of $A$ on a separable space (cf. \XRef{A.73}{A 73}).

% Source PDF physical page 164
\ResultLabel{Definition.}
\begin{itshape}
This equivalence class of representations is denoted by
\[
\int_{\widehat A}^\oplus\zeta\,d\mu(\zeta).
\]
\end{itshape}

\ResultHeading{8.6.4}{8.6.4. Proposition.}
\begin{itshape}
Let $A$ be a separable postliminal $C^{*}$-algebra, and let $\mu$ be a
positive measure on $\widehat A$. Let
$\zeta\mapsto\pi(\zeta)$, $\zeta\in\widehat A$, be a
$\mu$-measurable field of representations of $A$ on the spaces $H(\zeta)$
introduced in \XRef{8.6.1}{8.6.1}, such that $\pi(\zeta)$ belongs to the
class $\zeta$ for every $\zeta\in\widehat A$. Put
\[
\pi=\int_{\widehat A}^\oplus\pi(\zeta)\,d\mu(\zeta),
\]
and let $\mathfrak Z$ be the algebra of diagonalizable operators. Then
$\pi(A)'=\mathfrak Z$.
\end{itshape}

For every measurable subset $Y$ of $\widehat A$, denote by $E_Y$ the
diagonalizable projection corresponding to the characteristic function
of $Y$. Every open subset of $\widehat A$ is identified with the spectrum
$\widehat I$ of a closed two-sided ideal $I$ of $A$
(\XRef{3.2.2}{3.2.2}). Let $(x_1,x_2,\ldots)$ be an approximate identity
of the $C^{*}$-algebra $I$ (\XRef{1.7.2}{1.7.2}). If
$\zeta\in\widehat A\setminus\widehat I$, then
$\pi(\zeta)(x_n)=0$ for every $n$. If $\zeta\in\widehat I$, then
$\pi(\zeta)|_I$ is a nondegenerate representation of $I$, so
$\pi(\zeta)(x_n)$ tends strongly to $1$. Hence $\pi(x_n)$ tends strongly
to $E_{\widehat I}$ (\XRef{A.79}{A 79}). Thus $E_{\widehat I}$ belongs
to the von Neumann algebra $\mathcal A$ generated by $\pi(A)$. By
passage to the limit, it follows that $E_Y\in\mathcal A$ for every Borel
subset $Y$ of $\widehat A$. Thus $\mathfrak Z\subset\mathcal A$, and
therefore $\mathfrak Z'\supset\mathcal A'$. On the other hand, by
\XRef{8.5.1}{8.5.1}, $\mathfrak Z$ is a maximal commutative subalgebra of
$\mathcal A'$, that is,
$\mathfrak Z'\cap\mathcal A'\subset\mathfrak Z$. Since
$\mathfrak Z'\supset\mathcal A'$, it follows that
$\mathcal A'\subset\mathfrak Z$. The inclusion
$\mathfrak Z\subset\mathcal A'$ is clear, so
$\mathcal A'=\mathfrak Z$.

\ResultHeading{8.6.5}{8.6.5. Theorem.}
\begin{itshape}
Let $A$ be a separable postliminal $C^{*}$-algebra.
\begin{enumerate}[label=\textup{(\roman*)}]
\item For every positive measure $\mu$ on $\widehat A$,
\[
\int_{\widehat A}^\oplus\zeta\,d\mu(\zeta)
\]
is a multiplicity-free nondegenerate representation, and the measure
class on $\check A=\widehat A$ associated with this representation is the
class of $\mu$.
\item Let $\mu$ and $\mu_1$ be two positive measures on $\widehat A$.
Then
\[
\int_{\widehat A}^\oplus\zeta\,d\mu(\zeta)
\simeq
\int_{\widehat A}^\oplus\zeta\,d\mu_1(\zeta)
\]
if and only if $\mu$ and $\mu_1$ are equivalent.
\item Every multiplicity-free nondegenerate representation of $A$ on a
separable Hilbert space is equivalent to
\[
\int_{\widehat A}^\oplus\zeta\,d\mu(\zeta)
\]
for a bounded positive measure $\mu$ on $\widehat A$.
\end{enumerate}
\end{itshape}

Let
\[
\pi=\int_{\widehat A}^\oplus\zeta\,d\mu(\zeta).
\]
By \XRef{8.6.4}{8.6.4}, $\pi(A)'$ is commutative, so $\pi$ is
multiplicity-free. By \XRef{8.1.5}{8.1.5}, $\pi$ is nondegenerate. Since
$\mathfrak Z$ is the
% Source PDF physical page 165
center of the von Neumann algebra generated by $\pi(A)$
(\XRef{8.6.4}{8.6.4}), the class of $\mu$ is the measure class associated
with $\pi$.

If $\mu$ and $\mu_1$ are equivalent, then
\[
\int_{\widehat A}^\oplus\zeta\,d\mu(\zeta)
\simeq
\int_{\widehat A}^\oplus\zeta\,d\mu_1(\zeta)
\]
by \XRef{8.2.1}{8.2.1}. Conversely, if these representations are
equivalent, then $\mu$ and $\mu_1$ are equivalent by
\XRef{8.4.2}{8.4.2}(ii).

Finally, let $\pi$ be a multiplicity-free nondegenerate representation of
$A$ on a separable Hilbert space. Let $\mu$ be a bounded member of the
measure class associated with $\pi$, and put
\[
\rho=\int_{\widehat A}^\oplus\zeta\,d\mu(\zeta).
\]
By (i), $\rho$ is multiplicity-free, and the measure class associated
with $\rho$ is the class of $\mu$. Hence
$\rho\approx\pi$ (\XRef{8.4.4}{8.4.4}), and therefore
$\rho\simeq\pi$ (\XRef{5.4.6}{5.4.6}).

\ResultHeading{8.6.6}{8.6.6. Theorem.}
\begin{itshape}
Let $A$ be a separable postliminal $C^{*}$-algebra.
\begin{enumerate}[label=\textup{(\roman*)}]
\item Let $\pi$ be a nondegenerate representation of $A$ on a separable
Hilbert space. There exist pairwise mutually singular positive measures
$\mu_1,\mu_2,\ldots,\mu_{\aleph_0}$ on $\widehat A$ such that
\[
\pi\simeq
\int_{\widehat A}^\oplus\zeta\,d\mu_1(\zeta)
\oplus
2\int_{\widehat A}^\oplus\zeta\,d\mu_2(\zeta)
\oplus\cdots\oplus
\aleph_0\int_{\widehat A}^\oplus\zeta\,d\mu_{\aleph_0}(\zeta).
\]
\item If $\mu'_1,\mu'_2,\ldots,\mu'_{\aleph_0}$ have the same
properties, then $\mu_n$ and $\mu'_n$ are equivalent for every $n$.
\item The measure class on $\check A=\widehat A$ associated with $\pi$
is the class of
\[
\mu_1+\mu_2+\cdots+\mu_{\aleph_0}.
\]
\end{enumerate}
\end{itshape}

Since $H_\pi$ is separable,
\[
\pi=\pi_1\oplus\pi_2\oplus\cdots\oplus\pi_{\aleph_0},
\]
where $\pi_n$ has multiplicity $n$ (with the convention
$\infty=\aleph_0$), and the $\pi_n$ are pairwise disjoint
(\XRef{5.4.9}{5.4.9}, \XRef{5.5.2}{5.5.2}). Thus
\[
\pi\simeq\rho_1\oplus2\rho_2\oplus\cdots\oplus
\aleph_0\rho_{\aleph_0},
\]
where the $\rho_n$ are pairwise disjoint and multiplicity-free. By
\XRef{8.6.5}{8.6.5},
\[
\rho_n=\int_{\widehat A}^\oplus\zeta\,d\mu_n(\zeta),
\]
where $\mu_n$ is a positive measure on $\widehat A$. Since the $\rho_n$
are pairwise disjoint, the $\mu_n$ are pairwise mutually singular
(\XRef{8.4.7}{8.4.7}).

Let $\mu'_1,\mu'_2,\ldots,\mu'_{\aleph_0}$ have the properties of (i), and
put
\[
\rho'_n=\int_{\widehat A}^\oplus\zeta\,d\mu'_n(\zeta).
\]
We have
\[
\pi\simeq\rho'_1\oplus2\rho'_2\oplus\cdots\oplus
\aleph_0\rho'_{\aleph_0},
\]
% Source PDF physical page 166
and the $\rho'_n$ are pairwise disjoint (\XRef{8.4.7}{8.4.7}). By
\XRef{5.4.9}{5.4.9}, $n\rho_n\simeq n\rho'_n$. Thus $\mu_n$ and
$\mu'_n$ are equivalent for every $n$.

Let
\[
\mu=\mu_1+\mu_2+\cdots+\mu_{\aleph_0}.
\]
Then
\[
\pi\approx\int_{\widehat A}^\oplus\zeta\,d\mu(\zeta),
\]
so the measure class on $\widehat A$ associated with $\pi$ is the class
of $\mu$.

\ResultHeading{8.6.7}{8.6.7.}
We see that, to know a representation of $A$ up to equivalence, one must
know a measure class and certain \emph{multiplicities}. In central
disintegration, only the measure class mattered. This is because the
representations were studied up to quasi-equivalence and multiplicity
questions therefore no longer arose (recall that a representation is
quasi-equivalent to all its multiples).

\ResultHeading{8.6.8}{8.6.8. Proposition.}
\begin{itshape}
Let $A$ be a separable postliminal $C^{*}$-algebra. Let $\pi$ be a
nondegenerate representation of $A$ on a separable Hilbert space, let
$\mu$ be a member of the measure class associated with $\pi$, and let $S$
be the support of $\mu$ (the smallest closed subset of $\widehat A$ whose
complement is null). Let $\rho\in\widehat A$.
\begin{enumerate}[label=\textup{(\roman*)}]
\item $\rho$ is weakly contained in $\pi$ if and only if $\rho\in S$.
\item $\rho$ is contained in $\pi$ if and only if
$\mu(\{\rho\})>0$.
\end{enumerate}
\end{itshape}

Write
\[
\pi=\int_{\widehat A}^\oplus\pi(\zeta)\,d\mu(\zeta),
\]
where $\pi(\zeta)$ belongs to the class $\zeta$.

Suppose that $\rho\notin S$. Since $S$ is closed, there exists $x\in A$
such that $\rho(x)\ne0$ and $\sigma(x)=0$ for $\sigma\in S$. Then
$\pi(\zeta)(x)=0$ for almost every $\zeta$, and hence $\pi(x)=0$.
Therefore $\rho$ is not weakly contained in $\pi$.

Suppose that $\rho\in S$. Let $x\in A$ satisfy $\pi(x)=0$. Let $U$ be
the open subset of $\widehat A$ consisting of the $\zeta$ such that
$\pi(\zeta)(x)\ne0$. Since $\pi(\zeta)(x)=0$ almost everywhere, $U$ is
null, and is therefore disjoint from $S$. Thus $\rho(x)=0$, and $\rho$ is
weakly contained in $\pi$. This proves (i).

Let $\varepsilon_\rho$ be the Dirac measure at $\rho$. To say that
$\mu(\{\rho\})>0$ means that every $\mu$-null set is
$\varepsilon_\rho$-null; in other words, $\varepsilon_\rho$ is based on
$\mu$, or equivalently $\rho$ is quasi-equivalent to a subrepresentation
$\rho'$ of $\pi$ (\XRef{8.4.5}{8.4.5}). But every representation
quasi-equivalent to an irreducible representation contains it
(\XRef{5.2.1}{5.2.1}, \XRef{5.3.1}{5.3.1}). The preceding condition
therefore also means that $\rho$ is equivalent to a subrepresentation of
$\pi$.

\ResultHeading{8.6.9}{8.6.9.}
Retain the notation of \XRef{8.6.8}{8.6.8}. We see that the support of
$\pi$ (\XRef{3.4.6}{3.4.6}) is precisely the support of $\mu$. We shall
say that a subset $P$ of $\widehat A$ carries $\pi$ if $P$ carries
$\mu$, that is, if $\mu(\widehat A\setminus P)=0$. Naturally $S$ carries
$\pi$, but in general there are strictly
% Source PDF physical page 167
smaller sets than $S$ which carry $\pi$. Nevertheless, if
$\rho\in\widehat A$ is such that $\mu(\{\rho\})>0$, then every set
carrying $\pi$ necessarily contains $\rho$.

Bibliography: \cite{ref31,ref56,ref91}.

\BookSubsection{8.7}{An Interlude}

Since the decomposition of representations into irreducible
representations was one of our principal objectives, we shall pause for a
moment to summarize the main results obtained.

The situation is ideal for finite-dimensional representations: each one
decomposes as a finite sum of irreducible representations, and any two
such decompositions are isomorphic.

Let us pass to infinite-dimensional representations, restricting
ourselves (because little better is known) to separable $C^{*}$-algebras
and representations on separable spaces. The most elementary examples
(even commutative ones) show that the notion of Hilbert-space direct sum
must be replaced by that of direct integral. This corresponds to the
well-known fact that the decomposition of Hermitian operators in
infinite dimension involves \emph{infinitesimal} spectral projections.

If $A$ is a separable postliminal $C^{*}$-algebra (an important case for
applications to group theory), there is a fairly satisfactory theory.
The set $\widehat A$ of irreducible representations is a locally
quasi-compact space with a countable base, which is almost locally compact
(and is actually locally compact if $A$ is a continuous-trace
$C^{*}$-algebra). The Borel structure defined by the topology, which
equals the Mackey structure, is standard. Every representation is
uniquely a sum of representations of multiplicities
$1,2,\ldots,\infty$. Every representation of multiplicity $n$ is the sum
of $n$ equivalent multiplicity-free representations, whose class is
well determined. And every multiplicity-free nondegenerate representation
$\pi$ can be written as an integral
\[
\int_{\widehat A}^\oplus\zeta\,d\mu_\pi(\zeta),
\]
where the class of $\mu_\pi$ is completely determined by $\pi$. The map
\[
\pi\longmapsto\text{the class of }\mu_\pi
\]
defines a bijection from the set of classes of multiplicity-free
nondegenerate representations onto the set of measure classes on
$\widehat A$.

If $A$ is a separable non-postliminal $C^{*}$-algebra, matters are less
satisfactory. The topological space $\widehat A$, the Borel structure
underlying its topology, and the Mackey Borel structure (now distinct
from the preceding one) have pathological properties (cf.
\XRef{9.5.6}{9.5.6}). Every representation can indeed be written as an
integral of irreducible representations,
% Source PDF physical page 168
but uniqueness is completely lost. Moreover, since representations of
types II and III exist (\XRef{9.5.4}{9.5.4}), there can be no question of
reducing to multiplicity-free representations. Introduce the
quasi-spectrum $\check A$, which is not a topological space, but a Borel
space (and which reduces to the Borel space $\widehat A$ when $A$ is
postliminal). This Borel space is at least as pathological as
$\widehat A$, since it contains $\widehat A$ as a Borel subset. But every
representation $\pi$ can be written canonically as an integral over
$\check A$ with respect to a standard measure whose properties reflect
those of $\pi$ fairly well. This defines a bijective correspondence
between the quasi-equivalence classes of representations of $A$ and
certain classes (unfortunately not all) of standard measures on
$\check A$. We repeat that $\pi$ has a canonical decomposition over
$\check A$, but not a unique decomposition, since uniqueness already
fails over $\widehat A$. For example, if $\pi$ is a factor
representation, its canonical decomposition is written $\pi=\pi$; but
$\pi$ can sometimes be written in many ways as an integral of irreducible
representations.

It remains to ask whether a more satisfactory theory might be obtained
by restricting to traceable representations. Only very partial results
are known in this direction.

\BookSubsection{8.8}{Disintegration of a Positive Functional and of a Trace}

\ResultHeading{8.8.1}{8.8.1. Theorem.}
\begin{itshape}
Let $f$ be a positive functional on $A$, let $H$ be the space of
$\pi_f$, let $Z$ be a Borel space, let $\mu$ be a positive measure on
$Z$, let $\zeta\mapsto H(\zeta)$ be a measurable field of Hilbert spaces
over $Z$, and let $\zeta\mapsto\pi(\zeta)$ be a measurable field of
irreducible representations of $A$ on the $H(\zeta)$, such that
\[
H=\int_Z^\oplus H(\zeta)\,d\mu(\zeta),\qquad
\pi_f=\int_Z^\oplus\pi(\zeta)\,d\mu(\zeta).
\]
Let $\zeta\mapsto\xi(\zeta)$ be a vector field such that
\[
\xi_f=\int_Z^\oplus\xi(\zeta)\,d\mu(\zeta).
\]
\begin{enumerate}[label=\textup{(\roman*)}]
\item We have $\xi(\zeta)\ne0$ for almost every $\zeta$; let
$f(\zeta)$ be the pure positive functional, defined for almost every
$\zeta$ by $\pi(\zeta)$ and $\xi(\zeta)$.
\item For every $x\in A$, the function $\zeta\mapsto f(\zeta)(x)$ is
integrable, and
\[
f(x)=\int_Z f(\zeta)(x)\,d\mu(\zeta).
\]
\end{enumerate}
\end{itshape}

Let $Y\subset Z$ be the measurable set of $\zeta\in Z$ such that
$\xi(\zeta)\ne0$, and let $E_Y$ be the corresponding diagonalizable
projection. For every $x\in A$,
\[
E_Y\pi_f(x)\xi_f
=\int_Z^\oplus
\varphi_Y(\zeta)\pi(\zeta)(x)\xi(\zeta)\,d\mu(\zeta),
\]
% Source PDF physical page 169
where $\varphi_Y$ denotes the characteristic function of $Y$. Since
$\varphi_Y(\zeta)=0$ implies $\xi(\zeta)=0$, we see that
$E_Y\pi_f(x)\xi_f=\pi_f(x)\xi_f$. But $\xi_f$ is cyclic for $\pi_f$;
hence $E_Y=1$, and consequently $\xi(\zeta)\ne0$ almost everywhere.

We have
\[
f(\zeta)(x)
=\bigl(\pi(\zeta)(x)\xi(\zeta)\mid\xi(\zeta)\bigr),
\]
so the function $\zeta\mapsto f(\zeta)(x)$ is integrable, and its
integral is
\[
\int_Z\bigl(\pi(\zeta)(x)\xi(\zeta)\mid\xi(\zeta)\bigr)\,d\mu(\zeta)
=\bigl(\pi_f(x)\xi_f\mid\xi_f\bigr)=f(x).
\]

\ResultHeading{8.8.2}{8.8.2. Theorem.}
\begin{itshape}
Let $f$ be a semifinite lower semicontinuous trace on $A^+$, defining the
objects
\[
H_f,\ J_f,\ \lambda_f,\ \rho_f,\ \mathfrak U_f,\ \mathfrak V_f,\ t_f.
\]
\begin{enumerate}[label=\textup{(\roman*)}]
\item There exist:
\begin{enumerate}[label=\textup{\arabic*\textdegree}]
\item a standard positive measure $\mu$ on $\check A$;
\item for every $\zeta\in\check A$, a character $f(\zeta)$;
\item a measurable-field structure of Hilbert spaces on the field
$\zeta\mapsto H_{f(\zeta)}$,
\end{enumerate}
with the following properties:
\begin{enumerate}[label=\textit{\alph*.}]
\item $H_f$ is identified with
\[
\int_{\check A}^\oplus H_{f(\zeta)}\,d\mu(\zeta);
\]
the fields
\[
\begin{gathered}
\zeta\mapsto J_{f(\zeta)},\quad
\zeta\mapsto\lambda_{f(\zeta)},\quad
\zeta\mapsto\rho_{f(\zeta)},\\
\zeta\mapsto\mathfrak U_{f(\zeta)},\quad
\zeta\mapsto\mathfrak V_{f(\zeta)},\quad
\zeta\mapsto t_{f(\zeta)}
\end{gathered}
\]
are measurable, and
\begin{align*}
J_f&=\int_{\check A}^\oplus J_{f(\zeta)}\,d\mu(\zeta),&
\lambda_f&=\int_{\check A}^\oplus\lambda_{f(\zeta)}\,d\mu(\zeta),&
\rho_f&=\int_{\check A}^\oplus\rho_{f(\zeta)}\,d\mu(\zeta),\\
\mathfrak U_f&=\int_{\check A}^\oplus
  \mathfrak U_{f(\zeta)}\,d\mu(\zeta),&
\mathfrak V_f&=\int_{\check A}^\oplus
  \mathfrak V_{f(\zeta)}\,d\mu(\zeta),&
t_f&=\int_{\check A}^\oplus t_{f(\zeta)}\,d\mu(\zeta).
\end{align*}
The algebra $\mathfrak U_f\cap\mathfrak V_f$ is the algebra of
diagonalizable operators.
\item Almost everywhere,
\[
\lambda_{f(\zeta)}\in\zeta
\qquad\text{and}\qquad
\bar\rho^{\mathrm o}_{f(\zeta)}\in\zeta.
\]
\item For every $x\in A^+$, the function
$\zeta\mapsto f(\zeta)(x)$ is measurable and
\[
f(x)=\int_{\check A}f(\zeta)(x)\,d\mu(\zeta).
\]
\end{enumerate}
\item Suppose that a measure $\mu'$, characters $f'(\zeta)$, and a
measurable-field structure of Hilbert spaces on the field
$\zeta\mapsto H_{f'(\zeta)}$ are given, with the same properties as in
(i). Then $\mu$ and $\mu'$ are equivalent. Write $\mu'=d\mu$, where
$d(\zeta)>0$ almost everywhere. We have
\[
f(\zeta)=d(\zeta)f'(\zeta)
\quad\text{almost everywhere}.
\]
After changing the $f'(\zeta)$ on a null set, there is an isomorphism
from the measurable field $\zeta\mapsto H_{f(\zeta)}$ onto the measurable
field $\zeta\mapsto H_{f'(\zeta)}$ which carries
\[
\begin{gathered}
J_{f(\zeta)}\text{ to }J_{f'(\zeta)},\qquad
\lambda_{f(\zeta)}\text{ to }\lambda_{f'(\zeta)},\qquad
\rho_{f(\zeta)}\text{ to }\rho_{f'(\zeta)},\\
\mathfrak U_{f(\zeta)}\text{ to }\mathfrak U_{f'(\zeta)},\qquad
\mathfrak V_{f(\zeta)}\text{ to }\mathfrak V_{f'(\zeta)},\qquad
t_{f(\zeta)}\text{ to }d(\zeta)t_{f'(\zeta)}.
\end{gathered}
\]
\end{enumerate}
\end{itshape}

% Source PDF physical page 170
\ProofOf{Proof of (i).}
Let $B$ be the Hilbert algebra $\mathfrak n_f/N_f$, dense in $H_f$, and
let $B'$ be the corresponding completed Hilbert algebra. There exist
(\XRef{A.84}{A 84}, \XRef{A.94}{A 94}, and \XRef{A.95}{A 95}):
\begin{enumerate}[label=\textup{\arabic*\textdegree}]
\item a standard Borel space $S$;
\item a positive measure $\mu$ on $S$;
\item a measurable field $\zeta\mapsto H(\zeta)$ of nonzero Hilbert
spaces over $S$;
\item a measurable field $\zeta\mapsto B(\zeta)$ of completed Hilbert
algebras in the $H(\zeta)$,
\end{enumerate}
such that $H_f$ is identified with
\[
\int_S^\oplus H(\zeta)\,d\mu(\zeta),
\]
$\mathfrak U_f\cap\mathfrak V_f$ with the algebra of diagonalizable
operators, $B'$ with
\[
\int_S^\oplus B(\zeta)\,d\mu(\zeta),
\]
$\mathfrak U_f$ and $\mathfrak V_f$ with
\[
\int_S^\oplus\mathfrak U(\zeta)\,d\mu(\zeta)
\quad\text{and}\quad
\int_S^\oplus\mathfrak V(\zeta)\,d\mu(\zeta)
\]
respectively, where $\mathfrak U(\zeta)$ and $\mathfrak V(\zeta)$ denote
the left and right von Neumann algebras associated with $B(\zeta)$;
$J_f$ with
\[
\int_S^\oplus J(\zeta)\,d\mu(\zeta),
\]
where $J(\zeta)$ denotes the involution defined by $B(\zeta)$; and $t_f$
with
\[
\int_S^\oplus t(\zeta)\,d\mu(\zeta),
\]
where $t(\zeta)$ denotes the natural trace on
$\mathfrak U(\zeta)^+$ defined by $B(\zeta)$. Since the algebra of
diagonalizable operators is the common center of $\mathfrak U_f$ and
$\mathfrak V_f$, the algebras $\mathfrak U(\zeta)$ and
$\mathfrak V(\zeta)$ are factors for almost every $\zeta$
(\XRef{A.89}{A 89}).

By \XRef{8.3.1}{8.3.1}, there is a measurable field
$\zeta\mapsto\lambda(\zeta)$ (respectively,
$\zeta\mapsto\rho(\zeta)$) of representations of $A$ (respectively,
$A^{\mathrm o}$) on the $H(\zeta)$ such that
\[
\lambda_f=\int_S^\oplus\lambda(\zeta)\,d\mu(\zeta),\qquad
\rho_f=\int_S^\oplus\rho(\zeta)\,d\mu(\zeta).
\]
Let $(x_i)$ be a dense sequence in $A$. Since the $\lambda_f(x_i)$
generate the von Neumann algebra $\mathfrak U_f$, the
$\lambda(\zeta)(x_i)$ generate $\mathfrak U(\zeta)$ for almost every
$\zeta$ (\XRef{A.88}{A 88}). Thus, almost everywhere,
$\mathfrak U(\zeta)$ is the von Neumann algebra generated by
$\lambda(\zeta)(A)$. Similarly, almost everywhere,
$\mathfrak V(\zeta)$ is the von Neumann algebra generated by
$\rho(\zeta)(A)$.

If $x\in A^+$, the function
$\zeta\mapsto t(\zeta)(\lambda(\zeta)(x))$ is measurable, and
(\XRef{A.92}{A 92})
\[
f(x)=t_f(\lambda_f(x))
=\int_S t(\zeta)(\lambda(\zeta)(x))\,d\mu(\zeta)
=\int_S f(\zeta)(x)\,d\mu(\zeta),
\]
where
\[
f(\zeta)(x)=t(\zeta)(\lambda(\zeta)(x))
\qquad\text{for every }x\in A^+.
\]

Let $(y_i)$ be a sequence dense in $\mathfrak n_f$ for its pre-Hilbert
structure (\XRef{6.3.6}{6.3.6}). Since the $\lambda_f(y_i)$ generate the
von Neumann algebra
% Source PDF physical page 171
$\mathfrak U_f$, the $\lambda(\zeta)(y_i)$ generate, for almost every
$\zeta$, the von Neumann algebra $\mathfrak U(\zeta)$. On the other hand,
\[
\int_S f(\zeta)(y_i y_i^{\ast})\,d\mu(\zeta)
=f(y_i y_i^{\ast})<+\infty,
\]
so $f(\zeta)(y_i y_i^{\ast})<+\infty$ almost everywhere. Hence, for
almost every $\zeta$, $(\lambda(\zeta),t(\zeta))$ is a nonzero traced
factor representation of $A$. Consequently almost all the
$f(\zeta)$ are characters (\XRef{6.7.3}{6.7.3}).

Write
\[
\Lambda_fy_i=\int_S^\oplus y_i(\zeta)\,d\mu(\zeta).
\]
For arbitrary $i$ and $j$,
\begin{align*}
\int_S^\oplus\lambda(\zeta)(y_i)y_j(\zeta)\,d\mu(\zeta)
 &=\lambda_f(y_i)\Lambda_fy_j\\
 &=(\Lambda_fy_i)(\Lambda_fy_j)
 =\int_S^\oplus y_i(\zeta)y_j(\zeta)\,d\mu(\zeta).
\end{align*}
Therefore
\[
\lambda(\zeta)(y_i)y_j(\zeta)=y_i(\zeta)y_j(\zeta)
\]
almost everywhere. Thus $\lambda(\zeta)(y_i)$ is almost everywhere the
continuous operator extending left multiplication by $y_i(\zeta)$; an
analogous result holds for $\rho(\zeta)(y_i)$. It follows that, for
almost every $\zeta$ and every $i$,
\begin{align*}
(y_i(\zeta)\mid y_i(\zeta))
 &=t(\zeta)\!\left(
   \lambda(\zeta)(y_i)\lambda(\zeta)(y_i^{\ast})\right)\\
 &=t(\zeta)\!\left(\lambda(\zeta)(y_i y_i^{\ast})\right)
 =f(\zeta)(y_i y_i^{\ast}).
\end{align*}

If the $y_i$ have been chosen to form an involutive algebra over the
field of rational complex numbers, the equality
\[
(y_i(\zeta)\mid y_i(\zeta))=f(\zeta)(y_i y_i^{\ast})
\]
shows that $H(\zeta)$ is identified with a closed vector subspace of
$H_{f(\zeta)}$. We may also require the $y_i$ to be chosen so that every
product $x_i y_j$ is some $y_k$. Then the subspace $H(\zeta)$ of
$H_{f(\zeta)}$ is almost everywhere invariant under
$\lambda_{f(\zeta)}$ and $\rho_{f(\zeta)}$; since $f(\zeta)$ is a
character, $\lambda_{f(\zeta)}(A)$ and $\rho_{f(\zeta)}(A)$
generate the von Neumann algebra $\mathcal L(H_{f(\zeta)})$, and hence
$H(\zeta)=H_{f(\zeta)}$ almost everywhere. We have
\[
J(\zeta)y_i(\zeta)
=(J_f\Lambda_fy_i)(\zeta)
=y_i^{\ast}(\zeta)
=J_{f(\zeta)}y_i(\zeta)
\]
almost everywhere, and therefore
$J(\zeta)=J_{f(\zeta)}$ almost everywhere. Likewise,
\[
\lambda(\zeta)(y_i)y_j(\zeta)
=y_i(\zeta)y_j(\zeta)
=\lambda_{f(\zeta)}(y_i)y_j(\zeta).
\]
Thus $\lambda(\zeta)$ and $\lambda_{f(\zeta)}$ agree on
$\mathfrak n_f$ almost everywhere, and consequently agree on $A$, since
$\lambda(\zeta)$ and $\lambda_{f(\zeta)}$ are factor representations
that are nonzero on $\mathfrak n_f$
(\XRef{2.10.4}{2.10.4}, \XRef{2.11.1}{2.11.1}). Similarly,
$\rho(\zeta)=\rho_{f(\zeta)}$. It follows that
\[
\mathfrak U(\zeta)=\mathfrak U_{f(\zeta)},\qquad
\mathfrak V(\zeta)=\mathfrak V_{f(\zeta)}.
\]

% Source PDF physical page 172
Finally, for every $x\in A^+$,
\[
t(\zeta)(\lambda(\zeta)(x))
=f(\zeta)(x)
=t_{f(\zeta)}(\lambda_{f(\zeta)}(x))
=t_{f(\zeta)}(\lambda(\zeta)(x))
\]
almost everywhere. Writing this for $x=x_1,x=x_2,\ldots$ and applying
\XRef{A.29}{A 29}, we conclude that
$t(\zeta)=t_{f(\zeta)}$ almost everywhere.

After changing the $\lambda(\zeta)$ on a null set, the
$\lambda(\zeta)$ are pairwise disjoint (\XRef{8.4.1}{8.4.1}). Exactly as
in the proof of \XRef{8.4.2}{8.4.2}, we may replace $S$ by a standard
Borel subset of $\check A$ and arrange that
\[
\lambda_{f(\zeta)}=\lambda(\zeta)\in\zeta
\quad\text{almost everywhere}.
\]
Then
\[
\bar\rho^{\mathrm o}_{f(\zeta)}
\approx\lambda_{f(\zeta)}
\quad\Longrightarrow\quad
\bar\rho^{\mathrm o}_{f(\zeta)}\in\zeta,
\]
and all the assertions in (i) are proved.

\ProofOf{Proof of (ii).}
The equivalence of $\mu$ and $\mu'$ follows from
\XRef{8.4.2}{8.4.2}(ii). Write $\mu'=d\mu$, where $d(\zeta)>0$ almost
everywhere. For almost every $\zeta$, $f(\zeta)$ and $f'(\zeta)$ are the
characters of two quasi-equivalent representations, and hence there is
a number $k(\zeta)>0$ such that
\[
f'(\zeta)=k(\zeta)f(\zeta).
\]
For every $x\in A^+$,
\[
\int_{\check A} f(\zeta)(x)\,d\mu(\zeta)
=f(x)
=\int_{\check A} f'(\zeta)(x)\,d\mu'(\zeta)
=\int_{\check A} k(\zeta)f(\zeta)(x)d(\zeta)\,d\mu(\zeta).
\]
Hence, for $x=y_i y_i^{\ast}$,
\begin{equation}
\int_{\check A}(y_i(\zeta)\mid y_i(\zeta))\,d\mu(\zeta)
=\int_{\check A}
(y_i(\zeta)\mid y_i(\zeta))k(\zeta)d(\zeta)\,d\mu(\zeta).
\tag{1}\label{eq:8.8.2-1}
\end{equation}
The equalities
\[
y_1(\zeta)=y_2(\zeta)=\cdots=0
\]
define a $\mu$-null subset of $\check A$. Since the functions
\[
\zeta\longmapsto
(y_i(\zeta)\mid y_i(\zeta))k(\zeta)d(\zeta)
\]
are $\mu$-measurable, we first see that $kd$ is measurable. Since the
linear combinations of the functions
$\zeta\mapsto(y_i(\zeta)\mid y_i(\zeta))$ are dense in
$L^1(\mu)$, equality \eqref{eq:8.8.2-1} then proves that $kd=1$ almost
everywhere. Thus
\[
f(\zeta)=d(\zeta)f'(\zeta)
\quad\text{almost everywhere}.
\]
The map $x\mapsto d(\zeta)^{1/2}x$ on $\mathfrak n_{f(\zeta)}$ then
defines, almost everywhere, an isomorphism $V(\zeta)$ from
$H_{f(\zeta)}$ onto $H_{f'(\zeta)}$. This gives an isomorphism from the
field $\zeta\mapsto H_{f(\zeta)}$ onto the field
$\zeta\mapsto H_{f'(\zeta)}$ which plainly has the properties in (ii).

\ResultHeading{8.8.3}{8.8.3.}
Retain the notation of \XRef{8.8.2}{8.8.2}. Let $(z_i)$ be a sequence of
elements of $\mathfrak m_f^+$. The functions
$\zeta\mapsto f(\zeta)(z_i)$ are $\mu$-integrable, and hence finite almost
everywhere. Thus there exists a null subset $N$ of $\check A$ such that
\[
z_i\in\bigcap_{\zeta\in\check A\setminus N}
\mathfrak m_{f(\zeta)}
\qquad\text{for every }i.
\]

% Source PDF physical page 173
In particular, after changing the $f(\zeta)$ on a null set,
\[
\mathfrak m_f\cap
\left(\bigcap_{\zeta\in\check A}\mathfrak m_{f(\zeta)}\right)
\]
is dense in $\mathfrak m_f$ for the norm of $A$.

\ResultHeading{8.8.4}{8.8.4.}
The measure $\mu$ and the characters $f(\zeta)$ constructed in
\XRef{8.8.2}{8.8.2} have the following additional property. Let
$x\in\mathfrak n_f$, let $h$ be a bounded $\mu$-measurable function on
$\check A$, and let $\Delta_h$ be the corresponding diagonalizable
operator. Then
\[
(\Delta_h\Lambda_fx\mid\Lambda_fx)
=\int_{\check A}h(\zeta)f(\zeta)(x^{\ast}x)\,d\mu(\zeta).
\]
Indeed, write
\[
\Lambda_fx=\int_{\check A}^\oplus x(\zeta)\,d\mu(\zeta).
\]
Then
\[
(\Delta_h\Lambda_fx\mid\Lambda_fx)
=\int_{\check A}h(\zeta)(x(\zeta)\mid x(\zeta))\,d\mu(\zeta).
\]
In the proof of \XRef{8.8.2}{8.8.2}(i), however, we saw that
\[
(x(\zeta)\mid x(\zeta))
=f(\zeta)(x^{\ast}x)
\quad\text{almost everywhere}.
\]

\ResultHeading{8.8.5}{8.8.5.}
Let $A$ be a separable postliminal $C^{*}$-algebra. Let
$\zeta\mapsto K(\zeta)$ be the canonical field of Hilbert spaces over
$\widehat A$ (\XRef{8.6.1}{8.6.1}). Choose a field
$\zeta\mapsto\pi(\zeta)$ of representations of $A$ on the $K(\zeta)$
such that $\pi(\zeta)$ belongs to the class $\zeta$ for every $\zeta$ and
which is measurable for every positive measure on $\widehat A$
(\XRef{8.6.2}{8.6.2}). For simplicity, write $\zeta$ instead of
$\pi(\zeta)$. Let $\overline{K(\zeta)}$ be the Hilbert space conjugate to
$K(\zeta)$. Denote by $\zeta\otimes1$ the representation
\[
x\longmapsto\zeta(x)\otimes1
\]
of $A$ on $K(\zeta)\otimes\overline{K(\zeta)}$, and by
$1\otimes\overline\zeta$ the representation
\[
x\longmapsto1\otimes\overline{\zeta(x)}
\]
of $A$ on the same space. Let $J_\zeta$ be the canonical involution of
$K(\zeta)\otimes\overline{K(\zeta)}$, and let $t_\zeta$ be the trace
\[
T\otimes1\longmapsto\operatorname{Tr}T
\]
on
$\bigl(\mathcal L(K(\zeta))\otimes\mathbb C\bigr)^+$.

\ResultLabel{Theorem.}
\begin{itshape}
Let $f$ be a lower semicontinuous trace on $A^+$ such that
$\mathfrak m_f$ is dense in $A$. There exist a positive measure $\mu$ on
$\widehat A$ and an isomorphism $W$ from $H_f$ onto
\[
\int_{\widehat A}^\oplus
\bigl(K(\zeta)\otimes\overline{K(\zeta)}\bigr)\,d\mu(\zeta)
\]
with the following properties:
\begin{enumerate}[label=\textup{(\roman*)}]
\item $W$ carries
\begin{align*}
J_f&\ \text{to}\ \int_{\widehat A}^\oplus J_\zeta\,d\mu(\zeta),&
\lambda_f&\ \text{to}\ \int_{\widehat A}^\oplus
(\zeta\otimes1)\,d\mu(\zeta),\\
\bar\rho_f^{\mathrm o}&\ \text{to}\ \int_{\widehat A}^\oplus
(1\otimes\overline\zeta)\,d\mu(\zeta),&
\mathfrak U_f&\ \text{to}\ \int_{\widehat A}^\oplus
\bigl(\mathcal L(K(\zeta))\otimes\mathbb C\bigr)\,d\mu(\zeta),\\
\mathfrak V_f&\ \text{to}\ \int_{\widehat A}^\oplus
\bigl(\mathbb C\otimes\mathcal L(\overline{K(\zeta)})\bigr)\,d\mu(\zeta),&
t_f&\ \text{to}\ \int_{\widehat A}^\oplus t_\zeta\,d\mu(\zeta),
\end{align*}
and carries $\mathfrak U_f\cap\mathfrak V_f$ to the algebra of
diagonalizable operators.
% Source PDF physical page 174
\item If $x\in A^+$, the function
$\zeta\mapsto\operatorname{Tr}\zeta(x)$ on $\widehat A$ is lower
semicontinuous, and
\[
f(x)=\int_{\widehat A}\operatorname{Tr}\zeta(x)\,d\mu(\zeta).
\]
\end{enumerate}
\end{itshape}

Introduce all the notation of \XRef{8.8.2}{8.8.2}. Since $A$ is
postliminal, the Borel space $\check A$ is identified with the Borel
space $\widehat A$ (\XRef{7.3.7}{7.3.7}). For $x\in A^+$ and
$\zeta\in\widehat A$, put
\[
f'(\zeta)(x)=\operatorname{Tr}\zeta(x).
\]
By \XRef{8.8.2}{8.8.2}(i)b, for almost every $\zeta$ there is a number
$k(\zeta)>0$ such that
$f'(\zeta)=k(\zeta)f(\zeta)$. For every $x\in A^+$, the function
$\zeta\mapsto f(\zeta)(x)$ is measurable
(\XRef{8.8.2}{8.8.2}(i)c), while
$\zeta\mapsto f'(\zeta)(x)$ is lower semicontinuous
(\XRef{3.5.9}{3.5.9}). For every $\zeta_0\in\widehat A$ there exists
$x\in\mathfrak m_f$ such that $f'(\zeta_0)(xx^{\ast})>0$, and
therefore throughout an open neighborhood $V_{\zeta_0}$ of $\zeta_0$.
Moreover, the function $\zeta\mapsto f(\zeta)(xx^{\ast})$ is integrable
(\XRef{8.8.2}{8.8.2}(i)c), and hence finite almost everywhere. Thus
$\zeta\mapsto k(\zeta)$ is measurable on $V_{\zeta_0}$. Covering
$\widehat A$ by a sequence of sets $V_{\zeta_0}$
(\XRef{3.3.4}{3.3.4}), we see that $k$ is measurable. Replacing $\mu$ by
$k^{-1}\mu$, we may henceforth suppose that
$f'(\zeta)=f(\zeta)$ for every $\zeta$.

There is then an isomorphism from $H_{f(\zeta)}$ onto
$K(\zeta)\otimes\overline{K(\zeta)}$ which carries
\[
\lambda_{f(\zeta)}\text{ to }\zeta\otimes1,\qquad
\bar\rho^{\mathrm o}_{f(\zeta)}\text{ to }1\otimes\overline\zeta,
\]
\[
\mathfrak U_{f(\zeta)}\text{ to }
\mathcal L(K(\zeta))\otimes\mathbb C,\qquad
\mathfrak V_{f(\zeta)}\text{ to }
\mathbb C\otimes\mathcal L(\overline{K(\zeta)}),
\]
and $J_{f(\zeta)}$ to $J_\zeta$, $t_{f(\zeta)}$ to $t_\zeta$
(\XRef{6.7.7}{6.7.7}). In particular,
\[
\dim H_{f(\zeta)}
=\dim\bigl(K(\zeta)\otimes\overline{K(\zeta)}\bigr)
\qquad\text{for every }\zeta.
\]
It follows from \XRef{A.72}{A 72} that the measurable fields
$\zeta\mapsto H_{f(\zeta)}$ and
$\zeta\mapsto K(\zeta)\otimes\overline{K(\zeta)}$ are isomorphic. We may
therefore identify them. For every $\zeta\in\widehat A$, let
$\mathcal E(\zeta)$ be the set of automorphisms of the Hilbert space
$K(\zeta)\otimes\overline{K(\zeta)}$ which carry
$\lambda_{f(\zeta)}$ to $\zeta\otimes1$,
$\bar\rho^{\mathrm o}_{f(\zeta)}$ to $1\otimes\overline\zeta$, and
$J_{f(\zeta)}$ to $J_\zeta$ (and therefore carry
$\mathfrak U_{f(\zeta)}$ to
$\mathcal L(K(\zeta))\otimes\mathbb C$,
$\mathfrak V_{f(\zeta)}$ to
$\mathbb C\otimes\mathcal L(\overline{K(\zeta)})$, and
$t_{f(\zeta)}$ to $t_\zeta$). We have
$\mathcal E(\zeta)\ne\varnothing$. Hence there exists
(\XRef{A.82}{A 82}) a measurable operator field
$\zeta\mapsto W(\zeta)$ such that
$W(\zeta)\in\mathcal E(\zeta)$ for every $\zeta$. The identification of
the fields can therefore be made so that
\[
\lambda_{f(\zeta)}=\zeta\otimes1,\qquad
\bar\rho^{\mathrm o}_{f(\zeta)}=1\otimes\overline\zeta,\qquad
J_{f(\zeta)}=J_\zeta.
\]
The theorem now follows from \XRef{8.8.2}{8.8.2}(i).

\ResultHeading{8.8.6}{8.8.6. Theorem.}
\begin{itshape}
Retain the hypotheses and notation of \XRef{8.8.5}{8.8.5}. Let $\mu'$ be
a positive measure on $\widehat A$ such that, for every $x\in A^+$,
\[
f(x)=\int_{\widehat A}\operatorname{Tr}\zeta(x)\,d\mu'(\zeta).
\]
Then $\mu=\mu'$.
\end{itshape}

% Source PDF physical page 175
Let $G$ be an open subset of $\widehat A$, let $\varphi_G$ be its
characteristic function, let $I$ be the closed two-sided ideal of $A$ such
that $\widehat I=G$, and let $(u_1,u_2,\ldots)$ be an increasing
approximate identity of the $C^{*}$-algebra $I$
(\XRef{1.7.2}{1.7.2}). Let $x\in A^+$. If
$\zeta\in\widehat A\setminus G$, then
\[
\zeta(x^{1/2}u_nx^{1/2})=0
\qquad\text{for every }n;
\]
if $\zeta\in G$, then
\[
\zeta(x^{1/2}u_nx^{1/2})
=\zeta(x)^{1/2}\zeta(u_n)\zeta(x)^{1/2}
\]
increases strongly to $\zeta(x)$, and therefore
$\operatorname{Tr}\zeta(x^{1/2}u_nx^{1/2})$ increases to
$\operatorname{Tr}\zeta(x)$. But
\[
\int_{\widehat A}\operatorname{Tr}\zeta(x^{1/2}u_nx^{1/2})\,d\mu(\zeta)
=f(x^{1/2}u_nx^{1/2})
=\int_{\widehat A}\operatorname{Tr}\zeta(x^{1/2}u_nx^{1/2})\,
d\mu'(\zeta).
\]
Passing to the limit gives
\[
\int_{\widehat A}\varphi_G(\zeta)\operatorname{Tr}\zeta(x)\,d\mu(\zeta)
=\int_{\widehat A}\varphi_G(\zeta)\operatorname{Tr}\zeta(x)\,d\mu'(\zeta).
\]
If $x\in\mathfrak m_f^+$, this formula extends first to the case where
$G$ is a closed subset of $\widehat A$, and then, by passage to the
limit, to the case where $G$ is an arbitrary Borel subset of
$\widehat A$. Thus, for $x\in\mathfrak m_f^+$, the measures
\[
\operatorname{Tr}\zeta(x)\,d\mu(\zeta)
\quad\text{and}\quad
\operatorname{Tr}\zeta(x)\,d\mu'(\zeta)
\]
are equal. Let $(x_1,x_2,\ldots)$ be a dense sequence in
$\mathfrak m_f^+$. The set of $\zeta\in\widehat A$ such that
$\operatorname{Tr}\zeta(x_i)=+\infty$ for at least one $i$ is a Borel
subset of $\widehat A$ which is null for both $\mu$ and $\mu'$. On the
other hand, for every $\zeta_0\in\widehat A$ there is an $i$ such that
$\operatorname{Tr}\zeta_0(x_i)>0$, and hence throughout a
neighborhood of $\zeta_0$. Thus $\mu=\mu'$ in a neighborhood of every
$\zeta_0$, and consequently $\mu=\mu'$.

\ResultHeading{8.8.7}{8.8.7.}
Retain the hypotheses and notation of \XRef{8.8.5}{8.8.5}. Let
$x\in\mathfrak n_f$, let $h$ be a bounded $\mu$-measurable function on
$\widehat A$, and let $\Delta_h$ be the corresponding diagonalizable
operator. By \XRef{8.8.4}{8.8.4},
\[
(\Delta_h\Lambda_fx\mid\Lambda_fx)
=\int_{\widehat A}h(\zeta)
\operatorname{Tr}\bigl(\zeta(x)^{\ast}\zeta(x)\bigr)\,d\mu(\zeta).
\]

Bibliography:
\cite{ref52,ref53,ref57,ref94,ref95,ref96,ref101,ref108,ref110,ref137,ref148,ref161,ref162,ref163,ref164,ref171,ref179,ref186}.

\BookSubsection{8.9}{Complements}

\ResultHeading{8.9.1}{*8.9.1.}
Let $A$ be a separable $C^{*}$-algebra, $Z$ a standard Borel space,
$\mu$ a positive measure on $Z$, $\zeta\mapsto H(\zeta)$ a
$\mu$-measurable field of Hilbert spaces over $Z$, and
$\zeta\mapsto\pi(\zeta)$ a $\mu$-measurable field of representations of
$A$ on the $H(\zeta)$. Put
\[
\pi=\int_Z^\oplus\pi(\zeta)\,d\mu(\zeta),
\]
and let $\mathfrak Z$ be the algebra of diagonalizable operators. Then
the $\pi(\zeta)$ are homogeneous almost everywhere if and only if
$\mathfrak Z$ contains the ideal center $\mathfrak Y$ associated with
$\pi$ (\XRef{5.7.6}{5.7.6}). The $\pi(\zeta)$ are homogeneous and have
pairwise distinct kernels, after deletion of a null set, if and only if
$\mathfrak Z=\mathfrak Y$. \cite{ref20}

% Source PDF physical page 176
\ResultHeading{8.9.2}{8.9.2.}
Let $A$ be a separable $C^{*}$-algebra. If $A$ is not postliminal, there
exist standard measures on $\check A$ whose class is associated with no
representation of $A$ on a separable space. \cite{ref25,ref199}

\ResultHeading{8.9.3}{*8.9.3.}
Theorems \XRef{8.8.5}{8.8.5} and \XRef{8.8.6}{8.8.6} can be generalized
to nonseparable postliminal $C^{*}$-algebras. \cite{ref16}

\ResultHeading{8.9.4}{8.9.4.}
\begin{enumerate}[label=\textit{\alph*.}]
\item Let $A$ be a separable $C^{*}$-algebra. The canonical Borel field
of Hilbert spaces over $\widehat A$ is defined as in
\XRef{8.6.1}{8.6.1}. For every standard positive measure $\mu$ on
$\widehat A$, there is a Borel map
$\zeta\mapsto\pi(\zeta)$ from $\widehat A$ into
$\operatorname{Irr}(A)$ such that $\pi(\zeta)\in\zeta$ almost everywhere.
The class of
\[
\int_{\widehat A}^\oplus\pi(\zeta)\,d\mu(\zeta)
\]
depends only on the class $C$ of $\mu$; denote it by $\mathcal R(C)$.
\item For every class $R$ of type I representations of $A$, denote by
$\mathcal C(R)$ the associated class of standard measures on
$\widehat A$. If $R$ is a class of multiplicity-free representations of
$A$, then
\[
R=\mathcal R(\mathcal C(R)).
\]
\item If $C$ is a class of standard measures on $\widehat A$ and
$\mathcal R(C)$ is of type I, then $\mathcal R(C)$ is multiplicity-free
and
\[
\mathcal C(\mathcal R(C))=C.
\]
\item[\textit{*d.}] If $A$ is not postliminal, there exist two distinct classes
$C,C'$ of standard measures on $\widehat A$ such that
$\mathcal R(C)=\mathcal R(C')$ is a type II factor representation.
\cite{ref87,ref91,ref185,ref199}
\end{enumerate}

Problems: Which classes $C$ of standard measures on $\widehat A$ are
such that $\mathcal R(C)$ is of type I? Which equivalence classes of
representations have the form $\mathcal R(C)$, where $C$ is a class of
standard measures? Is every representation quasi-equivalent to an
element of a class $\mathcal R(C)$, where $C$ is a class of standard
measures?

\ResultHeading{8.9.5}{*8.9.5.}
Let $A$ be a separable $C^{*}$-algebra. Define the canonical Borel field
of Hilbert spaces over $\check A$ to be equal to the canonical field over
$\widehat A$, and over $\check A\setminus\widehat A$ to be the constant
field defined by a Hilbert space of dimension $\aleph_0$. If $\mu$ is a
standard positive measure on $\check A$, there is a Borel map
$\zeta\mapsto\pi(\zeta)$ from $\check A$ into
$\operatorname{Fac}(A)$ such that $\pi(\zeta)\in\zeta$ almost everywhere.
The quasi-equivalence class of
\[
\int_{\check A}^\oplus\pi(\zeta)\,d\mu(\zeta)
\]
depends only on the class $C$ of $\mu$; denote it by $\mathcal R'(C)$.
For every quasi-equivalence class $R$ of representations of $A$, denote
by $\mathcal C'(R)$ the associated class of standard measures on
$\check A$. We have
\[
\mathcal R'(\mathcal C'(R))=R.
\]
If $A$ is not postliminal, there exist classes $C$ of standard measures
on $\check A$ such that
\[
\mathcal C'(\mathcal R'(C))\ne C.
\]
\cite{ref20,ref25,ref91,ref199}

\ResultHeading{8.9.6}{8.9.6.}
Let $A$ be a separable $C^{*}$-algebra, let $\pi$ be a representation of
$A$ on a separable Hilbert space, and let
\[
\pi=\int_Z^\oplus\pi(\zeta)\,d\mu(\zeta)
\]
be a disintegration of $\pi$ into irreducible representations. If $\pi$
is a type I factor representation, then almost all the $\pi(\zeta)$ are
equivalent. \cite{ref90}

% Source PDF physical page 177
\BookSection[Type I C*-Algebras]{9}{Type I $C^{*}$-Algebras}

\BookSubsection{9.1}{Statement of the Theorem. Beginning of the Proof}

\ResultHeading{9.1}{Theorem.}
\begin{itshape}
Let $A$ be a separable $C^{*}$-algebra. The following conditions are
equivalent:
\begin{enumerate}[label=\textup{(\roman*)}]
\item $A$ is of type I;
\item every factor representation of $A$ is of type I;
\item $A$ is postliminal;
\item if $\pi$ is an irreducible representation of $A$, then $\pi(A)$
contains the set of compact operators on $H_\pi$;
\item if $\pi_1$ and $\pi_2$ are two irreducible representations of $A$
with the same kernel, then $\pi_1$ and $\pi_2$ are equivalent.
\end{enumerate}
\end{itshape}

(For the case in which $A$ is not separable, cf.
\XRef{9.5.9}{9.5.9}.)

The scheme of the proof is as follows:
\[
(\mathrm{iii})\Rightarrow(\mathrm{iv})\Rightarrow(\mathrm{v})
\Rightarrow(\mathrm{ii})\Rightarrow(\mathrm{iii})
\Rightarrow(\mathrm{i})\Rightarrow(\mathrm{ii}).
\]
\[
\begin{array}{rcl}
(\mathrm{iii})\Rightarrow(\mathrm{iv})&:&
  \text{cf. \XRef{4.3.7}{4.3.7}},\\
(\mathrm{iv})\Rightarrow(\mathrm{v})&:&
  \text{cf. \XRef{4.1.10}{4.1.10}},\\
(\mathrm{iii})\Rightarrow(\mathrm{i})&:&
  \text{cf. \XRef{5.5.2}{5.5.2}},\\
(\mathrm{i})\Rightarrow(\mathrm{ii})&:&
  \text{clear}.
\end{array}
\]

$(\mathrm{v})\Rightarrow(\mathrm{ii})$: suppose (v) holds. Let $\pi$ be
a factor representation of $A$, and let us prove that it is of type I.
We may suppose that $H_\pi=H$ is separable (\XRef{5.3.7}{5.3.7}). Apply
Theorem \XRef{8.5.2}{8.5.2}, whose notation we adopt. For every
$\mu$-measurable subset $Y$ of $Z$, let $E_Y$ be the corresponding
diagonalizable projection; it commutes with $\pi(A)$, and hence defines
a subrepresentation $\pi_Y$ of $\pi$. If $\mu(Y)\ne0$, then
$E_Y\ne0$, so $\pi_Y\approx\pi$ (\XRef{5.3.5}{5.3.5}), and therefore
\[
\lVert\pi(x)\rVert=\lVert\pi_Y(x)\rVert
=\lVert E_Y\pi(x)\rVert
\qquad\text{for every }x\in A
\]
by condition (ii) of \XRef{5.3.1}{5.3.1}. For
$n=1,2,\ldots$, let $Y_{n,x}$ be the set of $\zeta\in Z$ such that
\[
\lVert\pi(\zeta)(x)\rVert
\leq\lVert\pi(x)\rVert-\frac1n.
\]
We have
\[
\lVert E_{Y_{n,x}}\pi(x)\rVert
\leq\lVert\pi(x)\rVert-\frac1n
\]
(\XRef{A.77}{A 77}), so $Y_{n,x}$ is null by what precedes. Since this
is true for every $n$, it follows that
\[
\lVert\pi(\zeta)(x)\rVert=\lVert\pi(x)\rVert
\]
outside a null set $N_x$. Let $D$ be a countable dense subset of $A$.
Then
\[
N=\bigcup_{x\in D}N_x
\]
is null. If $\zeta\notin N$, then
$\lVert\pi(\zeta)(x)\rVert=\lVert\pi(x)\rVert$ for every $x\in D$, and
hence for every $x\in A$; consequently
$\ker\pi(\zeta)=\ker\pi$. By hypothesis (v), almost all the
$\pi(\zeta)$ are equivalent to one and the same irreducible
% Source PDF physical page 178
representation $\pi_0$. Thus $\pi$ is a multiple of $\pi_0$
(\XRef{8.1.7}{8.1.7}), and is therefore of type I
(\XRef{5.4.11}{5.4.11}).

It remains to prove $(\mathrm{ii})\Rightarrow(\mathrm{iii})$, which will
require fairly extensive preparations.

Bibliography: \cite{ref11,ref28,ref46,ref76,ref77}.

\BookSubsection{9.2}{Preliminaries on Systems of Matrix Units}

We give these preliminaries in greater detail than is needed below, in
order to illuminate somewhat a later, very technical proof.

\ResultHeading{9.2.1}{9.2.1.}
Equip $\mathbb C^p$ with its canonical structure as a Hilbert space of
dimension $p$:
\[
\bigl((z_1,\ldots,z_p)\mid(z'_1,\ldots,z'_p)\bigr)
=z_1\overline{z'_1}+\cdots+z_p\overline{z'_p}.
\]
Then the algebra $M_p(\mathbb C)$ of complex $p$-by-$p$ matrices is
identified with the $C^{*}$-algebra $\mathcal L(\mathbb C^p)$, which is
simple and has dimension $p^2$. Let $(e_{ij})$ be the canonical basis of
$M_p(\mathbb C)$. We have
\[
e_{kl}e_{ij}=\delta_l^i e_{kj},\qquad
e_{ij}^{\ast}=e_{ji},
\]
where $\delta_l^i$ is the Kronecker symbol.

\ResultHeading{9.2.2}{9.2.2.}
More generally, let $A$ be a $C^{*}$-algebra, and let
$(f_{ij})_{1\leq i,j\leq p}$ be a family of elements of $A$. If
\[
f_{kl}f_{ij}=\delta_l^i f_{kj},
\qquad
f_{ij}^{\ast}=f_{ji},
\]
this family is called a system of matrix units of order $p$. The linear
combinations of the $f_{ij}$ then form a $C^{*}$-subalgebra $B$ of $A$.
The linear map from $M_p(\mathbb C)$ onto $B$ which carries $e_{ij}$ to
$f_{ij}$ for $1\leq i,j\leq p$ is a homomorphism; since
$M_p(\mathbb C)$ is simple, this homomorphism is an isomorphism if
$B\ne0$.

Suppose that $A$ is realized as a $C^{*}$-algebra of operators on a
Hilbert space $H$. Since $f_{ii}=f_{ii}^{\ast}=f_{ii}^2$, the $f_{ii}$
are projections. One verifies that $f_{ij}$ is a partial isometry with
initial projection $f_{jj}$ and final projection $f_{ii}$. Put
$v_j=f_{j1}$, $j=1,\ldots,p$. Then
\[
f_{ij}=v_i v_j^{\ast}.
\]
The $v_j$ are partial isometries all having the initial projection
$v_1=f_{11}$; their final projections are pairwise orthogonal, that is,
$v_i^{\ast}v_j=0$ for $i\ne j$. Conversely, let
$v_1,\ldots,v_p$ be partial isometries in $A$ all having $v_1$ as initial
projection and satisfying $v_i^{\ast}v_j=0$ for $i\ne j$. It is readily
verified that, on putting $f_{ij}=v_i v_j^{\ast}$, the $f_{ij}$ form a
system of matrix units of order $p$. Geometrically, there are pairwise
orthogonal closed vector subspaces $H_1,\ldots,H_p$ of $H$ such that
$v_j$ is a partial isometry with initial subspace $H_1$ and final
subspace $H_j$, while $v_1$ is the projection onto $H_1$.

\ResultHeading{9.2.3}{9.2.3.}
Let $a_1,\ldots,a_n$ be variables, each of which can take the values $0$
and $1$. Let
\[
\bigl(v(a_1,\ldots,a_n)\bigr)_{a_1,\ldots,a_n\in\{0,1\}}
\]
be a system of partial isometries in $A$, all having $v(o_n)$ as initial
projection, where $o_n$ denotes the string $(0,\ldots,0)$ of $n$ zeros,
and such that
\[
v(b_1,\ldots,b_n)^{\ast}v(a_1,\ldots,a_n)=0
\quad\text{if}\quad
(b_1,\ldots,b_n)\ne(a_1,\ldots,a_n).
\]
Then the elements
\[
v(b_1,\ldots,b_n)v(a_1,\ldots,a_n)^{\ast}
\]
form a system of matrix units of order $2^n$. If
$H(a_1,\ldots,a_n)$ is the final subspace of
$v(a_1,\ldots,a_n)$, the spaces $H(a_1,\ldots,a_n)$ are pairwise
orthogonal, $v(a_1,\ldots,a_n)$ has $H(o_n)$ as its initial subspace, and
$v(o_n)$ is the projection onto $H(o_n)$.

% Source PDF physical page 179
\ResultHeading{9.2.4}{9.2.4.}
Let
\[
\bigl(v(a_1,\ldots,a_n)\bigr)_{a_1,\ldots,a_n\in\{0,1\}}
\quad\text{and}\quad
\bigl(v(a_1,\ldots,a_{n+1})\bigr)_{a_1,\ldots,a_{n+1}\in\{0,1\}}
\]
be two systems of elements of $A$ having the preceding properties.
Suppose that:
\begin{enumerate}[label=\textup{\arabic*\textdegree}]
\item the projection $v(o_n)$ is the sum of the projection
$v(o_{n+1})$ and the final projection of $v(o_n,1)$;
\item
\[
v(a_1,\ldots,a_{n+1})
=v(a_1,\ldots,a_n)v(o_n,a_{n+1}).
\]
\end{enumerate}
If $H(a_1,\ldots,a_n)$ and $H(a_1,\ldots,a_{n+1})$ are the final
subspaces of $v(a_1,\ldots,a_n)$ and $v(a_1,\ldots,a_{n+1})$, this
implies that $H(a_1,\ldots,a_n)$ is the Hilbert-space direct sum of
$H(a_1,\ldots,a_n,0)$ and $H(a_1,\ldots,a_n,1)$. We shall say that the
system of the $v(a_1,\ldots,a_{n+1})$ is obtained from the system of the
$v(a_1,\ldots,a_n)$ by bisection.

\ResultHeading{9.2.5}{9.2.5.}
For $n=0,1,2,\ldots$, let
\[
\bigl(v(a_1,\ldots,a_n)\bigr)_{a_1,\ldots,a_n\in\{0,1\}}
\]
be a system of elements of $A$ having the properties of
\XRef{9.2.3}{9.2.3}, each system being obtained from the preceding one
by bisection. [For $n=0$, the system reduces to the single element
$v(\varnothing)$.] In particular, the following relations hold:
\begin{enumerate}[label=\textup{(\roman*)}]
\item if $j\leq k$ and
$(a_1,\ldots,a_j)\ne(b_1,\ldots,b_j)$, then
\[
v(a_1,\ldots,a_j)^{\ast}v(b_1,\ldots,b_k)=0;
\]
\item if $k\geq1$, then
\[
v(a_1,\ldots,a_k)
=v(a_1,\ldots,a_{k-1})v(o_{k-1},a_k);
\]
\item if $j<k$, then
\[
v(a_1,\ldots,a_j)^{\ast}v(a_1,\ldots,a_j)
v(o_{k-1},a_k)=v(o_{k-1},a_k);
\]
\item
\[
v(o_k)\geq0.
\]
\end{enumerate}
If $H(a_1,\ldots,a_n)$ is the final subspace of
$v(a_1,\ldots,a_n)$, (i) means that $H(a_1,\ldots,a_j)$ and
$H(b_1,\ldots,b_k)$ are orthogonal whenever $j\leq k$ and
$(a_1,\ldots,a_j)\ne(b_1,\ldots,b_j)$; (iii) means that $H(o_j)$ contains
$H(o_k)$ and $H(o_{k-1},1)$ for $j<k$.

\ResultHeading{9.2.6}{9.2.6.}
In an arbitrary $C^{*}$-algebra one does not know how to construct
systems of nonzero elements having all the properties of
\XRef{9.2.5}{9.2.5}. Under certain fairly broad hypotheses, however
(cf. \XRef{9.3.7}{9.3.7}), we shall succeed in constructing in $A'$
systems of nonzero elements
\[
\bigl(v(a_1,\ldots,a_n)\bigr)_{a_1,\ldots,a_n\in\{0,1\}},
\qquad n=0,1,2,\ldots,
\]
satisfying conditions (i), (ii), (iii), and (iv). We shall now indicate
some consequences of the
% ===== ASSEMBLED TRANSLATION CHUNK 07: chunk_180_209.tex =====
% Source PDF physical page 180
properties (i), (ii), and (iii). We shall suppose that $A$ is realized as a
$C^{*}$-algebra of operators on a Hilbert space $H$. Put
\[
e(n)=\sum_{a_1,\ldots,a_n\in\{0,1\}}
v(a_1,\ldots,a_n)v(a_1,\ldots,a_n)^{*},
\]
which is an element of $A^{+}$. Let
\[
E(n)=\overline{e(n)(H)}=H\ominus\ker e(n).
\]
By (i), for fixed $n$ the subspaces
\[
v(a_1,\ldots,a_n)(H)
\qquad(a_1,\ldots,a_n\in\{0,1\})
\]
are pairwise orthogonal. Consequently $P_{E(n)}$ is the sum of the support
projections of the operators
\[
v(a_1,\ldots,a_n)v(a_1,\ldots,a_n)^{*}.
\]

\ResultHeading{9.2.7}{9.2.7.}
The spaces $E(n)$ form a decreasing sequence. Indeed, if $\xi\in H$ and
$e(n)\xi=0$, then
\[
\bigl(v(a_1,\ldots,a_n)v(a_1,\ldots,a_n)^{*}\xi\mid\xi\bigr)=0
\]
for all $a_1,\ldots,a_n$; hence
$v(a_1,\ldots,a_n)^{*}\xi=0$ for all $a_1,\ldots,a_n$, and therefore, by
(ii), $v(a_1,\ldots,a_{n+1})^{*}\xi=0$ for all
$a_1,\ldots,a_{n+1}$. Thus $e(n+1)\xi=0$.

\ResultHeading{9.2.8}{9.2.8.}
The operators
$v(a_1,\ldots,a_n)v(b_1,\ldots,b_n)^{*}$ leave $E(n+1)$ invariant. Indeed,
let $\eta\in H$ and $\xi=v(c_1,\ldots,c_{n+1})\eta$. Then
\begin{align}
&v(a_1,\ldots,a_n)v(b_1,\ldots,b_n)^{*}\xi \notag\\
&\quad=v(a_1,\ldots,a_n)v(b_1,\ldots,b_n)^{*}
v(c_1,\ldots,c_n)v(o_n,c_{n+1})\eta \notag\\
&\quad=\delta_{b_1}^{c_1}\cdots\delta_{b_n}^{c_n}
v(a_1,\ldots,a_n)v(o_n,c_{n+1})\eta \notag\\
&\quad=\delta_{b_1}^{c_1}\cdots\delta_{b_n}^{c_n}
v(a_1,\ldots,a_n,c_{n+1})\eta,
\tag{1}\label{eq:9.2.8.1}
\end{align}
which belongs to $E(n+1)$ by \XRef{9.2.6}{9.2.6}.

\ResultHeading{9.2.9}{9.2.9.}
The operators $v(a_1,\ldots,a_n)v(b_1,\ldots,b_n)^{*}$ induce on $E(n+1)$
a system of matrix units of order $2^n$ [thus
$e(n)|_{E(n+1)}$ is a projection, and hence
$e(n)|_{E(n+1)}=1$ by \XRef{9.2.7}{9.2.7}]. Indeed, by
\eqref{eq:9.2.8.1},
\begin{align*}
&v(a_1,\ldots,a_n)v(b_1,\ldots,b_n)^{*}
v(a'_1,\ldots,a'_n)v(b'_1,\ldots,b'_n)^{*}\xi\\
&\quad=\delta_{b_1}^{a'_1}\cdots\delta_{b_n}^{a'_n}
v(a_1,\ldots,a_n)v(b'_1,\ldots,b'_n)^{*}
v(c_1,\ldots,c_n)v(o_n,c_{n+1})\eta\\
&\quad=\delta_{b_1}^{a'_1}\cdots\delta_{b_n}^{a'_n}
\delta_{b'_1}^{c_1}\cdots\delta_{b'_n}^{c_n}
v(a_1,\ldots,a_n)v(o_n,c_{n+1})\eta\\
&\quad=\delta_{b_1}^{a'_1}\cdots\delta_{b_n}^{a'_n}
v(a_1,\ldots,a_n)v(b'_1,\ldots,b'_n)^{*}\xi.
\end{align*}
Consequently,
\begin{align*}
&\bigl[v(a_1,\ldots,a_n)v(b_1,\ldots,b_n)^{*}|_{E(n+1)}\bigr]
\bigl[v(a'_1,\ldots,a'_n)v(b'_1,\ldots,b'_n)^{*}|_{E(n+1)}\bigr]\\
&\qquad=\delta_{b_1}^{a'_1}\cdots\delta_{b_n}^{a'_n}
v(a_1,\ldots,a_n)v(b'_1,\ldots,b'_n)^{*}|_{E(n+1)}.
\end{align*}
On the other hand,
\[
v(a_1,\ldots,a_n)v(b_1,\ldots,b_n)^{*}
\quad\hbox{and}\quad
v(b_1,\ldots,b_n)v(a_1,\ldots,a_n)^{*}
\]
% Source PDF physical page 181
are adjoints of one another and both leave $E(n+1)$ invariant; their
restrictions to $E(n+1)$ are therefore adjoints of one another.

\ResultHeading{9.2.10}{9.2.10.}
One has
\begin{align*}
&v(a_1,\ldots,a_{n-1})v(b_1,\ldots,b_{n-1})^{*}|_{E(n+1)}\\
&\quad=\bigl[v(a_1,\ldots,a_{n-1},0)v(b_1,\ldots,b_{n-1},0)^{*}\\
&\hspace{7em}{}+v(a_1,\ldots,a_{n-1},1)v(b_1,\ldots,b_{n-1},1)^{*}\bigr]
|_{E(n+1)}.
\end{align*}
Indeed, with the notation of \XRef{9.2.8}{9.2.8},
\begin{align*}
&\bigl[v(a_1,\ldots,a_{n-1},0)v(b_1,\ldots,b_{n-1},0)^{*}
+v(a_1,\ldots,a_{n-1},1)v(b_1,\ldots,b_{n-1},1)^{*}\bigr]\xi\\
&\quad=\delta_{b_1}^{c_1}\cdots\delta_{b_{n-1}}^{c_{n-1}}
\delta_0^{c_n}v(a_1,\ldots,a_{n-1},0,c_{n+1})\eta\\
&\qquad{}+\delta_{b_1}^{c_1}\cdots\delta_{b_{n-1}}^{c_{n-1}}
\delta_1^{c_n}v(a_1,\ldots,a_{n-1},1,c_{n+1})\eta\\
&\quad=\delta_{b_1}^{c_1}\cdots\delta_{b_{n-1}}^{c_{n-1}}
v(a_1,\ldots,a_{n-1},c_n,c_{n+1})\eta,
\end{align*}
whereas
\begin{align*}
&v(a_1,\ldots,a_{n-1})v(b_1,\ldots,b_{n-1})^{*}\xi\\
&\quad=v(a_1,\ldots,a_{n-1})v(b_1,\ldots,b_{n-1})^{*}
v(c_1,\ldots,c_{n-1})v(o_{n-1},c_n)v(o_n,c_{n+1})\eta\\
&\quad=\delta_{b_1}^{c_1}\cdots\delta_{b_{n-1}}^{c_{n-1}}
v(a_1,\ldots,a_{n-1})v(o_{n-1},c_n)v(o_n,c_{n+1})\eta\\
&\quad=\delta_{b_1}^{c_1}\cdots\delta_{b_{n-1}}^{c_{n-1}}
v(a_1,\ldots,a_{n-1},c_n,c_{n+1})\eta.
\end{align*}

\ResultHeading{9.2.11}{9.2.11.}
The operators $v(a_1,\ldots,a_n)v(b_1,\ldots,b_n)^{*}$ leave $E(r)$
invariant for $r>n$ (which generalizes \XRef{9.2.8}{9.2.8}). We know this
when $r-n=1$. Suppose it has been established when $r-n=s\geq1$, and
consider $r=n+s+1\geq n+2$. By the induction hypothesis, the operators
$v(a_1,\ldots,a_{n+1})v(b_1,\ldots,b_{n+1})^{*}$ leave $E(r)$ invariant.
By \XRef{9.2.10}{9.2.10}, the operators
$v(a_1,\ldots,a_n)v(b_1,\ldots,b_n)^{*}$ act on $E(n+2)$ as linear
combinations of operators
$v(a_1,\ldots,a_{n+1})v(b_1,\ldots,b_{n+1})^{*}$, and hence leave
$E(r)\subset E(n+2)$ invariant.

\ResultHeading{9.2.12}{9.2.12.}
Let $M(n)$ be the set of linear combinations of the operators
\[
v(a_1,\ldots,a_n)v(b_1,\ldots,b_n)^{*},
\]
and let $N(n)$ be the set of linear combinations of elements of
$M(0)\cup M(1)\cup\cdots\cup M(n)$. Every element of $N(n)$ leaves
$E(n+1)$ invariant (\XRef{9.2.11}{9.2.11}). Let $\Phi_n$ be the map
$x\mapsto x|_{E(n+1)}$, where $x$ ranges over $N(n)$. By
\XRef{9.2.10}{9.2.10},
\[
\Phi_n(N(n))=\Phi_n(M(n)).
\]
Suppose in addition that $v(a_1,\ldots,a_n)\neq0$ for all $n$ and all
$a_1,\ldots,a_n$. Since
\[
v(a_1,\ldots,a_n)v(b_1,\ldots,b_n)^{*}v(b_1,\ldots,b_n,b_{n+1})
=v(a_1,\ldots,a_n,b_{n+1})
\]
[by \eqref{eq:9.2.8.1}, applied with $c_1=b_1,\ldots,c_n=b_n$], it follows
that
\[
v(a_1,\ldots,a_n)v(b_1,\ldots,b_n)^{*}|_{E(n+1)}\neq0.
\]

% Source PDF physical page 182
In view of \XRef{9.2.9}{9.2.9}, there is an isomorphism $\Psi_n$ from the
involutive algebra $\Phi_n(M(n))$ onto $M_{2^n}(\mathbb C)$.

For $x\in N(n)$, put
\[
f_n(x)=2^{-n}\operatorname{Tr}\bigl(\Psi_n(\Phi_n(x))\bigr).
\]
Then $f_n$ is a linear form on $N(n)$, and
\begin{align*}
f_n(x^{*})&=\overline{f_n(x)},&
f_n(x)&\geq0\quad(x\geq0),\\
f_n\bigl(v(a_1,\ldots,a_n)v(b_1,\ldots,b_n)^{*}\bigr)
&=2^{-n}\delta_{a_1}^{b_1}\cdots\delta_{a_n}^{b_n}.
\end{align*}
Successive applications of \XRef{9.2.10}{9.2.10} show that
\[
f_n\bigl(v(a_1,\ldots,a_p)v(b_1,\ldots,b_p)^{*}\bigr)
=2^{-p}\delta_{a_1}^{b_1}\cdots\delta_{a_p}^{b_p}
\qquad(p\leq n).
\]
We have $N(0)\subset N(1)\subset\cdots$, and the preceding formulas show
that the forms $f_n$ extend one another. They therefore define a linear form
$f$ on
\[
N=\bigcup_{n\geq0}N(n)
\]
with the following properties:
\begin{enumerate}[label=\textup{\arabic*\textdegree}]
\item $f(x^{*})=\overline{f(x)}$;
\item $f(x)\geq0$ for $x\geq0$;
\item For every $n,a_1,\ldots,a_n,b_1,\ldots,b_n$,
\[
f\bigl(v(a_1,\ldots,a_n)v(b_1,\ldots,b_n)^{*}\bigr)
=2^{-n}\delta_{a_1}^{b_1}\cdots\delta_{a_n}^{b_n}.
\]
\end{enumerate}

\ResultHeading{9.2.13}{9.2.13.}
For $n=0,1,\ldots,r$, let
$(v(a_1,\ldots,a_n))_{a_1,\ldots,a_n\in\{0,1\}}$ be a system of elements
of $A$ satisfying conditions (i), (ii), (iii), and (iv) of
\XRef{9.2.5}{9.2.5}. Choose arbitrary elements $v(o_{r+1})$ and $v(o_r,1)$
of $A$. For $a_1,\ldots,a_{r+1}\in\{0,1\}$ with
$(a_1,\ldots,a_r)\neq o_r$, put
\[
v(a_1,\ldots,a_{r+1})
=v(a_1,\ldots,a_r)v(o_r,a_{r+1}).
\]
This defines a system
$(v(a_1,\ldots,a_{r+1}))_{a_1,\ldots,a_{r+1}\in\{0,1\}}$. Let us determine
whether conditions (i)--(iv) still hold. It is necessary that
\begin{align}
v(a_1,\ldots,a_r)^{*}v(a_1,\ldots,a_r)v(o_r,a_{r+1})
&=v(o_r,a_{r+1}),
\tag{1}\label{eq:9.2.13.1}\\
v(o_r,1)^{*}v(o_{r+1})&=0,
\tag{2}\label{eq:9.2.13.2}\\
v(o_{r+1})&\geq0,
\tag{3}\label{eq:9.2.13.3}
\end{align}
by (iii), (i), and (iv), respectively. We show that
\eqref{eq:9.2.13.1}--\eqref{eq:9.2.13.3} are sufficient. Suppose they hold.
Let $j<k\leq r+1$ and prove (iii). If $k\leq r$, it already holds. If
$k=r+1$ and $j=r$, it follows from \eqref{eq:9.2.13.1}. Suppose that
$k=r+1$ and $j<r$. By \eqref{eq:9.2.13.1} and
\eqref{eq:9.2.13.3},
\[
v(o_r,a_{r+1})=v(o_r)^{*}v(o_r)v(o_r,a_{r+1})
=v(o_r)^2v(o_r,a_{r+1}).
\]
Using (iii) for systems with at most $r$ indices, we obtain
\begin{align*}
&v(a_1,\ldots,a_j)^{*}v(a_1,\ldots,a_j)v(o_r,a_{r+1})\\
&\quad=v(a_1,\ldots,a_j)^{*}v(a_1,\ldots,a_j)v(o_r)^2v(o_r,a_{r+1})\\
&\quad=v(o_r)v(o_r)v(o_r,a_{r+1})=v(o_r,a_{r+1}).
\end{align*}

% Source PDF physical page 183
Let us prove (ii). It suffices to do so when $k=r+1$, that is, to prove
\[
v(a_1,\ldots,a_{r+1})
=v(a_1,\ldots,a_r)v(o_r,a_{r+1}).
\]
This is precisely the definition of $v(a_1,\ldots,a_{r+1})$ when
$(a_1,\ldots,a_r)\neq o_r$. The relation already noted,
\[
v(o_r,a_{r+1})=v(o_r)^2v(o_r,a_{r+1}),
\]
implies
\[
v(o_r,a_{r+1})=v(o_r)v(o_r,a_{r+1}),
\]
because $v(o_r)\geq0$. Thus (ii) holds in every case. Let us prove (i). It
suffices to consider $k=r+1$. If $j<r+1$, then
\begin{align*}
&v(a_1,\ldots,a_j)^{*}v(b_1,\ldots,b_{r+1})\\
&\quad=v(a_1,\ldots,a_j)^{*}v(b_1,\ldots,b_r)v(o_r,b_{r+1}),
\end{align*}
since (ii) is now proved; one need only apply (i) to systems having at most
$r$ indices. Suppose $j=r+1$. By (ii),
\begin{align}
&v(a_1,\ldots,a_{r+1})^{*}v(b_1,\ldots,b_{r+1})\notag\\
&\quad=v(o_r,a_{r+1})^{*}v(a_1,\ldots,a_r)^{*}
v(b_1,\ldots,b_r)v(o_r,b_{r+1}).
\tag{4}\label{eq:9.2.13.4}
\end{align}
This vanishes if $(a_1,\ldots,a_r)\neq(b_1,\ldots,b_r)$. Finally, suppose
\[
(a_1,\ldots,a_r)=(b_1,\ldots,b_r),
\qquad a_{r+1}\neq b_{r+1}.
\]
By \eqref{eq:9.2.13.4} and (iii), which has already been proved,
\[
v(a_1,\ldots,a_{r+1})^{*}v(b_1,\ldots,b_{r+1})
=v(o_r,a_{r+1})^{*}v(o_r,b_{r+1}),
\]
and this vanishes by \eqref{eq:9.2.13.2}.

Bibliography: \cite{ref46}.

\BookSubsection{9.3}{Some Lemmas}

\ResultHeading{9.3.1}{9.3.1. Lemma.}
\begin{itshape}
Let $A$ be a unital $C^{*}$-algebra, let $\pi$ be a nonzero irreducible
representation of $A$, let $H'$ be a finite-dimensional vector subspace of
$H_\pi$, and let $L$ be the set of states
$x\mapsto(\pi(x)\xi\mid\xi)$ of $A$, where $\xi$ ranges over the unit sphere
of $H'$. Let $K$ be the weakly closed convex hull of $L$. Let $K_1$ be the
set of states of $A$ that take the value $1$ at every $x\in A$ such that
$0\leq x\leq1$ and $\pi(x)$ restricts to the identity on $H'$. Then
$K=K_1$.
\end{itshape}

Clearly $K_1$ is convex and weakly closed, and $L\subset K_1$; hence
$K\subset K_1$. Let $l$ be an extreme point of $K_1$. We shall show that
$l\in K$. Together with the Krein--Milman theorem, this will imply
$K_1\subset K$.

% Source PDF physical page 184
Suppose that $l=\tfrac12(l_1+l_2)$, where $l_1$ and $l_2$ are states of
$A$. If $x\in A$, $0\leq x\leq1$, and $\pi(x)$ restricts to the identity on
$H'$, then
\[
l_1(x)\leq1,\qquad l_2(x)\leq1,
\qquad1=l(x)=\tfrac12\bigl(l_1(x)+l_2(x)\bigr).
\]
Thus $l_1(x)=l_2(x)=1$, so $l_1,l_2\in K_1$ and hence $l_1=l_2=l$.
Therefore $l$ is a pure state. Suppose that $\pi_l$ is not equivalent to
$\pi$. There is then a Hermitian element $y\in A$ such that $\pi(y)$
restricts to the identity on $H'$ and $\pi_l(y)\xi_l=0$
(\XRef{2.8.3}{2.8.3}). Let $f:\mathbb R\to\mathbb R$ be the function for
which $f(t)$ is $0$, $t$, or $1$ according as $t<0$, $0\leq t\leq1$, or
$t>1$. Replacing $y$ by $f(y)$, we may suppose $0\leq y\leq1$. But then
\[
l(y)=(\pi_l(y)\xi_l\mid\xi_l)=0,
\]
which proves that $l\notin K_1$, a contradiction. Hence $\pi_l\sim\pi$,
and $l$ is the state defined by $\pi$ and a unit vector $\xi\in H_\pi$.
Write $\xi=\eta+\zeta$, with $\eta\in H'$ and
$\zeta\in H_\pi\ominus H'$. There is a Hermitian element $x\in A$ such
that $\pi(x)$ restricts to the identity on $H'$ and $\pi(x)\zeta=0$
(\XRef{2.8.3}{2.8.3}); arguing as before, we may suppose
$0\leq x\leq1$. Since $l\in K_1$,
\[
1=l(x)=(\pi(x)\xi\mid\xi)=(\eta\mid\xi)=(\eta\mid\eta).
\]
But
\[
1=\|\xi\|^2=\|\eta\|^2+\|\zeta\|^2,
\]
so $\zeta=0$, $\xi=\eta\in H'$, and $l\in K$.

\ResultHeading{9.3.2}{9.3.2. Lemma.}
\begin{itshape}
Retain the notation of Lemma~\XRef{9.3.1}{9.3.1}. Let $E$ be the set of
states of $A$. For every $\delta>0$ and every $x\in A$ such that
$0\leq x\leq1$ and $\pi(x)$ restricts to the identity on $H'$, denote by
$V_{x,\delta}$ the set of $f\in E$ such that $f(x)\geq1-\delta$. Then the
$V_{x,\delta}$ form a fundamental system of neighborhoods of $K$ in $E$
for the weak topology.
\end{itshape}

It is clear that $V_{x,\delta}$ is a closed neighborhood of $K$ in $E$.
Moreover, the intersection of all the $V_{x,\delta}$ is $K$ by
\XRef{9.3.1}{9.3.1}. Since $E$ is compact, it remains only to prove that
the family of the $V_{x,\delta}$ is downward directed. Let
$x_1,x_2\in A$ satisfy $0\leq x_1,x_2\leq1$, with $\pi(x_1)$ and
$\pi(x_2)$ both the identity on $H'$, and let $\delta_1,\delta_2>0$. Put
\[
x=\tfrac12(x_1+x_2)\in A,
\qquad \delta=\inf(\delta_1,\delta_2).
\]
Then $0\leq x\leq1$, and $\pi(x)$ is the identity on $H'$. If $f\in E$
satisfies $f(x)\geq1-\delta/2$, then
\[
f(x_1+x_2)\geq2-\delta,
\qquad f(x_1)\leq1,\qquad f(x_2)\leq1;
\]
hence $f(x_1),f(x_2)\geq1-\delta$, and therefore
\[
V_{x,\delta/2}\subset V_{x_1,\delta_1}\cap V_{x_2,\delta_2}.
\]

% Source PDF physical page 185
\ResultHeading{9.3.3}{9.3.3. Lemma.}
\begin{itshape}
Let $H$ be a Hilbert space, $T\in\mathcal L(H)$, $\xi\in H$, and
$\varepsilon>0$, with
\[
\|T\|\leq1,\qquad \|\xi\|\leq1,\qquad
\operatorname{Re}(T\xi\mid\xi)\geq1-\frac{\varepsilon}{2}.
\]
Then $\|T\xi-\xi\|^2\leq\varepsilon$.
\end{itshape}

Indeed,
\[
\|T\xi-\xi\|^2=(T\xi\mid T\xi)-2\operatorname{Re}(T\xi\mid\xi)
+(\xi\mid\xi)\leq1-2+\varepsilon+1=\varepsilon.
\]

\ResultHeading{9.3.4}{9.3.4. Lemma.}
\begin{itshape}
Let $\varepsilon>0$ and let $n$ be a positive integer. There is a number
$\delta(\varepsilon,n)>0$ with the following property: if $H$ is a Hilbert
space and $T_1,\ldots,T_n\in\mathcal L(H)$ and $\xi\in H$ satisfy
\[
0\leq T_1,\ldots,T_n\leq1,\qquad \|\xi\|\leq1,
\qquad
\operatorname{Re}(T_nT_{n-1}\cdots T_1\xi\mid\xi)
\geq1-\delta(\varepsilon,n),
\]
then
\[
(T_i\xi\mid\xi)\geq1-\varepsilon,
\qquad \|T_i\xi-\xi\|\leq\varepsilon
\quad(i=1,\ldots,n).
\]
\end{itshape}

Put
\[
\delta(\varepsilon,1)=\inf\left(1,\frac{\varepsilon^2}{2}\right),
\qquad
\delta(\varepsilon,n)=\frac{\delta(\varepsilon,n-1)^2}{16}.
\]
Then
\[
\delta(\varepsilon,n)\leq\frac{\delta(\varepsilon,n-1)}2,
\qquad
\delta(\varepsilon,1)\leq\frac{\varepsilon^2}{2},
\qquad
\delta(\varepsilon,1)\leq\varepsilon.
\]
The lemma for $n=1$ follows from \XRef{9.3.3}{9.3.3}. Let $n\geq2$ and
suppose the lemma true for $n-1$ (with the chosen values of $\delta$). If
\[
\operatorname{Re}(T_nT_{n-1}\cdots T_1\xi\mid\xi)
\geq1-\delta(\varepsilon,n),
\]
then
\begin{align*}
1-\delta(\varepsilon,n)
&\leq\operatorname{Re}(T_n\cdots T_1\xi\mid\xi)\\
&\leq
(T_nT_{n-1}\cdots T_1\xi\mid T_{n-1}\cdots T_1\xi)^{1/2}
(T_n\xi\mid\xi)^{1/2}\\
&\leq(T_n\xi\mid\xi)^{1/2}.
\end{align*}
Hence
\[
(T_n\xi\mid\xi)\geq1-2\delta(\varepsilon,n)\geq1-\varepsilon,
\]
and, by \XRef{9.3.3}{9.3.3},
\[
\|T_n\xi-\xi\|^2\leq4\delta(\varepsilon,n)\leq\varepsilon^2.
\]
It follows that
\begin{align*}
\operatorname{Re}(T_{n-1}\cdots T_1\xi\mid\xi)
&=\operatorname{Re}(T_n\cdots T_1\xi\mid\xi)
+\operatorname{Re}(T_{n-1}\cdots T_1\xi\mid(1-T_n)\xi)\\
&\geq\operatorname{Re}(T_n\cdots T_1\xi\mid\xi)-\|\xi-T_n\xi\|\\
&\geq1-\delta(\varepsilon,n)-2\delta(\varepsilon,n)^{1/2}\\
&\geq1-\tfrac12\delta(\varepsilon,n-1)
-\tfrac12\delta(\varepsilon,n-1)\\
&=1-\delta(\varepsilon,n-1).
\end{align*}
The induction hypothesis gives
$(T_i\xi\mid\xi)\geq1-\varepsilon$ and
$\|T_i\xi-\xi\|\leq\varepsilon$ for $i=1,\ldots,n-1$.

\ResultHeading{9.3.5}{9.3.5.}
For every $\varepsilon\in(0,1]$, let $f_\varepsilon$ denote the function
equal to $0$ on $(-\infty,1-\varepsilon]$, equal to $1$ on
$[1-\varepsilon/2,+\infty)$, and linear on
$[1-\varepsilon,1-\varepsilon/2]$. For $\varepsilon\in(0,1/2]$,
\[
f_\varepsilon f_{2\varepsilon}=f_\varepsilon.
\]

% Source PDF physical page 186
\ResultLabel{Lemma.}
\begin{itshape}
Let $A$ be an antiliminary $C^{*}$-algebra, let $d\in A^{+}$ with
$\|d\|=1$, and let $\tau\in(0,1]$. There exist $w,w',d'\in A$ such that:
\begin{enumerate}[label=\textup{(\roman*)}]
\item $\|w\|=\|w'\|=\|d'\|=1$, $w\geq0$, $d'\geq0$, and $w'^{*}w=0$;
\item $f_\tau(d)w=w$ and $f_\tau(d)w'=w'$;
\item $w^2d'=d'$ and $w'^{*}w'd'=d'$.
\end{enumerate}
\end{itshape}

Put $\sigma=\tau/8$. Choose $u\in A$ and $c\in A$ with
$\|u\|\leq1$ and $0\leq c\leq1$, and put
\[
d_0=f_{2\sigma}(d)c f_{2\sigma}(d),
\qquad d_1=f_{4\sigma}(d)-d_0.
\]
We have $0\leq d_0\leq1$, and hence $-1\leq d_1\leq1$. Moreover,
\[
f_{4\sigma}(d)d_0=d_0=d_0f_{4\sigma}(d),
\]
so $d_0$ and $f_{4\sigma}(d)$ generate a commutative $C^{*}$-subalgebra
$B$ of $A$ containing $d_1$. If $\rho$ is any character of $B$ such that
$\rho(d_0)\neq0$, then
\[
\rho(f_{4\sigma}(d))=1,
\qquad \rho(d_1)=1-\rho(d_0);
\]
also $0\leq\rho(d_0)\leq1$. Thus, if $g:\mathbb R\to\mathbb R$ is a
continuous function vanishing on $[0,1/2]$, then
\[
g(\rho(d_0))g(\rho(d_1))=0,
\qquad\text{that is,}\qquad
\rho\bigl(g(d_0)g(d_1)\bigr)=0.
\]
It follows that $g(d_0)g(d_1)=0$.

Put
\[
v=f_\sigma(d_1)u f_\sigma(d_0).
\]
Then
\[
\|v\|\leq1,
\qquad
v^{*}v=f_\sigma(d_0)u^{*}f_\sigma(d_1)^2u f_\sigma(d_0).
\]
Consequently,
\[
f_{2\sigma}(d_0)v^{*}v=v^{*}v,
\qquad v^{*}f_{2\sigma}(d_1)=v^{*},
\]
and hence
\[
v^{*}(v^{*}v)
=v^{*}f_{2\sigma}(d_1)f_{2\sigma}(d_0)v^{*}v=0.
\]
Furthermore,
\[
f_{8\sigma}(d)d_0=d_0,
\qquad f_{8\sigma}(d)d_1=d_1;
\]
therefore $f_{8\sigma}(d)p(d_0)=p(d_0)$ for every polynomial $p$ without
constant term, and hence for every continuous function $p$ vanishing at the
origin; likewise $f_{8\sigma}(d)p(d_1)=p(d_1)$. Thus
\[
f_{8\sigma}(d)v=v,
\qquad f_{8\sigma}(d)v^{*}=v^{*}.
\]
Finally put
\[
d'=f_{1/4}(v^{*}v),
\qquad w=f_{1/2}(v^{*}v)^{1/2},
\qquad w'=v\,k(v^{*}v),
\]
where $k:\mathbb R\to\mathbb R$ is the function equal to
$\bigl(f_{1/2}(t)t^{-1}\bigr)^{1/2}$ for $t\neq0$, and equal to $0$ for
$t=0$.
% Source PDF physical page 187
We have $0\leq d'\leq1$ and $0\leq w\leq1$. Since
$v^{*}(v^{*}v)=0$, we have $w'^{*}w=0$. Since
$f_{8\sigma}(d)v=v$ and $f_{8\sigma}(d)v^{*}=v^{*}$,
\[
f_{8\sigma}(d)w=w,
\qquad f_{8\sigma}(d)w'=w'.
\]
Moreover,
\[
w^2=f_{1/2}(v^{*}v),
\qquad
w'^{*}w'=k^2(v^{*}v)v^{*}v=f_{1/2}(v^{*}v),
\]
so
\[
w^2d'=f_{1/2}(v^{*}v)f_{1/4}(v^{*}v)
=f_{1/4}(v^{*}v)=d',
\]
and similarly $w'^{*}w'd'=d'$. On the other hand,
\[
\|w'\|=\|w'^{*}w'\|^{1/2}\leq1.
\]

We shall prove that $c$ and $u$ can be chosen so that $\|d'\|\geq1$.
Together with (iii), this will imply
\[
\|w\|,\ \|w'\|\geq1,
\qquad\text{and hence}\qquad
\|w\|=\|w'\|=\|d'\|=1,
\]
and the proof will be complete. We have $f_\sigma(d)\neq0$. Since $A$ is
antiliminary, there is an irreducible representation $\pi$ of $A$ such that
$\pi(f_\sigma(d))$ is not compact (\XRef{4.2.6}{4.2.6}). We can therefore
find two orthogonal unit vectors $\xi$ and $\eta$ in the range of
$\pi(f_\sigma(d))$. There is a Hermitian element $c\in A$ such that
$\pi(c)\xi=\xi$ and $\pi(c)\eta=0$ (\XRef{2.8.3}{2.8.3}); arguing as in
\XRef{9.3.1}{9.3.1}, we may suppose $0\leq c\leq1$. There is a unitary
element $u\in A$ such that $\pi(u)\xi=\eta$ (\XRef{2.8.3}{2.8.3}). Since
\[
f_{2\sigma}(d)f_\sigma(d)=f_\sigma(d)
\]
and $\xi,\eta\in\pi(f_\sigma(d))(H_\pi)$, we have
\[
\pi(f_{2\sigma}(d))\xi=\xi,
\qquad \pi(f_{2\sigma}(d))\eta=\eta,
\]
and likewise $\pi(f_{4\sigma}(d))\eta=\eta$. It follows successively that
\begin{gather*}
\pi(d_0)\xi=\xi,
\qquad \pi(d_0)\eta=0,
\qquad \pi(d_1)\eta=\eta,\\
\begin{aligned}
\pi(v^{*}v)\xi
&=f_\sigma(\pi(d_0))\pi(u)^{-1}f_\sigma(\pi(d_1))^2
\pi(u)f_\sigma(\pi(d_0))\xi\\
&=f_\sigma(\pi(d_0))\pi(u)^{-1}f_\sigma(\pi(d_1))^2\pi(u)\xi\\
&=f_\sigma(\pi(d_0))\pi(u)^{-1}f_\sigma(\pi(d_1))^2\eta\\
&=f_\sigma(\pi(d_0))\pi(u)^{-1}\eta
=f_\sigma(\pi(d_0))\xi=\xi,
\end{aligned}\\
\pi(d')\xi=\xi,
\qquad \|\pi(d')\|\geq1,
\qquad \|d'\|\geq1.
\end{gather*}

\ResultHeading{9.3.6}{9.3.6. Lemma.}
\begin{itshape}
Let $A$ be a unital $C^{*}$-algebra, and let $v_1,\ldots,v_p,c$ be
norm-one elements of $A$ such that
$c\geq0$, $v_i^{*}v_i c=c$, and $v_i^{*}v_j=0$ for $i\neq j$. Let $s$ be a
Hermitian element of $A$, and let $\varepsilon>0$. There exist a Hermitian
element $t$ of $A$, which is a linear combination of the $v_iv_j^{*}$, an
element $b\in A$ with $0\leq b\leq1$, and a number $\gamma>0$ such that:
\begin{enumerate}[label=\textup{(\roman*)}]
\item if $f$ is a state of $A$ such that $f(b)\geq1-\gamma$, then
$|f(s-t)|\leq\varepsilon$;
\item
\[
\left\|c\prod_{i=1}^{p}v_i^{*}bv_i\right\|=1.
\]
\end{enumerate}
\end{itshape}

% Source PDF physical page 188
The commutative $C^{*}$-subalgebra of $A$ generated by $1$ and $c$ has a
pure state taking the value $1$ at $c$. This pure state can be extended to a
pure state $m$ of $A$ (\XRef{2.10.1}{2.10.1}). Put
\[
\alpha_{ij}=m(v_i^{*}sv_j),
\qquad
t=\sum_{i,j}\alpha_{ij}v_iv_j^{*}\in A.
\]
We have $\overline{\alpha_{ij}}=\alpha_{ji}$, so $t$ is Hermitian. Moreover,
\[
(\pi_m(c)\xi_m\mid\xi_m)=m(c)=1
\geq\|\pi_m(c)\xi_m\|\,\|\xi_m\|,
\]
and consequently $\pi_m(c)\xi_m=\xi_m$. Hence
\[
\pi_m(v_i^{*}v_i)\xi_m
=\pi_m(v_i^{*}v_i c)\xi_m
=\pi_m(c)\xi_m=\xi_m.
\]
Let $H'$ be the finite-dimensional vector subspace of $H_{\pi_m}$ spanned
by the vectors $\pi_m(v_i)\xi_m$. We have
\begin{align*}
&(\pi_m(t)\pi_m(v_k)\xi_m\mid\pi_m(v_l)\xi_m)\\
&\quad=\sum_{i,j}\alpha_{ij}
(\pi_m(v_j^{*}v_k)\xi_m\mid\pi_m(v_i^{*}v_l)\xi_m)\\
&\quad=\alpha_{lk}(\xi_m\mid\xi_m)=\alpha_{lk}\\
&\quad=(\pi_m(v_l^{*}sv_k)\xi_m\mid\xi_m)
=(\pi_m(s)\pi_m(v_k)\xi_m\mid\pi_m(v_l)\xi_m),
\end{align*}
so
\[
P_{H'}\pi_m(t)P_{H'}=P_{H'}\pi_m(s)P_{H'}.
\]
Let $E$ be the set of states of $A$, and let $W$ be the set of $f\in E$
such that $|f(s-t)|\leq\varepsilon$. Let $K$ be the weakly closed convex
hull of the states of $A$ of the form
$x\mapsto(\pi_m(x)\eta\mid\eta)$, where $\eta$ ranges over the unit sphere
of $H'$. By what precedes, $W$ is a weak neighborhood of $K$ in $E$. By
\XRef{9.3.2}{9.3.2}, there exist $b\in A$, $0\leq b\leq1$, and $\gamma>0$
such that $\pi_m(b)$ restricts to the identity on $H'$ and $W$ contains
every $f\in E$ for which $f(b)\geq1-\gamma$. Furthermore, it is clear that
\[
\left\|c\prod_{i=1}^{p}v_i^{*}bv_i\right\|\leq1.
\]
% Source PDF physical page 189
Finally,
\begin{align*}
&\pi_m(c)\prod_{i=1}^{p}
\pi_m(v_i^{*})\pi_m(b)\pi_m(v_i)\xi_m\\
&\qquad=\pi_m(c)\prod_{i=1}^{p}\pi_m(v_i^{*}v_i)\xi_m
=\pi_m(c)\xi_m=\xi_m,
\end{align*}
and hence
\[
\left\|\pi_m\left(c\prod_{i=1}^{p}v_i^{*}bv_i\right)\right\|\geq1,
\qquad
\left\|c\prod_{i=1}^{p}v_i^{*}bv_i\right\|\geq1.
\]

\ResultHeading{9.3.7}{9.3.7. Lemma.}
\begin{itshape}
Let $A$ be a unital antiliminary $C^{*}$-algebra. Let
$(s_0,s_1,\ldots)$ be a sequence of Hermitian elements of $A$. There exist
nonzero elements $v(a_1,\ldots,a_n)$ in the unit ball of $A$
($a_1,\ldots,a_n\in\{0,1\}$; $n=0,1,2,\ldots$) with the following
properties:
\begin{enumerate}[label=\textup{(\roman*)}]
\item if $j\leq k$ and
$(a_1,\ldots,a_j)\neq(b_1,\ldots,b_j)$, then
$v(a_1,\ldots,a_j)^{*}v(b_1,\ldots,b_k)=0$;
\item if $k\geq1$, then
$v(a_1,\ldots,a_k)=v(a_1,\ldots,a_{k-1})v(o_{k-1},a_k)$;
\item if $j<k$, then
\[
v(a_1,\ldots,a_j)^{*}v(a_1,\ldots,a_j)v(o_{k-1},a_k)
=v(o_{k-1},a_k);
\]
\item $v(\varnothing)=1$ and $v(o_k)\geq0$;
\item if
\[
e(j)=\sum_{a_1,\ldots,a_j\in\{0,1\}}
v(a_1,\ldots,a_j)v(a_1,\ldots,a_j)^{*},
\]
then, for every $j\geq0$, there is a linear combination $t_j$ of the
operators
\[
v(a_1,\ldots,a_j)v(b_1,\ldots,b_j)^{*}
\]
such that
\[
\|e(j+1)(s_j-t_j)e(j+1)\|\leq\frac1{j+1}.
\]
\end{enumerate}
\end{itshape}

We shall construct not only elements $v(a_1,\ldots,a_n)$ having these
properties, but also elements $b(n)\in A^{+}$ such that $\|b(n)\|=1$ and
\begin{equation*}
\textup{(vi)}\qquad
v(a_1,\ldots,a_n)^{*}v(a_1,\ldots,a_n)b(n)=b(n).
\end{equation*}
For $n=0$, put $v(\varnothing)=b(0)=1$. Suppose that the nonzero elements
$v(a_1,\ldots,a_j)$ in the unit ball of $A$ and the norm-one elements
$b(j)\in A^{+}$ have been constructed for $j\leq n$, so that properties
(i)--(vi) hold [property (v), of course, only for $j\leq n-1$].

By \XRef{9.3.6}{9.3.6}, there exist a Hermitian element $t_n$ of $A$ that
is a linear combination of the
$v(a_1,\ldots,a_n)v(b_1,\ldots,b_n)^{*}$, an element $b\in A$ with
$0\leq b\leq1$, and a number $\gamma>0$ such that: first, if $f$ is a state
of $A$ and $f(b)\geq1-\gamma$, then
\[
|f(s_n-t_n)|\leq\frac1{n+1};
\]
second, $\|a\|=1$, where
\[
a=b(n)\prod_{a_1,\ldots,a_n}
v(a_1,\ldots,a_n)^{*}b\,v(a_1,\ldots,a_n),
\]
the elements $v(a_1,\ldots,a_n)$ being arranged in any order.

% Source PDF physical page 190
Apply \XRef{9.3.5}{9.3.5} with $d=aa^{*}$ and with $\tau$, for the moment,
arbitrary in $(0,1]$. Denote by $v(o_{n+1})$, $v(o_n,1)$, and $b(n+1)$ the
elements $w,w',d'$ constructed in \XRef{9.3.5}{9.3.5}. By
\XRef{9.3.5}{9.3.5}(ii),
\begin{align*}
&v(a_1,\ldots,a_n)^{*}v(a_1,\ldots,a_n)v(o_n,a_{n+1})\\
&\quad=v(a_1,\ldots,a_n)^{*}v(a_1,\ldots,a_n)
f_\tau(aa^{*})v(o_n,a_{n+1}).
\end{align*}
But (vi) implies
\[
v(a_1,\ldots,a_n)^{*}v(a_1,\ldots,a_n)aa^{*}=aa^{*},
\]
and hence
\[
v(a_1,\ldots,a_n)^{*}v(a_1,\ldots,a_n)f_\tau(aa^{*})
=f_\tau(aa^{*}).
\]
Thus, using \XRef{9.3.5}{9.3.5}(ii) once more,
\begin{align*}
&v(a_1,\ldots,a_n)^{*}v(a_1,\ldots,a_n)v(o_n,a_{n+1})\\
&\quad=f_\tau(aa^{*})v(o_n,a_{n+1})=v(o_n,a_{n+1}).
\end{align*}
Furthermore, $v(o_n,1)^{*}v(o_{n+1})=0$ by
\XRef{9.3.5}{9.3.5}(i), and $v(o_{n+1})\geq0$. By
\XRef{9.2.13}{9.2.13}, if we put
\[
v(a_1,\ldots,a_{n+1})
=v(a_1,\ldots,a_n)v(o_n,a_{n+1})
\quad\text{when }(a_1,\ldots,a_n)\neq o_n,
\]
conditions (i)--(iv) of \XRef{9.3.7}{9.3.7} still hold. Moreover,
$b(n+1)\in A^{+}$ and $\|b(n+1)\|=1$ by
\XRef{9.3.5}{9.3.5}(i). We have
\[
v(o_n,a_{n+1})^{*}v(o_n,a_{n+1})b(n+1)=b(n+1)
\]
by \XRef{9.3.5}{9.3.5}(iii); and, when
$(a_1,\ldots,a_n)\neq o_n$, condition (iii) of
\XRef{9.3.7}{9.3.7}, which has already been established, gives
\begin{align*}
&v(a_1,\ldots,a_{n+1})^{*}v(a_1,\ldots,a_{n+1})b(n+1)\\
&\quad=v(o_n,a_{n+1})^{*}v(a_1,\ldots,a_n)^{*}v(a_1,\ldots,a_n)
v(o_n,a_{n+1})b(n+1)\\
&\quad=v(o_n,a_{n+1})^{*}v(o_n,a_{n+1})b(n+1)=b(n+1).
\end{align*}
Thus condition (vi) of \XRef{9.3.7}{9.3.7} still holds; in particular,
the $v(a_1,\ldots,a_{n+1})$ are nonzero.

It remains to prove (v) for $j=n$. Suppose that $A$ is realized as a
$C^{*}$-algebra of operators on a Hilbert space $H$. Let $\xi\in H$ and let
\[
\eta=v(a_1,\ldots,a_{n+1})\xi
\]
be a unit vector. By (ii) and (iii),
\begin{align*}
v(a_1,\ldots,a_n)^{*}\eta
&=v(a_1,\ldots,a_n)^{*}v(a_1,\ldots,a_n)
v(o_n,a_{n+1})\xi\\
&=v(o_n,a_{n+1})\xi.
\end{align*}
Since $t\geq f_\tau(t)-\tau/2$ for every $t\geq0$,
\begin{align*}
&\bigl(aa^{*}v(a_1,\ldots,a_n)^{*}\eta
\mid v(a_1,\ldots,a_n)^{*}\eta\bigr)\\
&\quad=\bigl(aa^{*}v(o_n,a_{n+1})\xi
\mid v(o_n,a_{n+1})\xi\bigr)\\
&\quad\geq\bigl(f_\tau(aa^{*})v(o_n,a_{n+1})\xi
\mid v(o_n,a_{n+1})\xi\bigr)-\frac\tau2,
\end{align*}
% Source PDF physical page 191
which, by \XRef{9.3.5}{9.3.5}(ii), is
\begin{align*}
&\bigl(v(o_n,a_{n+1})\xi\mid v(o_n,a_{n+1})\xi\bigr)-\frac\tau2\\
&\quad\geq\|v(a_1,\ldots,a_n)v(o_n,a_{n+1})\xi\|^2-\frac\tau2
=\|\eta\|^2-\frac\tau2=1-\frac\tau2.
\end{align*}
Apply Lemma~\XRef{9.3.4}{9.3.4}, choosing $\tau$ so that
\[
\frac\tau2\leq
\delta\left(\frac{\gamma}{2^{2n+3}},2^{n+1}+2\right).
\]
We obtain
\begin{align*}
&\bigl(v(a_1,\ldots,a_n)^{*}b\,v(a_1,\ldots,a_n)
v(a_1,\ldots,a_n)^{*}\eta
\mid v(a_1,\ldots,a_n)^{*}\eta\bigr)\\
&\qquad\geq1-\frac{\gamma}{2^{2n+3}}.
\end{align*}
Since
\[
v(a_1,\ldots,a_n)v(a_1,\ldots,a_n)^{*}\eta
=v(a_1,\ldots,a_n)v(o_n,a_{n+1})\xi=\eta,
\]
this reads
\[
(b\eta\mid\eta)\geq1-\frac{\gamma}{2^{2n+3}}.
\]
Likewise, if $\eta'$ is a unit vector in the range of
$v(b_1,\ldots,b_{n+1})$, then
\[
(b\eta'\mid\eta')\geq1-\frac{\gamma}{2^{2n+3}}.
\]
If $(a_1,\ldots,a_{n+1})\neq(b_1,\ldots,b_{n+1})$, then $\eta$ and
$\eta'$ are orthogonal by (i), and therefore
\begin{align*}
|(b\eta\mid\eta')|^2
&=|((1-b)\eta\mid\eta')|^2\\
&\leq((1-b)\eta\mid\eta)((1-b)\eta'\mid\eta')
\leq\left(\frac{\gamma}{2^{2n+3}}\right)^2.
\end{align*}
Now let $\sigma\in e(n+1)(H)$ be a unit vector. By the definition of
$e(n+1)$,
\[
\sigma=\sum_{a_1,\ldots,a_{n+1}\in\{0,1\}}
\sigma(a_1,\ldots,a_{n+1}),
\]
where $\sigma(a_1,\ldots,a_{n+1})$ lies in the range of
$v(a_1,\ldots,a_{n+1})$. It follows that
\begin{align*}
(b\sigma\mid\sigma)
&\geq\sum_{a_1,\ldots,a_{n+1}}
\left(1-\frac\gamma2\right)
\|\sigma(a_1,\ldots,a_{n+1})\|^2\\
&\quad-\frac{\gamma}{2^{2n+3}}
\sum_{(a_1,\ldots,a_{n+1})\neq(b_1,\ldots,b_{n+1})}
\|\sigma(a_1,\ldots,a_{n+1})\|
\|\sigma(b_1,\ldots,b_{n+1})\|\\
&\geq1-\frac\gamma2-\frac\gamma2=1-\gamma.
\end{align*}
% Source PDF physical page 192
Consequently,
\[
|((s_n-t_n)\sigma\mid\sigma)|\leq\frac1{n+1}.
\]
Thus
\[
-\frac1{n+1}\leq e(n+1)(s_n-t_n)e(n+1)\leq\frac1{n+1},
\]
and
\[
\|e(n+1)(s_n-t_n)e(n+1)\|\leq\frac1{n+1}.
\]

Bibliography: \cite{ref46}.

\BookSubsection{9.4}{Completion of the Proof of the Theorem}

To prove that (ii) implies (iii) in Theorem~\XRef{9.1}{9.1}, we shall show
that if $A$ is a separable nonpostliminal $C^{*}$-algebra, then $A$ has
factorial representations that are not of type~I. Since $A$ has a nonzero
antiliminary quotient (\XRef{4.3.6}{4.3.6}), this will follow from the more
precise result below.

\ResultHeading{9.4.1}{Proposition.}
\begin{itshape}
Let $A$ be a nonzero separable antiliminary $C^{*}$-algebra. Then $A$ has a
factorial representation of type~II.
\end{itshape}

By \XRef{4.3.9}{4.3.9}, we may suppose that $A$ has an identity. Let
$(s_0,s_1,s_2,\ldots)$ be a dense sequence in the Hermitian part of $A$,
and let the $v(a_1,\ldots,a_n)$ have the properties of
\XRef{9.3.7}{9.3.7}. Let $M(n)$ be the set of linear combinations of the
$v(a_1,\ldots,a_n)v(b_1,\ldots,b_n)^{*}$, let $N(n)$ be the set of linear
combinations of elements of $M(0)\cup M(1)\cup\cdots\cup M(n)$, and put
\[
N=\bigcup_{n\geq0}N(n).
\]
By \XRef{9.2.12}{9.2.12}, there is a linear form $f$ on $N$ such that
$f(x^{*})=\overline{f(x)}$, $f(x)\geq0$ for $x\geq0$, and
\[
f\bigl(v(a_1,\ldots,a_n)v(b_1,\ldots,b_n)^{*}\bigr)
=2^{-n}\delta_{a_1}^{b_1}\cdots\delta_{a_n}^{b_n},
\qquad f(1)=1.
\]
By \XRef{2.10.1}{2.10.1}(i), there is a state $g$ of $A$ extending $f$.
We shall prove that $\pi_g$ is factorial of type~II; that is, the von Neumann
algebra $\mathfrak B$ generated by $\pi_g(A)$ is a factor of type~II.

The operators $\pi_g(v(a_1,\ldots,a_n))$ have properties (i), (ii), and
(iii) of \XRef{9.2.5}{9.2.5}. Put
\[
F(n)=\overline{\pi_g(e(n))(H)}.
\]
By \XRef{9.2.7}{9.2.7}, the spaces $F(n)$ form a decreasing sequence. It is
clear that $F(n)$ is invariant under every operator in $\pi_g(A)'$. Hence
$F=\bigcap_nF(n)$ is invariant under every operator in $\pi_g(A)'$, and so
$P_F\in\pi_g(A)''=\mathfrak B$. Moreover, $F$ is invariant under every
operator in $\pi_g(N)$ by \XRef{9.2.11}{9.2.11}. For every $n$,
\begin{align*}
\|\pi_g(e(n))\xi_g\|\,\|\xi_g\|
&\geq(\pi_g(e(n))\xi_g\mid\xi_g)=g(e(n))=f(e(n))\\
&=2^n2^{-n}=1
\geq\|\pi_g(e(n))\xi_g\|\,\|\xi_g\|,
\end{align*}
% Source PDF physical page 193
and therefore $\pi_g(e(n))\xi_g=\xi_g$, so $\xi_g\in F(n)$. Thus
$\xi_g\in F$. Now $\xi_g$ is cyclic for $\pi_g(A)$ and, a fortiori, for
$\mathfrak B$. Hence $P_F\in\mathfrak B$ has central support $1$, and
$\mathfrak B'$ is isomorphic to $\mathfrak B_F'$ (\XRef{A.20}{A~20}). It
suffices to prove that $\mathfrak B_F'$ is a factor of type~II, and hence
that $\mathfrak B_F$ is a factor of type~II (\XRef{A.52}{A~52}).

By \XRef{9.2.9}{9.2.9} and \XRef{9.2.11}{9.2.11}, the operators
\[
\pi_g\bigl(v(a_1,\ldots,a_n)v(b_1,\ldots,b_n)^{*}\bigr)
\]
induce on $F$ a system of matrix units of order $2^n$. Since
$\pi_g(e(n))\xi_g=\xi_g$, this system is nonzero. There is therefore an
isomorphism $\Psi_n$ from the involutive algebra
\[
\pi_g(M(n))|_F=\pi_g(N(n))|_F
\]
onto $M_{2^n}(\mathbb C)$. First, $\mathfrak B_F$ is infinite-dimensional
as a vector space. We shall prove: first,
\[
(aa'\xi_g\mid\xi_g)=(a'a\xi_g\mid\xi_g)
\qquad(a,a'\in\mathfrak B_F);
\]
second, $\mathfrak B_F$ is a factor. This factor will be finite
(\XRef{A.33}{A~33}), and it cannot be of type $I_q$ for finite $q$
(\XRef{A.49}{A~49}); it will therefore be of type $II_1$, and the proof
will be complete.

Let $s$ be a Hermitian element of $A$. For every positive integer $n$, there
is a $j\geq2n$ such that $\|s-s_j\|\leq1/(2n)$, and then a $t_j\in M(j)$
such that
\[
\|e(j+1)(s_j-t_j)e(j+1)\|\leq\frac1{2n}.
\]
It follows that
\[
\|P_F\pi_g(s-t_j)P_F\|
\leq\|\pi_g(e(j+1)(s-t_j)e(j+1))\|\leq\frac1n,
\]
which shows that $\pi_g(N)P_F$ is norm dense in
$P_F\pi_g(A)P_F$.

If $x\in\pi_g(N(n))$, then
\[
(x\xi_g\mid\xi_g)=2^{-n}\operatorname{Tr}\bigl(\Psi_n(x|_F)\bigr).
\]
It is enough to verify this when
$x=\pi_g(v(a_1,\ldots,a_n)v(b_1,\ldots,b_n)^{*})$, in which case
\[
(x\xi_g\mid\xi_g)
=g\bigl(v(a_1,\ldots,a_n)v(b_1,\ldots,b_n)^{*}\bigr)
=2^{-n}\delta_{a_1}^{b_1}\cdots\delta_{a_n}^{b_n}.
\]
Thus, for $x,x'\in\pi_g(N(n))$,
\begin{equation}
(xx'\xi_g\mid\xi_g)=(x'x\xi_g\mid\xi_g),
\tag{1}\label{eq:9.4.1}
\end{equation}
because
\[
\operatorname{Tr}\bigl(\Psi_n(x|_F)\Psi_n(x'|_F)\bigr)
=\operatorname{Tr}\bigl(\Psi_n(x'|_F)\Psi_n(x|_F)\bigr).
\]
Since \eqref{eq:9.4.1} holds for every $n$, it holds for
$x,x'\in\pi_g(N)$, hence for $x,x'\in P_F\pi_g(A)P_F$ by the norm density
just proved, and finally for $x,x'\in\mathfrak B_F$ because
$P_F\pi_g(A)P_F$ is strongly dense in $P_F\mathfrak BP_F$.

Finally, let $r$ be a projection in the center of $\mathfrak B_F$. Put
$\eta=r\xi_g$ and $\zeta=(1-r)\xi_g$. Since
$a\mapsto(a\xi_g\mid\xi_g)$ is a trace on $\mathfrak B_F$,
% Source PDF physical page 194
it follows at once that $a\mapsto(a\eta\mid\eta)$ and
$a\mapsto(a\zeta\mid\zeta)$ are traces $t_1,t_2$ on $\mathfrak B_F$.
Because $\pi_g(N(n))|_F$ is isomorphic to $M_{2^n}(\mathbb C)$, the
restrictions of $t_1,t_2$ to $\pi_g(N(n))_F$ [the set of operators induced
on $F$ by those in $\pi_g(N(n))$] are proportional. Since $t_1,t_2$ are
weakly continuous, they are proportional on $\mathfrak B_F$. We have
$\eta\neq0$ or $\zeta\neq0$. If, for example, $\eta\neq0$, then
$(r\eta\mid\eta)\neq0$, whereas $(r\zeta\mid\zeta)=0$; hence $t_2=0$.
Therefore
\[
0=t_2(1-r)=((1-r)\zeta\mid\zeta)=\|\zeta\|^2,
\qquad \xi_g=r\xi_g.
\]
Since $\xi_g$ is cyclic for $\mathfrak B_F$, it follows that $r=1$. This
proves that $\mathfrak B_F$ is a factor.

Bibliography: \cite{ref46}.

\BookSubsection{9.5}{Supplements}

\ResultHeading{9.5.1}{9.5.1.}
Let $A$ be a separable $C^{*}$-algebra. If $A/I$ is postliminal for every
primitive ideal $I$ of $A$, then $A$ is postliminal. (Apply
\XRef{9.1}{9.1}.)

\ResultHeading{9.5.2}{9.5.2.}
Let $A$ be a separable $C^{*}$-algebra. Then $A$ is postliminal if and only
if $\widehat A$ is a $T_0$-space. (Apply \XRef{9.1}{9.1}.)

\ResultHeading{9.5.3}{9.5.3.}
Let $A$ be a separable $C^{*}$-algebra. Then $A$ is liminal if and only if
every point of $\widehat A$ is closed. (Apply \XRef{9.5.2}{9.5.2} and
\XRef{4.7.15}{4.7.15}.)

\ResultHeading{9.5.4}{* 9.5.4.}
Let $A$ be a separable $C^{*}$-algebra. The following conditions are
equivalent:
\begin{enumerate}[label=\textup{(\roman*)}]
\item $A$ is antiliminary;
\item $A$ has an injective representation of type~II;
\item $A$ has an injective representation of type~III;
\item there is a family $(\pi_i)$ of irreducible representations of $A$ such
that $\bigoplus_i\pi_i$ is injective and
$\pi_i(A)\cap\mathcal C(H_{\pi_i})=0$;
\item there are families $(\pi_i)_{i\in I}$ and $(\rho_i)_{i\in I}$ of
irreducible representations of $A$ such that $\bigoplus\pi_i$ and
$\bigoplus\rho_i$ are injective, $\ker\pi_i=\ker\rho_i$ for every $i$, and
$\rho_i$ is not equivalent to $\pi_i$ for any $i$.
\end{enumerate}
\cite{ref46}.

\ResultHeading{9.5.5}{* 9.5.5.}
\begin{enumerate}[label=\textup{\alph*.}]
\item Let $G$ be the group with two elements, let $X$ be the compact group
$G\times G\times G\times\cdots$, and let $X'$ be the dense subgroup of
$X$ consisting of the elements $(g_1,g_2,\ldots)$ all but finitely many of
whose components equal $e$. Let $Y$ be the quotient Borel space $X/X'$.
There is no sequence of Borel subsets of $Y$ separating the points of $Y$.
\item Let $A$ be a separable nonpostliminal $C^{*}$-algebra. There is a
Borel injection of $Y$ into $\widehat A$, where $\widehat A$ is endowed
with its Mackey structure. \cite{ref46}
\end{enumerate}

\ResultHeading{9.5.6}{9.5.6.}
Let $A$ be a separable $C^{*}$-algebra. The following conditions are
equivalent:
\begin{enumerate}[label=\textup{(\roman*)}]
\item $A$ is postliminal;
\item for the Borel structure underlying its topology, $\widehat A$ is
standard;
\item for the Mackey structure, there is a sequence of Borel subsets of
$\widehat A$ separating its points;
\item the Mackey structure of $\widehat A$ is identical with the Borel
structure underlying its topology.
\end{enumerate}

% Source PDF physical page 195
Indeed, (i) implies (ii) and (iv) by \XRef{4.6.1}{4.6.1}; (ii) implies (iii)
by \XRef{3.8.3}{3.8.3}; (iv) implies (iii) by \XRef{3.3.4}{3.3.4} and
\XRef{3.8.4}{3.8.4}; and (iii) implies (i) by \XRef{9.5.5}{9.5.5}.
\cite{ref12,ref46}.

\ResultHeading{9.5.7}{* 9.5.7.}
Let $H,H'$ be two infinite-dimensional real vector spaces, and let $B$ be a
bilinear form on $H\times H'$ putting $H$ and $H'$ in duality. Let $K$ be a
Hilbert space. For every $x\in H$ (respectively, $x'\in H'$), let $p(x)$
[respectively, $q(x')$] be a self-adjoint, not necessarily bounded, operator
on $K$. Suppose that the $p(x)$ commute pairwise, that
$p(\lambda_1x_1+\lambda_2x_2)$ is the smallest closed extension of
$\lambda_1p(x_1)+\lambda_2p(x_2)$, and that the analogous properties hold
for $q$. Suppose that
\[
e^{ip(x)}e^{iq(x')}=e^{iB(x,x')}e^{iq(x')}e^{ip(x)}
\qquad(x\in H, x'\in H').
\]
If $M,M'$ are vector subspaces of $H,H'$, let $\mathcal A(M,M')$ be the von
Neumann algebra generated by the $e^{ip(x)}$ and $e^{iq(x')}$
($x\in M$, $x'\in M'$). Let $A$ be the $C^{*}$-algebra generated by the
$\mathcal A(M,M')$ as $M$ (respectively, $M'$) ranges over the
finite-dimensional vector subspaces of $H$ (respectively, $H'$). Then $A$
depends, up to isomorphism, only on $(H,H',B)$. The $C^{*}$-algebra $A$ is
antiliminary, has no closed two-sided ideals other than $A$ and $0$, and
has no traceable representation. If $H$ is a separable real Hilbert space,
$H'=H$, and $B(x,x')=(x\mid x')$, then $A$ has factorial representations
of type $II_\infty$ and factorial representations of type~III. Irreducible
representations of $A$ can be constructed explicitly
\cite{ref40,ref46,ref139}.

This is the physicists' problem of representing the commutation relations.
The representation of the anticommutation relations \cite{ref39} was
translated by Mackey \cite{ref92} into a problem concerning group
representations.

\ResultHeading{9.5.8}{9.5.8.}
Let $A$ be a separable $C^{*}$-algebra. If there is a Borel subset of
$\operatorname{Irr}(A)$ that meets each equivalence class in exactly one
point, then $A$ is postliminal. (Use \XRef{7.2.3}{7.2.3} and
\XRef{9.5.6}{9.5.6}.) \cite{ref91}.

\ResultHeading{9.5.9}{9.5.9.}
Let $A$ be a $C^{*}$-algebra. Conditions (i), (ii), (iii), and (iv) of
\XRef{9.1}{9.1} are equivalent, but it is not known whether they are
equivalent to (v). \cite{ref253,ref254,ref292}.

\BookSection[Continuous Fields of C*-Algebras]{10}{Continuous Fields of
$C^{*}$-Algebras}

Let $A$ be a $C^{*}$-algebra, $T$ a topological space, and $B$ the set of
continuous maps $f:T\to A$ such that
$\sup_{t\in T}\|f(t)\|<+\infty$. If $f,g\in B$ and $\lambda\in\mathbb C$,
define $f+g,\lambda f,fg,f^{*}\in B$ and $\|f\|$ by
\begin{align*}
(f+g)(t)&=f(t)+g(t),& (\lambda f)(t)&=\lambda f(t),\\
(fg)(t)&=f(t)g(t),& f^{*}(t)&=f(t)^{*}
\end{align*}
for every $t\in T$, and
\[
\|f\|=\sup_{t\in T}\|f(t)\|.
\]
It is easy to see that $B$ thereby becomes a $C^{*}$-algebra.

% Source PDF physical page 196
If $T$ is discrete, $B$ is simply the product $C^{*}$-algebra
$\prod_{t\in T}A_t$, where $A_t=A$ for every $t$. We have already
introduced the general product $\prod_{t\in T}A_t$ of $C^{*}$-algebras,
where $T$ is a set and $A_t$ is a $C^{*}$-algebra varying with $t$. We now
define a construction of $C^{*}$-algebras that includes the preceding
constructions: $T$ will be a topological space, not necessarily discrete,
and the $A_t$ will vary with $t$.

The principal interest of this construction is that fairly broad classes of
$C^{*}$-algebras can be obtained in this way with very simple $A_t$. This
will enable us to solve, very partially, the second problem stated at the
beginning of \SectionRef{3}{Section~3}.

\BookSubsection{10.1}{Continuous Fields of Banach Spaces}

\ResultHeading{10.1.1}{10.1.1.}
Let $T$ be a topological space and $(E(t))_{t\in T}$ a family of complex
Banach spaces. A \emph{vector field} is any element of
$\prod_{t\in T}E(t)$; that is, any function $x$ defined on $T$ such that
$x(t)\in E(t)$ for every $t\in T$. More generally, if $Y\subset T$, a
vector field on $Y$ is an element of $\prod_{t\in Y}E(t)$.

\ResultHeading{10.1.2}{10.1.2. Definition.}
\begin{itshape}
Let $T$ be a topological space. A \emph{continuous field $\mathcal E$ of
Banach spaces over $T$} is a family $(E(t))_{t\in T}$ of Banach spaces,
equipped with a set $\Gamma\subset\prod_{t\in T}E(t)$ of vector fields such
that:
\begin{enumerate}[label=\textup{(\roman*)}]
\item $\Gamma$ is a complex vector subspace of $\prod_{t\in T}E(t)$;
\item for every $t\in T$, the set of $x(t)$, where $x\in\Gamma$, is dense in
$E(t)$;
\item for every $x\in\Gamma$, the function $t\mapsto\|x(t)\|$ is
continuous;
\item if $x\in\prod_{t\in T}E(t)$ is a vector field such that, for every
$t\in T$ and every $\varepsilon>0$, there is an $x'\in\Gamma$ for which
$\|x(t')-x'(t')\|\leq\varepsilon$ on a neighborhood of $t$, then
$x\in\Gamma$.
\end{enumerate}
\end{itshape}
The elements of $\Gamma$ are called the \emph{continuous vector fields} of
$\mathcal E$. Definition~\XRef{10.1.2}{10.1.2} is analogous to the
definition of measurable fields.

\ResultHeading{10.1.3}{10.1.3.}
Let
\[
\mathcal E=((E(t))_{t\in T},\Gamma),
\qquad
\mathcal E'=((E'(t))_{t\in T},\Gamma')
\]
% Source PDF physical page 197
be two continuous fields of Banach spaces over $T$. An \emph{isomorphism}
of $\mathcal E$ onto $\mathcal E'$ is a family
$\varphi=(\varphi_t)_{t\in T}$ such that each $\varphi_t$ is an isometric
isomorphism of $E(t)$ onto $E'(t)$ and $\varphi$ carries $\Gamma$ onto
$\Gamma'$.

\ResultHeading{10.1.4}{10.1.4. Example.}
Let $E$ be a Banach space, let $T$ be a topological space, and let $\Gamma$
be the set of continuous maps from $T$ into $E$. Put $E(t)=E$ for every
$t\in T$. Then
$\mathcal E=((E(t))_{t\in T},\Gamma)$ is a continuous field of Banach
spaces over $T$, called the \emph{constant field over $T$ defined by $E$}.
(The reader will do well to keep this example constantly in mind in what
follows.) A field isomorphic to a constant field is called \emph{trivial}.

\ResultHeading{10.1.5}{10.1.5. Example.}
In Definition~\XRef{10.1.2}{10.1.2}, suppose $T$ is discrete. Axioms (ii)
and (iv) then imply that necessarily
\[
\Gamma=\prod_{t\in T}E(t).
\]

\ResultHeading{10.1.6}{10.1.6.}
Let $T$ be a topological space and
$\mathcal E=((E(t))_{t\in T},\Gamma)$ a continuous field of Banach spaces
over $T$. Let $Y\subset T$ and $t_0\in Y$. A vector field $x$ on $Y$ is
said to be \emph{continuous at $t_0$} if, for every $\varepsilon>0$, there
is an $x'\in\Gamma$ such that
$\|x(t)-x'(t)\|\leq\varepsilon$ on a neighborhood of $t_0$. It is
\emph{continuous on $Y$} if it is continuous at every point of $Y$. When
$Y=T$, this agrees with the definition already adopted for continuous
vector fields, by axiom~(iv).

\ResultHeading{10.1.7}{10.1.7.}
Let $\Gamma_Y$ be the set of continuous vector fields on $Y$. It is
immediate that $((E(t))_{t\in Y},\Gamma_Y)$ is a continuous field of Banach
spaces over $Y$, called the field induced by $\mathcal E$ on $Y$ and
denoted by $\mathcal E|Y$.

\ResultHeading{10.1.8}{10.1.8.}
If every point of $T$ has a neighborhood $V$ such that $\mathcal E|V$ is
trivial, then $\mathcal E$ is said to be \emph{locally trivial}.

\ResultHeading{10.1.9}{10.1.9. Proposition.}
\begin{itshape}
Let $((E(t))_{t\in T},\Gamma)$ be a continuous field of Banach spaces.
\begin{enumerate}[label=\textup{(\roman*)}]
\item Let $t_0\in T$, let $x$ be a vector field on $T$ continuous at $t_0$,
and let $f:T\to\mathbb C$ be a function continuous at $t_0$. Then the
vector field
\[
t\longmapsto(fx)(t)=f(t)x(t)
\]
is continuous at $t_0$.
\item If $y\in\Gamma$ and $g:T\to\mathbb C$ is continuous on $T$, then
$gy\in\Gamma$.
\end{enumerate}
\end{itshape}

Let $\varepsilon>0$. There is an $x'\in\Gamma$ such that
$\|x(t)-x'(t)\|\leq\varepsilon$ near $t_0$. Near $t_0$ we have
\[
|f(t)-f(t_0)|\leq\varepsilon,
\qquad |f(t)|\leq|f(t_0)|+1,
\qquad \|x'(t)\|\leq\|x'(t_0)\|+1.
\]
% Source PDF physical page 198
Therefore, near $t_0$,
\begin{align*}
\|f(t)x(t)-f(t_0)x'(t)\|
&\leq|f(t)|\,\|x(t)-x'(t)\|
+|f(t)-f(t_0)|\,\|x'(t)\|\\
&\leq\varepsilon\bigl(|f(t_0)|+\|x'(t_0)\|+2\bigr).
\end{align*}
This proves (i), and (ii) follows from (i).

\ResultHeading{10.1.10}{10.1.10. Proposition.}
\begin{itshape}
For every $t_0\in T$ and every $\xi\in E(t_0)$, there is an $x\in\Gamma$
such that $x(t_0)=\xi$.
\end{itshape}

In view of axiom~(ii) of Definition~\XRef{10.1.2}{10.1.2}, there is a
sequence $(\xi_n)$ of nonzero vectors of $E(t_0)$ with the following
properties:
\begin{enumerate}[label=\textup{\alph*.}]
\item $\sum\|\xi_n\|<+\infty$;
\item $\sum\xi_n=\xi$;
\item for every $n$ there is an $x_n\in\Gamma$ such that
$x_n(t_0)=\xi_n$.
\end{enumerate}
Let $f_n$ be the function on $T$ equal to $1$ when
$\|x_n(t)\|\leq\|\xi_n\|$, and equal to
$\|\xi_n\|/\|x_n(t)\|$ when $\|x_n(t)\|>\|\xi_n\|$. This function is
continuous, so $y_n=f_nx_n\in\Gamma$ by
\XRef{10.1.9}{10.1.9}(ii). We have $y_n(t_0)=\xi_n$ and
$\|y_n(t)\|\leq\|\xi_n\|$ for every $t\in T$; hence
$\sum y_n(t)$ converges for every $t$ to an element $x(t)\in E(t)$. Also
$x(t_0)=\sum\xi_n=\xi$. For every $\varepsilon>0$, there is a positive
integer $p$ such that
\[
\left\|x(t)-\sum_{n=1}^{p}y_n(t)\right\|\leq\varepsilon
\quad\text{on all of }T.
\]
Thus $x\in\Gamma$ by axioms (i) and (iv) of
Definition~\XRef{10.1.2}{10.1.2}.

\ResultHeading{10.1.11}{10.1.11. Lemma.}
\begin{itshape}
Let $T$ be paracompact, let $Y$ be a closed subset of $T$, let
$((E(t))_{t\in T},\Gamma)$ be a continuous field of Banach spaces over
$T$, let $x$ (respectively, $y$) be a continuous vector field on $T$
(respectively, on $Y$), and let $a>0$ satisfy
$\|x(t)-y(t)\|\leq a$ on $Y$. There exists $x'\in\Gamma$ such that
$\|x'(t)-y(t)\|\leq a/2$ on $Y$ and
$\|x'(t)-x(t)\|\leq2a$ on $T$.
\end{itshape}

There is a cover $(V_i)$ of $Y$ by open subsets of $T$ and, for every $i$,
an $x_i\in\Gamma$ such that $\|x_i(t)-y(t)\|\leq a/2$ on $V_i$.
Replacing $(V_i)$ by a refinement, we may suppose it locally finite. Let
$(\eta_i)$ be a family of nonnegative continuous functions on $T$
subordinate to $(V_i)$, such that $\sum\eta_i=1$ on $Y$. Put
$w=\sum\eta_ix_i$, which is continuous at every point of $T$, and hence
belongs to $\Gamma$. If $t\in Y$, only the indices $i$ for which
$\eta_i(t)>0$ contribute to $\sum\eta_i(t)x_i(t)$, and for those indices
\[
\|x_i(t)-y(t)\|\leq\frac a2;
\qquad\text{hence}\qquad
\left\|\sum\eta_i(t)x_i(t)-y(t)\right\|\leq\frac a2.
\]
% Source PDF physical page 199
Thus $\|w(t)-x(t)\|\leq2a$ on a neighborhood $Z$ of $Y$. Let
$f:T\to[0,1]$ be continuous, equal to $1$ on $Y$ and to $0$ on $T-Z$.
Put $x'=fw+(1-f)x$. On $Y$,
\[
\|x'(t)-y(t)\|=\|w(t)-y(t)\|\leq\frac a2.
\]
On $T$,
\[
\|x'(t)-x(t)\|=|f(t)|\,\|w(t)-x(t)\|\leq2a.
\]

\ResultHeading{10.1.12}{10.1.12. Proposition.}
\begin{itshape}
Let $T$ be paracompact, let $Y$ be a closed subset of $T$, let
$((E(t))_{t\in T},\Gamma)$ be a continuous field of Banach spaces, and let
$y$ be a continuous vector field on $Y$. There is an $x\in\Gamma$ extending
$y$.
\end{itshape}

Arguing as in Lemma~\XRef{10.1.11}{10.1.11}, first construct
$x_0\in\Gamma$ such that $\|x_0(t)-y(t)\|\leq1$ on $Y$. Then, applying
Lemma~\XRef{10.1.11}{10.1.11} successively, obtain a sequence
$(x_n)\subset\Gamma$ such that
\[
\|x_n(t)-y(t)\|\leq2^{-n}\quad\text{on }Y,
\qquad
\|x_n(t)-x_{n-1}(t)\|\leq2^{-n+2}\quad\text{on }T.
\]
The $x_n$ converge uniformly to an $x\in\Gamma$, and $x(t)=y(t)$ on $Y$.

\ResultHeading{10.1.13}{10.1.13. Proposition.}
\begin{itshape}
Let $T$ be completely regular, let $(Y_i)_{i\in I}$ be an open cover of
$T$, and put $Y_{ij}=Y_i\cap Y_j$. For every $i\in I$, let
$\mathcal E_i=((E_i(t))_{t\in Y_i},\Gamma_i)$ be a continuous field of
Banach spaces over $Y_i$. For all $i,j\in I$, let $g_{ij}$ be an isomorphism
of $\mathcal E_j|Y_{ij}$ onto $\mathcal E_i|Y_{ij}$. Suppose that
$g_{ij}g_{jk}=g_{ik}$ for all $i,j,k\in I$. Then there exists, uniquely up
to isomorphism, a continuous field $\mathcal E$ of Banach spaces over $T$
with the following property: for every $i\in I$ there is an isomorphism
$h_i$ of $\mathcal E_i$ onto $\mathcal E|Y_i$ such that
$g_{ij}=h_i^{-1}h_j$ for all $i,j\in I$.
\end{itshape}

Uniqueness is easy; we prove existence. For every $t\in T$, let
$I(t)\neq\varnothing$ be the set of $i\in I$ such that $t\in Y_i$. If
$i,j\in I(t)$, then $g_{ij}(t)$ is an isomorphism of $E_j(t)$ onto $E_i(t)$,
and
\[
g_{ij}(t)g_{jk}(t)=g_{ik}(t)
\qquad(i,j,k\in I(t)).
\]
There is therefore a Banach space $E(t)$ and, for every $i\in I(t)$, an
isomorphism $h_i(t)$ of $E_i(t)$ onto $E(t)$ such that
\[
g_{ij}(t)=h_i(t)^{-1}h_j(t)
\qquad(i,j\in I(t)).
\]
[For example, take $E(t)=E_{i_0}(t)$, where $i_0\in I(t)$, and
$h_i(t)=g_{i_0i}(t)$.] There is a unique set $\Delta_i$ of vector fields on
$Y_i$, with values in the $E(t)$, such that $(h_i(t))_{t\in Y_i}$ is an
isomorphism of $\mathcal E_i$ onto
$\mathcal F_i=((E(t))_{t\in Y_i},\Delta_i)$. Moreover,
$\mathcal F_i|Y_{ij}=\mathcal F_j|Y_{ij}$ for all $i,j\in I$.

% Source PDF physical page 200
Now let $\Gamma$ be the set of all
$x\in\prod_{t\in T}E(t)$ such that $x|Y_i\in\Delta_i$ for every $i\in I$.
It is immediate that $\Gamma$ satisfies axioms (i), (iii), and (iv) of
Definition~\XRef{10.1.2}{10.1.2}. We show that it also satisfies axiom
(ii). Let $t_0\in T$ and $\xi\in E(t_0)$. Choose $i\in I$ with
$t_0\in Y_i$. Let $V\subset Y_i$ be a neighborhood of $t_0$ that is closed
in $T$, and let $\eta:T\to\mathbb R$ be a continuous function equal to $1$
at $t_0$ and to $0$ on $T-V$. Choose $y\in\Delta_i$ with $y(t_0)=\xi$.
Let $x\in\prod_{t\in T}E(t)$ be $0$ on $T-V$ and $\eta y$ on $V$.
Then $x(t_0)=\xi$. We shall show that $x\in\Gamma$. Let $j\in I$; we must
show that $x|Y_j\in\Delta_j$. It is enough to show that, for every
$t_1\in Y_j$, $x$ agrees near $t_1$ with an element of $\Delta_j$. This is
immediate if $t_1\notin V$. Suppose $t_1\in V$. Then $t_1\in Y_{ij}$.
Since $\mathcal F_i|Y_{ij}=\mathcal F_j|Y_{ij}$ and $x$ agrees near $t_1$
with an element of $\Delta_i$, the assertion follows.

We have thus proved that $\mathcal E=((E(t))_{t\in T},\Gamma)$ is a
continuous field of Banach spaces over $T$. It is easy to see that
$\mathcal E|Y_i=\mathcal F_i$, which proves the proposition. We say that
$\mathcal E$ is obtained by \emph{gluing} the $\mathcal E_i$ by means of the
$g_{ij}$.

Bibliography: \cite{ref19,ref29,ref52,ref77,ref108,ref161,ref169}.

\BookSubsection{10.2}{Total Subsets}

\ResultHeading{10.2.1}{10.2.1. Definition.}
\begin{itshape}
Let $\mathcal E=((E(t))_{t\in T},\Gamma)$ be a continuous field of Banach
spaces over $T$. A subset $\Lambda\subset\Gamma$ is called \emph{total} if,
for every $t\in T$, the set of $x(t)$ with $x\in\Lambda$ is total in
$E(t)$. The field $\mathcal E$ is called \emph{separable} if $\Gamma$
contains a countable total subset.
\end{itshape}

If $\Lambda$ is a countable total subset of $\Gamma$, the set $\Lambda'$ of
linear combinations of elements of $\Lambda$ with rational complex
coefficients is again countable, and $\Lambda'(t)$ is dense in $E(t)$ for
every $t$.

\ResultHeading{10.2.2}{10.2.2. Proposition.}
\begin{itshape}
Let $((E(t))_{t\in T},\Gamma)$ be a continuous field of Banach spaces over
$T$. Let $\Lambda\subset\Gamma$ be a total subset, and let
$\overline\Lambda$ be the vector subspace of $\Gamma$ spanned by
$\Lambda$. For a vector field $x\in\prod_{t\in T}E(t)$, the following
properties are equivalent:
\begin{enumerate}[label=\textup{(\roman*)}]
\item[\textup{(i)}] $x\in\Gamma$;
\item[\textup{(ii)}] for every $t_0\in T$ and every $\varepsilon>0$, there is an
$x'\in\Gamma$ such that $\|x(t)-x'(t)\|\leq\varepsilon$ near $t_0$;
\item[\textup{(ii$'$)}] for every $t_0\in T$ and every $\varepsilon>0$,
there is an $x'\in\overline\Lambda$ such that
$\|x(t)-x'(t)\|\leq\varepsilon$ near $t_0$;
\item[\textup{(iii)}] for every $x'\in\Gamma$, the function
$t\mapsto\|x(t)-x'(t)\|$ is continuous;
\item[\textup{(iii$'$)}] for every $x'\in\overline\Lambda$, the function
$t\mapsto\|x(t)-x'(t)\|$ is continuous.
\end{enumerate}
\end{itshape}

The implications (ii$'$)$\Rightarrow$(ii)$\Rightarrow$(i)$\Rightarrow$
(iii)$\Rightarrow$(iii$'$) are immediate.

% Source PDF physical page 201
For (iii$'$)$\Rightarrow$(ii$'$), suppose (iii$'$) holds. Let $t_0\in T$
and $\varepsilon>0$. Since $\overline\Lambda(t_0)$ is dense in $E(t_0)$,
there is an $x'\in\overline\Lambda$ such that
$\|x(t_0)-x'(t_0)\|<\varepsilon$. By (iii$'$),
$\|x(t)-x'(t)\|<\varepsilon$ near $t_0$, proving (ii$'$).

\ResultHeading{10.2.3}{10.2.3. Proposition.}
\begin{itshape}
Let $T$ be a topological space, let $(E(t))_{t\in T}$ be a family of Banach
spaces, and let $\Lambda\subset\prod_{t\in T}E(t)$ satisfy axioms (i), (ii),
and (iii) of Definition~\XRef{10.1.2}{10.1.2}, with $\Gamma$ replaced by
$\Lambda$. There is a unique subset
$\Gamma\subset\prod_{t\in T}E(t)$ containing $\Lambda$ and satisfying
axioms (i)--(iv) of Definition~\XRef{10.1.2}{10.1.2}: namely, the set of
vector fields $x$ satisfying condition (ii$'$) of
Proposition~\XRef{10.2.2}{10.2.2}.
\end{itshape}

Uniqueness follows from \XRef{10.2.2}{10.2.2}: if $\Gamma$ exists, its
elements are characterized by (ii$'$). Conversely, the set of vector fields
$x$ satisfying (ii$'$) plainly satisfies axioms (i)--(iv) of
Definition~\XRef{10.1.2}{10.1.2} and contains $\Lambda$.

\ResultHeading{10.2.4}{10.2.4. Proposition.}
\begin{itshape}
Let $\mathcal E=((E(t))_{t\in T},\Gamma)$ and
$\mathcal E'=((E'(t))_{t\in T},\Gamma')$ be two continuous fields of Banach
spaces over $T$, and let $\Lambda$ be a total subset of $\Gamma$. For each
$t\in T$, let $\varphi_t$ be an isometric isomorphism of $E(t)$ onto
$E'(t)$. Then $\varphi=(\varphi_t)_{t\in T}$ is an isomorphism of
$\mathcal E$ onto $\mathcal E'$ if and only if
$\varphi(\Lambda)\subset\Gamma'$.
\end{itshape}

Necessity is clear. Suppose $\varphi(\Lambda)\subset\Gamma'$ and prove that
$\varphi$ is an isomorphism. We may suppose that $\Lambda$ is a vector
subspace of $\Gamma$. If $x\in\Gamma$, then $x$ is a limit of elements of
$\Lambda$ for local uniform convergence [\XRef{10.2.2}{10.2.2}(ii$'$)];
hence $\varphi(x)$ is a local uniform limit of elements of $\Gamma'$, and
so $\varphi(x)\in\Gamma'$. Conversely, let $x'\in\Gamma'$, put
$x=\varphi^{-1}(x')$, and let $t_0\in T$, $\varepsilon>0$. There is a
$y\in\Gamma$ such that $y(t_0)=x(t_0)$. We have
\[
\varphi(y)\in\Gamma',
\qquad
\varphi(y)(t_0)=\varphi_{t_0}(x(t_0))=x'(t_0),
\]
and hence $\|\varphi(y)(t)-x'(t)\|\leq\varepsilon$ near $t_0$, so
$\|y(t)-x(t)\|\leq\varepsilon$ near $t_0$. Thus $x\in\Gamma$, and
$\varphi^{-1}(\Gamma')\subset\Gamma$.

\ResultHeading{10.2.5}{10.2.5. Proposition.}
\begin{itshape}
Let $T$ be locally compact, let $((E(t))_{t\in T},\Gamma)$ be a continuous
field of Banach spaces over $T$, let $x\in\Gamma$, let $K$ be a compact
subset of $T$, let $\varepsilon>0$, and let $\Lambda$ be a total subset of
$\Gamma$. There exist compactly supported continuous complex functions
$f_1,\ldots,f_n$ on $T$ and elements $x_1,\ldots,x_n\in\Lambda$ such that
\[
\|x(t)-f_1(t)x_1(t)-\cdots-f_n(t)x_n(t)\|\leq\varepsilon
\quad\text{on }K.
\]
\end{itshape}

% Source PDF physical page 202
Let $\overline\Lambda$ be the vector subspace of $\Gamma$ spanned by
$\Lambda$. By \XRef{10.2.2}{10.2.2}(ii$'$), there is a cover $(V_i)$ of
$K$ by relatively compact open subsets of $T$, and elements
$x'_i\in\overline\Lambda$, such that
$\|x(t)-x'_i(t)\|\leq\varepsilon$ in $V_i$. Since $K$ is compact, we may
suppose the cover is finite, say $(V_1,\ldots,V_p)$. Let
$(\eta_1,\ldots,\eta_{p+1})$ be a continuous partition of unity on $T$
subordinate to the open cover $(V_1,\ldots,V_p,T-K)$. Then
\[
\|x(t)-\eta_1(t)x'_1(t)-\cdots-\eta_p(t)x'_p(t)\|\leq\varepsilon
\]
throughout $K$. Expressing the $x'_i$ as linear combinations of elements
of $\Lambda$ proves the proposition.

\ResultHeading{10.2.6}{10.2.6.}
Let $T$ be locally compact, let
$\mathcal E=((E(t))_{t\in T},\Gamma)$ be a continuous field of Banach
spaces over $T$, and let $T'$ be the compact space obtained by adjoining a
point at infinity $\omega$ to $T$. Put $E'(t)=E(t)$ for $t\in T$ and
$E'(\omega)=0$. Let $\Gamma'$ be the set of
$x\in\prod_{t\in T'}E'(t)$ such that $x|T\in\Gamma$ and
$\|x(t)\|\to0$ at infinity on $T$.

\ResultLabel{Proposition.}
\begin{itshape}
$\mathcal E'=((E'(t))_{t\in T'},\Gamma')$ is a continuous field of Banach
spaces over $T'$, and $\mathcal E'|T=\mathcal E$.
\end{itshape}

It is immediate that $\Gamma'$ satisfies axioms (i), (iii), and (iv) of
Definition~\XRef{10.1.2}{10.1.2}. Let $t_0\in T$ and $\xi\in E(t_0)$.
There is an $x\in\Gamma$ such that $x(t_0)=\xi$. Multiplying $x$ by a
compactly supported function equal to $1$ at $t_0$, we may suppose that
$\|x(t)\|\to0$ at infinity on $T$; hence there is at once an $x'\in\Gamma'$
with $x'(t_0)=\xi$. Let $\Gamma_0$ be the set of restrictions to $T$ of
elements of $\Gamma'$; that is, the set of elements of $\Gamma$ tending to
$0$ at infinity. What precedes shows that $\Gamma_0$ is a total subset of
$\Gamma$. For every $t\in T$, let $\varphi_t$ be the identity map on
$E(t)$, and put $\varphi=(\varphi_t)_{t\in T}$. For every
$x\in\Gamma_0$, $\varphi(x)$ is a continuous vector field relative to
$\mathcal E'|T$. Thus $\varphi$ is an isomorphism of $\mathcal E$ onto
$\mathcal E'|T$ by \XRef{10.2.4}{10.2.4}.

\ResultHeading{10.2.7}{10.2.7. Proposition.}
\begin{itshape}
Let $T$ be a separable metrizable space and
$\mathcal E=((E(t))_{t\in T},\Gamma)$ a locally trivial continuous field of
Banach spaces over $T$. If every $E(t)$ is separable, then $\mathcal E$ is
separable.
\end{itshape}

There is a countable open cover $(U_n)$ of $T$ such that every
$\mathcal E|U_n$ is trivial, and then an open cover $(V_n)$ of $T$ such
that $\overline{V_n}\subset U_n$ for every $n$. For every $n$, there is a
sequence $(x_{nm})_m$ of bounded continuous vector fields on $U_n$ such
that, for every $t\in U_n$, the $x_{nm}(t)$ are dense in $E(t)$ (indeed,
$\mathcal E|U_n$ is identified with the constant field defined by a
separable space). Let $\eta_n:T\to[0,1]$ be continuous, equal to $1$ on
$V_n$, and to $0$ on $T-U_n$. Let $y_{nm}$ be the vector field on $T$ equal
to $\eta_nx_{nm}$ on $U_n$ and to $0$ on $T-U_n$. Then
$y_{nm}\in\Gamma$: if $t\in U_n$, $y_{nm}$ is continuous at $t$; if
$t\in T-U_n$ and $\varepsilon>0$, there is
% Source PDF physical page 203
a neighborhood of $t$ on which $\|y_{nm}(t')\|\leq\varepsilon$. Let
$t\in T$ and $\xi\in E(t)$. There is an $n$ such that $t\in V_n$. Then
$\xi$ is a limit of vectors $x_{nm}(t)=y_{nm}(t)$. Thus the $y_{nm}$ form
a total subset of $\Gamma$.

Bibliography: \cite{ref19,ref29,ref52,ref77,ref108,ref161,ref169}.

\BookSubsection[Continuous Fields of C*-Algebras]{10.3}{Continuous Fields
of $C^{*}$-Algebras}

\ResultHeading{10.3.1}{10.3.1. Definition.}
\begin{itshape}
Let $T$ be a topological space. A \emph{continuous field of
$C^{*}$-algebras over $T$} is a continuous field
$((A(t))_{t\in T},\Theta)$ of Banach spaces over $T$, each $A(t)$ being
equipped with a multiplication and an involution that make it a
$C^{*}$-algebra, and $\Theta$ being stable under multiplication and
involution.
\end{itshape}

When we speak of isomorphisms of continuous fields of $C^{*}$-algebras, we
shall of course mean isomorphisms compatible with the involutive-algebra
structure on the fibers. The same convention applies to triviality and
local triviality.

A continuous field $((A(t))_{t\in T},\Theta)$ is called a
\emph{continuous field of elementary $C^{*}$-algebras} if every $A(t)$ is
an elementary $C^{*}$-algebra (\XRef{4.1.1}{4.1.1}).

\ResultHeading{10.3.2}{10.3.2. Proposition.}
\begin{itshape}
Let $\mathcal A=((A(t))_{t\in T},\Theta)$ be a continuous field of Banach
spaces over $T$, each $A(t)$ being a $C^{*}$-algebra. For $\mathcal A$ to
be a continuous field of $C^{*}$-algebras, it is enough that there be a
total subset $\Lambda$ of $\Theta$ stable under multiplication and
involution.
\end{itshape}

Indeed, let $\overline\Lambda$ be the vector subspace of $\Theta$ spanned
by $\Lambda$; it is stable under multiplication and involution. Every
element of $\Theta$ is a local uniform limit of elements of
$\overline\Lambda$. Hence, if $x,y\in\Theta$, the vector field $xy$ is a
local uniform limit of elements of $\overline\Lambda$, so $xy\in\Theta$.
Likewise $x^{*}\in\Theta$.

\ResultHeading{10.3.3}{10.3.3. Proposition.}
\begin{itshape}
Let $((A(t))_{t\in T},\Theta)$ be a continuous field of $C^{*}$-algebras
over $T$, let $x\in\Theta$ satisfy $xx^{*}=x^{*}x$, and let
$f:\mathbb C\to\mathbb C$ be continuous with $f(0)=0$. Then
$t\mapsto f(x(t))$ belongs to $\Theta$.
\end{itshape}

Let $t_0\in T$. There is a neighborhood $V$ of $t_0$ and a constant
$M>0$ such that $\|x(t)\|\leq M$ on $V$. Let $(p_n)$ be a sequence of
polynomials without constant terms in $z$ and $\overline z$, converging
uniformly to $f$ on the disk $|z|\leq M$. Then $p_n(x,x^{*})\in\Theta$, and
$p_n(x(t),x^{*}(t))$ converges uniformly on $V$ to $f(x(t))$. This proves
the proposition. The field $t\mapsto f(x(t))$ is denoted by $f(x)$.

\ResultHeading{10.3.4}{10.3.4. Corollary.}
\begin{itshape}
If $x\in\Theta$ and $x(t)\geq0$ for every $t$, then the field
$t\mapsto x(t)^{1/2}$ is an element $y$ of $\Theta$ satisfying
$y=y^{*}$ and $y^2=x$.
\end{itshape}

% Source PDF physical page 204
\ResultHeading{10.3.5}{10.3.5.}
If every $A(t)$ has an identity $1_t$, and if $t\mapsto1_t$ belongs to
$\Theta$, then the hypothesis $f(0)=0$ can be omitted in
\XRef{10.3.3}{10.3.3}.

\ResultHeading{10.3.6}{10.3.6. Proposition.}
\begin{itshape}
Let $((A(t))_{t\in T},\Theta)$ be a continuous field of $C^{*}$-algebras
over $T$, let $x\in\Theta$ satisfy $xx^{*}=x^{*}x$, let $t_0\in T$, and let
$U$ be a neighborhood of $\operatorname{Sp}'x(t_0)$ in $\mathbb C$. There
is a neighborhood $V$ of $t_0$ in $T$ such that
$\operatorname{Sp}'x(t)\subset U$ for every $t\in V$.
\end{itshape}

Let $f:\mathbb C\to\mathbb C$ be continuous, equal to $0$ on
$\operatorname{Sp}'x(t_0)$ and to $1$ on $\mathbb C-U$. Then
$f(x)(t_0)=0$, so $\|f(x)(t)\|<1$ in a neighborhood $V$ of $t_0$. If
$t\in V$, $\operatorname{Sp}'x(t)$ cannot meet $\mathbb C-U$.

Bibliography: \cite{ref19,ref29,ref77,ref108,ref169}.

\BookSubsection[The C*-Algebra Defined by a Continuous Field of
C*-Algebras]{10.4}{The $C^{*}$-Algebra Defined by a Continuous Field of
$C^{*}$-Algebras}

\ResultHeading{10.4.1}{10.4.1.}
Let $T$ be locally compact, let
$\mathcal A=((A(t))_{t\in T},\Theta)$ be a continuous field of
$C^{*}$-algebras over $T$, and let $A$ be the set of $x\in\Theta$ such that
$\|x(t)\|$ tends to $0$ at infinity on $T$. Then $A$ is an involutive
subalgebra of $\Theta$. For $x\in A$, put
\[
\|x\|=\sup_{t\in T}\|x(t)\|<+\infty.
\]
It is immediate that $A$ is a $C^{*}$-algebra for this norm. It is called
the \emph{$C^{*}$-algebra defined by $\mathcal A$}. For every $t\in T$, the
map $x\mapsto x(t)$ is a morphism of $A$ onto $A(t)$.

\ResultHeading{10.4.2}{10.4.2. Lemma.}
\begin{itshape}
Let $I$ be a closed left or right ideal of $A$. For every $t\in T$, let
$I(t)$ be the set of $x(t)$ with $x\in I$. Then $I$ is the set of all
$x\in A$ such that $x(t)\in I(t)$ for every $t\in T$.
\end{itshape}

We argue for a right ideal. If $x\in I$, then $x(t)\in I(t)$ for every $t$.
Conversely, let $x\in A$ satisfy $x(t)\in I(t)$ for every $t$, and let
$\varepsilon>0$. For every $t\in T$, there are $y_t\in I$ and $z_t\in A$
such that $x(t)=y_t(t)$ and
$\|y_t(t)-y_t(t)z_t(t)\|\leq\varepsilon$ (\XRef{1.7.2}{1.7.2}); hence
$\|x(t)-y_t(t)z_t(t)\|\leq\varepsilon$ near $t$. Let $K$ be the compact
set of $t\in T$ for which $\|x(t)\|\geq\varepsilon$. There is a cover
$(U_1,\ldots,U_n)$ of $K$ by open subsets of $T$, and, for each $i$,
elements $y_i\in I$, $z_i\in A$, such that
$\|x(t)-y_i(t)z_i(t)\|\leq\varepsilon$ on $U_i$. Put
$U_0=T-K$ and $y_0=z_0=0$. Let $(f_0,f_1,\ldots,f_n)$ be a continuous
partition of unity subordinate to $(U_0,U_1,\ldots,U_n)$. Since
$f_iz_i\in A$ by \XRef{10.1.9}{10.1.9}(ii), we have
$\sum f_i y_i z_i\in I$. Let $t\in T$, and let $i_1,\ldots,i_p$ be the
indices $i$ such that $t\in U_i$. Then
\[
\sum_i f_i(t)y_i(t)z_i(t)
=f_{i_1}(t)y_{i_1}(t)z_{i_1}(t)+\cdots
+f_{i_p}(t)y_{i_p}(t)z_{i_p}(t),
\]
and
\[
\|x(t)-y_{i_1}(t)z_{i_1}(t)\|,\ldots,
\|x(t)-y_{i_p}(t)z_{i_p}(t)\|\leq\varepsilon.
\]
% Source PDF physical page 205
Since $f_{i_1}(t)+\cdots+f_{i_p}(t)=1$, it follows that
\[
\left\|x(t)-\sum_i f_i(t)y_i(t)z_i(t)\right\|\leq\varepsilon.
\]
Thus $\|x-\sum_i f_i y_i z_i\|\leq\varepsilon$. Since $I$ is closed,
$x\in I$.

\ResultHeading{10.4.3}{10.4.3. Theorem.}
\begin{itshape}
Let $T$ be locally compact, let
$\mathcal A=((A(t))_{t\in T},\Theta)$ be a continuous field of
$C^{*}$-algebras over $T$, and let $A$ be the $C^{*}$-algebra defined by
$\mathcal A$. For $t\in T$ and $\pi\in\widehat{A(t)}$, let $\rho_\pi$ be
the irreducible representation $x\mapsto\pi(x(t))$ of $A$. Then
$\pi\mapsto\rho_\pi$ is a bijection from the disjoint union
$\coprod_{t\in T}\widehat{A(t)}$ onto $\widehat A$.
\end{itshape}

Let $\pi\in\widehat{A(t)}$ and $\pi'\in\widehat{A(t')}$. If $t\neq t'$,
there is a continuous complex function $f$ on $T$ such that $f(t)=1$ and
$f(t')=0$. Choose $x\in A$ with $\pi(x(t))\neq0$. Then
\[
\rho_\pi(fx)=\pi(x(t))\neq0,
\qquad
\rho_{\pi'}(fx)=\pi'(0)=0,
\]
so $\rho_\pi$ and $\rho_{\pi'}$ are inequivalent. Thus
$\pi\mapsto\rho_\pi$ is injective. Now let $\rho\in\widehat A$, and let
$I=\ker\rho$. Let $I(t)$ be the set of $x(t)$ with $x\in I$, and let $Y$
be the set of $t\in T$ for which $I(t)\neq A(t)$. By
\XRef{10.4.2}{10.4.2}, $Y\neq\varnothing$. Suppose $Y$ contains two
distinct points $t_1,t_2$, and choose disjoint neighborhoods $U_1,U_2$ of
them. Let $K_i$ be the closed two-sided ideal of $A$ consisting of the
$x\in A$ that vanish outside $U_i$. Then
$I\supset0=K_1K_2$, so, for example, $I\supset K_1$ by
\XRef{2.11.4}{2.11.4}. It follows that
$I(t_1)\supset K_1(t_1)=A(t_1)$, contrary to $t_1\in Y$. Hence $Y$
consists of one point $t_0\in T$. By \XRef{10.4.2}{10.4.2}, $I$ is the set
of $x\in A$ such that $x(t_0)\in I(t_0)$. Thus $\rho$ is the composite of
the morphism $x\mapsto x(t_0)$ and an irreducible representation of
$A(t_0)$.

\ResultHeading{10.4.4}{10.4.4. Corollary.}
\begin{itshape}
Let $T$ be locally compact, let
$\mathcal A=((A(t))_{t\in T},\Theta)$ be a continuous field of nonzero
elementary $C^{*}$-algebras over $T$, and let $A$ be the $C^{*}$-algebra
defined by $\mathcal A$. For each $t\in T$, let $\rho_t$ be the irreducible
representation of $A$ obtained by composing $x\mapsto x(t)$ with the
irreducible representation of $A(t)$, unique up to equivalence. Then
$t\mapsto\rho_t$ is a homeomorphism from $T$ onto $\widehat A$.
\end{itshape}

This map is bijective by \XRef{10.4.3}{10.4.3}. Let $Y\subset T$ and
$t_0\in T$. The intersection of the kernels of the $\rho_t$, $t\in Y$, is
the set $I_Y$ of $x\in A$ that vanish on $Y$. The point $\rho_{t_0}$ lies
in the closure of the set of $\rho_t$, $t\in Y$, if and only if
$I_{\{t_0\}}\supset I_Y$. If $t_0\in\overline Y$, this inclusion is clear.
If $t_0\notin\overline Y$, there is a continuous complex function on $T$
equal to $1$ at $t_0$ and to $0$ on $Y$, hence an $x\in A$ nonzero at
$t_0$ and zero on $Y$; thus $I_{\{t_0\}}\not\supset I_Y$. This proves the
corollary. We identify $T$ with $\widehat A$ by this homeomorphism.

% Source PDF physical page 206
\ResultHeading{10.4.5}{10.4.5. Corollary.}
\begin{itshape}
With the notation of Corollary~\XRef{10.4.4}{10.4.4}, $A$ is a liminal
$C^{*}$-algebra with spectrum $T$.
\end{itshape}

Indeed, for every $t\in T=\widehat A$, $\rho_t(A)$ consists of compact
operators, since $A(t)$ is elementary.

\ResultHeading{10.4.6}{10.4.6. Proposition.}
\begin{itshape}
Let $T$ be compact, let $\mathcal A=((A(t))_{t\in T},\Theta)$ be a
continuous field of $C^{*}$-algebras over $T$, and let $A=\Theta$ be the
$C^{*}$-algebra defined by $\mathcal A$. Suppose that every $A(t)$ has an
identity $1_t$ and that the field $t\mapsto1_t$, denoted by $1$, belongs to
$A$. Let $B$ be a $C^{*}$-subalgebra of $A$. Suppose that, for every pair
of distinct points $t_1,t_2\in T$, there is an $x\in B$ such that
$x(t_1)=1_{t_1}$ and $x(t_2)=0$ (and, if $T$ consists of one point, suppose
$1\in B$). Then, for every continuous complex function $f$ on $T$, the
field $t\mapsto f(t)1_t$ belongs to $B$.
\end{itshape}

Let $t_1,t_2\in T$, $t_1\neq t_2$. There is an $x\in B$ with
$x(t_1)=1_{t_1}$ and $x(t_2)=0$. Replacing $x$ by $xx^{*}$, we may suppose
$x\in B^{+}$. Let $p:\mathbb R\to\mathbb R$ be continuous, equal to $0$ in
a neighborhood of $0$ and to $1$ on $[1,+\infty)$. Put $y=p(x)$. Then
$y(t_1)=1_{t_1}$, and $y(t)$ vanishes in a neighborhood of $t_2$.

Let $K$ be a nonempty compact subset of $T$ not containing $t_1$. There
are open sets $U_1,\ldots,U_n$ covering $K$ and elements
$y_1,\ldots,y_n\in B$ such that $y_i(t_1)=1_{t_1}$ and $y_i(t)=0$ on
$U_i$. Forming the product $y_1\cdots y_ny_n^{*}\cdots y_1^{*}$, we see
that there is a $z\in B^{+}$ such that $z(t_1)=1_{t_1}$ and $z(t)=0$ on
$K$. The same conclusion is still true, and evident, when $K$ is empty.
Near $t_1$ we have $\|z(t)-1_t\|\leq1/2$, and hence
$z(t)\geq\tfrac12 1_t$. Now let $L$ be a compact subset of $T$ disjoint
from $K$. There are open sets $V_1,\ldots,V_p$ covering $L$ and elements
$z_1,\ldots,z_p\in B^{+}$ such that $z_i(t)=0$ on $K$ and
$z_i(t)\geq\tfrac12 1_t$ on $V_i$. Thus
$u=2(z_1+\cdots+z_p)\in B^{+}$, with $u(t)=0$ on $K$ and
$u(t)\geq1_t$ on $L$.

Let $f$ be a continuous real function on $T$, and let $\varepsilon>0$.
Choose an open cover $(W_1,\ldots,W_q)$ of $T$ such that the oscillation of
$f$ on each $W_i$ is at most $\varepsilon$, and choose a value $\lambda_i$
of $f$ on $W_i$. Choose an open cover $(W'_1,\ldots,W'_q)$ of $T$ such
that $\overline{W'_i}\subset W_i$ for every $i$. There is a
$v_i\in B^{+}$ such that $v_i(t)\geq1_t$ on $W'_i$ and $v_i(t)=0$ on
$T-W_i$. Put $v=v_1+\cdots+v_q$. Then $v(t)\geq1_t$ for every $t$, so by
functional calculus (\XRef{10.3.3}{10.3.3}) there is a $v'\in B^{+}$ such
that $(v'^2v)(t)=1_t$ for every $t$. Put $w_i=v'v_iv'\in B^{+}$, so that
$\sum w_i=v'vv'=1$. Let $t\in T$; for example, suppose
\[
t\in W_1,\ldots,W_m,
\qquad t\notin W_{m+1},\ldots,W_q.
\]
Then
\begin{align*}
f(t)1_t-\sum_{i=1}^{q}\lambda_iw_i(t)
&=f(t)\sum_{i=1}^{m}w_i(t)-\sum_{i=1}^{m}\lambda_iw_i(t)\\
&=\sum_{i=1}^{m}[f(t)-\lambda_i]w_i(t),
\end{align*}
% Source PDF physical page 207
and hence
\[
-\varepsilon1_t=-\varepsilon\sum_{i=1}^{m}w_i(t)
\leq\sum_{i=1}^{m}[f(t)-\lambda_i]w_i(t)
\leq\varepsilon\sum_{i=1}^{m}w_i(t)=\varepsilon1_t.
\]
Therefore
\[
\left\|f(t)1_t-\sum_{i=1}^{q}\lambda_iw_i(t)\right\|\leq\varepsilon
\qquad\text{for every }t.
\]
Since $\varepsilon$ is arbitrary, the field $t\mapsto f(t)1_t$ belongs to
$\overline B=B$, and hence to $B$.

\ResultHeading{10.4.7}{10.4.7. Proposition.}
\begin{itshape}
Let $T$ be a locally compact space with a countable base, let $\mathcal A$
be a separable continuous field of $C^{*}$-algebras over $T$, and let $A$
be the $C^{*}$-algebra defined by $\mathcal A$. Then $A$ is separable.
\end{itshape}

By \XRef{10.2.6}{10.2.6}, it is enough to treat the case in which $T$ is
compact with a countable base. Let $(f_1,f_2,\ldots)$ be a dense sequence
in the Banach space of continuous complex functions on $T$. Put
$\mathcal A=((A(t))_{t\in T},\Theta)$, and let $(x_1,x_2,\ldots)$ be a
total sequence in $\Theta$. The $f_ix_j$ are elements of $A$. We shall
prove that their linear combinations are dense in $A$, which establishes
the proposition. Let $x\in A$ and $\varepsilon>0$. By
\XRef{10.2.5}{10.2.5}, there are continuous complex functions
$g_1,\ldots,g_p$ on $T$ such that
$\|x-g_1x_1-\cdots-g_px_p\|\leq\varepsilon$. Choosing indices
$i_1,\ldots,i_p$ suitably gives
$\|x-f_{i_1}x_1-\cdots-f_{i_p}x_p\|\leq2\varepsilon$.

Bibliography: \cite{ref29,ref77,ref108,ref167,ref169}.

\BookSubsection[Continuous Fields of C*-Algebras Defined by Certain
C*-Algebras]{10.5}{Continuous Fields of $C^{*}$-Algebras Defined by
Certain $C^{*}$-Algebras}

\ResultHeading{10.5.1}{10.5.1.}
Let $A$ be a $C^{*}$-algebra whose spectrum $T$ is assumed Hausdorff. Then
$T$ is locally compact (\XRef{3.3.8}{3.3.8}). Two irreducible
representations with the same kernel are equivalent (\XRef{3.1.6}{3.1.6}),
so the elements of $T$ are identified with the primitive ideals of $A$.
For every $t\in T$, put $A(t)=A/t$, a nonzero $C^{*}$-algebra. Every
$x\in A$ defines a vector field $t\mapsto x'(t)\in A(t)$ on $T$, where
$x'(t)$ is the canonical image of $x$ in $A(t)$. Let $\Lambda$ be the set
of the fields $t\mapsto x'(t)$, $x\in A$. Then $\Lambda$ is an involutive
subalgebra of the set of all vector fields. For every $x\in A$, the
function $t\mapsto\|x'(t)\|$ is continuous (\XRef{3.3.9}{3.3.9}); and for
every $t\in T$, the set of $x'(t)$, $x\in A$, is all of $A(t)$.

It follows from \XRef{10.2.3}{10.2.3} that there is a unique subset
$\Gamma\subset\prod_{t\in T}A(t)$ containing $\Lambda$ such that
$\mathcal A=((A(t))_{t\in T},\Gamma)$ is a continuous field of Banach
spaces; $\Gamma$ is the set of vector fields that are local uniform limits
of elements of $\Lambda$. By \XRef{10.3.2}{10.3.2}, $\mathcal A$ is a
continuous field of $C^{*}$-algebras. We call $\mathcal A$ the
\emph{continuous field of $C^{*}$-algebras defined by $A$}. If $A$ is
liminal, $\mathcal A$ is a continuous field of elementary
$C^{*}$-algebras. If $A$ is separable, $\mathcal A$ is separable.

% Source PDF physical page 208
\ResultHeading{10.5.2}{10.5.2. Theorem.}
\begin{itshape}
Let $T$ be locally compact, let
$\mathcal A=((A(t))_{t\in T},\Theta)$ be a continuous field of nonzero
elementary $C^{*}$-algebras over $T$, and let $A$ be the liminal
$C^{*}$-algebra defined by $\mathcal A$, whose spectrum is $T$. Let
$\mathcal A'=((A'(t))_{t\in T},\Theta')$ be the continuous field of nonzero
elementary $C^{*}$-algebras defined by $A$. For every $t\in T$, let
$\varphi(t)$ be the canonical isomorphism of $A(t)$ onto $A'(t)$. Then
$(\varphi(t))_{t\in T}$ is an isomorphism of $\mathcal A$ onto
$\mathcal A'$.
\end{itshape}

The elements of $A$ form a total family of vector fields continuous relative
to $\mathcal A$, and their images under $(\varphi(t))_{t\in T}$ belong to
$\Theta'$ by the definition of $\mathcal A'$. It remains only to apply
\XRef{10.2.4}{10.2.4}.

\ResultHeading{10.5.3}{10.5.3.}
The following lemma will be greatly generalized later
(\XRef{11.5.3}{11.5.3}).

\ResultLabel{Lemma.}
\begin{itshape}
Let $T$ be locally compact, let
$\mathcal A=((A(t))_{t\in T},\Theta)$ be a continuous field of elementary
$C^{*}$-algebras over $T$, let $A$ be the $C^{*}$-algebra defined by
$\mathcal A$, and let $B$ be a $C^{*}$-subalgebra of $A$ with the following
property: whenever $t_1,t_2\in T$, $\xi_1\in A(t_1)$, and
$\xi_2\in A(t_2)$, with $\xi_1=\xi_2$ if $t_1=t_2$, there is an $x\in B$
such that $x(t_1)=\xi_1$ and $x(t_2)=\xi_2$. Then $B=A$.
\end{itshape}

By adjoining a point at infinity to $T$, we reduce to the case in which
$T$ is compact (\XRef{10.2.6}{10.2.6}). Then $A=\Theta$, and $B$ is a
total subset of $\Theta$. In view of \XRef{10.2.5}{10.2.5}, it is enough to
prove that $\varphi x\in B$ whenever $x\in B$ and $\varphi$ is a continuous
complex function on $T$. Let $\varepsilon>0$. We shall construct an element
of $B$ approximating $\varphi x$ to within
$2\varepsilon[\sup|\varphi|^2]^{1/2}$, which will complete the proof. Put
$y=x^{*}x$.

Let $t_0\in T$. There are numbers $\alpha,\beta$ such that
$0<\alpha<\beta<\varepsilon$ and the spectrum of $y(t_0)$ does not meet
$[\alpha,\beta]$ [because $A(t_0)$ is elementary]. Choose
$\alpha',\beta'$ with
$\alpha<\alpha'<\beta'<\beta$. There is a neighborhood of $t_0$ such
that, for $t$ in this neighborhood, $\operatorname{Sp}'y(t)$ does not meet
$[\alpha',\beta']$ (\XRef{10.3.6}{10.3.6}). Let
$p:\mathbb R\to\mathbb R$ be a nonnegative continuous function equal to
$0$ on $(-\infty,\alpha']$ and to $1$ on
$[\beta',+\infty)$. Put $e=p(y)\in B^{+}$. Then, for every $t\in T$, $e(t)$
is an element of $A(t)^{+}$ with finite spectrum; and, throughout a closed
neighborhood $U$ of $t_0$, $e(t)$ is idempotent and
$\|e(t)y(t)-y(t)\|\leq\varepsilon$.

Let $((A(t))_{t\in U},\Delta)$ be the continuous field of
$C^{*}$-algebras induced by $\mathcal A$ on $U$. For $z\in\Delta$, the
following conditions are equivalent:
\begin{enumerate}[label=\textup{\arabic*\textdegree}]
\item $z(t)\in A'(t)=e(t)A(t)e(t)$ for every $t\in U$;
\item there is a $z'\in\Delta$ such that
$z(t)=e(t)z'(t)e(t)$ for every $t\in U$.
\end{enumerate}
Let $\Delta'$ be the set of $z\in\Delta$ satisfying these conditions. It is
immediate that $\mathcal A'=((A'(t))_{t\in U},\Delta')$ is a continuous
field of $C^{*}$-algebras over $U$; for every $t\in U$, $e(t)$ is the
identity of $A'(t)$. Let $A'$ be the $C^{*}$-algebra defined by
$\mathcal A'$. Let $B'$ be the set of elements of $A'$ that are
restrictions to $U$ of elements of $B$. The map $x\mapsto x|U$ is a
morphism
% Source PDF physical page 209
from the $C^{*}$-algebra $A=\Theta$ onto the $C^{*}$-algebra $\Delta$
(\XRef{10.1.12}{10.1.12}); the image of $B$ under this morphism is a
$C^{*}$-subalgebra of $\Delta$, and hence $B'$ is a $C^{*}$-subalgebra of
$A'$. By the hypothesis on $B$, whenever $t_1,t_2\in U$,
$\xi_1\in A'(t_1)$, and $\xi_2\in A'(t_2)$, with $\xi_1=\xi_2$ if
$t_1=t_2$, there is an $x_1\in B$ such that
$x_1(t_1)=\xi_1$ and $x_1(t_2)=\xi_2$. Thus the restriction of $ex_1e$ to
$U$ is an element of $B'$ taking the values $\xi_1,\xi_2$ at $t_1,t_2$.
We may therefore apply \XRef{10.4.6}{10.4.6}: for every continuous complex
function $\psi'$ on $U$, the field $t\mapsto\psi'(t)e(t)$ on $U$ belongs
to $B'$. Consequently, for every continuous complex function $\psi$ on
$T$, the field $\psi e$ agrees on $U$ with an element of $B$.

Let $V$ be a closed subset of $T$ contained in the interior $\mathring U$ of
$U$. We show that there is an $f\in B$ such that $f(t)=e(t)$ on $V$ and
$f(t)=0$ on $T-\mathring U$. Let $t_1\in T-\mathring U$ and $t_2\in V$.
Since $e(t_1)$ has finite spectrum, the hypothesis on $B$ shows that there
is a Hermitian $d\in B$ such that $d(t_1)$ is idempotent and
\[
d(t_1)e(t_1)=e(t_1),
\qquad d(t_2)=0.
\]
Replacing $d$ by $q(d)$, where $q$ vanishes near $0$ and equals $1$ at
$1$, we may suppose that $d(t)$ vanishes near $t_2$. Covering $V$ by
finitely many such neighborhoods and multiplying the corresponding
elements $d$, we obtain a $d'\in B$ such that
\[
d'(t_1)e(t_1)=e(t_1),
\qquad d'(t)=0\quad\text{on }V.
\]
Then $d''=e-d'e\in B$ agrees with $e$ on $V$ and vanishes at $t_1$. By
functional calculus, we obtain $d'''\in B$ that agrees with $e$ on $V$ and
vanishes in a neighborhood of $t_1$. Covering $T-\mathring U$ by finitely
many such neighborhoods and multiplying the corresponding elements
$d'''$, we obtain $f$.

Let $\theta$ be a continuous complex function on $T$ vanishing outside
$V$. We have seen that there is an element of $B$ taking the value
$\theta(t)e(t)$ on $U$. Multiplying it by the preceding element $f$ gives
an element of $B$ taking the value $\theta(t)e(t)$ everywhere on $T$.

In summary, changing the notation slightly, we can associate with every
$t\in T$ an $e_t\in B^{+}$ and a closed neighborhood $U_t$ of $t$ with the
following properties:
\begin{enumerate}[label=\textup{\arabic*\textdegree}]
\item $\|e_t(t')y(t')-y(t')\|\leq\varepsilon$ on $U_t$;
\item if $\theta$ is a continuous complex function on $T$ whose support is
contained in $\mathring U_t$, then $\theta e_t\in B$.
\end{enumerate}
Let $(U_{t_1},\ldots,U_{t_r})$ be a finite cover of $T$. Choose an open
cover $(V_1,\ldots,V_r)$ of $T$ such that
$\overline V_1\subset\mathring U_{t_1},\ldots,
\overline V_r\subset\mathring U_{t_r}$, and let
$(\theta_1,\ldots,\theta_r)$ be a continuous partition of unity subordinate
to $(V_1,\ldots,V_r)$. Then $\varphi\theta_i e_{t_i}\in B$. Put
\[
z=\varphi x(\theta_1e_{t_1}+\cdots+\theta_re_{t_r})\in B.
\]
% ===== ASSEMBLED TRANSLATION CHUNK 08: chunk_210_239.tex =====
% Source PDF physical page 210
Then
\[
(z-\varphi x)^{*}(z-\varphi x)
=|\varphi|^{2}\sum_{i,j}\theta_i\theta_j
\bigl(e_{t_i}ye_{t_j}-e_{t_i}y-ye_{t_j}+y\bigr).
\]
Since
\[
\|e_{t_i}(t')y(t')-y(t')\|\leq\varepsilon
\qquad\text{on }U_{t_i},
\]
we have
\[
\|\bigl(e_{t_i}ye_{t_j}-e_{t_i}y-ye_{t_j}+y\bigr)(t')\|
\leq2\varepsilon
\qquad\text{on }U_{t_i},
\]
and therefore
\[
\|z-\varphi x\|^{2}\leq2\varepsilon\,[\sup|\varphi|^{2}].
\]

\ResultHeading{10.5.4}{10.5.4. Theorem.}
\begin{itshape}
Let $A$ be a liminal $C^{*}$-algebra whose spectrum $T$ is Hausdorff.
Let $\mathcal A=((A(t))_{t\in T},\Theta)$ be the continuous field of
nonzero elementary $C^{*}$-algebras over $T$ defined by $A$, and let $A'$
be the $C^{*}$-algebra defined by $\mathcal A$. For every $x\in A$, let
$x'$ be the element of $\Theta$ defined by $x$. Then $x'\in A'$, and
$x\mapsto x'$ is an isomorphism of $A$ onto $A'$.
\end{itshape}

If $x\in A$, the function $t\mapsto\|x'(t)\|$ tends to zero at infinity
on $T$ (\XRef{3.3.7}{3.3.7}), and hence $x'\in A'$. It is clear that
$x\mapsto x'$ is a morphism $\pi$ of $A$ into $A'$. By
\XRef{2.7.3}{2.7.3}, this morphism is isometric. If $t\in T$, the set of
all $x'(t)$, $x\in A$, is the whole of $A(t)$ by the definition of
$A(t)$. If $t_1,t_2$ are distinct points of $T$, and if
$\xi_1\in A(t_1)$ and $\xi_2\in A(t_2)$, there is an $x\in A$ such that
$x'(t_1)=\xi_1$ and $x'(t_2)=\xi_2$ (\XRef{4.2.5}{4.2.5}). Thus
$\pi(A)=A'$ by \XRef{10.5.3}{10.5.3}.

\ResultHeading{10.5.5}{10.5.5.}
Theorems \XRef{10.5.2}{10.5.2} and \XRef{10.5.4}{10.5.4} show that there
is a canonical one-to-one correspondence between liminal
$C^{*}$-algebras with Hausdorff spectrum (up to isomorphism) and
continuous fields of nonzero elementary $C^{*}$-algebras over locally
compact spaces (up to isomorphism, the notion of isomorphism used here
allowing a homeomorphism of the base space).

Let $A$ be a liminal $C^{*}$-algebra with Hausdorff spectrum, and let
$\mathcal A$ be the field defined by $A$. If $A$ is separable,
$\widehat A$ has a countable base (\XRef{3.3.4}{3.3.4}) and $\mathcal A$
is separable. The converse follows from \XRef{10.4.7}{10.4.7}.

\ResultHeading{10.5.6}{10.5.6. Corollary.}
\begin{itshape}
Let $A$ be a liminal $C^{*}$-algebra such that $\widehat A$ is Hausdorff,
let $x\in A$, and let $f$ be a bounded continuous complex function on
$\widehat A$. There is a $y\in A$ such that
$\pi(y)=f(\pi)\pi(x)$ for every $\pi\in\widehat A$.
\end{itshape}

With the notation of \XRef{10.5.4}{10.5.4}, we may replace $A$ by $A'$;
the corollary then follows from \XRef{10.1.9}{10.1.9(ii)}.

\ResultHeading{10.5.7}{10.5.7. Definition.}
\begin{itshape}
Let $T$ be a topological space and let
$\mathcal A=((A(t))_{t\in T},\Theta)$ be a continuous field of elementary
$C^{*}$-algebras over $T$. We say that $\mathcal A$ satisfies
\emph{Fell's condition} if, for every $t_0\in T$, there are a neighborhood
$V$ of $t_0$ and a vector field $p$ of $\mathcal A$, defined and continuous
on $V$, such that $p(t)$ is a rank-one projection for every $t\in V$
(cf.\ \XRef{4.1.9}{4.1.9}).
\end{itshape}

% Source PDF physical page 211
It is clear that a locally trivial continuous field of nonzero elementary
$C^{*}$-algebras satisfies Fell's condition. The converse, however, is
false.

\ResultHeading{10.5.8}{10.5.8. Proposition.}
\begin{itshape}
Let $A$ be a liminal $C^{*}$-algebra with Hausdorff spectrum, and let
$\mathcal A$ be the continuous field of nonzero elementary
$C^{*}$-algebras defined by $A$. Then $A$ has continuous trace if and only
if $\mathcal A$ satisfies Fell's condition.
\end{itshape}

If $A$ has continuous trace, $\mathcal A$ satisfies Fell's condition
[\XRef{4.5.3}{4.5.3(iii)}]. Conversely, suppose that $\mathcal A$
satisfies Fell's condition, and let $\pi_0\in\widehat A$. There are a
compact neighborhood $V$ of $\pi_0$ and a vector field $p$ of
$\mathcal A$, defined and continuous on $V$, such that $p(t)$ is a
rank-one projection for every $t\in V$. Let $f:T\to\mathbb R$ be a
continuous nonnegative function equal to $1$ on a neighborhood of
$\pi_0$ and to $0$ on a neighborhood of $T-V$. The vector field equal to
$0$ on $T-V$ and to $fp$ on $V$ is continuous and tends to zero at
infinity. It comes from an element $y\in A^{+}$ such that $\pi(y)$ is a
rank-one projection in a neighborhood of $\pi_0$. Thus $A$ has continuous
trace (\XRef{4.5.4}{4.5.4}).

Bibliography: \cite{ref19,ref29,ref77,ref108,ref167,ref169}.

\BookSubsection[Remarks on Elementary C*-Algebras]{10.6}{Remarks on
Elementary $C^{*}$-Algebras}

\ResultHeading{10.6.1}{10.6.1.}
Let $H$ be a Hilbert space. For $\xi,\eta\in H$, throughout the remainder
of this subsection we denote by $\theta_{\xi,\eta}$ the operator
$\zeta\mapsto(\zeta\mid\eta)\xi$ on $H$. It has rank at most $1$, and
every operator of rank at most $1$ is of this form. If
$\xi',\eta'\in H$, then
\begin{align*}
\theta_{\xi,\eta}\theta_{\xi',\eta'}(\zeta)
&=\theta_{\xi,\eta}\bigl((\zeta\mid\eta')\xi'\bigr)\\
&=(\zeta\mid\eta')(\xi'\mid\eta)\xi
=(\xi'\mid\eta)\theta_{\xi,\eta'}(\zeta),
\end{align*}
and hence
\[
\theta_{\xi,\eta}\theta_{\xi',\eta'}
=(\xi'\mid\eta)\theta_{\xi,\eta'}.
\]
On the other hand, if $\zeta'\in H$, then
\begin{align*}
(\theta_{\xi,\eta}^{*}(\zeta)\mid\zeta')
&=(\zeta\mid\theta_{\xi,\eta}(\zeta'))
=(\zeta\mid(\zeta'\mid\eta)\xi)\\
&=(\eta\mid\zeta')(\zeta\mid\xi)
=((\zeta\mid\xi)\eta\mid\zeta')
=(\theta_{\eta,\xi}(\zeta)\mid\zeta'),
\end{align*}
so
\[
\theta_{\xi,\eta}^{*}=\theta_{\eta,\xi}.
\]
We have $\theta_{\xi,\xi}=\|\xi\|^{2}P_{\mathbb C\xi}$; consequently
every finite-rank operator on $H$ is a linear combination of operators
$\theta_{\xi,\eta}$ (and even of operators $\theta_{\xi,\xi}$).

\ResultHeading{10.6.2}{10.6.2.}
Let $\xi_1,\ldots,\xi_{2n}$ be variable vectors in $H$. If
$\xi_1\to\eta_1,\ldots,\xi_{2n}\to\eta_{2n}$, it is clear that
\[
\theta_{\xi_1,\xi_2}+\cdots+\theta_{\xi_{2n-1},\xi_{2n}}
\longrightarrow
\theta_{\eta_1,\eta_2}+\cdots+\theta_{\eta_{2n-1},\eta_{2n}}
\]
uniformly. This being so, Lemma \XRef{3.5.6}{3.5.6} shows that, if
$(\xi_i\mid\xi_j)\to(\eta_i\mid\eta_j)$ for all $i,j$, then
\[
\|\theta_{\xi_1,\xi_2}+\cdots+
\theta_{\xi_{2n-1},\xi_{2n}}\|
\longrightarrow
\|\theta_{\eta_1,\eta_2}+\cdots+
\theta_{\eta_{2n-1},\eta_{2n}}\|.
\]
In other words,
% Source PDF physical page 212
\[
\|\theta_{\xi_1,\xi_2}+\cdots+
\theta_{\xi_{2n-1},\xi_{2n}}\|
\]
is a continuous function of the scalar products
$(\xi_i\mid\xi_j)$.

\ResultHeading{10.6.3}{10.6.3.}
If $H$ is a Hilbert space with a distinguished unit vector $\xi$, put
\[
\alpha(H,\xi)=(A,p),
\]
where $A$ is the elementary $C^{*}$-algebra $\mathcal{LC}(H)$ and
$p\in A$ is the projection onto $\mathbb C\xi$. Similarly, let
$(A',p')=\alpha(H',\xi')$. Every isomorphism
$U:(H,\xi)\to(H',\xi')$ (that is, every isomorphism of $H$ onto $H'$
carrying $\xi$ to $\xi'$) defines an isomorphism
\[
\alpha(U):(A,p)\longrightarrow(A',p')
\]
(that is, an isomorphism of $A$ onto $A'$ carrying $p$ to $p'$). If
$U_1,U_2$ are two such isomorphisms and
$\alpha(U_1)=\alpha(U_2)$, then $U_1=U_2$: indeed, by
\XRef{4.1.8}{4.1.8}, $U_1$ and $U_2$ differ by a scalar factor, and that
factor is $1$ because $U_1\xi=\xi'=U_2\xi$.

\ResultHeading{10.6.4}{10.6.4.}
Retain the preceding notation $H,\xi,A,p$. If $a\in Ap$, then
$a^{*}a\in pAp$; hence $a^{*}a$ leaves $\mathbb C\xi$ invariant and
vanishes on $H\ominus\mathbb C\xi$, so
\[
\|a^{*}a\|=(a^{*}a\xi\mid\xi)=\|a\xi\|^{2},
\qquad\text{whence}\qquad
\|a\|=\|a\xi\|.
\]
Thus $a\mapsto a\xi$ is an isometric isomorphism of the normed vector
space $Ap$ onto the normed vector space $H$.

It follows that $Ap$ may be regarded as a Hilbert subspace of $A$; it is
in fact well determined, since scalar products in a Hilbert space can be
computed from norms. If $a,b\in Ap$, then $b^{*}a\in pAp$, and therefore
\[
\operatorname{Tr}(b^{*}a)=(b^{*}a\xi\mid\xi)=(a\xi\mid b\xi).
\]
Thus the scalar product in the Hilbert space $Ap$ is given by
$(a\mid b)=\operatorname{Tr}(b^{*}a)$.

Notice also that $ApA$ is the set of elements of $A$ of rank at most $1$.
This follows from
\[
\theta_{\zeta,\eta}
=\theta_{\zeta,\xi}\theta_{\xi,\xi}\theta_{\xi,\eta}
\qquad(\|\xi\|=1).
\]

\ResultHeading{10.6.5}{10.6.5.}
If $A$ is an elementary $C^{*}$-algebra with a rank-one projection $p$,
put
\[
\beta(A,p)=(K,\eta),
\]
where $K$ is the Hilbert space $Ap$ (\XRef{10.6.4}{10.6.4}) and $\eta$
is the unit vector $p$ of $K$. Similarly, let
$(K',\eta')=\beta(A',p')$. Every isomorphism
$V:(A,p)\to(A',p')$ defines an isomorphism
\[
\beta(V):(K,\eta)\longrightarrow(K',\eta').
\]
If $V_1,V_2$ are two such isomorphisms and
$\beta(V_1)=\beta(V_2)$, then $V_1=V_2$: indeed, $V_1$ and $V_2$
coincide on $Ap$, hence on $pA=(Ap)^{*}$, hence on
$ApA=(Ap)(pA)$, and consequently on $A$ by \XRef{10.6.4}{10.6.4}.

\ResultHeading{10.6.6}{10.6.6. Lemma.}
\begin{itshape}
\begin{enumerate}[label=\textup{(\roman*)}]
\item Let $H$ be a Hilbert space and $\xi$ a unit vector of $H$. Put
$\alpha(H,\xi)=(A,p)$, so that
$\beta(\alpha(H,\xi))=(Ap,p)$. For every $a\in Ap$, put
$\varphi(a)=a\xi$. Then $\varphi$ is an isomorphism of the Hilbert space
$Ap$ onto the Hilbert space $H$, and $\varphi(p)=\xi$.

\item Let $A$ be an elementary $C^{*}$-algebra and $p$ a rank-one
projection of $A$. Put $\beta(A,p)=(H,\xi)$ and
$\alpha(H,\xi)=(A',p')$. For every $a\in A$, let $\psi(a)$ be the linear
map of $H=Ap$ into itself defined by $\psi(a)y=ay$. Then $\psi$ is an
isomorphism of the $C^{*}$-algebra $A$ onto the $C^{*}$-algebra $A'$,
and $\psi(p)=p'$.
\end{enumerate}
\end{itshape}

Part (i) was proved in \XRef{10.6.4}{10.6.4}.

% Source PDF physical page 213
Adopt the notation of (ii). There are a Hilbert space $H_0$ and a unit
vector $\xi_0\in H_0$ such that $(A,p)$ is identified with
$\alpha(H_0,\xi_0)$. Let $\varphi_0$ be the isomorphism of $(H,\xi)$
onto $(H_0,\xi_0)$ defined in (i). For every $a\in A$ and every $b\in H$,
\[
\varphi_0(\psi(a)b)=\varphi_0(ab)=(ab)\xi_0
=a(\varphi_0(b)).
\]
Thus $\psi(a)=\varphi_0^{-1}a\varphi_0$, and $\psi$ is indeed an
isomorphism of $A$ onto $A'$ carrying
$P_{\mathbb C\xi_0}=p$ to $P_{\mathbb C\xi}=p'$. This proves the lemma.

The isomorphisms $\varphi$ and $\psi$ will be called \emph{canonical}.

\ResultHeading{10.6.7}{10.6.7.}
Adopt the notation of \XRef{10.6.6}{10.6.6(ii)}. If $a,a'\in A$, then
$ap\in H$ and $a'^{*}p\in H$, so we may form
$\theta_{ap,a'^{*}p}$, which is an element of $A'$. We have
\[
\psi(apa')=\theta_{ap,a'^{*}p}.
\]
Indeed, for every $b\in A$,
\[
\psi(apa')(bp)=apa'bp=a[\operatorname{Tr}(pa'bp)]p,
\]
and by \XRef{10.6.4}{10.6.4} this is equal to
\[
(bp\mid a'^{*}p)ap=\theta_{ap,a'^{*}p}(bp).
\]

\ResultHeading{10.6.8}{10.6.8.}
Let $H,H_1$ be Hilbert spaces, and let $\xi,\xi_1$ be unit vectors of
$H,H_1$, respectively. Put
\[
\begin{array}{ll}
(A,p)=\alpha(H,\xi),& (H',\xi')=\beta(A,p),\\
(A_1,p_1)=\alpha(H_1,\xi_1),&
(H_1',\xi_1')=\beta(A_1,p_1).
\end{array}
\]
Let
\[
\varphi:(H',\xi')\longrightarrow(H,\xi),
\qquad
\varphi_1:(H_1',\xi_1')\longrightarrow(H_1,\xi_1)
\]
be the canonical isomorphisms. Let $U$ be an isomorphism of $(H,\xi)$
onto $(H_1,\xi_1)$. Then
$U'=\beta(\alpha(U))$ is an isomorphism of $(H',\xi')$ onto
$(H_1',\xi_1')$, and
\[
U\varphi=\varphi_1U'.
\]
Indeed, if $a\in H'=Ap$, then
\[
(U\varphi)(a)=U(a\xi)=UaU^{-1}\xi_1
=(\alpha(U)a)\xi_1
=\varphi_1(\alpha(U)(a))=\varphi_1U'(a).
\]
There is an analogous result if one starts with $(A,p),(A_1,p_1)$ and an
isomorphism of $(A,p)$ onto $(A_1,p_1)$.

Bibliography: \cite{ref19}.

\BookSubsection[Continuous Field of Elementary C*-Algebras Defined by a
Continuous Field of Hilbert Spaces]{10.7}{Continuous Field of Elementary
$C^{*}$-Algebras Defined by a Continuous Field of Hilbert Spaces}

\ResultHeading{10.7.1}{10.7.1.}
Let $T$ be a topological space, and let
$\mathcal H=((H(t))_{t\in T},\Gamma)$ be a continuous field of Hilbert
spaces over $T$. If $x,y\in\Gamma$, the function
$t\mapsto(x(t)\mid y(t))$ is continuous because
\begin{align*}
4(x(t)\mid y(t))
={}&\|x(t)+y(t)\|^{2}-\|x(t)-y(t)\|^{2}\\
&+i\|x(t)+iy(t)\|^{2}-i\|x(t)-iy(t)\|^{2}.
\end{align*}

% Source PDF physical page 214
\ResultHeading{10.7.2}{10.7.2.}
For every $t\in T$, let $A(t)=\mathcal{LC}(H(t))$. For $x,y\in\Gamma$,
define $\theta_{x,y}\in\prod_{t\in T}A(t)$ by
$\theta_{x,y}(t)=\theta_{x(t),y(t)}$. By \XRef{10.6.1}{10.6.1},
\[
\theta_{x,y}^{*}=\theta_{y,x},
\qquad
\theta_{x,y}\theta_{x',y'}=\theta_{z,y'},
\]
where $z(t)=(x'(t)\mid y(t))x(t)$. It follows that the set $\Lambda$ of
fields
\[
\theta_{x_1,x_2}+\theta_{x_3,x_4}+\cdots+
\theta_{x_{2n-1},x_{2n}},
\qquad x_1,x_2,\ldots,x_{2n}\in\Gamma,
\]
is an involutive subalgebra of $\prod_{t\in T}A(t)$. By
\XRef{10.6.1}{10.6.1}, $\Lambda(t)$ is the set of finite-rank operators
on $H(t)$ and is therefore dense in $A(t)$. By \XRef{10.6.2}{10.6.2},
\[
\|\theta_{x_1,x_2}(t)+\cdots+
\theta_{x_{2n-1},x_{2n}}(t)\|
\]
is a continuous function of the $(x_i(t)\mid x_j(t))$, and hence of $t$.
There is therefore (\XRef{10.2.3}{10.2.3}) a unique set
$\Theta\subset\prod_{t\in T}A(t)$ containing $\Lambda$ such that
$\mathcal A=((A(t))_{t\in T},\Theta)$ is a continuous field of Banach
spaces. Since $\Lambda$ is an involutive algebra, $\mathcal A$ is in fact
a continuous field of elementary $C^{*}$-algebras
(\XRef{10.3.2}{10.3.2}), called the field \emph{associated with}
$\mathcal H$. It will be denoted by $\mathcal A(\mathcal H)$. By
\XRef{10.2.4}{10.2.4}, it is immediate that, if $Y\subset T$, then
\[
\mathcal A(\mathcal H|Y)=\mathcal A(\mathcal H)|Y.
\]

\ResultHeading{10.7.3}{10.7.3.}
If $a\in\Theta$ and $x\in\Gamma$, the field
$ax:t\mapsto a(t)x(t)$ is an element of $\Gamma$. Indeed, since every
element of $\Theta$ is a local uniform limit of elements of $\Lambda$
(\XRef{10.2.3}{10.2.3}), it is enough to verify this for $a\in\Lambda$,
and hence for $a=\theta_{y,z}$, where $y,z\in\Gamma$; but
$\theta_{y,z}x$ is the field
$t\mapsto(x(t)\mid z(t))y(t)$.

\ResultHeading{10.7.4}{10.7.4.}
Let $(x_\alpha)$ be a family of elements of $\Gamma$ such that, for every
$t\in T$, the $x_\alpha(t)$ are dense in $H(t)$. Then, for every
$t\in T$, the $\theta_{x_\alpha,x_\beta}(t)$ are dense in the set of
operators of rank at most $1$ on $H(t)$. It follows that, if $\mathcal H$
is separable, $\mathcal A(\mathcal H)$ is separable.

\ResultHeading{10.7.5}{10.7.5.}
To study the correspondence between continuous fields of Hilbert spaces
and continuous fields of elementary $C^{*}$-algebras, it will be
convenient to make the following definitions.

Let $\mathcal H$ be a continuous field of Hilbert spaces over $T$, endowed
with a continuous vector field $x$ such that $\|x(t)\|=1$ for every $t$.
Put
\[
\alpha(\mathcal H,x)=(\mathcal A,p),
\]
where $\mathcal A=\mathcal A(\mathcal H)$ and $p=\theta_{x,x}$ (so that
$p$ is a continuous field, relative to $\mathcal A$, of rank-one
projections). Every isomorphism
$\varphi:(\mathcal H,x)\to(\mathcal H',x')$ defines, in the evident way,
an isomorphism
\[
\alpha(\varphi):\alpha(\mathcal H,x)
\longrightarrow\alpha(\mathcal H',x').
\]
It follows from \XRef{10.6.3}{10.6.3} that
$\alpha(\varphi_1)=\alpha(\varphi_2)$ implies
$\varphi_1=\varphi_2$.

% Source PDF physical page 215
Let $\mathcal A$ be a continuous field of elementary $C^{*}$-algebras
over $T$, endowed with a continuous field $p$ of rank-one projections.
Put
\[
\beta(\mathcal A,p)=(\mathcal H,x),
\]
where $\mathcal H$ and $x$ are defined as follows. Write
$\mathcal A=((A(t))_{t\in T},\Theta)$. Let $H(t)$ be the vector subspace
$A(t)p(t)$ of $A(t)$; this is a Hilbert space (\XRef{10.6.4}{10.6.4}).
Let $\Gamma$ be the set of $x\in\Theta$ such that $x(t)\in H(t)$ for every
$t\in T$. It is clear that
$\mathcal H=((H(t))_{t\in T},\Gamma)$ satisfies axioms (i), (iii), and
(iv) of \XRef{10.1.2}{10.1.2}. Moreover, if $a\in\Theta$, then $ap$ is
an element of $\Gamma$, so that
$\Gamma(t)=A(t)p(t)=H(t)$. Thus $\mathcal H$ is a continuous field of
Hilbert spaces over $T$. We put $x=p$; then $x\in\Gamma$ and
$\|x(t)\|=1$ for every $t$.

If $\mathcal A$ is separable, then $\mathcal H$ is separable. Indeed, if
$(a_1,a_2,\ldots)$ is a sequence in $\Theta$ such that the $a_i(t)$ are
dense in $A(t)$ for every $t$, then $(a_1p,a_2p,\ldots)$ is a sequence in
$\Gamma$ such that the $(a_ip)(t)$ are dense in $H(t)$ for every $t$.
Every isomorphism
$\varphi:(\mathcal A,p)\to(\mathcal A',p')$ defines, in the evident way,
an isomorphism
\[
\beta(\varphi):\beta(\mathcal A,p)
\longrightarrow\beta(\mathcal A',p').
\]
It follows from \XRef{10.6.5}{10.6.5} that
$\beta(\varphi_1)=\beta(\varphi_2)$ implies $\varphi_1=\varphi_2$.

\ResultHeading{10.7.6}{10.7.6. Lemma.}
\begin{itshape}
\begin{enumerate}[label=\textup{(\roman*)}]
\item Let $\mathcal H$ be a continuous field of Hilbert spaces over $T$,
endowed with a continuous vector field $x$ such that $\|x(t)\|=1$ for
every $t$. Let
$(\mathcal H',x')=\beta(\alpha(\mathcal H,x))$. Write
$\mathcal H=((H(t))_{t\in T},\Gamma)$ and
$\mathcal H'=((H'(t))_{t\in T},\Gamma')$. For every $t\in T$, let
$\varphi_t$ be the canonical isomorphism of $H'(t)$ onto $H(t)$
(\XRef{10.6.6}{10.6.6}). Then
$\varphi=(\varphi_t)_{t\in T}$ is an isomorphism of $\mathcal H'$ onto
$\mathcal H$ such that $\varphi(x')=x$.

\item Let $\mathcal A$ be a continuous field of elementary
$C^{*}$-algebras over $T$, endowed with a continuous field $p$ of
rank-one projections. Let
$(\mathcal A',p')=\alpha(\beta(\mathcal A,p))$. Write
$\mathcal A=((A(t))_{t\in T},\Theta)$ and
$\mathcal A'=((A'(t))_{t\in T},\Theta')$. For every $t\in T$, let
$\psi_t$ be the canonical isomorphism of $A(t)$ onto $A'(t)$
(\XRef{10.6.6}{10.6.6}). Then $\psi=(\psi_t)_{t\in T}$ is an isomorphism
of $\mathcal A$ onto $\mathcal A'$ such that $\psi(p)=p'$.
\end{enumerate}
\end{itshape}

Adopt the notation of (i). It is clear that $\varphi(x')=x$. Let
$y\in\Gamma$ and $y'=\theta_{y,x}\in\Gamma'$. Then
\[
\varphi_t(y'(t))=\theta_{y,x}(t)x(t)=y(t).
\]
Thus $\varphi^{-1}(y)=y'\in\Gamma'$, which proves that $\varphi^{-1}$ is
an isomorphism of $\mathcal H$ onto $\mathcal H'$
(\XRef{10.2.4}{10.2.4}).

Adopt the notation of (ii). It is clear that $\psi(p)=p'$. To show that
$\psi$ is an isomorphism of $\mathcal A$ onto $\mathcal A'$, it is enough
to show that $\psi(a)\in\Theta'$ for every $a$ of the form
$a'pa''$, with $a',a''\in\Theta$, since fields of this form constitute a
total subset of $\Theta$ (\XRef{10.6.4}{10.6.4}). But, if $a=a'pa''$,
then
$\psi(a)=\theta_{a'p,a''^{*}p}\in\Theta'$ by
\XRef{10.6.7}{10.6.7}.

% Source PDF physical page 216
\ResultHeading{10.7.7}{10.7.7. Proposition.}
\begin{itshape}
Let $T$ be a topological space and $\mathcal A$ a continuous field of
elementary $C^{*}$-algebras over $T$. The following conditions are
equivalent:
\begin{enumerate}[label=\textup{(\roman*)}]
\item for every $t_0\in T$, there are a neighborhood $V$ of $t_0$ and a
continuous field $\mathcal H$ of nonzero Hilbert spaces over $V$ such
that $\mathcal A|V$ is isomorphic to $\mathcal A(\mathcal H)$;
\item $\mathcal A$ satisfies Fell's condition.
\end{enumerate}
\end{itshape}

(i)$\Rightarrow$(ii). Suppose (i) holds. Let $t_0\in T$. Let $V_0$ be a
neighborhood of $t_0$ and let
$\mathcal H=((H(t))_{t\in V_0},\Gamma)$ be a continuous field of nonzero
Hilbert spaces over $V_0$ such that $\mathcal A|V_0$ is isomorphic to
$\mathcal A(\mathcal H)$. Choose a nonzero $\xi\in H(t_0)$, and a
continuous vector field $x$ of $\mathcal H$ such that $x(t_0)=\xi$. The
set $V$ of $t\in V_0$ for which $x(t)\ne0$ is a neighborhood of $t_0$;
put $y(t)=\|x(t)\|^{-1}x(t)$ for $t\in V$. The vector field
$\theta_{y,y}$ of $\mathcal A(\mathcal H)$ is defined and continuous on
$V$, and $\theta_{y,y}(t)$ is a rank-one projection for every $t\in V$.
There is therefore a vector field $p$ of $\mathcal A$, defined and
continuous on $V$, whose values $p(t)$ are rank-one projections.

(ii)$\Rightarrow$(i). Suppose that $\mathcal A$ satisfies Fell's
condition, and let $t_0\in T$. There are a neighborhood $V$ of $t_0$ and
a continuous field $p$ of rank-one projections of $\mathcal A$, defined
on $V$. Put $\beta(\mathcal A|V,p)=(\mathcal H,x)$. By
\XRef{10.7.6}{10.7.6}, $\mathcal A|V$ is isomorphic to
$\mathcal A(\mathcal H)$.

\ResultHeading{10.7.8}{10.7.8.}
Proposition \XRef{10.7.7}{10.7.7} is a strictly local result. Even if
$\mathcal A$ satisfies Fell's condition, in general there is no continuous
field $\mathcal H$ of Hilbert spaces over $T$ such that $\mathcal A$ is
isomorphic to $\mathcal A(\mathcal H)$. The existence of $\mathcal H$ is
obstructed by a cohomology class, which we shall now specify.

\ResultHeading{10.7.9}{10.7.9. Lemma.}
\begin{itshape}
Let $\mathcal H$ and $\mathcal H'$ be two continuous fields of nonzero
Hilbert spaces over $T$, let $\mathcal A$ and $\mathcal A'$ be their
associated fields of $C^{*}$-algebras, and let $f$ be an isomorphism of
$\mathcal A$ onto $\mathcal A'$. For every $t\in T$, there are a
neighborhood $V$ of $t$ and an isomorphism of $\mathcal H|V$ onto
$\mathcal H'|V$ that induces $f|V$.
\end{itshape}

There are a neighborhood $V$ of $t$ and a continuous field $p$ of
rank-one projections of $\mathcal A|V$ (\XRef{10.7.7}{10.7.7}). Then
$p'=f(p)$ is a continuous field of rank-one projections of
$\mathcal A'|V$. Choose a continuous vector field $x$ of $\mathcal H$
such that $p(t)x(t)\ne0$. In a neighborhood of $t$ we have
$p(t')x(t')\ne0$. Replacing $x$ by the field
\[
t'\longmapsto\|p(t')x(t')\|^{-1}p(t')x(t'),
\]
we obtain, in a neighborhood of $t$, a continuous vector field $y$ of
$\mathcal H$ such that $\|y(t')\|=1$ and $p$ is the field of projections
defined by $y$. Similarly, there is a continuous vector field $y'$ of
$\mathcal H'$ in a neighborhood of $t$ such that $\|y'(t')\|=1$ and
$p'$ is the field of projections defined by $y'$. Since the assertion to
be proved is local, we may henceforth suppose that $y,y',p,p'$ are all
defined on the whole of $T$. Then $\beta(f)$ is an isomorphism
% Source PDF physical page 217
of
\[
\beta(\mathcal A,p)=\beta(\alpha(\mathcal H,y))
\quad\text{onto}\quad
\beta(\mathcal A',p')=\beta(\alpha(\mathcal H',y')).
\]
By \XRef{10.7.6}{10.7.6} and \XRef{10.6.8}{10.6.8}, there is an
isomorphism $g$ of $(\mathcal H,y)$ onto $(\mathcal H',y')$ such that
$\beta(f)=\beta(\alpha(g))$. Hence $f=\alpha(g)$
(\XRef{10.7.5}{10.7.5}).

\ResultHeading{10.7.10}{10.7.10. Lemma.}
\begin{itshape}
Let $\mathcal H=((H(t))_{t\in T},\Gamma)$ be a continuous field of Hilbert
spaces over $T$, and let $\mathcal A=\mathcal A(\mathcal H)$. An
automorphism $\varphi$ of $\mathcal H$ induces the identity automorphism
of $\mathcal A$ if and only if it has the form
$(u(t)1_t)_{t\in T}$, where $1_t$ denotes the identity map of $H(t)$ and
$u$ is a continuous map of $T$ into $\mathbb U$, the group of complex
numbers of modulus $1$.
\end{itshape}

The condition is plainly sufficient. Conversely, if $\varphi$ induces
the identity automorphism of $\mathcal A$, then every $\varphi_t$ is a
homothety with ratio $u(t)\in\mathbb U$ (\XRef{4.1.8}{4.1.8}), and it
remains to prove that $u$ is continuous. For every $t_0\in T$ there is a
continuous vector field $x$ of $\mathcal H$ such that $x(t)\ne0$ on a
neighborhood $V$ of $t_0$. On $V$,
\[
u(t)=(\varphi_tx(t)\mid x(t))\|x(t)\|^{-2},
\]
so $u$ is continuous at $t_0$.

\ResultHeading{10.7.11}{10.7.11. Lemma.}
\begin{itshape}
Let $T$ be paracompact and let $\mathcal A$ be a continuous field of
elementary $C^{*}$-algebras over $T$ satisfying Fell's condition. There
exist:
\begin{enumerate}[label=\textup{(\roman*)}]
\item an open cover $(T_i)_{i\in I}$ of $T$;
\item for each $i\in I$, a continuous field $\mathcal H_i$ of Hilbert
spaces over $T_i$ and an isomorphism $h_i$ of
$\mathcal A(\mathcal H_i)$ onto $\mathcal A|T_i$;
\item for all $i,j\in I$, an isomorphism $g_{ij}$ of
$\mathcal H_j|T_{ij}$ onto $\mathcal H_i|T_{ij}$ that induces the
isomorphism $h_i^{-1}h_j$ of
$\mathcal A(\mathcal H_j|T_{ij})$ onto
$\mathcal A(\mathcal H_i|T_{ij})$.
\end{enumerate}
Here $T_i\cap T_j=T_{ij}$, and below similarly
$T_i\cap T_j\cap T_k=T_{ijk}$, and so forth.
\end{itshape}

There are an open cover $(U_\alpha)_{\alpha\in\Lambda}$ of $T$ and, for
every $\alpha\in\Lambda$, a continuous field $\mathcal K_\alpha$ of
Hilbert spaces over $U_\alpha$ and an isomorphism $k_\alpha$ of
$\mathcal A(\mathcal K_\alpha)$ onto $\mathcal A|U_\alpha$
(\XRef{10.7.7}{10.7.7}). Since $T$ is paracompact, we may suppose that
$(U_\alpha)$ is locally finite. There is an open cover
$(V_\alpha)_{\alpha\in\Lambda}$ of $T$ such that
$\overline V_\alpha\subset U_\alpha$ for every $\alpha$. Fix provisionally
$\alpha\in\Lambda$ and $t\in V_\alpha$. There are an open neighborhood
$W\subset V_\alpha$ of $t$ and a finite subset $\Lambda'$ of $\Lambda$
such that $W\cap\overline V_\beta=\varnothing$ for
$\beta\notin\Lambda'$. Let $\beta\in\Lambda'$. There are an open
neighborhood $W_\beta\subset W$ of $t$ and an isomorphism of
$\mathcal K_\beta|W_\beta\cap V_\beta$ onto
$\mathcal K_\alpha|W_\beta\cap V_\beta$ that induces
$k_\alpha^{-1}k_\beta$ on $W_\beta\cap V_\beta$. This is clear if
$t\notin\overline V_\beta$; if $t\in\overline V_\beta$, then
$t\in U_\beta\cap U_\alpha$, and the assertion follows from
\XRef{10.7.9}{10.7.9}. Since $\Lambda'$ is finite, we obtain an open
neighborhood $W_{\alpha,t}\subset V_\alpha$ of $t$ and, for every
$\beta\in\Lambda$, an isomorphism of
$\mathcal K_\beta|W_{\alpha,t}\cap V_\beta$ onto
$\mathcal K_\alpha|W_{\alpha,t}\cap V_\beta$ inducing
$k_\alpha^{-1}k_\beta$ there. Let the index set $I$ in the lemma be the
set of pairs $i=(\alpha,t)$
% Source PDF physical page 218
with $\alpha\in\Lambda$ and $t\in V_\alpha$, and put
\[
T_i=W_{\alpha,t},\qquad
\mathcal H_i=\mathcal K_\alpha|W_{\alpha,t},\qquad
h_i=k_\alpha|W_{\alpha,t}.
\]
The existence of the $g_{ij}$ follows at once from the preceding
construction.

In the remainder of this subsection we shall use the language of sheaves.
For this subject the reader may consult R. Godement, \emph{Th\'eorie des
faisceaux} (Act. Sci. Ind., no. 1252, Hermann, Paris, 1958).

\ResultHeading{10.7.12}{10.7.12. Lemma.}
\begin{itshape}
Retain the notation of \XRef{10.7.11}{10.7.11}.
\begin{enumerate}[label=\textup{(\roman*)}]
\item For all $i,j,k\in I$, there is a unique continuous function
$u_{ijk}:T_{ijk}\to\mathbb U$ such that
\[
g_{ij}g_{jk}=u_{ijk}g_{ik}.
\]
\item The $u_{ijk}$ define a $2$-cocycle of $(T_i)$ with values in the
sheaf $\mathcal U$ of germs of continuous maps of $T$ into $\mathbb U$.
\item The image $\gamma$ of this cocycle in the cohomology group
$H^{2}(T,\mathcal U)$ depends only on $\mathcal A$.
\end{enumerate}
\end{itshape}

If $i,j,k\in I$, then $g_{ij}g_{jk}g_{ik}^{-1}$ induces on
$\mathcal A(\mathcal H_i|T_{ijk})$ the automorphism
\[
h_i^{-1}h_jh_j^{-1}h_kh_k^{-1}h_i=1.
\]
Thus (i) follows from \XRef{10.7.10}{10.7.10}.

If $i,j,k,l\in I$, then on $T_{ijkl}$,
\begin{align*}
u_{jkl}u_{ikl}^{-1}u_{ijl}
&=(g_{jl}^{-1}g_{jk}g_{kl})
  (g_{kl}^{-1}g_{ik}^{-1}g_{il})
  (g_{il}^{-1}g_{ij}g_{jl})\\
&=g_{jl}^{-1}(g_{jk}g_{ik}^{-1})g_{ij}g_{jl}\\
&=g_{jl}^{-1}(u_{ijk}g_{ij}^{-1})g_{ij}g_{jl}
=u_{ijk},
\end{align*}
which proves (ii).

We prove (iii). Let $(T_\lambda)_{\lambda\in\Lambda}$,
$\mathcal K_\lambda,h_\lambda,g_{\lambda\mu}$ be data analogous to those
of \XRef{10.7.11}{10.7.11}, and let $(u_{\lambda\mu\nu})$ be the
corresponding cocycle. To prove that $(u_{ijk})$ and
$(u_{\lambda\mu\nu})$ define the same element of
$H^{2}(T,\mathcal U)$, we may replace $(T_i)$ by a refinement, replacing
$\mathcal H_i,h_i,g_{ij}$ by their restrictions to the members of the
new cover, and do the same for $(T_\lambda)$. This allows us to suppose,
as in the proof of \XRef{10.7.11}{10.7.11}, that for every
$i\in I$ and every $\lambda\in\Lambda$ there is an isomorphism
$g_{i\lambda}$ of $\mathcal H_i|T_{i\lambda}$ onto
$\mathcal K_\lambda|T_{i\lambda}$ inducing the isomorphism
$h_\lambda^{-1}h_i$ of the associated fields of $C^{*}$-algebras. Let
$\Lambda'$ be the disjoint union of $I$ and $\Lambda$. The preceding data
define a cover $(T_a)_{a\in\Lambda'}$, fields $\mathcal H_a$,
isomorphisms $h_a$, and transition maps $g_{ab}$ satisfying the
properties of \XRef{10.7.11}{10.7.11}, and hence functions
$u_{abc}$, $a,b,c\in\Lambda'$. The class in $H^{2}(T,\mathcal U)$ defined
by $(u_{abc})$ is at once the class defined by $(u_{ijk})$ and the class
defined by $(u_{\lambda\mu\nu})$. This proves the lemma.

\ResultHeading{10.7.13}{10.7.13.}
Let $\mathcal R$ be the sheaf of germs of continuous maps from $T$ into
$\mathbb R$. The map $t\mapsto\exp(2\pi it)$ of $\mathbb R$ onto
$\mathbb U$ defines a morphism $\mathcal R\to\mathcal U$. Its kernel is
the constant sheaf $\mathbb Z$. Since every continuous function
$g:T\to\mathbb U$ is locally of the form $\exp(2\pi if)$, where $f$ is a
continuous map of $T$ into $\mathbb R$, we have an exact sequence of
sheaves
\[
0\longrightarrow\mathbb Z\longrightarrow\mathcal R
\longrightarrow\mathcal U\longrightarrow1,
\]
% Source PDF physical page 219
and hence an exact cohomology sequence
\[
\cdots\longrightarrow H^{n}(T,\mathbb Z)
\longrightarrow H^{n}(T,\mathcal R)
\longrightarrow H^{n}(T,\mathcal U)
\longrightarrow H^{n+1}(T,\mathbb Z)\longrightarrow\cdots.
\]
The existence of continuous partitions of unity on $T$ shows that
$\mathcal R$ is a fine sheaf, so $H^{n}(T,\mathcal R)=0$ for $n>0$. It
follows in particular that the canonical morphism
\[
H^{2}(T,\mathcal U)\longrightarrow H^{3}(T,\mathbb Z)
\]
is an isomorphism.

\ResultHeading{10.7.14}{10.7.14. Definition.}
\begin{itshape}
Let $T$ be paracompact and let $\mathcal A$ be a continuous field of
elementary $C^{*}$-algebras over $T$ satisfying Fell's condition. Let
$\gamma$ be the element of $H^{2}(T,\mathcal U)$ defined in Lemma
\XRef{10.7.12}{10.7.12}. Its canonical image in
$H^{3}(T,\mathbb Z)$ will be denoted by $\delta(\mathcal A)$.
\end{itshape}

\ResultHeading{10.7.15}{10.7.15. Theorem.}
\begin{itshape}
Let $T$ be paracompact and let $\mathcal A$ be a continuous field of
elementary $C^{*}$-algebras over $T$ satisfying Fell's condition. There
exists a continuous field $\mathcal H$ of Hilbert spaces over $T$ such
that $\mathcal A$ is isomorphic to $\mathcal A(\mathcal H)$ if and only if
$\delta(\mathcal A)=0$.
\end{itshape}

Suppose first that $\mathcal A$ is isomorphic to
$\mathcal A(\mathcal H)$. In constructing $\delta(\mathcal A)$, we may
take as the cover of $T$ the cover consisting of the single open set $T$.
Every $2$-cocycle relative to this cover is cohomologous to zero, and so
$\delta(\mathcal A)=0$.

Conversely, suppose that $\delta(\mathcal A)=0$. There are
$T_i,\mathcal H_i,h_i,g_{ij},u_{ijk}$, with $i,j,k\in I$, having the
properties of \XRef{10.7.11}{10.7.11} and \XRef{10.7.12}{10.7.12}, and
continuous functions $v_{ij}:T_{ij}\to\mathbb U$ such that
\[
u_{ijk}=v_{jk}v_{ik}^{-1}v_{ij}
\qquad(i,j,k\in I).
\]
Put $g'_{ij}=v_{ij}^{-1}g_{ij}$. The isomorphisms
$g'_{ij}:\mathcal H_j|T_{ij}\to\mathcal H_i|T_{ij}$ satisfy
\[
g'_{ij}g'_{jk}
=v_{ij}^{-1}g_{ij}g_{jk}v_{jk}^{-1}
=u_{ijk}v_{ij}^{-1}v_{jk}^{-1}g_{ik}=g'_{ik}.
\]
We may therefore define a continuous field $\mathcal H$ of Hilbert spaces
over $T$ by gluing the $\mathcal H_i$ by means of the $g'_{ij}$
(\XRef{10.1.13}{10.1.13}). In other words, we may henceforth suppose that
$\mathcal H_i=\mathcal H|T_i$ and that the $g'_{ij}$ are the identity
automorphisms. Then
$h_i:\mathcal A(\mathcal H)|T_i\to\mathcal A|T_i$ and
$h_j:\mathcal A(\mathcal H)|T_j\to\mathcal A|T_j$ coincide on $T_{ij}$.
Thus the $h_i$ define an isomorphism of $\mathcal A(\mathcal H)$ onto
$\mathcal A$.

Bibliography: \cite{ref19,ref29}.

\BookSubsection[Locally Trivial Fields of Elementary C*-Algebras]{10.8}{Locally
Trivial Fields of Elementary $C^{*}$-Algebras}

\ResultHeading{10.8.1}{10.8.1. Lemma.}
\begin{itshape}
Let $H$ be an infinite-dimensional Hilbert space. There exist:
\begin{enumerate}[label=\textup{\alph*.}]
\item a unit vector $\xi\in H$;
\item for every $t\in[0,1]$, a closed vector subspace $H_t$ of $H$;
\item for every $t\in(0,1]$, a linear isometric map $U_t$ of $H_t$ onto
$H$;
\end{enumerate}
with the following properties:
\begin{enumerate}[label=\textup{(\roman*)}]
\item the projection $P_t=P_{H_t}$ is a strongly continuous function of
$t\in[0,1]$, and $H_1=H$, $H_0=\{0\}$;
% Source PDF physical page 220
\item the operators $U_tP_t$ and $U_t^{-1}$ are strongly continuous
functions of $t\in(0,1]$, and $U_1=1$;
\item $\|P_t\xi\|^{2}=t$.
\end{enumerate}
\end{itshape}

Suppose first that $\dim H=\aleph_0$, so that $H$ may be identified with
$L^{2}([0,1])$. Let $H_t$ be the set of $f\in L^{2}([0,1])$ such that
$f(x)=0$ for $x\geq t$. If $t\in(0,1]$ and $f\in H_t$, define
$U_tf\in H$ by
\[
(U_tf)(x)=\sqrt t\,f(tx).
\]
Finally, take for $\xi$ the constant function equal to $1$. Properties
(i) and (iii) are immediate. For $f\in H_t$,
\[
\|U_tf\|^{2}=\int_0^1 t|f(tx)|^{2}\,dx
=\int_0^t|f(x')|^{2}\,dx'=\|f\|^{2},
\]
and it is clear that $U_t$ is a linear isometric map of $H_t$ onto $H$;
if $g\in H$, then, for $0\leq x\leq t$,
\[
(U_t^{-1}g)(x)=\frac1{\sqrt t}\,g\!\left(\frac xt\right).
\]
For $f\in H$, the expression
\[
\|U_tP_tf-U_{t'}P_{t'}f\|^{2}
=\int_0^1\left|\sqrt t\,f(tx)
-\sqrt{t'}\,f(t'x)\right|^{2}\,dx
\]
tends to zero as $t'\to t$. On the other hand,
\begin{align*}
\|U_t^{-1}f-U_{t'}^{-1}f\|^{2}
={}&2\|f\|^{2}-2\operatorname{Re}(U_t^{-1}f\mid U_{t'}^{-1}f)\\
={}&2\|f\|^{2}-2\operatorname{Re}
\int_0^{\min(t,t')}\frac1{\sqrt t}f\!\left(\frac xt\right)
\frac1{\sqrt{t'}}\overline{f\!\left(\frac x{t'}\right)}\,dx.
\end{align*}
If, for example, $t\leq t'$, this gives
\[
\|U_t^{-1}f-U_{t'}^{-1}f\|^{2}
=2\|f\|^{2}-2\operatorname{Re}\int_0^1
\sqrt{\frac t{t'}}\,f(x')
\overline{f\!\left(\frac t{t'}x'\right)}\,dx',
\]
and this tends to zero as $t'\to t$. This proves (ii).

If now $\dim H$ is an arbitrary infinite cardinal, we may identify $H$
with a Hilbert space $K\otimes K'$, where $K,K'$ are Hilbert spaces and
$\dim K=\aleph_0$. There are in $K$ operators $P_t,U_t$ and a vector
$\xi$ having properties (i), (ii), and (iii); it is enough to consider on
$H$ the operators $P_t\otimes1$, $U_t\otimes1$, and the vector
$\xi\otimes\xi'$, where $\xi'$ is a unit vector of $K'$.

\ResultHeading{10.8.2}{10.8.2. Lemma.}
\begin{itshape}
Let $H$ be an infinite-dimensional Hilbert space, let $G$ be the unitary
group of $H$ with the strong topology, and let $S$ be the unit sphere of
$H$ with the norm topology. Then $G$ and $S$ are contractible.
\end{itshape}

Let $H_t,P_t,U_t,\xi$ have the properties of \XRef{10.8.1}{10.8.1}.
For $U\in G$ and $t\in(0,1]$, put
% Source PDF physical page 221
\[
\Phi(U,t)=(1-P_t)+U_t^{-1}UU_tP_t.
\]
The operator $\Phi(U,t)$ induces the identity map on $H\ominus H_t$ and,
on $H_t$, the unitary operator whose restriction is $U_t^{-1}UU_t$.
Thus $\Phi(U,t)\in G$, and $\Phi(U,1)=U$. Put $\Phi(U,0)=1$. We show that
$\Phi:G\times[0,1]\to G$ is continuous. On $G\times(0,1]$ this follows
at once from Lemma \XRef{10.8.1}{10.8.1}. Moreover, if
$(U',t')\in G\times(0,1]$ tends to $(U,0)\in G\times[0,1]$, then $P_{t'}$
converges strongly to zero and, for every $f\in H$,
\[
\|U_{t'}^{-1}U'U_{t'}P_{t'}f\|=\|P_{t'}f\|\longrightarrow0.
\]
Hence $\Phi(U',t')$ converges strongly to $1=\Phi(U,0)$. Thus $G$ is
contractible.

For $f\in S$ and $t\in(0,1]$, put
\[
\Psi(f,t)=\sqrt t\,U_t^{-1}f+(1-P_t)\xi.
\]
We have
\[
\|\Psi(f,t)\|^{2}
=t\|U_t^{-1}f\|^{2}+\|(1-P_t)\xi\|^{2}
=t+1-t=1.
\]
Thus $\Psi(f,t)\in S$, and $\Psi(f,1)=f$. Put $\Psi(f,0)=\xi$. The map
$\Psi:S\times[0,1]\to S$ is continuous on $S\times(0,1]$. Moreover, if
$(f',t')\in S\times(0,1]$ tends to $(f,0)$, then
$\|\sqrt{t'}U_{t'}^{-1}f'\|\to0$, and hence
$\Psi(f',t')\to\xi=\Psi(f,0)$. Thus $S$ is contractible.

\ResultHeading{10.8.3}{10.8.3.}
A continuous field $((A(t))_{t\in T},\Theta)$ of elementary
$C^{*}$-algebras is said to have \emph{rank $a$} if every $A(t)$ has rank
$a$ (\XRef{4.1.9}{4.1.9}).

\ResultHeading{10.8.4}{10.8.4. Theorem.}
\begin{itshape}
Let $T$ be a paracompact space and $a$ an infinite cardinal. The map
$\mathcal A\mapsto\delta(\mathcal A)$ defines a one-to-one correspondence
between $H^{3}(T,\mathbb Z)$ and the locally trivial continuous fields of
elementary $C^{*}$-algebras of rank $a$ over $T$, up to isomorphism.
\end{itshape}

Choose a Hilbert space $H$ of dimension $a$. Let $G$ be its unitary group,
let $\mathcal G$ be the sheaf of germs of continuous maps of $T$ into
$G$, and let $\mathcal U$ be the sheaf of germs of continuous maps of $T$
into $\mathbb U$.

$1^\circ$ Let $\mathcal A,\mathcal A'$ be two locally trivial continuous
fields of elementary $C^{*}$-algebras of rank $a$ over $T$. Suppose that
$\delta(\mathcal A)=\delta(\mathcal A')$. We shall prove that
$\mathcal A$ and $\mathcal A'$ are isomorphic.

There exist: (a) a locally finite open cover $(T_i)_{i\in I}$ of $T$;
(b) for every $i\in I$, continuous fields $\mathcal H_i,\mathcal H_i'$ of
Hilbert spaces over $T_i$, and isomorphisms
\[
h_i:\mathcal A(\mathcal H_i)\longrightarrow\mathcal A|T_i,
\qquad
h_i':\mathcal A(\mathcal H_i')\longrightarrow\mathcal A'|T_i;
\]
(c) for all $i,j\in I$, isomorphisms
\[
g_{ij}:\mathcal H_j|T_{ij}\longrightarrow\mathcal H_i|T_{ij},
\qquad
g'_{ij}:\mathcal H_j'|T_{ij}\longrightarrow\mathcal H_i'|T_{ij}
\]
% Source PDF physical page 222
that induce the isomorphisms $h_i^{-1}h_j$ and
$h_i'^{-1}h_j'$, respectively. If $u_{ijk},u'_{ijk}:T_{ijk}\to\mathbb U$
are defined by
\[
g_{ij}g_{jk}=u_{ijk}g_{ik},
\qquad
g'_{ij}g'_{jk}=u'_{ijk}g'_{ik},
\]
then $(u_{ijk})$ and $(u'_{ijk})$ both represent the classes
$\delta(\mathcal A)$ and $\delta(\mathcal A')$, respectively. Replacing
$(T_i)$ by a refinement, we may therefore suppose that there is a
$1$-cochain $(v_{ij})$ with values in $\mathcal U$ such that
\[
u'_{ijk}=u_{ijk}v_{ij}v_{jk}v_{ik}^{-1}.
\]
Replacing the $g'_{ij}$ by $v_{ij}^{-1}g'_{ij}$, we reduce to the case
$u'_{ijk}=u_{ijk}$.

Since $\mathcal A$ and $\mathcal A'$ are locally trivial, we may suppose,
after refining $(T_i)$ once more, that $\mathcal H_i$ and
$\mathcal H_i'$ are both the constant field over $T_i$ defined by $H$.
Then $g_{ij}$ and $g'_{ij}$ are continuous maps of $T_{ij}$ into $G$
with its strong topology, and hence define continuous sections
$\gamma_{ij},\gamma'_{ij}$ of $\mathcal G$ over $T_{ij}$. Choose an open
cover $(S_i)_{i\in I}$ of $T$ such that $\overline S_i\subset T_i$ for
every $i\in I$. Let $L$ be the set of pairs $(J,\beta)$, where
$J\subset I$, $\beta=(\beta_i)_{i\in J}$, and, for every $i\in J$,
$\beta_i$ is a continuous section of $\mathcal G$ over
$\overline S_i$, such that
\[
\gamma'_{ij}=\beta_i\gamma_{ij}\beta_j^{-1}
\quad\text{on }\overline S_i\cap\overline S_j
\qquad(i,j\in J).
\]
Order $L$ inductively, and let $(J,\beta)$ be a maximal element. We prove
that $J=I$. Suppose that $i\in I-J$, and put
\[
R=\overline S_i\cap\bigcup_{j\in J}\overline S_j.
\]
If $t\in R$, and $j_1,j_2\in J$ are such that
$t\in\overline S_{j_1}\cap\overline S_{j_2}$, then
\begin{align*}
\gamma'_{ij_2}(t)\beta_{j_2}(t)\gamma_{ij_2}(t)^{-1}
&=u_{ij_1j_2}(t)^{-1}\gamma'_{ij_1}(t)
  \gamma'_{j_1j_2}(t)\beta_{j_2}(t)
  \gamma_{j_1j_2}(t)^{-1}\gamma_{ij_1}(t)^{-1}
  u_{ij_1j_2}(t)\\
&=\gamma'_{ij_1}(t)
  \bigl(\gamma'_{j_1j_2}(t)\beta_{j_2}(t)
  \gamma_{j_1j_2}(t)^{-1}\bigr)\gamma_{ij_1}(t)^{-1}\\
&=\gamma'_{ij_1}(t)\beta_{j_1}(t)\gamma_{ij_1}(t)^{-1}.
\end{align*}
It follows that one can define a continuous section $\zeta$ of
$\mathcal G$ over $R$ such that
$\gamma'_{ij}=\zeta\gamma_{ij}\beta_j^{-1}$ for every $j\in J$. Since
$G$ is contractible (\XRef{10.8.2}{10.8.2}), $\mathcal G$ is soft; hence
$\zeta$ extends to a continuous section $\beta_i$ of $\mathcal G$ over
$\overline S_i$. This contradicts the maximality of $(J,\beta)$. Thus
$J=I$.

Each $\beta_i$ now defines an automorphism $a_i$ of
$\mathcal A(\mathcal H_i|S_i)$, and
\[
h_i'^{-1}h_j'=a_i h_i^{-1}h_j a_j^{-1}
\quad\text{on }S_{ij},
\]
or equivalently
\[
h_i'a_ih_i^{-1}=h_j'a_jh_j^{-1}
\quad\text{on }S_{ij}.
\]
Thus the $h_i'a_ih_i^{-1}$ define an isomorphism of $\mathcal A$ onto
$\mathcal A'$.

$2^\circ$ Let $(T_i)_{i\in I}$ be a locally finite open cover of $T$, and
let $(u_{ijk})$ be an alternating $2$-cocycle of this cover with values in
$\mathcal U$. We shall construct
% Source PDF physical page 223
a locally trivial continuous field $\mathcal A$ of elementary
$C^{*}$-algebras of rank $a$ over $T$ such that $\delta(\mathcal A)$ is
the element of $H^{3}(T,\mathbb Z)$ defined by $(u_{ijk})$. This will
complete the proof of the theorem.

Choose an open cover $(T_i')_{i\in I}$ such that
$\overline{T_i'}\subset T_i$. Let $L$ be the set of pairs $(J,\gamma)$,
where $J\subset I$, $\gamma=(\gamma_{ij})_{i,j\in J}$, and for all
$i,j\in J$, $\gamma_{ij}$ is a continuous section of $\mathcal G$ over
$\overline{T_i'}\cap\overline{T_j'}$, satisfying
\[
\gamma_{ij}\gamma_{jk}=\gamma_{ik}u_{ijk}
\qquad(i,j,k\in J).
\]
Order $L$ inductively, and let $(J,\gamma)$ be a maximal element. We show
that $J=I$.

Suppose that $i\in I-J$. Let $M$ be the set of pairs $(K,\gamma')$, where
$K\subset J$, $\gamma'=(\gamma'_{ij})_{j\in K}$, and, for every $j\in K$,
$\gamma'_{ij}$ is a continuous section of $\mathcal G$ over
$\overline{T_i'}\cap\overline{T_j'}$, such that
\[
\gamma'_{ij}\gamma_{jk}=\gamma'_{ik}u_{ijk}
\qquad(j,k\in K).
\]
Let $(K,\gamma')$ be a maximal element of $M$. We prove that $K=J$. If
$j\in J-K$, one can define a section $\beta$ of $\mathcal G$ over
\[
\overline{T_i'}\cap\overline{T_j'}\cap
\bigcup_{k\in K}\overline{T_k'}
\]
such that
\[
\beta\gamma_{jk}=\gamma'_{ik}u_{ijk}
\qquad(k\in K).
\]
Indeed, if $k,l\in K$, define on the corresponding intersections
$\beta_k,\beta_l$ by
\[
\beta_k\gamma_{jk}=\gamma'_{ik}u_{ijk},
\qquad
\beta_l\gamma_{jl}=\gamma'_{il}u_{ijl}.
\]
On $\overline{T_i'}\cap\overline{T_j'}\cap
\overline{T_k'}\cap\overline{T_l'}$ one obtains
\begin{align*}
\beta_l
&=\gamma'_{il}u_{ijl}\gamma_{jl}^{-1}\\
&=\gamma'_{ik}\gamma_{kl}u_{ikl}^{-1}u_{ijl}
  \gamma_{kl}^{-1}\gamma_{jk}^{-1}u_{jkl}\\
&=\gamma'_{ik}\gamma_{jk}^{-1}
  u_{ijl}u_{ikl}^{-1}u_{jkl}
=\gamma'_{ik}\gamma_{jk}^{-1}u_{ijk}=\beta_k.
\end{align*}
Since $\mathcal G$ is soft, $\beta$ extends to a continuous section of
$\mathcal G$ over $\overline{T_i'}\cap\overline{T_j'}$, contradicting the
maximality of $(K,\gamma')$. Thus $K=J$. Consequently there is a family
$(\gamma'_{ij})_{j\in J}$ such that each $\gamma'_{ij}$ is a continuous
section of $\mathcal G$ over
$\overline{T_i'}\cap\overline{T_j'}$ and
\[
\gamma'_{ij}\gamma_{jk}=\gamma'_{ik}u_{ijk}
\qquad(j,k\in J).
\]
The sections $\gamma_{ij}$ and $\gamma'_{ij}$ then contradict the
maximality of $(J,\gamma)$. Hence $J=I$.

Let $\mathcal H_i$ be the constant field over $T_i'$ defined by $H$, and
put $\mathcal A_i=\mathcal A(\mathcal H_i)$. For $i,j\in I$, the section
$\gamma_{ij}$ defines an isomorphism
$f_{ij}:\mathcal A_j|T_{ij}'\to\mathcal A_i|T_{ij}'$, and
$f_{ij}f_{jk}=f_{ik}$. There is therefore a continuous field
$\mathcal A$ of elementary $C^{*}$-algebras over $T$ and, for every
$i\in I$, an isomorphism $h_i$ of $\mathcal A_i$ onto
$\mathcal A|T_i'$ such that $f_{ij}=h_i^{-1}h_j$
(\XRef{10.1.13}{10.1.13}). To compute $\delta(\mathcal A)$, we may use
the cover $(T_i')$, the $\mathcal H_i$, the $h_i$, and the restrictions
of the $\gamma_{ij}$ to $T_{ij}'$. It follows that
$\delta(\mathcal A)$ is the element of $H^{3}(T,\mathbb Z)$ defined by
$(u_{ijk})$.

\ResultHeading{10.8.5}{10.8.5.}
We shall now see that the local-triviality hypothesis made in Theorem
\XRef{10.8.4}{10.8.4} is sometimes automatically satisfied.

\ResultLabel{Lemma.}
\begin{itshape}
Let $T$ be a topological space, let
$\mathcal H=((H(t))_{t\in T},\Gamma)$ be a separable continuous field of
Hilbert spaces over $T$, and let $H_0$ be a Hilbert space of dimension
% Source PDF physical page 224
$\aleph_0$. For every $t\in T$ there exist a closed vector subspace
$K(t)$ of $H_0$ and an isomorphism $U(t)$ of the Hilbert space $H(t)$
onto the Hilbert space $K(t)$, with the following properties:
\begin{enumerate}[label=\textup{\alph*.}]
\item if $\Delta$ is the set of continuous maps $x:T\to H_0$ such that
$x(t)\in K(t)$ for every $t\in T$, then
$((K(t))_{t\in T},\Delta)$ is a continuous field $\mathcal K$ of Hilbert
spaces over $T$;
\item $(U(t))_{t\in T}$ is an isomorphism of $\mathcal H$ onto
$\mathcal K$.
\end{enumerate}
\end{itshape}

Let $(x_1,x_2,\ldots)$ be a total subset of $\Gamma$. Multiplying the
$x_i$ by suitable continuous functions, we may suppose that
\[
\|x_i(t)\|\leq\frac1i
\qquad\text{for every }i\text{ and every }t.
\]
Let $(\varepsilon_1,\varepsilon_2,\ldots)$ be an orthonormal basis of
$H_0$. For every $t\in T$, there is a continuous linear map
$S(t):H_0\to H(t)$ such that $S(t)\varepsilon_i=x_i(t)$ for every $i$.
We have
\[
\|S(t)\|^{2}\leq\operatorname{Tr}(S(t)^{*}S(t))
\leq\sum_{i=1}^{\infty}i^{-2},
\]
and $S(t)(H_0)$ is dense in $H(t)$. Let
$t\mapsto x(t)=\sum_{i=1}^{\infty}\lambda_i(t)\varepsilon_i$ be a
continuous map of $T$ into $H_0$. Then
\[
t\longmapsto S(t)x(t)=\sum_{i=1}^{\infty}\lambda_i(t)x_i(t)
\]
belongs to $\Gamma$, since the $\lambda_i$ are continuous functions and
\[
\left\|\sum_{i=n+1}^{\infty}\lambda_i(t)x_i(t)\right\|
\leq\sum_{i=n+1}^{\infty}|\lambda_i(t)|\|x_i(t)\|
\leq\|x(t)\|\left(\sum_{i=n+1}^{\infty}i^{-2}\right)^{1/2}.
\]
We show that $t\mapsto S(t)^{*}S(t)x(t)$ is a continuous map of $T$ into
$H_0$. It is enough to prove this when $x(t)=\varepsilon_i$ for every
$t$. For every $z\in H_0$,
\begin{align*}
\|S(t)^{*}S(t)\varepsilon_i-z\|^{2}
={}&\|S(t)^{*}S(t)\varepsilon_i\|^{2}+\|z\|^{2}
-2\operatorname{Re}(S(t)^{*}S(t)\varepsilon_i\mid z)\\
={}&\sum_{j=1}^{\infty}|(S(t)^{*}S(t)\varepsilon_i\mid\varepsilon_j)|^{2}
+\|z\|^{2}-2\operatorname{Re}(S(t)^{*}S(t)\varepsilon_i\mid z)\\
={}&\sum_{j=1}^{\infty}|(x_i(t)\mid x_j(t))|^{2}
+\|z\|^{2}-2\operatorname{Re}(x_i(t)\mid S(t)z).
\end{align*}
This is a continuous function of $t$ by what precedes, and it follows
that $t\mapsto S(t)^{*}S(t)\varepsilon_i$ is continuous.

% Source PDF physical page 225
Let $z\in\Gamma$. For every $t_0\in T$ and every $\varepsilon>0$, there
is a continuous map $x:T\to H_0$ such that
\[
\|S(t_0)x(t_0)-z(t_0)\|<\varepsilon,
\]
and hence, in a neighborhood of $t_0$,
\[
\|S(t)^{*}S(t)x(t)-S(t)^{*}z(t)\|
<\left(\sum_{i=1}^{\infty}i^{-2}\right)^{1/2}\varepsilon.
\]
This proves that $t\mapsto S(t)^{*}z(t)$ is continuous. Let $K(t)$ be the
closure of $S(t)^{*}(H(t))$ in $H_0$. What precedes proves that the set
$\Delta$ in the statement of the lemma satisfies axiom (ii) of
\XRef{10.1.2}{10.1.2}; axioms (i), (iii), and (iv) are evident, and this
establishes (a). The map $t\mapsto S(t)^{*}S(t)$ from $T$ into
$\mathcal L(H_0)$ is strongly continuous, and therefore so is
\[
t\longmapsto(S(t)^{*}S(t))^{1/2}=|S(t)|.
\]
Let $S(t)=V(t)|S(t)|$ be the polar decomposition of $S(t)$. Since
$|S(t)|(H_0)$ and $S(t)^{*}(H(t))$ are dense in $K(t)$, while
$S(t)(H_0)$ is dense in $H(t)$, we see that $V(t)|K(t)$ is an isomorphism
of $K(t)$ onto $H(t)$. Thus $U(t)=V(t)^{-1}$ carries a total family in
$\Delta$ onto a total family in $\Gamma$, and (b) follows.

\ResultHeading{10.8.6}{10.8.6.}
We recall a result of E. Michael (\emph{Ann. of Math.}, vol. 64, 1956,
pp. 562--580, Theorem 1.2). Let $T$ be a paracompact space of finite
dimension (one says that $T$ has dimension at most $n$ if every finite
open cover of $T$ has a finite open refinement of order at most $n$).
Let $A$ be a closed subset of $T$, let $Y$ be a complete metric space,
and let $\varphi$ be a map of $T$ into the set of nonempty closed subsets
of $Y$. Suppose that $\varphi$ is lower semicontinuous, that is, whenever
$\varphi(t_0)$ meets an open subset $Y'$ of $Y$, $\varphi(t)$ meets $Y'$
for all $t$ sufficiently near $t_0$. Suppose that $\varphi(t)$ is
contractible for every $t\in T$. Finally, suppose that there is an
$\varepsilon_0>0$ such that, for every $t\in T$ and every open ball
$\beta$ in $Y$ of radius at most $\varepsilon_0$,
$\varphi(t)\cap\beta$ is contractible. Then every continuous map
$f:A\to Y$ such that $f(t)\in\varphi(t)$ for every $t\in A$ extends to a
continuous map $f':T\to Y$ such that $f'(t)\in\varphi(t)$ for every
$t\in T$.

\ResultHeading{10.8.7}{10.8.7. Lemma.}
\begin{itshape}
Let $T$ be a paracompact space of finite dimension, and let
$\mathcal H=((H(t))_{t\in T},\Gamma)$ be a separable continuous field of
Hilbert spaces over $T$, each $H(t)$ having dimension $\aleph_0$. Then
$\mathcal H$ is trivial.
\end{itshape}

Let $H_0$ be a Hilbert space of dimension $\aleph_0$. By
\XRef{10.8.5}{10.8.5}, we may suppose that every $H(t)$ is a closed vector
subspace of $H_0$ and that $\Gamma$ is the set of continuous maps
$x:T\to H_0$ such that $x(t)\in H(t)$ for every $t$. Let
$(y_1,y_2,\ldots)$ be a sequence in $\Gamma$ such that, for every
$t\in T$, the $y_i(t)$ are dense in $H(t)$. We shall
% Source PDF physical page 226
construct a sequence $(x_1,x_2,\ldots)$ in $\Gamma$ such that, for every
$t\in T$, $(x_i(t))$ is an orthonormal basis of $H(t)$. Then, if $U(t)$
is the isomorphism of $H(t)$ onto $H_0$ carrying $(x_i(t))$ onto a fixed
orthonormal basis of $H_0$, $(U(t))_{t\in T}$ will be an isomorphism of
$\mathcal H$ onto the constant field defined by $H_0$
(\XRef{10.2.4}{10.2.4}).

Suppose that elements $x_1,\ldots,x_n$ of $\Gamma$ have been constructed
with the following properties: $1^\circ$ for every $t\in T$,
$(x_1(t),x_2(t),\ldots,x_n(t))$ is an orthonormal system; $2^\circ$ there
are continuous complex functions $f_{i1},f_{i2},\ldots,f_{in}$ on $T$
such that
\[
\|y_i(t)-f_{i1}(t)x_1(t)-\cdots-f_{in}(t)x_n(t)\|
\leq\frac1i
\qquad(i=1,2,\ldots,n).
\]
We shall construct $x_{n+1}\in\Gamma$ so that
$(x_1,\ldots,x_{n+1})$ has the same properties with $n$ replaced by
$n+1$. This permits us to construct $(x_1,x_2,\ldots)$ inductively. The
system $(x_1(t),x_2(t),\ldots)$ will indeed be total in $H(t)$, since its
linear combinations approximate $y_i(t)$ to within $1/i$ and the
$y_i(t)$ are dense in $H(t)$.

Let $K(t)$ be the closed vector subspace of $H(t)$ orthogonal to
$x_1(t),\ldots,x_n(t)$, and let $\varphi(t)$ be the unit sphere of $K(t)$.
Since $\dim K(t)=\aleph_0$, $\varphi(t)$ is contractible
(\XRef{10.8.2}{10.8.2}). Moreover, the intersection of $\varphi(t)$ with
every open ball of radius less than $1$ in $H_0$ is contractible, since it
is homeomorphic to an open ball of $H_0$ by stereographic projection. We
show that the map $t\mapsto\varphi(t)$ from $T$ into the set of nonempty
closed subsets of $H_0$ is lower semicontinuous. Let $U$ be an open
subset of $H_0$ meeting $\varphi(t_0)$, and choose
$\xi\in\varphi(t_0)\cap U$. There is an $x\in\Gamma$ such that
$x(t_0)=\xi$. Define $x'$ by
\[
x'(t)=x(t)-\sum_{i=1}^{n}(x(t)\mid x_i(t))x_i(t)\in K(t).
\]
We have $x'(t_0)=\xi$, so $x'(t)\ne0$ for $t$ sufficiently near $t_0$.
Then
\[
t\longmapsto x''(t)=\frac{x'(t)}{\|x'(t)\|}
\]
is a continuous map of a neighborhood $V$ of $t_0$ into $H_0$, with
$x''(t)\in\varphi(t)$ for $t\in V$ and $x''(t_0)=\xi$. For $t$ sufficiently
near $t_0$, we have $x''(t)\in U$, and hence
$\varphi(t)\cap U\ne\varnothing$. We may therefore apply
\XRef{10.8.6}{10.8.6}.

Define $z\in\Gamma$ by
\[
y_{n+1}(t)=\sum_{i=1}^{n}(y_{n+1}(t)\mid x_i(t))x_i(t)+z(t).
\]
% Source PDF physical page 227
Then $z(t)\in K(t)$ for every $t$. Let $T'$ be the closed set of
$t\in T$ such that
\[
\|z(t)\|\geq\frac1{2(n+1)}.
\]
By \XRef{10.8.6}{10.8.6}, the map
$t\mapsto z(t)/\|z(t)\|$ from $T'$ into $H_0$ extends to a continuous map
$x_{n+1}:T\to H_0$ such that $x_{n+1}(t)\in\varphi(t)$ for every $t\in T$.
Let $f:T\to\mathbb R$ be the continuous function equal to $\|z(t)\|$ on
$T'$ and to $1/[2(n+1)]$ on $T-T'$. We have
\begin{align*}
&\left\|y_{n+1}(t)-\sum_{i=1}^{n}
(y_{n+1}(t)\mid x_i(t))x_i(t)-f(t)x_{n+1}(t)\right\|\\
&\hspace{7em}=\|z(t)-f(t)x_{n+1}(t)\|,
\end{align*}
and this is zero on $T'$ and at most
\[
\frac{2}{2(n+1)}=\frac1{n+1}
\]
on $T-T'$.

\ResultHeading{10.8.8}{10.8.8. Theorem.}
\begin{itshape}
Let $T$ be a paracompact space of finite dimension, and let $\mathcal A$
be a separable continuous field of elementary $C^{*}$-algebras over $T$,
of rank $\aleph_0$, satisfying Fell's condition. Then $\mathcal A$ is
locally trivial.
\end{itshape}

Let $t_0\in T$. There are a closed neighborhood $V$ of $t_0$ and a
continuous vector field $p$ over $V$, relative to $\mathcal A$, such that
$p(t)$ is everywhere on $V$ a rank-one projection. Put
$\beta(\mathcal A|V,p)=(\mathcal H,x)$ and
$\mathcal H=((H(t))_{t\in V},\Gamma)$. Then $\mathcal H$ is separable
(\XRef{10.7.5}{10.7.5}), every $H(t)$ has dimension $\aleph_0$, and $V$
is paracompact and finite-dimensional. Thus $\mathcal H$ is trivial
(\XRef{10.8.7}{10.8.7}). But $\mathcal A|V$ is isomorphic to
$\mathcal A(\mathcal H)$ (\XRef{10.7.6}{10.7.6}), and is therefore
trivial.

Bibliography: \cite{ref19}.

\BookSubsection[Application to C*-Algebras with Continuous Trace]{10.9}{Application
to $C^{*}$-Algebras with Continuous Trace}

\ResultHeading{10.9.1}{10.9.1.}
Let $A$ be a separable continuous-trace $C^{*}$-algebra, and let
$\mathcal A$ be the continuous field of elementary $C^{*}$-algebras over
$\widehat A$ defined by $A$. It satisfies Fell's condition
(\XRef{10.5.8}{10.5.8}). Moreover, $\widehat A$ is locally compact with a
countable base (\XRef{3.3.4}{3.3.4}, \XRef{3.3.8}{3.3.8}, and
\XRef{4.5.3}{4.5.3}), and is therefore paracompact. We may thus form the
element $\delta(\mathcal A)$ of $H^{3}(\widehat A,\mathbb Z)$
(\XRef{10.7.14}{10.7.14}). It will be denoted by $\delta(A)$.

\ResultHeading{10.9.2}{10.9.2.}
Let $T$ be a locally compact space and let $\mathcal H$ be a continuous
field of nonzero Hilbert spaces over $T$. Then
$\mathcal A(\mathcal H)$ satisfies Fell's condition
(\XRef{10.7.7}{10.7.7}). Let $A$ be the $C^{*}$-algebra defined by
$\mathcal A(\mathcal H)$ (\XRef{10.4.1}{10.4.1}). Then $A$ has spectrum
$T$ (\XRef{10.4.4}{10.4.4}) and continuous trace
(\XRef{10.5.2}{10.5.2} and \XRef{10.5.8}{10.5.8}). We shall say that
$A$ is \emph{defined by $\mathcal H$}.

\ResultHeading{10.9.3}{10.9.3. Theorem.}
\begin{itshape}
Let $A$ be a separable continuous-trace $C^{*}$-algebra. Then $A$ is
isomorphic to the $C^{*}$-algebra defined by a continuous field of
nonzero Hilbert spaces over $\widehat A$ if and only if $\delta(A)=0$.
\end{itshape}

Let $\mathcal H$ be a continuous field of nonzero Hilbert spaces over
$\widehat A$ such that $A$ is defined by $\mathcal A(\mathcal H)$. Then
$\delta(A)=\delta(\mathcal A(\mathcal H))$
(\XRef{10.5.2}{10.5.2}), and hence $\delta(A)=0$
(\XRef{10.7.15}{10.7.15}).

% Source PDF physical page 228
Conversely, suppose that $\delta(A)=0$. Let $\mathcal A$ be the continuous
field of $C^{*}$-algebras defined by $A$. There is a continuous field
$\mathcal H$ of Hilbert spaces over $\widehat A$ such that $\mathcal A$
is isomorphic to $\mathcal A(\mathcal H)$
(\XRef{10.7.15}{10.7.15}). Thus $A$ is isomorphic to the
$C^{*}$-algebra defined by $\mathcal A(\mathcal H)$
(\XRef{10.5.4}{10.5.4}).

\ResultHeading{10.9.4}{10.9.4. Definition.}
\begin{itshape}
Let $a$ be a cardinal. A $C^{*}$-algebra $A$ is called
\emph{homogeneous of degree $a$} if every irreducible representation of
$A$ has dimension $a$.
\end{itshape}

\ResultHeading{10.9.5}{10.9.5. Theorem.}
\begin{itshape}
Let $T$ be a locally compact space with a countable base and of finite
dimension.
\begin{enumerate}[label=\textup{(\roman*)}]
\item Let $A$ be a separable continuous-trace $C^{*}$-algebra,
homogeneous of degree $\aleph_0$, with spectrum $T$. The continuous field
of elementary $C^{*}$-algebras over $T$ defined by $A$ is locally trivial.
\item The map $A\mapsto\delta(A)$ defines a one-to-one correspondence
between $H^{3}(T,\mathbb Z)$ and the separable continuous-trace
$C^{*}$-algebras homogeneous of degree $\aleph_0$ with spectrum $T$, up
to isomorphism.
\end{enumerate}
\end{itshape}

Assertion (i) follows from \XRef{10.8.8}{10.8.8}.

Let $A,A'$ be two separable continuous-trace $C^{*}$-algebras,
homogeneous of degree $\aleph_0$, with spectrum $T$. Let
$\mathcal A,\mathcal A'$ be the continuous fields of $C^{*}$-algebras
over $T$ defined by $A,A'$. Suppose that $\delta(A)=\delta(A')$, that is,
$\delta(\mathcal A)=\delta(\mathcal A')$. The fields $\mathcal A$ and
$\mathcal A'$ have rank $\aleph_0$ and are locally trivial by (i). Hence
$\mathcal A$ and $\mathcal A'$ are isomorphic
(\XRef{10.8.4}{10.8.4}), and therefore $A$ and $A'$ are isomorphic.

Let $\delta\in H^{3}(T,\mathbb Z)$. There is a locally trivial continuous
field $\mathcal A$ of elementary $C^{*}$-algebras of rank $\aleph_0$ over
$T$ such that $\delta(\mathcal A)=\delta$
(\XRef{10.8.4}{10.8.4}). This field satisfies Fell's condition, and hence
defines a continuous-trace $C^{*}$-algebra $A$ with spectrum $T$,
homogeneous of degree $\aleph_0$ (\XRef{10.5.8}{10.5.8}). Moreover,
$\mathcal A$ is separable (\XRef{10.2.7}{10.2.7}), and hence $A$ is
separable (\XRef{10.4.7}{10.4.7}). The continuous field of
$C^{*}$-algebras defined by $A$ is isomorphic to $\mathcal A$
(\XRef{10.5.2}{10.5.2}), and therefore $\delta(A)=\delta$.

\ResultHeading{10.9.6}{10.9.6. Corollary.}
\begin{itshape}
Let $T$ be a locally compact space with a countable base, of finite
dimension, such that $H^{3}(T,\mathbb Z)=0$. Let $A_0$ be an elementary
$C^{*}$-algebra of rank $\aleph_0$, and let $A$ be the $C^{*}$-algebra of
continuous maps $t\mapsto x(t)$ from $T$ into $A_0$ such that
$\|x(t)\|$ tends to zero at infinity. Then every separable
continuous-trace $C^{*}$-algebra, homogeneous of degree $\aleph_0$ and
with spectrum $T$, is isomorphic to $A$.
\end{itshape}

Indeed, $A$ is separable (\XRef{10.4.7}{10.4.7}), has continuous trace
(\XRef{10.5.8}{10.5.8}), has spectrum $T$
(\XRef{10.4.4}{10.4.4}), and is homogeneous of degree $\aleph_0$
(\XRef{10.4.4}{10.4.4}). It is now enough to apply
\XRef{10.9.5}{10.9.5(ii)}.

Bibliography: \cite{ref19,ref29}.

% Source PDF physical page 229
\BookSubsection{10.10}{Complements}

\ResultHeading{10.10.1}{10.10.1.}
Let $(A_i)_{i\in I}$ be a family of $C^{*}$-algebras. The restricted
product of the $A_i$ is the $C^{*}$-algebra defined by a continuous field
of $C^{*}$-algebras over the discrete space $I$.

\ResultHeading{10.10.2}{10.10.2.}
Let $T$ be a locally compact space, let $B$ be a $C^{*}$-algebra, and let
$A$ be the $C^{*}$-algebra of continuous maps from $T$ into $B$ that tend
to zero at infinity. Then $\widehat A$ is canonically identified with
$T\times\widehat B$. More generally, in the situation of
\XRef{10.4.3}{10.4.3}, the topology of $\widehat A$ can be described
explicitly \cite{ref29}.

\ResultHeading{10.10.3}{10.10.3.}
Let $A$ be a $C^{*}$-algebra. Form the compact space $\mathcal X(A)$ and
the canonical injection of $\operatorname{Prim}(A)$ into $\mathcal X(A)$
(\XRef{1.9.13}{1.9.13}). Identify $\operatorname{Prim}(A)$ with a subset
of $\mathcal X(A)$ [in general the topology induced on
$\operatorname{Prim}(A)$ by that of $\mathcal X(A)$ is not the Jacobson
topology, because the latter is not always Hausdorff]. The closure $T$ of
$\operatorname{Prim}(A)$ in $\mathcal X(A)$ is called the
\emph{regularized spectrum} of $A$. For every $t\in T$, let $I(t)$ be the
closed two-sided ideal of $A$ consisting of those $x\in A$ on which the
seminorm $t$ vanishes; put $A(t)=A/I(t)$, and let $x(t)$ be the canonical
image of $x\in A$ in $A(t)$. The set $\Gamma$ of vector fields
$t\mapsto x(t)$, $x\in A$, has the following properties: $1^\circ$
$\Gamma$ is an involutive algebra of vector fields; $2^\circ$ the set of
$x(t)$, $x\in A$, is all of $A(t)$; $3^\circ$ for every $x\in A$, the
function $t\mapsto\|x(t)\|$ is continuous on $T$; $4^\circ$ the set of
these fields is complete for uniform convergence. For
$((A(t))_{t\in T},\Gamma)$ to be a continuous field of $C^{*}$-algebras,
it is necessary and sufficient that $\operatorname{Prim}(A)$ be Hausdorff
for the Jacobson topology \cite{ref29}.

\ResultHeading{10.10.4}{10.10.4.}
Let $T=\{0,1,2,3,\ldots,\infty\}$ with its usual topology. For $t\in T$,
define $A(t)$ as follows: if $t=\infty$, let $A(t)=\mathbb C$; if
$t\ne\infty$, let $A(t)$ be the $C^{*}$-algebra of two-rowed and
two-columned complex matrices. Let $\Theta$ be the set of vector fields
$t\mapsto x(t)\in A(t)$ having the following properties [write
$x(t)=(x_{ij}(t))_{i,j=1,2}$ for $t\ne\infty$]:
\[
\lim_{n\to\infty}x_{12}(n)
=\lim_{n\to\infty}x_{21}(n)
=\lim_{n\to\infty}x_{22}(2n)=0,
\]
\[
\lim_{n\to\infty}x_{11}(n)
=\lim_{n\to\infty}x_{22}(2n+1)=x(\infty).
\]
Then $\mathcal A=((A(t))_{t\in T},\Theta)$ is a continuous field of
$C^{*}$-algebras. Let $A=\Theta$ be the $C^{*}$-algebra defined by
$\mathcal A$. We have $\widehat A=T$. The ideal $\mathfrak m(A)$
(\XRef{4.5.2}{4.5.2}) is the set of $x\in A$ such that $x(\infty)=0$.
The algebra $A$ is liminal and has generalized continuous trace
(\XRef{4.7.12}{4.7.12}), but it does not have continuous trace;
$\mathcal A$ does not satisfy Fell's condition at $\infty$ \cite{ref27}.

\ResultHeading{10.10.5}{10.10.5.}
Let $H=L^{2}_{\mathbb C}([0,1])$. For every finite sequence
$s=(s_1,\ldots,s_n)$, where $n=1,2,\ldots$ and each $s_i$ is $0$ or $1$,
put
\[
k_s=\sum_{i=1}^{n}s_i/2^{i},
\qquad
I_s=[k_s,k_s+2^{-n}].
\]
% Source PDF physical page 230
If $s$ and $t$ are two such sequences of $n$ terms, denote by $Q_{s,t}$
the operator on $H$ that carries $f\in H$ into $g\in H$, where
\[
g(x)=0\quad\text{for }x\notin I_s,
\qquad
g(k_s+u)=f(k_t+u)\quad\text{if }0\leq u\leq2^{-n}.
\]
Let $A(n)$ be the $C^{*}$-algebra of operators generated by the $Q_{s,t}$,
with $s,t$ ranging over the sequences of $n$ terms. Then $A(n)$ is
isomorphic to the $C^{*}$-algebra of complex matrices with $2^n$ rows and
$2^n$ columns. We have $A(n)\subset A(n+1)$. Let $A$ be the
$C^{*}$-subalgebra of $\mathcal L(H)$ generated by the $A(n)$. Let
$A(\infty)$ be the commutative $C^{*}$-algebra of multiplication operators
by continuous complex functions on $[0,1]$. We have $A(\infty)\subset A$.

Let $T=\{1,2,\ldots,\infty\}$ with its usual topology. The $A(t)$,
$t\in T$, have now been defined. Let $\Theta$ be the set of continuous
maps $x:T\to A$ such that $x(t)\in A(t)$ for every $t\in T$. Then
$\mathcal A=((A(t))_{t\in T},\Theta)$ is a continuous field of
$C^{*}$-algebras. Let $A=\Theta$ be the $C^{*}$-algebra defined by
$\mathcal A$. The elements of $\widehat A$ are the
representations
\[
\pi_n:x\longmapsto x(n)\quad(n<\infty)
\]
and the representations
\[
\pi_\lambda:x\longmapsto\lambda(x(\infty)),
\qquad \lambda\in\widehat{A(\infty)}.
\]
The $C^{*}$-algebra $A$ is liminal and has generalized continuous
trace. The sequence $(\pi_n)_{n=1,2,\ldots}$ converges simultaneously to
all the $\pi_\lambda$ (Fell, unpublished).

\ResultHeading{10.10.6}{10.10.6.}
Let $A$ be a $C^{*}$-algebra.
\begin{enumerate}[label=\textup{\alph*.}]
\item The space $\operatorname{Prim}(A)$ is discrete if and only if $A$
is the restricted product of a family of $C^{*}$-algebras having no
nontrivial closed two-sided ideals.
\item The algebra $A$ is dual if and only if it is liminal and has
discrete spectrum \cite{ref74}.
\end{enumerate}

\ResultHeading{10.10.7}{10.10.7.}
Let $A$ be a liminal (respectively, postliminal) $C^{*}$-algebra, let $T$
be compact, and let $B$ be the $C^{*}$-algebra of continuous maps from
$T$ into $A$. Then $B$ is liminal (respectively, postliminal). [If $A$ is
liminal, apply \XRef{10.4.3}{10.4.3}. If $A$ is postliminal, let
$(I_\rho)$ be a composition series of $A$ such that the
$I_{\rho+1}/I_\rho$ are liminal. Let $J_\rho$ be the set of continuous
maps from $T$ into $I_\rho$. Then the $J_\rho$ form a composition series
of $B$, and $J_{\rho+1}/J_\rho$ is isomorphic to the $C^{*}$-algebra of
continuous maps from $T$ into $I_{\rho+1}/I_\rho$.] \cite{ref77}.

\ResultHeading{10.10.8}{10.10.8.}
Let $A$ be a liminal $C^{*}$-algebra, let $U$ be an open dense subset of
$\widehat A$, and let $I$ be the closed two-sided ideal of $A$ such that
$\widehat I=U$. Let
$\mathcal A=((A(t))_{t\in U},\Theta)$ be the continuous field of
elementary $C^{*}$-algebras over $U$ defined by $I$. For every $x\in A$
and $t\in U$, let $x(t)$ be the canonical image of $x$ in $A(t)$ [the
image of $A$ under the representation $t$ is the same as that of $I$,
namely $\mathcal{LC}(H_t)$]. For every $x\in A$, let $x'$ be the field
$t\mapsto x(t)$ over $U$. Then $x\mapsto x'$ is injective. In general the
field $x'$ is not an element of $\Theta$, but, for every $y\in I$, the
fields $x'y$ and $yx'$ belong to $\Theta$. Suppose that $A$ has
generalized continuous trace and that $I=J(A)$
(\XRef{4.7.12}{4.7.12}). Then $x'\in\Theta$ for every $x\in A$. [For
every $y\in I$, the function $t\mapsto\|x(t)-y(t)\|$ is continuous on
$U$, because the points of $U$ are separated in $\widehat A$
(\XRef{3.9.4}{3.9.4}).]

% Source PDF physical page 231
\ResultHeading{10.10.9}{10.10.9.}
Let $S$ be the two-dimensional sphere and let
$T=S\times S\times S\times\cdots$ be the compact product space. There is
over $T$ a continuous field
$\mathcal H=((H(t))_{t\in T},\Gamma)$ of Hilbert spaces of dimension
$\aleph_0$ which is not locally trivial at any point of $T$. Let
$A$ be the continuous-trace $C^{*}$-algebra defined by
$\mathcal A(\mathcal H)$. Then $A$ has no nonzero closed two-sided
ideal $I$ isomorphic to the $C^{*}$-algebra of continuous maps
$X\to A_0$ tending to zero at infinity, where $X$ is a locally compact
space and $A_0$ an elementary $C^{*}$-algebra \cite{ref19}.

\ResultHeading{10.10.10}{10.10.10.}
Let $T$ be a topological space, and let
\[
\mathcal A_1=((A_1(t))_{t\in T},\Theta_1),
\qquad
\mathcal A_2=((A_2(t))_{t\in T},\Theta_2)
\]
be two continuous fields of elementary $C^{*}$-algebras over $T$
satisfying Fell's condition. For every $t\in T$, let
$A(t)=A_1(t)\otimes A_2(t)$ (\XRef{2.12.15}{2.12.15}). Let $\Gamma$ be
the set of vector fields of the form
\[
t\longmapsto a_1(t)\otimes a_2(t)+b_1(t)\otimes b_2(t)+\cdots
+k_1(t)\otimes k_2(t),
\]
where $a_1,b_1,\ldots,k_1\in\Theta_1$ and
$a_2,b_2,\ldots,k_2\in\Theta_2$. There is a unique continuous field of
elementary $C^{*}$-algebras
$\mathcal A=((A(t))_{t\in T},\Theta)$ over $T$ such that
$\Theta\supset\Gamma$. We write $\mathcal A=\mathcal A_1\otimes
\mathcal A_2$. This field satisfies Fell's condition. If $T$ is
paracompact, then
\[
\delta(\mathcal A_1\otimes\mathcal A_2)
=\delta(\mathcal A_1)+\delta(\mathcal A_2).
\]
If $\mathcal A_1$ is separable and $\mathcal A_2$ is locally trivial of
infinite rank, then $\mathcal A_1\otimes\mathcal A_2$ is locally trivial
\cite{ref18}.

\ResultHeading{10.10.11}{10.10.11. Problem.}
Let $((A(t))_{t\in T},\Theta)$ be a continuous field of postliminal
$C^{*}$-algebras over $T$. For every $t\in T$, let $B(t)$ be the largest
liminal ideal of $A(t)$. Let $\Theta'$ be the set of $x\in\Theta$ such
that $x(t)\in B(t)$ for every $t$. Is $((B(t))_{t\in T},\Theta')$ a
continuous field of $C^{*}$-algebras?

\ResultHeading{10.10.12}{10.10.12.}
Let $A$ be a $C^{*}$-algebra, let $x\in A$, and let $f$ be a bounded
continuous complex function on $\widehat A$. There is a $y\in A$ such
that
\[
\pi(y)=f(\pi)\pi(x)
\qquad\text{for every }\pi\in\widehat A
\]
\cite{ref346}. Cf. also \cite{ref313}.

\BookSection[Extension to C*-Algebras of the Stone-Weierstrass
Theorem]{11}{Extension to $C^{*}$-Algebras of the Stone--Weierstrass
Theorem}

Let $A$ be the $C^{*}$-algebra of continuous complex functions tending to
zero at infinity on a locally compact space. The Stone--Weierstrass
theorem gives conditions under which a $C^{*}$-subalgebra $B$ of $A$ is
equal to $A$. But the $C^{*}$-algebras of the form $A$ are, up to
isomorphism, precisely the commutative $C^{*}$-algebras. In this section
we shall extend the Stone--Weierstrass theorem to arbitrary
$C^{*}$-algebras.

% Source PDF physical page 232
\BookSubsection[The Case of Postliminal C*-Algebras]{11.1}{The Case of
Postliminal $C^{*}$-Algebras}

\ResultHeading{11.1.1}{11.1.1. Definition.}
\begin{itshape}
Let $A$ be a $C^{*}$-algebra and $B$ a $C^{*}$-subalgebra of $A$. We say
that $B$ is a \emph{rich $C^{*}$-subalgebra} of $A$ if the following
conditions hold:
\begin{enumerate}[label=\textup{(\roman*)}]
\item for every irreducible representation $\pi$ of $A$, the
restriction $\pi|B$ is irreducible;
\item if $\pi$ and $\pi'$ are inequivalent irreducible representations
of $A$, then $\pi|B$ and $\pi'|B$ are inequivalent.
\end{enumerate}
\end{itshape}

It follows that, if $\pi$ is a nonzero irreducible representation of
$A$, then $\pi|B$ is nonzero.

\ResultHeading{11.1.2}{11.1.2. Lemma.}
\begin{itshape}
Let $A$ be a $C^{*}$-algebra, let $B$ be a rich $C^{*}$-subalgebra of
$A$, and let $\pi$ be an irreducible representation of $B$. There is an
irreducible representation $\pi'$ of $A$ on $H_\pi$ such that
$\pi=\pi'|B$.
\end{itshape}

There are a Hilbert space $H'\supset H_\pi$ and an irreducible
representation $\pi'$ of $A$ on $H'$ such that
$\pi(x)=\pi'(x)|H_\pi$ for every $x\in B$ (\XRef{2.10.2}{2.10.2}). Since
$B$ is rich, $\pi'|B$ is irreducible; and $H_\pi$ is invariant under
$\pi'(B)$, so $H_\pi=H'$.

\ResultHeading{11.1.3}{11.1.3. Lemma.}
\begin{itshape}
Let $A$ be a $C^{*}$-algebra, let $B$ be a rich $C^{*}$-subalgebra of
$A$, and let $I$ be a closed two-sided ideal of $A$.
\begin{enumerate}[label=\textup{(\roman*)}]
\item $B/(B\cap I)$ is a rich $C^{*}$-subalgebra of $A/I$;
\item $B\cap I$ is a rich $C^{*}$-subalgebra of $I$.
\end{enumerate}
\end{itshape}

Assertion (i) follows at once from the definitions. We prove (ii). Let
$\pi$ be an irreducible representation of $I$ and let us prove that
$\pi|B\cap I$ is irreducible. We may suppose that $\pi\ne0$. There is an
irreducible representation $\rho$ of $A$ on $H_\pi$ extending $\pi$
(\XRef{2.10.4}{2.10.4}). Since $B$ is rich, $\rho|B$ is irreducible. As
$B\cap I$ is a closed two-sided ideal of $B$, $\pi|B\cap I$ is zero or
irreducible (\XRef{2.11.2}{2.11.2}). If $\pi|B\cap I=0$, then $\rho|B$
defines an irreducible representation of $B/(B\cap I)$ which, by (i) and
\XRef{11.1.2}{11.1.2}, extends to an irreducible representation of
$A/I$. In other words, $\rho|B$ extends to an irreducible representation
$\rho'$ of $A$ that vanishes on $I$. Since $B$ is rich, $\rho'$ is
equivalent to $\rho$, and hence $\pi=\rho|I=0$, contrary to the
hypothesis. Thus $\pi|B\cap I$ is irreducible.

Now let $\pi,\pi'$ be two irreducible representations of $I$ and suppose
that $\pi|B\cap I\simeq\pi'|B\cap I$. Let $\rho,\rho'$ be representations
of $A$ on $H_\pi,H_{\pi'}$ extending $\pi,\pi'$
(\XRef{2.10.4}{2.10.4}). If $\pi|B\cap I$ and $\pi'|B\cap I$ are zero,
the preceding argument shows that $\pi$ and $\pi'$ are zero. Otherwise,
$\rho|B\simeq\rho'|B$ (\XRef{2.11.2}{2.11.2}), so
$\rho\simeq\rho'$ because $B$ is rich, and hence $\pi\simeq\pi'$.

\ResultHeading{11.1.4}{11.1.4. Lemma.}
\begin{itshape}
Let $A$ be a liminal $C^{*}$-algebra with Hausdorff spectrum, and let $B$
be a rich $C^{*}$-subalgebra of $A$. Then $B=A$.
\end{itshape}

% Source PDF physical page 233
Let $\mathcal A$ be the continuous field of elementary $C^{*}$-algebras
over $\widehat A$ defined by $A$. Then $A$ is identified with the
$C^{*}$-algebra defined by $\mathcal A$ (\XRef{10.5.4}{10.5.4}). If
$\pi\in\widehat A$, then $\pi|B\in\widehat B$, and hence
\[
\pi(A)=\pi(B)=\mathcal{LC}(H_\pi).
\]
If $\pi$ and $\pi'$ are two inequivalent irreducible representations of
$A$, then $\pi|B$ and $\pi'|B$ are inequivalent. Thus, if
$T\in\mathcal{LC}(H_\pi)$ and $T'\in\mathcal{LC}(H_{\pi'})$, there is an
$x\in B$ such that $\pi(x)=T$ and $\pi'(x)=T'$
(\XRef{4.2.5}{4.2.5}). Hence $B=A$ by \XRef{10.5.3}{10.5.3}.

\ResultHeading{11.1.5}{11.1.5. Lemma.}
\begin{itshape}
Let $A$ be a $C^{*}$-algebra, let $I$ be its largest postliminal ideal,
and let $B$ be a rich $C^{*}$-subalgebra of $A$. Then $B\supset I$.
\end{itshape}

There is a composition series $(I_\rho)_{\rho\leq\alpha}$ of $I$ such
that the $I_{\rho+1}/I_\rho$ are continuous-trace $C^{*}$-algebras
(\XRef{4.5.5}{4.5.5}), and hence liminal $C^{*}$-algebras with Hausdorff
spectrum (\XRef{4.5.3}{4.5.3}). We have $I_0=0\subset B$. Suppose that
there is a least ordinal $\rho$ for which $I_\rho\not\subset B$. Then
$I_{\rho'}\subset B$ for $\rho'<\rho$, so $\rho$ is not a limit ordinal.
Let $\rho=\rho'+1$. By \XRef{11.1.3}{11.1.3(i)}, $B/I_{\rho'}$ is a rich
$C^{*}$-subalgebra of $A/I_{\rho'}$. By
\XRef{11.1.3}{11.1.3(ii)}, $(B\cap I_\rho)/I_{\rho'}$ is a rich
$C^{*}$-subalgebra of $I_\rho/I_{\rho'}$. By
\XRef{11.1.4}{11.1.4},
\[
(B\cap I_\rho)/I_{\rho'}=I_\rho/I_{\rho'},
\]
whence $B\supset I_\rho$, a contradiction. Therefore $I=I_\alpha\subset B$.

\ResultHeading{11.1.6}{11.1.6. Proposition.}
\begin{itshape}
In a postliminal $C^{*}$-algebra $A$, every rich $C^{*}$-subalgebra is
equal to $A$.
\end{itshape}

This is a special case of \XRef{11.1.5}{11.1.5}.

\ResultHeading{11.1.7}{11.1.7.}
Let $E,F$ be two vector spaces in duality, with $A\subset E$ and
$B\subset F$. We say that $A$ \emph{separates} $B$ if, for every pair
$(x,y)$ of distinct points of $B$, there is an $a\in A$ such that
$\langle a,x\rangle\ne\langle a,y\rangle$.

\ResultLabel{Lemma.}
\begin{itshape}
Let $A$ be a $C^{*}$-algebra and let $B$ be a $C^{*}$-subalgebra that
separates $P(A)\cup\{0\}$. Then $B$ is a rich $C^{*}$-subalgebra of $A$.
\end{itshape}

Let $\pi\in\widehat A$. Suppose that $H_\pi$ is the Hilbert direct sum of
two nonzero subspaces $H_1,H_2$ invariant under $\pi|B$. Choose nonzero
$\xi_1\in H_1$, $\xi_2\in H_2$ such that
$\|\xi_1+\xi_2\|=1$. Then
\[
x\longmapsto(\pi(x)(\xi_1+\xi_2)\mid\xi_1+\xi_2),
\qquad
x\longmapsto(\pi(x)(\xi_1-\xi_2)\mid\xi_1-\xi_2)
\]
are pure states $f_1,f_2$ of $A$. If $x\in B$, then
\begin{align*}
f_1(x)
&=(\pi(x)\xi_1\mid\xi_1)+(\pi(x)\xi_2\mid\xi_2)\\
&=(\pi(x)\xi_1\mid\xi_1)+(\pi(x)(-\xi_2)\mid-\xi_2)=f_2(x).
\end{align*}
Since $B$ separates $P(A)$, we have $f_1=f_2$. Thus
$\xi_1+\xi_2$ is proportional to $\xi_1-\xi_2$
(\XRef{2.5.7}{2.5.7}), so $\xi_1$ and $\xi_2$ are proportional, which is
absurd. Therefore $\pi|B$ is irreducible.

Let $\pi_1,\pi_2$ be two irreducible representations of $A$, and suppose
there is an isomorphism $U:H_{\pi_1}\to H_{\pi_2}$ carrying
$\pi_1|B$ into $\pi_2|B$. Let $\xi_1$ be a unit vector of $H_{\pi_1}$,
let $\xi_2=U(\xi_1)$, and let $f_i\in P(A)\cup\{0\}$ be the functional
% Source PDF physical page 234
defined by $\pi_i$ and $\xi_i$. We have $f_1(x)=f_2(x)$ for $x\in B$,
and hence $f_1=f_2$. Thus either $\pi_1$ and $\pi_2$ are both zero, or
they are both nonzero. In the latter case $\xi_i$ is cyclic for $\pi_i$,
and $f_1=f_2$ implies $\pi_1\simeq\pi_2$ (\XRef{2.4.1}{2.4.1}).

\ResultHeading{11.1.8}{11.1.8. Theorem.}
\begin{itshape}
Let $A$ be a postliminal $C^{*}$-algebra and $B$ a $C^{*}$-subalgebra of
$A$. If $B$ separates $P(A)\cup\{0\}$, then $B=A$.
\end{itshape}

This follows from \XRef{11.1.6}{11.1.6} and \XRef{11.1.7}{11.1.7}.

It is not known whether the hypothesis that $A$ be postliminal is
essential. For arbitrary $C^{*}$-algebras we shall obtain later
(\XRef{11.5.2}{11.5.2}) a less precise result than
\XRef{11.1.8}{11.1.8}.

Bibliography: \cite{ref45,ref77}.

This subsection follows an unpublished account by Fell.

\BookSubsection[Abundance of Pure States in Certain C*-Algebras]{11.2}{Abundance
of Pure States in Certain $C^{*}$-Algebras}

\ResultHeading{11.2.1}{11.2.1. Lemma.}
\begin{itshape}
Let $H$ be a Hilbert space, let $C=\mathcal{LC}(H)$, let $A$ be a
$C^{*}$-subalgebra of $\mathcal L(H)$ containing $1$, and let $f$ be a
state of $A$ that vanishes on $A\cap C$.
\begin{enumerate}[label=\textup{(\roman*)}]
\item $f$ is a weak limit of vector states of $A$ (\XRef{2.4.2}{2.4.2}).
\item If $A$ acts irreducibly on $H$, then $f$ is a weak limit of pure
states of $A$.
\end{enumerate}
\end{itshape}

Part (ii) follows at once from (i); we prove (i).

Suppose first that $A\supset C$. Then $f$ defines a state $f'$ of $A/C$,
itself a weak limit of convex combinations of pure states of $A/C$
(\XRef{2.5.6}{2.5.6}). Hence $f$ is a weak limit of states of the form
\[
\lambda_1f_1+\cdots+\lambda_nf_n,
\]
where $\lambda_1,\ldots,\lambda_n\geq0$,
$\lambda_1+\cdots+\lambda_n=1$, and $f_1,\ldots,f_n$ are pure states of
$A$ vanishing on $C$. It is therefore enough to consider the case in
which $f$ itself has this form. Let $x_1,\ldots,x_s$ be Hermitian
elements of $A$ with $x_1=1$. We shall construct unit vectors
$\xi_1,\ldots,\xi_n$ in $H$ such that
$(x_i\xi_j\mid\xi_k)=0$ for $j<k$ and
\[
|f_j(x_i)-\omega_{\xi_j}(x_i)|\leq1
\qquad\text{for all }i,j.
\]
Suppose that the $\xi_j$ have been constructed for $j<m$. Let $K$ be the
vector subspace of $H$ spanned by the $x_i\xi_j$, with
$1\leq i\leq s$ and $j<m$, and let $K'=H\ominus K$. We have
$P_K\in C\subset A$, and hence $P_{K'}AP_{K'}$ is a subalgebra of $A$.
Moreover, $f_m|P_{K'}AP_{K'}$ is a state of $P_{K'}AP_{K'}$, since
\[
f_m(P_{K'})=f_m(P_K)+f_m(P_{K'})=f_m(1)=1.
\]
If
\[
f_m|P_{K'}AP_{K'}=\lambda g_1+(1-\lambda)g_2,
\qquad0<\lambda<1,
\]
where $g_1,g_2$ are states of $P_{K'}AP_{K'}$, let $\widetilde g_i$ be
the state $x\mapsto g_i(P_{K'}xP_{K'})$ of $A$. Since $f_m$ vanishes on
$C$, we have
$f_m=\lambda\widetilde g_1+(1-\lambda)\widetilde g_2$; therefore
$\widetilde g_1=\widetilde g_2$ because $f_m$ is pure, and hence
$g_1=g_2$. Thus $f_m|P_{K'}AP_{K'}$ is a pure state of
$P_{K'}AP_{K'}$. By \XRef{3.4.1}{3.4.1}, it is a weak limit of forms
$\omega_{\xi_\alpha}$, where the $\xi_\alpha$ are unit vectors in $K'$.
Thus, for $1\leq i\leq s$,
\[
f_m(x_i)=f_m(P_{K'}x_iP_{K'})
\]
% Source PDF physical page 235
is a limit of
\[
(P_{K'}x_iP_{K'}\xi_\alpha\mid\xi_\alpha)
=(x_i\xi_\alpha\mid\xi_\alpha).
\]
We can therefore find $\xi_m\in K'$ such that
\[
|f_m(x_i)-\omega_{\xi_m}(x_i)|\leq1
\qquad(1\leq i\leq s).
\]
The construction of the $\xi_j$ can thus be continued inductively. Put
\[
\xi=\sqrt{\lambda_1}\xi_1+\cdots+\sqrt{\lambda_n}\xi_n.
\]
Since $x_1=1$, we have $(\xi_j\mid\xi_k)=0$ for $j\ne k$, so $\xi$ is a
unit vector. Since $x_i$ is Hermitian, we have
$(x_i\xi_j\mid\xi_k)=0$ for $j\ne k$, and hence
\begin{align*}
\left|\sum_{j=1}^{n}\lambda_jf_j(x_i)-\omega_\xi(x_i)\right|
&=\left|\sum_{j=1}^{n}\lambda_jf_j(x_i)
-\sum_{j,k=1}^{n}(x_i\sqrt{\lambda_j}\xi_j
 \mid\sqrt{\lambda_k}\xi_k)\right|\\
&=\left|\sum_{j=1}^{n}\lambda_jf_j(x_i)
-\sum_{j=1}^{n}\lambda_j(x_i\xi_j\mid\xi_j)\right|
\leq\sum_{j=1}^{n}\lambda_j=1.
\end{align*}
This proves that $\sum_{j=1}^{n}\lambda_jf_j$ is a weak limit of vector
states of $A$.

Now consider the general case. Let $A_1$ be the $C^{*}$-algebra $A+C$.
Then $A_1/C$ is identified with $A/(A\cap C)$, and hence $f$ defines a
state of $A_1/C$; in other words, $f$ extends to a state $f'$ of $A_1$
that vanishes on $C$. By what precedes, $f'$ is a weak limit of vector
states of $A_1$, and consequently $f$ is a weak limit of vector states of
$A$.

\ResultHeading{11.2.2}{11.2.2.}
Let $A$ be a $C^{*}$-algebra. For every representation $\theta$ of $A$,
temporarily denote by $C_\theta$ the set of $x\in A$ such that
$\theta(x)$ is compact. We have
\[
C_\theta=\theta^{-1}(\mathcal{LC}(H_\theta)),
\]
so $C_\theta$ is a closed two-sided ideal of $A$.

\ResultLabel{Lemma.}
\begin{itshape}
Let $A$ be an antiliminal $C^{*}$-algebra with identity, and let $P'(A)$
be the set of states that vanish on at least one $C_\theta$,
$\theta\in\widehat A$. Then $P(A)$ and $P'(A)$ have the same weak
closure.
\end{itshape}

Let $x$ be a Hermitian element of $A$ such that $f(x)\geq0$ for every
$f\in P'(A)$. Let $\theta\in\widehat A$, and let $\pi_\theta$ be a
representation of $A$ with kernel $C_\theta$. Every state associated with
$\pi_\theta$ belongs to $P'(A)$ and is therefore nonnegative on $x$;
hence $\pi_\theta(x)\geq0$. But
\[
\bigcap_{\theta\in\widehat A}C_\theta=0
\]
because $A$ is antiliminal (\XRef{4.2.6}{4.2.6}). Thus
$\bigoplus_{\theta\in\widehat A}\pi_\theta$ is isometric, and consequently
$x\geq0$. By \XRef{3.4.1}{3.4.1}, the weak closure of $P'(A)$ contains
$P(A)$.

% Source PDF physical page 236
If $g\in P'(A)$, there is a $\theta\in\widehat A$ such that
$g(C_\theta)=0$. Thus $g$ defines a state $g'$ of $\theta(A)$ that
vanishes on $\theta(A)\cap\mathcal{LC}(H_\theta)$. By
\XRef{11.2.1}{11.2.1}, $g'$ is a weak limit of pure states of $\theta(A)$,
and hence $g$ is a weak limit of pure states of $A$. Thus the weak closure
of $P(A)$ contains $P'(A)$.

\ResultHeading{11.2.3}{11.2.3.}
Let $A$ be a $C^{*}$-algebra. Throughout the remainder of this subsection,
$\overline{P(A)}$ will denote the weak closure of $P(A)$ in the dual of
$A$, and $E(A)$ will denote the set of states of $A$.

\ResultLabel{Lemma.}
\begin{itshape}
Let $A$ be an antiliminal $C^{*}$-algebra with identity, and let
$f\in\overline{P(A)}$. Every state of $A$ that vanishes on
$\ker\pi_f$ belongs to $\overline{P(A)}$.
\end{itshape}

By \XRef{3.4.3}{3.4.3}, it is enough to prove that the state
\[
x\longmapsto\lambda_1f(y_1^{*}xy_1)+\cdots+
\lambda_nf(y_n^{*}xy_n),
\]
where
\[
\lambda_1,\ldots,\lambda_n\geq0,
\qquad y_1,\ldots,y_n\in A,
\qquad \sum_i\lambda_if(y_i^{*}y_i)=1,
\]
belongs to $\overline{P(A)}$. But $f$ is a weak limit of states
$f_\alpha\in P'(A)$ (\XRef{11.2.2}{11.2.2}), and the state
\[
x\longmapsto\lambda_1f_\alpha(y_1^{*}xy_1)+\cdots+
\lambda_nf_\alpha(y_n^{*}xy_n)
\]
plainly belongs to $P'(A)$, and hence to $\overline{P(A)}$
(\XRef{11.2.2}{11.2.2}).

\ResultHeading{11.2.4}{11.2.4. Lemma.}
\begin{itshape}
Let $A$ be an antiliminal $C^{*}$-algebra with identity. Suppose that any
two nonzero closed two-sided ideals of $A$ have nonzero intersection.
Then $E(A)=\overline{P(A)}$.
\end{itshape}

The condition on the ideals of $A$ means that any two nonempty open
subsets of $\widehat A$ have nonempty intersection
(\XRef{3.2.2}{3.2.2}). Let $R$ be the equivalence relation on $P(A)$
defined by the canonical map $P(A)\to\widehat A$. If $f\in P(A)$, the
elements of $P(A)$ congruent to $f$ modulo $R$ are the states
$x\mapsto f(u^{*}xu)$, where $u$ belongs to the unitary group $G$ of $A$
(\XRef{2.8.6}{2.8.6}). Since $\widehat A$ is the quotient space of $P(A)$
by $R$, any two nonempty open subsets of $P(A)$ saturated for $R$ have
nonempty intersection. Let now $f_1,f_2\in P(A)$ and
$\lambda\in[0,1]$. It follows from what precedes that
$\lambda f_1+(1-\lambda)f_2$ is a weak limit of states
\begin{equation}
 x\longmapsto\lambda f(x)+(1-\lambda)f(u^{*}xu),
 \tag{1}\label{eq:11.2.4-1}
\end{equation}
where $f\in P(A)$ and $u\in G$. But the state \eqref{eq:11.2.4-1}
vanishes on $\ker\pi_f$, and hence belongs to $\overline{P(A)}$
(\XRef{11.2.3}{11.2.3}). Thus
$\lambda f_1+(1-\lambda)f_2\in\overline{P(A)}$. By induction, every
convex combination of pure states of $A$ belongs to $\overline{P(A)}$.
Therefore $E(A)\subset\overline{P(A)}$ (\XRef{2.5.6}{2.5.6}).

Bibliography: \cite{ref45}.

% Source PDF physical page 237
\BookSubsection{11.3}{Statement of the Theorem}

\ResultHeading{11.3.1}{11.3.1. Theorem.}
\begin{itshape}
Let $A$ be a $C^{*}$-algebra with identity and let $B$ be a
$C^{*}$-subalgebra of $A$ containing $1$. If $B$ separates
$\overline{P(A)}$, then $B=A$.
\end{itshape}

The proof will be given in \XRef{11.5.1}{11.5.1}.

\ResultHeading{11.3.2}{11.3.2. Lemma.}
\begin{itshape}
Let $A$ and $B$ be as in \XRef{11.3.1}{11.3.1}. If $B$ separates
$E(A)$, then $B=A$.
\end{itshape}

Let $A_h$ (respectively, $B_h$) be the set of Hermitian elements of $A$
(respectively, $B$). Suppose that $A\ne B$, and hence $A_h\ne B_h$.
By the Hahn--Banach theorem there is a nonzero continuous Hermitian linear
functional $f$ on $A$ such that $f(B)=0$. Let $f=g-g'$, where $g,g'$ are
positive functionals on $A$ (\XRef{2.6.4}{2.6.4}). We have
\[
0=f(1)=g(1)-g'(1),
\qquad\text{so}\qquad g(1)=g'(1)\ne0.
\]
Multiplying $f,g,g'$ by the same scalar, we may suppose that
$g(1)=g'(1)=1$, and hence $g,g'\in E(A)$. But then $B$ does not separate
$g$ and $g'$, although $g\ne g'$.

\ResultHeading{11.3.3}{11.3.3.}
By \XRef{11.2.4}{11.2.4} and \XRef{11.3.2}{11.3.2}, Theorem
\XRef{11.3.1}{11.3.1} is proved when $A$ is an antiliminal
$C^{*}$-algebra in which the intersection of any two nonzero closed
two-sided ideals is nonzero. We shall reduce the general case to this
special case.

Bibliography: \cite{ref45}.

\BookSubsection{11.4}{Some Lemmas}

\ResultHeading{11.4.1}{11.4.1. Lemma.}
\begin{itshape}
Let $A$ be a $C^{*}$-algebra with identity, and let $B$ be a
$C^{*}$-subalgebra of $A$ containing $1$ and separating
$\overline{P(A)}$. Let $K_1,K_2$ be two closed two-sided ideals of $A$,
and put
\[
K=B\cap(K_1+K_2),
\qquad
K'=(B\cap K_1)+(B\cap K_2).
\]
Then $K=K'$.
\end{itshape}

It is clear that $K\supset K'$. Let $f$ be a pure state of $B$ such that
$f(K')=0$. We shall show that $f(K)=0$, which will complete the proof
(\XRef{2.7.3}{2.7.3}).

There is an $f'\in P(A)$ extending $f$ [\XRef{2.10.1}{2.10.1(ii)}]; since
$B$ separates $P(A)$, this extension is unique. By
\XRef{2.10.1}{2.10.1(iii)}, for every Hermitian element $x$ of $A$,
\begin{equation}
 f'(x)=\sup_{\substack{y=y^{*}\in B\\y\leq x}}f(y).
 \tag{1}\label{eq:11.4.1-1}
\end{equation}
Let $x\in K_1^{+}$. If $y=y^{*}\in B$ and $y\leq x$, and if $\varphi$
denotes the canonical morphism of $A$ onto $A/K_1$, then
$\varphi(y)\leq\varphi(x)=0$. Hence
$\varphi(y^{+})=\varphi(y)^{+}=0$, so $y^{+}\in K_1$; therefore
$f(y^{+})=0$ and $f(y)\leq0$. Equation \eqref{eq:11.4.1-1} now shows that
$f'(x)=0$. Thus $f'$ vanishes on $K_1$, and similarly on $K_2$; hence
$f(K)=0$.

\ResultHeading{11.4.2}{11.4.2. Lemma.}
\begin{itshape}
Let $A$ and $B$ be as in \XRef{11.4.1}{11.4.1}, and let $\pi_1,\pi_2$ be
representations of $A$ such that
$\pi_1(A)=\pi_1(B)$ and $\pi_2(A)=\pi_2(B)$. Then
\[
(\pi_1\oplus\pi_2)(A)=(\pi_1\oplus\pi_2)(B).
\]
\end{itshape}

% Source PDF physical page 238
Put $K_i=\ker\pi_i$. We have
$A=B+K_1=B+K_2$, whence first
\[
K_1\subset\bigl(B\cap(K_1+K_2)\bigr)+K_2.
\]
In view of \XRef{11.4.1}{11.4.1},
\[
K_1\subset(B\cap K_1)+(B\cap K_2)+K_2
=(B\cap K_1)+K_2.
\]
It follows that
\[
K_1\subset(B\cap K_1)+(K_1\cap K_2),
\]
and hence
\[
A=B+K_1=B+(K_1\cap K_2),
\]
which proves the lemma.

\ResultHeading{11.4.3}{11.4.3. Lemma.}
\begin{itshape}
Let $A$ be an antiliminal $C^{*}$-algebra with identity, and let $B$ be a
$C^{*}$-subalgebra of $A$ containing $1$ and separating
$\overline{P(A)}$. If $f_1,f_2\in\overline{P(A)}$, then
\[
(\pi_{f_1}\oplus\pi_{f_2})(A)
=(\pi_{f_1}\oplus\pi_{f_2})(B).
\]
\end{itshape}

If $f\in\overline{P(A)}$, every state of $A$ that vanishes on
$\ker\pi_f$ belongs to $\overline{P(A)}$ (\XRef{11.2.3}{11.2.3}). Thus
$\pi_f(B)$ separates the states of $\pi_f(A)$, and hence
$\pi_f(B)=\pi_f(A)$ (\XRef{11.3.2}{11.3.2}). It is now enough to apply
\XRef{11.4.2}{11.4.2}.

\ResultHeading{11.4.4}{11.4.4. Lemma.}
\begin{itshape}
Let $H$ be a Hilbert space, let $T\in\mathcal L(H)$ satisfy
$0\leq T\leq1$, let $\varepsilon\in[0,1]$, and let $\xi$ be a unit vector
of $H$ such that $(T\xi\mid\xi)\geq1-\varepsilon$.
\begin{enumerate}[label=\textup{(\roman*)}]
\item If $f:\mathbb R\to\mathbb R$ is a nonnegative Borel function such
that $f\geq1$ on $[1-\varepsilon^{1/2},1]$, then
\[
(f(T)\xi\mid\xi)\geq1-\varepsilon^{1/2}.
\]
\item If $T'\in\mathcal L(H)$, then
\[
|(T'\xi\mid\xi)-(T'T\xi\mid\xi)|
\leq2\varepsilon^{1/2}\|T'\|.
\]
\item We have
\[
T+3\varepsilon^{1/2}\geq P_{\mathbb C\xi}.
\]
\item If $S$ is a Hermitian operator on $H$ such that $S\geq T$, and if
$g:\mathbb R\to\mathbb R$ is a nonnegative Borel function such that
$g\geq1$ on $[\varepsilon,+\infty)$, then
\[
(g(S)\xi\mid\xi)\geq1-4\varepsilon^{1/2}.
\]
\end{enumerate}
\end{itshape}

Let
\[
T=\int_0^1\lambda\,dE_\lambda
\]
be the spectral decomposition of $T$, and put
$1-\varepsilon^{1/2}=\alpha$. We have
\begin{align*}
1-\varepsilon^{1/2}(E_\alpha\xi\mid\xi)
&=(1-\varepsilon^{1/2})(E_\alpha\xi\mid\xi)
+((1-E_\alpha)\xi\mid\xi)\\
&\geq\int_0^\alpha\lambda\,d(E_\lambda\xi\mid\xi)
+\int_\alpha^1\lambda\,d(E_\lambda\xi\mid\xi)\\
&=(T\xi\mid\xi)\geq1-\varepsilon,
\end{align*}
% Source PDF physical page 239
whence $(E_\alpha\xi\mid\xi)\leq\varepsilon^{1/2}$, and
\[
(f(T)\xi\mid\xi)\geq((1-E_\alpha)\xi\mid\xi)
\geq1-\varepsilon^{1/2}.
\]
This proves (i).

If $T'\in\mathcal L(H)$, then, in view of \XRef{9.3.3}{9.3.3},
\begin{align*}
|(T'\xi\mid\xi)-(T'T\xi\mid\xi)|
&=|(T'(\xi-T\xi)\mid\xi)|\\
&\leq\|T'\|\,\|\xi-T\xi\|
\leq(2\varepsilon)^{1/2}\|T'\|,
\end{align*}
which proves (ii).

Let $\zeta\in H$ be a unit vector. Write
$\zeta=\lambda\xi_1+\mu\xi$, where $\lambda,\mu\in\mathbb C$ and $\xi_1$
is a unit vector of $H$ orthogonal to $\xi$. Let $P$ be the orthogonal
projection onto $\mathbb C\xi_1+\mathbb C\xi$, and let
\[
\begin{pmatrix}p&q\\\overline q&r\end{pmatrix}
\]
be the matrix of $PTP$ relative to the basis $(\xi_1,\xi)$. We have
\[
r=(PTP\xi\mid\xi)=(T\xi\mid\xi)\in[1-\varepsilon,1].
\]
Since $0\leq PTP\leq P$, we have $p\in[0,1]$ and
\[
|q|^{2}\leq(1-p)(1-r)\leq\varepsilon.
\]
Therefore
\begin{align*}
((T-P_{\mathbb C\xi})\zeta\mid\zeta)
&=(p\lambda+q\mu)\overline\lambda
+(\overline q\lambda+(r-1)\mu)\overline\mu\\
&\geq p|\lambda|^{2}+(r-1)|\mu|^{2}-2|q\lambda\mu|\\
&\geq0-\varepsilon-2\varepsilon^{1/2}
\geq-3\varepsilon^{1/2},
\end{align*}
which proves (iii).

Let
\[
S=\int_0^{+\infty}\lambda\,dF_\lambda
\]
be the spectral decomposition of $S$. Write
$\xi=\rho\xi_1+\sigma\xi_2$, where $\rho,\sigma\in\mathbb C$ and
$\xi_1,\xi_2$ are unit vectors of $(1-F_\varepsilon)(H)$ and
$F_\varepsilon(H)$, respectively. Let $Q$ be the orthogonal projection
onto $\mathbb C\xi_1+\mathbb C\xi_2$. We have
\[
QSQ=a_1P_{\mathbb C\xi_1}+a_2P_{\mathbb C\xi_2},
\qquad a_1\geq\varepsilon\geq a_2\geq0.
\]
Consequently, by (iii),
\begin{align*}
4\varepsilon^{1/2}
&\geq a_2+3\varepsilon^{1/2}
=((QSQ+3\varepsilon^{1/2})\xi_2\mid\xi_2)\\
&\geq((QTQ+3\varepsilon^{1/2})\xi_2\mid\xi_2)
\geq(P_{\mathbb C\xi}\xi_2\mid\xi_2)
=|(\xi\mid\xi_2)|^{2}=|\sigma|^{2}.
\end{align*}
It follows that
\[
(g(S)\xi\mid\xi)\geq((1-F_\varepsilon)\xi\mid\xi)
=|\rho|^{2}=1-|\sigma|^{2}\geq1-4\varepsilon^{1/2}.
\]

\ResultHeading{11.4.5}{11.4.5. Lemma.}
\begin{itshape}
Let $A$ and $B$ be as in \XRef{11.4.3}{11.4.3}, let $K$ be a closed
two-sided ideal of $A$, let $f\in\overline{P(A)}$, and let
$\varepsilon\in[0,1/4)$. Suppose that there is an $x\in K$ such that
$0\leq x\leq1$ and $f(x)\geq1-\varepsilon$. Then there is a
$y\in B\cap K$ such that $0\leq y\leq1$ and
$f(y)\geq1-\varepsilon^{1/2}$.
\end{itshape}
\par

% ===== ASSEMBLED TRANSLATION CHUNK 09: chunk_240_269.tex =====
% Source PDF physical page 240
By \XRef{11.4.3}{11.4.3}, there exists $x'\in B$ such that
$\pi_f(x)=\pi_f(x')$; replacing $x'$ by
$\frac12(x'+x'^{\ast})$, we may suppose that $x'$ is Hermitian.

If $x'-x\leq\frac12$, let $p\colon\mathbb R\to\mathbb R$ be the function
equal to $0$ on $(-\infty,\frac12]$, equal to $1$ on
$[1-\varepsilon^{1/2},+\infty)$, and linear on
$[\frac12,1-\varepsilon^{1/2}]$; put $y=p(x')\in B$. We have
$0\leq y\leq1$. On the other hand, if $\omega$ is the canonical
homomorphism of $A$ onto $A/K$, then
\[
 \omega(x')=\omega(x'-x)\leq\frac12,
 \qquad \omega(y)=p(\omega(x'))=0,
\]
so that $y\in K$. Finally,
\[
 \pi_f(y)=p(\pi_f(x'))=p(\pi_f(x)),
\]
and hence
\[
 f(y)=(\pi_f(y)\xi_f\mid\xi_f)\geq1-\varepsilon^{1/2}
\]
by \XRef{11.4.4}{11.4.4}(i).

If nothing is known about $x'-x$, we shall modify $x'$ so as to reduce to
the preceding case. If $g\in\overline{P(A)}$, there exists $t_g\in B$ such
that
\[
 \pi_f(t_g)=\pi_f((x'-x)^{+})=0,
 \qquad
 \pi_g(t_g)=\pi_g((x'-x)^{+})
 \qquad\text{(\XRef{11.4.3}{11.4.3})};
\]
replacing $t_g$ by $\frac12(t_g+t_g^{\ast})^{+}$, we may suppose that
$t_g\geq0$. Let $W_g$ be the set of $h\in\overline{P(A)}$ such that
\[
 h(t_g)>h(x'-x)-\frac12;
\]
this is a neighborhood of $g$ in $\overline{P(A)}$, since
\[
 g(t_g)=g((x'-x)^{+})\geq g(x'-x)>g(x'-x)-\frac12.
\]
Let $(W_{g_1},\ldots,W_{g_n})$ be a finite cover of
$\overline{P(A)}$, which is compact, and put
\[
 x''=x'-\sum_{i=1}^{n}t_{g_i}\in B;
\]
we have $\pi_f(x'')=\pi_f(x')=\pi_f(x)$. Moreover, if
$h\in\overline{P(A)}$, then $h\in W_{g_i}$ for some $i$, and therefore
\[
 \begin{aligned}
 h(x''-x)
 &=h(x'-x)-h(t_{g_1})-\cdots-h(t_{g_n})\\
 &\leq h(x'-x)-h(t_{g_i})<\frac12,
 \end{aligned}
\]
hence $x''-x\leq\frac12$, and the proof is complete.

\ResultHeading{11.4.6}{11.4.6. Lemma.}
\begin{itshape}
Let $A$ and $B$ be as in \XRef{11.4.3}{11.4.3}, let $K$ be a closed
two-sided ideal of $A$, let $x_1,\ldots,x_n$ be elements of $K$ such that
$0\leq x_i\leq1$, and let $\varepsilon\in(0,\frac14)$. Let $Q$ be the set
of $f\in\overline{P(A)}$ such that $f(x_i)\geq1-\varepsilon$ for at least
one $i$. There exists $y\in B\cap K$ such that $0\leq y\leq1$ and
\[
 f(y)\geq1-6\varepsilon^{1/4}\qquad(f\in Q).
\]
\end{itshape}

% Source PDF physical page 241
Let $f\in Q$. There is an index $i_f$ such that
$f(x_{i_f})\geq1-\varepsilon$, and then there is $y_f\in B\cap K$ such
that
\[
 0\leq y_f\leq1,\qquad f(y_f)\geq1-\varepsilon^{1/2}
 \qquad\text{(\XRef{11.4.5}{11.4.5})}.
\]
Let $W_f$ be the set of $g\in Q$ such that
$g(y_f)\geq1-2\varepsilon^{1/2}$; this is a neighborhood of $f$ in $Q$.
Let $(W_{f_1},\ldots,W_{f_s})$ be a finite cover of $Q$, and put
$x=y_{f_1}+\cdots+y_{f_s}$. Let $p\colon\mathbb R\to\mathbb R$ be the
function equal to $0$ on $(-\infty,0]$, equal to $1$ on
$[2\varepsilon^{1/2},+\infty)$, and linear on
$[0,2\varepsilon^{1/2}]$. Put $y=p(x)$. Then $y\in B\cap K$ and
$0\leq y\leq1$. If $f\in Q$, then, for example, $f\in W_{f_1}$, and
hence $f(y_{f_1})\geq1-2\varepsilon^{1/2}$. Consequently,
\[
 f(y)\geq1-4(2\varepsilon^{1/2})^{1/2}
       \geq1-6\varepsilon^{1/4}
\]
by \XRef{11.4.4}{11.4.4}(iv).

\ResultHeading{11.4.7}{11.4.7. Lemma.}
\begin{itshape}
Let $A$ and $B$ be as in \XRef{11.4.3}{11.4.3}, let $K$ be a closed
two-sided ideal of $A$, let $f_1,\ldots,f_n$ be states of $A$ such that
each $f_i|_K$ is a state of $K$, and let $\varepsilon>0$. There exists
$y\in B\cap K$ such that $0\leq y\leq1$ and
\[
 f_1(y)\geq1-\varepsilon,\ldots,f_n(y)\geq1-\varepsilon.
\]
\end{itshape}

We may suppose that $(\varepsilon/12)^4\in(0,\frac14)$. For
$1\leq i\leq n$, there exists $x_i\in K$ such that $0\leq x_i\leq1$ and
\[
 f_i(x_i)>1-(\varepsilon/12)^5.
\]
Let $Q$ be the set of $f\in\overline{P(A)}$ such that
$f(x_i)\geq1-(\varepsilon/12)^4$ for at least one $i$. By
\XRef{11.4.6}{11.4.6}, there exists $y\in B\cap K$ such that
$0\leq y\leq1$ and
\[
 f(y)\geq1-\frac{\varepsilon}{2}\qquad(f\in Q).
\]
There are convex combinations
\[
 \lambda_{i1}^{\alpha}f_{i1}^{\alpha}+\cdots+
 \lambda_{ip}^{\alpha}f_{ip}^{\alpha}
\]
of pure states of $A$ which converge weakly to $f_i$. Let $J_\alpha$ be
the set of indices $j\in\{1,2,\ldots,p\}$ such that
\[
 f_{ij}^{\alpha}(x_i)\geq1-(\varepsilon/12)^4,
 \qquad
 \beta_\alpha=\sum_{j\in J_\alpha}\lambda_{ij}^{\alpha}.
\]
Since
\[
 \bigl(\lambda_{i1}^{\alpha}f_{i1}^{\alpha}+\cdots+
 \lambda_{ip}^{\alpha}f_{ip}^{\alpha}\bigr)(x_i)
 \leq\beta_\alpha+\bigl(1-(\varepsilon/12)^4\bigr)(1-\beta_\alpha),
\]
we have, for sufficiently large $\alpha$,
\begin{align*}
 1-(\varepsilon/12)^5
 &\leq\beta_\alpha+\bigl(1-(\varepsilon/12)^4\bigr)(1-\beta_\alpha),\\
 (\varepsilon/12)^4\beta_\alpha
 &\geq(\varepsilon/12)^4-(\varepsilon/12)^5,\\
 \beta_\alpha&\geq1-\frac{\varepsilon}{12}.
\end{align*}
For $j\in J_\alpha$, we have $f_{ij}^{\alpha}\in Q$, and therefore
$f_{ij}^{\alpha}(y)\geq1-\varepsilon/2$. Consequently, for sufficiently
large $\alpha$,
\[
 \biggl(\sum_j\lambda_{ij}^{\alpha}f_{ij}^{\alpha}\biggr)(y)
 \geq\beta_\alpha\left(1-\frac{\varepsilon}{2}\right)
 \geq\left(1-\frac{\varepsilon}{12}\right)
       \left(1-\frac{\varepsilon}{2}\right)
 \geq1-\varepsilon,
\]
and, on passing to the limit, $f_i(y)\geq1-\varepsilon$.

Bibliography: \cite{ref45}.

% Source PDF physical page 242
\BookSubsection{11.5}{Proof of the Theorem}

\ResultHeading{11.5.1}{11.5.1.}
Let $A$ be a unital $C^{*}$-algebra, and let $B$ be a $C^{*}$-subalgebra
containing $1$ and separating $\overline{P(A)}$. Suppose that $B\ne A$.
We shall derive a contradiction.

Let $M$ be the set of continuous Hermitian linear functionals on $A$ which
vanish on $B$ and have norm at most $1$. By the Hahn--Banach theorem,
$M\ne0$. Moreover, $M$ is convex and weakly compact; by the
Krein--Milman theorem, $M$ has at least one extreme point $m\ne0$. Clearly
$\lVert m\rVert=1$. Let $I$ be the largest two-sided ideal of $A$
contained in $\ker m$; this ideal is closed. We have
$B+I\subset\ker m$, and hence $(B+I)/I\ne A/I$. Every state of $A/I$
which is a weak limit of pure states of $A/I$ defines a state of $A$ which
is a weak limit of pure states of $A$. Thus $(B+I)/I$ separates
$\overline{P(A/I)}$.

We shall show that $A/I$ is antiliminal. Let $J$ be a closed two-sided
ideal of $A$ containing $I$ such that $J/I$ is postliminal. By
\XRef{11.1.5}{11.1.5} and \XRef{11.1.7}{11.1.7},
\[
 J/I\subset(B+I)/I,
\]
whence $J\subset B+I\subset\ker m$, and therefore $J\subset I$ by the
definition of $I$. Thus $J/I=0$, proving that $A/I$ is antiliminal. This
first shows that it is enough to establish Theorem \XRef{11.3.1}{11.3.1}
for antiliminal $C^{*}$-algebras. We shall therefore add to our previous
hypotheses the assumption that $A$ is antiliminal.

We shall prove that, in $A/I$, the intersection of two nonzero closed
two-sided ideals is nonzero. By \XRef{11.3.3}{11.3.3}, this will give the
desired contradiction.

Let $K,L$ be two closed two-sided ideals of $A$ not contained in $I$, and
let us prove that $K\cap L\not\subset I$. There are positive functionals
$f_1,f_2$ on $A$ such that
\[
 f_1-f_2=m,\qquad
 \lVert f_1\rVert+\lVert f_2\rVert=\lVert m\rVert=1
 \qquad\text{(\XRef{2.6.4}{2.6.4})}.
\]
Moreover, for $i=1,2$, there are positive functionals $g_i,h_i$ on $A$
such that
\[
 f_i=g_i+h_i,\qquad h_i(K)=0,\qquad
 \lVert g_i|_K\rVert=\lVert g_i\rVert
 \qquad\text{(\XRef{2.11.7}{2.11.7})}.
\]
Let $x\in B$. For every $\varepsilon>0$, there exists $y\in B\cap K$
such that
\[
 \lvert g_i(x)-g_i(xy)\rvert\leq\varepsilon\qquad(i=1,2)
\]
by \XRef{11.4.7}{11.4.7} and \XRef{11.4.4}{11.4.4}(ii). Since
$xy\in K\cap B$, we have
\begin{align*}
 \lvert g_1(x)-g_2(x)\rvert
 &\leq\lvert g_1(xy)-g_2(xy)\rvert+2\varepsilon\\
 &=\lvert f_1(xy)-f_2(xy)\rvert+2\varepsilon=2\varepsilon.
\end{align*}
Thus $(g_1-g_2)(B)=0$, and consequently $(h_1-h_2)(B)=0$. On the other
hand,
\begin{align*}
 1=\lVert m\rVert
 &=\lVert g_1-g_2+h_1-h_2\rVert\\
 &\leq\lVert g_1-g_2\rVert+\lVert h_1-h_2\rVert\\
 &\leq\lVert g_1\rVert+\lVert h_1\rVert+
       \lVert g_2\rVert+\lVert h_2\rVert\\
 &=\lVert f_1\rVert+\lVert f_2\rVert=1.
\end{align*}
Hence, if
\[
 \lambda=\lVert g_1-g_2\rVert,\qquad
 \mu=\lVert h_1-h_2\rVert,
\]
then $\lambda+\mu=1$. Since $K\not\subset I$, we have
$K\not\subset\ker m$, and therefore $g_1-g_2$ cannot be proportional to
$h_1-h_2$, which vanishes on $K$. This proves first that $\lambda\ne0$.
If $\mu\ne0$, the equality
\[
 m=\lambda\bigl(\lambda^{-1}(g_1-g_2)\bigr)
   +\mu\bigl(\mu^{-1}(h_1-h_2)\bigr)
\]
% Source PDF physical page 243
contradicts the fact that $m$ is extreme in $M$. Thus $\mu=0$ and
$m=g_1-g_2$. Let $z\in L$ be such that $m(z)\ne0$. There exists $z'\in K$
such that
\[
 \lvert g_1(z)-g_1(zz')\rvert
 \quad\text{and}\quad
 \lvert g_2(z)-g_2(zz')\rvert
\]
are arbitrarily small, by \XRef{11.4.7}{11.4.7} and
\XRef{11.4.4}{11.4.4}(ii), and hence such that
\[
 m(zz')=g_1(zz')-g_2(zz')\ne0.
\]
But $zz'\in K\cap L$, so $K\cap L\not\subset I$.

\ResultHeading{11.5.2}{11.5.2. Corollary.}
\begin{itshape}
Let $A$ be a $C^{*}$-algebra and $B$ a $C^{*}$-subalgebra separating
$\overline{P(A)}\cup\{0\}$. Then $B=A$.
\end{itshape}

Let $\widetilde A=A+\mathbb C e$ be the $C^{*}$-algebra obtained from $A$
by adjoining an identity element $e$, and let
$\widetilde B=B+\mathbb C e$. Let $f\in\overline{P(\widetilde A)}$.
Then $f$ is a weak limit of pure states $f_\alpha$ of $\widetilde A$. If
$f|_A$ is nonzero, we may suppose that the $f_\alpha|_A$ are nonzero;
they are then pure states of $A$ (\XRef{2.11.8}{2.11.8}), and hence in
this case $f|_A\in\overline{P(A)}$. Thus in every case
$f|_A\in\overline{P(A)}\cup\{0\}$. If $f$ and $g$ are two distinct
elements of $\overline{P(\widetilde A)}$, then $f|_A$ and $g|_A$ are
distinct and belong to $\overline{P(A)}\cup\{0\}$, and are therefore
separated by $B$. Thus $\widetilde B$ separates
$\overline{P(\widetilde A)}$. Theorem \XRef{11.3.1}{11.3.1} shows that
$\widetilde B=\widetilde A$, whence $B=A$.

\ResultHeading{11.5.3}{11.5.3. Corollary.}
\begin{itshape}
Let $T$ be a locally compact space, let
$\mathcal A=((A(t)),\Theta)$ be a continuous field of $C^{*}$-algebras
over $T$, let $A$ be the $C^{*}$-algebra defined by $\mathcal A$, and let
$B$ be a $C^{*}$-subalgebra of $A$ such that, for all $t_1,t_2\in T$ and
$\xi_1\in A(t_1)$, $\xi_2\in A(t_2)$, with $\xi_1=\xi_2$ if $t_1=t_2$,
there exists $x\in B$ satisfying
\[
 x(t_1)=\xi_1,\qquad x(t_2)=\xi_2.
\]
Then $B=A$.
\end{itshape}

By adjoining a point at infinity to $T$, we may suppose that $T$ is
compact (\XRef{10.2.6}{10.2.6}).

Let $f\in\overline{P(A)}$. Then $f$ is a weak limit of pure states
$f_\alpha$ of $A$. By \XRef{10.4.3}{10.4.3}, there are
$t_\alpha\in T$ and a pure state $g_\alpha$ of $A(t_\alpha)$ such that
\[
 f_\alpha(x)=g_\alpha(x(t_\alpha))\qquad(x\in A).
\]
We may suppose that the $t_\alpha$ converge to a point $t\in T$. Let
$x\in A$ satisfy $x(t)=0$. Then
$\lVert x(t_\alpha)\rVert\to\lVert x(t)\rVert=0$, and hence
$g_\alpha(x(t_\alpha))\to0$, so $f(x)=0$. This proves that there is a
positive functional $g$ on $A(t)$ such that
\[
 f(x)=g(x(t))\qquad(x\in A).
\]

Now let $f_1,f_2$ be two distinct points of
$\overline{P(A)}\cup\{0\}$, and let us prove that $B$ separates $f_1$
and $f_2$; then \XRef{11.5.3}{11.5.3} follows from
\XRef{11.5.2}{11.5.2}. Let $t_1,t_2\in T$, and let $g_i$ be a positive
functional on $A(t_i)$, such that
\[
 f_i(x)=g_i(x(t_i))\qquad(x\in A).
\]
At least one of $g_1,g_2$ is nonzero, say $g_1$. If $t_1=t_2$, then
$g_1\ne g_2$, so there exists $\xi\in A(t_1)$ such that
$g_1(\xi)\ne g_2(\xi)$; then there is $x\in B$ with $x(t_1)=\xi$, and
$f_1(x)\ne f_2(x)$. If $t_1\ne t_2$, there is $y\in B$ such that
$y(t_2)=0$ and $g_1(y(t_1))\ne0$, whence $f_1(y)\ne f_2(y)$.

Bibliography: \cite{ref29,ref45,ref167}.

% Source PDF physical page 244
\BookSubsection{11.6}{Complements}

\ResultHeading{11.6.1}{*11.6.1.}
Let $H$ be a Hilbert space, let $A$ be a $C^{*}$-subalgebra of
$\mathcal L(H)$ such that $1\in A$ (respectively, $1\notin A$), and let
$V$ be the weak closure of the set of vector states of $A$. Then $V$ is
the set of functionals
\[
 a(\omega_\xi|_A)+bg,
\]
where $a,b\in[0,1]$, $a+b=1$ (respectively, $a+b\leq1$), $\xi$ is a unit
vector of $\overline{A(H)}$, and $g$ is a state of $A$ such that
$g(A\cap\mathcal{LC}(H))=0$. \cite{ref45,ref46}

\ResultHeading{11.6.2}{*11.6.2.}
Let $A$ be a unital $C^{*}$-algebra, let $K$ be its largest postliminal
two-sided ideal, let $B$ be a $C^{*}$-subalgebra of $A$ containing $1$,
let $S$ be the set of pure states $f$ of $A$ such that $f|_K$ is a
(pure) state of $K$, and let $T$ be the set of states which are weak
limits of pure states of $A$ vanishing on $K$. If $B$ separates $S\cup T$,
then $B=A$. \cite{ref45}

\ResultHeading{11.6.3}{*11.6.3.}
Let $A$ be a von Neumann algebra, let $B$ be an antiliminal
$C^{*}$-subalgebra of $A$, weakly dense in $A$ and containing the center
$\mathfrak Z$ of $A$. Then $\overline{P(B)}$ is the set of states $f$ of
$B$ such that $f|_{\mathfrak Z}$ is a character of $\mathfrak Z$. In
particular, if $F$ is a continuous factor, then
$\overline{P(F)}=E(F)$. \cite{ref45}

\ResultHeading{11.6.4}{*11.6.4.}
Let $A$ be a von Neumann algebra and $\mathfrak Z$ its center. Then
$\overline{P(A)}$ is the set of functionals
\[
 ag+(1-a)h,
\]
where $a\in[0,1]$, $g,h\in E(A)$, the restriction
$(ag+(1-a)h)|_{\mathfrak Z}$ is a character of $\mathfrak Z$, there is an
abelian projection $e$ of $A$ such that $ag(e)=a$, and $(1-a)h$ vanishes
on every abelian projection of $A$. \cite{ref45}\translatornote{The source
prints $f$ in the first occurrence of this convex combination, but
subsequently uses $h$; the latter is required by the stated variables.}

\ResultHeading{11.6.5}{*11.6.5.}
Let $A$ be a von Neumann algebra, and let $B$ be a $C^{*}$-subalgebra of
$A$ containing $1$ and weakly dense in $A$. Then $\overline{P(B)}$ is the
set of restrictions to $B$ of the elements of $\overline{P(A)}$.
\cite{ref45}

\ResultHeading{11.6.6}{11.6.6.}
Let $A$ be a unital $C^{*}$-algebra. The following conditions are
equivalent:

\begin{enumerate}[label=\textup{(\roman*)}]
 \item $\overline{P(A)}=E(A)$;
 \item $A$ is antiliminal and the zero ideal is prime.
\end{enumerate}

[The implication (ii)$\Rightarrow$(i) follows from
\XRef{11.2.4}{11.2.4}. Let $I,J$ be nonzero closed two-sided ideals of
$A$ such that $I\cap J=0$. Let $f$ (respectively, $g$) be a pure state of
$A$ whose restriction to $I$ (respectively, $J$) is nonzero. Then the
restriction of $g$ (respectively, $f$) to $I$ (respectively, $J$) is zero,
and $\frac12(f+g)\notin\overline{P(A)}$.] \cite{ref169}

\BookSectionHere[Enveloping von Neumann Algebra of a C*-Algebra]{12}%
{Enveloping von Neumann Algebra of a $C^{*}$-Algebra}

\BookSubsection[Bidual of a C*-Algebra]{12.1}{Bidual of a $C^{*}$-Algebra}

\ResultHeading{12.1.1}{12.1.1. Proposition.}
\begin{itshape}
Let $H$ be a Hilbert space, let $A$ be a $C^{*}$-subalgebra of
$\mathcal L(H)$, and let $B$ be the weak closure of $A$. Suppose that
$1\in B$ and that every norm-continuous linear functional $f$ on $A$ is
ultrastrongly continuous, and hence has a unique ultrastrongly continuous
extension $\widetilde f$ to $B$.
% Source PDF physical page 245
\begin{enumerate}[label=\textup{(\roman*)}]
 \item The map $f\mapsto\widetilde f$ is an isometric isomorphism from the
 Banach dual $A'$ of $A$ onto the predual of $B$.
 \item For each $y\in B$, let $\widehat y$ be the linear functional
 $f\mapsto\widetilde f(y)$ on $A'$. Then $y\mapsto\widehat y$ is an
 isometric isomorphism from $B$ onto the Banach dual $A''$ of $A'$, whose
 restriction to $A$ is the canonical injection of $A$ into $A''$.
\end{enumerate}
\end{itshape}

If $f\in A'$, then $\widetilde f$ belongs to the predual of $B$ by
hypothesis. Let $A_1,B_1$ be the unit balls of $A,B$. Since $A_1$ is
ultrastrongly dense in $B_1$ (\XRef{A.13}{A 13}),
\[
 \sup_{y\in B_1}\lvert\widetilde f(y)\rvert
 =\sup_{x\in A_1}\lvert f(x)\rvert,
\]
that is, $\lVert\widetilde f\rVert=\lVert f\rVert$. Conversely, every
element of the predual of $B$ is the ultrastrongly continuous extension of
its restriction to $A$. This proves (i). In view of (i) and
\XRef{A.23}{A 23}, $y\mapsto\widehat y$ is an isometric isomorphism of
$B$ onto $A''$; if $y\in A$, then $\widehat y$ is the functional
$f\mapsto f(y)$ on $A'$, that is, the canonical image of $A$ in $A''$.

\ResultHeading{12.1.2}{12.1.2. Corollary.}
\begin{itshape}
Let $H$ be a Hilbert space and let $A$ be the $C^{*}$-algebra of compact
operators on $H$.

\begin{enumerate}[label=\textup{(\roman*)}]
 \item Every linear functional $f\in A'$ is ultrastrongly continuous on
 $A$, and hence has a unique ultrastrongly continuous extension
 $\widetilde f$ to $\mathcal L(H)$.
 \item The map $f\mapsto\widetilde f$ is an isometric isomorphism from
 $A'$ onto the predual of $\mathcal L(H)$.
 \item For every $y\in\mathcal L(H)$, let $\widehat y$ be the linear
 functional $f\mapsto\widetilde f(y)$ on $A'$. Then
 $y\mapsto\widehat y$ is an isometric isomorphism from $\mathcal L(H)$
 onto $A''$, whose restriction to $A$ is the canonical injection of $A$
 into $A''$.
\end{enumerate}
\end{itshape}

Assertion (i) follows from \XRef{2.6.4}{2.6.4} and
\XRef{4.1.3}{4.1.3}, and (ii), (iii) are then special cases of
\XRef{12.1.1}{12.1.1}.

\ResultHeading{12.1.3}{12.1.3. Corollary.}
\begin{itshape}
Let $A$ be a $C^{*}$-algebra, let $\pi$ be the universal representation
of $A$ (\XRef{2.7.6}{2.7.6}), and let $B$ be the weak closure of $\pi(A)$,
which is a von Neumann algebra on $H_\pi$.

\begin{enumerate}[label=\textup{(\roman*)}]
 \item Every normal positive functional on $B$ is a functional
 $\omega_\xi$, with $\xi\in H_\pi$. Every ultrastrongly continuous linear
 functional on $B$ is weakly continuous.
 \item If $f\in A'$, there is one and only one weakly continuous linear
 functional $\widetilde f$ on $B$ such that
 \[
   \widetilde f(\pi(x))=f(x)\qquad(x\in A).
 \]
 \item The map $f\mapsto\widetilde f$ is an isometric isomorphism from the
 dual $A'$ of $A$ onto the predual of $B$; it transforms the positive
 functionals on $A$ into the normal positive functionals on $B$. For every
 $f\in A'$ one has
 \[
   \widetilde{f^{\ast}}=(\widetilde f)^{\ast}.
 \]
 \item For every $y\in B$, let $\widehat y$ be the linear functional
 $f\mapsto\widetilde f(y)$ on $A'$. Then $y\mapsto\widehat y$ is an
 isometric isomorphism from $B$ onto $A''$, whose composite with $\pi$ is
 the canonical injection of $A$ into $A''$.

% Source PDF physical page 246
 \item This isomorphism is bicontinuous for the weak operator topology on
 $B$ and the topology $\sigma(A'',A')$ on $A''$.
\end{enumerate}
\end{itshape}

Let $f$ be a normal positive functional on $B$. Then $f|_{\pi(A)}$ is a
positive functional on $\pi(A)$. By the construction of $\pi$, there is
$\xi\in H_\pi$ such that
\[
 f(\pi(x))=(\pi(x)\xi\mid\xi)\qquad(x\in A).
\]
The equality $f(y)=\omega_\xi(y)$, valid for $y\in\pi(A)$, remains true
by continuity for $y\in B$. Since every ultrastrongly continuous linear
functional on $B$ is a linear combination of normal positive functionals
on $B$ (\XRef{A.25}{A 25}), this proves (i).

We have just observed that every positive functional on $\pi(A)$ is
weakly continuous. The same therefore holds for every continuous linear
functional on $\pi(A)$ (\XRef{2.6.4}{2.6.4}), which proves (ii). By
\XRef{12.1.1}{12.1.1}, $f\mapsto\widetilde f$ is an isometric
isomorphism from $A'$ onto the predual of $B$.\translatornote{The source
prints $A$ here; $B$ is required by \XRef{12.1.1}{12.1.1} and by the
definition of the predual.} We have seen that, if
$f\geq0$, then $\widetilde f$ is a functional $\omega_\xi$, and is
therefore normal and positive, and that all normal positive functionals
on $B$ arise in this way. If $f\in A'$, then, for $x\in A$,
\[
 \widetilde{f^{\ast}}(\pi(x))
 =f^{\ast}(x)=\overline{f(x^{\ast})}
 =\overline{\widetilde f(\pi(x)^{\ast})}.
\]
Thus, by weak continuity,
\[
 \widetilde{f^{\ast}}(y)=\overline{\widetilde f(y^{\ast})}
 \qquad(y\in B),
\]
that is, $\widetilde{f^{\ast}}=(\widetilde f)^{\ast}$. This proves (iii),
and (iv) is a consequence of \XRef{12.1.1}{12.1.1}. By (i) and (iii), the
weak operator topology on $B$ is defined by the linear functionals
$y\mapsto\widetilde f(y)$, where $f\in A'$; the isomorphism
$y\mapsto\widehat y$ carries this topology into $\sigma(A'',A')$, since
$\widetilde f(y)=\widehat y(f)$. This proves (v).

\ResultHeading{12.1.4}{12.1.4. Definition.}
\begin{itshape}
The algebra $B$ is called the \emph{enveloping von Neumann algebra} of
$A$, and $\pi$ the canonical homomorphism of $A$ into $B$.
\end{itshape}

We sometimes identify $A$ with $\pi(A)$ and $B$ with $A''$.

\ResultHeading{12.1.5}{12.1.5.}
The pair $(B,\pi)$ solves a universal problem. Indeed:

\ResultLabel{Proposition.}
\begin{itshape}
Let $A$ be a $C^{*}$-algebra, let $B$ be the enveloping von Neumann algebra
of $A$, and let $\pi\colon A\to B$ be the canonical homomorphism. Let
$\rho$ be a representation of $A$. There exists one and only one normal
representation $\widetilde\rho$ of $B$ on $H_\rho$ such that
\[
 \widetilde\rho(\pi(x))=\rho(x)\qquad(x\in A).
\]
The image $\widetilde\rho(B)$ is the weak closure of $\rho(A)$.
\end{itshape}

A normal representation is ultraweakly continuous (\XRef{A.27}{A 27});
since $\pi(A)$ is ultraweakly dense in $B$, this proves the uniqueness of
$\widetilde\rho$. To prove its existence, we may suppose that $\rho$
admits a cyclic vector, since every nondegenerate representation of $A$
is a Hilbert direct sum of representations admitting cyclic vectors
(\XRef{2.2.7}{2.2.7}). Then there is a positive functional $f$ on $A$ such
that $\rho$ may be identified with $\pi_f$ (\XRef{2.4.6}{2.4.6}). By the
construction of $\pi$, there is a closed vector subspace $K$ of $H_\pi$,
invariant under $\pi(A)$, such that $\rho$ may be identified with the
representation $x\mapsto\pi(x)|_K$. We may therefore take
$\widetilde\rho$ to be the representation
% Source PDF physical page 247
$y\mapsto y|_K$ of $B$ on $K$. Finally, the image of $B$ under a normal
representation is weakly closed (\XRef{A.27}{A 27}), which proves the last
assertion of the proposition.

\ResultHeading{12.1.6}{12.1.6.}
Recall that, if $E$ is a Banach space and $E',E'',\ldots$ denote its
successive duals, the canonical injection of $E$ into $E''$ has as its
transpose a norm-one projection, called canonical, from $E'''$ onto $E'$.
The kernel of this projection is the set of continuous linear functionals
on $E''$ which vanish on $E$.

Now let $A$ be a von Neumann algebra, let $A'$ be its dual, and let $A''$
be its bidual, identified with the enveloping von Neumann algebra of $A$.
Recall that $A$ is identified with the dual $P'$ of its predual $P$, a
closed vector subspace of $A'$.

\ResultLabel{Proposition.}
\begin{itshape}
Let $\rho$ be the normal representation of $A''$ on $A$ extending the
identity map of $A$ (\XRef{12.1.5}{12.1.5}). If $A$ and $A''$ are
identified with $P'$ and $P'''$, respectively, then $\rho$ is the
canonical projection from $P'''$ onto $P'$. The kernel of this projection
is a weakly closed two-sided ideal of $A''$.
\end{itshape}

The canonical projection $\pi$ from $P'''$ onto $P'$ is continuous for
\[
 \sigma(P''',P'')=\sigma(A'',A')
 \qquad\text{and}\qquad
 \sigma(P',P)=\sigma(A,P).
\]
In view of \XRef{12.1.3}{12.1.3}(i),(v), $\rho$ is continuous for
$\sigma(A'',A')$ and $\sigma(A,P)$. Now $\pi$ and $\rho$ both reduce on
$A$ to the identity map, and $A$ is dense in $A''$ for
$\sigma(A'',A')$. Hence $\pi=\rho$. The kernel of $\rho$ is a two-sided
ideal of $A''$, ultraweakly closed and therefore weakly closed.

\ResultHeading{12.1.7}{12.1.7.}
Corollary \XRef{12.1.3}{12.1.3} reduces questions concerning continuous
linear functionals on a $C^{*}$-algebra to questions concerning
ultraweakly continuous linear functionals on a von Neumann algebra. The
latter questions are even, in certain respects, more general, since the
predual of a von Neumann algebra $A$ is not always the dual of a
$C^{*}$-subalgebra of $A$.

Bibliography: \cite{ref55,ref141,ref150}.

\BookSubsection{12.2}{Polar Decomposition of a Linear Functional}

\ResultHeading{12.2.1}{12.2.1.}
Let $A$ be a $C^{*}$-algebra, let $A'$ be its dual, let $f\in A'$, and let
$x\in A$. We denote by $x\mathbin{\cdot}f$ and
$f\mathbin{\cdot}x$ the linear functionals
\[
 y\longmapsto f(xy),\qquad y\longmapsto f(yx)
\]
on $A$, which are elements of $A'$. We have
\[
 \lVert x\mathbin{\cdot}f\rVert\leq\lVert x\rVert\lVert f\rVert,
 \qquad
 \lVert f\mathbin{\cdot}x\rVert\leq\lVert x\rVert\lVert f\rVert.
\]
The operators $f\mapsto f\mathbin{\cdot}x$ (respectively,
$f\mapsto x\mathbin{\cdot}f$) define a representation of $A$
(respectively, of the opposite algebra) in the non-Hilbertian vector space
$A'$. For $x,y\in A$ and $f\in A'$, we have
\[
 \begin{aligned}
 (x\mathbin{\cdot}f)^{\ast}(y)
 &=\overline{(x\mathbin{\cdot}f)(y^{\ast})}
  =\overline{f(xy^{\ast})}\\
 &=f^{\ast}(yx^{\ast})
  =(f^{\ast}\mathbin{\cdot}x^{\ast})(y),
 \end{aligned}
\]
% Source PDF physical page 248
and hence
\[
 (x\mathbin{\cdot}f)^{\ast}=f^{\ast}\mathbin{\cdot}x^{\ast}.
\]
If $A$ is a von Neumann algebra and $f$ is ultraweakly continuous
(respectively, weakly continuous), then so are
$x\mathbin{\cdot}f$ and $f\mathbin{\cdot}x$.

\ResultHeading{12.2.2}{12.2.2. Lemma.}
\begin{itshape}
Let $A$ be a unital $C^{*}$-algebra, let $A_1$ be the ball
$\lVert x\rVert\leq1$ in $A$, and let $t$ be an extreme point of $A_1$.
Then $t^{\ast}t$ is idempotent.
\end{itshape}

We have $0\leq t^{\ast}t\leq1$ and $\lVert t^{\ast}t\rVert=1$. Let $B$
be the commutative $C^{*}$-subalgebra of $A$ generated by $1$ and
$t^{\ast}t$, let $\Omega$ be its compact spectrum, let $B'$ be the algebra
of continuous complex-valued functions on $\Omega$, let $\varphi$ be the
Gelfand isomorphism of $B$ onto $B'$, and put
$f=\varphi(t^{\ast}t)$. Then $0\leq f\leq1$, and $f$ takes the value $1$.
Suppose that there is $\omega\in\Omega$ such that $0<f(\omega)<1$. There
is then a continuous function $g\colon\Omega\to\mathbb R$, with
$g\geq0$, such that
\[
 fg\ne0,\qquad
 \sup f(1+g)^2=\sup f(1-g)^2=1
\]
(it is enough to take the support of $g$ in a sufficiently small
neighborhood of $\omega$, and $\lvert g\rvert$ sufficiently small). Thus
there exists $u\in B^{+}$ such that
\[
 \lVert t^{\ast}t(1+u)^2\rVert
 =\lVert t^{\ast}t(1-u)^2\rVert=1.
\]
Hence $t(1+u),t(1-u)\in A_1$. The equality
\[
 t=\frac12\bigl(t(1+u)+t(1-u)\bigr)
\]
implies, since $t$ is extreme, that $t=t(1+u)=t(1-u)$, whence $tu=0$ and
$fg=0$, a contradiction. Therefore $f$ takes only the values $0$ and $1$,
so $t^{\ast}t$ is a projection.

\ResultHeading{12.2.3}{12.2.3. Lemma.}
\begin{itshape}
Let $A$ be a von Neumann algebra, let $e$ be a projection of $A$, and let
$f$ belong to the predual of $A$.

\begin{enumerate}[label=\textup{(\roman*)}]
 \item
 \[
   \lVert f\rVert^2\geq
   \lVert e\mathbin{\cdot}f\rVert^2+
   \lVert(1-e)\mathbin{\cdot}f\rVert^2.
 \]
 \item If $\lVert f\rVert=\lVert e\mathbin{\cdot}f\rVert$, then
 $f=e\mathbin{\cdot}f$.
\end{enumerate}
\end{itshape}

Assertion (ii) follows from (i). We prove (i). Let $H$ be the Hilbert
space on which $A$ acts. Let $P$ be the predual of $A$, let $Q$ be the
predual of $\mathcal L(H)$, and let $L$ be the orthogonal of $A$ in $Q$.
Since $A$ is closed in $\mathcal L(H)$ for $\sigma(\mathcal L(H),Q)$,
$A$ is the orthogonal of $L$ in $\mathcal L(H)$. Thus $A$ is canonically
identified with the dual of the Banach space $Q/L$. Let $f\in P$. By the
Hahn--Banach theorem, $f$ has an extension $f'\in Q$. The element $f'+L$
of $Q/L$ defines the functional $f$ on $A$. Hence $\lVert f\rVert$ is the
norm of $f'+L$ in $Q/L$, that is,
\[
 \lVert f\rVert=\inf_{g\in f'+L}\lVert g\rVert.
\]
There is therefore a functional $g\in Q$ extending $f$ and satisfying
$\lVert g\rVert\leq\lVert f\rVert+\varepsilon$ (this is, incidentally,
\XRef{12.5.2}{12.5.2}). If we prove that
\[
 \lVert g\rVert^2\geq
 \lVert e\mathbin{\cdot}g\rVert^2+
 \lVert(1-e)\mathbin{\cdot}g\rVert^2,
\]
we shall deduce
\[
 (\lVert f\rVert+\varepsilon)^2\geq
 \lVert e\mathbin{\cdot}f\rVert^2+
 \lVert(1-e)\mathbin{\cdot}f\rVert^2,
\]
and (i) will follow from the arbitrariness of $\varepsilon$. We are thus
reduced to the case $A=\mathcal L(H)$.
% Source PDF physical page 249
Since the weakly continuous linear functionals on $\mathcal L(H)$ are
norm-dense in $Q$, we may suppose that $f$ is weakly continuous. There
are therefore two orthonormal systems
\[
 (\xi_1,\ldots,\xi_n),\qquad(\xi'_1,\ldots,\xi'_n)
\]
in $H$, and numbers $\lambda_1\geq0,\ldots,\lambda_n\geq0$, such that
\[
 \lVert f\rVert=\lambda_1+\cdots+\lambda_n,
 \qquad
 f(x)=\sum_i\lambda_i(x\xi_i\mid\xi'_i)
 \qquad(x\in A)
\]
(\XRef{A.24}{A 24}). Then
\[
 (e\mathbin{\cdot}f)(x)
 =\sum_i\lambda_i(x\xi_i\mid e\xi'_i),
 \qquad
 ((1-e)\mathbin{\cdot}f)(x)
 =\sum_i\lambda_i(x\xi_i\mid(1-e)\xi'_i).
\]
Consequently,
\begin{align*}
 &\lVert e\mathbin{\cdot}f\rVert^2+
  \lVert(1-e)\mathbin{\cdot}f\rVert^2\\
 &\leq
 \left(\sum_i\lambda_i\lVert e\xi'_i\rVert\right)^2
 +\left(\sum_i\lambda_i\lVert(1-e)\xi'_i\rVert\right)^2\\
 &=\sum_{i=1}^{n}\lambda_i^2
   \bigl(\lVert e\xi'_i\rVert^2+
         \lVert(1-e)\xi'_i\rVert^2\bigr)\\
 &\quad+2\sum_{1\leq i<j\leq n}\lambda_i\lambda_j
 \bigl(\lVert e\xi'_i\rVert\lVert e\xi'_j\rVert+
       \lVert(1-e)\xi'_i\rVert\lVert(1-e)\xi'_j\rVert\bigr)\\
 &\leq\sum_{i=1}^{n}\lambda_i^2
   +2\sum_{1\leq i<j\leq n}\lambda_i\lambda_j
 =\left(\sum_{i=1}^{n}\lambda_i\right)^2
 =\lVert f\rVert^2.
\end{align*}

\ResultHeading{12.2.4}{12.2.4. Theorem.}
\begin{itshape}
Let $A$ be a von Neumann algebra and let $f$ belong to the predual of
$A$.

\begin{enumerate}[label=\textup{(\roman*)}]
 \item There exists a pair $(p,u)$ with the following properties:
 \begin{enumerate}[label=\textup{\alph*.}]
  \item $p$ is a normal positive functional on $A$, and
  $\lVert p\rVert=\lVert f\rVert$;
  \item $u$ is a partial isometry in $A$ whose final projection is the
  support of $p$ (\XRef{A.26}{A 26});
  \item $f=u\mathbin{\cdot}p$ and
  $p=u^{\ast}\mathbin{\cdot}f$.
 \end{enumerate}
 \item Let $p'$ be a normal positive functional on $A$ and $u'$ a partial
 isometry in $A$ whose final projection is majorized by the support of
 $p'$, with
 \[
   f=u'\mathbin{\cdot}p',\qquad
   p'=u'^{\ast}\mathbin{\cdot}f.
 \]
 Then $p'=p$ and $u'=u$.
\end{enumerate}
\end{itshape}

We may suppose that $\lVert f\rVert=1$. Let $B$ be the ball
$\lVert x\rVert\leq1$ in $A$, and let $B'$ be the set of $x\in B$ such
that $f(x)=1$. Since $B$ is ultraweakly compact and $f$ is ultraweakly
continuous, there are $x\in B$ such that $\lvert f(x)\rvert=1$.
Multiplying $x$ by a suitable scalar shows that $B'\ne\varnothing$;
moreover, $B'$ is convex and ultraweakly compact.
% Source PDF physical page 250
There is therefore an extreme point $v$ in $B'$. This point is also
extreme in $B$, for if $v=\frac12(s+t)$ with $s,t\in B$, then
\[
 \lvert f(s)\rvert\leq1,\qquad \lvert f(t)\rvert\leq1,
\]
and
\[
 1=f(v)=\frac12\bigl(f(s)+f(t)\bigr),
 \qquad\text{whence}\qquad f(s)=f(t)=1.
\]
Thus $s,t\in B'$, and $s=t=v$ because $v$ is extreme in $B'$. By
\XRef{12.2.2}{12.2.2}, $v^{\ast}v$ is a projection, so $v$ is a partial
isometry.

Put $p=v\mathbin{\cdot}f$. We have
$\lVert p\rVert\leq\lVert v\rVert\lVert f\rVert=1$. Since
$p(1)=f(v)=1$, $p$ is positive (\XRef{2.1.9}{2.1.9}), has norm $1$, and
is normal because it is ultraweakly continuous. Let $e$ be the support of
$p$. We have
\[
 p(v^{\ast}v)=f(vv^{\ast}v)=f(v)=1,
\]
and hence $v^{\ast}v\geq e$. Put $u=ev^{\ast}$. Then
\[
 uu^{\ast}=e(v^{\ast}v)e=e,
\]
so $u$ is a partial isometry with final projection $e$. For every $x\in A$,
\[
 (u^{\ast}\mathbin{\cdot}f)(x)
 =f(u^{\ast}x)=f(vex)=p(ex)=p(x),
\]
and therefore $p=u^{\ast}\mathbin{\cdot}f$. Put $u^{\ast}u=e'$. We have
$\lVert e'\mathbin{\cdot}f\rVert\leq1$ and
\[
 (e'\mathbin{\cdot}f)(u^{\ast})
 =f(u^{\ast}uu^{\ast})=f(u^{\ast})=f(ve)=p(e)=1.
\]
Thus $\lVert e'\mathbin{\cdot}f\rVert=1=\lVert f\rVert$, and hence
$e'\mathbin{\cdot}f=f$ by \XRef{12.2.3}{12.2.3}(ii). Consequently, for
every $x\in A$,
\[
 (u\mathbin{\cdot}p)(x)=p(ux)=f(u^{\ast}ux)
 =f(e'x)=(e'\mathbin{\cdot}f)(x)=f(x),
\]
and therefore $f=u\mathbin{\cdot}p$. This proves (i).

\ResultHeading{12.2.5}{12.2.5.}
To prove (ii), we first establish the following proposition.

\ResultLabel{Proposition.}
\begin{itshape}
Let $A$ be a unital $C^{*}$-algebra, let $f$ be a continuous linear
functional on $A$, let $p$ be a positive functional on $A$, and let $a$
be an element of $A$ such that
\[
 \lVert a\rVert\leq1,\qquad
 a\mathbin{\cdot}p=f,\qquad
 \lVert p\rVert=\lVert f\rVert.
\]
If $b\in A$ satisfies
\[
 \lVert b\rVert\leq1,\qquad
 b\mathbin{\cdot}f\geq0,\qquad
 \lVert b\mathbin{\cdot}f\rVert=\lVert f\rVert,
\]
then $b\mathbin{\cdot}f=p$.
\end{itshape}

We may suppose that $\lVert p\rVert=\lVert f\rVert=1$. We have
\begin{align*}
 (\xi_p\mid\pi_p(b^{\ast}a^{\ast})\xi_p)
 &=(\pi_p(ab)\xi_p\mid\xi_p)\\
 &=p(ab)=f(b)=(b\mathbin{\cdot}f)(1)
 =\lVert b\mathbin{\cdot}f\rVert=1.
\end{align*}
Since $\lVert b^{\ast}a^{\ast}\rVert\leq1$, equality in the
Cauchy--Schwarz inequality shows that
\[
 \pi_p(b^{\ast}a^{\ast})\xi_p=\xi_p.
\]
Thus, for every $x\in A$,
\begin{align*}
 (b\mathbin{\cdot}f)(x)
 &=f(bx)=p(abx)=(\pi_p(abx)\xi_p\mid\xi_p)\\
 &=(\pi_p(x)\xi_p\mid\pi_p(b^{\ast}a^{\ast})\xi_p)
  =(\pi_p(x)\xi_p\mid\xi_p)=p(x),
\end{align*}
which proves $b\mathbin{\cdot}f=p$.

% Source PDF physical page 251
\ResultHeading{12.2.6}{12.2.6.}
We return to the proof of \XRef{12.2.4}{12.2.4}. Let $(p',u')$ be a pair
having the properties in (ii). We have
$\lVert u'^{\ast}\rVert\leq1$, $u'^{\ast}\mathbin{\cdot}f\geq0$, and
\[
 \lVert p'\rVert
 =\lVert u'^{\ast}\mathbin{\cdot}f\rVert
 \leq\lVert f\rVert
 =\lVert u'\mathbin{\cdot}p'\rVert
 \leq\lVert p'\rVert,
\]
so $\lVert u'^{\ast}\mathbin{\cdot}f\rVert=\lVert f\rVert$. Proposition
\XRef{12.2.5}{12.2.5}, applied with $a=u$ and $b=u'^{\ast}$, gives
\[
 p=u'^{\ast}\mathbin{\cdot}f=p'.
\]
We shall prove that $u=u'$. Put $x=uu'^{\ast}$. Then
\[
 p(x)=(u\mathbin{\cdot}p)(u'^{\ast})=f(u'^{\ast})
 =(u'^{\ast}\mathbin{\cdot}f)(1)=p(1)=1.
\]
Thus $p(x^{\ast})=1$ and
\[
 1=p(x)^2\leq p(xx^{\ast})\leq1,
\]
so $p(xx^{\ast})=1$, and
\[
 p((e-x)(e-x)^{\ast})=1-1-1+1=0.
\]
Now $exe=x$, since the final projections of $u$ and $u'$ are majorized by
$e$. Since $p$ is faithful on $eAe$, it follows that $x=e$. Let $H$ be
the Hilbert space on which $A$ acts, and let $e_1$ be the initial
projection of $u$. For $\xi\in e(H)$,
\[
 \lVert u(u'^{\ast}\xi)\rVert=\lVert e\xi\rVert=\lVert\xi\rVert,
 \qquad\text{hence}\qquad
 \lVert u(u'^{\ast}\xi)\rVert\geq\lVert u'^{\ast}\xi\rVert.
\]
Thus $u'^{\ast}\xi\in e_1(H)$. Since $u$ maps $e_1(H)$ isometrically
onto $e(H)$, the equality
\[
 uu'^{\ast}\xi=\xi=uu^{\ast}\xi
\]
implies $u'^{\ast}\xi=u^{\ast}\xi$. Moreover, $u^{\ast}$ and
$u'^{\ast}$ vanish on $H\ominus e(H)$. Hence $u'^{\ast}=u^{\ast}$ and
$u'=u$.

\ResultHeading{12.2.7}{12.2.7. Definition.}
\begin{itshape}
With the notation of \XRef{12.2.4}{12.2.4}, $p$ is called the
\emph{absolute value} of $f$ and is denoted by $\lvert f\rvert$. The
equality $f=u\mathbin{\cdot}\lvert f\rvert$ is called the
\emph{polar decomposition} of $f$.
\end{itshape}

\ResultHeading{12.2.8}{12.2.8. Definition.}
\begin{itshape}
Let $A$ be a $C^{*}$-algebra, let $E$ be the enveloping von Neumann
algebra of $A$, and let $f$ be a continuous linear functional on $A$. The
polar decomposition $f=u\mathbin{\cdot}p$ of $f$, regarded as an element
of the predual of $E$, is called the \emph{enveloping polar decomposition}
of $f$. The functional $p$, regarded as a positive functional on $A$, is
called the \emph{absolute value} of $f$ and is denoted by $\lvert f\rvert$.
\end{itshape}

If $A$ is a von Neumann algebra and $f$ belongs to its predual, then $f$
has both a polar decomposition and an enveloping polar decomposition; in
general these are distinct (\XRef{12.5.7}{12.5.7}). We shall nevertheless
see that the two possible definitions of $\lvert f\rvert$ coincide. Let
$f=u\mathbin{\cdot}p$ be the polar decomposition of $f$, and identify $p$
with a normal positive functional on $E$ (\XRef{12.1.3}{12.1.3}(iii)). We
have
\[
 u^{\ast}\in A\subset E,\qquad
 \lVert u^{\ast}\rVert\leq1,\qquad
 \lVert p\rVert=\lVert f\rVert.
\]
The equality $p(x)=f(u^{\ast}x)$, valid for $x\in A$, extends by
continuity to every $x\in E$, where $f$ is identified with a weakly
continuous linear functional on $E$. Proposition \XRef{12.2.5}{12.2.5},
applied to $E$ and $b=u^{\ast}$, then proves that $p$ is the absolute
value of $f$ in the sense of \XRef{12.2.8}{12.2.8}.

% Source PDF physical page 252
\ResultHeading{12.2.9}{12.2.9. Proposition.}
\begin{itshape}
Let $A$ be a $C^{*}$-algebra, let $f$ be a continuous linear functional
on $A$, and let $p$ be a positive functional on $A$. In order that
$p=\lvert f\rvert$, it is necessary and sufficient that
\[
 \lVert p\rVert=\lVert f\rVert,
 \qquad
 \lvert f(x)\rvert^2\leq\lVert f\rVert p(x^{\ast}x)
 \qquad(x\in A).
\]
\end{itshape}

Let $f=u\mathbin{\cdot}\lvert f\rvert$ be the enveloping polar
decomposition of $f$. If $p=\lvert f\rvert$, then
$\lVert p\rVert=\lVert f\rVert$, and, for every $x\in A$,
\[
 \lvert f(x)\rvert^2
 =\lvert p(ux)\rvert^2
 \leq p(u^{\ast}u)p(x^{\ast}x)
 \leq\lVert p\rVert p(x^{\ast}x)
 =\lVert f\rVert p(x^{\ast}x).
\]

Conversely, suppose that $\lVert p\rVert=\lVert f\rVert$ and that
\[
 \lvert f(x)\rvert^2\leq\lVert f\rVert p(x^{\ast}x)
 \qquad(x\in A).
\]
Then
\[
 \lvert f(x)\rvert^2
 \leq\lVert f\rVert(\pi_p(x^{\ast}x)\xi_p\mid\xi_p)
 =\lVert f\rVert\lVert\pi_p(x)\xi_p\rVert^2
\]
for every $x\in A$, and hence for every $x$ in the enveloping von Neumann
algebra $E$ of $A$; here $\pi_p$ is extended to a normal representation
of $E$, the function $x\mapsto\lVert\pi_p(x)\xi_p\rVert$ is continuous
for the ultrastrong topology, and $A$ is ultrastrongly dense in $E$.
Applying this inequality to $u^{\ast}x$ gives
\[
 \lvert f\rvert(x)^2
 \leq\lVert f\rVert\lVert\pi_p(u^{\ast}x)\xi_p\rVert^2
 \leq\lVert f\rVert\lVert\pi_p(x)\xi_p\rVert^2
 \qquad(x\in E).
\]
There is therefore a vector $\xi$ in the representation space of $\pi_p$
such that $\lVert\xi\rVert\leq\lVert f\rVert^{1/2}$ and
\[
 \lvert f\rvert(x)=(\pi_p(x)\xi_p\mid\xi).
\]
Consequently,
\[
 (\xi_p\mid\xi)=\lvert f\rvert(1)=\lVert f\rVert
 =\lVert p\rVert^{1/2}\lVert f\rVert^{1/2}
 \geq\lVert\xi_p\rVert\lVert\xi\rVert,
\]
so $\xi=\xi_p$ and
\[
 \lvert f\rvert(x)=(\pi_p(x)\xi_p\mid\xi_p)=p(x).
\]

Bibliography: \cite{ref21,ref65,ref128,ref131,ref163}.

\BookSubsection{12.3}{Decomposition of a Hermitian Functional into Positive
and Negative Parts}

\ResultHeading{12.3.1}{12.3.1. Proposition.}
\begin{itshape}
Let $A$ be a von Neumann algebra, let $f,f'$ be two normal positive
functionals on $A$, let $e,e'$ be their supports, and put $g=f-f'$.

\begin{enumerate}[label=\textup{(\roman*)}]
 \item The equality
 \[
   \lVert g\rVert=\lVert f\rVert+\lVert f'\rVert
 \]
 holds if and only if $ee'=0$.
 \item Suppose this condition holds. Then
 \[
   \lvert g\rvert=f+f',
 \]
 the support of $\lvert g\rvert$ is $e+e'$, and
 $g=(e-e')\mathbin{\cdot}\lvert g\rvert$ is the polar decomposition of
 $g$; moreover,
 \[
   f=e\mathbin{\cdot}g,\qquad
   f'=-e'\mathbin{\cdot}g.
 \]
\end{enumerate}
\end{itshape}

Suppose that $ee'=0$. In order that a projection $d$ of $A$ satisfy
$(f+f')(d)=0$, it is necessary and sufficient that
$f(d)=f'(d)=0$, that is, that $de=de'=0$. Thus $f+f'$ has support
$e+e'$. On the other hand,
% Source PDF physical page 253
$e-e'$ is a Hermitian partial isometry with initial and final projection
$e+e'$. For every $x\in A$, we have
\begin{align*}
 ((e-e')\mathbin{\cdot}(f+f'))(x)
 &=(f+f')(ex-e'x)\\
 &=f(ex)-f'(e'x)=f(x)-f'(x)=g(x),\\
 ((e-e')^{\ast}\mathbin{\cdot}(f-f'))(x)
 &=(f-f')(ex-e'x)\\
 &=f(ex)+f'(e'x)=f(x)+f'(x).
\end{align*}
By \XRef{12.2.4}{12.2.4}(ii), this proves that
$\lvert g\rvert=f+f'$ and that
$g=(e-e')\mathbin{\cdot}\lvert g\rvert$ is the polar decomposition of
$g$. We have
\[
 e\mathbin{\cdot}g=e\mathbin{\cdot}f=f,
 \qquad
 e'\mathbin{\cdot}g=-e'\mathbin{\cdot}f'=-f'.
\]
Finally,
\[
 \lVert g\rVert=\lVert f+f'\rVert
 =\lVert f\rVert+\lVert f'\rVert.
\]

Conversely, suppose that
$\lVert g\rVert=\lVert f\rVert+\lVert f'\rVert$. The set $B$ of
$x\in A$ satisfying $x=x^{\ast}$ and $\lVert x\rVert\leq1$ is
ultraweakly compact, and $g$ is ultraweakly continuous. There is therefore
$x\in B$ such that $g(x)=\lVert g\rVert$. It follows that
\[
 f(x^{+})+f'(x^{-})-f(x^{-})-f'(x^{+})
 =g(x)=\lVert g\rVert=\lVert f\rVert+\lVert f'\rVert.
\]
But
\[
 f(x^{+})\leq\lVert f\rVert,\qquad
 f'(x^{-})\leq\lVert f'\rVert,\qquad
 -f(x^{-})-f'(x^{+})\leq0,
\]
and hence
\[
 f(x^{+})=\lVert f\rVert,\qquad
 f'(x^{-})=\lVert f'\rVert.
\]
Let $s,s'$ be the supports of $x^{+},x^{-}$. Then
\[
 \lVert f\rVert=f(x^{+})\leq f(s)\leq\lVert f\rVert,
\]
so $f(1-s)=0$ and $s\geq e$; similarly, $s'\geq e'$. Since $ss'=0$, we
have $ee'=0$.

\ResultHeading{12.3.2}{12.3.2. Definition.}
\begin{itshape}
If the conditions of \XRef{12.3.1}{12.3.1}(i) are satisfied, $f$ and
$f'$ are said to be \emph{mutually singular}.
\end{itshape}

\ResultHeading{12.3.3}{12.3.3. Theorem.}
\begin{itshape}
Let $A$ be a von Neumann algebra, and let $g$ be an ultraweakly continuous
Hermitian linear functional on $A$. There exists one and only one pair
$(f,f')$ of mutually singular normal positive functionals on $A$ such
that $g=f-f'$.
\end{itshape}

Let $P$ be the predual of $A$, let $A'$ be its dual, and let $A''$ be its
bidual, identified with its enveloping von Neumann algebra. By
\XRef{2.6.4}{2.6.4}, there are positive functionals $f,f'$ on $A$,
identifiable with normal positive functionals on $A''$, such that
\[
 g=f-f',\qquad
 \lVert g\rVert=\lVert f\rVert+\lVert f'\rVert.
\]
Applying \XRef{12.3.1}{12.3.1} to $A''$, we obtain
$\lvert g\rvert=f+f'$. But $\lvert g\rvert\in P$
(\XRef{12.2.8}{12.2.8}). Hence $f,f'\in P$, so that $f$ and $f'$ are
normal.

Suppose that $g=f_1-f'_1$, with $f_1,f'_1$ mutually singular normal
positive functionals. Let $e,e',e_1,e'_1$ be the supports of
$f,f',f_1,f'_1$. By \XRef{12.3.1}{12.3.1}(ii),
\[
 f_1+f'_1=\lvert g\rvert=f+f',
 \qquad\text{whence}\qquad f=f_1,\quad f'=f'_1.
\]

% Source PDF physical page 254
\ResultHeading{12.3.4}{12.3.4. Corollary.}
\begin{itshape}
Let $A$ be a $C^{*}$-algebra and let $g$ be a continuous Hermitian linear
functional on $A$. There exists one and only one pair $(f,f')$ of
positive functionals on $A$ such that
\[
 g=f-f',\qquad
 \lVert g\rVert=\lVert f\rVert+\lVert f'\rVert.
\]
\end{itshape}

Existence was already known. Uniqueness follows from
\XRef{12.3.3}{12.3.3} and \XRef{12.1.3}{12.1.3}(iii).

\ResultHeading{12.3.5}{12.3.5.}
The functionals $f,f'$ of \XRef{12.3.4}{12.3.4} are called the positive
and negative parts of $g$, and are denoted by $g^{+}$ and $g^{-}$. By
\XRef{12.3.1}{12.3.1},
\[
 \lvert g\rvert=g^{+}+g^{-}.
\]

\ResultHeading{12.3.6}{12.3.6.}
Let $A$ be a $C^{*}$-algebra, let $E$ be the enveloping von Neumann algebra
of $A$, and let $f$ be a positive functional on $A$. Then $f$, regarded
as a normal positive functional on $E$, has a support $e$, which is a
projection of $E$. This is called the \emph{enveloping support} of $f$.
[If $A$ is a von Neumann algebra and $f$ is normal, then $f$ has a support
and an enveloping support which are in general distinct
(\XRef{12.5.7}{12.5.7}).] If $f,f'$ are two positive functionals on the
$C^{*}$-algebra $A$, the equality
\[
 \lVert f-f'\rVert=\lVert f\rVert+\lVert f'\rVert
\]
is equivalent to the enveloping supports of $f$ and $f'$ being orthogonal
(\XRef{12.3.1}{12.3.1}(i)). One then says that $f$ and $f'$ are mutually
singular, which generalizes Definition \XRef{12.3.2}{12.3.2}.

Bibliography: \cite{ref55,ref70,ref131,ref150,ref151,ref152}.

\BookSubsection[Positive Part of an Ideal in a C*-Algebra]{12.4}%
{Positive Part of an Ideal in a $C^{*}$-Algebra}

\ResultHeading{12.4.1}{12.4.1. Proposition.}
\begin{itshape}
Let $A$ be a $C^{*}$-algebra, let $\mathcal I$ be the set of closed left
ideals of $A$, and let $\mathcal J$ be the set of closed subsets $J$ of
$A^{+}$ having the following properties:

\begin{enumerate}[label=\textup{(\roman*)}]
 \item $x,y\in J\Longrightarrow x+y\in J$;
 \item $x\in J$, $y\in A^{+}$, and $y\leq x\Longrightarrow y\in J$.
\end{enumerate}
Then the map $I\mapsto I^{+}=I\cap A^{+}$ is a bijection from
$\mathcal I$ onto $\mathcal J$.
\end{itshape}

First observe that, if $I\in\mathcal I$ and $x\in I^{+}$, then
$x^{1/2}\in I^{+}$, for $I\cap I^{\ast}$ is a $C^{*}$-subalgebra of $A$
and $I^{+}\subset I\cap I^{\ast}$.

Next observe that, if $J\in\mathcal J$ and $x\in J$, then
$x^{1/2}\in J$. Indeed, there is a sequence of continuous functions
$f_n\colon\operatorname{Sp}x\to\mathbb R$, with $f_n\geq0$, converging
uniformly to $t\mapsto\sqrt t$ and vanishing in a neighborhood of $0$.
There is a sequence of constants $\lambda_n\geq0$ such that
$f_n(t)\leq\lambda_nt$ on $\operatorname{Sp}x$. Thus
\[
 0\leq f_n(x)\leq\lambda_nx,\qquad f_n(x)\in J,\qquad x^{1/2}\in J.
\]

We interrupt the proof to establish the following lemma.

\ResultHeading{12.4.2}{12.4.2. Lemma.}
\begin{itshape}
Regard $A$ as extended inside its enveloping von Neumann algebra $E$. Let
$I\in\mathcal I$, and let $EI^{+}$ be the set of products $ax$, where
$a\in E$ and $x\in I^{+}$. Then
\[
I=(EI^{+})\cap A.
\]
\end{itshape}

% Source PDF physical page 255
Let $y\in I$. Then $y=u\lvert y\rvert$ for some $u\in E$. Moreover,
$\lvert y\rvert=(y^{\ast}y)^{1/2}\in I^{+}$ by an earlier observation.
Thus $y\in EI^{+}$, and hence $I\subset(EI^{+})\cap A$. Conversely, let
$a\in E$, $x\in I^{+}$, and suppose that $ax\in A$. We must prove that
$ax\in I$. Now $a$ is in the closure of $A$ for $\sigma(E,A')$, and so
$ax$ is in the closure of $Ax\subset I$ for $\sigma(A,A')$. Since $I$ is
a norm-closed vector subspace of $A$, it is closed for $\sigma(A,A')$.
Thus $ax\in I$.

\ResultHeading{12.4.3}{12.4.3.}
We now return to the proof of \XRef{12.4.1}{12.4.1}. Let
$I\in\mathcal I$. If $x,y\in I^{+}$, then $x+y\in I^{+}$. Let
$z\in A^{+}$ satisfy $z\leq x$. There exists $s\in E$ such that
$z^{1/2}=sx^{1/2}$ (\XRef{A.8}{A 8}); since $x^{1/2}\in I^{+}$ by an
earlier observation,
\[
 z=(sx^{1/2})^{\ast}(sx^{1/2})
  =(x^{1/2}s^{\ast}s)x^{1/2}\in EI^{+},
\]
and therefore $z\in I^{+}$ by \XRef{12.4.2}{12.4.2}. Thus
$I^{+}\in\mathcal J$. The map $I\mapsto I^{+}$ from $\mathcal I$ to
$\mathcal J$ is injective by \XRef{12.4.2}{12.4.2}.

We prove that it is surjective. Let $J\in\mathcal J$, and put
$I=(EJ)\cap A$. Plainly $I^{+}\supset J$. Let $a\in E$ and $x\in J$ be
such that $ax\in A^{+}$, and let us prove that $ax\in J$. We have
\[
 (ax)^2=xa^{\ast}ax\leq\lVert a\rVert^2x^2
 \leq\lVert a\rVert^2\lVert x\rVert x.
\]
Hence $(ax)^2\in J$, and therefore $ax\in J$ by an earlier observation.
Thus $I^{+}=J$.

We shall prove that $I$ is a left ideal of $A$. Clearly $AI\subset I$.
Let $x,y\in I$. Then $x^{\ast}x,y^{\ast}y\in I^{+}=J$, and hence
$x^{\ast}x+y^{\ast}y\in J$. Since
\[
 (x+y)^{\ast}(x+y)\leq2(x^{\ast}x+y^{\ast}y),
\]
it follows that
\[
 \lvert x+y\rvert^2=(x+y)^{\ast}(x+y)\in J,
 \qquad\text{and hence}\qquad \lvert x+y\rvert\in J.
\]
But $x+y=u\lvert x+y\rvert$ for some $u\in E$, so
$x+y\in(EJ)\cap A=I$. Finally, we show that $I$ is closed. If $(x_n)$ is
a sequence of points of $I$ converging to $x\in A$, then
$x_n^{\ast}x_n\in I^{+}=J$ converges to $x^{\ast}x$. Thus
$\lvert x\rvert^2=x^{\ast}x\in J$, $\lvert x\rvert\in J$, and
$x\in(EJ)\cap A=I$.

\ResultHeading{12.4.4}{12.4.4.}
Proposition \XRef{12.4.1}{12.4.1} remains valid if left ideals are
replaced by right ideals. Moreover, if $I$ is a left ideal of $A$, then
$I^{\ast}$ is a right ideal and
\[
 I^{+}=(I^{\ast})^{+}.
\]

Bibliography: \cite{ref21}.

\BookSubsection{12.5}{Complements}

\ResultHeading{12.5.1}{12.5.1.}
Let $A$ be a $C^{*}$-algebra. In the dual $A'$ of $A$, let $(f_\alpha)$
be a family indexed by a directed set and converging weakly to $f$. If
$\lVert f_\alpha\rVert\to\lVert f\rVert$, then
$\lvert f_\alpha\rvert$ converges weakly to $\lvert f\rvert$.
\cite{ref21}

\ResultHeading{12.5.2}{12.5.2.}
Let $H$ be a Hilbert space, let $A$ be a von Neumann algebra on $H$, and
let $f$ be an ultraweakly continuous linear functional on $A$. There is an
ultraweakly continuous linear functional $g$ on $\mathcal L(H)$ extending
$f$ and such that
% Source PDF physical page 256
\[
 \lVert f\rVert=\lVert g\rVert.
\]
[If $f\geq0$, there are vectors $\xi_i\in H$ such that
$\sum_i\lVert\xi_i\rVert^2<+\infty$ and
\[
 f(x)=\sum_i(x\xi_i\mid\xi_i)\qquad(x\in A).
\]
Put
\[
 g(x)=\sum_i(x\xi_i\mid\xi_i)\qquad(x\in\mathcal L(H));
\]
then $\lVert g\rVert=g(1)=f(1)=\lVert f\rVert$. In the general case, use
the polar decomposition of $f$.] \cite{ref128,ref131}

\ResultHeading{12.5.3}{12.5.3.}
Let $A$ be a $C^{*}$-algebra, let $E$ be the enveloping von Neumann
algebra of $A$, and let $f,g$ be two pure states of $A$. In order that
$\pi_f\simeq\pi_g$, it is necessary and sufficient that the supports of
$f$ and $g$, regarded as normal positive functionals on $E$, be equivalent
relative to $E$. \cite{ref168}

\ResultHeading{12.5.4}{*12.5.4.}
Let $A$ be a $C^{*}$-algebra. In order that the dual of $A$ be separable
in norm, it is necessary and sufficient that $A$ admit a countable
composition series $(I_\rho)$ such that the quotients $I_{\rho+1}/I_\rho$
are separable dual $C^{*}$-algebras. \cite{ref168} (See also
\emph{Mathematical Reviews}, 1963, p.~808.)

\ResultHeading{12.5.5}{12.5.5.}
In order that a factor, regarded as a Banach space, be the bidual of a
$C^{*}$-algebra, it is necessary and sufficient that it be of type I.
\cite{ref55,ref152}

\ResultHeading{12.5.6}{12.5.6.}
Let $A$ be a von Neumann algebra, let $A''$ be its enveloping von Neumann
algebra, let $P$ be the predual of $A$, and let $f\in P$ be identified
with a linear functional on $A''$. If $x\in A''$, then
\[
 x\mathbin{\cdot}f\in P,\qquad f\mathbin{\cdot}x\in P.
\]
[Let $\rho$ be the canonical projection from $A''$ onto $A$. If $y\in A$,
then, since $f$ vanishes on $\ker\rho$,
\[
 f(xy)=f(\rho(xy))=f(\rho(x)y),
\]
and therefore $x\mathbin{\cdot}f=\rho(x)\mathbin{\cdot}f\in P$.]
(Grothendieck, unpublished.)

\ResultHeading{12.5.7}{12.5.7.}
Let $H$ be a Hilbert space, let $(\xi_1,\xi_2,\ldots)$ be an orthonormal
basis of $H$, and let $A$ be the von Neumann algebra on $H$ generated by
the projections $P_i=P_{\mathbb C\xi_i}$. The positive functional
\[
 \sum_i\lambda_iP_i\longmapsto\sum_i i^{-2}\lambda_i
\]
is normal and has support $1$. Let $g$ be the positive functional
\[
 \sum_i\lambda_iP_i\longmapsto\lim_i\lambda_i,
\]
where a generalized Banach limit is used. Regard $f$ and $g$ as normal
positive functionals on the enveloping algebra $E$ of $A$. Put
\[
 Q_n=P_n+P_{n+1}+\cdots\in A\subset E.
\]
The $Q_n$ have a strong limit $Q$ in $E$, which is a projection of $E$.
We have $g(Q_n)=1$ for every $n$, so $g(Q)=1$ and $Q\ne0$. On the other
hand, $f(Q_n)\to0$, so $f(Q)=0$. Thus the support of $f$ relative to $E$
is not $1$.

The polar decomposition of $f$ therefore differs from its enveloping
polar decomposition.

\ResultHeading{12.5.8}{12.5.8.}
Let $A$ be a von Neumann algebra, let $f,f'$ be two mutually singular
normal positive functionals on $A$, and put $g=f-f'$. Then
\[
 \lVert f\rVert=
 \sup_{\substack{x\in A^{+}\\\lVert x\rVert\leq1}}g(x),
 \qquad
 \lVert f'\rVert=
 \sup_{\substack{x\in A^{+}\\\lVert x\rVert\leq1}}-g(x).
\]
% Source PDF physical page 257
[If $x\in A^{+}$ and $\lVert x\rVert\leq1$, then
$g(x)\leq f(x)\leq\lVert f\rVert$. On the other hand, if $e$ is the
support of $f$, then $g(e)=f(e)=\lVert f\rVert$.]

\ResultHeading{12.5.9}{12.5.9.}
Let $A$ be a von Neumann algebra. Let $\mathcal E$ be the set of weakly
closed left ideals of $A$. Let $\mathcal E'$ be the set of weakly closed
subsets $P$ of $A^{+}$ such that:

\begin{enumerate}[label=\textup{\arabic*\textdegree}]
 \item $x,y\in P\Longrightarrow x+y\in P$;
 \item $x\in P$ and $\lambda\geq0\Longrightarrow\lambda x\in P$;
 \item $x\in P$, $y\in A^{+}$, and $y\leq x\Longrightarrow y\in P$.
\end{enumerate}
Then $L\mapsto L\cap A^{+}$ is a bijection from $\mathcal E$ onto
$\mathcal E'$. \cite{ref21}

\ResultHeading{12.5.10}{12.5.10.}
Let $A$ be a $C^{*}$-algebra. Let $\mathcal E$ be the set of closed
two-sided ideals of $A$. Let $\mathcal E'$ be the set of closed subsets
$P$ of $A^{+}$ such that:

\begin{enumerate}[label=\textup{\arabic*\textdegree}]
 \item $x,y\in P\Longrightarrow x+y\in P$;
 \item $x\in P$ and $\lambda\geq0\Longrightarrow\lambda x\in P$;
 \item $x\in P$ and $y\in A\Longrightarrow y^{\ast}xy\in P$.
\end{enumerate}
Then $I\mapsto I\cap A^{+}$ is a bijection from $\mathcal E$ onto
$\mathcal E'$. \cite{ref21}

\ResultHeading{12.5.11}{12.5.11.}
Let $A$ be a $C^{*}$-algebra, let $B$ be its enveloping von Neumann
algebra, and let $\pi\colon A\to B$ be the canonical homomorphism. Define
$A_1,B_1,\pi_1$ similarly. Let $\varphi\colon A\to A_1$ be a
homomorphism. There exists one and only one homomorphism
$\psi\colon B\to B_1$ such that
\[
 \pi_1\circ\varphi=\psi\circ\pi.
\]
If $B,B_1$ are identified with $A'',A_1''$, then $\psi$ is the
bitranspose of $\varphi$.

% Source PDF physical page 258
\BookDivision{part-two}{Applications to Group Representations}
\BookSectionHere{13}{Unitary Representations of Locally Compact Groups}

\BookSubsection{13.1}{Elementary Definitions Concerning Representations}

\ResultHeading{13.1.1}{13.1.1. Definition.}
\begin{itshape}
Let $G$ be a topological group and $H$ a Hilbert space. A
\emph{continuous unitary representation} of $G$ on $H$ is a homomorphism
from $G$ into the unitary group of $H$ which is continuous for the strong
topology.
\end{itshape}

In other words, a continuous unitary representation of $G$ on $H$ is a
map $\pi$ from $G$ into the unitary group of $H$ such that
\[
 \pi(st)=\pi(s)\pi(t)
\]
and such that, for every $\xi\in H$, the map $s\mapsto\pi(s)\xi$ from
$G$ into $H$ is continuous for the strong topology of $H$. The identity
$\pi(st)=\pi(s)\pi(t)$ implies
\[
 \pi(e)=1,
 \qquad
 \pi(s^{-1})=\pi(s)^{-1}=\pi(s)^{\ast},
\]
where $e$ will always denote the identity element of the group under
consideration.

The functions
\[
 s\longmapsto\varphi_{\xi,\eta}(s)
  =(\pi(s)\xi\mid\eta)
 \qquad(s\in G;\ \xi,\eta\in H\text{ fixed})
\]
are called the \emph{coefficients} of $\pi$. We have
\[
 \varphi_{\xi,\eta}(s)
 =\overline{\varphi_{\eta,\xi}(s^{-1})}.
\]

\ResultHeading{13.1.2}{13.1.2.}
On the unitary group of $H$, the strong and weak topologies coincide. Thus
throughout the preceding discussion the strong topology may be replaced
by the weak topology. If, on the other hand, continuity for the norm
topology were imposed, a much more restrictive definition, of no interest
for what follows, would result. Of course, if
$\dim H<+\infty$, all these notions of continuity coincide.

\ResultHeading{13.1.3}{13.1.3.}
The Hilbert dimension of $H$ is called the dimension of $\pi$ and is
denoted by $\dim\pi$. The space $H$ is called the space of $\pi$ and is
denoted by $H_\pi$.

By analogy with the case of involutive algebras, the following notions
are also defined without difficulty: equivalent representations, class of
representations, intertwining operator, intertwining number, Hilbert
direct sum of representations, multiple of a representation,
subrepresentation, representation contained in a representation, and
cyclic vector. [The vector $\xi$ is called cyclic for $\pi$ if
$\pi(G)\xi$ is total in $H_\pi$.] We use the same notation as in the case
of algebras, and these notions have
% Source PDF physical page 259
the same elementary properties; we leave their verification to the
reader.

There is no need here to define the essential subspace of a
representation. Everything proceeds as if continuous unitary
representations were automatically nondegenerate.

\ResultHeading{13.1.4}{13.1.4.}
Let $\pi$ be a continuous unitary representation of $G$. The following
conditions are equivalent:

\begin{enumerate}[label=\textup{(\roman*)}]
 \item the only closed vector subspaces of $H_\pi$ invariant under
 $\pi(G)$ are $0$ and $H_\pi$;
 \item the commutant of $\pi(G)$ in $\mathcal L(H_\pi)$ consists only of
 the scalars;
 \item every nonzero vector of $H_\pi$ is cyclic for $\pi$.
\end{enumerate}
If these conditions hold and, moreover, $H_\pi\ne0$, then $\pi$ is said
to be \emph{topologically irreducible}, or simply \emph{irreducible}.
[We shall never consider the notion of algebraic irreducibility, except
when $\dim\pi<+\infty$, in which case it is equivalent to topological
irreducibility.] We denote by $\widehat G$ the set of classes of
irreducible representations of $G$.

Exactly as in \SectionRef{5}{Section 5}, we define disjoint, factor,
quasi-equivalent, type I, type II, and so on, multiplicity-free, and
multiplicity-$n$ representations. All the arguments and all the results of
\SectionRef{5}{Section 5} extend immediately to group representations.
[One may also use the correspondence between representations of $G$ and
representations of $L^1(G)$ which will be established in
\SectionRef{13.3}{13.3}, at least when $G$ is locally compact.]

\ResultHeading{13.1.5}{13.1.5.}
We shall now define for group representations operations which have no
meaning for representations of involutive algebras.

Let $\pi,\pi'$ be continuous unitary representations of $G$. For every
$s\in G$, form, on the Hilbert tensor product $H_\pi\otimes H_{\pi'}$, the
unitary operator $\pi(s)\otimes\pi'(s)$. It is immediate that
\[
 s\longmapsto\pi(s)\otimes\pi'(s)
\]
is a continuous unitary representation of $G$ on
$H_\pi\otimes H_{\pi'}$, called the \emph{tensor product} of $\pi$ and
$\pi'$, and denoted by $\pi\otimes\pi'$.

Let $\pi$ be a continuous unitary representation of $G$ on $H$, and let
$\overline H$ be the conjugate Hilbert space of $H$. Each $\overline{\pi(s)}$
is a unitary operator on $\overline H$, and it is immediate that
$s\mapsto\overline{\pi(s)}$ is also a continuous unitary representation
of $G$ on $\overline H$, called the \emph{conjugate representation} of
$\pi$ and denoted by $\overline\pi$.

\ResultHeading{13.1.6}{13.1.6.}
Let $G$ be a locally compact group equipped with a left Haar measure
$ds$. For every $s\in G$, let $\lambda(s)$ be the operator on $L^2(G)$
defined by
\[
 (\lambda(s)f)(x)=f(s^{-1}x)
 \qquad(f\in L^2(G),\ x\in G).
\]
One checks immediately that $\lambda$ is a continuous unitary
representation of $G$ on $L^2(G)$, called the \emph{left regular
representation}.
% Source PDF physical page 260
Let $\Delta\colon G\to\mathbb R$ be the modular function of $G$. For every
$s\in G$, let $\rho(s)$ be the operator on $L^2(G)$ defined by
\[
 (\rho(s)f)(x)=\Delta(s)^{1/2}f(xs)
 \qquad(f\in L^2(G),\ x\in G).
\]
One checks immediately that $\rho$ is a continuous unitary representation
of $G$ on $L^2(G)$, called the \emph{right regular representation}.

The representations $\lambda$ and $\rho$ are injective.

For every $f\in L^2(G)$, define $f'\in L^2(G)$ by
\[
 f'(x)=\Delta(x)^{-1/2}f(x^{-1}).
\]
Then $f\mapsto f'$ is an isomorphism of the Hilbert space $L^2(G)$ onto
itself, and, for every $s\in G$,
\begin{align*}
 (\lambda(s)f)'(x)
 &=\Delta(x)^{-1/2}f(s^{-1}x^{-1})\\
 &=\Delta(s)^{1/2}\Delta(xs)^{-1/2}f((xs)^{-1})
  =(\rho(s)f')(x).
\end{align*}
Thus the isomorphism $f\mapsto f'$ transforms $\lambda$ into $\rho$, so
$\lambda\simeq\rho$. One therefore sometimes speaks simply of the
regular representation of $G$.

For every $f\in L^2(G)$, let $\overline f\in\overline{L^2(G)}$ be the
complex conjugate function. Then $f\mapsto\overline f$ is an isomorphism
from the Hilbert space $L^2(G)$ onto its conjugate Hilbert space which
transforms $\lambda$ into $\overline\lambda$. Thus
\[
 \lambda\simeq\overline\lambda,
 \qquad
 \rho\simeq\overline\rho.
\]

\ResultHeading{13.1.7}{13.1.7.}
Let $a$ be a cardinal and let $H$ be a Hilbert space of Hilbert dimension
$a$. The representation $s\mapsto1$ of $G$ on $H$ is called the
\emph{trivial representation} of $G$ on $H$, or the trivial
representation of $G$ of dimension $a$. It is denoted by $1_H$, or simply
by $1$ when there is no ambiguity concerning $H$.

\ResultHeading{13.1.8}{13.1.8. Proposition.}
\begin{itshape}
Let $G_1,G_2$ be topological groups, let $G=G_1\times G_2$, and let
$\sigma$ be an irreducible continuous unitary representation of $G$.
Suppose that every continuous unitary factor representation of $G_1$ is
of type I. Then, for $i=1,2$, there is an irreducible continuous unitary
representation $\sigma_i$ of $G_i$ such that $\sigma$ is equivalent to the
representation
\[
 (s_1,s_2)\longmapsto\sigma_1(s_1)\otimes\sigma_2(s_2)
\]
on the Hilbert space $H_{\sigma_1}\otimes H_{\sigma_2}$.
\end{itshape}

Let $\mathfrak A_i$ be the von Neumann algebra generated by
$\sigma(G_i)$. Then $\mathfrak A_1$ and $\mathfrak A_2$ commute and
together generate the von Neumann algebra $\mathcal L(H_\sigma)$. Every
element of the center of $\mathfrak A_1$ commutes with $\mathfrak A_1$
and $\mathfrak A_2$, and is therefore a scalar. Hence $\mathfrak A_1$ is
a factor, of type I by the hypothesis on $G_1$. There are Hilbert spaces
$H_1,H_2$ such that $H_\sigma$ is identified with $H_1\otimes H_2$ and
\[
 \mathfrak A_1=\mathcal L(H_1)\otimes\mathbb C
 \qquad\text{(\XRef{A.36}{A 36})}.
\]
For every $s_1\in G_1$, the operator $\sigma(s_1)$ has the form
$\sigma_1(s_1)\otimes1$, where $\sigma_1(s_1)$ is a unitary operator on
$H_1$, and $\sigma_1$ is a continuous unitary representation of $G_1$ on
$H_1$. The operators $\sigma_1(s_1)$ generate the von Neumann algebra
$\mathcal L(H_1)$, so $\sigma_1$ is irreducible. We have
\[
 \mathfrak A_2\subset\mathbb C\otimes\mathcal L(H_2)
 \qquad\text{(\XRef{A.18}{A 18})},
\]
and therefore, for every $s_2\in G_2$, the operator $\sigma(s_2)$ has the
form $1\otimes\sigma_2(s_2)$, where $\sigma_2$ is a continuous unitary
representation of $G_2$ on $H_2$. If $T\in\mathcal L(H_2)$
% Source PDF physical page 261
commutes with the $\sigma_2(s_2)$, then $1\otimes T$ commutes with
$\mathfrak A_1$ and $\mathfrak A_2$, and is therefore a scalar. Thus $T$
is a scalar and $\sigma_2$ is irreducible.

Bibliography: \cite{ref89,ref108}.

\BookSubsection[The Involutive Algebra L1(G)]{13.2}%
{The Involutive Algebra $L^1(G)$}

\ResultHeading{13.2.1}{13.2.1.}
Let $G$ be a locally compact group, and let $M^1(G)$ be the algebra, for
the convolution product, of bounded complex measures on $G$. It is a
Banach algebra whose identity is the Dirac measure $\varepsilon_e$ at the
point $e$. If, for every $\mu\in M^1(G)$, $\mu^{\ast}$ is defined by
\[
 d\mu^{\ast}(s)=d\overline\mu(s^{-1}),
\]
then $M^1(G)$ becomes an involutive algebra and
$\lVert\mu^{\ast}\rVert=\lVert\mu\rVert$. Thus $M^1(G)$ is an
involutive Banach algebra.

\ResultHeading{13.2.2}{13.2.2.}
Fix once and for all a left Haar measure $ds$ on $G$. Associating with
each $f\in L^1(G)$ the measure
\[
 d\mu(s)=f(s)\,ds\in M^1(G)
\]
gives an isometric homomorphism $\Phi$ from the Banach algebra $L^1(G)$
into the Banach algebra $M^1(G)$. For $g\in\mathcal K(G)$, the set of
continuous compactly supported complex-valued functions on $G$, and with
$\Delta$ denoting the modular function of $G$, we have
\begin{align*}
 \int g(s)\,d\mu^{\ast}(s)
 &=\overline{\int\overline{g(s^{-1})}\,d\mu(s)}
  =\int g(s^{-1})\overline{f(s)}\,ds\\
 &=\int g(s)\overline{f(s^{-1})}\Delta(s^{-1})\,ds.
\end{align*}
For every complex-valued function $f$ on $G$, define $f^{\ast}$ by
\[
 f^{\ast}(s)=\overline{f(s^{-1})}\Delta(s^{-1}).
\]
The preceding calculation then shows that $d\mu^{\ast}(s)=f^{\ast}(s)ds$.
Thus $f\mapsto f^{\ast}$ is an isometric involution on $L^1(G)$, and
$\Phi$ is a homomorphism of the involutive algebra $L^1(G)$ into the
involutive algebra $M^1(G)$. We identify $L^1(G)$ with $\Phi(L^1(G))$.
In general, $L^1(G)$ is not a $C^{*}$-algebra.

\ResultHeading{13.2.3}{13.2.3.}
For every complex-valued function $f$ on $G$, define $\check f$ and
$\widetilde f$ by
\[
 \check f(s)=f(s^{-1}),
 \qquad
 \widetilde f(s)=\overline{f(s^{-1})}.
\]
Thus $f^{\ast}=\widetilde f\,\Delta^{-1}$. If $a\in G$, put
\[
 {}_af(s)=f(as),\qquad f_a(s)=f(sa).
\]
Writing $\varepsilon_a$ for the Dirac measure at $a$, we have
\[
 \varepsilon_a\ast f={}_{a^{-1}}f,
 \qquad
 f\ast\varepsilon_a=\Delta(a^{-1})f_{a^{-1}}.
\]
The formula
\[
 (\varepsilon_a\ast f)^{\ast}
 =f^{\ast}\ast\varepsilon_a^{\ast}
 =f^{\ast}\ast\varepsilon_{a^{-1}}
\]
therefore gives
\[
 ({}_{a^{-1}}f)^{\ast}=\Delta(a)(f^{\ast})_a.
\]

\ResultHeading{13.2.4}{13.2.4.}
If $G$ is discrete, $L^1(G)$ has an identity element. If $G$ is separable,
$L^1(G)$ is separable.

% Source PDF physical page 262
\ResultHeading{13.2.5}{13.2.5.}
Let $s\in G$. Let $I$ be the set of neighborhoods of $s$, ordered by the
relation opposite to inclusion. For each $i\in I$, let $u_i$ be a positive
function in $L^1(G)$, vanishing outside $i$ and having integral $1$. Then,
for every $f\in L^1(G)$,
\[
 \lVert u_i\ast f-\varepsilon_s\ast f\rVert_1\longrightarrow0,
 \qquad
 \lVert f\ast u_i-f\ast\varepsilon_s\rVert_1\longrightarrow0.
\]
[It is enough to verify this for $f\in\mathcal K(G)$, when it follows
immediately from uniform continuity.] Applying this with $s=e$ shows that
$L^1(G)$ has an approximate identity.

Bibliography: \cite{ref86,ref108}.

\BookSubsection[Representations of G and Representations of L1(G)]{13.3}%
{Representations of $G$ and Representations of $L^1(G)$}

\ResultHeading{13.3.1}{13.3.1. Proposition.}
\begin{itshape}
Let $\pi$ be a continuous unitary representation of $G$. For every
$\mu\in M^1(G)$, put
\[
 \pi(\mu)=\int\pi(s)\,d\mu(s)\in\mathcal L(H_\pi).
\]
Then $\mu\mapsto\pi(\mu)$ is a representation of the involutive algebra
$M^1(G)$ on $H_\pi$, whose restriction to $L^1(G)$ is nondegenerate.
\end{itshape}

The assertion that $\mu\mapsto\pi(\mu)$ is a representation of the
involutive algebra $M^1(G)$ follows from easy calculations. For example,
if $\mu\in M^1(G)$ and $\xi,\eta\in H_\pi$, then
\begin{align*}
 (\eta\mid\pi(\mu^{\ast})\xi)
 &=\overline{(\pi(\mu^{\ast})\xi\mid\eta)}
  =\overline{\int(\pi(s^{-1})\xi\mid\eta)\,d\overline\mu(s)}\\
 &=\int(\pi(s)\eta\mid\xi)\,d\mu(s)
  =(\pi(\mu)\eta\mid\xi)
  =(\eta\mid\pi(\mu)^{\ast}\xi),
\end{align*}
so $\pi(\mu^{\ast})=\pi(\mu)^{\ast}$. On the other hand, let $s\in G$,
$\xi\in H_\pi$, and $\varepsilon>0$. With the notation of
\XRef{13.2.5}{13.2.5}, there is $i\in I$ such that
\[
 \lVert\pi(s')\xi-\pi(s)\xi\rVert\leq\varepsilon
 \qquad(s'\in i),
\]
whence $\lVert\pi(u_i)\xi-\pi(s)\xi\rVert\leq\varepsilon$. Thus
\begin{equation}
 \pi(u_i)\longrightarrow\pi(s)
 \label{eq:13.3.1-1}
\end{equation}
for the strong operator topology. Applying this with $s=e$, we see that
the representation of $L^1(G)$ is nondegenerate.

\ResultHeading{13.3.2}{13.3.2.}
The representations obtained from $M^1(G)$ and $L^1(G)$ are said to be
associated with the given representation of $G$.

\ResultHeading{13.3.3}{13.3.3.}
For $s\in G$, we have $\pi(\varepsilon_s)=\pi(s)$. Thus, if
$f\in L^1(G)$,
\[
 \pi({}_af)=\pi(\varepsilon_{a^{-1}}\ast f)
 =\pi(a^{-1})\pi(f),
\]
and
\[
 \pi(f_a)
 =\pi(\Delta(a^{-1})f\ast\varepsilon_{a^{-1}})
 =\Delta(a^{-1})\pi(f)\pi(a^{-1}).
\]

\ResultHeading{13.3.4}{13.3.4. Proposition.}
\begin{itshape}
Let $H$ be a Hilbert space and let $\pi$ be a nondegenerate representation
of the involutive algebra $L^1(G)$ on $H$. Then $\pi$ is associated with
one and only one continuous unitary representation of $G$.
\end{itshape}

% Source PDF physical page 263
Uniqueness follows from formula \eqref{eq:13.3.1-1}. We prove existence.
Let $H'$ be the vector subspace of $H$ spanned by the $\pi(f)(H)$. By
hypothesis, $H'$ is dense in $H$. Let $s\in G$. With the notation of
\XRef{13.2.5}{13.2.5},
\[
 \lVert u_i\ast f-\varepsilon_s\ast f\rVert_1\longrightarrow0,
 \qquad
 \lVert\pi(u_i)\pi(f)-\pi(\varepsilon_s\ast f)\rVert\longrightarrow0.
\]
Thus $\pi(u_i)$ converges, for the topology of pointwise convergence on
$H'$, to an operator $\pi(s)$ on $H'$ satisfying
\begin{equation}
 \pi(s)\pi(f)=\pi(\varepsilon_s\ast f).
 \tag{1}\label{eq:13.3.4-1}
\end{equation}
Since $\lVert\pi(u_i)\rVert\leq\lVert u_i\rVert=1$, $\pi(s)$ extends
uniquely to a continuous linear operator on $H$, still denoted by
$\pi(s)$, such that $\lVert\pi(s)\rVert\leq1$. For $s,t\in G$ and
$f\in L^1(G)$, formula \eqref{eq:13.3.4-1} gives
\begin{align*}
 \pi(st)\pi(f)
 &=\pi(\varepsilon_{st}\ast f)
  =\pi(\varepsilon_s\ast\varepsilon_t\ast f)\\
 &=\pi(s)\pi(\varepsilon_t\ast f)
  =\pi(s)\pi(t)\pi(f).
\end{align*}
Hence $\pi(st)=\pi(s)\pi(t)$ on $H'$, and therefore on $H$. Formula
\eqref{eq:13.3.4-1} shows that $\pi(e)=1$ and that $\pi(s)\xi$ depends
continuously on $s$ for $\xi\in H'$, and hence for $\xi\in H$. Finally,
since $\pi(s)$ and $\pi(s^{-1})=\pi(s)^{-1}$ decrease norms, $\pi(s)$ is
unitary. This proves that $\pi$ is a continuous unitary representation of
$G$.

We now show that the associated representation of $L^1(G)$ is the
original representation. Let $f,g\in L^1(G)$. Since
\[
 f\ast g=\int f(s)(\varepsilon_s\ast g)\,ds
\]
in $L^1(G)$, we have
\begin{align*}
 \pi(f)\pi(g)=\pi(f\ast g)
 &=\int f(s)\pi(\varepsilon_s\ast g)\,ds\\
 &=\int f(s)\pi(s)\pi(g)\,ds
  =\left(\int f(s)\pi(s)\,ds\right)\pi(g).
\end{align*}
Thus
\[
 \pi(f)=\int f(s)\pi(s)\,ds.
\]

\ResultHeading{13.3.5}{13.3.5.}
Propositions \XRef{13.3.1}{13.3.1} and \XRef{13.3.4}{13.3.4} establish
a bijective correspondence between continuous unitary representations of
$G$ and nondegenerate representations of $L^1(G)$. Let $\pi$ be a
continuous unitary representation of $G$, and let $\pi'$ be the associated
representation of $L^1(G)$. The closed vector subspaces invariant under
$\pi(G)$ and $\pi'(L^1(G))$ are the same, and the cyclic vectors for
$\pi(G)$ and $\pi'(L^1(G))$ are the same, in particular by formula
\eqref{eq:13.3.1-1}. Thus $\pi(G)$ and $\pi'(L^1(G))$ have the same
commutant and generate the same von Neumann algebra. In particular,
topological irreducibility of $\pi'$ is equivalent to that of $\pi$. We
have therefore established a bijective correspondence between
% Source PDF physical page 264
irreducible continuous unitary representations of $G$ and topologically
irreducible representations of $L^1(G)$. Similarly, saying that $\pi$ is
a factor representation, of type I, and so on, is equivalent to saying
that $\pi'$ is a factor representation, of type I, and so on. If
$\pi_1,\pi_2$ are continuous unitary representations of $G$, and
$\pi'_1,\pi'_2$ the associated representations of $L^1(G)$, then the
intertwining operators of $\pi_1$ and $\pi_2$ are the same as those of
$\pi'_1$ and $\pi'_2$, and
\[
 \pi'_1\oplus\pi'_2=(\pi_1\oplus\pi_2)',
\]
and so forth. Henceforth the same letter will denote a continuous unitary
representation of $G$ and the associated representation of $L^1(G)$.

\ResultHeading{13.3.6}{13.3.6.}
Let $\lambda$ be the left regular representation of $G$ on $L^2(G)$. If
$f\in L^1(G)$, then $\lambda(f)$ is the operator
\[
 g\longmapsto f\ast g
\]
on $L^2(G)$. This representation of $L^1(G)$ is called the left regular
representation of $L^1(G)$ on $L^2(G)$. If $\lambda(f)=0$, then
$f\ast g=0$ for every $g\in\mathcal K(G)$, and hence $f=0$ by
\XRef{13.2.5}{13.2.5}. Thus the left regular representation of $L^1(G)$
on $L^2(G)$ is injective.

If $G$ is commutative and $\widehat G$ denotes the dual group of $G$, the
Plancherel isomorphism from $L^2(G)$ onto $L^2(\widehat G)$ transforms
$\lambda(f)$ into multiplication by $\mathcal Ff$, the Fourier transform
of $f$, on $L^2(\widehat G)$. Hence
\[
 \lVert\lambda(f)\rVert
 =\sup_{t\in\widehat G}\lvert(\mathcal Ff)(t)\rvert.
\]
In general,
\[
 \int\lvert f(s)\rvert\,ds
 \ne\sup\lvert\mathcal Ff\rvert,
\]
so $\lambda$ is not isometric; thus $L^1(G)$ is not a $C^{*}$-algebra.

Bibliography: \cite{ref86,ref108}.

\BookSubsection[Positive Forms on L1(G) and Functions of Positive Type]{13.4}%
{Positive Forms on $L^1(G)$ and Functions of Positive Type}

\ResultHeading{13.4.1}{13.4.1.}
Section \SectionRef{13.3}{13.3} leads us to study the representations of
$L^1(G)$, or, what amounts to the same thing by \SectionRef{2.4}{2.4}, the
continuous positive forms on $L^1(G)$. Now a continuous linear form on
$L^1(G)$ is defined by an element of $L^\infty(G)$. This leads us to
introduce certain bounded functions on $G$.

\ResultLabel{Definition.}
\begin{itshape}
A continuous complex-valued function $\varphi$ on $G$ is said to be of
positive type if, for arbitrary elements $s_1,\ldots,s_n$ of $G$, the
matrix
\[
 \bigl(\varphi(s_i^{-1}s_j)\bigr)_{1\leq i,j\leq n}
\]
is positive Hermitian. In other words, for arbitrary
$s_1,\ldots,s_n\in G$ and $\alpha_1,\ldots,\alpha_n\in\mathbb C$, one
has
\begin{equation}
 \sum_{i,j=1}^{n}\alpha_i\overline{\alpha_j}
   \varphi(s_i^{-1}s_j)\geq0.
 \tag{1}\label{eq:13.4.1-1}
\end{equation}
\end{itshape}

\ResultHeading{13.4.2}{13.4.2.}
The sum of two continuous functions of positive type is of positive
type. If $\varphi$ is continuous and of positive type and $\lambda\geq0$,
then $\lambda\varphi$ is of positive type. For examples of continuous
functions of positive type, see \XRef{13.6.3}{13.6.3}.

% Source PDF physical page 265
\ResultHeading{13.4.3}{13.4.3.}
In formula \eqref{eq:13.4.1-1}, take $n=2$, $s_1=e$, and $s_2=s\in G$.
The matrix
\[
 \begin{pmatrix}
  \varphi(e)&\varphi(s)\\
  \varphi(s^{-1})&\varphi(e)
 \end{pmatrix}
\]
must be positive Hermitian. This implies
\[
 \varphi(s^{-1})=\overline{\varphi(s)},
 \qquad
 \lvert\varphi(s)\rvert\leq\varphi(e)
\]
for every $s\in G$. In particular, $\varphi$ is bounded, and
$\lVert\varphi\rVert_\infty=\varphi(e)$.

\ResultHeading{13.4.4}{13.4.4. Proposition.}
\begin{itshape}
Let $\varphi$ be a continuous complex-valued function on $G$. The
following conditions are equivalent:

\begin{enumerate}[label=(\roman*)]
 \item $\varphi$ is of positive type;
 \item $\varphi$ is bounded and, for every bounded complex measure $\mu$
 on $G$, one has $\langle\varphi,\mu^{*}\ast\mu\rangle\geq0$, in other
 words,
 \[
  \iint \varphi(y^{-1}z)\,d\overline\mu(y)\,d\mu(z)\geq0;
 \]
 \item for every $f\in\mathcal K(G)$, one has
 $\langle\varphi,f^{*}\ast f\rangle\geq0$, in other words,
 \[
  \iint \varphi(y^{-1}z)\overline{f(y)}f(z)\,dy\,dz\geq0.
 \]
\end{enumerate}
\end{itshape}

$(\mathrm{ii})\Rightarrow(\mathrm{i})$.
Suppose condition (ii) holds. Let $s_1,\ldots,s_n\in G$ and
$\alpha_1,\ldots,\alpha_n\in\mathbb C$, and let $\mu$ be the measure on
$G$ having mass $\alpha_i$ at $s_i$. Then
\[
 0\leq
 \iint\varphi(y^{-1}z)\,d\overline\mu(y)\,d\mu(z)
 =\sum_{i,j=1}^{n}\overline{\alpha_i}\alpha_j
   \varphi(s_i^{-1}s_j),
\]
so $\varphi$ is of positive type.

$(\mathrm{i})\Rightarrow(\mathrm{iii})$.
Suppose that $\varphi$ is of positive type. Let $f\in\mathcal K(G)$.
The function
\[
 (y,z)\longmapsto\varphi(y^{-1}z)\overline{f(y)}f(z)
\]
on $G\times G$ is continuous and has compact support $S$; $S$ is
contained in a set $K\times K$, where $K$ is a compact subset of $G$.
The measure induced on $K$ by Haar measure is the vague limit of positive
finitely supported measures $\nu_i$ of bounded norms; hence the measure
induced on $K\times K$ by Haar measure on $G\times G$ is the vague limit
of the measures $\nu_i\otimes\nu_i$. Now, if $\nu_i$ is given by the
masses $\beta_1,\ldots,\beta_n$ at the points $s_1,\ldots,s_n$, then
\[
 \iint\varphi(y^{-1}z)\overline{f(y)}f(z)\,d\nu_i(y)\,d\nu_i(z)
 =\sum_{i,j=1}^{n}\varphi(s_i^{-1}s_j)
   \overline{f(s_i)}f(s_j)\beta_i\beta_j\geq0,
\]
and, on passing to the limit,
\[
 \iint\varphi(y^{-1}z)\overline{f(y)}f(z)\,dy\,dz\geq0.
\]

% Source PDF physical page 266
$(\mathrm{iii})\Rightarrow(\mathrm{ii})$.
Suppose condition (iii) holds. Let $\mu$ be a compactly supported complex
measure on $G$. For every $f\in\mathcal K(G)$, $\mu\ast f$ belongs to
$\mathcal K(G)$, and consequently
\begin{align*}
 0&\leq\langle\varphi,
       f^{*}\ast\mu^{*}\ast\mu\ast f\rangle\\
  &=\iiiint\varphi(xyzt)f^{*}(x)f(t)\,dx\,dt\,
       d\mu^{*}(y)\,d\mu(z)\\
  &=\iint d\mu^{*}(y)\,d\mu(z)
       \iint\varphi(xyzt)f^{*}(x)f(t)\,dx\,dt.
\end{align*}
Taking $f\geq0$, with integral $1$ and support in an increasingly small
neighborhood of $e$, the integral
\[
 \iint\varphi(xyzt)f^{*}(x)f(t)\,dx\,dt
\]
tends to $\varphi(yz)$ uniformly on every compact subset of $G\times G$.
It follows on passing to the limit that
$\langle\varphi,\mu^{*}\ast\mu\rangle\geq0$. Arguing as in
$(\mathrm{ii})\Rightarrow(\mathrm{i})$, we first see that $\varphi$ is of
positive type, and hence bounded. Now let $\mu$ be any bounded measure on
$G$. Then $\mu$ is the norm limit of compactly supported measures $\nu$.
Since $\varphi$ is bounded,
$\langle\varphi,\mu^{*}\ast\mu\rangle$ is the limit of
$\langle\varphi,\nu^{*}\ast\nu\rangle$, and is therefore nonnegative by
what precedes.

\ResultHeading{13.4.5}{13.4.5. Theorem.}
\begin{itshape}
\begin{enumerate}[label=(\roman*)]
 \item Let $\varphi\in L^\infty(G)$, and let $\omega$ be the continuous
 linear form on $L^1(G)$ defined by $\varphi$. Then $\omega$ is positive
 if and only if $\varphi$ is locally almost everywhere equal to a
 continuous function of positive type.
 \item A complex-valued function $\psi$ on $G$ is continuous and of
 positive type if and only if there exist a continuous unitary
 representation $\pi$ of $G$ and a vector $\xi\in H_\pi$ (which may be
 assumed cyclic for $\pi$) such that
 \[
  \psi(s)=(\pi(s)\xi\mid\xi).
 \]
 One then has $\lVert\psi\rVert_\infty=(\xi\mid\xi)$.
 \item Let $\pi$ and $\pi'$ be continuous unitary representations of
 $G$, and let $\xi$ (respectively, $\xi'$) be a cyclic vector for $\pi$
 (respectively, $\pi'$). If
 \[
  (\pi(s)\xi\mid\xi)=(\pi'(s)\xi'\mid\xi')
  \qquad(s\in G),
 \]
 there exists an isomorphism from $H_\pi$ onto $H_{\pi'}$ which
 transforms $\pi$ into $\pi'$ and $\xi$ into $\xi'$.
\end{enumerate}
\end{itshape}

Suppose that $\varphi$ is locally almost everywhere equal to a continuous
function of positive type. Then, for every $f\in L^1(G)$,
\[
 \int\varphi(s)(f^{*}\ast f)(s)\,ds\geq0
\]
by \XRef{13.4.4}{13.4.4}. Thus $\omega$ is a positive form on $L^1(G)$.
Conversely, suppose that $\omega\geq0$. Form the representation
$\pi_\omega$ of $L^1(G)$ and the vector $\xi_\omega$. By
\XRef{13.3.4}{13.3.4}, $\pi_\omega$ is associated with a continuous
unitary representation of $G$ on the space of $\pi_\omega$, which we
again denote by $\pi_\omega$.
% Source PDF physical page 267
Then, for every $f\in L^1(G)$,
\[
 \int\varphi(s)f(s)\,ds
 =\omega(f)
 =(\pi_\omega(f)\xi_\omega\mid\xi_\omega)
 =\int(\pi_\omega(s)\xi_\omega\mid\xi_\omega)f(s)\,ds,
\]
and therefore
\begin{equation}
 \varphi(s)=(\pi_\omega(s)\xi_\omega\mid\xi_\omega)
 \quad\text{locally almost everywhere}.
 \tag{1}\label{eq:13.4.5-1}
\end{equation}
If $\varphi$ is continuous, equality \eqref{eq:13.4.5-1} holds
everywhere, since its right-hand side is a continuous function of $s$.
To complete the proof of (i) and (ii), it is enough to prove that, if
$\pi$ is a continuous unitary representation of $G$ and $\xi\in H_\pi$,
then $s\mapsto(\pi(s)\xi\mid\xi)$ is of positive type. But if
$\mu\in M^1(G)$, then
\[
 0\leq\lVert\pi(\mu)\xi\rVert^2
 =(\pi(\mu^{*}\ast\mu)\xi\mid\xi)
 =\int(\pi(s)\xi\mid\xi)\,d(\mu^{*}\ast\mu)(s),
\]
which proves the assertion by \XRef{13.4.4}{13.4.4}.

Now adopt the hypotheses and notation of (iii). We have
\[
 (\pi(f)\xi\mid\xi)=(\pi'(f)\xi'\mid\xi')
 \qquad\text{for every }f\in L^1(G).
\]
Consequently there exists an isomorphism $\Phi$ from $H_\pi$ onto
$H_{\pi'}$ which transforms $\xi$ into $\xi'$ and $\pi(f)$ into
$\pi'(f)$ for every $f\in L^1(G)$ \XRef{2.4.1}{[2.4.1(ii)]}. By
\XRef{13.3.4}{13.3.4}, $\Phi$ transforms $\pi(s)$ into $\pi'(s)$ for
every $s\in G$.

\ResultHeading{13.4.6}{13.4.6.}
Let $\pi$ be a continuous unitary representation of $G$, and let
$\xi\in H_\pi$. The function $s\mapsto(\pi(s)\xi\mid\xi)$ is called the
function of positive type defined by $\pi$ and $\xi$. With $\pi$ fixed
and $\xi$ varying in $H_\pi$, one obtains the functions of positive type
associated with $\pi$. If $S$ is a set of representations of $G$, the
functions of positive type associated with $S$ are the functions of
positive type associated with the various elements of $S$.

Let $\varphi$ be a continuous function of positive type on $G$. It
defines a positive form $\omega$ on $L^1(G)$, and hence a pair
$(\pi_\omega,\xi_\omega)$, where $\pi_\omega$ may be regarded as a
continuous unitary representation of $G$. This pair is also denoted by
$(\pi_\varphi,\xi_\varphi)$, and is said to be defined by $\varphi$. Up
to isomorphism, it is characterized by the facts that
\[
 \varphi(s)=(\pi_\varphi(s)\xi_\varphi\mid\xi_\varphi)
 \qquad(s\in G)
\]
and that $\xi_\varphi$ is cyclic for $\pi_\varphi$.

\ResultHeading{13.4.7}{13.4.7. Proposition.}
\begin{itshape}
Let $\varphi$ be a continuous function of positive type on $G$. If
$s,t\in G$, then
\[
 \lvert\varphi(s)-\varphi(t)\rvert^2
 \leq2\varphi(e)\bigl[\varphi(e)-\operatorname{Re}
     \varphi(s^{-1}t)\bigr].
\]
\end{itshape}

% Source PDF physical page 268
Indeed, put $\pi=\pi_\varphi$ and $\xi=\xi_\varphi$. Then
\begin{align*}
 \lvert\varphi(s)-\varphi(t)\rvert^2
 &=\lvert((\pi(s)-\pi(t))\xi\mid\xi)\rvert^2\\
 &\leq\lVert\xi\rVert^2
       \lVert\pi(s)\xi-\pi(t)\xi\rVert^2\\
 &=\varphi(e)\bigl[\lVert\pi(s)\xi\rVert^2
       +\lVert\pi(t)\xi\rVert^2
       -2\operatorname{Re}(\pi(s)\xi\mid\pi(t)\xi)\bigr]\\
 &=\varphi(e)\bigl[2\lVert\xi\rVert^2
       -2\operatorname{Re}(\pi(s^{-1}t)\xi\mid\xi)\bigr]\\
 &=2\varphi(e)\bigl[\varphi(e)-\operatorname{Re}
       \varphi(s^{-1}t)\bigr].
\end{align*}

\ResultHeading{13.4.8}{13.4.8. Proposition.}
\begin{itshape}
Every translate of a continuous function of positive type is a linear
combination of four continuous functions of positive type.
\end{itshape}

Let $\varphi$ be a continuous function of positive type on $G$, put
$\pi=\pi_\varphi$ and $\xi=\xi_\varphi$, and let $a,b\in G$. Set
$\pi(a)\xi=\eta$ and $\pi(b)^{-1}\xi=\zeta$. Then
\begin{align*}
 4\varphi(bsa)
 &=4(\pi(bsa)\xi\mid\xi)
  =4(\pi(s)\eta\mid\zeta)\\
 &=(\pi(s)(\eta+\zeta)\mid\eta+\zeta)
   -(\pi(s)(\eta-\zeta)\mid\eta-\zeta)\\
 &\quad+i(\pi(s)(\eta+i\zeta)\mid\eta+i\zeta)
   -i(\pi(s)(\eta-i\zeta)\mid\eta-i\zeta).
\end{align*}

\ResultHeading{13.4.9}{13.4.9. Proposition.}
\begin{itshape}
Let $\varphi,\varphi'$ be two continuous functions of positive type on
$G$. Then $\varphi\varphi'$ is of positive type. More precisely, put
$\pi=\pi_\varphi$, $\xi=\xi_\varphi$, $\pi'=\pi_{\varphi'}$, and
$\xi'=\xi_{\varphi'}$. Then
\[
 \varphi(s)\varphi'(s)
 =\bigl((\pi\otimes\pi')(s)(\xi\otimes\xi')
       \mid\xi\otimes\xi'\bigr).
\]
\end{itshape}

Indeed,
\begin{align*}
 \bigl((\pi\otimes\pi')(s)(\xi\otimes\xi')
       \mid\xi\otimes\xi'\bigr)
 &=(\pi(s)\xi\otimes\pi'(s)\xi'
       \mid\xi\otimes\xi')\\
 &=(\pi(s)\xi\mid\xi)(\pi'(s)\xi'\mid\xi')
  =\varphi(s)\varphi'(s).
\end{align*}

\ResultHeading{13.4.10}{13.4.10. Proposition.}
\begin{itshape}
Let $\varphi$ be a continuous function of positive type on $G$, and put
$\pi=\pi_\varphi$. Every continuous function of positive type associated
with $\pi$ is the uniform limit on $G$ of functions of the form
\[
 s\longmapsto
 \sum_{i,j=1}^{n}\lambda_i\overline{\lambda_j}
   \varphi(s_j^{-1}ss_i),
 \qquad
 s_1,\ldots,s_n\in G,\quad
 \lambda_1,\ldots,\lambda_n\in\mathbb C.
\]
\end{itshape}

Let $\xi=\xi_\varphi$, $\eta\in H_\pi$, and $\varepsilon>0$. There exist
$s_1,\ldots,s_n\in G$ and $\lambda_1,\ldots,\lambda_n\in\mathbb C$ such
that
\[
 \left\lVert\eta-
   \sum_{i=1}^{n}\lambda_i\pi(s_i)\xi\right\rVert\leq\varepsilon.
\]
Then, for every $s\in G$,
\begin{multline*}
 \left\lvert(\pi(s)\eta\mid\eta)
 -\left(\pi(s)\left(\sum_{i=1}^{n}\lambda_i\pi(s_i)\xi\right)
       \mathrel{\Big|}
       \sum_{i=1}^{n}\lambda_i\pi(s_i)\xi\right)\right\rvert\\
 \leq\varepsilon\lVert\eta\rVert
   +\varepsilon(\lVert\eta\rVert+\varepsilon).
\end{multline*}
% Source PDF physical page 269
But the left-hand side is
\[
 \left\lvert(\pi(s)\eta\mid\eta)
 -\sum_{i,j=1}^{n}\lambda_i\overline{\lambda_j}
    \varphi(s_j^{-1}ss_i)\right\rvert.
\]

\ResultHeading{13.4.11}{13.4.11.}
Let $f\in L^2(G)$. If $\lambda$ denotes the left regular representation
of $G$, then
\begin{align*}
 (\overline{\lambda(s)f}\mid\overline f)
 &=\int\overline{f(s^{-1}t)}f(t)\,dt\\
 &=\int f(t)\widetilde f(t^{-1}s)\,dt
  =(f\ast\widetilde f)(s).
\end{align*}
Thus $f\ast\widetilde f$ is a continuous function of positive type
associated with $\overline\lambda$. If $f\in\mathcal K(G)$, then
$f\ast\widetilde f\in\mathcal K(G)$; by passage to the limit, if
$f\in L^2(G)$, then $f\ast\widetilde f$ is the uniform limit of functions
in $\mathcal K(G)$, that is, it tends to zero at infinity.

Bibliography: \cite{ref43,ref50,ref63,ref108,ref133,ref134}.

\BookSubsectionToc[Weak Convergence and Compact Convergence of Continuous Functions of Positive Type]{13.5}%
{Weak Convergence and Compact Convergence of Continuous Functions of Positive Type}%
{Weak Convergence and Compact Convergence of Continuous\RaggedTitleBreak
 Functions of Positive Type}

\ResultHeading{13.5.1}{13.5.1. Lemma.}
\begin{itshape}
Let $A$ be a bounded subset of $L^\infty(G)$, and let $f\in L^1(G)$. If
$\varphi\in A$ converges weakly to $\varphi_0\in A$, then
$f\ast\varphi$ converges to $f\ast\varphi_0$ for the topology of compact
convergence.
\end{itshape}

Indeed,
\[
 (f\ast\varphi)(s)
 =\int f(t)\varphi(t^{-1}s)\,dt
 =\int f(st)\varphi(t^{-1})\,dt
 =\langle\check\varphi,{}_s f\rangle.
\]
As $\varphi$ tends weakly to $\varphi_0$ while remaining in $A$,
$\langle\check\varphi,g\rangle$ tends to
$\langle\check\varphi_0,g\rangle$ uniformly on every compact subset of
$L^1(G)$. But, as $s$ ranges over a compact subset of $G$, the set of
the ${}_s f$ is strongly compact in $L^1(G)$, since the map
$s\mapsto{}_s f$ from $G$ into $L^1(G)$ is strongly continuous.

\ResultHeading{13.5.2}{13.5.2. Theorem.}
\begin{itshape}
Let $P_1$ be the set of continuous functions $\varphi$ of positive type
on $G$ such that $\varphi(e)=1$. On $P_1$, the weak topology
$\sigma(L^\infty(G),L^1(G))$ coincides with the topology of compact
convergence.
\end{itshape}

Since $\lVert\varphi\rVert_\infty=1$ for every $\varphi\in P_1$, it is
clear that, on $P_1$, the topology of compact convergence is finer than
the weak topology.

Now let $\varphi_0\in P_1$, let $K$ be a compact subset of $G$, and let
$\varepsilon>0$. We shall prove that if $\varphi\in P_1$ belongs to a
suitable weak neighborhood of $\varphi_0$, then
\[
 \lvert\varphi(s)-\varphi_0(s)\rvert
 \leq\varepsilon+4\sqrt\varepsilon
 \qquad(s\in K).
\]
This will complete the proof.

There exists a compact neighborhood $V$ of
$e$ in $G$ such that
\[
 \lvert\varphi_0(s)-1\rvert
 =\lvert\varphi_0(s)-\varphi_0(e)\rvert
 \leq\varepsilon
 \qquad(s\in V).
\]
% ===== ASSEMBLED TRANSLATION CHUNK 10: chunk_270_299.tex =====
% Source PDF physical page 270
Let $\chi$ be the characteristic function of $V$, and let $a>0$ be the
measure of $V$. Let $\mathcal V$ be the weak neighborhood of $\varphi_0$
in $P_1$ defined by the condition
\[
\bigl|\langle\varphi-\varphi_0,\chi\rangle\bigr|\leq\varepsilon a,
\]
that is, by the condition
\[
\left|\int_V\bigl(\varphi(s)-\varphi_0(s)\bigr)\,ds\right|
\leq\varepsilon a.
\]
For $\varphi\in\mathcal V$, we have
\begin{align*}
\left|\int_V(1-\varphi(s))\,ds\right|
&\leq
\left|\int_V(1-\varphi_0(s))\,ds\right|
+\left|\int_V(\varphi(s)-\varphi_0(s))\,ds\right|\\
&\leq2\varepsilon a.
\end{align*}
On the other hand, for $\varphi\in\mathcal V$ and $s\in G$,
\begin{align*}
\bigl|(a^{-1}\chi\ast\varphi)(s)-\varphi(s)\bigr|
&=\left|a^{-1}\int\chi(t)\varphi(t^{-1}s)\,dt-\varphi(s)\right|\\
&=\left|a^{-1}\int_V\varphi(t^{-1}s)\,dt
      -a^{-1}\int_V\varphi(s)\,dt\right|\\
&\leq a^{-1}\int_V|\varphi(t^{-1}s)-\varphi(s)|\,dt.
\end{align*}
By \XRef{13.4.7}{13.4.7}, this is at most
\begin{align*}
a^{-1}\int_V\sqrt2\,(1-\operatorname{Re}\varphi(t))^{1/2}\,dt
&\leq\sqrt2\,a^{-1}
 \left(\int_V(1-\operatorname{Re}\varphi(t))\,dt\right)^{1/2}
 \left(\int_V1\,dt\right)^{1/2}\\
&\leq\sqrt2\,a^{-1}\sqrt{2\varepsilon a}\sqrt a
=2\sqrt\varepsilon.
\end{align*}
There now exists (\XRef{13.5.1}{13.5.1}) a weak neighborhood
$\mathcal V'$ of $\varphi_0$ in $P_1$ such that
$\varphi\in\mathcal V'$ implies
\[
\bigl|(a^{-1}\chi\ast\varphi_0)(s)
      -(a^{-1}\chi\ast\varphi)(s)\bigr|\leq\varepsilon
\qquad\text{for every }s\in K.
\]
Thus, for $\varphi\in\mathcal V\cap\mathcal V'$,
\[
|\varphi(s)-\varphi_0(s)|\leq\varepsilon+4\sqrt\varepsilon
\qquad\text{for every }s\in K.
\]

Bibliography: \cite{ref121,ref187}.

\BookSubsection{13.6}{Pure Functions of Positive Type}

\ResultHeading{13.6.1}{13.6.1. Definition.}
\begin{itshape}
A continuous function $\varphi$ of positive type on $G$ is called
\emph{pure} if $\pi_\varphi$ is irreducible.
\end{itshape}

\ResultHeading{13.6.2}{13.6.2.}
This is equivalent to saying that the positive form on $L^1(G)$ defined
by $\varphi$ is pure (\XRef{2.5.4}{2.5.4}). In view of
\XRef{13.4.5}{13.4.5}, it is also equivalent to saying that, in every
decomposition $\varphi=\varphi_1+\varphi_2$ of $\varphi$ as the sum of
two continuous functions of positive type, $\varphi_1$ and $\varphi_2$
are proportional to $\varphi$.

\ResultHeading{13.6.3}{13.6.3.}
If $G$ is commutative, the continuous irreducible unitary representations
of $G$ are the characters of $G$. The pure continuous functions of
% Source PDF physical page 271
positive type on $G$ are therefore the functions of the form
$\lambda\chi$, where $\lambda\geq0$ and $\chi$ is a character of $G$.

\ResultHeading{13.6.4}{13.6.4. Theorem.}
\begin{itshape}
Let $\varphi$ be a continuous function of positive type on $G$ such that
$\varphi(e)=1$. Then, for the topology of compact convergence,
$\varphi$ is a limit of functions of the form
\[
\lambda_1\varphi_1+\cdots+\lambda_n\varphi_n,
\]
where $\varphi_1,\ldots,\varphi_n$ are pure continuous functions of
positive type equal to $1$ at $e$, and
$\lambda_1,\ldots,\lambda_n$ are numbers $\geq0$ such that
$\lambda_1+\cdots+\lambda_n=1$.
\end{itshape}

The positive form $\psi$ on $L^1(G)$ defined by $\varphi$ is the weak
limit of convex combinations $\psi_i$ of pure states of $A$ and the zero
form (\XRef{2.5.5}{2.5.5}). Since
\[
\liminf_i\|\psi_i\|\geq\|\psi\|=1
\qquad\text{and}\qquad
\|\psi_i\|\leq1,
\]
we may even suppose, after multiplying the $\psi_i$ by suitable scalars,
that $\|\psi_i\|=1$. Then $\psi_i$ is the positive form on $L^1(G)$
defined by a function
$\lambda_1\varphi_1+\cdots+\lambda_n\varphi_n$, where the
$\varphi_j$ are pure continuous functions of positive type,
$\varphi_1(e)=\cdots=\varphi_n(e)=1$, the $\lambda_j$ are nonnegative,
and $\lambda_1+\cdots+\lambda_n=1$. It remains only to apply
\XRef{13.5.2}{13.5.2}.

\ResultHeading{13.6.5}{13.6.5. Corollary.}
\begin{itshape}
Every continuous complex-valued function on $G$ is, for the topology of
compact convergence, a limit of linear combinations of pure functions
of positive type.
\end{itshape}

It plainly suffices to prove the corollary for a function
$f\in\mathcal K(G)$. Such a function is a uniform limit on $G$ of
functions $f\ast g$, with $g\in\mathcal K(G)$. On the other hand,
\begin{align*}
4f\ast g
={}&(f+g)\ast(f+g)^{\sim}-(f-g)\ast(f-g)^{\sim}\\
&+i(f+ig)\ast(f+ig)^{\sim}
-i(f-ig)\ast(f-ig)^{\sim}.
\end{align*}
Finally, every function of the form $h\ast h^{\sim}$, with
$h\in\mathcal K(G)$, is a continuous function of positive type
(\XRef{13.4.11}{13.4.11}), to which \XRef{13.6.4}{13.6.4} can be
applied.

\ResultHeading{13.6.6}{13.6.6. Corollary.}
\begin{itshape}
For every $s\in G$ distinct from $e$, there exists a continuous
irreducible unitary representation $\pi$ of $G$ such that $\pi(s)\ne1$.
\end{itshape}

There is a continuous complex-valued function on $G$ taking distinct
values at $s$ and $e$, and hence, by \XRef{13.6.5}{13.6.5}, a pure
continuous function $\varphi$ of positive type on $G$ taking distinct
values at $s$ and $e$. We have
\[
(\pi_\varphi(s)\xi_\varphi\mid\xi_\varphi)
=\varphi(s)\ne\varphi(e)
=(\xi_\varphi\mid\xi_\varphi),
\]
so $\pi_\varphi(s)\ne1$, and $\pi_\varphi$ is irreducible.

\ResultHeading{13.6.7}{13.6.7.}
Corollary \XRef{13.6.6}{13.6.6}, due to Gelfand and Raikov, says that a
locally compact group has ``enough'' continuous irreducible unitary
representations. This corollary would have failed had we restricted
ourselves to finite-dimensional continuous unitary representations. This
justifies considering infinite-dimensional continuous unitary
representations.

% Source PDF physical page 272
\ResultHeading{13.6.8}{13.6.8.}
We saw in \XRef{13.6.4}{13.6.4} how to reconstruct all continuous
functions of positive type from the pure continuous functions of positive
type. We shall now give a result of the same kind, but in integral form.

Suppose that $G$ is separable. Let $B$ be the convex set of continuous
functions of positive type on $G$ that are at most $1$ at $e$. For the
weak topology, $B$ is compact and is separable because $L^1(G)$ is
separable (\XRef{B.7}{B 7}). For every $s\in G$, the function
$\varphi\mapsto\varphi(s)$ on $B$, which is bounded in modulus by $1$,
is Borel, because it is the limit of a sequence of continuous functions
\[
\varphi\longmapsto\int_G\varphi(t)f_n(t)\,dt;
\]
here the $f_n$ are chosen in $\mathcal K(G)$, nonnegative, of integral
$1$, and with supports in successively smaller neighborhoods of $s$.
On the other hand, the set $P$ of pure continuous functions of positive
type equal to $1$ at $e$ is the set of extreme points of $B$ with $0$
removed (\XRef{2.5.5}{2.5.5}); it is therefore a $G_\delta$ in $B$
(\XRef{B.13}{B 13}). With this understood, we have the following result.

\ResultLabel{Proposition.}
\begin{itshape}
Let $\varphi$ be a continuous function of positive type on $G$ such that
$\varphi(e)=1$. There exists a positive measure of norm $1$ on $B$,
concentrated on $P$, such that
\[
\varphi(s)=\int_P\zeta(s)\,d\mu(\zeta)
\qquad\text{for every }s\in G.
\]
\end{itshape}

There exists a positive measure $\mu$ of norm $1$ on $B$, concentrated
on $P\cup\{0\}$, such that
\[
\varphi=\int_{P\cup\{0\}}\zeta\,d\mu(\zeta),
\]
where the integral is taken in the weak sense (\XRef{B.13}{B 13}). We
may plainly suppose that $\mu(\{0\})=0$, and hence that $\mu$ is
concentrated on $P$. Let $f\in L^1(G)$. We have
\begin{equation}
\int_G\varphi(s)f(s)\,ds
=\int_Bd\mu(\zeta)\int_G\zeta(s)f(s)\,ds.
\tag{1}\label{eq:13.6.8-1}
\end{equation}
The function $(\zeta,s)\mapsto\zeta(s)$ is continuous on $P\times G$ by
\XRef{13.5.2}{13.5.2}, and is therefore measurable on $B\times G$ for
the product measure $d\mu(\zeta)\,ds$; moreover it is bounded in modulus
by $1$. Thus $(\zeta,s)\mapsto\zeta(s)f(s)$ is integrable for
$d\mu(\zeta)\,ds$, and \eqref{eq:13.6.8-1} may be written
\[
\int_G\varphi(s)f(s)\,ds
=\int_Gf(s)\,ds\int_B\zeta(s)\,d\mu(\zeta).
\]
Consequently,
\begin{equation}
\varphi(s)=\int_B\zeta(s)\,d\mu(\zeta)
\tag{2}\label{eq:13.6.8-2}
\end{equation}
almost everywhere on $G$. On the other hand, if
$(s_1,s_2,\ldots)$ is a sequence of elements of $G$ tending to $s$, then
$\zeta(s_n)\to\zeta(s)$ for every $\zeta\in B$, and hence
\[
\int_B\zeta(s_n)\,d\mu(\zeta)
\longrightarrow\int_B\zeta(s)\,d\mu(\zeta)
\]
by the Lebesgue theorem. Thus $\int_B\zeta(s)\,d\mu(\zeta)$ depends
continuously on $s$, so \eqref{eq:13.6.8-2} holds everywhere on $G$.

% Source PDF physical page 273
Bibliography: \cite{ref43,ref50,ref52,ref108,ref134,ref185}.

\BookSubsection{13.7}{Measures of Positive Type}

\ResultHeading{13.7.1}{13.7.1. Definition.}
\begin{itshape}
A complex measure $\mu$ on $G$ is said to be of positive type if
\begin{equation}
\langle\mu,f\ast f^{\sim}\rangle\geq0
\tag{1}\label{eq:13.7.1-1}
\end{equation}
for every $f\in\mathcal K(G)$. We then write $\mu\gg0$.
\end{itshape}

\ResultHeading{13.7.2}{13.7.2.}
If $\mu\gg0$, then $\mu^{\ast}=\mu$. Indeed,
\[
\langle\mu^{\ast},f\ast f^{\sim}\rangle
=\overline{\langle\mu,(f\ast f^{\sim})^{\sim}\rangle}
=\overline{\langle\mu,f\ast f^{\sim}\rangle}
=\langle\mu,f\ast f^{\sim}\rangle
\]
for every $f\in\mathcal K(G)$. Hence, by polarization,
\[
\langle\mu^{\ast},f\ast g^{\sim}\rangle
=\langle\mu,f\ast g^{\sim}\rangle
\]
for all $f,g\in\mathcal K(G)$, and therefore, by passage to the limit,
$\langle\mu^{\ast},h\rangle=\langle\mu,h\rangle$ for every
$h\in\mathcal K(G)$.

\ResultHeading{13.7.3}{13.7.3.}
Relation \eqref{eq:13.7.1-1} may be written
\[
\iint f(s)\overline{f(t^{-1}s)}\,ds\,d\mu(t)\geq0
\]
for every $f\in\mathcal K(G)$, or, on replacing $f$ by $\overline f$,
\[
\int_G(\mu\ast f)(s)\overline{f(s)}\,ds\geq0
\]
for every $f\in\mathcal K(G)$.

\ResultHeading{13.7.4}{13.7.4.}
Suppose that $\mu$ is bounded. If $\lambda$ is the left regular
representation of $G$, the operator $\lambda(\mu)$ is defined. By
\XRef{13.7.3}{13.7.3}, saying that $\mu\gg0$ means that
$(\lambda(\mu)f\mid f)\geq0$ for every $f\in\mathcal K(G)$. Since
$\mathcal K(G)$ is dense in $L^2(G)$,
\[
\mu\gg0\quad\Longleftrightarrow\quad\lambda(\mu)\geq0.
\]

\ResultHeading{13.7.5}{13.7.5.}
Let $\Delta$ be the modular function of $G$, and let
$f\in\mathcal K(G)$. We have
\begin{align*}
(f\ast f^{\ast})(s)\Delta(s)^{1/2}
&=\int f(t)\overline{f(s^{-1}t)}
       \Delta(s^{-1}t)\Delta(s)^{1/2}\,dt\\
&=\int f(t)\Delta(t)^{1/2}\overline{f(s^{-1}t)}
       \Delta(s^{-1}t)^{1/2}\,dt\\
&=\bigl((f\Delta^{1/2})\ast(f\Delta^{1/2})^{\sim}\bigr)(s).
\end{align*}
Thus, on replacing $f\Delta^{1/2}$ by $f$, condition
\eqref{eq:13.7.1-1} is also equivalent to
\[
\langle\Delta^{1/2}\mu,f\ast f^{\ast}\rangle\geq0
\qquad\text{for every }f\in\mathcal K(G).
\]

% Source PDF physical page 274
\ResultHeading{13.7.6}{13.7.6. Definition.}
\begin{itshape}
A locally integrable function $\varphi$ on $G$ is said to be of positive
type if the measure $\Delta^{-1/2}(s)\varphi(s)\,ds$ on $G$ is of
positive type. We then write $\varphi\gg0$.
\end{itshape}

By \XRef{13.7.5}{13.7.5}, this is equivalent to
\[
\langle\varphi,f\ast f^{\ast}\rangle\geq0
\qquad\text{for every }f\in\mathcal K(G).
\]
For continuous $\varphi$, this gives precisely the earlier notion
(\XRef{13.4.4}{13.4.4}). If $\varphi$ and $\psi$ are two locally
integrable functions on $G$ such that $\psi-\varphi\gg0$, we write
$\psi\gg\varphi$.

\ResultHeading{13.7.7}{13.7.7.}
If $\varphi\gg0$, then the measures satisfy
\[
\Delta^{-1/2}(s)\varphi(s)\,ds
=\bigl(\Delta^{-1/2}(s)\varphi(s)\,ds\bigr)^{\ast}
=\Delta^{1/2}(s)\overline{\varphi(s^{-1})}\,d(s^{-1})
=\Delta^{-1/2}(s)\overline{\varphi(s^{-1})}\,ds,
\]
and hence
\[
\varphi=\varphi^{\sim}
\qquad\text{locally almost everywhere}.
\]

\ResultHeading{13.7.8}{13.7.8.}
Let $\varphi$ be a complex-valued function on $G$ such that
$\Delta^{-1/2}\varphi\in L^1(G)$. By \XRef{13.7.4}{13.7.4},
$\varphi\gg0$ if and only if
\[
\lambda(\Delta^{-1/2}\varphi)\geq0,
\]
where $\lambda$ denotes the left regular representation of $G$.

\ResultHeading{13.7.9}{13.7.9.}
Let $\mu$ be a measure of positive type on $G$. For
$f,g\in\mathcal K(G)$, put
\[
(f\mid g)_\mu
=\langle\mu,g^{\sim}\ast f\rangle
=\iint f(t^{-1}s)\overline{g(t^{-1})}\,dt\,d\mu(s).
\]
Then $\mathcal K(G)$ becomes a pre-Hilbert space. Let $H_\mu$ be the
separated Hilbert-space completion of this pre-Hilbert space. For
$s\in G$ and $f\in\mathcal K(G)$, put
\[
\sigma(s)f=(\varepsilon_s\ast f)\Delta(s)^{1/2},
\]
that is,
\[
(\sigma(s)f)(x)=f(s^{-1}x)\Delta(s)^{1/2}.
\]
We have
\begin{align*}
\|\sigma(s)f\|_\mu^2
&=\iint\overline{(\sigma(s)f)(t^{-1})}
          (\sigma(s)f)(t^{-1}u)\,dt\,d\mu(u)\\
&=\iint\overline{f(s^{-1}t^{-1})}\Delta(s)^{1/2}
          f(s^{-1}t^{-1}u)\Delta(s)^{1/2}\,dt\,d\mu(u)\\
&=\iint\overline{f(t^{-1})}f(t^{-1}u)
          \Delta(s)\Delta(s)^{-1}\,dt\,d\mu(u)\\
&=\iint\overline{f(t^{-1})}f(t^{-1}u)\,dt\,d\mu(u)
=\|f\|_\mu^2.
\end{align*}
% Source PDF physical page 275
Thus $\sigma(s)$ defines a unitary operator on $H_\mu$, which we again
denote by $\sigma(s)$. It is clear that $\sigma$ is a unitary
representation of $G$ on $H_\mu$. Moreover, as $s\to s_0$,
$\sigma(s)f$ tends uniformly to $\sigma(s_0)f$, while its support remains
contained in a fixed compact set. It follows readily that the unitary
representation $\sigma$ is continuous.

\ResultHeading{13.7.10}{13.7.10. Proposition.}
\begin{itshape}
Let $V$ be a neighborhood of $e$, and let $\mu$ be a measure of positive
type on $G$ such that
\begin{equation}
\langle\mu,f^{\sim}\ast f\rangle
\leq K\left(\int f(x)\,dx\right)^2
\tag{1}\label{eq:13.7.10-1}
\end{equation}
for all nonnegative $f\in\mathcal K(G)$ vanishing outside $V$, where $K$
is independent of $f$. Then
\[
d\mu(s)=\varphi(s)\Delta^{-1/2}(s)\,ds,
\]
where $\varphi$ is a continuous function of positive type.
\end{itshape}

There is a family $(f_i)$ of nonnegative functions in $\mathcal K(G)$,
indexed by an increasing directed set, vanishing outside $V$, such that
$\int f_i(x)\,dx=1$ and
\[
\int f_i(x)g(x)\,dx\longrightarrow g(e)
\qquad\text{for every }g\in\mathcal K(G).
\]
By \eqref{eq:13.7.10-1}, $\|f_i\|_\mu^2\leq K$ for every $i$. On the
other hand, for every $g\in\mathcal K(G)$,
\begin{align*}
(f_i\mid g)_\mu
&=\iint\overline{g(t^{-1})}f_i(t^{-1}s)\,dt\,d\mu(s)\\
&=\int d\mu(s)\int\overline{g(t^{-1}s^{-1})}f_i(t^{-1})\,dt
\longrightarrow\int\overline{g(s^{-1})}\,d\mu(s).
\end{align*}
Hence, if $\Lambda$ denotes the canonical map from $\mathcal K(G)$ into
$H_\mu$, $\Lambda f_i$ converges weakly to an element $\xi\in H_\mu$
such that
\begin{equation}
(\xi\mid\Lambda g)=\int\overline{g(s^{-1})}\,d\mu(s)
\qquad(g\in\mathcal K(G)).
\tag{2}\label{eq:13.7.10-2}
\end{equation}
Put
\[
\varphi(s)=(\sigma(s)\xi\mid\xi)
\qquad(s\in G).
\]
Then $\varphi$ is a continuous function of positive type on $G$
(\XRef{13.4.5}{13.4.5}). Moreover, \eqref{eq:13.7.10-2} implies, for
every $t\in G$,
\begin{align*}
(\sigma(t)\xi\mid\Lambda g)
&=(\xi\mid\sigma(t^{-1})\Lambda g)\\
&=\int\overline{(\sigma(t^{-1})g)(s^{-1})}\,d\mu(s)\\
&=\int\overline{g(ts^{-1})}\Delta^{-1/2}(t)\,d\mu(s).
\end{align*}
% Source PDF physical page 276
Thus, for every $f\in\mathcal K(G)$,
\begin{align*}
\int f(t)(\sigma(t)\xi\mid\Lambda g)\,dt
&=\iint f(t)\overline{g(ts^{-1})}\Delta^{-1/2}(t)\,dt\,d\mu(s)\\
&=\iint f(t^{-1})\overline{g(t^{-1}s^{-1})}
       \Delta^{-1/2}(t)\,dt\,d\mu(s)\\
&=\iint\overline{g(t^{-1})}f(t^{-1}s)
       \Delta^{1/2}(s^{-1}t)\,dt\,d\mu(s)\\
&=\langle\mu,g^{\sim}\ast(\Delta^{1/2}f)\rangle
=(\Delta^{1/2}f\mid g)_\mu.
\end{align*}
It follows that
\[
\int f(t)(\sigma(t)\xi\mid\xi)\,dt
=(\Delta^{1/2}f\mid\xi).
\]
In view of \eqref{eq:13.7.10-2} and $\mu^{\ast}=\mu$, this says
\[
\int f(t)\varphi(t)\,dt
=\int\Delta^{1/2}(t^{-1})f(t^{-1})\,d\bar\mu(t)
=\int\Delta^{-1/2}(t)f(t)\,d\mu(t).
\]
Since this holds for every $f\in\mathcal K(G)$, we conclude that
\[
d\mu(t)=\Delta^{-1/2}(t)\varphi(t)\,dt.
\]

\ResultHeading{13.7.11}{13.7.11. Corollary.}
\begin{itshape}
Let $\varphi$ be a continuous function of positive type, and let $\psi$
be a locally integrable function of positive type such that
$\varphi\gg\psi$. Then $\psi$ is locally almost everywhere equal to a
continuous function of positive type.
\end{itshape}

Indeed, for every $f\in\mathcal K(G)$,
\[
\int\psi(s)\Delta^{-1/2}(s)(f^{\sim}\ast f)(s)\,ds
\leq
\int\varphi(s)\Delta^{-1/2}(s)(f^{\sim}\ast f)(s)\,ds.
\]
If $f\in\mathcal K(G)$ is nonnegative and vanishes outside a compact
symmetric neighborhood $V$ of $e$, the right-hand side is at most
\begin{align*}
&\sup_{s\in V^2}|\varphi(s)\Delta^{-1/2}(s)|
 \iint f(t^{-1})f(t^{-1}s)\,dt\,ds\\
&\qquad=
\sup_{s\in V^2}|\varphi(s)\Delta^{-1/2}(s)|
\sup_{t\in V}\Delta(t)
\left(\int f(t)\,dt\right)^2,
\end{align*}
and it remains only to apply \XRef{13.7.10}{13.7.10}.

Bibliography: \cite{ref43,ref50}.

\BookSubsection{13.8}{Square-Integrable Functions of Positive Type}

\ResultHeading{13.8.1}{13.8.1.}
Throughout \S13.8, $\lambda$ denotes the left regular representation of
$G$. Let $f\in L^2(G)$. For every $g\in\mathcal K(G)$,
$\lambda(g)f\in L^2(G)$. If there is a finite constant $M$ such that
\[
\|\lambda(g)f\|_2\leq M\|g\|_2
\qquad\text{for every }g\in\mathcal K(G),
\]
% Source PDF physical page 277
we shall say that $f$ is a \emph{bounded element} of $L^2(G)$. The map
$g\mapsto\lambda(g)f$ then extends uniquely to a continuous linear
operator on $L^2(G)$, denoted by $\rho(f)$. For every
$g\in\mathcal K(G)$,
\begin{align*}
(\rho(f)\overline g\mid\overline g)
&=(\lambda(\overline g)f\mid\overline g)
=\int(\overline g\ast f)(s)g(s)\,ds\\
&=\iint\overline{g(t)}f(t^{-1}s)g(s)\,ds\,dt
=\langle f,g^{\sim}\ast g\rangle.
\end{align*}
Thus $f\gg0$ if and only if $\rho(f)\geq0$
(\XRef{13.7.6}{13.7.6}).

\ResultHeading{13.8.2}{13.8.2.}
Now let $f$ be an element of $L^2(G)$ of positive type, but not
necessarily bounded. The operator $g\mapsto\lambda(g)f$, defined on
$\mathcal K(G)$ and taking values in $L^2(G)$, is nonnegative by the same
calculation as in \XRef{13.8.1}{13.8.1}. We denote by $\rho(f)$ its
Friedrichs extension (\XRef{B.23}{B 23}), which is self-adjoint and
nonnegative. When $f$ is bounded, this is the operator $\rho(f)$ defined
in \XRef{13.8.1}{13.8.1}.

\ResultHeading{13.8.3}{13.8.3. Lemma.}
\begin{itshape}
Let $f$ be an element of $L^2(G)$ of positive type.
\begin{enumerate}[label=\textup{(\roman*)}]
\item For every $s\in G$, $\rho(f)$ commutes with $\lambda(s)$.
\item If $h$ belongs to the domain $D$ of $\rho(f)$, then
$\rho(f)h=h\ast f$.
\end{enumerate}
\end{itshape}

For $g\in\mathcal K(G)$ and $s\in G$, we have
\[
\varepsilon_s\ast g\in\mathcal K(G),
\qquad
\varepsilon_s\ast(g\ast f)=(\varepsilon_s\ast g)\ast f,
\]
that is, $\lambda(s)\rho(f)g=\rho(f)\lambda(s)g$. Hence
\[
\rho(f)|_{\mathcal K(G)}
=\lambda(s)\rho(f)\lambda(s)^{-1}|_{\mathcal K(G)},
\]
and consequently
\[
\rho(f)=\lambda(s)\rho(f)\lambda(s)^{-1}
\]
by \XRef{B.23}{B 23}.

On the other hand, on $D$ the operator $\rho(f)$ coincides with the
adjoint of $\rho(f)|_{\mathcal K(G)}$ (\XRef{B.23}{B 23}). Thus, if
$h\in D$ and $g\in\mathcal K(G)$,
\begin{equation}
(\rho(f)h\mid g)=(h\mid\rho(f)g)
=\int h(s)\,ds\int\overline{g(t)}f(t^{-1}s)\,dt.
\tag{1}\label{eq:13.8.3-1}
\end{equation}
Now $h(s)\overline{g(t)}f(t^{-1}s)$ is measurable on $G\times G$ for
$ds\,dt$, and vanishes outside a countable union of sets integrable for
$ds\,dt$. Moreover,
\[
\int^{\ast}ds\int^{\ast}|h(s)g(t)f(t^{-1}s)|\,dt
\leq\int|h(s)|\bigl(|g|\ast|f|\bigr)(s)\,ds<+\infty,
\]
because $|h|\in L^2(G)$ and $|g|\ast|f|\in L^2(G)$. Thus the function
is integrable for $ds\,dt$, and \eqref{eq:13.8.3-1} becomes
\[
(\rho(f)h\mid g)
=\int\left[\int h(s)f(s^{-1}t)\,ds\right]\overline{g(t)}\,dt
=\int(h\ast f)(t)\overline{g(t)}\,dt,
\]
whence $\rho(f)h=h\ast f$ almost everywhere.

% Source PDF physical page 278
\ResultHeading{13.8.4}{13.8.4. Lemma.}
\begin{itshape}
Let $a,b\in L^2(G)$ be two bounded elements of positive type such that
$\rho(a)$ and $\rho(b)$ commute. Then $a\ast b$ is a continuous
function of positive type and a bounded element of $L^2(G)$. We have
\[
\rho(a\ast b)=\rho(a)\rho(b),
\qquad
(a\mid b)\geq0.
\]
If, in addition, $a\ll b$, then
\[
\|b-a\|_2^2\leq\|b\|_2^2-\|a\|_2^2.
\]
\end{itshape}

The convolution product of two functions in $L^2(G)$ is continuous on
$G$. Let $(a_n)$ be a sequence of functions in $\mathcal K(G)$ tending
to $a$ in $L^2(G)$. Then
$a_n\ast b=\rho(b)a_n$ tends to $\rho(b)a$ in $L^2(G)$. On the other
hand,
\[
\|a_n\ast b-a\ast b\|_\infty\longrightarrow0,
\]
and hence $a\ast b=\rho(b)a\in L^2(G)$.

Let $f\in\mathcal K(G)$. We have $f\ast a_n\in\mathcal K(G)$, and
$f\ast a_n$ tends to $f\ast a$ in $L^2(G)$. Therefore
\[
(f\ast a_n)\ast b=\rho(b)(f\ast a_n)
\]
tends to $\rho(b)(f\ast a)=\rho(b)\rho(a)f$ in $L^2(G)$. On the other
hand, $f\ast(a_n\ast b)$ tends to $f\ast(a\ast b)$ in $L^2(G)$.
Thus
\[
f\ast(a\ast b)=\rho(b)\rho(a)f
=\rho(a)\rho(b)f
\qquad\text{for every }f\in\mathcal K(G),
\]
which proves that $a\ast b$ is bounded and that
$\rho(a\ast b)=\rho(a)\rho(b)$. Since $\rho(a)$ and $\rho(b)$ are
nonnegative and commute, $\rho(a)\rho(b)$ is nonnegative; hence
$a\ast b$ is of positive type. We have
\[
(a\mid b)=\int a(s)b(s^{-1})\,ds=(a\ast b)(e)\geq0.
\]
Finally, if $a\ll b$, put $c=b-a\gg0$. Then $c$ is bounded, and by what
precedes,
\[
(a\mid a)\leq(a\mid a)+(a\mid c)=(a\mid b),
\]
so
\[
\|b-a\|_2^2
=\|b\|_2^2+\|a\|_2^2-2(a\mid b)
\leq\|b\|_2^2-\|a\|_2^2.
\]

\ResultHeading{13.8.5}{13.8.5. Lemma.}
\begin{itshape}
Let $a_1,a_2,\ldots$ be bounded elements of positive type in $L^2(G)$
such that $a_1\ll a_2\ll\cdots$ and the $\rho(a_i)$ commute pairwise. If
\[
\sup_i\|a_i\|_2<+\infty,
\]
then the $a_i$ have a strong limit in $L^2(G)$.
\end{itshape}

By \XRef{13.8.4}{13.8.4}, the numbers $\|a_i\|_2$ form an increasing,
and therefore convergent, sequence. Another application of
\XRef{13.8.4}{13.8.4} shows that the $a_i$ form a Cauchy sequence.

\ResultHeading{13.8.6}{13.8.6. Theorem.}
\begin{itshape}
Let $\varphi$ be a continuous square-integrable function of positive type
on $G$. There exists a square-integrable function $\psi$ of positive type
on $G$ such that
\[
\varphi=\psi\ast\psi=\psi\ast\psi^{\sim}.
\]
\end{itshape}

First suppose that $\varphi$ is bounded. Multiplying $\varphi$ by a
positive constant if necessary, we may suppose that
$0\leq\rho(\varphi)\leq1$. Let $p_1,p_2,\ldots$ be
% Source PDF physical page 279
a sequence of nonnegative polynomials on $[0,1]$, vanishing at $0$, and
converging uniformly and increasingly to $\sqrt t$ on $[0,1]$. By
\XRef{13.8.4}{13.8.4}, we may form
\[
\psi_1=p_1(\varphi),\quad
\psi_2=p_2(\varphi),\quad\ldots,
\]
where the multiplication used is convolution. These elements belong to
$L^2(G)$ and are bounded, and
\[
\rho(\psi_1)=p_1(\rho(\varphi)),\qquad
\rho(\psi_2)=p_2(\rho(\varphi)),\quad\ldots.
\]
We see that
\[
0\leq\rho(\psi_1)\leq\rho(\psi_2)\leq\cdots,
\]
and hence $0\ll\psi_1\ll\psi_2\ll\cdots$; moreover, the
$\rho(\psi_n)$ commute pairwise. Since $p_n(t)^2\leq t$ on $[0,1]$,
we have $\rho(\psi_n)^2\leq\rho(\varphi)$, and hence
$\psi_n\ast\psi_n\ll\varphi$. Consequently,
\[
(\psi_n\ast\psi_n)(e)=\|\psi_n\|_2^2\leq\varphi(e).
\]
Thus the $\psi_n$ converge strongly to an element $\psi\in L^2(G)$
(\XRef{13.8.5}{13.8.5}). At the same time,
$\rho(\psi_n)=p_n(\rho(\varphi))$ converges in norm to
$\rho(\varphi)^{1/2}$. For $f\in\mathcal K(G)$,
$f\ast\psi_n=\rho(\psi_n)f$ tends in $L^2(G)$ both to
$f\ast\psi$ and to $\rho(\varphi)^{1/2}f$. Hence $\psi$ is bounded and
of positive type, and
\[
\rho(\psi)=\rho(\varphi)^{1/2}.
\]
This implies
\[
\rho(\varphi)=\rho(\psi)^2=\rho(\psi\ast\psi),
\]
and therefore $\varphi=\psi\ast\psi$.

We now pass to the general case. Let
\[
\rho(\varphi)=\int_0^{+\infty}\zeta\,dE_\zeta
\]
be the spectral decomposition of $\rho(\varphi)$. By
\XRef{13.8.3}{13.8.3}(i), the projections $E_\zeta$ commute with the
$\lambda(s)$, $s\in G$. Put $\varphi_\zeta=E_\zeta\varphi$. For every
$g\in\mathcal K(G)$,
\begin{align*}
(g\ast\varphi_\zeta^{\sim})(t)
&=\int g(s)\varphi_\zeta^{\sim}(s^{-1}t)\,ds
=\int g(s)\overline{\varphi_\zeta(t^{-1}s)}\,ds\\
&=(g\mid\lambda(t)E_\zeta\varphi)
=(E_\zeta g\mid\lambda(t)\varphi)\\
&=\int(E_\zeta g)(s)\overline{\varphi(t^{-1}s)}\,ds
=\int(E_\zeta g)(s)\varphi(s^{-1}t)\,ds\\
&=(E_\zeta g\ast\varphi)(t).
\end{align*}
By \XRef{13.8.3}{13.8.3}(ii), this is equal to
$(\rho(\varphi)E_\zeta g)(t)$. Hence $\varphi_\zeta^{\sim}$ is bounded
and
\[
\rho(\varphi_\zeta^{\sim})=\rho(\varphi)E_\zeta\geq0,
\]
so $\varphi_\zeta^{\sim}=\varphi_\zeta\gg0$. Moreover,
$\rho(\varphi_\zeta)\leq\rho(\varphi)$; hence
$\varphi_\zeta\ll\varphi$, and therefore $\varphi_\zeta$ is continuous
by \XRef{13.7.11}{13.7.11}. By the first part of the proof there is a
bounded element $\omega_\zeta$ of positive type such that
\[
\varphi_\zeta=\omega_\zeta\ast\omega_\zeta.
\]
Take in particular $\zeta=1,2,\ldots,n,\ldots$. The
$\rho(\varphi_n)$ commute pairwise and increase; consequently so do
\[
\rho(\omega_n)=\rho(\varphi_n)^{1/2}.
\]
Thus $\omega_1\ll\omega_2\ll\cdots$. On the other hand,
\[
\|\omega_n\|_2^2=\varphi_n(e)\leq\varphi(e),
\]
so the $\omega_n$ have a strong limit $\omega\gg0$
(\XRef{13.8.5}{13.8.5}). Then
$\varphi_n=\omega_n\ast\omega_n$ converges uniformly on $G$ to
$\omega\ast\omega$, while $\varphi_n=E_n\varphi$ converges strongly to
$\varphi$ in $L^2(G)$. Hence $\varphi=\omega\ast\omega$.

Bibliography: \cite{ref50}.

\BookSubsection[C*-Algebra of a Locally Compact Group]{13.9}
  {$C^{*}$-Algebra of a Locally Compact Group}

\ResultHeading{13.9.1}{13.9.1.}
Since $L^1(G)$ is an involutive Banach algebra with an approximate
identity, we may form its enveloping $C^{*}$-algebra
% Source PDF physical page 280
(\XRef{2.7.2}{2.7.2}). This $C^{*}$-algebra is called the
$C^{*}$-algebra of $G$ and is denoted by $C^{*}(G)$.

For $f\in L^1(G)$, put
\[
 \|f\|'=\sup_\pi\|\pi(f)\|\leq\|f\|_1,
\]
where $\pi$ ranges over the nondegenerate representations of $L^1(G)$ or,
equivalently, over the continuous unitary representations of $G$. Then
$f\mapsto\|f\|'$ is a seminorm on
$L^1(G)$ (\XRef{2.7.1}{2.7.1}), and indeed a norm, since $L^1(G)$ has
an injective representation (\XRef{13.3.6}{13.3.6}). The
$C^{*}$-algebra of $G$ is precisely the completion of $L^1(G)$ for this
norm.

\ResultHeading{13.9.2}{13.9.2.}
If $G$ is discrete, $C^{*}(G)$ has an identity element. If $G$ is
separable, $C^{*}(G)$ is separable (\XRef{13.2.4}{13.2.4}).

\ResultHeading{13.9.3}{13.9.3.}
By \XRef{2.7.4}{2.7.4} and \XRef{13.3.5}{13.3.5}, there is a
bijective correspondence between continuous unitary representations of
$G$ and nondegenerate representations of $C^{*}(G)$. Everything said in
\XRef{13.3.5}{13.3.5} remains valid with $L^1(G)$ replaced by
$C^{*}(G)$.

The left regular representation of $G$ corresponds to a representation
of $C^{*}(G)$ called the left regular representation of $C^{*}(G)$ on
$L^2(G)$.

\ResultHeading{13.9.4}{13.9.4.}
The group $G$ is called liminal, postliminal, antiliminal, or of type I,
according as $C^{*}(G)$ is liminal, postliminal, antiliminal, or of type I.

For $G$ to be liminal, it is necessary and sufficient that, for every
irreducible continuous unitary representation $\pi$ of $G$ and every
$f\in L^1(G)$, the operator $\pi(f)$ be compact.

Suppose that $G$ is postliminal. Then, for every irreducible continuous
unitary representation $\pi$ of $G$, the norm closure of
$\pi(L^1(G))$ contains $\mathcal{LC}(H_\pi)$ (\XRef{4.3.7}{4.3.7}).
The converse is true if $G$ is separable (\XRef{9.1}{9.1}).

For $G$ to be of type I, it is necessary and sufficient that, for every
continuous unitary representation $\pi$ of $G$, the von Neumann algebra
generated by $\pi(G)$ be of type I. If $G$ is postliminal, then $G$ is
of type I (\XRef{5.5.2}{5.5.2}). If $G$ is separable, the following
conditions are equivalent:
\begin{enumerate}[label=\arabic*${}^{\circ}$]
 \item $G$ is of type I;
 \item for every factor continuous unitary representation $\pi$ of
 $G$, the factor generated by $\pi(G)$ is of type I;
 \item $G$ is postliminal.
\end{enumerate}
This follows from \XRef{9.1}{9.1}.

Bibliography: \cite{ref27,ref77}.

\BookSubsection{13.10}{Hilbert Algebra of a Locally Compact Group}

Throughout this section, $G$ denotes a unimodular locally compact group.

\ResultHeading{13.10.1}{13.10.1.}
Under convolution and the involution $f\mapsto f^{\ast}$,
$\mathcal K(G)$ is an involutive algebra. Endow $\mathcal K(G)$ with the
scalar product
% Source PDF physical page 281
\[
 (f\mid g)=\int f(s)\overline{g(s)}\,ds.
\]
It is easy to see that $\mathcal K(G)$ then becomes a Hilbert algebra; its
completed Hilbert space is precisely $L^2(G)$. The completed Hilbert
algebra $A$ of bounded elements (\XRef{A.57}{A 57}) is called the
Hilbert algebra of $G$; one has
\[
 \mathcal K(G)\subset A\subset L^2(G).
\]
It is clear that $A$ is stable under $f\mapsto\overline f$, and hence also
under $f\mapsto\tilde f=(\overline f)^{\ast}$.

\ResultHeading{13.10.2}{13.10.2.}
A continuous linear operator on $L^2(G)$ commutes with the left
translation operators if and only if it commutes with the operators of
left convolution by the elements of $\mathcal K(G)$ (this follows, for
example, from \XRef{13.3.5}{13.3.5} applied to the left regular
representation of $G$). Thus $\mathfrak U(A)$ is the von Neumann algebra
generated by the left translation operators on $L^2(G)$. Similarly,
$\mathfrak V(A)$ is the von Neumann algebra generated by the right
translation operators on $L^2(G)$.

\ResultHeading{13.10.3}{13.10.3.}
If $f\in A$, recall that $U_f,V_f$ denote the continuous linear operators
on $L^2(G)$ extending left and right multiplication by $f$ in $A$. If
$f\in A$ and $g\in L^2(G)$, then
\[
 U_fg=f\ast g,\qquad V_fg=g\ast f.
\]
Indeed, first suppose that $g\in\mathcal K(G)$, and let $(f_n)$ be a
sequence of elements of $\mathcal K(G)$ converging to $f$ in mean square.
Then $V_gf_n=f_n\ast g$ converges in mean square both to
$V_gf=U_fg$ and to $f\ast g$; hence $U_fg=f\ast g$. In the general
case, let $(g_n)$ be a sequence of elements of $\mathcal K(G)$ converging
to $g$ in mean square. Then $U_fg_n$ converges to $U_fg$ in mean square,
and $f\ast g_n$ converges uniformly on $G$ to $f\ast g$. Since
$U_fg_n=f\ast g_n$ by the first part of the argument, we obtain
$U_fg=f\ast g$.

\ResultHeading{13.10.4}{13.10.4.}
Recall that the map $f\mapsto f^{\ast}$ on $L^2(G)$ commutes with every
Hermitian element of $\mathfrak U(A)\cap\mathfrak V(A)$, the common center
of $\mathfrak U(A)$ and $\mathfrak V(A)$ (\XRef{A.54}{A 54}). Recall also
that $\mathfrak U(A)$ and $\mathfrak V(A)$ are each other's commutants and
are semifinite von Neumann algebras (\XRef{A.60}{A 60}).

\ResultHeading{13.10.5}{13.10.5. Proposition.}
\begin{itshape}
Let $G$ be a unimodular locally compact group and $A$ its Hilbert algebra.
The following conditions are equivalent:
\begin{enumerate}[label=(\roman*)]
 \item[\normalfont(i)] The von Neumann algebra $\mathfrak U(A)$ is
 finite;
 \item[\normalfont(i$'$)] the von Neumann algebra $\mathfrak V(A)$ is
 finite;
 \item[\normalfont(ii)] there exists in $G$ a fundamental system of
 compact neighborhoods of $e$ invariant under the inner automorphisms of
 $G$.
\end{enumerate}
\end{itshape}

Conditions (i) and (i$'$) are equivalent for every Hilbert algebra
(\XRef{A.63}{A 63}).

Suppose that there exists a fundamental system $(V_i)_{i\in I}$ of compact
neighborhoods of $e$ invariant under the inner automorphisms of $G$. Let
$f_i$ be
% Source PDF physical page 282
the characteristic function of $V_i$. The $f_i$ are central elements of
$A$ [they belong to the center of $L^1(G)$]. On the other hand, every
$f\in L^2(G)$ is a strong limit of elements $f_i\ast f$. Thus the
characteristic projection of $A$ is $1$ (\XRef{A.62}{A 62}). Hence
$\mathfrak U(A)$ and $\mathfrak V(A)$ are finite von Neumann algebras
(\XRef{A.63}{A 63}).

Conversely, suppose that $\mathfrak U(A)$ and $\mathfrak V(A)$ are finite
von Neumann algebras. The characteristic projection of $A$ is $1$
(\XRef{A.63}{A 63}). Let $s\in G$, $s\ne e$. The operator
$f\mapsto{}_s f$ on $L^2(G)$ is not the identity; it commutes with right
translations, while the set of right translates of the elements of
$L^2(G)$ central relative to $A$ is total in $L^2(G)$
(\XRef{A.62}{A 62}). Hence there exists an $f\in L^2(G)$ central
relative to $A$ such that ${}_s f\ne f$. Since $f$ is central,
${}_t f=f_t$ for every $t\in G$ (\XRef{A.62}{A 62}); hence $f$ is
invariant under the inner automorphisms of $G$. Consider the function
\[
 t\longmapsto g(t)=({}_t f\mid f)
\]
on $G$. It is continuous, invariant under inner automorphisms, tends to
zero at infinity, and $g(s)\ne g(e)$. If $W$ is a compact neighborhood
of $g(e)$ in $\mathbb C$ containing neither $0$ nor $g(s)$, then the
condition $g(t)\in W$ defines a compact neighborhood of $e$, invariant
under inner automorphisms and not containing $s$. Thus the intersection
of the compact neighborhoods of $e$ invariant under inner automorphisms
is $\{e\}$. Consequently every neighborhood of $e$ contains a compact
neighborhood invariant under inner automorphisms.

Bibliography: \cite{ref53,ref54,ref95,ref96,ref136,ref137}.

\BookSubsection{13.11}{Complements}

\ResultHeading{13.11.1}{13.11.1.}
Let $G$ be a locally compact group and $f\in L^1(G)$. If $f(s)\geq0$
for every $s\in G$, the norm of $f$ in $L^1(G)$ is the same as its norm
in $C^{*}(G)$. (Consider the trivial one-dimensional representation of
$G$.) \cite{ref27}.

\ResultHeading{13.11.2}{13.11.2.}
Let $G$ be a locally compact group, $\pi$ a continuous unitary
representation of $G$, $S$ a set of continuous unitary representations
of $G$, and $K$ the set of $\xi\in H_\pi$ such that the function
\[
 s\longmapsto(\pi(s)\xi\mid\xi)
\]
is, uniformly on each compact set, a limit of sums of functions of
positive type associated with $S$. Then $K$ is a closed vector subspace
of $H_\pi$ invariant under $\pi(G)$. \cite{ref33}.

\ResultHeading{13.11.3}{13.11.3.}
Let $G$ be a locally compact group, $\lambda$ the left regular
representation of $G$, and $\pi$ a continuous unitary representation of
$G$. Then
\[
 \lambda\otimes\pi\simeq(\dim\pi)\lambda.
\]
[Consider the isomorphism from $L^2(G)\otimes H_\pi$ onto
$L^2(G)\otimes H_\pi=L^2_{H_\pi}(G)$ which sends
$f\otimes\xi$ to the function $s\mapsto f(s)\pi(s^{-1})\xi$.]
\cite{ref31}.

\ResultHeading{13.11.4}{13.11.4.}
Let $G$ be the locally compact group of affine transformations of
$\mathbb R$, and let $G'$ be the closed normal subgroup of translations.
If $\pi$ is a continuous unitary representation of $G$ and $\chi$ is a
nontrivial character of $G'$, then $\pi|_{G'}$ does not contain $\chi$.
(If $\pi|_{G'}$ contained $\chi$, then, by the action of the inner
automorphisms of $G$ on $G'$, $\pi|_{G'}$ would contain all the nontrivial
characters of $G'$; the corresponding subspaces of $H_\pi$ would be
pairwise orthogonal.) It follows that a nontrivial character of $G'$ cannot
% Source PDF physical page 283
be extended to a continuous function of positive type on $G$. (Douady,
unpublished.)

\ResultHeading{13.11.5}{13.11.5.}
Let $G$ be a locally compact group and $\Delta$ its modular function. If a
continuous function of positive type on $G$ is integrable with respect to
\[
 \Delta(s)^{-1/2}\,ds,
\]
then it is square-integrable with respect to $ds$. \cite{ref50}.

\ResultHeading{13.11.6}{*13.11.6.}
Let $G$ be a locally compact group, and let $E$ be the set of linear
combinations of continuous functions of positive type. For
$\varphi\in E$, let $K_\varphi^{g}$ (respectively
$K_\varphi^{d}$) be the closed convex hull of the set of left
(respectively right) translates of $\varphi$ in the space of bounded
continuous complex-valued functions on $G$, endowed with the uniform
norm. Then $K_\varphi^{g}$ and $K_\varphi^{d}$ contain one and the
same constant $M(\varphi)$, and no other constant. For a continuous
function $\varphi$ of positive type to satisfy $M(\varphi)=0$, it is
necessary and sufficient that $\pi_\varphi$ have no nonzero
finite-dimensional subrepresentation. \cite{ref50}.

\ResultHeading{13.11.7}{13.11.7.}
Let $G_1,G_2$ be two topological groups.
\begin{enumerate}[label=\alph*.]
 \item Let $\pi_1$ be an irreducible continuous unitary representation
 of $G_1$, and let $\pi_2,\pi'_2$ be continuous unitary representations
 of $G_2$. If the representations
 \[
  (s_1,s_2)\longmapsto\pi_1(s_1)\otimes\pi_2(s_2)
  \quad\text{and}\quad
  (s_1,s_2)\longmapsto\pi_1(s_1)\otimes\pi'_2(s_2)
 \]
 are equivalent, then $\pi_2$ and $\pi'_2$ are equivalent.
 \item Let $\pi$ be a continuous unitary representation of
 $G_1\times G_2$. If $\pi$ is a factor representation, then
 $\pi|_{G_1}$ and $\pi|_{G_2}$ are factor representations.
 \item Let $\pi$ be a continuous unitary representation of
 $G_1\times G_2$, and put $\pi_1=\pi|_{G_1}$ and
 $\pi_2=\pi|_{G_2}$. Suppose that $\pi_1$ is a type I factor
 representation. Then there exist $\pi'_1\simeq\pi_1$ and
 $\pi'_2\simeq\pi_2$ such that $\pi$ is equivalent to the
 representation
 \[
  (s_1,s_2)\longmapsto\pi'_1(s_1)\otimes\pi'_2(s_2).
 \]
 \item If $G_1$ and $G_2$ are of type I, then $G_1\times G_2$ is of
 type I. \cite{ref89}
\end{enumerate}

\ResultHeading{13.11.8}{13.11.8.}
Let $G$ be a locally compact group and $\pi$ a unitary representation of
$G$. Suppose that, for all $\xi,\eta\in H_\pi$, the function
$s\mapsto(\pi(s)\xi\mid\eta)$ is measurable for Haar measure on
$G$. We may define
\[
 \pi(f)=\int f(s)\pi(s)\,ds
\]
for every $f\in L^1(G)$. Let $K\subset H_\pi$ be the essential
subspace for the representation $\pi$ of $L^1(G)$. Then $K$ and
$H_\pi\ominus K$ are invariant under $\pi(G)$. The representation
$s\mapsto\pi(s)|_K$ of $G$ is continuous. For
$\xi,\eta\in H_\pi\ominus K$, one has
\[
 (\pi(s)\xi\mid\eta)=0
\]
locally almost everywhere on $G$. The space $H_\pi\ominus K$ is zero
or nonseparable. It may be nonzero. \cite{ref140}.

\ResultHeading{13.11.9}{13.11.9.}
Let $G$ be a topological group and $\pi$ a continuous unitary
representation of $G$. In
\[
 H_\pi\otimes H_\pi\otimes\cdots\otimes H_\pi
 \qquad(n\text{ factors}),
\]
the symmetric tensors span a closed vector subspace $K$ invariant under
$\rho=\pi\otimes\pi\otimes\cdots\otimes\pi$. The
subrepresentation $\rho|_K$ of $\rho$ is called the $n$th symmetric
tensor power of $\pi$. The antisymmetric tensor powers are defined
similarly.

% Source PDF physical page 284
\ResultHeading{13.11.10}{*13.11.10.}
Let $G$ be a connected real Lie group and $\pi$ an irreducible unitary
representation of $G$ continuous for the norm topology on operators. Then
$\pi$ is finite-dimensional. \cite{ref143}.

\ResultHeading{13.11.11}{13.11.11.}
Let $G$ be a topological group. A unitary representation $\pi$ of $G$ is
called \emph{real} if there exists a closed real vector subspace $K$ of
$H_\pi$ such that $H_\pi$ is the direct sum of $K$ and $iK$, the scalar
product is real on $K$, and $\pi(G)$ leaves $K$ invariant (in other words,
$H_\pi$ is the complexification of a real Hilbert space $K$, and $\pi$
is the complexification of a representation of $G$ on $K$). Equivalently,
there exists a conjugation $J$ on $H_\pi$ which commutes with $\pi(G)$.
If $\pi$ is real, then $\pi\simeq\overline\pi$, but the converse is
false.

\ResultHeading{13.11.12}{*13.11.12.}
A connected real semisimple Lie group and a connected real nilpotent Lie
group are liminal (cf. \XRef{15.5.6}{15.5.6}). A real algebraic linear
group is postliminal. Fairly broad classes of postliminal solvable Lie
groups are known, but Mautner has given an example of a solvable Lie
group which is not postliminal. The noncommutative connected solvable real
Lie group of dimension $2$ is postliminal but not liminal. A countable
discrete group is of type I if and only if it is an extension of a finite
group by an abelian group. Problems: Is a $p$-adic algebraic group
postliminal? Is the $C^{*}$-algebra of a liminal real Lie group of
generalized continuous trace?
\cite{ref5,ref7,ref10,ref30,ref59,ref79,ref81,ref146,ref155,ref209}.

\ResultHeading{13.11.13}{*13.11.13.}
Let $G$ be a locally compact group and $\pi$ an irreducible continuous
unitary representation of $G$. Then $\pi(L^1(G))$ is not, in general,
algebraically irreducible.

\ResultHeading{13.11.14}{13.11.14.}
\begin{enumerate}[label=\alph*.]
 \item Let $G$ be a discrete group, $A$ its Hilbert algebra,
 $\mathfrak U=\mathfrak U(A)$, which is a finite von Neumann algebra, and
 $t$ the natural trace on $\mathfrak U^+$ defined by $A$. Since the
 operator $1\in\mathfrak U$ corresponds to the characteristic function of
 $e$, one has $\mathfrak m_t=\mathfrak n_t=\mathfrak U$, and every element
 of $\mathfrak U$ is of the form $U_x$, where $x\in L^2(G)$. We shall again
 denote by $t$ the linear extension of $t$ to $\mathfrak U$. If
 $U_x\in\mathfrak U$, with $x\in L^2(G)$, then
 \[
  t(U_x)=x(e).
 \]
 \item Let $E_{\mathrm I_i}$ (respectively $E_{\mathrm{II}}$) be the
 largest projection in the center of $\mathfrak U$ such that the
 corresponding induced algebra is of type $\mathrm I_i$ (respectively of
 type $\mathrm{II}$), where $i=1,2,\ldots$. One has
 \[
  E_{\mathrm{II}}+\sum_i E_{\mathrm I_i}=1.
 \]
 Put $t(E_{\mathrm{II}})=r$ and $t(E_{\mathrm I_i})=r_i$; then
 \[
  r+\sum_i r_i=1.
 \]
 \item Let $C$ be the commutator subgroup of $G$, and let
 $\varphi_C$ be its characteristic function. If $C$ is infinite, then
 $r_1=0$. If $C$ is finite, then
 \[
  r_1=(\operatorname{Card}C)^{-1},
 \]
 and $E_{\mathrm I_1}$ is the element of $\mathfrak U$ defined by the
 function $(\operatorname{Card}C)^{-1}\varphi_C$ on $G$.
 \item If $G$ is infinite and $C$ coincides with the center of $G$ and
 has prime order $p$, then $\mathfrak U$ is the product of an abelian von
 Neumann algebra and $p-1$ factors of type $\mathrm{II}_1$. One has
 \[
  r_1=\frac1p,\qquad r=\frac{p-1}{p}.
 \]

% Source PDF physical page 285
 \item Let $(G_i)$ be an infinite family of noncommutative finite groups.
 Let $G$ be the set of elements of $\prod_iG_i$ all but finitely many of
 whose components are equal to $e$, and regard $G$ as a discrete group.
 Then $\mathfrak U$ is of type $\mathrm{II}_1$.
 \item Let $G$ be a discrete group, and let $G_0$ be the subgroup of
 $G$ which is the union of the finite conjugacy classes of $G$. If
 $G/G_0$ is infinite, then $\mathfrak U$ is of type $\mathrm{II}_1$. If
 $G_0=\{e\}$ and $G$ is infinite, then $\mathfrak U$ is a factor of type
 $\mathrm{II}_1$. \cite{ref76,ref90,ref98,ref99}
\end{enumerate}

\BookSection{14}{Irreducible Square-Integrable Representations}

Throughout this section, $G$ denotes a unimodular locally compact group.

\BookSubsection{14.1}{Definition of Square-Integrable Representations}

\ResultHeading{14.1.1}{14.1.1. Lemma.}
\begin{itshape}
Let $\varphi$ be a continuous square-integrable function of positive type
on $G$. The representation $\pi_\varphi$ is contained in the left regular
representation.
\end{itshape}

There exists $\psi\in L^2(G)$ such that
$\varphi=\overline\psi\ast\check\psi$
(\XRef{13.8.6}{13.8.6}). In other words,
\[
 \varphi(s)=\int\overline{\psi(t)}\psi(s^{-1}t)\,dt
 =(\lambda(s)\psi\mid\psi),
\]
where $\lambda$ denotes the left regular representation of $G$. Hence
$\pi_\varphi$ is equivalent to the subrepresentation of $\lambda$ which
has $\psi$ as a cyclic vector [\XRef{13.4.5}{13.4.5(iii)}].

\ResultHeading{14.1.2}{14.1.2. Lemma.}
\begin{itshape}
Let $\lambda$ be the left regular representation of $G$, and let $\pi$ be
a representation contained in $\lambda$, with $H_\pi\ne0$. There exists
a nonzero square-integrable function of positive type associated with
$\pi$.
\end{itshape}

Let $A$ be the Hilbert algebra of $G$, and let $K$ be a closed vector
subspace of $L^2(G)$ invariant under $\lambda(G)$. Then
\[
 P_K\in\mathfrak U(A)'=\mathfrak V(A).
\]
Consequently $P_K(A)\subset A$ (\XRef{A.59}{A 59}). If $K\ne0$, there
is therefore a nonzero $\psi\in K\cap A$. The subrepresentation of
$\lambda$ corresponding to $K$ has as an associated function of positive
type
\[
 s\longmapsto(\pi(s)\psi\mid\psi)
 =\overline\psi\ast\check\psi.
\]
Now $\overline\psi\in A$, and hence
$\overline\psi\ast\check\psi\in L^2(G)$.

\ResultHeading{14.1.3}{14.1.3. Definition.}
\begin{itshape}
An irreducible continuous unitary representation $\pi$ of $G$ is called
square-integrable if there exists $\xi\in H_\pi$, $\xi\ne0$, such that
the coefficient
\[
 s\longmapsto(\pi(s)\xi\mid\xi)
\]
of $\pi$ is square-integrable on $G$.
\end{itshape}

[We shall see that in this case all the coefficients of $\pi$ are
square-integrable (\XRef{14.3.2}{14.3.2}).]

By \XRef{14.1.1}{14.1.1} and \XRef{14.1.2}{14.1.2}, these
representations are precisely the irreducible representations contained
in the regular representation.

% Source PDF physical page 286
\ResultHeading{14.1.4}{14.1.4.}
Suppose that $G$ is commutative. An element $\chi$ of $\widehat G$ is a
character of $G$, and $|\chi|^2=1$. Thus, if $G$ is compact, every
element of $\widehat G$ is square-integrable (cf. \SectionRef{15}{\S15}); and
if $G$ is not compact, no element of $\widehat G$ is square-integrable.

Bibliography: \cite{ref51,ref60}.

\BookSubsection[Square-Integrable Representations and Minimal Bi-Invariant Subspaces of L2(G)]{14.2}
  {Square-Inte\-grable Rep\-re\-sen\-ta\-tions and Minimal
   Bi-In\-var\-i\-ant Sub\-spaces of $L^2(G)$}

\ResultHeading{14.2.1}{14.2.1.}
Let $A$ be the Hilbert algebra of $G$, and let $K$ be a closed vector
subspace of $L^2(G)$. To say that $K$ is invariant under the left
(respectively right) regular representation of $G$ is equivalent to
saying that
\[
 P_K\in\mathfrak V(A)
 \qquad(\text{respectively }P_K\in\mathfrak U(A)).
\]
To say that $K$ is invariant under both the left and right regular
representations, or, as we shall say, \emph{bi-invariant}, is equivalent
to saying that $P_K$ belongs to
$\mathfrak U(A)\cap\mathfrak V(A)$, the common center of
$\mathfrak U(A)$ and $\mathfrak V(A)$. We shall establish a bijective
correspondence between the classes of square-integrable irreducible
representations of $G$ and certain minimal bi-invariant subspaces of
$L^2(G)$, in other words, certain minimal projections of
$\mathfrak U(A)\cap\mathfrak V(A)$.

\ResultHeading{14.2.2}{14.2.2. Proposition.}
\begin{itshape}
Let $\lambda$ be the left regular representation of $G$, and let $H$ be a
closed vector subspace of $L^2(G)$ minimally invariant under $\lambda$, so
that the subrepresentation $\sigma$ of $\lambda$ on $H$ is irreducible.
Let $E=P_H\in\mathfrak V(A)$, and let $F$ be the central support of $E$ in
$\mathfrak U(A)\cap\mathfrak V(A)$.
\begin{enumerate}[label=(\roman*)]
 \item $F$ is a minimal projection of
 $\mathfrak U(A)\cap\mathfrak V(A)$.
 \item $\mathfrak U(A)_F$ and $\mathfrak V(A)_F$ are type I factors
 commuting with one another.
 \item Every element of $F(L^2(G))$ belongs to $A$.
\end{enumerate}
\end{itshape}

Assertion (i) follows from \XRef{5.2.7}{5.2.7}. By condition (v) of
\XRef{5.3.1}{5.3.1}, the subrepresentation of $\lambda$ defined by $F$
is quasi-equivalent to $\sigma$, and is therefore a type I factor
representation (\XRef{5.4.11}{5.4.11}), which proves (ii). Finally,
(iii) follows from \XRef{A.65}{A 65} and \XRef{A.67}{A 67}.

\ResultHeading{14.2.3}{14.2.3. Proposition.}
\begin{itshape}
Retain the notation of \XRef{14.2.2}{14.2.2}. Let $\sigma'$ be another
irreducible subrepresentation of $\lambda$, and let $F'$ be the
corresponding minimal projection of
$\mathfrak U(A)\cap\mathfrak V(A)$. The representations $\sigma$ and
$\sigma'$ are equivalent (respectively inequivalent) if and only if $F$
and $F'$ are equal (respectively orthogonal).
\end{itshape}

If $F=F'$, then $\sigma$ and $\sigma'$ are quasi-equivalent [condition
(v) of \XRef{5.3.1}{5.3.1}], and hence equivalent
[\XRef{5.3.3}{5.3.3(ii)}]. If $F$ and $F'$ are orthogonal, then
$\sigma$ and $\sigma'$ are disjoint [condition (iii) of
\XRef{5.2.1}{5.2.1}].

\ResultHeading{14.2.4}{14.2.4.}
We call $F$ [respectively $F(L^2(G))$] the central projection
(respectively the bi-invariant subspace) associated with $\sigma$, or with
any equivalent representation.

% Source PDF physical page 287
\ResultHeading{14.2.5}{14.2.5. Proposition.}
\begin{itshape}
Let $F$ be a minimal projection of
$\mathfrak U(A)\cap\mathfrak V(A)$ such that $\mathfrak U(A)_F$ and
$\mathfrak V(A)_F$, which are factors, are of type I. There exists at
least one minimal projection $E$ of $\mathfrak V(A)$ dominated by $F$.
The subrepresentation $\sigma$ of $\lambda$ defined by $E$ is irreducible
and square-integrable; its class does not depend on the choice of $E$,
and the associated central projection is $F$.
\end{itshape}

Every type I factor has minimal projections (\XRef{A.36}{A 36}), which
proves the existence of $E$. Since $\sigma$ is an irreducible
subrepresentation of $\lambda$, it is square-integrable
(\XRef{14.1.3}{14.1.3}); its class does not depend on the choice of $E$
(\XRef{14.2.3}{14.2.3}); and its associated central projection is
evidently $F$.

Bibliography: \cite{ref51,ref60}.

\BookSubsection{14.3}{Coefficients of Square-Integrable Representations}

\ResultHeading{14.3.1}{14.3.1. Proposition.}
\begin{itshape}
Let $\sigma$ be an irreducible subrepresentation of the left regular
representation of $G$. For $\xi,\eta\in H_\sigma$, let
\[
 s\longmapsto\varphi^\sigma_{\xi,\eta}(s)
 =(\sigma(s)\xi\mid\eta)
\]
be the corresponding coefficient of $\sigma$. Then
\[
 \check\varphi^\sigma_{\xi,\eta}
 =\xi\ast\eta^{\ast}\in L^2(G).
\]
The closed vector subspace of $L^2(G)$ spanned by the
$\check\varphi^\sigma_{\xi,\eta}$, $\xi,\eta\in H_\sigma$, is the
bi-invariant subspace associated with $\sigma$.
\end{itshape}

Indeed,
\begin{align*}
 \check\varphi^\sigma_{\xi,\eta}(s)
 &=(\xi\mid\sigma(s)\eta)
 =\int\xi(t)\overline{\eta(s^{-1}t)}\,dt\\
 &=\int\xi(t)\eta^{\ast}(t^{-1}s)\,dt
 =(\xi\ast\eta^{\ast})(s).
\end{align*}
Let $K$ be the bi-invariant subspace associated with $\sigma$. Then
$\xi$ and $\eta$, and hence $\eta^{\ast}$, belong to $K$, and therefore
to the Hilbert algebra $A$ of $G$ [\XRef{14.2.2}{14.2.2(iii)}]. Thus
\[
 \xi\ast\eta^{\ast}=V_{\eta^{\ast}}\xi\in K.
\]
On the other hand, since $H_\sigma$ is invariant under
$\mathfrak U(A)$, the set of the $\xi\ast\eta^{\ast}$,
$\xi,\eta\in H_\sigma$, is stable under left or right multiplication by
elements of $A$. If $K'$ denotes the closed vector subspace spanned by
this set, then $K'$ is a closed bi-invariant vector subspace of $L^2(G)$
contained in $K$. Since $K$ is a minimal closed bi-invariant vector
subspace [\XRef{14.2.2}{14.2.2(i)}], one has $K'=K$.

\ResultHeading{14.3.2}{14.3.2. Corollary.}
\begin{itshape}
Let $\sigma$ be a square-integrable irreducible representation of $G$.
All the coefficients of $\sigma$ are square-integrable.
\end{itshape}

\ResultHeading{14.3.3}{14.3.3. Theorem.}
\begin{itshape}
Let $\sigma$ be a square-integrable irreducible representation of $G$.
For $\xi,\eta\in H_\sigma$, put
$\varphi_{\xi,\eta}(s)=(\sigma(s)\xi\mid\eta)$, $s\in G$. There exists a
unique constant $d_\sigma$, $0<d_\sigma<+\infty$, such that
\begin{align}
 \int\varphi_{\xi,\eta}(s)
  \overline{\varphi_{\xi',\eta'}(s)}\,ds
 &=d_\sigma^{-1}(\xi\mid\xi')
   \overline{(\eta\mid\eta')},
 \tag{1}\label{eq:14.3.3-1}\\
 \varphi_{\xi,\eta}\ast
 \varphi_{\xi',\eta'}
 &=d_\sigma^{-1}(\xi\mid\eta')
   \varphi_{\xi',\eta},
\tag{2}\label{eq:14.3.3-2}
\end{align}
% Source PDF physical page 288
for all $\xi,\eta,\xi',\eta'\in H_\sigma$. In particular, if
$\varphi$ is a function of positive type associated with $\sigma$ and
equal to $1$ at $e$, then
\begin{align}
 \|\varphi\|_2^2&=d_\sigma^{-1},
 \tag{3}\label{eq:14.3.3-3}\\
 \varphi\ast\varphi
 &=d_\sigma^{-1}\varphi.
 \tag{4}\label{eq:14.3.3-4}
\end{align}
\end{itshape}

Let $A$ be the Hilbert algebra of $G$. We may suppose that $\sigma$ is the
subrepresentation of the left regular representation defined by a minimal
projection $E$ of $\mathfrak V(A)$. Using the equality
$\check\varphi_{\xi,\eta}=\xi\ast\eta^{\ast}
=V_{\eta^{\ast}}\xi$ of \XRef{14.3.1}{14.3.1}, we have
\begin{align*}
 \int\varphi_{\xi,\eta}(s)
  \overline{\varphi_{\xi',\eta'}(s)}\,ds
 &=\int\check\varphi_{\xi,\eta}(s^{-1})
  \overline{\check\varphi_{\xi',\eta'}(s^{-1})}\,ds\\
 &=(V_{\eta^{\ast}}\xi\mid V_{\eta'^{\ast}}\xi')
  =(V_{\eta'}V_{\eta^{\ast}}\xi\mid\xi')\\
 &=(EV_{\eta'}V_{\eta^{\ast}}E\xi\mid\xi').
\end{align*}
Since $E$ is a minimal projection of $\mathfrak V(A)$, there exists
$\lambda_{\eta,\eta'}\in\mathbb C$ such that
\[
 EV_{\eta'}V_{\eta^{\ast}}E=\lambda_{\eta,\eta'}E.
\]
Consequently
\[
 \int\varphi_{\xi,\eta}(s)
  \overline{\varphi_{\xi',\eta'}(s)}\,ds
 =\lambda_{\eta,\eta'}(\xi\mid\xi').
\]
By \XRef{13.1.1}{13.1.1}, the integral is also equal to
\[
 \int\varphi_{\eta',\xi'}(s^{-1})
  \overline{\varphi_{\eta,\xi}(s^{-1})}\,ds
 =\lambda_{\xi',\xi}(\eta'\mid\eta).
\]
Thus
$\lambda_{\eta,\eta'}=\lambda_\sigma(\eta'\mid\eta)$, where
$\lambda_\sigma$ is a finite constant. Taking
$\xi=\xi'=\eta=\eta'\ne0$ shows that $\lambda_\sigma>0$. It is then
enough to put $\lambda_\sigma^{-1}=d_\sigma$ to obtain
\eqref{eq:14.3.3-1}.

In particular,
\begin{align*}
 (\varphi_{\xi,\eta}\ast
  \varphi_{\xi',\eta'})(s)
 &=\int(\sigma(t)\xi\mid\eta)
   (\sigma(t^{-1}s)\xi'\mid\eta')\,dt\\
 &=\int(\sigma(t)\xi\mid\eta)
   \overline{(\sigma(t)\eta'\mid\sigma(s)\xi')}\,dt\\
 &=d_\sigma^{-1}(\xi\mid\eta')
   (\sigma(s)\xi'\mid\eta)
 =d_\sigma^{-1}(\xi\mid\eta')
  \varphi_{\xi',\eta}(s).
\end{align*}

\ResultHeading{14.3.4}{14.3.4. Definition.}
\begin{itshape}
The constant $d_\sigma$ is called the formal dimension of $\sigma$.
\end{itshape}

Its value depends on the choice of Haar measure on $G$. We shall see in
\XRef{15.2.3}{15.2.3} that, if $G$ is compact and the chosen Haar
measure has total mass $1$, then $d_\sigma$ is the dimension of $\sigma$
in the usual sense. This explains the terminology adopted.

\ResultHeading{14.3.5}{14.3.5. Proposition.}
\begin{itshape}
Let $\lambda,\rho$ be the left and right regular representations of $G$,
let $\sigma$ be an irreducible subrepresentation of $\lambda$, let
$K\supset H_\sigma$ be the bi-invariant subspace of $L^2(G)$ associated
with $\sigma$, let $\lambda_1,\rho_1$ be the subrepresentations of
$\lambda,\rho$ defined by $K$, and let $\overline H_\sigma$ be the
conjugate Hilbert space of $H_\sigma$.

% Source PDF physical page 289
\begin{enumerate}[label=(\roman*)]
 \item There exists one and only one isomorphism $\Phi$ from the Hilbert
 space $H_\sigma\otimes\overline H_\sigma$ onto $K$ which sends
 $\xi\otimes\overline\eta$ to
 $d_\sigma^{1/2}\check\varphi_{\xi,\eta}$ for every
 $\xi,\eta\in H_\sigma$.
 \item Let $1_{H_\sigma}$ and $1_{\overline H_\sigma}$ be the trivial
 representations of $G$ on $H_\sigma$ and $\overline H_\sigma$.
 Then $\Phi$ transforms $\sigma\otimes1_{\overline H_\sigma}$ into
 $\lambda_1$, and $1_{H_\sigma}\otimes\overline\sigma$ into $\rho_1$.
 \item $\lambda_1\simeq(\dim\sigma)\sigma$ and
 $\rho_1\simeq(\dim\sigma)\overline\sigma$.
\end{enumerate}
\end{itshape}

The map
\[
 (\xi,\eta)\longmapsto
 d_\sigma^{1/2}\check\varphi_{\xi,\eta}
\]
is linear in $\xi$ and conjugate-linear in $\eta$. Hence there is a
linear map from the algebraic tensor product of $H_\sigma$ and
$\overline H_\sigma$ onto a dense vector subspace of $K$ which sends
$\xi\otimes\overline\eta$ to
$d_\sigma^{1/2}\check\varphi_{\xi,\eta}$
(\XRef{14.3.1}{14.3.1}). If
$\xi,\xi',\eta,\eta'\in H_\sigma$, then, by
\XRef{14.3.3}{14.3.3},
\[
 (\xi\otimes\overline\eta\mid
   \xi'\otimes\overline{\eta'})
 =(\xi\mid\xi')\overline{(\eta\mid\eta')}
 =\bigl(d_\sigma^{1/2}\check\varphi_{\xi,\eta}
  \bigm|
  d_\sigma^{1/2}\check\varphi_{\xi',\eta'}\bigr).
\]
Thus the preceding linear map extends to an isomorphism $\Phi$ from
$H_\sigma\otimes\overline H_\sigma$ onto $K$. The uniqueness in (i) is
immediate, since the $\xi\otimes\overline\eta$ form a total set in
$H_\sigma\otimes\overline H_\sigma$. For
$s\in G$ and $\xi,\eta\in H_\sigma$, one has
\[
 \Phi(\sigma(s)\xi\otimes\overline\eta)
 =\check\varphi_{\sigma(s)\xi,\eta}
\]
and
\begin{align*}
 \check\varphi_{\sigma(s)\xi,\eta}(t)
 &=(\sigma(t^{-1})\sigma(s)\xi\mid\eta)
  =(\sigma(t^{-1}s)\xi\mid\eta)\\
 &=\check\varphi_{\xi,\eta}(s^{-1}t)
  =(\lambda(s)\check\varphi_{\xi,\eta})(t).
\end{align*}
Thus $\Phi$ transforms $\sigma\otimes1_{\overline H_\sigma}$ into
$\lambda_1$; similarly, it transforms
$1_{H_\sigma}\otimes\overline\sigma$ into $\rho_1$. Finally, (iii)
follows from (ii).

\ResultHeading{14.3.6}{14.3.6. Proposition.}
\begin{itshape}
Retain the notation of \XRef{14.3.5}{14.3.5}. If
$H_\sigma\otimes\overline H_\sigma$ is canonically identified with the
set of Hilbert--Schmidt operators on $H_\sigma$, then
\[
 \Phi(u^{\ast})=\Phi(u)^{\ast},\qquad
 \Phi(uv)=d_\sigma^{1/2}\Phi(u)\ast\Phi(v)
\]
for all $u,v\in H_\sigma\otimes\overline H_\sigma$.
\end{itshape}

It is enough to verify this for
$u=\xi\otimes\overline\eta$ and
$v=\xi'\otimes\overline{\eta'}$, where
$\xi,\eta,\xi',\eta'\in H_\sigma$. Now
\[
 \Phi((\xi\otimes\overline\eta)^{\ast})
 =\Phi(\eta\otimes\overline\xi)
 =d_\sigma^{1/2}\check\varphi_{\eta,\xi}
 =\bigl(d_\sigma^{1/2}\check\varphi_{\xi,\eta}\bigr)^{\ast},
\]
and
\begin{align*}
 \Phi\bigl((\xi\otimes\overline\eta)
  (\xi'\otimes\overline{\eta'})\bigr)
 &=\Phi\bigl((\xi'\mid\eta)
    \xi\otimes\overline{\eta'}\bigr)\\
 &=d_\sigma^{1/2}(\xi'\mid\eta)
   \check\varphi_{\xi,\eta'}\\
 &=d_\sigma^{1/2}d_\sigma
   \check\varphi_{\xi,\eta}\ast
   \check\varphi_{\xi',\eta'}\\
 &=d_\sigma^{1/2}\Phi(\xi\otimes\overline\eta)
   \ast\Phi(\xi'\otimes\overline{\eta'}).
\end{align*}

% Source PDF physical page 290
\ResultHeading{14.3.7}{14.3.7. Theorem.}
\begin{itshape}
Let $\sigma,\sigma'$ be inequivalent square-integrable irreducible
representations of $G$. For $\xi,\eta\in H_\sigma$ and
$\xi',\eta'\in H_{\sigma'}$, one has
\[
 \int\varphi^\sigma_{\xi,\eta}(s)
  \overline{\varphi^{\sigma'}_{\xi',\eta'}(s)}\,ds=0,
 \qquad
 \check\varphi^\sigma_{\xi,\eta}\ast
 \check\varphi^{\sigma'}_{\xi',\eta'}=0.
\]
\end{itshape}

The bi-invariant subspaces $K$ and $K'$ of $L^2(G)$ associated with
$\sigma$ and $\sigma'$ are orthogonal (\XRef{14.2.3}{14.2.3}). Moreover,
$\check\varphi^\sigma_{\xi,\eta}\in K$ and
$\check\varphi^{\sigma'}_{\xi',\eta'}\in K'$
(\XRef{14.3.1}{14.3.1}), whence
\[
 (\varphi^\sigma_{\xi,\eta}\mid
   \varphi^{\sigma'}_{\xi',\eta'})
 =(\check\varphi^\sigma_{\xi,\eta}\mid
   \check\varphi^{\sigma'}_{\xi',\eta'})=0.
\]
Finally, $K$ and $K'$ are invariant under $\mathfrak U(A)$ and
$\mathfrak V(A)$, where $A$ denotes the Hilbert algebra of $G$; hence
\[
 \check\varphi^{\sigma'}_{\xi',\eta'}\ast
 \check\varphi^\sigma_{\xi,\eta}\in K\cap K'=0,
 \qquad
 \check\varphi^\sigma_{\xi,\eta}\ast
 \check\varphi^{\sigma'}_{\xi',\eta'}=0.
\]

Bibliography: \cite{ref51,ref60}.

\BookSubsection{14.4}{Formal Dimension and Trace}

\ResultHeading{14.4.1}{14.4.1. Lemma.}
\begin{itshape}
Let $H$ be a Hilbert space, $B$ the Hilbert algebra of Hilbert--Schmidt
operators on $H$, $b\mapsto U_b$ the canonical map from $B$ into
$\mathfrak U(B)$, $E$ a minimal projection of $\mathfrak V(B)$, and
$\sigma(b)$ the operator induced by $U_b$ on $E(B)$. For every $b\in B$,
\[
 \|b\|^2=\operatorname{Tr}\bigl(\sigma(b)\sigma(b)^{\ast}\bigr).
\]
\end{itshape}

The algebra $\mathfrak U(B)$ is a type I factor, and $b\mapsto U_b$ is
an isomorphism from the Hilbert algebra $B$ onto the Hilbert algebra of
Hilbert--Schmidt operators relative to $\mathfrak U(B)$
(\XRef{A.61}{A 61}). On the other hand, $T\mapsto T_E$ is an isomorphism
from $\mathfrak U(B)$ onto $\mathcal L(E(B))$. Hence
$b\mapsto\sigma(b)$ is an isomorphism from the Hilbert algebra $B$ onto
the Hilbert algebra of Hilbert--Schmidt operators on $E(B)$. Therefore
\[
 \operatorname{Tr}\bigl(\sigma(b)\sigma(b)^{\ast}\bigr)
 =\operatorname{Tr}(bb^{\ast})=\|b\|^2.
\]

\ResultHeading{14.4.2}{14.4.2. Proposition.}
\begin{itshape}
Let $\sigma$ be a square-integrable irreducible representation of $G$,
and let $F$ be the central projection associated with $\sigma$. For every
$f\in L^1(G)\cap L^2(G)$, one has
\[
 \int |(Ff)(s)|^2\,ds
 =d_\sigma\operatorname{Tr}\bigl(\sigma(f)\sigma(f^{\ast})\bigr).
\]
\end{itshape}

Let $A$ be the Hilbert algebra of $G$, and put $K=F(L^2(G))$. Then
$P_K(A)=A\cap K=K$ (\XRef{14.2.2}{14.2.2}), and $K$ is a Hilbert
algebra whose associated von Neumann algebras are
$\mathfrak U(A)_F$ and $\mathfrak V(A)_F$. Let $E$ be a minimal
projection of $\mathfrak V(A)$ dominated by $F$, so that $\sigma$ is
identified with the subrepresentation of the left regular representation
defined by $E$. Every $f\in A$ defines an operator
$U_f\in\mathfrak U(A)$, and $U_f$ induces on
% Source PDF physical page 291
$H_\sigma=E(L^2(G))$ an operator which we shall denote by $\sigma(f)$;
this agrees with the usual notation if $f\in L^1(G)\cap L^2(G)$. In
fact, we shall prove the equality in the proposition for every $f\in A$.
One has $U_{Ff}=FU_f$, and hence $\sigma(Ff)=\sigma(f)$. Neither side of
the equality to be proved is therefore changed if $f$ is replaced by
$Ff$. We shall henceforth suppose that $f\in K$.

Resume the notation of \XRef{14.3.5}{14.3.5}. There exists
$u\in H_\sigma\otimes\overline H_\sigma$ such that $f=\Phi(u)$, and
$\|f\|_2=\|u\|$. Moreover, if
$v\in H_\sigma\otimes\overline H_\sigma$, then
\[
 \Phi(u)\ast\Phi(v)=d_\sigma^{-1/2}\Phi(uv)
 \qquad\text{(\XRef{14.3.6}{14.3.6})}.
\]
Thus $\Phi^{-1}$ transforms left multiplication by $f$ on $K$ into left
multiplication by $d_\sigma^{-1/2}u$ on
$H_\sigma\otimes\overline H_\sigma$. In view of
\XRef{14.4.1}{14.4.1}, we have
\[
 \operatorname{Tr}\bigl(\sigma(f)\sigma(f)^{\ast}\bigr)
 =\|d_\sigma^{-1/2}u\|^2
 =d_\sigma^{-1}\|u\|^2
 =d_\sigma^{-1}\|f\|_2^2.
\]

\ResultHeading{14.4.3}{14.4.3. Proposition.}
\begin{itshape}
Let $f,f'\in L^1(G)\cap L^2(G)$ and put $g=f\ast f'$. For
$t\in G$, let $g^t$ denote the function $s\mapsto g(tst^{-1})$ on $G$.
Let $\sigma$ be a square-integrable irreducible representation of $G$.
\begin{enumerate}[label=(\roman*)]
 \item The operator $\sigma(g)$ is of trace class.
 \item For every unit vector $\xi$ of $H_\sigma$, the function
 $t\mapsto(\sigma(g^t)\xi\mid\xi)$ on $G$ is integrable, and
 \[
  \int(\sigma(g^t)\xi\mid\xi)\,dt
  =d_\sigma^{-1}\operatorname{Tr}\sigma(g).
 \]
\end{enumerate}
\end{itshape}

By linearity, it is enough to consider the case $f'=f^{\ast}$. Since
$\sigma(f)$ is a Hilbert--Schmidt operator (\XRef{14.4.2}{14.4.2}),
$\sigma(g)=\sigma(f)\sigma(f)^{\ast}$ is a positive trace-class operator.
Let $(\xi_i)$ be an orthonormal basis of $H_\sigma$ consisting of
eigenvectors of $\sigma(g)$. Write
$\sigma(g)\xi_i=\lambda_i\xi_i$, so that
$\sum_i\lambda_i<+\infty$. For every $t\in G$,
\begin{align*}
 (\sigma(g^t)\xi\mid\xi)
 &=(\sigma(t^{-1})\sigma(g)\sigma(t)\xi\mid\xi)
  =(\sigma(g)\sigma(t)\xi\mid\sigma(t)\xi)\\
 &=\sum_i(\sigma(g)\sigma(t)\xi\mid\xi_i)
   \overline{(\sigma(t)\xi\mid\xi_i)}
  =\sum_i\lambda_i|(\sigma(t)\xi\mid\xi_i)|^2.
\end{align*}
Consequently, by \XRef{14.3.3}{14.3.3},
\[
 \int(\sigma(g^t)\xi\mid\xi)\,dt
 =\sum_i\lambda_i\int|(\sigma(t)\xi\mid\xi_i)|^2\,dt
 =\sum_i\lambda_i d_\sigma^{-1}
 =d_\sigma^{-1}\operatorname{Tr}\sigma(g).
\]

Bibliography: \cite{ref60}.

% Source PDF physical page 292
\BookSubsection{14.5}{Integrable Representations}

\ResultHeading{14.5.1}{14.5.1. Proposition.}
\begin{itshape}
Let $\sigma$ be an irreducible continuous unitary representation of $G$.
The following conditions are equivalent:
\begin{enumerate}[label=(\roman*)]
 \item There exists $\xi\in H_\sigma$, $\xi\ne0$, such that the function
 $s\mapsto(\sigma(s)\xi\mid\xi)$ is integrable on $G$.
 \item There exists a dense vector subspace $H'$ of $H_\sigma$ such that,
 for $\eta,\zeta\in H'$, the function
 $s\mapsto(\sigma(s)\eta\mid\zeta)$ is integrable on $G$.
\end{enumerate}
\end{itshape}

Suppose that (i) holds. Let $g,h\in\mathcal K(G)$. For every $s\in G$,
\begin{align*}
 (\sigma(g)\xi\mid\sigma(s)\sigma(h)\xi)
 &=\iint g(t)\overline{h(u)}
   (\sigma(t)\xi\mid\sigma(s)\sigma(u)\xi)\,dt\,du\\
 &=\iint g(t)\overline{h(u)}
   (\xi\mid\sigma(t^{-1}su)\xi)\,dt\,du.
\end{align*}
Hence
\begin{align*}
 &\int|(\sigma(g)\xi\mid\sigma(s)\sigma(h)\xi)|\,ds\\
 &\quad\leq
 \iiint|g(t)h(u)|\,|(\xi\mid\sigma(t^{-1}su)\xi)|\,ds\,dt\,du\\
 &\quad=
 \iiint|g(t)h(u)|\,|(\xi\mid\sigma(s)\xi)|\,ds\,dt\,du\\
 &\quad=
 \int|g(t)|\,dt\int|h(u)|\,du
 \int|(\xi\mid\sigma(s)\xi)|\,ds<+\infty.
\end{align*}
Thus (ii) holds with $H'=\sigma(\mathcal K(G))\xi$. [This subspace is
nonzero and invariant under $\sigma(G)$, and is therefore dense in
$H_\sigma$, since $\sigma$ is irreducible.]

\ResultHeading{14.5.2}{14.5.2. Definition.}
\begin{itshape}
An irreducible continuous unitary representation of $G$ is called
integrable if it satisfies the equivalent conditions of
\XRef{14.5.1}{14.5.1}.
\end{itshape}

With the notation of \XRef{14.5.1}{14.5.1}, the function
$s\mapsto(\sigma(s)\xi\mid\xi)$ on $G$, being bounded and integrable, is
square-integrable. Thus an integrable irreducible representation is
square-integrable. The converse is false.

Bibliography: \cite{ref60}.

\BookSubsection{14.6}{Complements}

\ResultHeading{14.6.1}{14.6.1.}
Let $G$ be a locally compact group, $\pi$ a continuous unitary
representation of $G$, and $p,q\in[1,+\infty]$ such that
$p^{-1}+q^{-1}=1$. The following conditions are equivalent:
\begin{enumerate}[label=\alph*.]
 \item There exists a finite constant $M$ such that
 \[
  \|\pi(f)\|\leq M\|f\|_p
  \qquad\text{for }f\in L^1(G)\cap L^p(G),
 \]
% Source PDF physical page 293
 where $\|f\|_p$ denotes the norm in $L^p(G)$.
 \item For all $\xi,\eta\in H_\pi$, the function
 $s\mapsto(\pi(s)\xi\mid\eta)$ belongs to $L^q(G)$.
\end{enumerate}
\cite{ref84}.

\ResultHeading{14.6.2}{14.6.2.}
Let $G$ be a unimodular locally compact group, let $\sigma$ be a
square-integrable irreducible continuous unitary representation of $G$,
and let $K_\sigma$ be the associated bi-invariant subspace. Then
$\overline\sigma$ is square-integrable and irreducible, and its associated
bi-invariant subspace is the image of $K_\sigma$ under the map
$f\mapsto\overline f$, and also under the map $f\mapsto\check f$ from
$L^2(G)$ onto $L^2(G)$.

\ResultHeading{14.6.3}{*14.6.3.}
The square-integrable irreducible continuous unitary representations of
connected semisimple real Lie groups are beginning to be well understood.
Connected complex semisimple Lie groups and simply connected nilpotent
real Lie groups have no square-integrable irreducible continuous unitary
representations. \cite{ref9,ref60,ref61}.

\BookSection{15}{Representations of Compact Groups}

Throughout this section, $G$ denotes a compact group. The Haar measure
chosen on $G$ will always be the one of total mass $1$. One has
\[
 L^1(G)\supset L^2(G),\qquad
 \|f\|_2\geq\|f\|_1\quad\text{for every }f\in L^2(G).
\]

\BookSubsection{15.1}{Complete Reducibility}

\ResultHeading{15.1.1}{15.1.1. Lemma.}
\begin{itshape}
There exists in $G$ a fundamental system of neighborhoods of $e$ invariant
under inner automorphisms.
\end{itshape}

It is enough to show that, if $s\in G$ and $s\ne e$, there exists a
compact neighborhood $V$ of $e$, invariant under inner automorphisms,
such that $s\notin V$. Let $W$ be a compact neighborhood of $s$ such that
$e\notin W$. The image of $G\times W$ under the map
\[
 (t,w)\longmapsto twt^{-1}
\]
is a compact neighborhood $W'$ of $s$, invariant under inner
automorphisms, such that $e\notin W'$. The complement of $W'$ is therefore
an open neighborhood of $e$, invariant under inner automorphisms, whose
closure does not contain $s$.

\ResultHeading{15.1.2}{15.1.2. Lemma.}
\begin{itshape}
The Hilbert algebra $A$ of $G$ is equal to $L^2(G)$. Let $(F_i)$ be the
family of pairwise orthogonal minimal projections of
$\mathfrak U(A)\cap\mathfrak V(A)$. Then $A$ is the Hilbert sum of the
$A_i=F_i(A)$; the $A_i$ are pairwise annihilating self-adjoint two-sided
ideals of $A$; and each $A_i$ is isomorphic to the Hilbert algebra of
Hilbert--Schmidt operators on a finite-dimensional Hilbert space.
\end{itshape}

One has
\[
 A\supset L^1(G)\cap L^2(G)=L^2(G),
\qquad\text{and hence }A=L^2(G).
\]
By \XRef{15.1.1}{15.1.1} and \XRef{13.10.5}{13.10.5},
$\mathfrak U(A)$ and $\mathfrak V(A)$ are finite von Neumann algebras.
The lemma now follows from \XRef{A.65}{A 65} and \XRef{A.68}{A 68}.

% Source PDF physical page 294
\ResultHeading{15.1.3}{15.1.3. Theorem.}
\begin{itshape}
Let $\pi$ be a continuous unitary representation of $G$. Then $\pi$ is a
Hilbert sum of finite-dimensional irreducible representations.
\end{itshape}

Retain the notation of \XRef{15.1.2}{15.1.2}. Let $H_i$ be the closed
vector subspace of $H_\pi$ spanned by the
$\pi(f)\xi$, $f\in A_i$, $\xi\in H_\pi$. Since $A=L^2(G)$ is dense in
$L^1(G)$, the sum of the $H_i$ is dense in $H_\pi$. If $H_\pi\ne0$,
there is an $i$ such that $H_i\ne0$. For $j\ne i$, one has
$A_jA_i=0$, and therefore
\[
 \pi(A_j)\pi(A_i)(H_\pi)=0,\qquad
 \pi(A_j)(H_i)=0.
\]
Let $\xi$ be a nonzero element of $H_i$. The finite-dimensional subspace
$\pi(A_i)\xi$ is equal to $\pi(A)\xi$, and hence to
$\pi(L^1(G))\xi$; it is nonzero and invariant under $\pi(G)$. This
subspace contains a nonzero subspace invariant under $\pi(G)$ on which
the corresponding subrepresentation of $G$ is irreducible
(\XRef{2.3.5}{2.3.5}).

Let $(K_\alpha)$ be a maximal family of pairwise orthogonal,
finite-dimensional subspaces of $H_\pi$, invariant under $\pi(G)$, such
that the corresponding subrepresentations of $G$ are irreducible. Then
$H_\pi\ominus\bigl(\bigoplus_\alpha K_\alpha\bigr)$ is invariant under
$\pi(G)$ and is therefore zero, by what precedes and the maximality of
the family $(K_\alpha)$. Hence
\[
 H_\pi=\bigoplus_\alpha K_\alpha.
\]

\ResultHeading{15.1.4}{15.1.4. Corollary.}
\begin{itshape}
Every irreducible continuous unitary representation of $G$ is
finite-dimensional.
\end{itshape}

\ResultHeading{15.1.5}{15.1.5. Corollary.}
\begin{itshape}
The $C^{*}$-algebra of $G$ is liminal.
\end{itshape}

This follows from \XRef{15.1.4}{15.1.4}, since every linear operator on
a finite-dimensional Hilbert space is compact.

\ResultHeading{15.1.6}{15.1.6. Corollary.}
\begin{itshape}
Let $(\sigma_\alpha)$ be the family of finite-dimensional irreducible
continuous unitary representations of $G$. Then
\[
 \bigcap_\alpha\ker\sigma_\alpha=\{e\}.
\]
\end{itshape}

This follows from \XRef{13.6.6}{13.6.6} and
\XRef{15.1.4}{15.1.4}.

\ResultHeading{15.1.7}{15.1.7. Corollary.}
\begin{itshape}
$G$ is a projective limit of compact Lie groups.
\end{itshape}

Let $(\rho_\alpha)$ be the family of finite-dimensional continuous
unitary representations of $G$. The family $(\ker\rho_\alpha)$ is
decreasingly directed, and
\[
 \bigcap_\alpha\ker\rho_\alpha=\{e\}
 \qquad\text{(\XRef{15.1.6}{15.1.6})}.
\]
Thus $G$ is the projective limit of the groups
$G/\ker\rho_\alpha$, themselves isomorphic to $\rho_\alpha(G)$. Since
$\dim\rho_\alpha<+\infty$, the unitary group of $H_{\rho_\alpha}$ is a
Lie group; since $\rho_\alpha(G)$ is a closed subgroup of this group,
$\rho_\alpha(G)$ is a Lie group.

\ResultHeading{15.1.8}{15.1.8. Corollary.}
\begin{itshape}
Let $\pi$ be a continuous unitary representation of $G$.
\begin{enumerate}[label=(\roman*)]
 \item There exist cardinals $n_\sigma$, $\sigma\in\widehat G$, such that
 \[
  \pi=\bigoplus_{\sigma\in\widehat G}n_\sigma\sigma.
 \]
 \item The subspace $H_{n_\sigma\sigma}$ of $H_\pi$ is independent of
 the chosen decomposition of $\pi$. It is the closed vector subspace of
 $H_\pi$ spanned by the spaces of the subrepresentations equivalent to
 $\sigma$.

% Source PDF physical page 295
 \item The cardinals $n_\sigma$ are independent of the chosen
 decomposition of $\pi$.
\end{enumerate}
\end{itshape}

Assertion (i) follows at once from \XRef{15.1.3}{15.1.3}. Let $\sigma'$
be a subrepresentation of $\pi$ equivalent to
$\sigma\in\widehat G$. Then $\sigma'$ is disjoint from $\tau$ for
$\tau\in\widehat G$, $\tau\ne\sigma$ (\XRef{5.2.2}{5.2.2}); consequently
$H_{\sigma'}$ is orthogonal to $H_{n_\tau\tau}$ for $\tau\ne\sigma$
[condition (iii) of \XRef{5.2.1}{5.2.1}]. Hence
$H_{\sigma'}\subset H_{n_\sigma\sigma}$, which proves (ii). One has
\[
 \dim H_{n_\sigma\sigma}=n_\sigma\dim\sigma,
\]
so (iii) follows from (ii). Another intrinsic description of the
$H_{n_\sigma\sigma}$ will be given in \XRef{15.3.12}{15.3.12}.

\ResultHeading{15.1.9}{15.1.9.}
With the notation of \XRef{15.1.8}{15.1.8}, we say that $\sigma$ occurs
$n_\sigma$ times in $\pi$.

Bibliography: \cite{ref58,ref82,ref86,ref104}.

\BookSubsection{15.2}{Irreducible Representations of a Compact Group}

\ResultHeading{15.2.1}{15.2.1.}
By what precedes, the study of the continuous unitary representations of
$G$ reduces to that of the irreducible representations. The latter are
obviously square-integrable, so the results of \SectionRef{14}{\S14} apply to
them. In particular, they are all subrepresentations of the regular
representation.

\ResultHeading{15.2.2}{15.2.2. Theorem.}
\begin{itshape}
Let $A=L^2(G)$ be the Hilbert algebra of $G$, let $(F_i)$ be the family
of minimal projections of $\mathfrak U(A)\cap\mathfrak V(A)$, let
$\lambda,\rho$ be the left and right regular representations of $G$, let
$\lambda_i,\rho_i$ be the subrepresentations of $\lambda,\rho$ defined
by $F_i$, let $\sigma_i$ be the class of irreducible representations
whose associated central projection is $F_i$, and put
$\delta_i=\dim\sigma_i$.
\begin{enumerate}[label=(\roman*)]
 \item $\lambda$ is the Hilbert sum of the $\lambda_i$, and $\rho$ is
 the Hilbert sum of the $\rho_i$.
 \item $\lambda_i$ is the Hilbert sum of $\delta_i$ representations of
 class $\sigma_i$, and $\rho_i$ is the Hilbert sum of $\delta_i$
 representations of class $\overline{\sigma_i}$.
\end{enumerate}
\end{itshape}

We know that $L^2(G)$ is the Hilbert sum of the $F_i(L^2(G))$
(\XRef{15.1.2}{15.1.2}), which proves (i); and (ii) is a special case
of \XRef{14.3.5}{14.3.5(iii)}.

\ResultHeading{15.2.3}{15.2.3. Proposition.}
\begin{itshape}
Let $\sigma$ be an irreducible continuous unitary representation of $G$.
The formal dimension $d_\sigma$ of $\sigma$ is equal to $\dim\sigma$.
\end{itshape}

The algebra $\sigma(\mathcal K(G))$ is irreducible on $H_\sigma$, and
hence equals $\mathcal L(H_\sigma)$, since $\dim H_\sigma<+\infty$. Let
$f\in\mathcal K(G)$ be such that $\sigma(f)=1$, and put
$g=f\ast f^{\ast}$. For every $t\in G$, let $g^t$ be the function
$s\mapsto g(tst^{-1})$ on $G$. Then
\[
 \sigma(g^t)=\sigma(t)^{-1}\sigma(f\ast f^{\ast})\sigma(t)=1.
\]
Let $\xi$ be a unit vector of $H_\sigma$. By
\XRef{14.4.3}{14.4.3},
\[
 d_\sigma^{-1}(\dim\sigma)
 =d_\sigma^{-1}\operatorname{Tr}\sigma(g)
 =\int(\sigma(g^t)\xi\mid\xi)\,dt=1,
\]
whence $d_\sigma=\dim\sigma$.

% Source PDF physical page 296
\ResultHeading{15.2.4}{15.2.4. Theorem.}
\begin{itshape}
For every $f\in L^2(G)$, one has
\[
 \int_G|f(s)|^2\,ds
 =\sum_{\sigma\in\widehat G}(\dim\sigma)
  \operatorname{Tr}\bigl(\sigma(f)\sigma(f^{\ast})\bigr).
\]
\end{itshape}

For every $\sigma\in\widehat G$, let $F_\sigma$ be the central projection
on $L^2(G)$ associated with $\sigma$. By
\XRef{15.1.2}{15.1.2}, one has
\[
 L^2(G)=\bigoplus_{\sigma\in\widehat G}F_\sigma(L^2(G)),
 \qquad
 \|f\|_2^2=\sum_{\sigma\in\widehat G}\|F_\sigma f\|_2^2.
\]
By \XRef{14.4.2}{14.4.2} and \XRef{15.2.3}{15.2.3},
\[
 \|F_\sigma f\|_2^2
 =(\dim\sigma)\operatorname{Tr}
   \bigl(\sigma(f)\sigma(f)^{\ast}\bigr).
\]

\ResultHeading{15.2.5}{15.2.5.}
For every $\sigma\in\widehat G$, choose an orthonormal basis
\[
 (\xi_1,\xi_2,\ldots,\xi_{\dim\sigma})
\]
of $H_\sigma$, and put
\[
 \varphi_{i,j,\sigma}(s)=(\sigma(s)\xi_i\mid\xi_j).
\]
Then, by \XRef{15.1.2}{15.1.2} and
\XRef{14.3.5}{14.3.5(i)}, the functions
\[
 (\dim\sigma)^{1/2}\varphi_{i,j,\sigma}
 \quad
 (\sigma\in\widehat G,\quad
 i,j=1,2,\ldots,\dim\sigma)
\]
form an orthonormal basis of $L^2(G)$. By
\XRef{14.3.3}{14.3.3} and \XRef{14.3.7}{14.3.7}, one has
\[
 \varphi_{i,j,\sigma}\ast\varphi_{i',j',\sigma}
 =(\dim\sigma)^{-1}\delta_i^{\,j'}\varphi_{i',j,\sigma},
 \qquad
 \varphi_{i,j,\sigma}\ast\varphi_{i',j',\sigma'}=0
 \quad\text{if }\sigma\ne\sigma'.
\]

Bibliography: \cite{ref86,ref108,ref181}.

\BookSubsection{15.3}{Characters of Compact Groups}

\ResultHeading{15.3.1}{15.3.1. Definition.}
\begin{itshape}
Let $\sigma$ be an irreducible continuous unitary representation of $G$.
The \emph{character} of $\sigma$, denoted by $\chi_\sigma$, is the
continuous function
\[
 s\longmapsto\operatorname{Tr}\sigma(s)
\]
on $G$. The function $(\dim\sigma)^{-1}\chi_\sigma$ is called the
\emph{normalized character}.
\end{itshape}

When $G$ is compact and commutative, this recovers the usual notion of a
character.

\ResultHeading{15.3.2}{15.3.2.}
The character of $\sigma$ depends only on the class of $\sigma$. If
$(\xi_1,\ldots,\xi_n)$ is an orthonormal basis of $H_\sigma$, then
\[
 \chi_\sigma(s)
 =(\sigma(s)\xi_1\mid\xi_1)+\cdots+
  (\sigma(s)\xi_n\mid\xi_n),
\]
which proves that $\chi_\sigma$ is a function of positive type
[\XRef{13.4.5}{13.4.5(ii)}], equal to $\dim\sigma$ at $e$.

We shall say that a function on a group is \emph{central} if it is
invariant under the inner automorphisms of that group. Definition
\XRef{15.3.1}{15.3.1} immediately shows that $\chi_\sigma$ is central.

\ResultHeading{15.3.3}{15.3.3.}
By \XRef{15.2.5}{15.2.5}, if $\sigma,\sigma'\in\widehat G$ are
inequivalent, then
\begin{align}
 (\chi_\sigma\mid\chi_{\sigma'})&=0,
 \tag{1}\label{eq:15.3.3-1}\\
 \chi_\sigma\ast\chi_{\sigma'}&=0.
 \tag{2}\label{eq:15.3.3-2}
\end{align}
% Source PDF physical page 297
On the other hand, retaining the notation of \XRef{15.3.2}{15.3.2},
\[
 (\chi_\sigma\ast\chi_\sigma)(s)
 =n^{-1}\bigl[(\sigma(s)\xi_1\mid\xi_1)+\cdots+
  (\sigma(s)\xi_n\mid\xi_n)\bigr],
\]
whence
\begin{align}
 \chi_\sigma\ast\chi_\sigma
 &=(\dim\sigma)^{-1}\chi_\sigma,
 \tag{3}\label{eq:15.3.3-3}\\
 \|\chi_\sigma\|_2&=1.
 \tag{4}\label{eq:15.3.3-4}
\end{align}

\ResultHeading{15.3.4}{15.3.4.}
Let $K_\sigma$ be the bi-invariant subspace of $L^2(G)$ associated with
$\sigma$. As a Hilbert algebra, $K_\sigma$ is isomorphic to the algebra
of Hilbert--Schmidt operators on a finite-dimensional Hilbert space
(\XRef{15.1.2}{15.1.2}). This algebra has a one-dimensional center. Now
\[
 \check\chi_\sigma=\overline{\chi_\sigma}=\chi_{\overline\sigma}
\]
is a central function on $G$ belonging to $K_\sigma$
(\XRef{14.3.1}{14.3.1}). Thus the center of $K_\sigma$ is
$\mathbb C\chi_{\overline\sigma}$. By \XRef{15.1.2}{15.1.2}, the center
of the Hilbert algebra $L^2(G)$ is the Hilbert sum of the
$\mathbb C\chi_\sigma$. In view of formula
\eqref{eq:15.3.3-4}, the $\chi_\sigma$ form an orthonormal basis of the
center of $L^2(G)$.

\ResultHeading{15.3.5}{15.3.5.}
By formula \eqref{eq:15.3.3-3}, the operator of convolution by
$(\dim\sigma)\chi_{\overline\sigma}$ on $L^2(G)$ is a nonzero
projection. Since
$\chi_{\overline\sigma}=\chi_\sigma^{\ast}$, this projection is
Hermitian. It belongs to $\mathfrak U(A)\cap\mathfrak V(A)$, where $A$
is the Hilbert algebra of $G$, and its range is contained in $K_\sigma$.
Since $P_{K_\sigma}$ is a minimal projection of
$\mathfrak U(A)\cap\mathfrak V(A)$, the operator of convolution by
$(\dim\sigma)\chi_{\overline\sigma}$ is $P_{K_\sigma}$. Thus, for every
$f\in L^2(G)$,
\[
 f=\sum_{\sigma\in\widehat G}
  (\dim\sigma)(f\ast\chi_{\overline\sigma}),
\]
the series converging in $L^2(G)$. If $f$ is central, this expansion
reduces to the expansion of $f$ in the orthonormal basis
$(\chi_\sigma)$ of the center of $L^2(G)$:
\[
 f=\sum_{\sigma\in\widehat G}(f\mid\chi_\sigma)\chi_\sigma.
\]

Let $f$ be a continuous central complex-valued function on $G$. In view
of \XRef{15.1.1}{15.1.1}, $f$ is a uniform limit on $G$ of functions of
the form $f\ast g$, where $g$ is central and belongs to $L^2(G)$. Since
$f$ and $g$ are limits in $L^2(G)$ of finite linear combinations of
characters of $G$, every continuous central complex-valued function on
$G$ is a uniform limit on $G$ of finite linear combinations of
characters of $G$.

For $G$ equal to the group of real numbers modulo $1$, these facts reduce
to familiar facts about Fourier series.

\ResultHeading{15.3.6}{15.3.6.}
For every finite-dimensional continuous unitary representation $\pi$ of
$G$, one can define the character $\chi_\pi$ of $\pi$ as in
\XRef{15.3.1}{15.3.1}.

% Source PDF physical page 298
Let
\[
 \pi=n_1\sigma_1\oplus\cdots\oplus n_p\sigma_p
\]
be the decomposition of $\pi$ into irreducible representations, where
the $\sigma_i$ are pairwise inequivalent and the $n_i$ are integers. Then
\[
 (\chi_\pi\mid\chi_{\sigma_i})
 =n_i(\chi_{\sigma_i}\mid\chi_{\sigma_i})=n_i.
\]
Thus knowledge of $\chi_\pi$ determines the $n_i$. In other words, two
finite-dimensional continuous unitary representations of $G$ are
equivalent if and only if they have the same character.

If $\pi$ and $\pi'$ are two finite-dimensional continuous unitary
representations of $G$, then
\[
 \chi_{\pi\otimes\pi'}=\chi_\pi\chi_{\pi'}.
\]

Henceforth, when we speak of characters of a compact group, we shall mean
only characters of irreducible representations.

\ResultHeading{15.3.7}{15.3.7. Proposition.}
\begin{itshape}
Let $\mu$ be a central measure on $G$, that is, one which commutes with
every measure under convolution. For $\mu\gg0$, it is necessary and
sufficient that $\mu(\chi)\geq0$ for every character $\chi$ of $G$.
\end{itshape}

If $\mu\gg0$, then
\[
 \mu(\chi)=(\dim\sigma)\mu(\check\chi\ast\chi),
\]
where $\sigma$ is a representation with character $\chi$. Now
$\mu(\check\chi\ast\chi)=\mu(\chi\ast\chi^{\ast})\geq0$.

Conversely, suppose that $\mu(\chi)\geq0$ for every character $\chi$.
Let $\lambda$ be the left regular representation of $G$. To prove that
$\mu\gg0$, it is enough to prove that $\lambda(\mu)\geq0$
(\XRef{13.7.4}{13.7.4}). Let $A$ be the Hilbert algebra of $G$, let
$(F_i)$ be the family of minimal projections of
$\mathfrak U(A)\cap\mathfrak V(A)$, and put $K_i=F_i(L^2(G))$. Since
$\mu$ is central,
\[
 \lambda(\mu)\in\mathfrak U(A)\cap\mathfrak V(A),
\]
and the restriction of $\lambda(\mu)$ to each $K_i$ is a scalar
$\lambda_i$. Let $\chi_i$ be the character belonging to $K_i$. If
$\sigma_i$ is a representation with character $\chi_i$, then
\begin{align*}
 \lambda_i(\chi_i\mid\chi_i)
 &=(\lambda(\mu)\chi_i\mid\chi_i)
  =(\mu\ast\chi_i\mid\chi_i)\\
 &=(\mu\ast\chi_i\ast\check\chi_i)(e)
  =(\dim\sigma_i)^{-1}(\mu\ast\chi_i)(e)\\
 &=(\dim\sigma_i)^{-1}\mu(\check\chi_i)\geq0.
\end{align*}

\ResultHeading{15.3.8}{15.3.8. Proposition.}
\begin{itshape}
Let $f\in L^2(G)$ be a central function on $G$. Write
\[
 f=\sum_{\sigma\in\widehat G}\lambda_\sigma\chi_\sigma,
 \qquad
 \sum_{\sigma\in\widehat G}|\lambda_\sigma|^2<+\infty
 \quad\text{(\XRef{15.3.4}{15.3.4})}.
\]
For $f$ to be continuous and of positive type, it is necessary and
sufficient that $\lambda_\sigma\geq0$ for every
$\sigma\in\widehat G$ and
\[
 \sum_{\sigma\in\widehat G}\lambda_\sigma\dim\sigma<+\infty.
\]
\end{itshape}

The condition is sufficient because $\chi_\sigma$ is of positive type
and is bounded in absolute value by $\dim\sigma$. Suppose that $f$ is
continuous and of positive type. Then $f=g\ast g^{\ast}$ for some
$g\in L^2(G)$ (\XRef{13.8.6}{13.8.6}). For every
$\sigma\in\widehat G$, let $K_\sigma$ be the bi-invariant subspace of
$L^2(G)$ associated with $\sigma$. One has
\[
 g=\sum_{\sigma\in\widehat G}\mu_\sigma g_\sigma,
\]
% Source PDF physical page 299
where
\[
 g_\sigma\in K_\sigma,\qquad
 \|g_\sigma\|_2=1,\qquad
 \sum_{\sigma\in\widehat G}|\mu_\sigma|^2<+\infty.
\]
Since $g_\sigma\ast g_{\sigma'}^{\ast}=0$ for
$\sigma\ne\sigma'$, and
\[
 \|u\ast v\|_\infty\leq\|u\|_2\|v\|_2
 \qquad(u,v\in L^2(G)),
\]
$f$ is the sum of the uniformly convergent series
\[
 \sum_{\sigma\in\widehat G}
 |\mu_\sigma|^2(g_\sigma\ast g_\sigma^{\ast}).
\]
Moreover, $g_\sigma\ast g_\sigma^{\ast}\in K_\sigma$, so that
$|\mu_\sigma|^2(g_\sigma\ast g_\sigma^{\ast})$ is the orthogonal
projection of $f$ onto $K_\sigma$ and is therefore central. We may
suppose that $g_\sigma$ is central if $\mu_\sigma=0$. Thus
$g_\sigma\ast g_\sigma^{\ast}$ is central, and
\[
 (g_\sigma\ast g_\sigma^{\ast})(e)
 =\|g_\sigma\|_2^2=1.
\]
By \XRef{15.3.4}{15.3.4},
\[
 g_\sigma\ast g_\sigma^{\ast}
 =(\dim\sigma)^{-1}\chi_{\overline\sigma},
\]
and hence
\[
 \lambda_{\overline\sigma}
 =|\mu_\sigma|^2(\dim\sigma)^{-1},
\]
which completes the proof.

\ResultHeading{15.3.9}{15.3.9. Corollary.}
\begin{itshape}
In the convex set of continuous central functions $f$ of positive type on
$G$ such that $f(e)\leq1$, the extreme points are $0$ and the normalized
characters.
\end{itshape}

\ResultHeading{15.3.10}{15.3.10. Proposition.}
\begin{itshape}
Let $\psi$ be a continuous complex-valued function on $G$. For $\psi$ to
be a normalized character, it is necessary and sufficient that
$\psi\ne0$ and
\begin{equation}
 \int\psi(xsx^{-1}t)\,dx=\psi(s)\psi(t)
 \qquad(s,t\in G).
 \tag{1}\label{eq:15.3.10-1}
\end{equation}
\end{itshape}

Suppose that there exists $\sigma\in\widehat G$ such that
$\psi=(\dim\sigma)^{-1}\chi_\sigma$. Put
\[
 S=\int\sigma(xsx^{-1})\,dx\in\mathcal L(H_\sigma).
\]
For every $t\in G$, one has
\[
 \sigma(t)S
 =\int\sigma(txsx^{-1})\,dx
 =\int\sigma\bigl(xs(t^{-1}x)^{-1}\bigr)\,dx
 =S\sigma(t).
\]
Since $\sigma$ is irreducible, $S$ is scalar. Moreover,
\[
 \operatorname{Tr}S
 =\int\chi_\sigma(xsx^{-1})\,dx
 =\int\chi_\sigma(s)\,dx
 =\chi_\sigma(s).
\]
Consequently
\begin{equation}
 \int\sigma(xsx^{-1})\,dx
 =(\dim\sigma)^{-1}\chi_\sigma(s)\,1.
 \tag{2}\label{eq:15.3.10-2}
\end{equation}
It follows that
\[
 \int\sigma(xsx^{-1}t)\,dx
 =\left(\int\sigma(xsx^{-1})\,dx\right)\sigma(t)
 =(\dim\sigma)^{-1}\chi_\sigma(s)\sigma(t).
\]
Taking traces of both sides gives \eqref{eq:15.3.10-1}.

% ===== ASSEMBLED TRANSLATION CHUNK 11: chunk_300_329.tex =====
% Source PDF physical page 300
Suppose now that $\psi\ne0$ and that $\psi$ satisfies \eqref{eq:15.3.10-1}.
Since $\psi\ne0$, there is a continuous irreducible unitary representation
$\sigma$ of $G$ such that $\bar\sigma(\psi)\ne0$. We have
\begin{align*}
\psi(s)\bar\sigma(\psi)
 &=\int\psi(s)\psi(t)\bar\sigma(t)\,dt
  =\iint\psi(xsx^{-1}t)\bar\sigma(t)\,dx\,dt\\
 &=\iint\psi(t)\bar\sigma(xs^{-1}x^{-1}t)\,dx\,dt\\
 &=\int\psi(t)\left(\int\bar\sigma(xs^{-1}x^{-1})\,dx\right)
       \bar\sigma(t)\,dt.
\end{align*}
By \eqref{eq:15.3.10-2},
\begin{align*}
\psi(s)\bar\sigma(\psi)
 &=\int\psi(t)(\dim\sigma)^{-1}\chi_{\bar\sigma}(s^{-1})
       \bar\sigma(t)\,dt\\
 &=(\dim\sigma)^{-1}\chi_\sigma(s)\bar\sigma(\psi),
\end{align*}
and hence $\psi=(\dim\sigma)^{-1}\chi_\sigma$.

\ResultHeading{15.3.11}{15.3.11. Lemma.}
\begin{itshape}
Let $\sigma$ and $\sigma'$ be continuous irreducible unitary representations
of $G$. If $\sigma$ and $\sigma'$ are inequivalent, then
$\sigma'(\chi_{\bar\sigma})=0$. If $\sigma$ and $\sigma'$ are equivalent,
then
\[
\sigma'(\chi_{\bar\sigma})=(\dim\sigma)^{-1}1.
\]
\end{itshape}

Let $\lambda$ be the left regular representation of $G$. Then
$\lambda((\dim\sigma)\chi_{\bar\sigma})$ is the projection onto the
bi-invariant subspace $K_\sigma$ of $L^2(G)$ associated with $\sigma$
(\XRef{15.3.5}{15.3.5}). On the other hand, we may suppose that $\sigma'$
is a subrepresentation of $\lambda$ (\XRef{15.2.1}{15.2.1}), and
$H_{\sigma'}$ is orthogonal to $K_\sigma$, or is contained in $K_\sigma$,
according as $\sigma'$ is inequivalent or equivalent to $\sigma$
(\XRef{14.2.3}{14.2.3}).

\ResultHeading{15.3.12}{15.3.12. Theorem.}
\begin{itshape}
Let $\pi$ be a continuous unitary representation of $G$, and let
$\pi=\bigoplus_{\sigma\in\widehat G}n_\sigma\sigma$ be its decomposition
into irreducible representations. For every $\sigma\in\widehat G$, the
orthogonal projection of $H_\pi$ onto $H_{n_\sigma\sigma}$ is
$\pi((\dim\sigma)\chi_{\bar\sigma})$.
\end{itshape}

This follows at once from \XRef{15.3.11}{15.3.11}.

Bibliography: \cite{ref50,ref86,ref108,ref181}.

\BookSubsection{15.4}{Representations of Finite Groups}

\ResultHeading{15.4.1}{15.4.1. Proposition.}
\begin{itshape}
Let $G$ be a finite group of order $n$. Let $r$ be the number of conjugacy
classes of elements of $G$. Then $\widehat G$ has $r$ elements, and
\[
\sum_{\sigma\in\widehat G}(\dim\sigma)^2=n.
\]
\end{itshape}

For each $\sigma\in\widehat G$, let $K_\sigma$ be the bi-invariant
subspace of $L^2(G)$ associated with $\sigma$. We have
\[
\dim K_\sigma=(\dim\sigma)^2,
\qquad
n=\dim L^2(G)=\sum_{\sigma\in\widehat G}\dim K_\sigma,
\]
% Source PDF physical page 301
by \XRef{15.2.1}{15.2.1} and \XRef{15.2.2}{15.2.2}. Thus
$\sum_{\sigma\in\widehat G}(\dim\sigma)^2=n$. On the other hand, the
number of elements of $\widehat G$ equals the number of characters of $G$,
that is, the dimension of the center $Z$ of $L^2(G)$
(\XRef{15.3.4}{15.3.4}); the characteristic functions of the conjugacy
classes of elements of $G$ form a basis of the vector space $Z$. Hence
$\dim Z=r$.

\ResultHeading{15.4.2}{15.4.2.}
Let $C_1$, $C_2$, \ldots, $C_r$ be the distinct conjugacy classes in $G$,
and let $\sigma_1$, $\sigma_2$, \ldots, $\sigma_r$ be the distinct elements of
$\widehat G$. The \emph{character table} of $G$ is the matrix
$(\chi_i^j)_{1\leq i,j\leq r}$, where $\chi_i^j$ is the value on $C_i$ of
the character of $\sigma_j$. Let $h_i$ be the number of elements of $C_i$.
By formulas (1) and (4) of \XRef{15.3.3}{15.3.3},
\[
\frac1n\sum_i h_i\chi_i^j\overline{\chi_i^{j'}}
=
\begin{cases}
0,&j\ne j',\\
1,&j=j'.
\end{cases}
\tag{1}\label{eq:15.4.2-1}
\]
In other words, the matrix
$\bigl(\sqrt{h_i/n}\,\chi_i^j\bigr)$ is unitary. Consequently
\[
\sum_j\sqrt{\frac{h_i}{n}}\,\chi_i^j
       \sqrt{\frac{h_{i'}}{n}}\,\overline{\chi_{i'}^j}
=\delta_i^{i'},
\]
or
\[
\frac1n\sum_j\chi_i^j\overline{\chi_{i'}^j}
=
\begin{cases}
0,&i\ne i',\\[2pt]
1/h_i,&i=i'.
\end{cases}
\tag{2}\label{eq:15.4.2-2}
\]

\ResultHeading{15.4.3}{15.4.3.}
The following proposition requires very different methods.

\ResultLabel{Proposition.}
\begin{itshape}
The dimension of an irreducible representation of $G$ divides the order of
$G$.
\end{itshape}

Retain the notation of \XRef{15.4.2}{15.4.2}. Put
$d_j=\dim\sigma_j$.

\emph{a.} Let $s\in C_i$. Then
$\chi_i^j=\operatorname{Tr}\sigma_j(s)$ is the sum of the eigenvalues of
$\sigma_j(s)$. Since $\sigma_j(s)$ has finite order, these eigenvalues are
roots of unity. Thus $\chi_i^j$ is an algebraic integer.

\emph{b.} Let $f_i$ be the characteristic function of $C_i$. We have
\[
f_i\ast f_{i'}=\sum_{i''=1}^{r}c_{ii'}^{i''}f_{i''},
\]
where $c_{ii'}^{i''}$ is the number of ways in which an element of
$C_{i''}$ can be written as the product of an element of $C_i$ and an
element of $C_{i'}$. Thus the $c_{ii'}^{i''}$ are rational integers. On the
other hand, for every $\sigma\in\widehat G$,
\[
\sigma(f_i)\sigma(f_{i'})
=\sum_{i''=1}^{r}c_{ii'}^{i''}\sigma(f_{i''}).
\]
% Source PDF physical page 302
Identifying $\sigma(f_i)$, $\sigma(f_{i'})$, and $\sigma(f_{i''})$ with
scalars, we see that $\sigma(f_i)$ is an eigenvalue of the matrix
$(c_{ii'}^{i''})_{1\leq i',i''\leq r}$, and hence is an algebraic integer.
Now
\[
\sigma_j(f_i)=\frac1{d_j}\operatorname{Tr}\sigma_j(f_i)
=\frac1{d_j}\sum_{s\in C_i}\operatorname{Tr}\sigma_j(s)
=\frac{h_i}{d_j}\chi_i^j.
\]
Thus $(h_i/d_j)\chi_i^j$ is an algebraic integer.

\emph{c.} By formula \eqref{eq:15.4.2-1},
\[
\frac{n}{d_j}
=\sum_i\left(\frac{h_i}{d_j}\chi_i^j\right)\overline{\chi_i^j}
\]
is an algebraic integer; on the other hand it is rational, and is therefore a
rational integer.

The bibliography concerning representations of finite groups is enormous.
See, for example, M. Hall, \emph{The Theory of Groups}, The Macmillan
Company, New York, 1959.

\BookSubsection{15.5}{Use of Compact Subgroups of Arbitrary Groups}

\ResultHeading{15.5.1}{15.5.1.}
Let $G'$ be a locally compact group and $G$ a compact subgroup of $G'$.
If $f\in L^1(G)$, we may identify $f$ with a measure on $G$, hence with a
compactly supported measure on $G'$. We shall do so in this subsection.

\ResultLabel{Lemma.}
\begin{itshape}
Let $\pi$ be a continuous irreducible unitary representation of $G'$, put
$\rho=\pi|G$, and let
$\rho=\bigoplus_{\sigma\in\widehat G}n_\sigma\sigma$ be the decomposition
of $\rho$ into irreducible representations. Suppose that $n_\sigma$ is
finite for every $\sigma\in\widehat G$.
\begin{enumerate}[label=\textup{(\roman*)}]
\item If $\chi$ is a character of $G$ and $f\in L^1(G')$, then
$\pi(f\ast\chi)$ is an operator of finite rank.
\item The linear combinations of the $f\ast\chi$, where $f\in L^1(G')$
and $\chi$ is a character of $G$, are dense in $L^1(G')$.
\item $\pi(C^{*}(G'))=\mathcal{LC}(H_\pi)$.
\end{enumerate}
\end{itshape}

Let $\sigma\in\widehat G$. By \XRef{15.3.12}{15.3.12},
$\rho(\chi_{\bar\sigma})$ has rank $n_\sigma\dim\sigma$. Thus, for every
$f\in L^1(G')$,
$\pi(f\ast\chi_{\bar\sigma})=\pi(f)\rho(\chi_{\bar\sigma})$ has finite
rank. This proves (i).

There are arbitrarily small neighborhoods $V_i$ of $e$ in $G$ invariant
under the inner automorphisms of $G$. Let $\varphi_i$ be the characteristic
function of $V_i$ in $G$. By \XRef{15.3.5}{15.3.5}, $\varphi_i$ is a
limit in $L^2(G)$, hence in $M^1(G)$, of finite linear combinations of
characters of $G$. If $f\in L^1(G')$, then $f\ast\varphi_i$ can be made
arbitrarily close to $f$ in $L^1(G')$. This proves (ii).

Finally, (iii) follows from (i), (ii), and \XRef{4.1.11}{4.1.11}.

% Source PDF physical page 303
\ResultHeading{15.5.2}{15.5.2. Theorem.}
\begin{itshape}
Let $G'$ be a locally compact group and $G$ a compact subgroup of $G'$.
Suppose that, for every $\pi\in\widehat{G'}$ and every
$\sigma\in\widehat G$, $\sigma$ occurs only finitely many times in
$\pi|G$. Then $G'$ is liminal.
\end{itshape}

This follows at once from \XRef{15.5.1}{15.5.1 (iii)}.

\ResultHeading{15.5.3}{15.5.3. Lemma.}
\begin{itshape}
Let $\Gamma$ be a topological group, let $\sigma$ be a representation of
$\Gamma$ of finite dimension $d$, and let $n$ be an integer greater than
$d$. Then $n\sigma$ has no cyclic vector.
\end{itshape}

Indeed, $n\sigma=\sigma\otimes\tau$, where $\tau$ denotes the trivial
representation of $\Gamma$ of dimension $n$. Every vector of
$H_\sigma\otimes H_\tau$ can be written
$\xi=\xi_1\otimes\eta_1+\cdots+\xi_d\otimes\eta_d$, where
$(\xi_1,\ldots,\xi_d)$ is a basis of $H_\sigma$. Let $H$ be the vector
subspace of $H_\tau$ spanned by $\eta_1,\ldots,\eta_d$. Then
$H\ne H_\tau$ and
$[(\sigma\otimes\tau)(\Gamma)](\xi)\subset H_\sigma\otimes H$; hence
$\xi$ is not cyclic.

\ResultHeading{15.5.4}{15.5.4.}
We shall call a topological group \emph{linear} if it admits a continuous
injective linear representation, not necessarily unitary, in a
finite-dimensional complex vector space.

\ResultHeading{15.5.5}{15.5.5.}
Let $n$ be a positive integer, and let $M_n$ be the algebra of complex
$n$-rowed, $n$-columned matrices. Denote by $r(n)$ the least integer $r$
for which the identity
\[
\sum_{\alpha\in\mathfrak S_r}\varepsilon_\alpha
X_{\alpha(1)}\cdots X_{\alpha(r)}=0
\tag{1}\label{eq:15.5.5-1}
\]
holds in $M_n$ (cf. \XRef{3.6.1}{3.6.1}). Recall that
$r(n+1)>r(n)$.

\ResultLabel{Lemma.}
\begin{itshape}
Let $G'$ be a connected linear semisimple real Lie group, let $G$ be a
maximal compact subgroup of $G'$, let $\sigma\in\widehat G$, and put
$d=\dim\sigma$. Let $\chi$ be the character of $\sigma$ and
$\psi=d\chi$, so that $\psi$ is an idempotent of the algebra
$\mathcal K(G)$ and
$\psi\ast\mathcal K(G')\ast\psi$ is a subalgebra $A$ of
$\mathcal K(G')$. Then identity \eqref{eq:15.5.5-1} holds in $A$, with
$r=r(d^2)$.
\end{itshape}

Let $\tau$ be a linear representation, not necessarily unitary, of $G'$ in
a finite-dimensional complex vector space $V$. If $\xi\in V$ and
$\xi'\in V'$ (the dual of $V$), the function
$s\mapsto\langle\tau(s)\xi,\xi'\rangle$ is called a coefficient of $\tau$.
Let $C$ be the set of linear combinations of coefficients of continuous
finite-dimensional linear representations of $G'$. By considering the
conjugate of such a representation and the tensor product of two such
representations, we see that $C$ is an involutive algebra of continuous
functions on $G'$ [for the product $(f,g)\mapsto fg$ and the involution
$f\mapsto\bar f$]. Since $G'$ is linear, this algebra separates the points
of $G'$. By the Stone--Weierstrass theorem, every continuous complex
function on $G'$ is a uniform limit on every compact set of elements of
$C$. Thus, if $f\in\mathcal K(G')$ and $f\ne0$, there is a $g\in C$ such
that $\int f(s)g(s)\,ds\ne0$.
% Source PDF physical page 304
Consequently, if $f\in\mathcal K(G')$ and $f\ne0$, there is a continuous
finite-dimensional linear representation $\tau$ of $G'$, not necessarily
unitary, such that $\tau(f)\ne0$; and, since $G'$ is connected and
semisimple, we may suppose $\tau$ irreducible. To prove the lemma it is
therefore enough to establish the following: if $\tau$ is an irreducible
finite-dimensional representation of $G'$, not necessarily unitary, and
$r=r(d^2)$, then identity \eqref{eq:15.5.5-1} holds in $\tau(A)$.

The space $V$ of $\tau$ may be given a Hilbert-space structure invariant
under $\tau|G$ (\XRef{B.34}{B 34}). Let
$\tau|G=\bigoplus_{\rho\in\widehat G}n_\rho\rho$ be the decomposition of
$\tau|G$ into irreducible representations. Then $\tau(\psi)$ is the
projection onto the space of $n_\sigma\sigma$
(\XRef{15.3.12}{15.3.12}), and so has rank $n_\sigma d$. Consequently
$\tau(A)$ may be identified with a subalgebra of the algebra of
$n_\sigma d$ by $n_\sigma d$ matrices, and it remains only to prove that
$n_\sigma\leq d$.

There is a closed connected solvable subgroup $N$ of $G'$ such that every
element of $G'$ is the product of an element of $G$ and an element of $N$.
By Lie's theorem, there is a nonzero $\xi\in V$ such that $\tau|N$ leaves
$\mathbb C\xi$ invariant. Since $\xi$ is cyclic for $\tau$, we see that
$\xi$ is cyclic for $\tau|G$. Hence $n_\sigma\sigma$, which is a
subrepresentation of $\tau|G$, has a cyclic vector. Thus
$n_\sigma\leq\dim\sigma=d$ by \XRef{15.5.3}{15.5.3}.

\ResultHeading{15.5.6}{15.5.6. Theorem.}
\begin{itshape}
Let $G'$ be a connected linear semisimple real Lie group and $G$ a maximal
compact subgroup of $G'$.
\begin{enumerate}[label=\textup{(\roman*)}]
\item If $\pi\in\widehat{G'}$ and $\sigma\in\widehat G$, then $\sigma$
occurs at most $\dim\sigma$ times in $\pi|G$.
\item $G'$ is liminal.
\end{enumerate}
\end{itshape}

Let $\pi\in\widehat{G'}$, and write
$\pi|G=\bigoplus_{\sigma\in\widehat G}n_\sigma\sigma$. Let
$\sigma\in\widehat G$, put $d=\dim\sigma$, and let $E\subset H_\pi$ be
the space of $n_\sigma\sigma$. By \XRef{15.5.5}{15.5.5}, identity
\eqref{eq:15.5.5-1} holds in the algebra
$P_E\pi(\mathcal K(G'))P_E$, with $r=r(d^2)$. Since $\pi$ is irreducible,
$\pi(\mathcal K(G'))$ is strongly dense in $\mathcal L(H_\pi)$. Hence
$P_E\pi(\mathcal K(G'))P_E$ is strongly dense in
$P_E\mathcal L(H_\pi)P_E$, and identity \eqref{eq:15.5.5-1} holds in
$\mathcal L(E)$, with $r=r(d^2)$. Since $r(d^2+1)>r(d^2)$, $E$ cannot
contain a vector subspace of dimension $d^2+1$. Thus $\dim E\leq d^2$,
and consequently $n_\sigma\leq d$. This proves (i), and (ii) follows from
(i) and \XRef{15.5.2}{15.5.2}.

Bibliography: \cite{ref54,ref59,ref146}.

\BookSubsection{15.6}{Complements}

In this subsection $G$ denotes a compact group.

\ResultHeading{15.6.1}{15.6.1.}
Every continuous unitary representation of $G$ that has a cyclic vector is
a subrepresentation of the regular representation \cite{ref90}.

\ResultHeading{15.6.2}{15.6.2.}
A continuous complex function $f$ on $G$ is called \emph{almost invariant}
if the linear combinations of its translates ${}_sf$ form a
finite-dimensional vector space. The coefficients of finite-dimensional
% Source PDF physical page 305
representations are such functions. Every continuous complex function on
$G$ is a uniform limit of linear combinations of these coefficients
\cite{ref86}.

\ResultHeading{15.6.3}{15.6.3.}
Let $H$ be a Hilbert space, and let $\pi$ be a representation in $H$ of the
involutive subalgebra $\mathcal K(G)$ of $L^1(G)$. There is a unique
continuous unitary representation $\widetilde\pi$ of $G$ in $H$ such that
\[
\pi(f)=\int\widetilde\pi(s)f(s)\,ds
\]
for every $f\in\mathcal K(G)$.

\ResultHeading{15.6.4}{15.6.4.}
Let $G$ be a compact group and $\varphi$ a continuous function of positive
type on $G$. There are continuous pure functions of positive type
$\varphi_1,\varphi_2,\ldots$ on $G$, and constants
$\lambda_1,\lambda_2,\ldots\geq0$, such that
\[
\sum_n\lambda_n<+\infty,
\qquad
\varphi_1(e)=\varphi_2(e)=\cdots=1,
\qquad
\varphi=\sum_n\lambda_n\varphi_n,
\]
where the series converges uniformly on $G$ \cite{ref50}.

\ResultHeading{15.6.5}{15.6.5.}
All connected compact real Lie groups are known. For such a group, the list
of continuous irreducible unitary representations is also known: these
representations are classified by means of a finite number of integers.
The simplest of the noncommutative connected compact real Lie groups is the
group $G$ of rotations of ordinary three-dimensional space. If $\widetilde
G$ is its universal two-sheeted covering, then $\widetilde G$ has a
continuous irreducible unitary representation $\pi$ of dimension $2$, and
the other continuous irreducible unitary representations of $\widetilde G$
are the symmetric tensor powers of $\pi$.

The list of finite groups is very far from being known, so the results are
less complete in that case. But a great many particular results are known.

\BookSection{16}{Almost-Periodic Functions}

In this section, $G$ denotes a topological group. We shall see that the
search for finite-dimensional continuous unitary representations of $G$
reduces to that for finite-dimensional continuous unitary representations
of a compact group $\Sigma$ associated with $G$. In truth, even for very
simple groups $G$, the group $\Sigma$ is often complicated, so this will
not give us a practical method of determining the finite-dimensional
continuous unitary representations of $G$. But the properties of
\SectionRef{15}{Section~15}, applied to $\Sigma$, will imply remarkable
properties of a certain class of functions on $G$, called almost-periodic
functions.

\BookSubsection{16.1}{Compact Group Associated with a Topological Group}

\ResultHeading{16.1.1}{16.1.1. Theorem.}
\begin{itshape}
Let $G$ be a topological group. There exist a compact group $\Sigma$ and a
continuous homomorphism $\alpha:G\to\Sigma$ having the following property:
% Source PDF physical page 306
for every compact group $\Sigma'$ and every continuous homomorphism
$\alpha':G\to\Sigma'$, there is a unique continuous homomorphism
$\beta:\Sigma\to\Sigma'$ such that $\alpha'=\beta\circ\alpha$.

Moreover, this property determines the pair $(\Sigma,\alpha)$ up to
isomorphism.
\end{itshape}

Let $(\pi_i)$ be the family of finite-dimensional continuous unitary
representations of $G$. Let $U_i$ be the unitary group of $H_{\pi_i}$,
let $U$ be the compact product group of the $U_i$, let $\alpha$ be the
continuous homomorphism $s\mapsto(\pi_i(s))$ of $G$ into $U$, and let
$\Sigma$ be the closure of $\alpha(G)$ in $U$, which is a compact subgroup
of $U$. We show that $(\Sigma,\alpha)$ has the property in the theorem.
Let $\Sigma'$ be a compact group and $\alpha':G\to\Sigma'$ a continuous
homomorphism. The uniqueness of a continuous homomorphism
$\beta:\Sigma\to\Sigma'$ such that $\alpha'=\beta\circ\alpha$ is immediate,
since $\alpha(G)$ is dense in $\Sigma$. We prove the existence of $\beta$.
The intersection of the kernels of the finite-dimensional unitary
representations of $\Sigma'$ is reduced to $e$ (\XRef{15.1.6}{15.1.6});
hence $\Sigma'$ can be identified with a subgroup of a group
$\prod_jV_j$, where each $V_j$ is the unitary group of a finite-dimensional
Hilbert space. Then $\alpha'$ is identified with a homomorphism
$s\mapsto(\rho_j(s))$, where each $\rho_j$ is a continuous homomorphism of
$G$ into $V_j$. Since $\rho_j$ is equivalent to one of the representations
$\pi_i$, there is a continuous homomorphism $\beta_j$ of $\Sigma$ into
$V_j$ such that $\rho_j=\beta_j\circ\alpha$. The $\beta_j$ define a
continuous homomorphism $\beta$ of $\Sigma$ into $\prod_jV_j$ such that
$\alpha'=\beta\circ\alpha$, whence $\beta(\Sigma)\subset\Sigma'$.

The uniqueness of the pair $(\Sigma,\alpha)$ up to isomorphism is an easy
general property of universal problems. Let $(\Sigma_1,\alpha_1)$ be
another solution. There are continuous homomorphisms
$\beta:\Sigma\to\Sigma_1$ and $\beta_1:\Sigma_1\to\Sigma$ such that
$\alpha_1=\beta\circ\alpha$ and
$\alpha=\beta_1\circ\alpha_1$; hence
$\alpha=(\beta_1\circ\beta)\circ\alpha$. Thus $\beta_1\circ\beta$ is the
identity map of $\Sigma$ (by uniqueness of $\beta$ when $\alpha$ and
$\alpha'$ are given). Similarly, $\beta\circ\beta_1$ is the identity map of
$\Sigma_1$. Hence $\beta$ is an isomorphism of $\Sigma$ onto $\Sigma_1$
that carries $\alpha$ to $\alpha_1$.

\ResultHeading{16.1.2}{16.1.2. Definition.}
\begin{itshape}
The group $\Sigma$ is called the compact group associated with $G$, and
$\alpha$ is called the canonical homomorphism of $G$ into $\Sigma$.
\end{itshape}

Thus $\overline{\alpha(G)}=\Sigma$. With the notation of
\XRef{16.1.1}{16.1.1}, $\beta(\Sigma)$ is the closure of $\alpha'(G)$. If
$G$ is compact, we may take $\Sigma=G$ and $\alpha$ to be the identity map
of $G$.

\ResultHeading{16.1.3}{16.1.3. Proposition.}
\begin{itshape}
For every finite-dimensional continuous unitary representation $\rho$ of
$\Sigma$, put $\rho'=\rho\circ\alpha$, which is a finite-dimensional
continuous unitary representation of $G$. This defines a bijection from the
set of classes of finite-dimensional continuous unitary representations of
$\Sigma$ onto the set of classes of finite-dimensional continuous unitary
representations of $G$.
\end{itshape}

If $\rho'$ is equivalent to $\rho_1'$, then $\rho$ is equivalent to
$\rho_1$ (because $\overline{\alpha(G)}=\Sigma$), so the map in question is
injective. On the other hand, if $\pi$ is a finite-dimensional continuous
unitary representation of $G$, there is, by \XRef{16.1.1}{16.1.1},
% Source PDF physical page 307
a continuous unitary representation $\rho$ of $\Sigma$ in $H_\pi$ such
that $\pi=\rho\circ\alpha$; hence the map in question is surjective.

\ResultHeading{16.1.4}{16.1.4.}
When $G$ is a locally compact commutative group, the pair
$(\Sigma,\alpha)$ is given by the following proposition.

\ResultLabel{Proposition.}
\begin{itshape}
Let $G$ be a locally compact commutative group and $\widehat G$ its locally
compact dual group. Let $G'$ be the compact commutative group whose dual
$\widehat{G'}$ is the group $\widehat G$ endowed with the discrete topology.
Let $\varphi$ be the continuous homomorphism of $G$ into $G'$ whose dual
$\widehat\varphi$ is the identity homomorphism of $\widehat{G'}$ into
$\widehat G$. Then $G'$ may be identified with the compact group associated
with $G$, and $\varphi$ with the canonical homomorphism of $G$ into $G'$.
\end{itshape}

If $\chi'\in\widehat{G'}$ is trivial on $\varphi(G)$, then
$\widehat\varphi(\chi')$ is trivial on $G$; hence
$\widehat\varphi(\chi')=1$ and $\chi'=1$. This proves that
$\overline{\varphi(G)}=G'$. Now let $H$ be a compact group and $\psi$ a
continuous homomorphism of $G$ into $H$. We shall prove that there is a
unique continuous homomorphism $\beta$ of $G'$ into $H$ such that
$\psi=\beta\circ\varphi$; this will establish the proposition. Replacing
$H$ by $\psi(H)$, we reduce to the case where $H$ is compact and
commutative. Then the dual $\widehat\psi$ of $\psi$ is a homomorphism from
the discrete group $\widehat H$ into $\widehat G$. We are reduced to proving
that there is a unique homomorphism $\widehat\beta$ from $\widehat H$ into
$\widehat{G'}$ such that
$\widehat\psi=\widehat\varphi\circ\widehat\beta$; this is evident because
$\widehat\varphi$ is the identity map of $\widehat{G'}$ onto $\widehat G$.

Bibliography: \cite{ref86,ref181}.

\BookSubsection{16.2}{Almost-Periodic Functions}

\ResultHeading{16.2.1}{16.2.1. Theorem.}
\begin{itshape}
Let $G$ be a topological group, $\Sigma$ the associated compact group,
$\alpha$ the canonical homomorphism of $G$ into $\Sigma$, $E(G)$ the
vector space of bounded continuous complex functions on $G$, furnished
with the norm of uniform convergence, and $f\in E(G)$. The following
conditions are equivalent:
\begin{enumerate}[label=\textup{(\roman*)}]
\item The set of ${}_sf$, $s\in G$, is relatively compact in $E(G)$.
\item The set of $f_t$, $t\in G$, is relatively compact in $E(G)$.
\item The set of ${}_sf_t$, $s,t\in G$, is relatively compact in $E(G)$.
\item There is a continuous complex function $g$ on $\Sigma$ such that
$f=g\circ\alpha$.
\item $f$ is a uniform limit on $G$ of linear combinations of coefficients
of finite-dimensional continuous irreducible unitary representations of
$G$.
\end{enumerate}
\end{itshape}

[Since $E(G)$ is complete, ``relatively compact'' may be replaced by
``precompact'' in (i), (ii), and (iii).]

(iv) $\Rightarrow$ (iii): Let $g\in\mathcal K(\Sigma)$ and
$f=g\circ\alpha$. The map $(\sigma,\tau)\mapsto{}_{\sigma}g_\tau$ from
$\Sigma\times\Sigma$ into the Banach space $\mathcal K(\Sigma)$ is
continuous; hence the set
% Source PDF physical page 308
of the ${}_{\sigma}g_\tau$ is compact. Its image in $E(G)$ under the
isometric map $h\mapsto h\circ\alpha$ is compact and contains the set of
the ${}_sf_t$, $s,t\in G$.

(iii) $\Rightarrow$ (i) and (ii): This is evident.

(i) $\Rightarrow$ (iv) [the proof of (ii) $\Rightarrow$ (iv) is analogous]:
For every $g\in E(G)$ and every $s\in G$, put
$\psi(s)g={}_{s^{-1}}g$. Then $\psi(s)$ is an isometric linear map of
$E(G)$ onto $E(G)$, and $\psi(ss')=\psi(s)\psi(s')$. Let $A$ be the set
of the ${}_sf$, $s\in G$. It is clear that every $\psi(s)$ induces a
bijection of $A$ onto $A$, hence also a bijection $\varphi(s)$ of
$\overline A$ onto $\overline A$. The $\varphi(s)$ are isometries, and
$\overline A$ is compact if condition (i) holds. Thus $\varphi(G)$ is
relatively compact in the set
$\mathcal C(\overline A,\overline A)$ of continuous maps of $\overline A$
into itself, endowed with the topology of uniform convergence (Ascoli's
theorem). Consequently, the closure of $\varphi(G)$ in
$\mathcal C(\overline A,\overline A)$ is a compact group $\Gamma$ of
homeomorphisms of $\overline A$ (\XRef{B.16}{B 16}). The map $\varphi$
is a homomorphism of $G$ into $\Gamma$. We show that it is continuous.
It is enough to prove that, as $s\to e$, $\varphi(s)$ tends to
$\varphi(e)$ in the topology of uniform convergence on $\overline A$, or
even only in the topology of pointwise convergence on $\overline A$
[which is equivalent here because $\varphi(G)$ is equicontinuous]. Let
$g\in\overline A$ and let $V$ be an open neighborhood of $g$ in
$\overline A$. Then $\overline A\setminus V$ is compact. If
$h\in\overline A\setminus V$, then $|g(u)-h(u)|>0$ at some point $u$ of
$G$; hence there are an open neighborhood $V_h$ of $h$ in $\overline A$
and a neighborhood $W_h$ of $e$ in $G$ such that
$\varphi(s)g\notin V_h$ for $s\in W_h$. We can cover
$\overline A\setminus V$ by finitely many neighborhoods
$V_{h_1},\ldots,V_{h_n}$; then, for
$s\in W_{h_1}\cap\cdots\cap W_{h_n}$, we have
$\varphi(s)g\notin V_{h_1}\cup\cdots\cup V_{h_n}$, and hence
$\varphi(s)g\in V$. Thus $\varphi$ is a continuous homomorphism of $G$
into $\Gamma$. By the definition of $\Sigma$ and $\alpha$, there is a
continuous homomorphism $\beta$ of $\Sigma$ into $\Gamma$ such that
$\varphi=\beta\circ\alpha$. The function
$\sigma\mapsto(\beta(\sigma)f)(e)$ on $\Sigma$ is continuous, and
\[
f(s^{-1})=(\varphi(s)f)(e)=(\beta(\alpha(s))f)(e),
\]
so condition (iv) holds.

(iv) $\Rightarrow$ (v): Suppose that $f=g\circ\alpha$, where
$g\in\mathcal K(\Sigma)$. Then $g$ is a uniform limit on $\Sigma$ of
linear combinations of coefficients of continuous irreducible unitary
representations, necessarily finite-dimensional, of $\Sigma$
(\XRef{13.6.5}{13.6.5} and \XRef{15.1.4}{15.1.4}). Hence $f$ is a
uniform limit on $G$ of linear combinations of coefficients of
finite-dimensional continuous irreducible unitary representations of $G$.

(v) $\Rightarrow$ (iv): Suppose that, for every $n$, there is a linear
combination $f_n$ of coefficients of finite-dimensional continuous
irreducible unitary representations of $G$ such that
$\|f-f_n\|_\infty\leq1/n$. There is a function $g_n$ on $\Sigma$, a linear
combination of coefficients of continuous irreducible unitary
representations of $\Sigma$,
% Source PDF physical page 309
such that $f_n=g_n\circ\alpha$ (\XRef{16.1.3}{16.1.3}). We have
$\|g_m-g_n\|_\infty=\|f_m-f_n\|_\infty\to0$ as $m,n\to+\infty$; hence
the $g_n$ converge uniformly to a continuous function $g$ on $\Sigma$, and
$f=g\circ\alpha$.

\ResultHeading{16.2.2}{16.2.2. Definition.}
\begin{itshape}
An almost-periodic function on $G$ is a bounded continuous complex function
on $G$ satisfying the equivalent conditions of \XRef{16.2.1}{16.2.1}.
\end{itshape}

We denote the set of these functions by $\operatorname{PP}(G)$.

\ResultHeading{16.2.3}{16.2.3.}
By condition (iv) of \XRef{16.2.1}{16.2.1}, a function in
$\operatorname{PP}(G)$ is uniformly continuous for the right and left
uniform structures of $G$. The sum, product, and so forth of two
almost-periodic functions is almost-periodic; $\operatorname{PP}(G)$ is
closed in $E(G)$.

Bibliography: \cite{ref86,ref181}.

\BookSubsection{16.3}{Mean of an Almost-Periodic Function}

\ResultHeading{16.3.1}{16.3.1. Theorem.}
\begin{itshape}
Let $f\in\operatorname{PP}(G)$. Let $K$ be the closed convex hull in
$E(G)$ (notation of \XRef{16.2.1}{16.2.1}) of the set of the ${}_sf$ as
$s$ ranges over $G$. Then $K$ contains exactly one constant $M(f)$. If
$f'$ is the function on $\Sigma$ corresponding to $f$, then
\[
M(f)=\int_\Sigma f'(s)\,ds,
\]
the Haar measure on $\Sigma$ being chosen to have total mass $1$.
\end{itshape}

Through the canonical isomorphism of $\operatorname{PP}(G)$ onto
$\mathcal K(\Sigma)$, we reduce to the case $G=\Sigma$. Let therefore
$f\in\mathcal K(\Sigma)$ and let $\varepsilon>0$. There is a neighborhood
$V$ of $e$ in $\Sigma$ such that $st^{-1}\in V$ implies
$|f(s)-f(t)|\leq\varepsilon$. Let $(V_is_i)_{1\leq i\leq n}$ be a finite
covering of $\Sigma$, and let $(h_i)_{1\leq i\leq n}$ be a partition of
unity subordinate to this covering. Put
$c_i=\int_\Sigma h_i(s)\,ds$. Then
\[
\left|\int_\Sigma f(s)\,ds-\sum_i c_i f(s_it)\right|
=\left|\sum_i\int_\Sigma h_i(s)[f(st)-f(s_it)]\,ds\right|,
\]
and this is bounded by $\varepsilon$, since
$|f(st)-f(s_it)|\leq\varepsilon$ whenever $h_i(s)\ne0$. Moreover,
$\sum_i c_i=1$. Thus the constant $\int_\Sigma f(s)\,ds$ belongs to $K$.
Finally, all translates of $f$ have the same integral on $\Sigma$, and
hence all functions of $K$ have the same integral on $\Sigma$; if a
constant belongs to $K$, its value is necessarily
$\int_\Sigma f(s)\,ds$.

\ResultHeading{16.3.2}{16.3.2. Definition.}
\begin{itshape}
The constant $M(f)$ of \XRef{16.3.1}{16.3.1} is called the mean of $f$.
\end{itshape}

% Source PDF physical page 310
\ResultHeading{16.3.3}{16.3.3.}
The map $f\mapsto M(f)$ is a linear form on $\operatorname{PP}(G)$. We
have $M(f)\geq0$ for $f\geq0$, $M(1)=1$, and
$M({}_sf)=M(f_s)=M(f)$ for every $s\in G$. All this follows from the
equality $M(f)=\int_\Sigma f'(s)\,ds$ in \XRef{16.3.1}{16.3.1}.

\ResultHeading{16.3.4}{16.3.4.}
For $f,g\in\operatorname{PP}(G)$, put
$(f\mid g)=M(f\bar g)$. Then $\operatorname{PP}(G)$ becomes a separated
pre-Hilbert space, canonically isomorphic to $\mathcal K(\Sigma)$ regarded
as a vector subspace of $L^2(\Sigma)$. Let $C$ be the set of classes of
finite-dimensional continuous irreducible unitary representations of $G$.
For each $\sigma\in C$, choose an orthonormal basis
$(\xi_1,\ldots,\xi_{\dim\sigma})$ of $H_\sigma$, and put
\[
\varphi_{i,j,\sigma}(s)=(\sigma(s)\xi_i\mid\xi_j).
\]
Then the functions
$(\dim\sigma)^{1/2}\varphi_{i,j,\sigma}$, where $\sigma\in C$ and
$1\leq i,j\leq\dim\sigma$, form an orthonormal basis of the pre-Hilbert
space $\operatorname{PP}(G)$. This follows from \XRef{15.2.5}{15.2.5}
and \XRef{16.1.3}{16.1.3}.

Bibliography: \cite{ref86,ref181}.

\BookSubsection{16.4}{Groups Embeddable in a Compact Group}

\ResultHeading{16.4.1}{16.4.1. Lemma.}
\begin{itshape}
The following three subgroups of $G$ are equal:
\begin{enumerate}[label=\textup{\arabic*\textsuperscript{o}}]
\item the kernel of the canonical homomorphism $G\to\Sigma$;
\item the intersection of the kernels of all continuous homomorphisms of
$G$ into compact groups;
\item the intersection of the kernels of the finite-dimensional continuous
irreducible unitary representations of $G$.
\end{enumerate}
\end{itshape}

Denote these three sets by $N_1,N_2,N_3$. We have $N_2\supset N_1$ by
the definition of $\Sigma$ and of the homomorphism $G\to\Sigma$. Clearly
$N_3\supset N_2$. Finally, if $s\in N_3$, its canonical image in $\Sigma$
belongs to the kernel of every continuous irreducible unitary
representation of $\Sigma$ (\XRef{16.1.3}{16.1.3}), and hence equals $e$
(\XRef{13.6.6}{13.6.6}); thus $N_3\subset N_1$.

\ResultHeading{16.4.2}{16.4.2. Definition.}
\begin{itshape}
A topological group is said to be embeddable in a compact group if the
subgroup in \XRef{16.4.1}{16.4.1} is reduced to $e$.
\end{itshape}

\ResultHeading{16.4.3}{16.4.3.}
A subgroup of a group embeddable in a compact group is embeddable in a
compact group. A product of groups embeddable in compact groups is
embeddable in a compact group. Hence an inverse limit of groups embeddable
in compact groups is embeddable in a compact group. The quotient of an
arbitrary topological group by the normal subgroup of \XRef{16.4.1}{16.4.1}
is embeddable in a compact group. A locally compact commutative group is
embeddable in a compact group, since the subgroup defined by condition
$3^\circ$ of \XRef{16.4.1}{16.4.1} is reduced to $e$.

\ResultHeading{16.4.4}{16.4.4. Proposition.}
\begin{itshape}
Let $G$ be a topological group generated by a compact neighborhood of $e$,
and suppose that $G$ is embeddable in a compact group. Then $G$ is an
% Source PDF physical page 311
inverse limit of Lie groups locally isomorphic to compact groups.
\end{itshape}

Let $\mathcal F$ be the filter base on $G$ consisting of the kernels of the
finite-dimensional continuous unitary representations of $G$. Since $G$ is
embeddable in a compact group, the intersection of the elements of
$\mathcal F$ is $\{e\}$. By hypothesis, there is a compact symmetric
neighborhood $U$ of $e$ that generates $G$. The trace of $\mathcal F$ on
the compact set $U^3$ converges to $e$. Hence there is an
$N_0\in\mathcal F$ such that $U^3\cap N_0\subset U$. For
$N\in\mathcal F$, $N\subset N_0$, put
\[
K_N=U^2\cap N\subset U.
\]
The $K_N$ have the following properties.

\emph{a.} $K_N$ is a compact subgroup of $G$, since
$K_NK_N^{-1}\subset U^2\cap N=K_N$.

\emph{b.} $K_N$ is normal in $U^\infty=G$, for, if $s\in U$, then
\[
sK_Ns^{-1}\subset U^3\cap N\subset U\cap N\subset K_N.
\]
Since $\bigcap_NK_N=\{e\}$, $G$ is the inverse limit of the $G/K_N$.

\emph{c.} $K_N$ is open in $N$, so $N/K_N$ is a discrete subgroup of
$G/K_N$; consequently $G/K_N$ is locally isomorphic to $G/N$. Now $G/N$
has a finite-dimensional continuous injective unitary representation, and
is therefore a Lie group. Moreover, the Lie algebra $\mathfrak g$ of $G/N$
is isomorphic to a subalgebra of a compact Lie algebra. Hence
$\mathfrak g$ is a product of a commutative Lie algebra and a compact Lie
algebra, and is therefore the Lie algebra of a compact group $C$; thus
$G/K_N$ is locally isomorphic to $C$.

\ResultHeading{16.4.5}{16.4.5. Lemma.}
\begin{itshape}
Let $G$ be a connected locally compact group having a fundamental system
of neighborhoods of $e$ invariant under inner automorphisms. Let $A$ and
$B$ be closed subgroups of $G$ such that: $1^\circ$ $A$ is contained
in the center of $G$; $2^\circ$ $A\subset B$, and $B/A$ is contained
in the center of $G/A$. Then $B$ is contained in the center of $G$.
\end{itshape}

Let $x\in B$, $y\in G$, let $\chi$ be a character of the commutative group
$A$, and let $N$ be the kernel of $\chi$. We have
$xyx^{-1}y^{-1}\in A$. We shall show that
$\chi(xyx^{-1}y^{-1})=1$. Since $\chi$ is arbitrary, this will prove that
$xyx^{-1}y^{-1}=e$ and establish the lemma. There is a neighborhood $V$
of $e$ in $G$ such that $|\chi(u)-1|\leq1$ for $u\in V\cap A$, and a
neighborhood $V'$ of $e$ in $G$, invariant under inner automorphisms, such
that $V'V'^{-1}\subset V$. For every $z\in V'$, the element
$xzx^{-1}z^{-1}$ is an element $a$ of $A$, and the relation
$xzx^{-1}=az$ gives by induction $x^nzx^{-n}=a^nz$. Hence
\[
a^n=(x^nzx^{-n})z^{-1}\in V'V'^{-1}\subset V,
\]
% Source PDF physical page 312
and consequently $a^n\in V\cap A$ and $|\chi(a^n)-1|\leq1$. Since this
holds for every $n$, it follows that $\chi(a)=1$; that is, $x$ and $z$
commute modulo $N$. Since $G$ is connected, $y$ is a product of elements
of $V'$, so $x$ and $y$ commute modulo $N$.

\ResultHeading{16.4.6}{16.4.6. Theorem.}
\begin{itshape}
Let $G$ be a connected locally compact group, and let $Z$ be its center.
The following properties are equivalent:
\begin{enumerate}[label=\textup{(\roman*)}]
\item $G$ is embeddable in a compact group;
\item $G$ is the product of a compact group and a group $\mathbb R^n$;
\item $G$ is an inverse limit of Lie groups locally isomorphic to compact
groups;
\item $G/Z$ is compact;
\item $G$ has a fundamental system of neighborhoods of $e$ invariant under
inner automorphisms.
\end{enumerate}
Moreover, if $G$ has these properties, every connected group locally
isomorphic to $G$ also has them.
\end{itshape}

(i) $\Rightarrow$ (iii): This follows from \XRef{16.4.4}{16.4.4}.

(iii) $\Rightarrow$ (ii): Suppose condition (iii) holds. There is in $G$
a filter base converging to $e$, formed of closed normal subgroups $N$ such
that the groups $G/N=G_N$ are Lie groups locally isomorphic to compact
groups. Moreover, $G_N$ is connected. Let $G_N^1,G_N^2$ be the connected
closed normal subgroups of $G_N$ corresponding to the center and the
derived algebra of the reductive Lie algebra of $G_N$. Then
$G_N=G_N^1G_N^2$, the groups $G_N^1$ and $G_N^2$ commute, and $G_N^2$ is
compact. If $N'\subset N$, the canonical map of $G_{N'}$ onto $G_N$ maps
$G_{N'}^1$ onto $G_N^1$ and $G_{N'}^2$ onto $G_N^2$. Let $A\subset G$ be
the inverse limit of the $G_N^1$ and $K\subset G$ the inverse limit of the
$G_N^2$. Then $A$ is commutative and connected, $K$ is compact, $A$ and
$K$ commute, and $G=AK$. We have $A=\mathbb R^n\times K_1$, with $K_1$
compact (\XRef{B.25}{B 25}). Thus $G=\mathbb R^nK_2$, with a compact
subgroup $K_2$ commuting with $\mathbb R^n$. Necessarily
$K_2\cap\mathbb R^n=\{e\}$, and hence $G$ is isomorphic to
$\mathbb R^n\times K_2$.

(ii) $\Rightarrow$ (v): This follows from \XRef{15.1.1}{15.1.1}.

(v) $\Rightarrow$ (iv): Suppose condition (v) holds. For each $s\in G$,
let $j(s)$ be the inner automorphism of $G$ defined by $s$. The hypothesis
implies that the right and left uniform structures of $G$ coincide and that
$j(G)$ is an equicontinuous group of homeomorphisms of $G$. There is a
compact neighborhood $U$ of $e$ invariant under inner automorphisms. For
each $s\in G$, there is a positive integer $n$ such that $s\in U^n$, since
$G$ is connected; the conjugacy class of $s$ is then contained in $U^n$,
and is therefore relatively compact. Let $\mathcal T$ be the topology of
compact convergence on the set $\mathcal C(G,G)$ of continuous maps of
$G$ into itself. By Ascoli's theorem, $j(G)$ is relatively compact in
$\mathcal C(G,G)$ for $\mathcal T$. Thus the closure of $j(G)$ is a group
$\Gamma$ of homeomorphisms of $G$, compact
% Source PDF physical page 313
for $\mathcal T$ (\XRef{B.16}{B 16}). Finally, $j$ is continuous as a map
of $G$ into $\Gamma$, and therefore the injective homomorphism of $G/Z$
into $\Gamma$ induced by $j$ is a topological isomorphism onto its image.
Thus $G/Z$ is embeddable in a compact group. By the implication
(i) $\Rightarrow$ (ii), already established, we conclude that
$G/Z=A\times K$, with $A$ commutative and $K$ compact. By
\XRef{16.4.5}{16.4.5}, $A$ is reduced to $\{e\}$; hence $G/Z$ is compact.

(iv) $\Rightarrow$ (i): If $G/Z$ is compact, then $Z$ is generated by a
compact neighborhood of $e$ (\XRef{B.24}{B 24}), and hence is of the form
$K\times L$, with $K$ compact and
$L=\mathbb R^n\times\mathbb Z^p$ (\XRef{B.25}{B 25}). Let $s$ be an
element of $G$ distinct from $e$. We shall prove, thereby establishing
(iv) $\Rightarrow$ (i), that there is a continuous homomorphism $\varphi$
of $G$ into a compact group such that $\varphi(s)\ne e$. There is a closed
subgroup $M$ of $L$ such that $s\notin M$ and $L/M$ is compact. Then
$G/M$, which is an extension of $G/Z$ by $K\times(L/M)$, is compact, and
it is enough to take for $\varphi$ the canonical homomorphism of $G$ onto
$G/M$.

Finally, if $G$ has property (v), every connected group locally isomorphic
to $G$ also has it.

In place of ``embeddable in a compact group,'' one also says
``representable in a compact group'' or ``maximally almost-periodic.''

Bibliography: \cite{ref2,ref181}.

\BookSubsection{16.5}{Complements}

\ResultHeading{16.5.1}{16.5.1.}
A group embeddable in a compact group may have infinite-dimensional
continuous irreducible unitary representations \cite{ref50}.

\ResultHeading{16.5.2}{* 16.5.2.}
Let $G$ be a locally compact group. We use again the notation of
\XRef{13.11.6}{13.11.6}.

\emph{a.} For $\varphi,\psi\in E$ and $x\in G$, the function
$t\mapsto\psi(t)\varphi(t^{-1}x)$ is an element $\omega_x$ of $E$. Put
$M(\omega_x)=(\varphi\times\psi)(x)$. Then
$\varphi\times\psi\in\operatorname{PP}(G)$.

\emph{b.} Every continuous function of positive type can be written
$\varphi_1+\varphi_2$, where $\varphi_1$ is of positive type and
almost-periodic, and $\varphi_2$ is of positive type and satisfies
$M(\varphi_2\overline{\varphi_2})=0$.

\emph{c.} Every almost-periodic function of positive type can be written
$\sum\lambda_i\varphi_i$, where $\lambda_i>0$,
$\sum\lambda_i<+\infty$, and $\varphi_i$ is a pure function of positive
type corresponding to a finite-dimensional representation \cite{ref50}.

\ResultHeading{16.5.3}{16.5.3.}
\emph{a.} Let $G$ be a locally compact group and $G_0$ the connected
component of $e$. Suppose that $G/G_0$ is compact. The following conditions
are equivalent: (i) $G$ is embeddable in a compact group; (ii) $G_0$ is
embeddable in a compact group; (iii) $G$ is a semidirect product of a
compact subgroup $K$ and a normal subgroup $V$ isomorphic to
$\mathbb R^n$, every element of $V$ commuting with every element of the
identity component of $K$; (iv) $G$ has a fundamental system of compact
neighborhoods of $e$ invariant under inner automorphisms.

% Source PDF physical page 314
\emph{b.} There are totally disconnected locally compact groups embeddable
in compact groups that do not satisfy condition (iv) of part a
\cite{ref85,ref103}.

\ResultHeading{16.5.4}{16.5.4.}
Problem: prove the equivalence of conditions (i), (ii), (iv), and (v) of
\XRef{16.4.6}{16.4.6} without passing through Lie groups.

\BookSection{17}{Characters of a Locally Compact Group}

\BookSubsection{17.1}{Definitions}

\ResultHeading{17.1.1}{17.1.1.}
Let $G$ be a locally compact group. A \emph{trace} (respectively a
\emph{bitrace}, a \emph{character}) of $G$ means a semifinite lower
semicontinuous trace (respectively a maximal bitrace, a character) of
$C^{*}(G)$.

\ResultHeading{17.1.2}{17.1.2.}
A \emph{traced representation} of $G$ means a pair $(\pi,t)$ with the
following properties: (i) $\pi$ is a continuous unitary representation of
$G$; (ii) $t$ is a faithful normal trace on $\mathfrak U^{+}$, where
$\mathfrak U$ denotes the von Neumann algebra generated by $\pi(G)$;
(iii) $\pi(C^{*}(G))\cap\mathfrak m_t$ generates the von Neumann algebra
$\mathfrak U$.

Traced representations of $G$ are identified with traced representations
of $C^{*}(G)$. Quasi-equivalence of two traced representations is defined
in the evident manner (cf. \XRef{6.6.2}{6.6.2}). With every traced
representation of $G$ there is canonically associated a trace of $G$.

\ResultHeading{17.1.3}{17.1.3.}
A trace $f$ of $G$ defines the objects
$\mathfrak m_f,\allowbreak\mathfrak n_f,\allowbreak N_f,\allowbreak
\Lambda_f,\allowbreak H_f,\allowbreak J_f,\allowbreak\lambda_f,\allowbreak
\rho_f,\allowbreak\mathfrak U_f,\allowbreak
\mathfrak V_f,\allowbreak t_f$
(\SectionRef{6.2}{Section~6.2} and \XRef{6.4.5}{6.4.5}). The
representation $\lambda_f$ of $C^{*}(G)$ is identified with a
representation of $G$, also denoted by $\lambda_f$. The representation
$\bar\rho_f^{\mathrm o}$ of $C^{*}(G)$ in $\bar H_f$ is the transform of
$\lambda_f$ by the isomorphism $J_f$ of $H_f$ onto $\bar H_f$
(\XRef{6.2.2}{6.2.2}); it is identified with a representation of $G$, also
denoted by $\bar\rho_f^{\mathrm o}$, which is the transform by $J_f$ of the
representation $\lambda_f$ of $G$. Thus
$\bar\rho_f^{\mathrm o}\simeq\lambda_f$. A trace of $G$ is said to be of
type I, type II, and so on, if $\lambda_f$ is of type I, type II, and so on.

The preceding discussion and \XRef{17.1.2}{17.1.2} define a bijection from
the set of traces of $G$ onto the set of quasi-equivalence classes of
traced representations of $G$.

\ResultHeading{17.1.4}{17.1.4.}
Let $\pi$ be a continuous unitary representation of $G$, and let
$\mathfrak U$ be the von Neumann algebra generated by $\pi(G)$. The
representation $\pi$ is called \emph{traceable} if there is a trace $t$ on
$\mathfrak U^{+}$ such that $(\pi,t)$ is a traced representation.
Traceable representations of $G$ are identified with traceable
representations of $C^{*}(G)$.

% Source PDF physical page 315
\ResultHeading{17.1.5}{17.1.5.}
Let $\pi$ be a continuous factor unitary representation of $G$, and let
$\mathfrak F$ be the factor generated by $\pi(G)$. The following
conditions are equivalent: (i) $\pi$ is traceable; (ii) $\mathfrak F$ is
semifinite and $\pi(C^{*}(G))$ contains a nonzero Hilbert--Schmidt operator
relative to $\mathfrak F$ (\XRef{6.7.2}{6.7.2}).

Let $f$ be a trace of $G$. The corresponding representation of $G$ is a
factor representation if and only if $f$ is a character
(\XRef{6.7.3}{6.7.3}). Hence there is a canonical bijection between:
\emph{a.} the set of characters of $G$ considered up to a positive scalar
factor; \emph{b.} the set of quasi-equivalence classes of traceable factor
representations of $G$ (\XRef{6.7.4}{6.7.4}).

\ResultHeading{17.1.6}{17.1.6.}
Suppose $G$ is postliminal. Every continuous irreducible unitary
representation $\pi$ of $G$ is traceable (\XRef{6.7.5}{6.7.5}) and has a
normalized character $f_\pi$, defined by
$f_\pi(x)=\operatorname{Tr}\pi(x)$ for $x\in C^{*}(G)^{+}$. If $\pi'$ is
another continuous irreducible unitary representation of $G$, then $\pi$
and $\pi'$ are equivalent if and only if $f_\pi=f_{\pi'}$
(\XRef{6.7.6}{6.7.6}). Every character of $G$ has the form
$\lambda f_\pi$, where $\lambda>0$ and $\pi$ is a continuous irreducible
unitary representation of $G$ (\XRef{6.7.6}{6.7.6}). Let $f=f_\pi$.
Then $H_f,\lambda_f,\bar\rho_f^{\mathrm o},\mathfrak U_f,\mathfrak V_f$
are canonically identified with
$H_\pi\otimes\overline{H_\pi}$, $\pi\otimes1$,
$1\otimes\bar\pi$, $\mathcal L(H_\pi)\otimes\mathbb C$, and
$\mathbb C\otimes\mathcal L(\overline{H_\pi})$, respectively;
$J_f$ is identified with the canonical involution of
$H_\pi\otimes\overline{H_\pi}$; and $t_f$ is identified with the trace
\[
T\otimes1\longmapsto\operatorname{Tr}T
\quad\text{on }(\mathcal L(H_\pi)\otimes\mathbb C)^{+}
\]
(\XRef{6.7.7}{6.7.7}).

Bibliography: \cite{ref54}.

\BookSubsection{17.2}{Character Defined by a Measure or a Distribution}

\ResultHeading{17.2.1}{17.2.1. Lemma.}
\begin{itshape}
Let $A$ be a $C^{*}$-algebra, let $A'$ be a dense involutive subalgebra of
$A$, and let $s$ be a complex-valued function on $A'\times A'$ satisfying
the following conditions:
\begin{enumerate}[label=\textup{(\roman*)}]
\item $s$ is a positive Hermitian sesquilinear form;
\item $s(y,x)=s(x^{\ast},y^{\ast})$ for $x,y\in A'$;
\item $s(zx,y)=s(x,z^{\ast}y)$ for $x,y,z\in A'$;
\item $s(zx,zx)\leq\|z\|^2s(x,x)$ for $x,z\in A'$;
\item the elements $xy$, where $x,y\in A'$, are dense in $A'$ for the
pre-Hilbert structure defined by $s$.
\end{enumerate}
Then there is a unique maximal bitrace of $A$ extending $s$.
\end{itshape}

Uniqueness follows from \XRef{6.5.3}{6.5.3}. We prove existence. Let $N$
be the set of $x\in A'$ such that $s(x,x)=0$. As in
\XRef{6.2.2}{6.2.2}, one verifies that $N$ is a self-adjoint two-sided
ideal of $A'$ and that $A'/N$ is a Hilbert algebra for the inner product
$(\,\cdot\mid\cdot\,)$ induced by $s$. Let $H$ be the Hilbert-space
completion of $A'/N$, and let $\Lambda$ be the canonical map of $A'$ into
$H$. Hypothesis (iv) implies that, for every $z\in A'$, there is a unique
bounded linear operator $\lambda(z)$ on $H$ such that
$\lambda(z)\Lambda x=\Lambda(zx)$ for every $x\in A'$; moreover
$\|\lambda(z)\|\leq\|z\|$. It is immediate that $\lambda$ is a
representation of $A'$
% Source PDF physical page 316
in $H$; it extends uniquely to a representation, still denoted by
$\lambda$, of $A$ in $H$. Then $\lambda(A)$ and $\lambda(A')$ generate the
same von Neumann algebra, namely $\mathfrak U(A'/N)$. Let $\theta$ be the
natural trace on $\mathfrak U(A'/N)^{+}$ defined by the Hilbert algebra
$A'/N$. Then $\mathfrak m_\theta\supset\lambda(A')$, so $(\lambda,\theta)$
is a traced representation of $A$. Let $s_1$ be the corresponding maximal
bitrace of $A$. For $x\in A'$, we have $\lambda(x)\in\mathfrak n_\theta$,
so $x\in\mathfrak n_{s_1}$, and
\[
s_1(x,x)=\theta\bigl(\lambda(x)\lambda(x)^{\ast}\bigr)=s(x,x).
\]
Thus $s_1$ extends $s$.

\ResultHeading{17.2.2}{17.2.2.}
If $G$ is a locally compact group, $\mathcal K(G)$ is an involutive
subalgebra of $L^1(G)$, dense in $L^1(G)$ and hence in $C^{*}(G)$.

\ResultLabel{Proposition.}
\begin{itshape}
Let $G$ be a unimodular locally compact group and $\mu$ a measure on $G$.
For $f,g\in\mathcal K(G)$, put
$s(f,g)=\mu(g^{\ast}\ast f)$. Then $s$ satisfies conditions (i)--(v) of
\XRef{17.2.1}{17.2.1} if and only if $\mu$ is a central measure of positive
type.
\end{itshape}

It is clear that $s$ is sesquilinear, so condition (i) of
\XRef{17.2.1}{17.2.1} holds if and only if
$\mu(f^{\ast}\ast f)\geq0$ for every $f\in\mathcal K(G)$, that is, if and
only if $\mu\gg0$. Suppose henceforth that this is so. Condition (ii)
holds if and only if
$\mu(g^{\ast}\ast f)=\mu(f\ast g^{\ast})$ for all
$f,g\in\mathcal K(G)$, that is, if and only if $\mu$ is central.
Condition (iii) reads
\[
\mu(g^{\ast}\ast(h\ast f))
=\mu((h^{\ast}\ast g)^{\ast}\ast f)
\]
for all $f,g,h\in\mathcal K(G)$; it holds automatically. Let
$f\in\mathcal K(G)$ and let $K$ be a compact neighborhood of the support
of $f$. There is a $g\in\mathcal K(G)$ such that $f\ast g$ has support in
$K$ and $\|f-f\ast g\|_\infty$ is arbitrarily small; it follows at once
that condition (v) holds. We saw in \XRef{13.7.9}{13.7.9} that $\mu$
defines a unitary representation $\sigma$ of $G$. The space $H_\sigma$ is
obtained precisely by furnishing $\mathcal K(G)$ with the inner product
$s(f,g)$ and passing to the separated completed Hilbert space. If $\Lambda$
denotes the canonical map of $\mathcal K(G)$ into $H_\sigma$, then
$\sigma(x)\Lambda f=\Lambda(\varepsilon_x\ast f)$ for every $x\in G$ and
$f\in\mathcal K(G)$. It follows that
$\sigma(g)\Lambda f=\Lambda(g\ast f)$ for $f,g\in\mathcal K(G)$, and hence
\begin{align*}
s(g\ast f,g\ast f)
 &=(\sigma(g)\Lambda f\mid\sigma(g)\Lambda f)\\
 &\leq\|\sigma(g)\|^2(\Lambda f\mid\Lambda f)
 \leq\|g\|_1^2s(f,f),
\end{align*}
so condition (iv) holds.

\ResultHeading{17.2.3}{17.2.3.}
Let $\mu$ be a central measure of positive type on the unimodular group
$G$. By \XRef{17.2.1}{17.2.1} and \XRef{17.2.2}{17.2.2}, the form
$(f,g)\mapsto\mu(g^{\ast}\ast f)$ on
$\mathcal K(G)\times\mathcal K(G)$ extends to a unique maximal bitrace
$\mu'$ of $C^{*}(G)$. Let $\sigma'$ be the representation of $C^{*}(G)$,
or of $G$, defined by $\mu'$, and retain the notation $\sigma$ of
\XRef{17.2.2}{17.2.2}. By \XRef{6.3.6}{6.3.6},
$H_{\sigma'}=H_\sigma$; moreover, for $f\in\mathcal K(G)$,
$\sigma'(f)$ is induced by the operator of left convolution by $f$ in
$\mathcal K(G)$, so $\sigma'(f)=\sigma(f)$. Thus $\sigma'=\sigma$.

% Source PDF physical page 317
In fact, $\mu'$ defines a traced representation $(\sigma,t)$ of $G$. For
$f\in\mathcal K(G)$, we have $\sigma(f)\in\mathfrak n_t$ and
\[
t\bigl(\sigma(f)\sigma(f)^{\ast}\bigr)
=\mu'(f,f)=\mu(f^{\ast}\ast f).
\]
The map $\mu\mapsto\mu'$ is injective, since knowledge of $\mu$ on the
functions $g\ast f$, $f,g\in\mathcal K(G)$, determines $\mu$ completely.

One sometimes identifies $\mu$, $\mu'$, and the associated trace, so it is
meaningful to say that a central measure $\mu$ of positive type on $G$ is a
character of $G$. By the preceding discussion, this is so if and only if
the continuous unitary representation of $G$ obtained from $\mu$ by the
procedure of \XRef{13.7.9}{13.7.9} is a factor representation. More
particularly, it is meaningful to say that a continuous central function
$\varphi$ of positive type on $G$ is a character [we identify $\varphi$
with the measure $\varphi(x)\,dx$].

\ResultHeading{17.2.4}{17.2.4.}
We spell out a special case important for applications.

\ResultLabel{Proposition.}
\begin{itshape}
Let $G$ be a unimodular postliminal group, let $\sigma$ be a continuous
irreducible unitary representation of $G$, and let $\mu$ be a central
measure of positive type on $G$. The representation $\sigma$ has normalized
character $\mu$ if and only if $\sigma(f)$ is a Hilbert--Schmidt operator
for every $f\in\mathcal K(G)$ and, for $f,g\in\mathcal K(G)$,
\[
\operatorname{Tr}\bigl(\sigma(g)^{\ast}\sigma(f)\bigr)
=\mu(g^{\ast}\ast f).
\]
If this is so, $\sigma$ is quasi-equivalent to the representation of $G$
obtained from $\mu$ by the procedure of \XRef{13.7.9}{13.7.9}.
\end{itshape}

\ResultHeading{17.2.5}{17.2.5.}
Suppose $G$ is unimodular. The Dirac measure $\varepsilon_e$ at $e$ is
central and of positive type. It defines a maximal bitrace $s$ of
$C^{*}(G)$. Then $\mathcal K(G)$ is dense in $\mathfrak n_s$ for the
pre-Hilbert structure of $\mathfrak n_s$; for
$f,g\in\mathcal K(G)$,
\[
s(f,g)=\int f(x)\overline{g(x)}\,dx.
\]
The space $H_s$ is $L^2(G)$, $\lambda_s$ is the left regular
representation, $\bar\rho_s^{\mathrm o}$ is the conjugate of the right
regular representation, and $J_s$ is the map $f\mapsto f^{\ast}$ of
$L^2(G)$ onto $\overline{L^2(G)}$. The completed Hilbert algebra of
$\mathfrak n_s/N_s$ is the Hilbert algebra $A$ of $G$; the von Neumann
algebra $\mathfrak U_s$ is $\mathfrak U(A)$, that is, the von Neumann
algebra generated by the left-translation operators on $L^2(G)$;
$\mathfrak U_s$ has a natural trace $t_s$, and $(\lambda_s,t_s)$ is a
traced representation of $G$. If $f\in L^1(G)\cap L^2(G)$, then
\[
t_s\bigl(\lambda_s(f)\lambda_s(f)^{\ast}\bigr)<+\infty;
\]
and, for $f,g\in L^1(G)\cap L^2(G)$,
\[
t_s\bigl(\lambda_s(f)\lambda_s(g)^{\ast}\bigr)
=\int f(x)\overline{g(x)}\,dx.
\]

% Source PDF physical page 318
\ResultHeading{17.2.6}{17.2.6.}
Let $G$ be a real Lie group. Denote by $\mathcal D(G)$ the set of
infinitely differentiable compactly supported complex functions on $G$.
This set is dense in $L^1(G)$, and hence in $C^{*}(G)$.

A distribution $\mu$ on $G$ is called \emph{central} if it is invariant
under the inner automorphisms of $G$, or, equivalently, if
$\mu(f\ast g)=\mu(g\ast f)$ for all $f,g\in\mathcal D(G)$. A distribution
$\mu$ on $G$ is said to be of \emph{positive type} if
$\mu(f^{\ast}\ast f)\geq0$ for every $f\in\mathcal D(G)$.

\ResultLabel{Proposition.}
\begin{itshape}
Let $G$ be a unimodular real Lie group and $\mu$ a distribution on $G$.
For $f,g\in\mathcal D(G)$, put
$s(f,g)=\mu(g^{\ast}\ast f)$. Then $s$ satisfies conditions (i)--(v) of
\XRef{17.2.1}{17.2.1} if and only if $\mu$ is a central distribution of
positive type.
\end{itshape}

The proof, entirely analogous to that of \XRef{17.2.2}{17.2.2}, is left to
the reader.

\ResultHeading{17.2.7}{17.2.7.}
Retain the hypotheses of the proposition in \XRef{17.2.6}{17.2.6}. By
\XRef{17.2.1}{17.2.1}, the form $s$ extends to a unique maximal bitrace
$\mu'$ of $C^{*}(G)$, which defines a traced representation $(\sigma,t)$
of $G$. The representation $\sigma$ is obtained by furnishing
$\mathcal D(G)$ with the inner product $s(f,g)$, passing to the separated
completed Hilbert space, and making $G$ act by left translation on
$\mathcal D(G)$. For $f\in\mathcal D(G)$, we have
$\sigma(f)\in\mathfrak n_t$ and
\[
t\bigl(\sigma(f)\sigma(f)^{\ast}\bigr)
=\mu'(f,f)=\mu(f^{\ast}\ast f).
\]
The map $\mu\mapsto\mu'$ is injective.

One sometimes identifies $\mu$, $\mu'$, and the associated trace, so it is
meaningful to say that a central distribution of positive type on $G$ is a
character of $G$. In particular:

\ResultHeading{17.2.8}{17.2.8. Proposition.}
\begin{itshape}
Let $G$ be a unimodular postliminal Lie group, let $\sigma$ be a continuous
irreducible unitary representation of $G$, and let $\mu$ be a central
distribution of positive type on $G$. The representation $\sigma$ has
normalized character $\mu$ if and only if $\sigma(f)$ is a
Hilbert--Schmidt operator for every $f\in\mathcal D(G)$ and, for
$f,g\in\mathcal D(G)$,
\[
\operatorname{Tr}\bigl(\sigma(g)^{\ast}\sigma(f)\bigr)
=\mu(g^{\ast}\ast f).
\]
If this is so, $\sigma$ is quasi-equivalent to the representation of $G$
obtained from $\mu$ by the procedure of \XRef{17.2.7}{17.2.7}.
\end{itshape}

Bibliography: \cite{ref54,ref115}.

\BookSubsection{17.3}{Characters of Finite Type}

\ResultHeading{17.3.1}{17.3.1. Proposition.}
\begin{itshape}
Let $G$ be a unimodular locally compact group, let $s$ be a character of
$G$, and let $(\pi,t)$ be a traced factor representation of $G$ associated
with $s$.

% Source PDF physical page 319
\begin{enumerate}[label=\textup{(\roman*)}]
\item The character $s$ is of finite type if and only if it is defined by a
continuous central function $\chi$ of positive type on $G$.
\item If this is so, then
$\chi(x)=t'(\pi(x))$, where $t'$ denotes the linear extension of $t$ to the
von Neumann algebra generated by $\pi(G)$.
\end{enumerate}
\end{itshape}

The character $s$ is of finite type if and only if it extends to a central
positive linear form on $C^{*}(G)$ (\XRef{6.8.1}{6.8.1} and
\XRef{6.8.5}{6.8.5}). The central positive linear forms on $C^{*}(G)$
correspond bijectively to the continuous central positive linear forms on
$L^1(G)$ (\XRef{2.7.5}{2.7.5}), and hence also to the continuous central
functions of positive type on $G$. Thus, if $s$ is of finite type, there is
a continuous central function $\chi$ of positive type on $G$ such that
\[
s(f\ast f^{\ast})=\langle\chi,f\ast f^{\ast}\rangle
\qquad(f\in\mathcal K(G));
\]
in this sense $s$ is defined by $\chi$. The converse is evident. Suppose
$s$ is defined by $\chi$. The function $x\mapsto t'(\pi(x))$ on $G$ is
continuous. [Indeed, $t'$ is normal, hence ultra-strongly continuous
(\XRef{A.25}{A 25}), and therefore strongly continuous on the unit ball
(\XRef{A.1}{A 1}).] For all $f,g\in\mathcal K(G)$, we have
\begin{align*}
\int t'(\pi(x))(g^{\ast}\ast f)(x)\,dx
 &=t'\left(\int\pi(x)(g^{\ast}\ast f)(x)\,dx\right)\\
 &=t'(\pi(g^{\ast}\ast f))
  =\int\chi(x)(g^{\ast}\ast f)(x)\,dx.
\end{align*}
Thus the measures $t'(\pi(x))\,dx$ and $\chi(x)\,dx$ are equal, and hence
$t'(\pi(x))=\chi(x)$ for every $x\in G$.

We shall identify the finite-type characters of $G$ with continuous central
functions of positive type on $G$.

\ResultHeading{17.3.2}{17.3.2.}
Suppose $G$ is compact. Every factor representation of $G$ is a multiple
of an irreducible representation $\pi$ (\XRef{15.1.8}{15.1.8}), and
$\dim\pi<+\infty$. By \XRef{17.3.1}{17.3.1}, every character of $G$ in
the sense of \XRef{17.1.1}{17.1.1} is identified with a function
$s\mapsto\lambda\operatorname{Tr}\pi(s)$ on $G$, where $\lambda>0$ and
$\pi$ is an irreducible representation of $G$. Up to the scalar factor
$\lambda$, this is precisely the definition in \XRef{15.3.1}{15.3.1}.

\ResultHeading{17.3.3}{17.3.3.}
Suppose $G$ is commutative. Every factor representation of $G$ is a
multiple of a one-dimensional representation. Thus the characters of $G$
in the sense of \XRef{17.1.1}{17.1.1} are identified with the functions
$\lambda\chi$, where $\lambda>0$ and $\chi$ is a character of $G$ in the
usual sense.

\ResultHeading{17.3.4}{17.3.4. Proposition.}
\begin{itshape}
Let $G$ be a unimodular locally compact group. There is a canonical
bijection between:
\begin{enumerate}[label=\textup{\alph*.}]
\item the set of quasi-equivalence classes of continuous finite-type factor
unitary representations of $G$;
\item the set of finite-type characters of $G$ equal to $1$ at $e$.
\end{enumerate}
\end{itshape}

This follows from \XRef{6.8.6}{6.8.6}.

% Source PDF physical page 320
\ResultHeading{17.3.5}{17.3.5. Proposition.}
\begin{itshape}
Let $G$ be a unimodular locally compact group, and let $C$ be the set of
continuous central functions $\varphi$ of positive type on $G$ such that
$\varphi(e)\leq1$.
\begin{enumerate}[label=\textup{(\roman*)}]
\item $C$ is convex and compact for the weak topology
$\sigma(L^\infty(G),L^1(G))$.
\item The extreme points of $C$ are $0$ and the finite-type characters of
$G$ equal to $1$ at $e$.
\item $C$ is the weakly closed convex hull of $0$ and of the set of
finite-type characters equal to $1$ at $e$.
\end{enumerate}
\end{itshape}

This follows from \XRef{6.8.7}{6.8.7}.

\ResultHeading{17.3.6}{17.3.6. Proposition.}
\begin{itshape}
Suppose that $G$ has a fundamental system of compact neighborhoods of $e$
invariant under inner automorphisms.
\begin{enumerate}[label=\textup{(\roman*)}]
\item For every $s\in G$, $s\ne e$, there is a finite-type character
$\chi$ of $G$ such that $\chi(s)\ne\chi(e)$.
\item For every $s\in G$, $s\ne e$, there is a continuous finite-type
factor unitary representation $\pi$ of $G$ such that $\pi(s)\ne1$.
\end{enumerate}
\end{itshape}

Let $s\in G$, $s\ne e$. There is a compact symmetric neighborhood $V$ of
$e$, invariant under inner automorphisms, such that $s\notin V^2$. Let $f$
be the characteristic function of $V$. Then
$g=f\ast f^{\ast}$ is a continuous central function of positive type on
$G$, equal to $0$ at $s$ and to
$\lambda=\|f\|_2^2>0$ at $e$. By \XRef{17.3.5}{17.3.5},
$\lambda^{-1}g$ is a weak limit of positive linear combinations $h_i$ of
finite-type characters such that $\|h_i\|_\infty\leq1$. Since
\[
\|\lambda^{-1}g\|_\infty\leq\liminf_i\|h_i\|_\infty,
\]
we may suppose that $\|h_i\|_\infty=1$ for every $i$; then $h_i$ tends to
$\lambda^{-1}g$ uniformly on every compact set (\XRef{13.5.2}{13.5.2}).
Thus $h_i(e)\ne h_i(s)$ for some $i$, proving (i).

Let $\chi$ be a finite-type character of $G$ such that
$\chi(s)\ne\chi(e)$, and let $(\pi,t)$ be the corresponding finite-type
traced factor representation. Denote by $t'$ the linear extension of $t$ to
the von Neumann algebra generated by $\pi(G)$. We have
\[
t'(\pi(s))=\chi(s)\ne\chi(e)=t'(\pi(e));
\]
hence $\pi(s)\ne\pi(e)=1$.

\ResultHeading{17.3.7}{17.3.7.}
It is not known whether the converse of \XRef{17.3.6}{17.3.6} is true.
However:

\ResultLabel{Proposition.}
\begin{itshape}
\begin{enumerate}[label=\textup{(\roman*)}]
\item Let $G$ be a locally compact group. Suppose that, for every
$s\in G$, $s\ne e$, there is a continuous finite-type factor unitary
representation $\pi$ of $G$ such that $\pi(s)\ne1$. Then:

$(\star)$ if $K$ is a compact subset of $G$ with $e\notin K$, there is a
neighborhood of $e$ disjoint from $K$ and invariant under inner
automorphisms.
\item If a locally compact group $G$ satisfies $(\star)$ and is generated
by a compact set, then $G$ has a fundamental system of compact
neighborhoods of $e$ invariant under inner automorphisms.
\end{enumerate}
\end{itshape}

% Source PDF physical page 321
Let $s$ be a point of $G$ distinct from $e$. There is a continuous
finite-type factor unitary representation $\pi$ of $G$ such that
$\pi(s)\ne1$. Let $t$ be the canonical trace on the factor generated by
$\pi(G)$. The function
\[
x\longmapsto\varphi(x)
=t\bigl[(\pi(x)^{\ast}-1)(\pi(x)-1)\bigr]
=t\bigl(2-\pi(x)-\pi(x^{-1})\bigr)
\]
on $G$ is continuous (cf. the proof of \XRef{17.3.1}{17.3.1}) and
central; $\varphi(e)=0$ and $\varphi(s)\ne0$. Hence there is a closed
neighborhood $V_s$ of $s$, invariant under inner automorphisms, such that
$e\notin V_s$. If $K$ is a compact subset of $G$ with $e\notin K$, then
$K$ can be covered by finitely many sets $V_{s_1},\ldots,V_{s_n}$, and the
complement of $V_{s_1}\cup\cdots\cup V_{s_n}$ is a neighborhood of $e$
disjoint from $K$ and invariant under inner automorphisms. This proves (i).
Assertion (ii) follows at once from the next lemma.

\ResultHeading{17.3.8}{17.3.8. Lemma.}
\begin{itshape}
Let $G$ be a locally compact group satisfying condition $(\star)$ of
\XRef{17.3.7}{17.3.7}. Let $V$ be a neighborhood of $e$ and $K$ a compact
subset of $G$. There is a neighborhood of $e$ contained in $V$ and
invariant under the inner automorphisms defined by the elements of $K$.
\end{itshape}

Shrinking $V$, we may suppose it open and relatively compact. Then,
enlarging $K$, we may suppose $V\subset K$ and $K=K^{-1}$. Put
$K'=K^3\cap(G\setminus V)$, which is compact. We have $e\notin K'$. By
$(\star)$ there is a neighborhood $W$ of $e$, disjoint from $K'$ and
invariant under inner automorphisms. We shall show, and this will complete
the proof, that if $s\in K$, then $V\cap W$ is invariant under the
automorphism $g\mapsto sgs^{-1}$ of $G$. Let $x\in V\cap W$. We have
$sxs^{-1}\in W$ and
$sxs^{-1}\in KVK^{-1}\subset K^3$. If $sxs^{-1}\notin V$, then
$sxs^{-1}\in K^3\cap(G\setminus V)=K'$, a contradiction. Hence
$sxs^{-1}\in V\cap W$.

\ResultHeading{17.3.9}{17.3.9.}
A connected group is generated by every neighborhood of the identity.
Thus, among locally compact groups having sufficiently many continuous
finite-type factor unitary representations, the only connected groups are,
by \XRef{17.3.7}{17.3.7} and \XRef{16.4.6}{16.4.6}, products of groups
$\mathbb R^n$ with compact groups. These are not very interesting examples,
since in fact the continuous irreducible unitary representations of these
groups are finite-dimensional (\XRef{13.1.8}{13.1.8}). But note that, by
\XRef{17.3.6}{17.3.6}, every discrete group has sufficiently many
continuous finite-type factor unitary representations.

Later, in \XRef{18.7.9}{18.7.9}, we shall obtain results concerning the
existence of characters that are not necessarily of finite type.

Bibliography: \cite{ref53,ref54}.

\BookSubsection{17.4}{Complements}

\ResultHeading{17.4.1}{17.4.1.}
\emph{a.} Let $G$ be a locally compact group having a fundamental system
$(V_i)$ of compact neighborhoods of $e$ invariant under inner
automorphisms. Let $f$ be a character of $G$ such that $\mathfrak m_f$ is
dense in $C^{*}(G)$.
% Source PDF physical page 322
Then $f$ is of finite type. [Let $\varphi_i$ be the characteristic function
of $V_i$. The operator $\lambda_f(\varphi_i)$ is scalar and nonzero for
some $i$. Since $\varphi_i$ is a limit in $C^{*}(G)$ of elements of
$\mathfrak m_f$, the operator $1$ is a norm limit of trace-class operators
relative to $\mathfrak U_f$. Thus $\mathfrak U_f$ is a finite factor.]
\cite{ref54}.

\emph{* b.} A countable discrete group may have characters that are not of
finite type \cite{ref57}.

\ResultHeading{17.4.2}{* 17.4.2.}
There are countable discrete antiliminal groups for which the space of
finite-type characters of norm $1$ is not locally compact \cite{ref57}.

\ResultHeading{17.4.3}{17.4.3.}
Let $G$ be a connected locally compact group, let $\pi$ be a continuous
unitary representation of $G$, and let $A$ be the von Neumann algebra
generated by $\pi(G)$. Then the largest projection $E$ in the center of
$A$ such that $A_E$ is of type $\mathrm{II}_1$ is $0$
\cite{ref72,ref140}.

\ResultHeading{17.4.4}{17.4.4.}
Let $G$ be a connected locally compact group having a family of unitary
representations $(\pi_i)$ with the following properties:
$1^\circ$ for every $s\in G$, there is an $i$ such that
$\pi_i(s)\ne1$; $2^\circ$ the maps $s\mapsto\pi_i(s)$ are
norm-continuous. Then $G$ is a product of a compact group and a group
$\mathbb R^n$. (Use \XRef{16.4.6}{16.4.6} and
\XRef{17.3.8}{17.3.8}.) \cite{ref72,ref143}.

\ResultHeading{17.4.5}{17.4.5.}
Let $G$ be a connected noncompact simple real Lie group. The only
continuous unitary representations $\pi$ of $G$ for which $\pi(G)$ is
contained in a finite factor are the trivial representations. This applies
in particular to finite-dimensional representations \cite{ref140}.

\ResultHeading{17.4.6}{* 17.4.6.}
Let $G$ be a connected semisimple real Lie group and $\pi$ a continuous
irreducible unitary representation of $G$. If $f$ is a compactly supported
element of $L^2(G)$, then $\pi(f)$ is a Hilbert--Schmidt operator. If $f$
is infinitely differentiable and compactly supported on $G$, then $\pi(f)$
is a trace-class operator, and $f\mapsto\operatorname{Tr}\pi(f)$ is a
distribution $\chi_\pi$ on $G$. Thus $\pi$ has a character, namely the
distribution $\chi_\pi$. Harish-Chandra has recently announced that
$\chi_\pi$ is defined by a locally integrable function on $G$, analytic on
the set of regular elements of $G$. For the classical complex simple
groups, $\chi_\pi$ has been calculated by Gelfand and Naimark
\cite{ref59,ref60,ref61,ref115}.

\ResultHeading{17.4.7}{* 17.4.7.}
Let $G$ be a simply connected nilpotent real Lie group and $\pi$ a
continuous irreducible unitary representation of $G$. If $f$ is infinitely
differentiable on $G$ and rapidly decreasing at infinity, then $\pi(f)$ is
a trace-class operator, and $f\mapsto\operatorname{Tr}\pi(f)$ is a
tempered distribution $\chi_\pi$ on $G$. Thus $\pi$ has a character,
namely the distribution $\chi_\pi$. In general this distribution is not
defined by a function, or even by a measure
\cite{ref6,ref7,ref8,ref53,ref81}.

\ResultHeading{17.4.8}{17.4.8.}
Let $G$ be a locally compact group in which $e$ has a fundamental system
of compact neighborhoods invariant under inner automorphisms. If $\mu$ is
a central measure of positive type on $G$, the corresponding unitary
representation of $G$ is of finite type \cite{ref53}.

% Source PDF physical page 323
\BookSection{18}{Dual of a Locally Compact Group}

\BookSubsection{18.1}{Definition of the Dual}

\ResultHeading{18.1.1}{18.1.1.}
Let $G$ be a locally compact group. There is a canonical bijection from
$C^{*}(G)^{\widehat{\ }}$ onto $\widehat G$ (\XRef{13.9.3}{13.9.3}). By
transporting the topology of $C^{*}(G)^{\widehat{\ }}$ under this
bijection, one obtains a topology on $\widehat G$.

\ResultLabel{Definition.}
\begin{itshape}
The topological space so obtained is called the dual of $G$.
\end{itshape}

Henceforth $\widehat G$ will denote this topological space. The spaces
$\widehat G$ and $C^{*}(G)^{\widehat{\ }}$ are often identified.

\ResultHeading{18.1.2}{18.1.2.}
The space $\widehat G$ is a locally quasi-compact Baire space
(\XRef{3.4.13}{3.4.13} and \XRef{3.3.8}{3.3.8}). If $G$ is discrete,
$C^{*}(G)$ has an identity, so $\widehat G$ is quasi-compact
(\XRef{3.1.8}{3.1.8}). If $G$ is separable, $\widehat G$ is separable
(\XRef{3.3.4}{3.3.4}). If $G$ is postliminal, $\widehat G$ has a dense
open locally compact subset (\XRef{4.4.5}{4.4.5}).

\ResultHeading{18.1.3}{18.1.3.}
Let $\pi$ be a continuous unitary representation of $G$, and let $S$ be a
set of continuous unitary representations of $G$. We say that $\pi$ is
\emph{weakly contained} in $S$ if $\pi$, regarded as a representation of
$C^{*}(G)$, is weakly contained in $S$, regarded as a set of
representations of $C^{*}(G)$.

We shall now translate \XRef{3.4.4}{3.4.4}, \XRef{3.4.9}{3.4.9}, and
\XRef{3.4.10}{3.4.10} into the language of groups. For this translation we
use: $1^\circ$ Proposition \XRef{2.7.5}{2.7.5}, which passes from
positive forms on $C^{*}(G)$ to continuous positive forms on $L^1(G)$;
$2^\circ$ Theorems \XRef{13.4.5}{13.4.5} and
\XRef{13.5.2}{13.5.2}, which pass from positive forms on $L^1(G)$ to
continuous functions of positive type on $G$. We obtain:

\ResultHeading{18.1.4}{18.1.4. Proposition.}
\begin{itshape}
Let $\pi$ be a continuous unitary representation of $G$, and let $S$ be a
set of continuous unitary representations of $G$. The following conditions
are equivalent:
\begin{enumerate}[label=\textup{(\roman*)}]
\item $\pi$ is weakly contained in $S$;
\item every function of positive type associated with $\pi$ is a uniform
limit on every compact set of sums of functions of positive type associated
with $S$.
\end{enumerate}
If $\pi$ has a cyclic vector $\xi$, these conditions are also equivalent to:

\textup{(ii$'$)} the function
$s\mapsto(\pi(s)\xi\mid\xi)$ is a uniform limit on every compact set of
sums of functions of positive type associated with $S$.
\end{itshape}

% Source PDF physical page 324
\ResultHeading{18.1.5}{18.1.5. Proposition.}
\begin{itshape}
Let $\pi\in\widehat G$ and $S\subset\widehat G$. The following conditions
are equivalent:
\begin{enumerate}[label=\textup{(\roman*)}]
\item $\pi\in\overline S$;
\item $\pi$ is weakly contained in $S$;
\item one of the nonzero functions of positive type associated with $\pi$
is a uniform limit on every compact set of functions of positive type
associated with $S$;
\item every function of positive type associated with $\pi$ is a uniform
limit on every compact set of functions of positive type associated with
$S$.
\end{enumerate}
\end{itshape}

\ResultHeading{18.1.6}{18.1.6.}
The equivalence (i) $\Leftrightarrow$ (iii) of \XRef{18.1.5}{18.1.5}
proves in particular that, if $G$ is commutative, the topology defined in
\XRef{18.1.1}{18.1.1} is the usual topology of the dual group of $G$.

\ResultHeading{18.1.7}{18.1.7. Definition.}
\begin{itshape}
If $\pi$ is a continuous unitary representation of $G$, the \emph{support}
of $\pi$ is the set of $\sigma\in\widehat G$ that are weakly contained in
$\pi$. If $\varphi$ is a continuous function of positive type on $G$, the
\emph{analyzer} of $\varphi$ is the support of $\pi_\varphi$.
\end{itshape}

If $S$ is the support of $\pi$, then $S$ is closed in $\widehat G$, and
$\pi$ is weakly contained in $S$ (\XRef{3.4.6}{3.4.6}).

\ResultHeading{18.1.8}{18.1.8. Proposition.}
\begin{itshape}
Let $\varphi$ be a continuous function of positive type on $G$, and let
$A\subset\widehat G$ be its analyzer.
\begin{enumerate}[label=\textup{(\roman*)}]
\item $\varphi$ is a uniform limit on every compact set of sums of pure
functions of positive type associated with the elements of $A$.
\item Every function of positive type associated with an element of $A$ is
a uniform limit on every compact set of functions
\[
s\longmapsto\sum_{i,j=1}^{n}\lambda_i\overline{\lambda_j}\,
\varphi(s_j^{-1}ss_i),
\]
where $s_1,\ldots,s_n\in G$ and $\lambda_1,\ldots,\lambda_n\in\mathbb C$.
\end{enumerate}
\end{itshape}

Since $\pi_\varphi$ is weakly contained in $A$, (i) follows from
\XRef{18.1.4}{18.1.4}. Let $\pi\in A$. If $\pi$ and $\pi_\varphi$ are
regarded as representations of $C^{*}(G)$, then
$\ker\pi\supset\ker\pi_\varphi$; hence every state associated with $\pi$
is a weak limit of states associated with $\pi_\varphi$
[\XRef{3.4.2}{3.4.2 (ii)}]. Assertion (ii) now follows from
\XRef{13.4.10}{13.4.10}.

\ResultHeading{18.1.9}{18.1.9.}
Let $n$ be a cardinal, let $H_n$ be the standard Hilbert space of dimension
$n$, and let $\operatorname{Rep}_n(G)$ be the set of continuous unitary
representations of $G$ in $H_n$. There is a canonical bijection from
$\operatorname{Rep}_n(G)$ onto $\operatorname{Rep}'_n(C^{*}(G))$, the set
of nondegenerate representations of $C^{*}(G)$ in $H_n$. The inverse image
under this bijection of the topology of
$\operatorname{Rep}'_n(C^{*}(G))$ (\XRef{3.5.2}{3.5.2}) is a certain
topology on $\operatorname{Rep}_n(G)$, which we now describe directly. By
\XRef{3.5.8}{3.5.8}, this will give a new definition of the topology of
$\widehat G_n$, the subspace of $\widehat G$ formed by those
$\sigma\in\widehat G$ for which $\dim\sigma=n$.

% Source PDF physical page 325
\ResultLabel{Proposition.}
\begin{itshape}
Let $\pi_i,\pi\in\operatorname{Rep}_n(G)$, and let $\pi_i',\pi'$ be the
corresponding elements of $\operatorname{Rep}'_n(C^{*}(G))$. The following
conditions are equivalent:
\begin{enumerate}[label=\textup{(\roman*)}]
\item $\pi_i'\to\pi'$;
\item for every $\xi\in H_n$ and every compact subset $K$ of $G$,
\[
\|\pi_i(s)\xi-\pi(s)\xi\|\longrightarrow0
\quad\text{uniformly on }K;
\]
\item for every $\xi,\eta\in H_n$ and every compact subset $K$ of $G$,
\[
\bigl|(\pi_i(s)\xi\mid\eta)-(\pi(s)\xi\mid\eta)\bigr|
\longrightarrow0
\quad\text{uniformly on }K.
\]
\end{enumerate}
\end{itshape}

(ii) $\Rightarrow$ (iii): This is evident.

(iii) $\Rightarrow$ (i): Suppose (iii) holds. For every
$f\in\mathcal K(G)$,
\[
\bigl|(\pi_i'(f)\xi\mid\eta)-(\pi'(f)\xi\mid\eta)\bigr|
\leq\int|f(s)|\,
\bigl|(\pi_i(s)\xi\mid\eta)-(\pi(s)\xi\mid\eta)\bigr|\,ds,
\]
and the right-hand side tends to zero. Since $\mathcal K(G)$ is dense in
$C^{*}(G)$, $\pi_i'\to\pi'$ by \XRef{3.5.4}{3.5.4}.

(i) $\Rightarrow$ (ii): Suppose (i) holds, and prove (ii). It is enough to
do so when $\xi$ has the form $\pi(f)\eta$, with $f\in L^1(G)$ and
$\eta\in H_n$, since vectors of this form are dense in $H_n$. Now
\begin{align*}
\|\pi_i(s)\pi(f)\eta-\pi(s)\pi(f)\eta\|
&\leq\|\pi_i(s)\pi_i(f)\eta-\pi(s)\pi(f)\eta\|\\
&\quad+\|\pi_i(s)\pi(f)\eta-\pi_i(s)\pi_i(f)\eta\|\\
&=\|\pi_i({}_{s^{-1}}f)\eta-\pi({}_{s^{-1}}f)\eta\|
  +\|\pi_i(f)\eta-\pi(f)\eta\|.
\end{align*}
The set of maps $g\mapsto\pi_i(g)\eta$ from $L^1(G)$ into $H_n$ is
equicontinuous for fixed $\eta$. On the compact set of translates
${}_sf$, $s\in K$, pointwise convergence of these maps is therefore
equivalent to uniform convergence. Thus
$\|\pi_i({}_{s^{-1}}f)\eta-\pi({}_{s^{-1}}f)\eta\|\to0$ uniformly on $K$.
On the other hand, $\|\pi_i(f)\eta-\pi(f)\eta\|\to0$.

\ResultHeading{18.1.10}{18.1.10.}
We denote by $\operatorname{Irr}_n(G)$ the set of continuous irreducible
unitary representations of $G$ in $H_n$. It is furnished with the topology
induced by that of $\operatorname{Rep}_n(G)$. If $G$ is separable,
$\operatorname{Rep}_n(G)$ and $\operatorname{Irr}_n(G)$ are Polish spaces
(\XRef{3.7.4}{3.7.4} and \XRef{3.7.5}{3.7.5}).

Bibliography: \cite{ref27,ref50,ref91}.

\BookSubsection{18.2}{Fourier Transform}

\ResultHeading{18.2.1}{18.2.1. Definition.}
\begin{itshape}
Let $\mu\in M^1(G)$. For every $\pi\in\widehat G$, put
$(\mathcal F\mu)(\pi)=\pi(\mu)$, which is a bounded linear operator on
$H_\pi$. The function $\pi\mapsto(\mathcal F\mu)(\pi)$ on $\widehat G$
is called the Fourier transform of $\mu$.
\end{itshape}

% Source PDF physical page 326
Recall that the space of $\pi$ is defined only up to isomorphism, which is
rather inconvenient.

If $G$ is commutative, $(\mathcal F\mu)(\pi)$ is scalar multiplication by
$\int\pi(s)\,d\mu(s)$ on the one-dimensional space $H_\pi$, and the usual
definition of the Fourier transform is recovered by identifying this
scalar multiplication with its ratio.

\ResultHeading{18.2.2}{18.2.2.}
We have immediately
\begin{align*}
\mathcal F(\alpha\mu)&=\alpha\mathcal F\mu
 &&(\alpha\in\mathbb C,\ \mu\in M^1(G)),\\
\mathcal F(\mu+\mu')&=\mathcal F\mu+\mathcal F\mu'
 &&(\mu,\mu'\in M^1(G)),\\
\mathcal F(\mu\ast\mu')&=(\mathcal F\mu)(\mathcal F\mu')
 &&(\mu,\mu'\in M^1(G)),\\
\mathcal F(\mu^{\ast})&=(\mathcal F\mu)^{\ast}
 &&(\mu\in M^1(G)),\\
\sup_{\pi\in\widehat G}\|(\mathcal F\mu)(\pi)\|&\leq\|\mu\|
 &&(\mu\in M^1(G)).
\end{align*}

\ResultHeading{18.2.3}{18.2.3.}
The Fourier transform is injective. Indeed, let $\mu$ be a nonzero element
of $M^1(G)$. There is an $f\in\mathcal K(G)$ such that
$g=f\ast\mu$ is a nonzero element of $L^1(G)$. Hence there is a
$\pi\in\widehat G$ such that
$\pi(f)\pi(\mu)=\pi(f\ast\mu)\ne0$ (\XRef{2.7.3}{2.7.3}), and therefore
$(\mathcal F\mu)(\pi)=\pi(\mu)\ne0$.

\ResultHeading{18.2.4}{18.2.4.}
If $f\in L^1(G)$, we denote by $\mathcal Ff$ the Fourier transform of the
measure $f(s)\,ds$. In other words,
$(\mathcal Ff)(\pi)=\pi(f)$. If $\varepsilon>0$, the function
$\pi\mapsto\|(\mathcal Ff)(\pi)\|$ is less than $\varepsilon$ outside a
quasi-compact subset of $\widehat G$ (\XRef{3.3.7}{3.3.7}).

Bibliography: \cite{ref22,ref23,ref24,ref53,ref83,ref101,ref111,ref137,
ref145,ref147,ref148,ref162,ref171}.

\BookSubsection{18.3}{Reduced Dual}

\ResultHeading{18.3.1}{18.3.1. Definition.}
\begin{itshape}
The reduced dual of $G$ is the support of the regular representation of
$G$.
\end{itshape}

This reduced dual is a closed subset $\widehat G_r$ of $\widehat G$.

\ResultHeading{18.3.2}{18.3.2.}
Let $\lambda$ be the left regular representation of $G$. Regarded as a
representation of $C^{*}(G)$, it has a certain kernel $N$. Then
$\widehat G_r$ is the set of $\sigma\in\widehat G=C^{*}(G)^{\widehat{\ }}$
whose kernel contains $N$. In other words, $\widehat G_r$ is the spectrum
of the $C^{*}$-algebra $C^{*}(G)/N$, itself isomorphic to the
$C^{*}$-algebra $\lambda(C^{*}(G))$. This $C^{*}$-algebra
$\lambda(C^{*}(G))$ is the norm closure in $\mathcal L(L^2(G))$ of the set
of operators of left convolution by elements of $L^1(G)$.

Recall that $N\cap L^1(G)=0$ (\XRef{13.3.6}{13.3.6}).

\ResultHeading{18.3.3}{18.3.3.}
Suppose $G$ is unimodular. Every square-integrable irreducible
representation $\pi$ of $G$ is contained in the regular representation,
and therefore belongs to $\widehat G_r$. In particular, if $G$ is compact,
$\widehat G=\widehat G_r$.

% Source PDF physical page 327
\ResultHeading{18.3.4}{18.3.4.}
Suppose $G$ is commutative. Let $f\in L^1(G)$, let $\chi$ be a character
of $G$, and let $\lambda$ be the regular representation of $G$. We have
$|\chi(f)|\leq\|\lambda(f)\|$ (\XRef{13.3.6}{13.3.6}), and this inequality
extends by continuity to $f\in C^{*}(G)$. Thus $\chi$ is weakly contained
in $\lambda$, so $\widehat G_r=\widehat G$.

For other examples, see \XRef{18.3.9}{18.3.9},
\XRef{18.9.8}{18.9.8}, and \XRef{18.9.11}{18.9.11}.

\ResultHeading{18.3.5}{18.3.5. Proposition.}
\begin{itshape}
\emph{a.} Let $\varphi$ be a continuous function of positive type on $G$.
The following conditions are equivalent:
\begin{enumerate}[label=\textup{(\roman*)}]
\item $\varphi$ is a uniform limit on every compact set of functions of the
form $f\ast\widetilde f$, where $f\in\mathcal K(G)$;
\item $\varphi$ is a uniform limit on every compact set of continuous
compactly supported functions of positive type;
\item $\varphi$ is a uniform limit on every compact set of continuous
square-integrable functions of positive type;
\item $\varphi$ is associated with a representation weakly contained in
the regular representation.
\end{enumerate}
\emph{b.} The functions satisfying (i)--(iv) form a weakly closed convex
cone $Q$ in $L^\infty(G)$.

\emph{c.} If $\varphi\in Q$ and $\psi$ is a continuous function of positive
type, then $\varphi\psi\in Q$.
\end{itshape}

(i) $\Rightarrow$ (ii) $\Rightarrow$ (iii): This is evident.

(iii) $\Rightarrow$ (i): It is enough to show that if $\varphi$ is a
continuous square-integrable function of positive type, then $\varphi$ is a
uniform limit on $G$ of functions $f\ast\widetilde f$ with
$f\in\mathcal K(G)$. There is a $\psi\in L^2(G)$ such that
$\varphi=\psi\ast\widetilde\psi$ (\XRef{13.8.6}{13.8.6}). Let $(f_n)$ be
a sequence in $\mathcal K(G)$ tending to $\psi$ in $L^2(G)$. Then
$f_n\ast\widetilde f_n$ tends uniformly to
$\psi\ast\widetilde\psi=\varphi$.

(i) $\Rightarrow$ (iv): Let $\lambda$ be the left regular representation
of $G$. If $f\in\mathcal K(G)$, then $f\ast\widetilde f$ is a function of
positive type associated with $\lambda$ (\XRef{13.4.11}{13.4.11}). If
$\varphi$ satisfies (i), then $\pi_\varphi$ is weakly contained in
$\lambda$ (\XRef{18.1.4}{18.1.4}).

(iv) $\Rightarrow$ (ii): If $\varphi$ satisfies (iv), then $\varphi$ is a
uniform limit on every compact set of sums of functions of positive type
associated with $\lambda$ (\XRef{18.1.4}{18.1.4}). A function of positive
type associated with $\lambda$ has the form
$\psi\ast\widetilde\psi$, with $\psi\in L^2(G)$, and hence is a uniform
limit on $G$ of functions of the form $f\ast\widetilde f$, with
$f\in\mathcal K(G)$. A sum of such functions is a uniform limit on $G$ of
continuous compactly supported functions of positive type.

This proves a.

It is clear that the functions satisfying (i)--(iv) form a convex cone $Q$
in $L^\infty(G)$. Let $Q_1$ be the intersection of $Q$ with the unit ball
of $L^\infty(G)$. To prove that $Q$ is weakly closed, it is enough to prove
that $Q_1$ is weakly closed (\XRef{B.8}{B 8}). Let $\varphi$ be a function
in $L^\infty(G)$ weakly adherent to $Q_1$, and let $\omega$ be the positive
form on $C^{*}(G)$ defined by $\varphi$. Then $\omega$ is a weak limit of
positive forms on $C^{*}(G)$ defined by elements
% Source PDF physical page 328
of $Q_1$ [\XRef{2.7.5}{2.7.5 (iii)}]. Thus $\pi_\omega$ is weakly
contained in $\lambda$ (\XRef{3.4.9}{3.4.9}), and hence
$\varphi\in Q_1$.

Finally, c follows from characterization (ii) of the elements of $Q$ and
from the fact that the product of two continuous functions of positive type
is of positive type (\XRef{13.4.9}{13.4.9}).

\ResultHeading{18.3.6}{18.3.6. Proposition.}
\begin{itshape}
Let $G$ be a locally compact group. The following conditions are
equivalent:
\begin{enumerate}[label=\textup{(\roman*)}]
\item the dual of $G$ equals the reduced dual;
\item[\textup{(i$'$)}] the one-dimensional trivial representation of $G$
belongs to the reduced dual;
\item[\textup{(ii)}] every continuous function of positive type on $G$ is
a uniform limit on every compact set of functions of the form
$f\ast\widetilde f$, where $f\in\mathcal K(G)$;
\item[\textup{(ii$'$)}] the function $1$ is a uniform limit on every compact
set of functions of the form $f\ast\widetilde f$, where
$f\in\mathcal K(G)$;
\item[\textup{(iii)}] for every bounded measure $\mu$ of positive type on
$G$ and every continuous function $\varphi$ of positive type on $G$,
\[
\int\varphi(x)\,d\mu(x)\geq0;
\]
\item[\textup{(iii$'$)}] for every bounded measure $\mu$ of positive type
on $G$,
\[
\int d\mu(x)\geq0.
\]
\end{enumerate}
\end{itshape}

The equivalence (i) $\Leftrightarrow$ (ii) follows from
\XRef{18.3.5}{18.3.5}. The definition of measures of positive type shows
that (ii) $\Rightarrow$ (iii). Suppose (iii) holds. Let $\varphi$ be a
continuous function of positive type on $G$, and use the notation $Q$ of
\XRef{18.3.5}{18.3.5}. To show that $\varphi\in Q$, it is enough to prove
that $\varphi$ belongs to the bipolar of $Q$. Let $f\in L^1(G)$ belong to
the polar of $Q$. Since $Q$ is a cone,
$\operatorname{Re}\langle f,q\rangle\leq0$ for every $q\in Q$, and in
particular
$\operatorname{Re}\langle f,g\ast\widetilde g\rangle\leq0$ for every
$g\in\mathcal K(G)$. Thus the measure $-(f+f^{\ast})(s)\,ds$ is of
positive type. By (iii),
$\operatorname{Re}\int f(s)\varphi(s)\,ds\leq0$, which proves that
$\varphi$ belongs to the bipolar of $Q$. Hence (iii) $\Rightarrow$ (ii).

The equivalences (i$'$) $\Leftrightarrow$ (ii$'$)
$\Leftrightarrow$ (iii$'$) are established similarly. Clearly
(iii) $\Rightarrow$ (iii$'$). Finally, if $1\in Q$, every continuous
function of positive type belongs to $Q$ by \XRef{18.3.5}{18.3.5 (c)}.

\ResultHeading{18.3.7}{18.3.7. Lemma.}
\begin{itshape}
If $G$ is a semidirect product of two groups $H$ and $K$ satisfying the
conditions of \XRef{18.3.6}{18.3.6}, then $G$ also satisfies these
conditions.
\end{itshape}

Let $C$ be a compact subset of $G$. There are compact subsets $C'$ and
$C''$ of $H$ and $K$, respectively, such that $C\subset C'C''$, and we may
suppose that $e\in C'$. Let $\varepsilon>0$. There is an
$x\in\mathcal K(H)$ such that
$|x\ast\widetilde x-1|\leq\varepsilon$ on $C'$. Let $C'_0$ be the support
of $x$. The set $C''_0$ of elements $h^{-1}kh$, where
$h\in C'_0$ and $k\in C''$, is compact in the
% Source PDF physical page 329
normal subgroup $K$. There is a $y\in\mathcal K(K)$ such that
$|y\ast\widetilde y-1|\leq\varepsilon$ on $C''_0$. Define
$z\in\mathcal K(G)$ by
\[
z(hk)=x(h)y(k)\qquad(h\in H,\ k\in K).
\]
For $h'\in H$ and $k'\in K$, we have
\begin{align*}
(z\ast\widetilde z)(h'k')
 &=\int z(h'k'g)\overline{z(g)}\,dg\\
 &=\iint_{H\times K}x(h'h)y((h^{-1}k'h)k)
       \overline{x(h)}\,\overline{y(k)}\,dh\,dk\\
 &=\int_H(y\ast\widetilde y)(h^{-1}k'h)
       x(h'h)\overline{x(h)}\,dh.
\end{align*}
Consequently,
\begin{align*}
\bigl|(z\ast\widetilde z)(h'k')-1\bigr|
&\leq\left|\int_H(y\ast\widetilde y)(h^{-1}k'h)
  x(h'h)\overline{x(h)}\,dh
  -\int_Hx(h'h)\overline{x(h)}\,dh\right|\\
&\quad+\left|\int_Hx(h'h)\overline{x(h)}\,dh-1\right|\\
&\leq\int_{C'_0}\bigl|(y\ast\widetilde y)(h^{-1}k'h)-1\bigr|
  |x(h'h)|\,|x(h)|\,dh
  +\bigl|(x\ast\widetilde x)(h')-1\bigr|.
\end{align*}
If $h'k'\in C$, then $h'\in C'$ and $k'\in C''$, so
$h^{-1}k'h\in C''_0$ for $h\in C'_0$. Therefore
\begin{align*}
\bigl|(z\ast\widetilde z)(h'k')-1\bigr|
&\leq\varepsilon\int|x(h'h)|\,|x(h)|\,dh+\varepsilon\\
&\leq\varepsilon\int|x(h)|^2\,dh+\varepsilon\\
&=\varepsilon\bigl[(x\ast\widetilde x)(e)+1\bigr]
 \leq\varepsilon(2+\varepsilon).
\end{align*}

\ResultHeading{18.3.8}{18.3.8. Lemma.}
\begin{itshape}
Let $G$ be a locally compact group satisfying the conditions of
\XRef{18.3.6}{18.3.6}, and let $N$ be a closed normal subgroup of $G$.
Then $G/N$ satisfies the conditions of \XRef{18.3.6}{18.3.6}.
\end{itshape}

Let $G^{*}=G/N$, and let $\varphi$ be the canonical homomorphism of $G$
onto $G^{*}$. Let $C^{*}$ be a compact subset of $G/N$ and let
$\varepsilon>0$. There is a compact subset $C$ of $G$ such that
$\varphi(C)\supset C^{*}$, and we may suppose that $e\in C$. Then there is
an $f\in\mathcal K(G)$ such that
$|f\ast\widetilde f-1|\leq\varepsilon$ on $C$ and
$(f\ast\widetilde f)(e)=1$, that is, $\|f\|_2=1$. Let $F$ be the function
on $G^{*}$ defined by
\[
F(s^{*})=\left(\int_N|f(sn)|^2\,dn\right)^{1/2}.
\tag{1}\label{eq:18.3.8-1}
\]
% ===== ASSEMBLED TRANSLATION CHUNK 12: chunk_330_359.tex =====
% Source PDF physical page 330
where $dn$ denotes a left Haar measure on $N$, and where
$s^{*}=\varphi(s)$ [the right-hand side of \eqref{eq:18.3.8-1} plainly
depends only on $s^{*}$]. Then $F\in\mathcal K(G^{*})$ and
\begin{equation}
\|F\|_2^2
=\int_{G^{*}}ds^{*}\int_N|f(sn)|^2\,dn
=\int_G|f(s)|^2\,ds=1.
\tag{2}\label{eq:18.3.8-2}
\end{equation}
In view of \eqref{eq:18.3.8-2}, for every $s'\in G$ we have
\begin{align*}
\bigl|(F\ast\widetilde F)(s'^{*})-1\bigr|^2
&=\left|\int_{G^{*}}F(s'^{*}s^{*})\overline{F(s^{*})}\,ds^{*}-1\right|^2\\
&=\left|\int_{G^{*}}
 \bigl(F(s'^{*}s^{*})-F(s^{*})\bigr)\overline{F(s^{*})}\,ds^{*}\right|^2\\
&\leq\int_{G^{*}}|F(s'^{*}s^{*})-F(s^{*})|^2\,ds^{*}\\
&=\int_{G^{*}}\left|
 \left(\int_N|f(s'sn)|^2\,dn\right)^{1/2}
 -\left(\int_N|f(sn)|^2\,dn\right)^{1/2}\right|^2ds^{*}.
\end{align*}
Applying in $L^2(N)$ the inequality
$(\|a\|-\|b\|)^2\leq\|a-b\|^2$, the last expression is at most
\begin{align*}
\int_{G^{*}}ds^{*}\int_N|f(s'sn)-f(sn)|^2\,dn
 &=\int_G|f(s's)-f(s)|^2\,ds\\
 &=2-2\operatorname{Re}\int_Gf(s's)\overline{f(s)}\,ds\\
 &=2\bigl[1-\operatorname{Re}(f\ast\widetilde f)(s')\bigr].
\end{align*}
Thus, if $s'\in C$, then
$|(F\ast\widetilde F)(s'^{*})-1|^2\leq2\varepsilon$. Consequently
$|F\ast\widetilde F-1|^2\leq2\varepsilon$ on $C^{*}$.

\ResultHeading{18.3.9}{18.3.9. Proposition.}
\begin{itshape}
Let $G$ be a connected real Lie group and $R$ its radical. If $G/R$ is
compact, then $\widehat G=\widehat G_r$.
\end{itshape}

Let $G'$ be the universal covering group of $G$; it is the semidirect
product of a compact group $C$ and a simply connected solvable group $S$.
Since $G$ is a quotient of $G'$, it is enough to prove that
$\widehat{G'}=\widehat{G'}_r$ (\XRef{18.3.8}{18.3.8}). Now
$\widehat C=\widehat C_r$ (\XRef{18.3.3}{18.3.3}), and by
\XRef{18.3.7}{18.3.7} we are reduced to showing that a simply connected
solvable group $S$ satisfies $\widehat S=\widehat S_r$. We argue by
induction on the dimension $n$ of $S$. The assertion is evident for
$n=0$. Let $n>0$, and suppose the assertion established when
$\dim S<n$. Then $S$ is the semidirect product of a group isomorphic to
$\mathbb R$ and a simply connected solvable group of dimension $n-1$;
it now suffices to apply \XRef{18.3.4}{18.3.4},
\XRef{18.3.7}{18.3.7}, and the induction hypothesis.

Bibliography: \cite{ref4,ref27,ref154}.

% Source PDF physical page 331
\BookSubsection{18.4}{Reduced Dual and Integrable Representations}

\ResultHeading{18.4.1}{18.4.1. Proposition.}
\begin{itshape}
Let $G$ be a locally compact unimodular group and $\sigma$ a
square-integrable irreducible representation of $G$. Then $\{\sigma\}$ is
closed in $\widehat G_r$.
\end{itshape}

For $f\in\mathcal K(G)$, $\sigma(f)$ is a Hilbert--Schmidt operator
(\XRef{14.4.2}{14.4.2}), and therefore
$\sigma(C^{*}(G))\subset\mathcal{LC}(H_\sigma)$. Consequently the kernel of
$\sigma$ in $C^{*}(G)$ is a maximal closed two-sided ideal $I$ of
$C^{*}(G)$, and every irreducible representation of $C^{*}(G)$ having
kernel $I$ is equivalent to $\sigma$ (\XRef{4.1.10}{4.1.10} and
\XRef{4.1.11}{4.1.11}); hence $\{I\}$ is closed in
$\operatorname{Prim}C^{*}(G)$ (\XRef{3.1.4}{3.1.4}), and
$\{\sigma\}$ is closed in $(C^{*}(G))^{\widehat{\ }}=\widehat G$.
Moreover, $\sigma\in\widehat G_r$ (\XRef{18.3.3}{18.3.3}).

\ResultHeading{18.4.2}{18.4.2. Proposition.}
\begin{itshape}
Let $G$ be a locally compact unimodular group and $\sigma$ an integrable
irreducible representation of $G$. Then $\{\sigma\}$ is open and closed in
$\widehat G_r$.
\end{itshape}

We already know that $\{\sigma\}$ is closed in $\widehat G_r$
(\XRef{18.4.1}{18.4.1}). Let $\lambda$ be the left regular representation
of $G$; then $\lambda(C^{*}(G))=B$ is a $C^{*}$-algebra of operators on
$L^2(G)$ whose spectrum is $\widehat G_r$ (\XRef{18.3.2}{18.3.2}). We may
identify $\sigma$ with a subrepresentation of $\lambda$, since $\sigma$ is
square-integrable. Let $K\supset H_\sigma$ be the bi-invariant subspace of
$L^2(G)$ associated with $\sigma$, and let $\lambda_1$ be the
subrepresentation of $\lambda$ defined by $K$. For every $f\in L^1(G)$,
$\lambda_1(f)$ is the restriction of $\lambda(f)$ to $K$; thus, if
$\lambda_1$ is identified with a representation of $B$, it is simply the
map $T\mapsto T|K$ from $B$ into $\mathcal L(K)$. The kernel of $\sigma$ in
$B$ equals that of $\lambda_1$, since
$\lambda_1\simeq(\dim\sigma)\cdot\sigma$
[\XRef{14.3.5}{14.3.5 (iii)}]; this kernel is therefore the set $I$ of
$T\in B$ such that $T|K=0$. Let $J$ be the set of $T\in B$ such that
$T|\bigl(L^2(G)\ominus K\bigr)=0$. Let $\xi$ be a nonzero vector of
$H_\sigma$ such that the function
$s\mapsto\varphi(s)=(\lambda(s)\xi\mid\xi)$ on $G$ is integrable. We have
$\check\varphi\in K$ (\XRef{14.3.1}{14.3.1}). Since
$\check\varphi\in L^1(G)$, we can form $\lambda(\check\varphi)$, which is a
nonzero element of $B$. For every $g\in L^2(G)$,
$\lambda(\check\varphi)g=\check\varphi\ast g\in K$, because $K$ is
bi-invariant; hence
$\lambda(\check\varphi)|\bigl(L^2(G)\ominus K\bigr)=0$. It follows that
$J\ne0$.

Now let $\tau\in\widehat G_r$, regarded as an irreducible representation
of $B$. Its kernel in $B$ is a primitive ideal containing
$0=I\cap J$, and hence contains $I$ or $J$ (\XRef{2.11.4}{2.11.4}). If
$\ker\tau\supset I$, then $\tau$ is adherent to $\{\sigma\}$ and thus is
equal to $\sigma$. This proves that
\[
\bigcap_{\tau\in\widehat G_r,\,\tau\ne\sigma}\ker\tau\supset J\ne0.
\]
Thus $\widehat G_r\setminus\{\sigma\}$ is not dense in $\widehat G_r$.
Consequently $\{\sigma\}$ contains a nonempty open subset of
$\widehat G_r$, and $\{\sigma\}$ is open in $\widehat G_r$.

\ResultHeading{18.4.3}{18.4.3. Corollary.}
\begin{itshape}
If $G$ is compact, $\widehat G$ is discrete.
\end{itshape}

Since every irreducible representation of $G$ is integrable, this follows
from \XRef{18.4.2}{18.4.2}.

Bibliography: \cite{ref13,ref50}.

% Source PDF physical page 332
\BookSubsection{18.5}{Mackey Borel Structure}

\ResultHeading{18.5.1}{18.5.1.}
Let $G$ be a separable locally compact group. For every cardinal
$n\leq\aleph_0$, we equip $\operatorname{Rep}_n(G)$ and
$\operatorname{Irr}_n(G)$ with the Borel structures underlying their
topologies. With the notation of \XRef{18.1.9}{18.1.9}, the bijections
\[
\operatorname{Rep}_n(G)\longrightarrow\operatorname{Rep}'_n(C^{*}(G)),
\qquad
\operatorname{Irr}_n(G)\longrightarrow\operatorname{Irr}'_n(C^{*}(G))
\]
are Borel isomorphisms. We then define $\operatorname{Rep}(G)$ and
$\operatorname{Irr}(G)$ as the Borel sum spaces of the
$\operatorname{Rep}_n(G)$ and the $\operatorname{Irr}_n(G)$; all these
spaces are standard.

\ResultHeading{18.5.2}{18.5.2.}
The functions
$\pi\mapsto(\pi(s)\xi\mid\eta)$ on $\operatorname{Rep}_n(G)$, where
$s\in G$ and $\xi,\eta\in H_n$, define on $\operatorname{Rep}_n(G)$ a
Borel structure identical with that of \XRef{18.5.1}{18.5.1}. Indeed,
these functions are continuous (\XRef{18.1.9}{18.1.9}), and hence this
structure is coarser than that of \XRef{18.5.1}{18.5.1}. It is then enough
to observe that the functions
$\pi\mapsto(\pi(s_k)\xi_p\mid\eta_q)$, where $(s_k)$ is dense in $G$ and
$(\xi_p),(\eta_q)$ are dense in $H_n$, separate the points of
$\operatorname{Rep}_n(G)$, and to apply \XRef{B.21}{B 21}.

\ResultHeading{18.5.3}{18.5.3.}
The \emph{Mackey Borel structure} on $\widehat G$ is the quotient structure
of that on $\operatorname{Irr}(G)$ under the canonical map
$\operatorname{Irr}(G)\to\widehat G$. This structure is identified with
the Mackey Borel structure of $(C^{*}(G))^{\widehat{\ }}$. It is finer than
the Borel structure underlying the topology of $\widehat G$
(\XRef{3.8.3}{3.8.3}), and is identical with it if $G$ is postliminal
(\XRef{4.6.1}{4.6.1}). The points of $\widehat G$ are Borel sets for the
Mackey structure (\XRef{3.8.4}{3.8.4}).

Bibliography: \cite{ref91}.

\BookSubsection{18.6}{Quasi-Dual}

\ResultHeading{18.6.1}{18.6.1.}
Let $G$ be a separable locally compact group, let $n$ be a cardinal, and
let $H_n$ be the model Hilbert space of dimension $n$. We denote by
$\operatorname{Fac}_n(G)$ the set of continuous unitary factor
representations of $G$ on $H_n$, by $\operatorname{Fac}(G)$ the union of
the $\operatorname{Fac}_n(G)$ for $n=1,2,\ldots,\aleph_0$, by
$\operatorname{FacI}_n(G)$ the set of $\pi\in\operatorname{Fac}(G)$ for
which $\pi(G)''$ is a factor of type $I_n$, and by
$\operatorname{FacI}(G)$ the set of $\pi\in\operatorname{Fac}(G)$ which
are of type I. All these sets are Borel subsets of
$\operatorname{Rep}(G)$ and hence are standard Borel spaces
(\XRef{7.1.4}{7.1.4} and \XRef{7.3.4}{7.3.4}). They are identified with
$\operatorname{Fac}_n(C^{*}(G))$, and so on.

\ResultHeading{18.6.2}{18.6.2.}
We denote by $\check G$ the set of quasi-equivalence classes of continuous
unitary factor representations of $G$. The Mackey Borel structure on
$\check G$ is the quotient structure of that on $\operatorname{Fac}(G)$
under the canonical map $\operatorname{Fac}(G)\to\check G$. The Borel
space $\check G$ is called the \emph{quasi-dual} of $G$. It is identified
with $(C^{*}(G))^{\check{\ }}$. Every point of $\check G$ is a Borel set
in $\check G$ (\XRef{7.2.4}{7.2.4}).

% Source PDF physical page 333
\ResultHeading{18.6.3}{18.6.3.}
Let $\check G_I$ be the set of quasi-equivalence classes of the
$\pi\in\operatorname{Fac}(G)$ which are of type I. Then $\check G_I$ is
identified with $(C^{*}(G))^{\check{\ }}_I$. It is a Borel subset of
$\check G$, and the canonical bijection from $\widehat G$ onto
$\check G_I$ is a Borel isomorphism (\XRef{7.3.6}{7.3.6}).

\ResultHeading{18.6.4}{18.6.4.}
Let $\operatorname{Facf}_n(G)$ be the set of
$\pi\in\operatorname{Fac}_n(G)$ which are of finite type, let
$\operatorname{Facf}(G)$ be the union of the $\operatorname{Facf}_n(G)$,
and let $\check G_f$ be the canonical image of $\operatorname{Facf}(G)$ in
$\check G$. The sets $\operatorname{Facf}(G)$ and
$\operatorname{Facf}_n(G)$ are Borel in $\operatorname{Fac}(G)$,
$\check G_f$ is Borel in $\check G$, all these spaces are standard
(\XRef{7.4.3}{7.4.3} and \XRef{7.4.4}{7.4.4}), and they are identified
with $\operatorname{Facf}(C^{*}(G))$,
$\operatorname{Facf}_n(C^{*}(G))$, and $(C^{*}(G))^{\check{\ }}_f$.

Let $C$ be the set of continuous central functions $\varphi$ of positive
type on $G$ such that $\varphi(e)\leq1$. Then $C$ is compact metrizable for
the weak topology. The set $D$ of finite-type characters of $G$ equal to
$1$ at $e$ is a $G_\delta$ in $C$ (\XRef{17.3.5}{17.3.5} and
\XRef{B.13}{B 13}), and hence is a Polish space. Equip $\check G_f$ with
the topology transported from $D$ by the canonical bijection of $D$ onto
$\check G_f$ (\XRef{17.3.4}{17.3.4}). Then $\check G_f$ becomes a Polish
space, and the Borel structure underlying its topology is the structure
induced by that of $\check G$ (\XRef{7.4.3}{7.4.3}).

Bibliography: \cite{ref25,ref57}.

\BookSubsection{18.7}{Integration and Disintegration of Representations}

\ResultHeading{18.7.1}{18.7.1.}
Let $Z$ be a Borel space, let $\mu$ be a positive measure on $Z$, and let
$\zeta\mapsto H(\zeta)$ be a $\mu$-measurable field of Hilbert spaces on
$Z$. For every $\zeta\in Z$, let $\pi(\zeta)$ be a continuous unitary
representation of $G$ on $H(\zeta)$; we say that
$\zeta\mapsto\pi(\zeta)$ is a field of continuous unitary representations
of $G$. This field is called measurable if, for every $s\in G$, the field
of operators $\zeta\mapsto\pi(\zeta)(s)$ is measurable.

\ResultHeading{18.7.2}{18.7.2.}
Suppose that $G$ is separable. Suppose that $Z$ is the union of pairwise
disjoint Borel sets $Z_1,Z_2,\ldots,Z_{\aleph_0}$, and that on $Z_n$ the
field $\zeta\mapsto H(\zeta)$ reduces to the constant field
$\zeta\mapsto H_n$, where $H_n$ is the model Hilbert space of dimension
$n$. Let $\zeta\mapsto\pi(\zeta)$ be a field of continuous unitary
representations of $G$ on the $H(\zeta)$. This field is measurable if and
only if $\zeta\mapsto\pi(\zeta)$ is almost everywhere equal to a Borel map
of $Z$ into $\operatorname{Rep}(G)$. The proof is exactly the same as that
of \XRef{8.1.8}{8.1.8}.

\ResultHeading{18.7.3}{18.7.3. Proposition.}
\begin{itshape}
Let $\zeta\mapsto\pi(\zeta)$ be a field of continuous unitary
representations of $G$. For every $\zeta\in Z$, let $\pi'(\zeta)$ be the
corresponding representation of $C^{*}(G)$. Suppose $G$ is separable. The
field $\zeta\mapsto\pi(\zeta)$ is measurable if and only if the field
$\zeta\mapsto\pi'(\zeta)$ is measurable.
\end{itshape}

First reduce to the case in which the field $\zeta\mapsto H(\zeta)$ is the
constant field defined by a fixed Hilbert space $H$ (\XRef{A.72}{A 72}).
One may then
% Source PDF physical page 334
suppose that $H$ is a model Hilbert space $H_n$. To say that the field
$\zeta\mapsto\pi(\zeta)$ [respectively,
$\zeta\mapsto\pi'(\zeta)$] is measurable means that it is almost
everywhere equal to a Borel map of $Z$ into
$\operatorname{Rep}_n(G)$ [respectively, into
$\operatorname{Rep}_n(C^{*}(G))$], by \XRef{18.7.2}{18.7.2}
[respectively, \XRef{8.1.8}{8.1.8}]. The proposition now follows from
\XRef{18.5.1}{18.5.1}.

\ResultHeading{18.7.4}{18.7.4. Proposition.}
\begin{itshape}
Retain the hypotheses of \XRef{18.7.3}{18.7.3}, and suppose the fields
$\zeta\mapsto\pi(\zeta)$ and $\zeta\mapsto\pi'(\zeta)$ are measurable.
Let
\[
\pi'=\int_Z^\oplus\pi'(\zeta)\,d\mu(\zeta).
\]
Let $\pi$ be the continuous unitary representation of $G$ corresponding
to $\pi'$. Then, for every $s\in G$,
\[
\pi(s)=\int_Z^\oplus\pi(\zeta)(s)\,d\mu(\zeta).
\]
\end{itshape}

Let $(V_1,V_2,\ldots)$ be a sequence of compact neighborhoods of $s$ whose
intersection is $\{s\}$. Let $f_n$ be a nonnegative continuous function,
supported in $V_n$, with integral $1$. Then $\pi'(f_n)$ converges strongly
to $\pi(s)$, $\pi'(\zeta)(f_n)$ converges strongly to $\pi(\zeta)(s)$ for
every $\zeta$, and
\[
\pi'(f_n)=\int_Z^\oplus\pi'(\zeta)(f_n)\,d\mu(\zeta).
\]
Hence
$\pi(s)=\int_Z^\oplus\pi(\zeta)(s)\,d\mu(\zeta)$
(\XRef{A.79}{A 79}).

\ResultHeading{18.7.5}{18.7.5.}
With the preceding notation, $\pi$ is called the Hilbert integral of the
$\pi(\zeta)$, and we write
\[
\pi=\int_Z^\oplus\pi(\zeta)\,d\mu(\zeta).
\]

\ResultHeading{18.7.6}{18.7.6.}
This having been established, all of \S~8, from Proposition
\XRef{8.1.6}{8.1.6} through \SectionRef{8.7}{8.7}, may be copied verbatim, with
$A$ replaced everywhere by $G$ (a separable locally compact group) and
``representation'' by ``continuous unitary representation.''

\ResultHeading{18.7.7}{18.7.7.}
We shall not translate \XRef{8.8.1}{8.8.1}. Applying \XRef{8.8.2}{8.8.2}
gives, in particular, the following:

\ResultLabel{Proposition.}
\begin{itshape}
Let $G$ be a separable locally compact unimodular group; let $\lambda$ and
$\rho$ be the left and right regular representations of $G$; let
$\mathfrak U$ and $\mathfrak V$ be the von Neumann algebras on $L^2(G)$
generated by $\lambda(G)$ and $\rho(G)$; let $J$ be the map $f\mapsto f^*$
of $L^2(G)$ onto itself; and let $t$ be the natural trace on
$\mathfrak U^+$ defined by $\varepsilon_e$ (\XRef{17.2.5}{17.2.5}).

\begin{enumerate}[label=\textup{(\roman*)}]
\item There exist: (1) a standard positive measure $\mu$ on $\check G$;
(2) for every $\zeta\in\check G$, a character $f(\zeta)$; and (3) a
measurable-field structure on the field
$\zeta\mapsto H_{f(\zeta)}$, with the following properties:

\emph{a.} $L^2(G)$ is identified with
$\int_{\check G}^\oplus H_{f(\zeta)}\,d\mu(\zeta)$; the fields
\[
\begin{array}{lll}
\zeta\mapsto J_{f(\zeta)},&
\zeta\mapsto\lambda_{f(\zeta)},&
\zeta\mapsto\bar\rho^{\mathrm o}_{f(\zeta)},\\
\zeta\mapsto\mathfrak U_{f(\zeta)},&
\zeta\mapsto\mathfrak V_{f(\zeta)},&
\zeta\mapsto t_{f(\zeta)}.
\end{array}
\]
% Source PDF physical page 335
are measurable, and
\begin{align*}
J&=\int_{\check G}^\oplus J_{f(\zeta)}\,d\mu(\zeta),&
\lambda&=\int_{\check G}^\oplus\lambda_{f(\zeta)}\,d\mu(\zeta),&
\rho&=\int_{\check G}^\oplus\bar\rho^{\mathrm o}_{f(\zeta)}\,d\mu(\zeta),\\
\mathfrak U&=\int_{\check G}^\oplus\mathfrak U_{f(\zeta)}\,d\mu(\zeta),&
\mathfrak V&=\int_{\check G}^\oplus\mathfrak V_{f(\zeta)}\,d\mu(\zeta),&
t&=\int_{\check G}^\oplus t_{f(\zeta)}\,d\mu(\zeta).
\end{align*}
The algebra $\mathfrak U\cap\mathfrak V$ is the algebra of diagonalizable
operators.

\emph{b.} $\lambda_{f(\zeta)}\in\zeta$ and
$\bar\rho^{\mathrm o}_{f(\zeta)}\in\zeta$ almost everywhere.

\emph{c.} For every $x\in C^{*}(G)^+$, the function
$\zeta\mapsto f(\zeta)(x)$ is measurable and
\[
f(x)=\int_{\check G}f(\zeta)(x)\,d\mu(\zeta).
\]

\item Suppose a measure $\mu'$, characters $f'(\zeta)$, and a measurable
field structure on $\zeta\mapsto H_{f'(\zeta)}$ are given, with the same
properties as in (i). Then $\mu$ and $\mu'$ are equivalent; write
$\mu'=d\cdot\mu$, where $d(\zeta)>0$ almost everywhere. We have
$f(\zeta)=d(\zeta)f'(\zeta)$ almost everywhere. After modifying the
$f'(\zeta)$ on a null set, there is an isomorphism from the measurable
field $\zeta\mapsto H_{f(\zeta)}$ onto the measurable field
$\zeta\mapsto H_{f'(\zeta)}$ which transforms $J_{f(\zeta)}$ into
$J_{f'(\zeta)}$, $\lambda_{f(\zeta)}$ into $\lambda_{f'(\zeta)}$,
$\bar\rho^{\mathrm o}_{f(\zeta)}$ into
$\bar\rho^{\mathrm o}_{f'(\zeta)}$,
$\mathfrak U_{f(\zeta)}$ into $\mathfrak U_{f'(\zeta)}$,
$\mathfrak V_{f(\zeta)}$ into $\mathfrak V_{f'(\zeta)}$, and
$t_{f(\zeta)}$ into $d(\zeta)t_{f'(\zeta)}$.
\end{enumerate}
\end{itshape}

\ResultHeading{18.7.8}{18.7.8. Corollary.}
\begin{itshape}
Retain the preceding notation. Let $u\in L^1(G)\cap L^2(G)$. For almost
every $\zeta\in\check G$, $u\in\mathfrak n_{f(\zeta)}$; the function
\[
\zeta\longmapsto
t_{f(\zeta)}\bigl(\lambda_{f(\zeta)}(u)
\lambda_{f(\zeta)}(u)^*\bigr)
\]
is $\mu$-integrable, and
\[
\int_G|u(s)|^2\,ds
=\int_{\check G}t_{f(\zeta)}\bigl(\lambda_{f(\zeta)}(u)
\lambda_{f(\zeta)}(u)^*\bigr)\,d\mu(\zeta).
\]
\end{itshape}

We have $u\in C^{*}(G)$ and $u\ast u^*\in C^{*}(G)^+$. Hence the function
\[
\zeta\longmapsto f(\zeta)(u\ast u^*)
=t_{f(\zeta)}\bigl(\lambda_{f(\zeta)}(u)
\lambda_{f(\zeta)}(u)^*\bigr)
\]
is measurable, and its integral with respect to $\mu$ is
\[
f(u\ast u^*)=\varepsilon_e(u\ast u^*)
=\int_G|u(s)|^2\,ds<+\infty.
\]
Thus, for almost every $\zeta$,
\[
t_{f(\zeta)}\bigl(\lambda_{f(\zeta)}(u)
\lambda_{f(\zeta)}(u)^*\bigr)<+\infty,
\]
and consequently $u\in\mathfrak n_{f(\zeta)}$.

\ResultHeading{18.7.9}{18.7.9. Corollary.}
\begin{itshape}
Let $G$ be a separable locally compact unimodular group, and let
$s\in G$, $s\ne e$. There exists a traceable continuous unitary factor
representation $\tau$ of $G$ such that $\tau(s)\ne1$.
\end{itshape}

% Source PDF physical page 336
With the notation of \XRef{18.7.7}{18.7.7}, we have $\lambda(s)\ne1$, and
hence $\lambda_{f(\zeta)}(s)\ne1$ for a non-null set of values of $\zeta$.
But the $\lambda_{f(\zeta)}$ are traceable factor representations.

\ResultHeading{18.7.10}{18.7.10.}
Let $(u_i)$ be a sequence of elements of $\mathcal K(G)$. By
\XRef{8.8.3}{8.8.3}, the representation $\tau$ of
\XRef{18.7.9}{18.7.9} may be chosen so that, for every $i$, $\tau(u_i)$ is
a Hilbert--Schmidt operator relative to the factor generated by
$\tau(G)$.

Bibliography: \cite{ref25,ref52,ref53,ref57,ref87,ref88,ref89,ref90,
ref91,ref95,ref96,ref107,ref108,ref109,ref110,ref137,ref162}.

\BookSubsection{18.8}{Plancherel Measure}

\ResultHeading{18.8.1}{18.8.1.}
Let $G$ be a separable postliminal locally compact group.

Let $\zeta\mapsto K(\zeta)$ be the canonical field of Hilbert spaces over
$\widehat G$ (\XRef{8.6.1}{8.6.1}). Choose a field
$\zeta\mapsto\pi(\zeta)$ of continuous unitary representations of $G$ on
the $K(\zeta)$ such that $\pi(\zeta)$ belongs to the class $\zeta$ for
every $\zeta$, and which is measurable for every positive measure on
$\widehat G$ (\XRef{8.6.2}{8.6.2}). For simplicity we shall write $\zeta$
in place of $\pi(\zeta)$. Let $\overline{K(\zeta)}$ be the Hilbert space
conjugate to $K(\zeta)$, let $J_\zeta$ be the canonical involution of
$K(\zeta)\otimes\overline{K(\zeta)}$, and let $t_\zeta$ be the trace
$T\otimes1\mapsto\operatorname{Tr}T$ on
$(\mathcal L(K(\zeta))\otimes\mathbb C)^+$.

\ResultLabel{Theorem.}
\begin{itshape}
Suppose in addition that $G$ is unimodular. Let $\lambda$ and $\rho$ be
the left and right regular representations of $G$, let $\mathfrak U$ and
$\mathfrak V$ be the von Neumann algebras on $L^2(G)$ generated by
$\lambda(G)$ and $\rho(G)$, let $J$ be the map $f\mapsto f^*$ on $L^2(G)$,
and let $t$ be the natural trace on $\mathfrak U^+$ defined by
$\varepsilon_e$ (\XRef{17.2.5}{17.2.5}). There exist a positive measure
$\mu$ on $\widehat G$ and an isomorphism $W$ from $L^2(G)$ onto
\[
\int_{\widehat G}^\oplus
\bigl(K(\zeta)\otimes\overline{K(\zeta)}\bigr)\,d\mu(\zeta)
\]
with the following properties:
\begin{enumerate}[label=\textup{(\roman*)}]
\item $W$ transforms
\begin{align*}
J&\quad\text{into}\quad
 \int_{\widehat G}^\oplus J_\zeta\,d\mu(\zeta),&
\lambda&\quad\text{into}\quad
 \int_{\widehat G}^\oplus(\zeta\otimes1)\,d\mu(\zeta),\\
\rho&\quad\text{into}\quad
 \int_{\widehat G}^\oplus(1\otimes\overline\zeta)\,d\mu(\zeta),&
\mathfrak U&\quad\text{into}\quad
 \int_{\widehat G}^\oplus
 (\mathcal L(K(\zeta))\otimes\mathbb C)\,d\mu(\zeta),\\
\mathfrak V&\quad\text{into}\quad
 \int_{\widehat G}^\oplus
 (\mathbb C\otimes\mathcal L(\overline{K(\zeta)}))\,d\mu(\zeta),&
t&\quad\text{into}\quad
 \int_{\widehat G}^\oplus t_\zeta\,d\mu(\zeta),
\end{align*}
and $\mathfrak U\cap\mathfrak V$ into the algebra of diagonalizable
operators.

% Source PDF physical page 337
\item If $x\in C^{*}(G)^+$, the function
$\zeta\mapsto\operatorname{Tr}\zeta(x)$ on $\widehat G$ is lower
semicontinuous, and
\[
t(x)=\int_{\widehat G}\operatorname{Tr}\zeta(x)\,d\mu(\zeta).
\]
In particular, if $u\in L^1(G)\cap L^2(G)$, then
\begin{equation}
\int_G|u(s)|^2\,ds
=\int_{\widehat G}\operatorname{Tr}
 \bigl(\zeta(u)\zeta(u)^*\bigr)\,d\mu(\zeta).
\tag{1}\label{eq:18.8.1-1}
\end{equation}
\end{enumerate}
\end{itshape}

This follows from \XRef{8.8.5}{8.8.5}, applied to the trace $t$.
Formula \eqref{eq:18.8.1-1} is called the \emph{Plancherel formula}.

\ResultHeading{18.8.2}{18.8.2. Theorem.}
\begin{itshape}
Let $G$ be a separable unimodular postliminal locally compact group. There
exists one and only one positive measure $\mu$ on $\widehat G$ such that,
for every $u\in L^1(G)\cap L^2(G)$,
\[
\int_G|u(s)|^2\,ds
=\int_{\widehat G}\operatorname{Tr}
 \bigl(\zeta(u)\zeta(u)^*\bigr)\,d\mu(\zeta).
\]
\end{itshape}

The measure $\mu$ of \XRef{18.8.1}{18.8.1} has this property. Let $\mu'$
be another positive measure on $\widehat G$ having this property. For
$x\in C^{*}(G)^+$, put
\[
f(x)=\int_{\widehat G}\operatorname{Tr}\zeta(x)\,d\mu'(\zeta).
\]
It is clear that $f$ is a trace on $C^{*}(G)^+$. If $(x_1,x_2,\ldots)$ is
a sequence of elements of $C^{*}(G)^+$ tending to $x$, then
$\|\zeta(x_n)-\zeta(x)\|\to0$ for every $\zeta\in\widehat G$; hence, by
Fatou's lemma,
\begin{align*}
\liminf f(x_n)
&\geq\int_{\widehat G}\liminf\operatorname{Tr}\zeta(x_n)\,d\mu'(\zeta)\\
&\geq\int_{\widehat G}\operatorname{Tr}\zeta(x)\,d\mu'(\zeta)=f(x).
\end{align*}
Thus $f$ is lower semicontinuous. We have $f(u\ast u^*)<+\infty$ for
every $u\in L^1(G)\cap L^2(G)$, so $\mathfrak m_f$ is dense in
$C^{*}(G)$. The bitraces corresponding to $f$ and $t$ agree on
$L^1(G)\cap L^2(G)$, and hence agree everywhere (\XRef{6.5.3}{6.5.3}).
Consequently
\[
t(x)=\int_{\widehat G}\operatorname{Tr}\zeta(x)\,d\mu'(\zeta)
\]
for every $x\in C^{*}(G)^+$, whence $\mu=\mu'$
(\XRef{8.8.6}{8.8.6}).

\ResultHeading{18.8.3}{18.8.3. Definition.}
\begin{itshape}
The measure $\mu$ of \XRef{18.8.2}{18.8.2} is called the Plancherel
measure on $\widehat G$ associated with the Haar measure of $G$.
\end{itshape}

Suppose the Haar measure is multiplied by a scalar $k>0$. Let
$u\in L^1(G)\cap L^2(G)$. For every $\zeta\in\widehat G$, $\zeta(u)$ is
multiplied by $k$, and hence
$\operatorname{Tr}(\zeta(u)\zeta(u)^*)$ is multiplied by $k^2$. We see
that the Plancherel measure is therefore multiplied by $k^{-1}$.

% Source PDF physical page 338
\ResultHeading{18.8.4}{18.8.4.}
The support of the Plancherel measure $\mu$ equals the support of the
regular representation of $G$ (\XRef{8.6.9}{8.6.9} and
\XRef{18.8.1}{18.8.1}), that is, the reduced dual of $G$
(\XRef{18.3.1}{18.3.1}).

If $K$ is a quasi-compact subset of $\widehat G$, then
$\mu(K)<+\infty$. Indeed, for every $\pi_0\in K$ there is a
$u\in L^1(G)\cap L^2(G)$ such that $\pi_0(u)\ne0$; hence there is an
$\alpha>0$ such that
$\operatorname{Tr}(\pi(u)\pi(u)^*)\geq\alpha$ on an open neighborhood
$V_{\pi_0}$ of $\pi_0$. By the Plancherel formula,
$\mu(V_{\pi_0})<+\infty$. But $K$ can be covered by finitely many of the
sets $V_{\pi_0}$.

\ResultHeading{18.8.5}{18.8.5. Proposition.}
\begin{itshape}
Let $G$ be a separable unimodular postliminal locally compact group,
equipped with a Haar measure, and let $\mu$ be the corresponding
Plancherel measure on $\widehat G$. Let $\zeta_0\in\widehat G$. The
representation $\zeta_0$ is square-integrable if and only if
$\mu(\{\zeta_0\})>0$, and $\mu(\{\zeta_0\})$ then equals the formal
dimension of $\zeta_0$.
\end{itshape}

If $\mu(\{\zeta_0\})>0$, then $\zeta_0\otimes1$ is contained in the
regular representation $\lambda$ [\XRef{8.6.8}{8.6.8 (ii)} and
\XRef{18.8.1}{18.8.1}]; thus $\zeta_0$ is contained in $\lambda$, and so
$\zeta_0$ is square-integrable (\XRef{14.1.3}{14.1.3}).

Suppose that $\zeta_0$ is square-integrable. Use the notation of
\XRef{18.8.1}{18.8.1}. Let $F$ be the central projection associated with
$\zeta_0$ (\XRef{14.2.4}{14.2.4}). It is diagonalizable; let $Y$ be the
$\mu$-measurable subset of $\widehat G$ such that $F$ corresponds to the
characteristic function $\varphi_Y$ of $Y$. Then the subrepresentation of
$\lambda$ defined by $F$ is, on the one hand, quasi-equivalent to
$\zeta_0$ (\XRef{14.3.5}{14.3.5}) and, on the other hand, equivalent to
\[
\int_{\widehat G}^\oplus
 (\zeta\otimes1)\varphi_Y(\zeta)\,d\mu(\zeta)
\]
(\XRef{18.8.1}{18.8.1}). Thus $\varphi_Y\cdot\mu$ is equivalent to the
Dirac measure at $\zeta_0$ (\XRef{8.4.4}{8.4.4}), and hence
$\mu(\{\zeta_0\})>0$. Furthermore, for
$f\in L^1(G)\cap L^2(G)$, by \XRef{8.8.7}{8.8.7},
\begin{align*}
\int_G|Ff(s)|^2\,ds
 &=\int_{\widehat G}\varphi_Y(\zeta)
   \operatorname{Tr}\bigl(\zeta(f)^*\zeta(f)\bigr)\,d\mu(\zeta)\\
 &=\mu(\{\zeta_0\})
   \operatorname{Tr}\bigl(\zeta_0(f)^*\zeta_0(f)\bigr).
\end{align*}
On the other hand, if $d$ denotes the formal dimension of $\zeta_0$, then
\[
\int_G|Ff(s)|^2\,ds
=d\operatorname{Tr}\bigl(\zeta_0(f)^*\zeta_0(f)\bigr)
\qquad(\XRef{14.4.2}{14.4.2}).
\]
We may choose $f\in L^1(G)\cap L^2(G)$ such that $\zeta_0(f)\ne0$;
therefore $\mu(\{\zeta_0\})=d$.

\ResultHeading{18.8.6}{18.8.6.}
If $G$ is separable and commutative, the classical Plancherel theorem
shows that the Plancherel measure is precisely the suitably normalized
Haar measure of the dual group $\widehat G$.

Bibliography: \cite{ref16,ref53,ref60,ref62,ref81,ref94,ref95,ref96,
ref101,ref107,ref137,ref148,ref162,ref171,ref179}.

% Source PDF physical page 339
\BookSubsection{18.9}{Supplements}

\ResultHeading{18.9.1}{18.9.1.}
Let $G$ be a separable unimodular postliminal locally compact group, let
$\widehat G_r$ be its reduced dual, and let $\pi\in\widehat G$.

\emph{a.} If $\{\pi\}$ is open in $\widehat G_r$, then $\pi$ is
square-integrable. [For $\pi$ belongs to the support of the regular
representation; hence, if $\mu$ denotes Plancherel measure,
$\mu(\{\pi\})>0$.]

\emph{b.} It may happen that $\pi\in\widehat G_r$ is square-integrable,
that $\{\pi\}$ is open in $\widehat G$, and that $\pi$ is not integrable.
It may happen that $\pi$ is square-integrable and that $\{\pi\}$ is not
open in $\widehat G$. It may happen that $\{\pi\}$ is open in
$\widehat G$ and that $\pi\notin\widehat G_r$.

\emph{c.} Problems: If $\pi$ is integrable, is $\{\pi\}$ open in
$\widehat G$? If $\pi$ is square-integrable, is $\{\pi\}$ open in
$\widehat G_r$? \cite{ref13,ref283}.

\ResultHeading{18.9.2}{*18.9.2.}
The Plancherel formula can be generalized to nonseparable unimodular
postliminal locally compact groups. \cite{ref16}.

\ResultHeading{18.9.3}{18.9.3.}
Let $G$ be a locally compact group. Every continuous complex-valued
function on $G$ is a uniform limit on every compact set of linear
combinations of functions $f\ast\widetilde f$, $f\in\mathcal K(G)$; these
functions are associated with the regular representation. Compare this
with \XRef{18.1.4}{18.1.4} and with the fact that, for some $G$, not every
element of $\widehat G$ is weakly contained in the regular representation.
\cite{ref27}.

\ResultHeading{18.9.4}{*18.9.4.}
Let $G$ be a connected semisimple real Lie group. The set of
$f\in L^1(G)$ for which the rank of $\pi(f)$ remains finite and bounded as
$\pi$ ranges over $\widehat G$ is a self-adjoint two-sided ideal of
$L^1(G)$ dense in $L^1(G)$. Thus \XRef{4.7.11}{4.7.11} may be
applied. \cite{ref27}.

\ResultHeading{18.9.5}{18.9.5.}
Let $G$ be a locally compact group. If the regular representation of $G$
weakly contains a finite-dimensional continuous unitary representation,
then $\widehat G=\widehat G_r$. \cite{ref33}.

\ResultHeading{18.9.6}{18.9.6.}
Let $G$ be a separable unimodular postliminal locally compact group, let
$\mu$ be the Plancherel measure on $\widehat G$, and let $E$ be the set of
$\pi\in\widehat G$ such that
$\pi(C^{*}(G))=\mathcal{LC}(H_\pi)$, that is, such that $\{\pi\}$ is closed
in $\widehat G$. Then $\mu(\widehat G\setminus E)=0$. What can be said
about $E$ topologically?

\ResultHeading{18.9.7}{18.9.7.}
Let $G$ be a separable locally compact group in which $e$ has a
neighborhood basis consisting of compact neighborhoods invariant under
inner automorphisms. In \XRef{18.7.7}{18.7.7}, almost all the characters
$f(\zeta)$ are of finite type. (Use \XRef{8.8.3}{8.8.3} and
\XRef{17.4.1}{17.4.1}.)

\ResultHeading{18.9.8}{*18.9.8.}
Let $G$ be the free group on two generators, equipped with the discrete
topology. Then $\widehat G_r\ne\widehat G$; the set of pure functions of
positive type is dense in the set of functions of positive
type for compact convergence, and the disintegration of the regular
representation $\lambda$ into
% Source PDF physical page 340
irreducible representations can be carried out in two completely distinct
ways. The factor representation $\lambda$ admits a disintegration (over a
standard Borel space) into representations which are not factor
representations almost everywhere. (One may observe that the von Neumann
algebras associated with the Hilbert algebra of $G$ are factors of type
$II_1$, and apply \XRef{5.7.5}{5.7.5}, \XRef{5.7.6}{5.7.6}, and
\XRef{11.2.4}{11.2.4}.) \cite{ref4,ref185,ref199}.

\ResultHeading{18.9.9}{18.9.9.}
Let $G$ be a locally compact group, let $H$ be a closed normal subgroup of
$G$, let $\pi$ be a continuous unitary factor representation of $G$, and
put $\rho=\pi|H$. Then $\rho$ is either of type I, of type $II_1$, of type
$II_\infty$, or of type III. If $\rho$ is of type I, it has multiplicity
$n$ for a uniquely determined $n$. The measure class on $\widehat H$
associated with $\rho$ is invariant and ergodic under the action of $G$ on
$\widehat H$ (the group $G$ acts on $H$ by inner automorphisms and hence on
$\widehat H$ by transport of structure). \cite{ref92}.

\ResultHeading{18.9.10}{18.9.10.}
Let $G$ be a locally compact group and $H$ an open subgroup of $G$. If
$\widehat G=\widehat G_r$, then $\widehat H=\widehat H_r$.
\cite{ref154}.

\ResultHeading{18.9.11}{*18.9.11.}
Let $G$ be a connected noncompact simple real Lie group. Then
$\widehat G\ne\widehat G_r$. \cite{ref154}.

\ResultHeading{18.9.12}{18.9.12.}
Let $G$ be a locally compact group. The map $\pi\mapsto\overline\pi$ is a
homeomorphism of $\widehat G$ onto $\widehat G$.

\ResultHeading{18.9.13}{18.9.13.}
The topological space $\widehat G$ has been computed explicitly in only a
small number of cases. Here are some examples. If $G$ is the noncommutative
simply connected nilpotent real Lie group of dimension $3$, then
$\widehat G$ may be identified with
$(\mathbb R\setminus\{0\})\cup\mathbb R^2$; the topology is that of the
sum space of $\mathbb R\setminus\{0\}$ and $\mathbb R^2$, except that a
point of $\mathbb R\setminus\{0\}$ which tends to zero in the ordinary
sense tends in $\widehat G$ to every point of $\mathbb R^2$. If
$G=\operatorname{SL}(2,\mathbb C)$, then $\widehat G$ may be identified
with the subset of $\mathbb R^2$ formed by the pairs
$(n,\nu)$, where $n=1,2,\ldots$ and $\nu\in\mathbb R$, the pairs
$(0,\nu)$, where $\nu\geq0$, and the pairs $(x,0)$, where
$-1\leq x\leq0$; the topology is that induced by the topology of
$\mathbb R^2$, except that a point of $\widehat G$ which tends to
$(-1,0)$ in the ordinary sense tends in $\widehat G$ to both $(-1,0)$ and
$(2,0)$. \cite{ref8,ref27}.

For $G=\operatorname{SL}(2,\mathbb C)$, the structure of $C^{*}(G)$ has
been completely determined. \cite{ref29}.

For most groups, one is still at the stage of determining the set
$\widehat G$. This work is nearly complete for the complex semisimple Lie
groups, and in that case the Plancherel measure is also known; less
progress has been made for real semisimple Lie groups. The fundamental
work here is that of Gelfand, Graev, Harish-Chandra, and Naimark.

\ResultHeading{18.9.14}{18.9.14.}
If $G$ is commutative, $\widehat G$ carries a group structure. The
extension of this fact to the noncommutative case is the existence, for
every pair $(\pi,\pi')$ of elements of $\widehat G$, of a measure class and
multiplicities on $\widehat G$, namely the measure class and
multiplicities associated with $\pi\otimes\pi'$ (when $G$ is separable and
postliminal). This measure class and these multiplicities have been
computed in only a small number of cases.

% Source PDF physical page 341
\ResultHeading{18.9.15}{18.9.15.}
Let $G$ be a locally compact group, let $\pi_1,\pi_2$ be continuous
unitary representations of $G$, and let $S_1,S_2$ be sets of continuous
unitary representations of $G$. If $\pi_k$ is weakly contained in $S_k$
for $k=1,2$, then $\pi_1\otimes\pi_2$ is weakly contained in the set of
$\rho_1\otimes\rho_2$ with $\rho_1\in S_1$ and $\rho_2\in S_2$.
\cite{ref33}.

\ResultHeading{18.9.16}{*18.9.16.}
Let $G$ be a locally compact group, let $\pi$ be a finite-dimensional
irreducible continuous unitary representation of $G$, let $S$ be a set of
continuous unitary representations of $G$, and let $T$ be the set of all
$\rho\otimes\overline\pi$ with $\rho\in S$. The set $S$ weakly contains
$\pi$ if and only if $T$ weakly contains the one-dimensional trivial
representation. Both parts of this assertion may fail if
$\dim\pi=+\infty$. \cite{ref33}.

% Source PDF physical page 342
\BookAppendix{A}{Von Neumann Algebras}

This Appendix is a list of results concerning von Neumann algebras which
are used in the book. No proofs are given, but only references to
\emph{Les alg\`ebres d'op\'erateurs dans l'espace hilbertien}, second
edition, Gauthier--Villars, Paris, 1969 (designated by \AvNTarget).

Throughout, $H$ denotes a Hilbert space.

\ResultHeading{A.1}{A 1.}
The \emph{norm topology} on $\mathcal L(H)$ is the topology induced by the
norm $T\mapsto\|T\|$. The \emph{strong}, \emph{weak},
\emph{ultrastrong}, and \emph{ultraweak} topologies are the four topologies
defined, respectively, by the following seminorms:
\[
T\longmapsto\|T\xi\|\qquad(\xi\in H),
\]
for the strong topology;
\[
T\longmapsto |(T\xi\mid\eta)|\qquad(\xi,\eta\in H),
\]
for the weak topology;
\[
T\longmapsto\left(\sum_{i=1}^{\infty}\|T\xi_i\|^2\right)^{1/2}
\quad
\left(\xi_1,\xi_2,\ldots\in H,\quad
\sum_i\|\xi_i\|^2<+\infty\right),
\]
for the ultrastrong topology; and
\[
T\longmapsto\left|\sum_{i=1}^{\infty}(T\xi_i\mid\eta_i)\right|
\quad
\left(\begin{array}{c}
\xi_1,\xi_2,\ldots,\eta_1,\eta_2,\ldots\in H,\\
\sum_i\|\xi_i\|^2<+\infty,\quad
\sum_i\|\eta_i\|^2<+\infty
\end{array}\right),
\]
for the ultraweak topology.

On a bounded subset of $\mathcal L(H)$, the strong and ultrastrong
topologies coincide, and the weak and ultraweak topologies coincide.
(\AvNRef, pp.~30--34.)

\ResultHeading{A.2}{A 2.}
Let $M$ be a subset of $\mathcal L(H)$. The \emph{commutant} of $M$,
denoted $M'$, is the set of elements of $\mathcal L(H)$ which commute with
the elements of $M$. Put $(M')'=M''$ (the \emph{bicommutant} of $M$).
The set $M'$ is a subalgebra of $\mathcal L(H)$ containing $1$. The
inclusion $M\subset N$ implies $M'\supset N'$. We have
$M'=M'''$ and $M\subset M''$.

\ResultHeading{A.3}{A 3.}
Let $A$ be an involutive subalgebra of $\mathcal L(H)$ and let $A_1$ be
the unit ball of $A$. The following eight conditions are equivalent:
$A$ (respectively, $A_1$) is weakly closed, $A$ (respectively, $A_1$) is
strongly closed, $A$ (respectively, $A_1$) is ultrastrongly closed, and
$A$ (respectively, $A_1$) is ultraweakly closed. When these conditions
hold, $A$ contains a projection $E$ greater than every other projection
in $A$.
% Source PDF physical page 343
For every $T\in A$, one has $ET=TE=T$. The elements of $A''$ are precisely
the operators $T+\lambda1$, where $T\in A$ and $\lambda\in\mathbb C$.
(\AvNRef, p.~41, Theorem~2.)

\ResultHeading{A.4}{A 4.}
Let $A$ be an involutive subalgebra of $\mathcal L(H)$. The following nine
conditions are equivalent: (1) each of the eight conditions in
\XRef{A.3}{A 3}, with the additional hypothesis $1\in A$; (2) $A=A''$.
When they hold, $A$ is called a \emph{von Neumann algebra on $H$}
(\AvNRef, p.~42). The set $\mathcal L(H)$ is a von Neumann algebra on
$H$. The set of scalar operators on $H$ is a von Neumann algebra on $H$,
denoted $\mathbb C_H$, or simply $\mathbb C$.

\ResultHeading{A.5}{A 5.}
If $M$ is a self-adjoint subset of $\mathcal L(H)$, then $M'$ is a von
Neumann algebra (\AvNRef, p.~2).

\ResultHeading{A.6}{A 6.}
If $\mathcal A$ is a von Neumann algebra, then $\mathcal A$ and
$\mathcal A'$ are each the commutant of the other. Their common center is
$\mathcal A\cap\mathcal A'$. If this center consists only of the scalar
operators, $\mathcal A$ is called a \emph{factor}. The algebra
$\mathcal A$ is a factor if and only if $\mathcal A'$ is a factor
(\AvNRef, p.~3).

\ResultHeading{A.7}{A 7.}
Let $\mathcal A$ be a von Neumann algebra on $H$, and let $T$ be a
Hermitian element of $\mathcal L(H)$. The operator $T$ belongs to
$\mathcal A$ if and only if its spectral projections belong to
$\mathcal A$ (\AvNRef, p.~3).

\ResultHeading{A.8}{A 8.}
Let $\mathcal A$ be a von Neumann algebra, and let $S,T\in\mathcal A^+$
with $S\leq T$. There exists an $R\in\mathcal A$ such that
$S^{1/2}=RT^{1/2}$ (\AvNRef, p.~11, Lemma~2).

\ResultHeading{A.9}{A 9.}
Let $\mathcal A$ be a von Neumann algebra and $\mathfrak m$ a two-sided
ideal of $\mathcal A$. Then $\mathfrak m$ is self-adjoint and is the set
of linear combinations of elements of
$\mathfrak m^+=\mathfrak m\cap\mathcal A^+$. The set of $T\in\mathcal A$
such that $TT^*\in\mathfrak m$ is a two-sided ideal $\mathfrak n$ of
$\mathcal A$ for which the ideal $\mathfrak n^2$ equals $\mathfrak m$; it
is denoted $\mathfrak m^{1/2}$ (\AvNRef, pp.~10--12).

Every weakly closed two-sided ideal of $\mathcal A$ has the form
$E\mathcal A$, where $E$ is a projection of
$\mathcal A\cap\mathcal A'$ (\AvNRef, p.~42, Corollary~3).

\ResultHeading{A.10}{A 10.}
Let $\mathcal A$ be a von Neumann algebra and $E$ a projection of
$\mathcal A$. Among the projections of $\mathcal A\cap\mathcal A'$ which
majorize $E$, there is a least one. It is called the \emph{central support}
of $E$ (\AvNRef, p.~6).

\ResultHeading{A.11}{A 11.}
Let $M$ be a subset of $\mathcal L(H)$. There is a least von Neumann
algebra $\mathcal A$ on $H$ which contains $M$; it is called the von
Neumann algebra generated by $M$. If $M=M^*$, then $\mathcal A=M''$
(\AvNRef, p.~2).

\ResultHeading{A.12}{A 12.}
Let $A$ be an involutive subalgebra of $\mathcal L(H)$. The weak, strong,
ultraweak, and ultrastrong closures of $A$ coincide; denote this closure by
$\mathcal B$. Let $K$ be the closed vector subspace of $H$ spanned by the
% Source PDF physical page 344
vectors $T\xi$, where $T\in A$ and $\xi\in H$. Then $P_K$ is the greatest
projection in $\mathcal B$. If $K=H$, then $\mathcal B=A''$ is the von
Neumann algebra generated by $A$ (\AvNRef, p.~42).

\ResultHeading{A.13}{A 13.}
Let $A$ and $B$ be involutive subalgebras of $\mathcal L(H)$, with
$A\subset B$. If $A$ is strongly dense in $B$, then the unit ball $A_1$ of
$A$ is strongly dense in the unit ball $B_1$ of $B$ (\AvNRef, p.~43,
Theorem~3). If, in addition, $H$ is separable, every element of $B_1$ is
the strong limit of a sequence of elements of $A_1$ (\AvNRef, p.~32).

\ResultHeading{A.14}{A 14.}
Let $A$ be an involutive subalgebra of $\mathcal L(H)$ and let $\xi\in H$.
The vector $\xi$ is called \emph{separating} for $A$ if $T\in A$ and
$T\xi=0$ imply $T=0$. It is called \emph{cyclic} for $A$ if $A\xi$ is
dense in $H$. If $\xi$ is cyclic for $A$, then $\xi$ is separating for
$A'$. The converse holds if $A$ is a von Neumann algebra
(\AvNRef, pp.~5--6).

\ResultHeading{A.15}{A 15.}
Let $\mathcal A$ be a von Neumann algebra on $H$, and let $E=P_K$ be a
projection of $\mathcal A$. Let $\mathcal B$ be the set of
$T\in\mathcal A$ such that $ET=TE=T$. The operators $T|K$, with
$T\in\mathcal B$, form a von Neumann algebra on $K$, called the
\emph{reduction} of $\mathcal A$, and denoted $\mathcal A_E$ or
$\mathcal A_K$. Let $E'=P_{K'}$ be a projection of $\mathcal A'$. The
operators $T|K'$, with $T\in\mathcal A$, form a von Neumann algebra on
$K'$, said to be \emph{induced} by $\mathcal A$ on $K'$, and denoted
$\mathcal A_{E'}$ or $\mathcal A_{K'}$. The notations
$\mathcal A_E,\mathcal A_{E'}$ are consistent when the projection belongs
to $\mathcal A\cap\mathcal A'$. The map $T\mapsto T|K'$, $T\in\mathcal A$,
is called the \emph{induction} of $\mathcal A$ on $\mathcal A_{E'}$. The
von Neumann algebras $(\mathcal A_{E'})'$ and $(\mathcal A')_{E'}$ are
equal and are denoted simply by $\mathcal A'_{E'}$
(\AvNRef, pp.~16--18). If $\mathfrak Z$ is the center of $\mathcal A$,
the center of $\mathcal A_E$ is $\mathfrak Z_E$
(\AvNRef, p.~18, Corollary).

\ResultHeading{A.16}{A 16.}
Let $(H_i)$ be a family of Hilbert spaces, let $H$ be their Hilbert sum,
and let $\mathcal A_i$ be a von Neumann algebra on $H_i$. For every family
$(T_i)$ such that $T_i\in\mathcal A_i$ for all $i$ and
$\sup_i\|T_i\|<+\infty$, form the operator
$(x_i)\mapsto(T_ix_i)$ on $H$. The set of all operators so obtained is a
von Neumann algebra on $H$, called the \emph{product} of the
$\mathcal A_i$ and denoted $\prod_i\mathcal A_i$ (\AvNRef, p.~19).

\ResultHeading{A.17}{A 17.}
Let $H_1,H_2$ be two Hilbert spaces. On the algebraic tensor product
$H_0$ of $H_1$ and $H_2$ there is a unique pre-Hilbert space structure
such that
\[
(\xi_1\otimes\xi_2\mid\eta_1\otimes\eta_2)
=(\xi_1\mid\eta_1)(\xi_2\mid\eta_2)
\]
for $\xi_1,\eta_1\in H_1$ and $\xi_2,\eta_2\in H_2$. This structure is
separated. The Hilbert space $H$ obtained by completing $H_0$ is called the
\emph{Hilbert tensor product} of $H_1$ and $H_2$ and is denoted
$H_1\otimes H_2$. If $T_1\in\mathcal L(H_1)$ and
$T_2\in\mathcal L(H_2)$, their algebraic tensor product acts continuously
on $H_0$. Its continuous extension to $H$ is denoted $T_1\otimes T_2$
(\AvNRef, p.~21).

There is a family $(K_i)_{i\in I}$ of closed vector subspaces of
$H_1\otimes H_2$, of which $H_1\otimes H_2$ is the Hilbert sum, and
isomorphisms $U_i:H_1\to K_i$, with the following properties: (1) for
every $T\in\mathcal L(H_1)$,
$T\otimes1$ leaves every $K_i$ invariant and
$(T\otimes1)|K_i=U_iTU_i^{-1}$; (2) $\operatorname{Card}I=\dim H_2$
(\AvNRef, pp.~22--24).

% Source PDF physical page 345
\ResultHeading{A.18}{A 18.}
Let $\mathcal A_1$ be a von Neumann algebra on $H_1$ and
$\mathcal A_2$ a von Neumann algebra on $H_2$. Operators of the form
\[
R_1\otimes R_2+S_1\otimes S_2+\cdots+T_1\otimes T_2,
\]
where $R_1,\ldots,T_1\in\mathcal A_1$ and
$R_2,\ldots,T_2\in\mathcal A_2$, form an involutive subalgebra $A_0$ of
$\mathcal L(H_1\otimes H_2)$. The von Neumann algebra on
$H_1\otimes H_2$ generated by $A_0$ is called the \emph{von Neumann
tensor-product algebra} of $\mathcal A_1$ and $\mathcal A_2$, and is
denoted $\mathcal A_1\otimes\mathcal A_2$ (\AvNRef, p.~24). If $A$ is an
involutive subalgebra of $\mathcal L(H_1)$, the von Neumann algebra
generated by $A\otimes\mathbb C_{H_2}$ has commutant
$A'\otimes\mathcal L(H_2)$ (\AvNRef, p.~24).

\ResultHeading{A.19}{A 19.}
Let $\mathcal A_1,\mathcal A_2$ be von Neumann algebras on $H_1,H_2$.
An isomorphism from $\mathcal A_1$ onto $\mathcal A_2$ is simply an
isomorphism of the underlying involutive algebras; such an isomorphism is
automatically isometric. Of course, an isomorphism $H_1\to H_2$ of Hilbert
spaces defines an isomorphism (called \emph{spatial}) from every von
Neumann algebra on $H_1$ onto a von Neumann algebra on $H_2$. In general,
however, an isomorphism need not be spatial (\AvNRef, p.~9).

\ResultHeading{A.20}{A 20.}
Let $\mathcal A$ be a von Neumann algebra and $E'$ a projection of
$\mathcal A'$. The induction $T\mapsto T_{E'}$ from $\mathcal A$ onto
$\mathcal A_{E'}$ is an isomorphism if and only if the central support of
$E'$ is $1$ (\AvNRef, p.~18, Proposition~2).

\ResultHeading{A.21}{A 21.}
Let $\mathcal A$ be a von Neumann algebra on $H$, and let $H'$ be a Hilbert
space. The map $T\mapsto T\otimes1$ is an isomorphism of $\mathcal A$ onto
the von Neumann algebra
$\mathcal A\otimes\mathbb C_{H'}\subset\mathcal L(H\otimes H')$. Such an
isomorphism is called an \emph{amplification} (\AvNRef, p.~24).

\ResultHeading{A.22}{A 22.}
Let $\mathcal A_1,\mathcal A_2$ be von Neumann algebras on $H_1,H_2$, and
let $\Phi$ be an isomorphism of $\mathcal A_1$ onto $\mathcal A_2$. There
exist a Hilbert space $K$ and a projection
$E'\in(\mathcal A_1\otimes\mathbb C_K)'$ of central support $1$ such that
$\Phi$ is the composite of the amplification
$T\mapsto T\otimes1_K$, the induction
$T\otimes1_K\mapsto(T\otimes1_K)_{E'}$, and a spatial isomorphism. If
$H_1$ and $H_2$ are separable, $K$ may be taken separable
(\AvNRef, pp.~55--56).

\ResultHeading{A.23}{A 23.}
Let $\mathcal A$ be a von Neumann algebra on $H$, and let
$\mathcal A'$ denote the Banach dual of the Banach space $\mathcal A$.
For $f\in\mathcal A'$, ultraweak continuity and ultrastrong continuity of
$f$ are equivalent; the set of such forms is a vector subspace $P$ of
$\mathcal A'$, closed in the norm of $\mathcal A'$. Likewise, for
$f\in\mathcal A'$, weak continuity and strong continuity are equivalent;
the set of such forms is a vector subspace of $P$ dense in $P$ in the norm
of $\mathcal A'$. The canonical bilinear form on
$\mathcal A\times\mathcal A'$ induces a bilinear form on
$\mathcal A\times P$ for which $\mathcal A$, with its norm, is the Banach
dual of $P$. The space $P$ is called the \emph{predual} of $\mathcal A$
(\AvNRef, pp.~34--40).

% Source PDF physical page 346
\ResultHeading{A.24}{A 24.}
Let $f$ be a weakly continuous linear form on $\mathcal L(H)$. There are
two orthonormal systems $(\xi_1,\ldots,\xi_n)$ and
$(\xi'_1,\ldots,\xi'_n)$ in $H$, and numbers
$\lambda_1\geq0,\ldots,\lambda_n\geq0$, such that
$\|f\|=\lambda_1+\cdots+\lambda_n$ and
\[
f(T)=\sum_{i=1}^n\lambda_i(T\xi_i\mid\xi'_i)
\qquad(T\in\mathcal L(H)).
\]
(\AvNRef, pp.~36--37.)

\ResultHeading{A.25}{A 25.}
Let $\mathcal A$ be a von Neumann algebra. Every bounded increasing
directed subset of $\mathcal A^+$ has a supremum in $\mathcal A^+$
(\AvNRef, p.~4). For a positive form $f$ on $\mathcal A$, the following
conditions are equivalent: (1) for every increasing directed set
$\mathcal F\subset\mathcal A^+$ with supremum $T\in\mathcal A^+$,
$f(T)$ is the supremum of $f(\mathcal F)$; (2) $f$ belongs to the predual
of $\mathcal A$. Such a form $f$ is called \emph{normal}. Every element of
the predual of $\mathcal A$ is a linear combination of positive normal
forms (\AvNRef, pp.~50--51).

\ResultHeading{A.26}{A 26.}
Let $\mathcal A$ be a von Neumann algebra, and let $f$ be a positive normal
form on $\mathcal A$. There is a greatest projection $F$ of $\mathcal A$
such that $f(F)=0$. One has $f(TF)=f(FT)=0$ for every $T\in\mathcal A$.
The projection $E=1-F$ is called the \emph{support} of $f$. For every
$T\in\mathcal A$,
$f(T)=f(ET)=f(TE)=f(ETE)$ (\AvNRef, p.~57).

\ResultHeading{A.27}{A 27.}
Let $\mathcal A$ and $\mathcal B$ be two von Neumann algebras, and let
$\varphi$ be a morphism from $\mathcal A$ into $\mathcal B$. The morphism
$\varphi$ is called \emph{normal} if, for every increasing directed set
$\mathcal F\subset\mathcal A^+$ with supremum $T\in\mathcal A^+$,
$\varphi(T)$ is the supremum of $\varphi(\mathcal F)$. Then $\varphi$ is
ultrastrongly and ultraweakly continuous, and $\varphi(\mathcal A)$ is
weakly closed (\AvNRef, p.~53, Theorem~2, and p.~54, Corollary~2).

\ResultHeading{A.28}{A 28.}
Let $\mathcal A$ be a von Neumann algebra. A \emph{trace} on
$\mathcal A^+$ is a function $\varphi$ on $\mathcal A^+$, with values
$\geq0$ finite or infinite, which has the following properties:
\begin{enumerate}[label=\textup{(\roman*)}]
\item if $S,T\in\mathcal A^+$, then
$\varphi(S+T)=\varphi(S)+\varphi(T)$;
\item if $S\in\mathcal A^+$ and $\lambda\geq0$, then
$\varphi(\lambda S)=\lambda\varphi(S)$, with the convention
$0\cdot(+\infty)=0$;
\item if $S\in\mathcal A^+$ and $U$ is a unitary operator in $\mathcal A$,
then $\varphi(USU^{-1})=\varphi(S)$.
\end{enumerate}
The trace $\varphi$ is called \emph{finite} if
$\varphi(S)<+\infty$ for every $S\in\mathcal A^+$. It is called
\emph{semifinite} if, for every $S\in\mathcal A^+$, $\varphi(S)$ is the
supremum of the numbers $\varphi(T)$ for $T\in\mathcal A^+$ such that
$T\leq S$ and $\varphi(T)<+\infty$.

(Everything preceding is a special case of \XRef{6.1.1}{6.1.1}.)

The trace $\varphi$ is called \emph{faithful} if $S\in\mathcal A^+$ and
$\varphi(S)=0$ imply $S=0$. It is called \emph{normal} if, for every
increasing directed set $\mathcal F\subset\mathcal A^+$ with supremum
$S\in\mathcal A^+$, $\varphi(S)$ is the supremum of $\varphi(\mathcal F)$.

% Source PDF physical page 347
\ResultHeading{A.29}{A 29.}
Let $\mathcal A$ be a von Neumann algebra, let $\omega,\omega'$ be normal
traces defined on $\mathcal A^+$, and let $B$ be an involutive subalgebra
of $\mathcal A$ weakly dense in the von Neumann algebra $\mathcal A$.
Suppose that
\[
\omega(TT^*)=\omega'(TT^*)<+\infty\qquad(T\in B).
\]
Then $\omega=\omega'$ (\AvNRef, p.~203, Lemma~1).

\ResultHeading{A.30}{A 30.}
Let $(e_i)$ be an orthonormal basis of $H$. For $T\in\mathcal L(H)^+$,
put
\[
\operatorname{Tr}(T)=\sum_i(Te_i\mid e_i).
\]
Then $\operatorname{Tr}$ is a faithful normal semifinite trace on
$\mathcal L(H)^+$, independent of the orthonormal basis chosen
(\AvNRef, p.~94, Theorem~5).

\ResultHeading{A.31}{A 31.}
A von Neumann algebra $\mathcal A$ is called \emph{finite}
(respectively, \emph{semifinite}) if, for every nonzero
$T\in\mathcal A^+$, there exists a finite (respectively, semifinite)
normal trace $f$ on $\mathcal A^+$ such that $f(T)\ne0$. It is called
\emph{properly infinite} (respectively, \emph{purely infinite}) if there
is no finite (respectively, semifinite) normal trace on $\mathcal A^+$
other than $0$ (\AvNRef, p.~97, Definition~5).

On a semifinite von Neumann algebra there exists a faithful normal
semifinite trace (\AvNRef, p.~99, Proposition~9).

\ResultHeading{A.32}{A 32.}
Let $\mathcal F$ be a semifinite factor, and let $f$ be a faithful normal
semifinite trace on $\mathcal F^+$. Every normal trace on $\mathcal F^+$
has the form $\lambda f$, where $\lambda\in[0,+\infty]$, with the
convention $0\cdot(+\infty)=0$ (\AvNRef, p.~81, Corollary~2, and p.~90,
Corollary). The ideals $\mathfrak m_f,\mathfrak n_f$ are independent of
the choice of $f$. The elements of $\mathfrak m_f$ (respectively,
$\mathfrak n_f$) are called the \emph{trace elements} (respectively, the
\emph{Hilbert--Schmidt elements}) relative to the semifinite factor
$\mathcal F$ (\AvNRef, p.~101). When $\mathcal F=\mathcal L(H)$, one
simply speaks of trace-class operators and Hilbert--Schmidt operators. If
$T,T'\in\mathcal L(H)$ and $T$ is trace-class, then
\[
|\operatorname{Tr}(TT')|\leq(\operatorname{Tr}|T|)\,\|T'\|
\]
(\AvNRef, p.~106).

\ResultHeading{A.33}{A 33.}
Let $\mathcal F$ be a factor. If $\mathcal F^+$ carries a nonzero finite
normal trace $f$, then $\mathcal F$ is finite (for $f$ is faithful by
\AvNRef, p.~83, Corollary~2).

\ResultHeading{A.34}{A 34.}
Let $\mathcal A$ be a von Neumann algebra. A projection $E$ of
$\mathcal A$ is called \emph{minimal} if $E\ne0$ and the only projections
of $\mathcal A$ majorized by $E$ are $0$ and $E$. Equivalently, $E\ne0$
and the reduced algebra $\mathcal A_E$ contains only scalar operators
(\AvNRef, p.~123).

A projection $E$ of $\mathcal A$ is called \emph{abelian} if the reduced
algebra $\mathcal A_E$ is commutative
(\AvNRef, p.~123, Definition~3).

\ResultHeading{A.35}{A 35.}
Let $\mathcal A$ be a von Neumann algebra. The following conditions on
$\mathcal A$ are equivalent: (i) $\mathcal A$ is isomorphic to a von
Neumann algebra $\mathcal B$ whose commutant $\mathcal B'$ is commutative;
(ii) every nonzero projection of $\mathcal A$ majorizes a nonzero abelian
projection; (iii) $\mathcal A$ has an abelian projection whose central
support is $1$. If these conditions hold,
% Source PDF physical page 348
$\mathcal A$ is said to be of \emph{type I}. Every type I algebra is
semifinite (\AvNRef, p.~122, Proposition~2; p.~123, Theorem~1; p.~126,
no.~4).

\ResultHeading{A.36}{A 36.}
Let $\mathcal A$ be a von Neumann algebra on $H$. The following conditions
are equivalent: (i) $\mathcal A$ is a type I factor; (ii) $\mathcal A$ is
a factor and has minimal projections; (iii) there are Hilbert spaces
$H_1,H_2$ and an isomorphism of $H$ onto $H_1\otimes H_2$ which transforms
$\mathcal A$ into $\mathcal L(H_1)\otimes\mathbb C_{H_2}$; (iv) there are
a family $(K_i)$ of closed vector subspaces whose Hilbert sum is $H$ and
isomorphisms $U_i$ from one Hilbert space $K$ onto the $K_i$, with the
following property: $\mathcal A$ is the set of $T\in\mathcal L(H)$ reduced
by the $K_i$ for which all the $T|K_i$ are transported, by the $U_i$, from
one and the same element of $\mathcal L(K)$
(\AvNRef, p.~124, Corollary~3, and pp.~21--23).

\ResultHeading{A.37}{A 37.}
Every automorphism of a type I factor $\mathcal A$ is inner and is defined
by a unitary element of $\mathcal A$ (\AvNRef, p.~241, Corollary~2; this
may also be regarded as a consequence of \XRef{4.1.5}{4.1.5}).

\ResultHeading{A.38}{A 38.}
A von Neumann algebra $\mathcal A$ is called \emph{continuous} if there is
no nonzero projection $E\in\mathcal A\cap\mathcal A'$ such that
$\mathcal A_E$ is of type I. It is called of \emph{type II} if it is
semifinite and continuous. A purely infinite algebra is continuous; such
an algebra is also called of \emph{type III}. If the space on which
$\mathcal A$ acts is nonzero, $\mathcal A$ cannot be of two different
types (\AvNRef, pp.~121, 122, 126).

\ResultHeading{A.39}{A 39.}
Let $\mathcal A$ be a von Neumann algebra. There are projections
$E_I,E_{II},E_{III}$ in $\mathcal A\cap\mathcal A'$, characterized by the
following properties: $E_I,E_{II},E_{III}$ are mutually orthogonal and
have sum $1$, while $\mathcal A_{E_I}$ is of type I,
$\mathcal A_{E_{II}}$ of type II, and $\mathcal A_{E_{III}}$ of type III.

The projection $E_I$ (respectively, $E_{II},E_{III}$) is the greatest
projection $E\in\mathcal A\cap\mathcal A'$ for which $\mathcal A_E$ is of
type I (respectively, II, III); $E_I+E_{II}$ is the greatest projection
$E\in\mathcal A\cap\mathcal A'$ for which $\mathcal A_E$ is semifinite;
and $E_{II}+E_{III}$ is the greatest projection
$E\in\mathcal A\cap\mathcal A'$ for which $\mathcal A_E$ is continuous
(\AvNRef, pp.~98 and 122).

\ResultHeading{A.40}{A 40.}
A finite type II algebra is called a \emph{type $II_1$ von Neumann
algebra}; a properly infinite type II algebra is called a \emph{type
$II_\infty$ von Neumann algebra}. If $\mathcal A$ is a von Neumann
algebra, there are projections $E_{II_1},E_{II_\infty}$ characterized as
follows: they are orthogonal and have sum $E_{II}$, and the corresponding
algebras induced by $\mathcal A$ are of types $II_1$ and $II_\infty$
(\AvNRef, pp.~98 and 122).

\ResultHeading{A.41}{A 41.}
Let $\mathcal A$ be a von Neumann algebra and let $E,F$ be two projections
of $\mathcal A$. They are called \emph{equivalent relative to
$\mathcal A$}, written $E\sim F$, if there is a $U\in\mathcal A$ such that
$U^*U=E$ and $UU^*=F$. We write $E\prec F$, or $F\succ E$, if there is a
projection of $\mathcal A$, equivalent to $E$, which is majorized by $F$
(\AvNRef, p.~215, Definition~1).

% Source PDF physical page 349
\ResultHeading{A.42}{A 42.}
The relation $E\sim F$ is an equivalence relation. The relation
$E\prec F$ is a preorder. If $E\prec F$ and $F\prec E$, then $E\sim F$
(\AvNRef, pp.~215--216).

\ResultHeading{A.43}{A 43.}
Equivalent projections have the same central support
(\AvNRef, p.~215).

\ResultHeading{A.44}{A 44.}
Let $E,F$ be two projections of $\mathcal A$, and let $E_1,F_1$ be their
central supports. If $E_1$ and $F_1$ are not orthogonal, there are
equivalent projections $E',F'$ of $\mathcal A$ majorized respectively by
$E,F$ (\AvNRef, p.~217, Lemma~1).

\ResultHeading{A.45}{A 45.}
Let $H$ be the space on which $\mathcal A$ acts, and let
$T\in\mathcal A$. The projections onto $H\ominus\ker T$ and onto
$\overline{T(H)}$ belong to $\mathcal A$ and are equivalent
(\AvNRef, p.~216, Proposition~2).

\ResultHeading{A.46}{A 46.}
If $\mathcal A$ is a factor and $E,F$ are projections of $\mathcal A$,
then either $E\prec F$ or $F\prec E$
(\AvNRef, p.~218, Corollary~1).

\ResultHeading{A.47}{A 47.}
A von Neumann algebra $\mathcal A$ is called \emph{homogeneous} if it
contains a family $(E_i)$ of abelian, equivalent, pairwise orthogonal
projections whose sum is $1$. Such an algebra is of type I. The cardinal
$n$ of the family $(E_i)$ depends only on $\mathcal A$, and $\mathcal A$
is said to be of \emph{type $I_n$}. Equivalently,
$\mathcal A=\mathcal B\otimes\mathcal L(K)$, where $\mathcal B$ is a
commutative von Neumann algebra and $K$ is a Hilbert space of dimension
$n$ (\AvNRef, pp.~239--240).

\ResultHeading{A.48}{A 48.}
A von Neumann algebra $\mathcal A$ is continuous if and only if every
projection of $\mathcal A$ is the sum of two orthogonal equivalent
projections of $\mathcal A$ (\AvNRef, p.~219, Corollary~3).

\ResultHeading{A.49}{A 49.}
A factor is of exactly one of the types $I$, $II_1$, $II_\infty$, and
$III$ (except on the space $H=0$). If it is of type I, it is of type $I_n$
for a uniquely determined $n$. The only factors which are
finite-dimensional as vector spaces are those of type $I_n$ with
$n<+\infty$. (This follows, for example, from \XRef{A.48}{A 48}.)

\ResultHeading{A.50}{A 50.}
Let $\mathcal A$ be a type I von Neumann algebra. There is a family
$(E_i)_{i\in I}$ of nonzero, pairwise orthogonal projections of
$\mathcal A\cap\mathcal A'$, with sum $1$, such that the
$\mathcal A_{E_i}$ are of types $I_{n_i}$ with the $n_i$ pairwise
distinct. The $E_i$ are uniquely determined up to a bijection of the
index set (\AvNRef, p.~240).

\ResultHeading{A.51}{A 51.}
Every isomorphism between type III von Neumann algebras acting on
separable spaces is spatial (\AvNRef, p.~301, Corollary~7).

\ResultHeading{A.52}{A 52.}
Let $\mathcal A$ be a von Neumann algebra. The projections $E_I$
(respectively, $E_{II},E_{III}$) associated with $\mathcal A$ and with
$\mathcal A'$ are equal. The algebra $\mathcal A$ is of type I
% Source PDF physical page 350
(respectively, of type II, of type III, semifinite, continuous) if and
only if $\mathcal A'$ is of type I (respectively, of type II, of type III,
semifinite, continuous) (\AvNRef, p.~104, Corollaries~1, 2, 3; p.~123,
Theorem~1; p.~124, Corollary~1).

\ResultHeading{A.53}{A 53.}
Let $\mathcal A$ be a von Neumann algebra and $E$ a projection of
$\mathcal A$ or of $\mathcal A'$. If $\mathcal A$ is of type I
(respectively, of type II, of type III, semifinite, continuous), then
$\mathcal A_E$ is of type I (respectively, of type II, of type III,
semifinite, continuous) (\AvNRef, p.~103, Proposition~11; p.~104,
Corollary~4; p.~125, Proposition~4).

\ResultHeading{A.54}{A 54.}
A \emph{Hilbert algebra} is an involutive algebra $A$ equipped with an
inner product which makes $A$ a separated pre-Hilbert space and for which
the following axioms hold:
\begin{enumerate}[label=\textup{(\roman*)}]
\item $(x\mid y)=(y^*\mid x^*)$ for $x,y\in A$;
\item $(xy\mid z)=(y\mid x^*z)$ for $x,y,z\in A$;
\item for every $x\in A$, the map $y\mapsto xy$ from $A$ into $A$ is
continuous;
\item the elements $xy$, $x,y\in A$, are dense in $A$.
\end{enumerate}
One has $(yx\mid z)=(y\mid zx^*)$ for $x,y,z\in A$, and the map
$y\mapsto yx$ from $A$ into $A$ is continuous.

Let $H$ be the Hilbert completion of $A$, let $J$ be the involution of $H$
which extends the map $x\mapsto x^*$ on $A$, and let $U_x,V_x$ be the
elements of $\mathcal L(H)$ which extend left and right multiplication by
$x$ on $A$. Then $x\mapsto U_x$ (respectively, $x\mapsto V_x$) is a
nondegenerate representation of $A$ (respectively, of the opposite
algebra) on $H$. The $U_x$ commute with the $V_y$. For every $x\in A$,
$JU_xJ=V_{x^*}$. The weak closure of the set of $U_x$ (respectively,
$V_x$) is a von Neumann algebra $\mathfrak U(A)$ (respectively,
$\mathfrak V(A)$), called the left-associated (respectively,
right-associated) algebra of $A$. One has
\[
\mathfrak V(A)=\mathfrak U(A)'=J\mathfrak U(A)J,\qquad
\mathfrak U(A)=\mathfrak V(A)'=J\mathfrak V(A)J.
\]
The maps $x\mapsto U_x$ and $x\mapsto V_x$ are called the canonical maps
of $A$ into $\mathfrak U(A)$ and $\mathfrak V(A)$. If
$C\in\mathfrak U(A)\cap\mathfrak V(A)$, then $JCJ=C^*$
(\AvNRef, pp.~69--73).

\ResultHeading{A.55}{A 55.}
Let $A$ be a Hilbert algebra and $A'$ an involutive subalgebra of $A$.
Then $A'$ is a Hilbert algebra.

If the $U_x$, $x\in A'$, are strongly dense in $\mathfrak U(A)$, then
$A'$ is dense in $A$.

If $A'$ is dense in $A$, then
$\mathfrak U(A')=\mathfrak U(A)$ and
$\mathfrak V(A')=\mathfrak V(A)$
(\AvNRef, p.~72, Proposition~1).

\ResultHeading{A.56}{A 56.}
Retain the preceding notation. For $a\in H$, the following conditions are
equivalent: (i) the map $x\mapsto V_xa$ from $A$ into $H$ is continuous;
(ii) the map $x\mapsto U_xa$ from $A$ into $H$ is continuous. Then $a$ is
called \emph{bounded} relative to $A$. Let $U_a$ (respectively, $V_a$) be
the continuous extension to $H$ of the map $x\mapsto V_xa$
(respectively, $x\mapsto U_xa$) from $A$ into $H$. If $a\in A$, then $a$
is bounded and the notations $U_a,V_a$ agree with the preceding ones. An
element $a\in H$ is bounded if and only if $Ja$ is bounded. The Hilbert
algebra $A$ is called \emph{completed} if every bounded element belongs to
$A$ (\AvNRef, pp.~72--74).

% Source PDF physical page 351
\ResultHeading{A.57}{A 57.}
Let $A$ be a Hilbert algebra dense in the Hilbert space $H$. If $a,b\in H$
are bounded, set
\[
U_ab=V_ba=ab,\qquad Ja=a^*.
\]
Then the set of bounded elements becomes a Hilbert algebra $A'$ containing
$A$ as a dense Hilbert subalgebra. The Hilbert algebra $A'$ is completed,
and
$\mathfrak U(A')=\mathfrak U(A)$,
$\mathfrak V(A')=\mathfrak V(A)$
(\AvNRef, p.~74).

\ResultHeading{A.58}{A 58.}
The maps $a\mapsto U_a$ and $a\mapsto V_a$, $a\in A'$, are injective
(\AvNRef, p.~70).

\ResultHeading{A.59}{A 59.}
The operators $U_a$ (respectively, $V_a$), where $a\in A'$, form a
self-adjoint two-sided ideal of $\mathfrak U(A)$ (respectively,
$\mathfrak V(A)$). For $a\in A'$, $T\in\mathfrak U(A)$, and
$T'\in\mathfrak V(A)$, one has
\[
Ta\in A',\qquad T'a\in A',\qquad
U_{Ta}=TU_a,\qquad V_{T'a}=T'V_a
\]
(\AvNRef, p.~72, Proposition~2).

\ResultHeading{A.60}{A 60.}
For $S\in\mathfrak U(A)^+$, put $f(S)=(a\mid a)$ if
$S^{1/2}=U_a$ with $a\in A'$, and put $f(S)=+\infty$ otherwise. Then $f$
is a faithful normal semifinite trace on $\mathfrak U(A)^+$. The two-sided
ideal of $\mathfrak U(A)$ formed by the $U_a$, $a\in A'$, is the set of
$T\in\mathfrak U(A)$ such that $f(T^*T)<+\infty$, that is,
$\mathfrak n_f$. If $a,b\in A'$, then
\[
(a\mid b)=f'(U_b^*U_a),
\]
where $f'$ denotes the linear extension of the trace to
$\mathfrak m_f$. There are analogous results for $\mathfrak V(A)$. The
von Neumann algebras $\mathfrak U(A)$ and $\mathfrak V(A)$ are
semifinite. The trace $f$ on $\mathfrak U(A)^+$ and the analogous trace
on $\mathfrak V(A)^+$ are called the \emph{natural traces} defined by $A$.
These traces are unchanged when $A$ is replaced by $A'$
(\AvNRef, pp.~85--88).

\ResultHeading{A.61}{A 61.}
Let $\mathcal A$ be a von Neumann algebra, let $f$ be a faithful normal
semifinite trace on $\mathcal A^+$, and let $f'$ be the linear extension
of $f$ to $\mathfrak m_f$. With the inner product
$(S\mid T)=f'(T^*S)$, $\mathfrak n_f$ is a completed Hilbert algebra. Let
$H_f$ be the Hilbert completion of $\mathfrak n_f$, and let $J_f$ be the
continuous involution of $H_f$ extending that of $\mathfrak n_f$. For
$R\in\mathcal A$, left and right multiplication by $R$ on
$\mathfrak n_f$ extend to continuous linear operators
$\lambda_f(R),\rho_f(R)$, and
$\rho_f(R)=J_f\lambda_f(R)J_f$. The map $\lambda_f$ is an isomorphism of
$\mathcal A$ onto $\mathfrak U(\mathfrak n_f)$ which extends the
canonical map of $\mathfrak n_f$ into $\mathfrak U(\mathfrak n_f)$ and
transforms $f$ into the natural trace defined on
$\mathfrak U(\mathfrak n_f)^+$ by $\mathfrak n_f$
(\AvNRef, p.~88, Theorem~2).

\ResultHeading{A.62}{A 62.}
Let $A$ be a Hilbert algebra dense in a Hilbert space $H$, and let $a\in H$.
The following conditions are equivalent:
(i) $(a\mid xy)=(a\mid yx)$ for all $x,y\in A$;
(ii) $Ta=JT^*Ja$ for every $T\in\mathfrak U(A)$.
If $a$ satisfies these conditions, it is called \emph{central relative to
$A$}. Elements of $A$ which commute with every element of $A$ are central
relative to $A$. Let $Z$ be the set of elements central relative to $A$.
The transforms of the elements of $Z$ by $\mathfrak U(A)$
(respectively, $\mathfrak V(A)$) span a closed vector subspace $Y$
(respectively, $Y'$) of $H$.
% Source PDF physical page 352
One has $Y=Y'$, and
$P_Y\in\mathfrak U(A)\cap\mathfrak V(A)$ is called the
\emph{characteristic projection} of $A$ (\AvNRef, pp.~75--76).

\ResultHeading{A.63}{A 63.}
With the notation of \XRef{A.62}{A 62}, the following conditions are
equivalent: (i) $\mathfrak U(A)$ is a finite von Neumann algebra;
(ii) $\mathfrak V(A)$ is a finite von Neumann algebra; (iii) $P_Y=1$
(\AvNRef, p.~100, Theorem~6).

\ResultHeading{A.64}{A 64.}
Let $(A_i)$ be a family of Hilbert algebras, and let $H_i$ be the Hilbert
completion of $A_i$. Let $A$ be the direct sum of the $A_i$, which is in
the evident way an involutive algebra and a pre-Hilbert space. Then $A$ is
a Hilbert algebra, called the \emph{direct-sum Hilbert algebra} of the
$A_i$. The Hilbert completion $H$ of $A$ is the Hilbert sum of the $H_i$.
One has
\[
\mathfrak U(A)=\prod_i\mathfrak U(A_i),\qquad
\mathfrak V(A)=\prod_i\mathfrak V(A_i).
\]
If $E_i$ is the projection of $H$ onto $H_i$, then
\[
E_i\in\mathfrak U(A)\cap\mathfrak V(A),\qquad
\mathfrak U(A_i)=\mathfrak U(A)_{E_i},\qquad
\mathfrak V(A_i)=\mathfrak V(A)_{E_i}
\]
(\AvNRef, p.~76).

\ResultHeading{A.65}{A 65.}
Let $A$ be a completed Hilbert algebra with Hilbert completion $H$, let
$(E_i)$ be a family of pairwise orthogonal projections in
$\mathfrak U(A)\cap\mathfrak V(A)$ with sum $1$, and put
$H_i=E_i(H)$ and $A_i=E_i(A)=H_i\cap A$. Then $A_i$ is an involutive
subalgebra of $A$ and a completed Hilbert algebra dense in $H_i$. The
$A_i$ annihilate one another and hence are two-sided ideals of $A$. One
has
\[
\mathfrak U(A_i)=\mathfrak U(A)_{E_i},\qquad
\mathfrak V(A_i)=\mathfrak V(A)_{E_i}.
\]
The direct-sum Hilbert algebra of the $A_i$ is a dense involutive
subalgebra of $A$. The natural trace of $\mathfrak U(A)^+$ induces on
$\mathfrak U(A)_{E_i}^+=\mathfrak U(A_i)^+$ the natural trace of
$\mathfrak U(A_i)^+$ (\AvNRef, p.~77).

\ResultHeading{A.66}{A 66.}
Let $A$ be the involutive algebra of Hilbert--Schmidt operators on the
Hilbert space $H$. With the inner product
$(T\mid T')=\operatorname{Tr}(T'^*T)$, it is a complete Hilbert algebra.
Let $\overline H$ be the Hilbert space conjugate to $H$. There is a unique
isomorphism $\Phi$ from the Hilbert space $A$ onto
$H\otimes\overline H$ which transforms the operator
$\xi\mapsto(\xi\mid\eta)\zeta$, $\zeta,\eta\in H$, into
$\zeta\otimes\overline\eta$. Identifying $A$ with
$H\otimes\overline H$ by $\Phi$, one has, for $\xi,\eta\in H$ and
$S\in\mathcal L(H)$,
\[
(\xi\otimes\overline\eta)^*=\eta\otimes\overline\xi,\qquad
S\cdot(\xi\otimes\overline\eta)=S\xi\otimes\overline\eta,\qquad
(\xi\otimes\overline\eta)\cdot S
=\xi\otimes\overline{S^*\eta}.
\]
(\AvNRef, pp.~95--97.)

\ResultHeading{A.67}{A 67.}
Let $A$ be a completed Hilbert algebra such that $\mathfrak U(A)$ is a
factor. The following conditions are equivalent:
\begin{enumerate}[label=\textup{(\roman*)}]
\item $A$ is complete;

% Source PDF physical page 353
\item $\mathfrak U(A)$ is of type I;
\item after multiplication of the norm of $A$ by a constant, $A$ is
isomorphic to the Hilbert algebra of Hilbert--Schmidt operators on a
suitable Hilbert space.
\end{enumerate}
The preceding Hilbert space is determined by $A$ up to isomorphism
(\AvNRef, p.~127).

\ResultHeading{A.68}{A 68.}
Let $A$ be a complete Hilbert algebra, and let $(E_i)$ be the family of
minimal projections of $\mathfrak U(A)\cap\mathfrak V(A)$.
\begin{enumerate}[label=\textup{(\roman*)}]
\item The Hilbert space $A$ is the Hilbert sum of the $A_i=E_i(A)$.
\item The $A_i$ are pairwise annihilating self-adjoint two-sided ideals of
$A$.
\item After multiplication of the norm of $A_i$ by a constant, $A_i$ is
isomorphic to the Hilbert algebra of Hilbert--Schmidt operators on a
suitable Hilbert space.
\item If $\mathfrak U(A)$ and $\mathfrak V(A)$ are finite von Neumann
algebras, each $A_i$ is finite-dimensional.
\end{enumerate}
(\AvNRef, p.~127, Proposition~7.)

\ResultHeading{A.69}{A 69.}
Let $Z$ be a Borel space (\XRef{B.19}{B 19}) and let $\mu$ be a positive
measure on $Z$ (\XRef{B.30}{B 30}).

A \emph{$\mu$-measurable field of Hilbert spaces} on $Z$ is a pair
\[
\mathcal E=\bigl((H(\zeta))_{\zeta\in Z},\Gamma\bigr),
\]
where $(H(\zeta))_{\zeta\in Z}$ is a family of Hilbert spaces indexed by
$Z$ and $\Gamma$ is a set of vector fields satisfying the following
conditions:
\begin{enumerate}[label=\textup{(\roman*)}]
\item $\Gamma$ is a vector subspace of $\prod_{\zeta\in Z}H(\zeta)$;
\item there is a sequence $(x_1,x_2,\ldots)$ of elements of $\Gamma$ such
that, for every $\zeta\in Z$, the $x_n(\zeta)$ form a total sequence in
$H(\zeta)$;
\item for every $x\in\Gamma$, the function
$\zeta\mapsto\|x(\zeta)\|$ is $\mu$-measurable;
\item let $x$ be a vector field; if, for every $y\in\Gamma$, the function
$\zeta\mapsto(x(\zeta)\mid y(\zeta))$ is $\mu$-measurable, then
$x\in\Gamma$.
\end{enumerate}
The elements of $\Gamma$ are then called the \emph{measurable vector
fields} of $\mathcal E$. If $x,y\in\Gamma$, the function
$\zeta\mapsto(x(\zeta)\mid y(\zeta))$ is measurable
(\AvNRef, p.~142).

% Source PDF physical page 354
\ResultHeading{A.70}{A 70.}
Let $Z,Z'$ be Borel spaces; let $\mu,\mu'$ be positive measures on
$Z,Z'$; and let
\[
\mathcal E=((H(\zeta)),\Gamma),\qquad
\mathcal E'=((H'(\zeta')),\Gamma')
\]
be $\mu$- and $\mu'$-measurable fields of Hilbert spaces. Let
$\eta:Z\to Z'$ be a Borel isomorphism which transforms $\mu$ into $\mu'$.
An \emph{$\eta$-isomorphism} from $\mathcal E$ onto $\mathcal E'$ is a
family $(V(\zeta))_{\zeta\in Z}$ having the following properties:
(i) for every $\zeta\in Z$, $V(\zeta)$ is an isomorphism of $H(\zeta)$
onto $H'(\eta(\zeta))$; (ii) a vector field
$\zeta\mapsto x(\zeta)\in H(\zeta)$ on $Z$ is $\mu$-measurable if and only
if the field
\[
\eta(\zeta)\longmapsto V(\zeta)x(\zeta)\in H'(\eta(\zeta))
\]
on $Z'$ is $\mu'$-measurable (\AvNRef, p.~143).

If $Z=Z'$ and $\mu=\mu'$, an $\eta$-isomorphism for which $\eta$ is the
identity map of $Z$ is simply called an isomorphism from $\mathcal E$ onto
$\mathcal E'$ (\AvNRef, p.~143).

\ResultHeading{A.71}{A 71.}
Let $Z$ be a Borel space, let $\mu$ be a positive measure on $Z$, and let
$H_0$ be a separable Hilbert space. Put $H(\zeta)=H_0$ for every
$\zeta\in Z$. Let $\Gamma$ be the set of measurable maps of $Z$ into
$H_0$, that is, the maps $x:Z\to H_0$ such that
$\zeta\mapsto(x(\zeta)\mid a)$ is measurable for every $a\in H_0$. Then
$\mathcal E=((H(\zeta)),\Gamma)$ is a measurable field of Hilbert spaces
on $Z$, called the \emph{constant field} defined by $H_0$
(\AvNRef, p.~143). A field isomorphic to a constant field
(cf. \XRef{A.70}{A 70}) is called \emph{trivial} (\AvNRef, p.~143).

\ResultHeading{A.72}{A 72.}
Retain the notation of \XRef{A.69}{A 69}. Let
$p=1,2,\ldots,\aleph_0$. The set
\[
Z_p=\{\zeta\in Z:\dim H(\zeta)=p\}
\]
is measurable. Let $\mathcal E_p$ be the restriction, defined in the
evident way, of $\mathcal E$ to $Z_p$. Then every $\mathcal E_p$ is
trivial (\AvNRef, pp.~144--145, Propositions~1 and 3).

\ResultHeading{A.73}{A 73.}
Let $\mathcal E=((H(\zeta)),\Gamma)$ be a $\mu$-measurable field of Hilbert
spaces on $Z$. A vector field $x$ is called \emph{square-integrable} if
$x\in\Gamma$ and
\[
\int_Z\|x(\zeta)\|^2\,d\mu(\zeta)<+\infty.
\]
If $x,y$ are square-integrable, so are $x+y$ and $\lambda x$, and the
function $\zeta\mapsto(x(\zeta)\mid y(\zeta))$ is integrable; put
\[
(x\mid y)=\int_Z(x(\zeta)\mid y(\zeta))\,d\mu(\zeta).
\]
The square-integrable vector fields then form a Hilbert space $H$, called
the \emph{Hilbert integral} of the $H(\zeta)$ and denoted
\[
\int_Z^\oplus H(\zeta)\,d\mu(\zeta).
\]
(Fields equal almost everywhere are identified.) An element $x$ of this
space is also written
$\int_Z^\oplus x(\zeta)\,d\mu(\zeta)$. If $\mu$ is standard, $H$ is
separable (\AvNRef, pp.~145--146 and p.~149, Corollary).

% Source PDF physical page 355
\ResultHeading{A.74}{A 74.}
Use the notation of \XRef{A.70}{A 70}. The map which transforms the field
\[
x\in H=\int_Z^\oplus H(\zeta)\,d\mu(\zeta)
\]
into the field
\[
\eta(\zeta)\longmapsto V(\zeta)x(\zeta)
\in H'=\int_{Z'}^\oplus H'(\zeta')\,d\mu'(\zeta')
\]
is an isomorphism of $H$ onto $H'$, denoted
$\int_Z^\oplus V(\zeta)\,d\mu(\zeta)$
(\AvNRef, p.~159, Definition~2).

\ResultHeading{A.75}{A 75.}
Let $\mathcal E=((H(\zeta)),\Gamma)$ be a $\mu$-measurable field of Hilbert
spaces on $Z$, let $\widetilde\mu$ be a positive measure on $Z$ equivalent
to $\mu$, and put
\[
\rho(\zeta)=\frac{d\widetilde\mu(\zeta)}{d\mu(\zeta)}.
\]
The map which assigns to every
$x\in H=\int_Z^\oplus H(\zeta)\,d\mu(\zeta)$ the field
\[
\zeta\longmapsto\rho(\zeta)^{-1/2}x(\zeta)
\in\widetilde H=\int_Z^\oplus H(\zeta)\,d\widetilde\mu(\zeta)
\]
is an isomorphism from $H$ onto $\widetilde H$, called the
\emph{canonical isomorphism} (\AvNRef, p.~148).

\ResultHeading{A.76}{A 76.}
Let $Z$ be a Borel space, let $\mu$ be a positive measure on $Z$, let
$H_0$ be a separable Hilbert space, and let
$\mathcal E=((H(\zeta)),\Gamma)$ be the constant field on $Z$ defined by
$H_0$. There is a unique isomorphism (called canonical) from
$L^2(Z,\mu)\otimes H_0$ onto
$\int_Z^\oplus H(\zeta)\,d\mu(\zeta)$ which, for every
$f\in L^2(Z,\mu)$ and every $\xi\in H_0$, transforms $f\otimes\xi$ into
the field $\zeta\mapsto f(\zeta)\xi$
(\AvNRef, p.~153, Corollary).

\ResultHeading{A.77}{A 77.}
Let $\mathcal E=((H(\zeta)),\Gamma)$ be a measurable field of Hilbert
spaces on $Z$. For every $\zeta\in Z$, let
$T(\zeta)\in\mathcal L(H(\zeta))$. The field
$\zeta\mapsto T(\zeta)$ is called a \emph{measurable field of operators}
if, for every $x\in\Gamma$, the field
$\zeta\mapsto T(\zeta)x(\zeta)$ is measurable. In that case the function
$\zeta\mapsto\|T(\zeta)\|$ is measurable. Suppose, in addition, that this
function is essentially bounded; the field $\zeta\mapsto T(\zeta)$ is
then called \emph{essentially bounded}. Let
$H=\int_Z^\oplus H(\zeta)\,d\mu(\zeta)$. For every $x\in H$, the field
$Tx:\zeta\mapsto T(\zeta)x(\zeta)$ belongs to $H$. One has
$T\in\mathcal L(H)$, and $\|T\|$ is the essential supremum of
$\zeta\mapsto\|T(\zeta)\|$. Put
\[
T=\int_Z^\oplus T(\zeta)\,d\mu(\zeta).
\]
The $T(\zeta)$ are determined by $T$ up to null sets. Operators on $H$ of
the form $\int_Z^\oplus T(\zeta)\,d\mu(\zeta)$ are called
\emph{decomposable} (\AvNRef, pp.~156--159).

% Source PDF physical page 356
\ResultHeading{A.78}{A 78.}
If
\[
S=\int_Z^\oplus S(\zeta)\,d\mu(\zeta),\qquad
T=\int_Z^\oplus T(\zeta)\,d\mu(\zeta)
\]
are decomposable, then
\begin{align*}
S+T&=\int_Z^\oplus\bigl(S(\zeta)+T(\zeta)\bigr)\,d\mu(\zeta),&
ST&=\int_Z^\oplus S(\zeta)T(\zeta)\,d\mu(\zeta),\\
\lambda S&=\int_Z^\oplus\lambda S(\zeta)\,d\mu(\zeta),&
S^*&=\int_Z^\oplus S(\zeta)^*\,d\mu(\zeta).
\end{align*}
(\AvNRef, p.~159, Proposition~3.)

\ResultHeading{A.79}{A 79.}
Let
\[
T_i=\int_Z^\oplus T_i(\zeta)\,d\mu(\zeta)\quad(i=1,2,\ldots),
\qquad
T=\int_Z^\oplus T(\zeta)\,d\mu(\zeta)
\]
be decomposable operators. If $T_i$ converges strongly to $T$, there is a
subsequence $(T_{n_i})$ such that $T_{n_i}(\zeta)$ converges strongly to
$T(\zeta)$ almost everywhere. If $T_i(\zeta)$ converges strongly to
$T(\zeta)$ almost everywhere and $\sup_i\|T_i\|<+\infty$, then $T_i$
converges strongly to $T$ (\AvNRef, p.~160, Proposition~4).

\ResultHeading{A.80}{A 80.}
Operators of the form
$\int_Z^\oplus T(\zeta)\,d\mu(\zeta)$, where $T(\zeta)$ is scalar for
every $\zeta$, are called \emph{diagonalizable operators}. Let
$\mathfrak Z$ be the set of diagonalizable operators. Then $\mathfrak Z$
is a commutative von Neumann algebra, and $\mathfrak Z'$ is the set of
decomposable operators (\AvNRef, p.~162, Definition~3; p.~163,
Proposition~7; p.~164, Corollary).

\ResultHeading{A.81}{A 81.}
Suppose that $\mathcal E=((H(\zeta)),\Gamma)$ is the constant field
defined by $H_0$. Let $S\in\mathcal L(H_0)$ and put $S(\zeta)=S$ for
every $\zeta$. The canonical isomorphism from
$L^2(Z,\mu)\otimes H_0$ onto
$\int_Z^\oplus H(\zeta)\,d\mu(\zeta)$ transforms $1\otimes S$ into
$\int_Z^\oplus S(\zeta)\,d\mu(\zeta)$
(\AvNRef, pp.~169--170).

\ResultHeading{A.82}{A 82.}
Let $Z,Z'$ be Borel spaces and let $\mu,\mu'$ be positive measures on
$Z,Z'$. Let
\[
\mathcal E=((H(\zeta))_{\zeta\in Z},\Gamma),\qquad
\mathcal E'=((H'(\zeta'))_{\zeta'\in Z'},\Gamma')
\]
be $\mu$- and $\mu'$-measurable fields of Hilbert spaces. Let $I$ be a
countable set, and let $\eta$ be a Borel isomorphism of $Z$ onto $Z'$
which transforms $\mu$ into $\mu'$. For every $i\in I$, let
$\zeta\mapsto T_i(\zeta)$ be a measurable field of operators relative to
$\mathcal E$, and let $\zeta'\mapsto T'_i(\zeta')$ be a measurable field
of operators relative to $\mathcal E'$.
% Source PDF physical page 357
Suppose that, for every $\zeta\in Z$, there is an isomorphism from
$H(\zeta)$ onto $H'(\eta(\zeta))$ which, for every $i\in I$, transforms
$T_i(\zeta)$ into $T'_i(\eta(\zeta))$. Suppose $\mu$ and $\mu'$ are
standard. Under these hypotheses, there is an $\eta$-isomorphism
(\XRef{A.70}{A 70}) from $\mathcal E$ onto $\mathcal E'$ which, for every
$i\in I$, transforms the field $\zeta\mapsto T_i(\zeta)$ into the field
$\zeta\mapsto T'_i(\eta(\zeta))$.

(This reduces at once to the case $Z=Z'$, $\mu=\mu'$, and $\eta$ the
identity map of $Z$. Then see \AvNRef, p.~166, Lemma~2.)

\ResultHeading{A.83}{A 83.}
Let $\mathcal E=((H(\zeta)),\Gamma)$ be a $\mu$-measurable field of Hilbert
spaces on $Z$, let
\[
T_i=\int_Z^\oplus T_i(\zeta)\,d\mu(\zeta)\qquad(i\in I)
\]
be a family of decomposable operators on
$\int_Z^\oplus H(\zeta)\,d\mu(\zeta)$, let $H_0$ be a separable Hilbert
space, and let $(S_i)_{i\in I}$ be a family of elements of
$\mathcal L(H_0)$. Suppose that, for every $\zeta\in Z$, there is an
isomorphism of $H(\zeta)$ onto $H_0$ which transforms every
$T_i(\zeta)$ into $S_i$. Suppose $\mu$ is standard. Then there is an
isomorphism from
$\int_Z^\oplus H(\zeta)\,d\mu(\zeta)$ onto
$L^2(Z,\mu)\otimes H_0$ which, for every $i$, transforms $T_i$ into
$1\otimes S_i$ (\AvNRef, p.~167, Theorem~2).

\ResultHeading{A.84}{A 84.}
Let $H$ be a separable Hilbert space and let
$\mathfrak Z\subset\mathcal L(H)$ be a commutative von Neumann algebra.
There exist a standard Borel space $Z$, a bounded positive measure $\mu$
on $Z$, a measurable field $\zeta\mapsto H(\zeta)$ of nonzero Hilbert
spaces on $Z$, and an isomorphism from $H$ onto
$\int_Z^\oplus H(\zeta)\,d\mu(\zeta)$ which transforms $\mathfrak Z$ into
the algebra of diagonalizable operators
(\AvNRef, p.~210, Theorem~2).

\ResultHeading{A.85}{A 85.}
Let $Z$ be a Borel space, let $\mu$ be a standard positive measure on
$Z$, let $\mathcal E=(H(\zeta))$ be a $\mu$-measurable field of nonzero
Hilbert spaces on $Z$, put
\[
H=\int_Z^\oplus H(\zeta)\,d\mu(\zeta),
\]
and let $\mathfrak Z$ be the algebra of diagonalizable operators on $H$.
Define similarly $Z_1,\mu_1,\mathcal E_1=(H_1(\zeta_1)),H_1,\mathfrak Z_1$.
Let $U$ be an isomorphism from $H$ onto $H_1$ which transforms
$\mathfrak Z$ into $\mathfrak Z_1$. There exist: (i) a Borel
$\mu$-null set $N$ in $Z$ and a Borel $\mu_1$-null set $N_1$ in $Z_1$;
(ii) a Borel isomorphism $\eta$ from $Z\setminus N$ onto
$Z_1\setminus N_1$ which transforms $\mu$ into a measure
$\widetilde\mu_1$ equivalent to $\mu_1$; and (iii) an $\eta$-isomorphism
$(V(\zeta))$ from $\mathcal E|_{Z\setminus N}$ onto
$\mathcal E_1|_{Z_1\setminus N_1}$, such that $U$ is the composite of
$\int^\oplus V(\zeta)\,d\mu(\zeta)$ and the canonical isomorphism from
\[
\int_{Z_1}^\oplus H_1(\zeta_1)\,d\widetilde\mu_1(\zeta_1)
\quad\text{onto}\quad H_1
\]
(\AvNRef, p.~212, Theorem~4).

% Source PDF physical page 358
\ResultHeading{A.86}{A 86.}
Let $\mathcal E=((H(\zeta)),\Gamma)$ be a measurable field of Hilbert
spaces, put
\[
H=\int_Z^\oplus H(\zeta)\,d\mu(\zeta),
\qquad
T_i=\int_Z^\oplus T_i(\zeta)\,d\mu(\zeta),
\]
where $(T_i)$ is a sequence of decomposable operators, let $\mathcal A$ be
the von Neumann algebra generated by the $T_i$, let $\mathcal A(\zeta)$
be the von Neumann algebra generated by the $T_i(\zeta)$, and let
$\mathfrak Z$ be the algebra of diagonalizable operators. The algebra
$\mathfrak Z$ is a maximal commutative von Neumann subalgebra of
$\mathcal A'$ if and only if
$\mathcal A(\zeta)=\mathcal L(H(\zeta))$ almost everywhere
(\AvNRef, p.~172, Corollary~1).

\ResultHeading{A.87}{A 87.}
Let $\mathcal E=((H(\zeta)),\Gamma)$ be a measurable field of Hilbert
spaces on $Z$. For every $\zeta\in Z$, let $\mathcal A(\zeta)$ be a von
Neumann algebra on $H(\zeta)$. The field
$\zeta\mapsto\mathcal A(\zeta)$ is called a \emph{measurable field of von
Neumann algebras} on $Z$ if there is a sequence
\[
\zeta\mapsto T_1(\zeta),\quad
\zeta\mapsto T_2(\zeta),\quad\ldots
\]
of measurable fields of operators such that, for every $\zeta\in Z$,
$\mathcal A(\zeta)$ is the von Neumann algebra generated by the
$T_i(\zeta)$. The set of decomposable operators
\[
T=\int_Z^\oplus T(\zeta)\,d\mu(\zeta)
\quad\text{such that}\quad
T(\zeta)\in\mathcal A(\zeta)\ \text{for every }\zeta
\]
is a von Neumann algebra $\mathcal A$ on
$H=\int_Z^\oplus H(\zeta)\,d\mu(\zeta)$, denoted
\[
\mathcal A=\int_Z^\oplus\mathcal A(\zeta)\,d\mu(\zeta).
\]
The $\mathcal A(\zeta)$ are determined by $\mathcal A$ up to null sets.
Von Neumann algebras of this form are called \emph{decomposable}. Let
$\mathfrak Z$ be the algebra of diagonalizable operators. Suppose that
$\mu$ is standard. A von Neumann algebra $\mathcal A$ is decomposable if
and only if
\[
\mathfrak Z\subset\mathcal A\subset\mathfrak Z'
\]
(\AvNRef, pp.~173--174).

\ResultHeading{A.88}{A 88.}
Let $\mathcal E=((H(\zeta)),\Gamma)$ be a measurable field of Hilbert
spaces, let
$T_i=\int_Z^\oplus T_i(\zeta)\,d\mu(\zeta)$ be a sequence of decomposable
operators, let $\mathfrak Z$ be the algebra of diagonalizable operators,
let $\mathcal A(\zeta)$ be the von Neumann algebra generated by the
$T_i(\zeta)$, and let $\mathcal A$ be the von Neumann algebra generated
by $\mathfrak Z$ and the $T_i$. Then
\[
\mathcal A=\int_Z^\oplus\mathcal A(\zeta)\,d\mu(\zeta)
\]
(\AvNRef, p.~171, Theorem~1).

\ResultHeading{A.89}{A 89.}
Let
$\mathcal A=\int_Z^\oplus\mathcal A(\zeta)\,d\mu(\zeta)$ be a
decomposable von Neumann algebra, and let $\mathfrak Z$ be the algebra of
diagonalizable operators. If $\mathfrak Z$ is the center of $\mathcal A$,
then $\mathcal A(\zeta)$ is a factor almost everywhere
(\AvNRef, p.~174, Theorem~3).

\ResultHeading{A.90}{A 90.}
Let
$\mathcal A=\int_Z^\oplus\mathcal A(\zeta)\,d\mu(\zeta)$ be a
decomposable von Neumann algebra, and suppose $\mu$ is standard. The
algebra $\mathcal A$ is of type I if and only if $\mathcal A(\zeta)$ is
of type I almost everywhere
(\AvNRef, p.~183, Corollary~2 to Proposition~7).

% Source PDF physical page 359
\ResultHeading{A.91}{A 91.}
Let
\[
\mathcal A=\int_Z^\oplus\mathcal A(\zeta)\,d\mu(\zeta),\qquad
\mathcal A_1=\int_Z^\oplus\mathcal A_1(\zeta)\,d\mu(\zeta)
\]
be two decomposable von Neumann algebras. For every $\zeta\in Z$, let
$\Phi(\zeta)$ be an isomorphism from $\mathcal A(\zeta)$ onto
$\mathcal A_1(\zeta)$. The field of isomorphisms
$\zeta\mapsto\Phi(\zeta)$ is called \emph{measurable} if, for every
$T=\int_Z^\oplus T(\zeta)\,d\mu(\zeta)\in\mathcal A$, the field
$\zeta\mapsto\Phi(\zeta)(T(\zeta))$ is measurable. If
\[
\Phi(T)=\int_Z^\oplus\Phi(\zeta)(T(\zeta))\,d\mu(\zeta),
\]
then $\Phi$ is an isomorphism from $\mathcal A$ onto $\mathcal A_1$,
denoted $\int_Z^\oplus\Phi(\zeta)\,d\mu(\zeta)$
(\AvNRef, pp.~183 and 185).

\ResultHeading{A.92}{A 92.}
Let
$\mathcal A=\int_Z^\oplus\mathcal A(\zeta)\,d\mu(\zeta)$ be a
decomposable von Neumann algebra. For every $\zeta\in Z$, let $f(\zeta)$
be a trace on $\mathcal A(\zeta)^+$. The field of traces
$\zeta\mapsto f(\zeta)$ is called \emph{measurable} if, for every
measurable field of operators
$\zeta\mapsto T(\zeta)\in\mathcal A(\zeta)^+$, the function
$\zeta\mapsto f(\zeta)(T(\zeta))$ is measurable. If
$T=\int_Z^\oplus T(\zeta)\,d\mu(\zeta)\in\mathcal A^+$, put
\[
f(T)=\int_Z f(\zeta)(T(\zeta))\,d\mu(\zeta).
\]
Then $f$ is a trace on $\mathcal A^+$, denoted
$\int_Z^\oplus f(\zeta)\,d\mu(\zeta)$
(\AvNRef, p.~198).

\ResultHeading{A.93}{A 93.}
Let $\mathcal E=((H(\zeta)),\Gamma)$ be a measurable field of Hilbert
spaces on $Z$. For every $\zeta\in Z$, let $A(\zeta)$ be a Hilbert
algebra whose Hilbert completion is $H(\zeta)$. The field
$\zeta\mapsto A(\zeta)$ of Hilbert algebras is called
\emph{$\mu$-measurable} when the following conditions hold:
\begin{enumerate}[label=\textup{(\roman*)}]
\item if $x\in\Gamma$ and $x(\zeta)\in A(\zeta)$ for every $\zeta$, then
the field $x^*:\zeta\mapsto x(\zeta)^*$ is measurable;
\item if $x,y\in\Gamma$ and $x(\zeta),y(\zeta)\in A(\zeta)$ for every
$\zeta$, then the field
$xy:\zeta\mapsto x(\zeta)y(\zeta)$ is measurable;
\item there is a sequence $(x_1,x_2,\ldots)$ of elements of $\Gamma$ such
that $x_i(\zeta)\in A(\zeta)$ for every $i$ and every $\zeta$, and such
that $(x_1(\zeta),x_2(\zeta),\ldots)$ is total in $H(\zeta)$ for every
$\zeta$.
\end{enumerate}
(\AvNRef, p.~187, Definition~1.)

\ResultHeading{A.94}{A 94.}
Retain the notation of \XRef{A.93}{A 93}, and put
$H=\int_Z^\oplus H(\zeta)\,d\mu(\zeta)$. Let $A$ be the set of $x\in H$
such that $x(\zeta)\in A(\zeta)$ for every $\zeta$ and the field of
operators $\zeta\mapsto U_{x(\zeta)}$, which is automatically measurable,
is essentially bounded. Then, with the operations $x\mapsto x^*$,
$(x,y)\mapsto xy$, and the inner product induced by that of $H$, $A$ is a
Hilbert algebra dense in $H$, denoted
\[
A=\int_Z^\oplus A(\zeta)\,d\mu(\zeta).
\]
One then has
\[
\mathfrak U(A)=\int_Z^\oplus\mathfrak U(A(\zeta))\,d\mu(\zeta),
\qquad
\mathfrak V(A)=\int_Z^\oplus\mathfrak V(A(\zeta))\,d\mu(\zeta).
\]
% ===== ASSEMBLED TRANSLATION CHUNK 13: chunk_360_384.tex =====
% Source PDF physical page 360
If $J$ (respectively, $J(\zeta)$) is the involution of $H$
(respectively, $H(\zeta)$) defined by $A$ (respectively, $A(\zeta)$),
then
\[
 J=\int_Z^\oplus J(\zeta)\,d\mu(\zeta).
\]
The natural trace on $\mathfrak U(A)^+$ defined by $A$ is
$\int_Z^\oplus t(\zeta)\,d\mu(\zeta)$, where $t(\zeta)$ denotes the
natural trace on $\mathfrak U(A(\zeta))^+$ defined by $A(\zeta)$
\cite[pp.~188--190 and p.~199, Theorem~1]{appBref1}.

\ResultHeading{A.95}{A 95.}
Let $\mathcal E=((H(\zeta)),\Gamma)$ be a measurable field of Hilbert
spaces over $Z$, let
\[
 H=\int_Z^\oplus H(\zeta)\,d\mu(\zeta),
\]
let $\mathfrak D$ be the algebra of diagonalizable operators, and let $A$
be a completed Hilbert algebra having $H$ as its completed Hilbert space.
If $\mathfrak U(A)\subset\mathfrak D'$, there exists a measurable field
$\zeta\mapsto A(\zeta)$ of completed Hilbert algebras having the
$H(\zeta)$ as their completed Hilbert spaces, such that
\[
 A=\int_Z^\oplus A(\zeta)\,d\mu(\zeta)
\]
\cite[p.~194, Theorem~1]{appBref1}.

\ResultHeading{A.96}{A 96.}
Let $Z$ be a Borel space. A Borel field $\mathcal E$ of Hilbert spaces
over $Z$ is a family $(H(\zeta))_{\zeta\in Z}$ of Hilbert spaces together
with a set $\Gamma$ of vector fields satisfying conditions (i)--(iv) of
\XRef{A.69}{A 69}, with the word ``$\mu$-measurable'' replaced everywhere
by ``Borel'' \cite[pp.~143--144]{appBref1}. Such fields have properties
analogous to those of measurable fields, proved in the same way. In
particular, \XRef{A.71}{A 71} and \XRef{A.72}{A 72} may be repeated with
``measurable'' replaced everywhere by ``Borel.''

\ResultHeading{A.97}{A 97.}
Let $\mathcal E=((H(\zeta)),\Gamma)$ be a Borel field of Hilbert spaces
over $Z$, and let $\mu$ be a positive measure on $Z$. There exists one and
only one set $\Gamma'\supset\Gamma$ of vector fields such that
$((H(\zeta)),\Gamma')$ is a $\mu$-measurable field of Hilbert spaces;
$\Gamma'$ is the set of vector fields $x$ such that, for every
$y\in\Gamma$, the function
\[
 \zeta\longmapsto(x(\zeta)\mid y(\zeta))
\]
is $\mu$-measurable. The $\mu$-measurable field
$((H(\zeta)),\Gamma')$ is said to be derived from the Borel field
$((H(\zeta)),\Gamma)$ \cite[p.~144]{appBref1}.

\ResultHeading{A.98}{A 98.}
Let $Z$ be a Borel space, let $\zeta\mapsto H(\zeta)$ be a field of
Hilbert spaces over $Z$, and let
$\zeta\mapsto x_i(\zeta)$, $i=1,2,\ldots$, be a sequence of vector fields
with the following properties:

\begin{enumerate}[label=(\roman*)]
 \item for every $\zeta\in Z$, the sequence of the $x_i(\zeta)$ is total
 in $H(\zeta)$;
 \item the functions
 $\zeta\mapsto(x_i(\zeta)\mid x_j(\zeta))$ are Borel.
\end{enumerate}
% Source PDF physical page 361
Then there exists one and only one Borel-field structure on
$\zeta\mapsto H(\zeta)$ for which the $\zeta\mapsto x_i(\zeta)$ are
Borel vector fields \cite[p.~145, Proposition~4, and p.~146,
Remark]{appBref1}.

% Source PDF physical page 362
\BookAppendix{B}{Miscellaneous Results}

This Appendix is an unordered list of very varied results used in the
course of the text. The references are to the special bibliography at the
end of the Appendix.

\ResultHeading{B.1}{B 1.}
Let $A$ be a complex Banach algebra. For every $x\in A$,
$\lVert x^n\rVert^{1/n}$ has a limit $\rho(x)$ as $n\to+\infty$, called
the spectral radius of $x$. One has $\rho(x)\leq\lVert x\rVert$, and
$\rho(x)$ is the supremum of $\lvert z\rvert$ as $z$ ranges over
$\operatorname{Sp}_A x$ \cite[Theorems~1.4.1 and 1.6.4]{appBref11}.
Every maximal regular left ideal of $A$ is closed
\cite[Corollary~2.1.4]{appBref11}.

\ResultHeading{B.2}{B 2.}
Let $A$ be a unital Banach algebra, let $B$ be a closed subalgebra
containing $1$, and let $x\in B$. Then
$\operatorname{Sp}_B x\supset\operatorname{Sp}_A x$, and every boundary
point of $\operatorname{Sp}_B x$ relative to $\mathbb C$ belongs to
$\operatorname{Sp}_A x$ \cite[Theorem~1.6.12]{appBref11}.

\ResultHeading{B.3}{B 3.}
Let $A$ be a commutative Banach algebra. Every character of $A$ has norm
at most $1$. The set of characters of $A$, endowed with the weak
topology, is a locally compact space $S$ called the spectrum of $A$. For
every $x\in A$, the function $\chi\mapsto\chi(x)$ on $S$ is called the
Gelfand transform of $x$ and is denoted by $\mathcal F x$. The functions
$\mathcal F x$ ($x\in A$) are continuous, separate the points of $S$,
and, for every $\chi\in S$, there exists $x\in A$ such that
$(\mathcal F x)(\chi)\ne0$. If $A$ has an identity element, the set of
values of $\mathcal F x$ is $\operatorname{Sp}_A x$
\cite[p.~110, Theorem~3.1.6, Corollary~3.1.7, and
Theorem~3.1.20]{appBref11}.

\ResultHeading{B.4}{B 4.}
Let $A$ be a unital Banach algebra, let $x\in A$, and let $f$ be a
function holomorphic in a neighborhood of $\operatorname{Sp}_A x$. There
exists an element $f(x)\in A$ with the following properties:

\begin{enumerate}[label=(\roman*)]
 \item $\operatorname{Sp}_A f(x)=f(\operatorname{Sp}_A x)$;
 \item if $f(z)=\sum_{n=0}^{\infty}c_nz^n$ is an entire function, then
 $f(x)=\sum_{n=0}^{\infty}c_nx^n$, the series converging normally.
\end{enumerate}
See \cite[pp.~157--158 and the references cited there]{appBref11}. The
hypothesis made in \cite{appBref11} that $A$ is commutative is not really
used. Moreover, at the only place in the present book
(\XRef{1.3.9}{1.3.9}) where these results are applied, it would be easy
to reduce matters to the case of commutative algebras.

% Source PDF physical page 363
\ResultHeading{B.5}{B 5.}
Let $E$ be a real locally convex space, let $C$ be a closed convex cone
with vertex $0$, and let $x$ be a point of $E$ not belonging to $C$.
There exists a continuous linear form on $E$ which is nonnegative on $C$
and negative at $x$. [Indeed, there exist a continuous linear form $f$ on
$E$ and a real number $\alpha$ such that $f(y)\geq\alpha$ on $C$ and
$f(x)<\alpha$. We have $0=f(0)\geq\alpha$, and hence $f(x)<0$. If
$f(y)<0$ at a point of $C$, then $f(\lambda y)<\alpha$ for sufficiently
large $\lambda>0$, which is absurd.]

\ResultHeading{B.6}{B 6.}
Let $E$ be a separated locally convex space, let $P$ be a convex cone in
$E$ such that $P\cap(-P)=0$, let $M$ be a vector subspace of $E$, and let
$f$ be a linear form on $M$ which is nonnegative on $M\cap P$. Suppose
that $M$ contains an interior point of $P$. Then $f$ extends to a linear
form on $E$ which is nonnegative on $P$ \cite[p.~64,
Corollary]{appBref4}.

\ResultHeading{B.7}{B 7.}
Let $E$ be a separable Banach space. The unit ball of the dual of $E$ is
compact metrizable for the weak topology \cite[p.~112,
Proposition~2]{appBref5}.

\ResultHeading{B.8}{B 8.}
Let $E$ be a Banach space, let $E'$ be its dual, and let $K$ be a convex
subset of $E'$. Then $K$ is weakly closed if and only if its intersection
with every closed ball of $E'$ is weakly closed \cite[p.~74,
Theorem~5]{appBref5}.

\ResultHeading{B.9}{B 9.}
Let $H$ be a separable Hilbert space, and let $B$ be the unit ball of
$\mathcal L(H)$ endowed with the weak topology. Then $B$ is compact
metrizable \cite[p.~25, Proposition~6, and p.~65,
Corollary~3]{appBref5}.

\ResultHeading{B.10}{B 10.}
Let $H$ be a separable Hilbert space. The set of unitary operators on $H$
is a Polish space for the weak topology \cite[p.~280,
Lemma~4]{appBref7}.

\ResultHeading{B.11}{B 11.}
Let $E,F$ be two locally convex spaces. Suppose that $E$ is barrelled and
that $F$ is separated and quasi-complete. Then $\mathcal L(E,F)$, endowed
with the topology of pointwise convergence, is separated and
quasi-complete \cite[p.~31, Corollary~2]{appBref5}.

\ResultHeading{B.12}{B 12.}
Let $E,F$ be two topological vector spaces. Suppose that $F$ is separated
and quasi-complete. Then, in $\mathcal L(E,F)$ endowed with the topology
of pointwise convergence, every closed equicontinuous subset is a
complete uniform subspace \cite[p.~30, Theorem~4]{appBref5}.

\ResultHeading{B.13}{B 13.}
Let $E$ be a separated locally convex space, and let $C$ be a metrizable
compact convex subset of $E$. The set $A$ of extreme points of $C$ is a
$G_\delta$ subset of $C$ \cite[pp.~118--119]{appBref8}. Let $x\in C$.
There exists a positive measure $\mu$ of total mass $1$ on $C$,
concentrated on $A$ (that is, such that $C\setminus A$ is
$\mu$-negligible), and such that
\[
 x=\int_C y\,d\mu(y),
\]
in other words,
% Source PDF physical page 364
\[
 f(x)=\int_C f(y)\,d\mu(y)
\]
for every continuous linear form $f$ on $E$ (G.~Choquet's theorem;
\cite[Theorem]{appBref9}).

\ResultHeading{B.14}{B 14.}
Let $E$ be a separated locally convex space, let $C$ be a compact convex
subset of $E$, and let $A$ be the set of extreme points of $C$. Then $C$
is the closed convex hull of $A$ (the Krein--Milman theorem). If a subset
$A'$ of $C$ has closed convex hull equal to $C$, then $\overline{A'}$
contains $A$.

We shall prove that $A$ is a Baire space. (This unpublished result, as
well as the proof which follows, is due to G.~Choquet.) We may suppose
that $E$ is real. For every continuous linear form $f$ on $E$ and every
real number $\alpha$, let $U_{f,\alpha}$ (respectively,
$F_{f,\alpha}$) denote the set of $x\in C$ such that $f(x)<\alpha$
(respectively, $f(x)\leq\alpha$).

Let $x\in A$. We first show that the set of the $F_{f,\alpha}$ such that
$x\in U_{f,\alpha}$ is a fundamental system of neighborhoods of $x$. By
the Hahn--Banach theorem, the intersection of these $F_{f,\alpha}$ is
$\{x\}$; since the $F_{f,\alpha}$ are compact, it is enough to show that
the family of the $F_{f,\alpha}$ such that $x\in U_{f,\alpha}$ is
downward directed. Let $f_1,\alpha_1,f_2,\alpha_2$ be such that
\[
 x\in U_{f_1,\alpha_1}\cap U_{f_2,\alpha_2}.
\]
Let $C_1,C_2$ be the complements in $C$ of $U_{f_1,\alpha_1}$ and
$U_{f_2,\alpha_2}$. Since $C_1$ and $C_2$ are compact, their convex hull
$C_3$ is compact. It does not contain $x$, because $x$ is extreme. Hence
there exist a continuous linear form $f$ on $E$ and a real number
$\alpha$ such that $f(x)<\alpha$ and $f(y)>\alpha$ for $y\in C_3$.
Then $F_{f,\alpha}$ meets neither $C_1$ nor $C_2$; that is,
\[
 F_{f,\alpha}\subset
 U_{f_1,\alpha_1}\cap U_{f_2,\alpha_2},
 \qquad x\in U_{f,\alpha}.
\]

Let $(V_n)$ be a sequence of open dense subsets of $A$. We must prove
that $\bigcap_nV_n$ is dense in $A$, that is, that it meets every nonempty
open subset $V$ of $A$. Let $U_n,U$ be open subsets of $C$, with $U_n$
dense in $C$, such that $U_n\cap A=V_n$ and $U\cap A=V$. We may suppose
that the $U_n$ and the $V_n$ are decreasing. We may also suppose that
$U$ is of the form $U_{f_1,\alpha_1}$.

With this understood, we shall prove the existence of
$(f_1,\alpha_1),(f_2,\alpha_2),\ldots$ with the following properties:
\[
 F_{f_{n+1},\alpha_{n+1}}
 \subset U_{f_n,\alpha_n}\cap U_{n+1},
 \qquad
 U_{f_n,\alpha_n}\cap A\ne\varnothing
 \quad(n=1,2,\ldots).
\]

The initial pair $(f_1,\alpha_1)$ has already been chosen. Suppose that the
pairs $(f_i,\alpha_i)$ have been constructed for $1\leq i\leq n$. There is
$x\in U_{f_n,\alpha_n}\cap A\cap U_{n+1}$. Then
$U_{f_n,\alpha_n}\cap U_{n+1}$ is a neighborhood of $x$ in $C$, and so
there exists $(f_{n+1},\alpha_{n+1})$ such that
\[
 x\in U_{f_{n+1},\alpha_{n+1}}
 \subset F_{f_{n+1},\alpha_{n+1}}
 \subset U_{f_n,\alpha_n}\cap U_{n+1}.
\]
Since $x\in A$, one has
$U_{f_{n+1},\alpha_{n+1}}\cap A\ne\varnothing$, and the induction may
continue.

The sets $F_{f_n,\alpha_n}$ form a decreasing sequence of nonempty compact
sets, and hence their intersection $F$ is nonempty. We have
\[
 F\subset U_{f_1,\alpha_1},
 \qquad
 F\subset\bigcap_nU_n.
\]
Finally, $F$ is compact and convex, and its complement in $C$ is convex.
An easy lemma then proves that $F$ contains at least one extreme point
$y$ of $C$. (The set $F$ has an extreme point $x$. If $x$ is extreme in
$C$,
% Source PDF physical page 365
we are done. Otherwise, let $\delta$ be a line through $x$ such that $x$
is an interior point of the segment $C\cap\delta$; one then shows that
one of the endpoints of $F\cap\delta$ is an extreme point of $C$.) The
point $y$ belongs to $A\cap U_n=V_n$ for every $n$, and to
$A\cap U_{f_1,\alpha_1}=V$.

\ResultHeading{B.15}{B 15.}
Every $G_\delta$ subset of a Polish space is a Polish space
\cite[p.~123, Theorem~1]{appBref2}.

\ResultHeading{B.16}{B 16.}
Let $E$ be a locally compact space, let $G$ be a group of homeomorphisms
of $E$ onto $E$, and let $\mathcal C(E,E)$ be the set of continuous maps
of $E$ into $E$, endowed with the topology of compact convergence.
Suppose that $G$ is relatively compact in $\mathcal C(E,E)$. Then the
closure of $G$ in $\mathcal C(E,E)$ is a compact group of homeomorphisms
of $E$ onto $E$ \cite[p.~52, Corollary]{appBref3}.

\ResultHeading{B.17}{B 17.}
Let $E$ be a Polish space, and let $R$ be an equivalence relation on $E$.
Suppose that the saturation in $E$ of every open subset is Borel and that
the equivalence classes are closed. Then there exists a Borel subset of
$E$ which meets each equivalence class in one and only one point
\cite[p.~279, Lemma~2]{appBref7}.

\ResultHeading{B.18}{B 18.}
Let $E$ be a Baire space, and let $f\colon E\to\mathbb R$ be a lower
semicontinuous function. Then $f$ has a point of continuity. Indeed, on
replacing $f$ by $f(1+\lvert f\rvert)^{-1}$, we may suppose that $f$ is
bounded. For every $x\in E$, let $\omega(x)$ be the oscillation of $f$ at
$x$. Then $\omega$ is upper semicontinuous, and hence the set $E_n$ of
points at which $\omega(x)\geq1/n$ is closed. Suppose that $E_n$ contains
a nonempty open subset $U$. Put
\[
 \alpha=\sup_{x\in U}f(x).
\]
There exists $x_0\in U$ such that
$f(x_0)>\alpha-1/(2n)$. Throughout some neighborhood of $x_0$ one has
\[
 \alpha-\frac1{2n}<f(x)\leq\alpha,
\]
and hence $\omega(x_0)\leq1/(2n)$, a contradiction. Thus $E\setminus E_n$
is open and dense, so $\bigcap_n(E\setminus E_n)$ is nonempty. If
$x\in\bigcap_n(E\setminus E_n)$, then $\omega(x)=0$, and therefore $f$ is
continuous at $x$.

\ResultHeading{B.19}{B 19.}
A Borel space is a set $E$ together with a collection $\mathcal B$ of
subsets of $E$ having the following properties:
$E\in\mathcal B$, $\varnothing\in\mathcal B$, and $\mathcal B$ is stable
under countable unions, countable intersections, and passage to the
complement. The elements of $\mathcal B$ are called the Borel subsets of
$E$.

Let $E$ be a set and let $\mathcal A$ be a collection of subsets of $E$.
Among the collections $\mathcal B$ of subsets of $E$ such that
$\mathcal B\supset\mathcal A$ and $(E,\mathcal B)$ is a Borel space,
there is one, $\mathcal B_0$, smaller than all the others. The Borel
structure of $(E,\mathcal B_0)$ is said to be generated by $\mathcal A$.

% Source PDF physical page 366
Let $(E,\mathcal B_1)$ and $(E,\mathcal B_2)$ be two Borel spaces. If
$\mathcal B_2\supset\mathcal B_1$, the Borel structure of
$(E,\mathcal B_2)$ is said to be finer than that of
$(E,\mathcal B_1)$.

Let $(E_i)_{i\in I}$ be a family of Borel spaces, let $E$ be a set, and
let $f_i\colon E\to E_i$ be maps. The Borel structure on $E$ defined by
the $f_i$ is the structure generated by the sets $f_i^{-1}(A)$, where
$i$ ranges over $I$ and $A$ ranges over the Borel subsets of $E_i$.

Let $E,F$ be two Borel spaces. A map $f\colon E\to F$ is said to be
Borel if the inverse image under $f$ of every Borel subset of $F$ is a
Borel subset of $E$.

Let $E$ be a Borel space and let $E'$ be a subset of $E$. The
intersections with $E'$ of the Borel subsets of $E$ define a Borel
structure on $E'$. The Borel space $E'$ is called a Borel subspace of
$E$.

Let $E$ be a Borel space and let $R$ be an equivalence relation on $E$.
The subsets of $E/R$ whose inverse images in $E$ are Borel define a Borel
structure on $E/R$. The Borel space $E/R$ is called the quotient Borel
space of $E$ by $R$.

Let $(E_\alpha)$ be a family of Borel spaces, and let $E$ be the disjoint
union of the $E_\alpha$. The subsets of $E$ whose intersections with the
$E_\alpha$ are Borel define a Borel structure on $E$. The Borel space
$E$ is called the sum Borel space of the $E_\alpha$.

Let $E$ be a topological space. The subsets of $E$ which are Borel for
the topology define a Borel structure on $E$, called the structure
underlying the topology. If $E'\subset E$, the Borel structure induced on
$E'$ by that of $E$ is the structure underlying the topology induced on
$E'$ by the topology of $E$.

For all the preceding material, see \cite[pp.~136--137]{appBref10}.

\ResultHeading{B.20}{B 20.}
A Borel space $E$ is said to be standard if its Borel structure underlies
a topology of a Polish space. If $E$ is countable and standard, every
subset of $E$ is Borel. If $E$ is uncountable and standard, $E$ is
isomorphic to the Borel space $[0,1]$ endowed with the Borel structure
underlying its usual topology. Every Borel subset of a standard space is
standard. The sum of a sequence of standard spaces is standard
\cite[p.~138]{appBref10}.

\ResultHeading{B.21}{B 21.}
Let $E_1,E_2$ be two Borel spaces, and let $f\colon E_1\to E_2$ be an
injective Borel map. Suppose that:

\begin{enumerate}[label=(\roman*)]
 \item $E_1$ is standard;
 \item there exists a sequence of Borel subsets of $E_2$ separating the
 points of $E_2$ and generating the Borel structure of $E_2$ (as is the
 case if $E_2$ is standard).
\end{enumerate}
Then $f(E_1)$ is Borel in $E_2$, and $f$ is an isomorphism from $E_1$
onto $f(E_1)$ \cite[p.~139, Theorem~3.2]{appBref10}.

\ResultHeading{B.22}{B 22.}
Let $E_1,E_2$ be two Borel spaces, and let $f\colon E_1\to E_2$ be a
bijective Borel map. Suppose that:

\begin{enumerate}[label=(\roman*)]
 \item $E_2$ is standard;
 \item $E_1$ is a quotient
% Source PDF physical page 367
of a standard Borel space.
\end{enumerate}
Then $f$ is an isomorphism from $E_1$ onto $E_2$
\cite[p.~140, Theorem~4.2, and p.~141, Corollary]{appBref10}.

\ResultHeading{B.23}{B 23.}
Let $H$ be a Hilbert space, and let $S$ be a linear operator defined on a
vector subspace $D_S$ of $H$ dense in $H$, with values in $H$, such that
$(Sf\mid f)\geq0$ for every $f\in D_S$. For $f,g\in D_S$, put
\[
 (f\mid g)_S=(Sf\mid g)+(f\mid g).
\]
Every Cauchy sequence for $(\mathord\cdot\mid\mathord\cdot)_S$ is a
Cauchy sequence in $H$, and two Cauchy sequences equivalent for
$(\mathord\cdot\mid\mathord\cdot)_S$ are equivalent in $H$. Thus the
completion $K$ of $D_S$ for $(\mathord\cdot\mid\mathord\cdot)_S$ is
identified with a vector subspace of $H$. Let $D$ be the intersection of
$K$ with the domain of definition of $S^{*}$, and let $S'$ be the
restriction of $S^{*}$ to $D$. Then $S'$ is self-adjoint and
nonnegative. It is called the Friedrichs extension of $S$. If $U$ is a
unitary operator on $H$ such that $USU^{-1}=S$, then
$US'U^{-1}=S'$ \cite[pp.~35--36]{appBref12}.

\ResultHeading{B.24}{B 24.}
Let $G$ be a locally compact group and let $H$ be a closed subgroup such
that the homogeneous space $G/H$ is compact. If $G$ is generated by a
compact neighborhood of $e$, the same is true of $H$
\cite[p.~78, Lemma~3]{appBref6}.

\ResultHeading{B.25}{B 25.}
Let $G$ be a commutative locally compact group generated by a compact
neighborhood of $e$. Then $G$ is isomorphic to a group
$K\times\mathbb R^n\times\mathbb Z^p$, where $K$ is a compact
commutative group \cite[p.~110]{appBref13}.

\ResultHeading{B.26}{B 26.}
Let $A$ be an algebra and let $a,b\in A$. Then
\[
 \operatorname{Sp}'(ab)=\operatorname{Sp}'(ba)
\]
(this is due to N.~Jacobson). Indeed, let $\widetilde A$ be the algebra
obtained from $A$ by adjoining an identity element $e$. It is enough to
prove that, if $\lambda\ne0$ is such that $ab-\lambda e$ has an inverse
$u$ in $\widetilde A$, then $ba-\lambda e$ is invertible in
$\widetilde A$. But
\begin{align*}
 (ba-\lambda e)(bua-e)
 &=b(abu)a-ba-\lambda bua+\lambda e\\
 &=b(\lambda u+e)a-ba-\lambda bua+\lambda e
  =\lambda e,
\end{align*}
and similarly $(bua-e)(ba-\lambda e)=\lambda e$. Since $\lambda\ne0$,
this indeed shows that $ba-\lambda e$ is invertible.

\ResultHeading{B.27}{B 27.}
Let $A$ be an algebra over a field $k$ of characteristic zero. Let
$\pi,\pi'$ be irreducible representations of $A$ on finite-dimensional
vector spaces over $k$. If
\[
 \operatorname{Tr}\pi(x)=\operatorname{Tr}\pi'(x)
 \qquad(x\in A),
\]
then $\pi$ and $\pi'$ are equivalent \cite[p.~136,
Proposition~3]{appBref1}.

\ResultHeading{B.28}{B 28.}
Let $T$ be a locally compact space, and let $\mu$ be a bounded complex
measure on $T$. If $\lVert\mu\rVert=\mu(1)$, then $\mu\geq0$. Indeed,
put $\nu=\lvert\mu\rvert$. We have
$\lVert\mu\rVert=\lVert\nu\rVert=\nu(1)$ and $\mu=f\mathbin{\cdot}\nu$,
where $f$ is a complex $\nu$-measurable function such that
$\lvert f(t)\rvert=1$ almost everywhere. Let $\alpha<1$, and let
$T_\alpha$ be the set of $t\in T$ such that
$\alpha\leq\operatorname{Re}f(t)\leq1$. Then
% Source PDF physical page 368
\begin{align*}
 \lVert\nu\rVert=\mu(1)
 &=\operatorname{Re}\int_{T\setminus T_\alpha}f(t)\,d\nu(t)
   +\operatorname{Re}\int_{T_\alpha}f(t)\,d\nu(t)\\
 &\leq\alpha\int_{T\setminus T_\alpha}d\nu(t)
   +\int_{T_\alpha}d\nu(t)\\
 &=\int_Td\nu(t)
   -(1-\alpha)\int_{T\setminus T_\alpha}d\nu(t)\\
 &=\lVert\nu\rVert
   -(1-\alpha)\int_{T\setminus T_\alpha}d\nu(t).
\end{align*}
Thus $T\setminus T_\alpha$ is $\nu$-negligible. Since this is true for
every $\alpha<1$, one has $\operatorname{Re}f(t)\geq1$ almost
everywhere, hence $f(t)=1$ almost everywhere and $\mu=\nu\geq0$.

\ResultHeading{B.29}{B 29.}
Let $A$ be a normed algebra. An approximate identity of $A$ is a family
$(u_i)_{i\in I}$ of elements of $A$, indexed by an upward directed set,
with the following properties:

\begin{enumerate}[label=\alph*.]
 \item $\lVert u_i\rVert\leq1$ for every $i$;
 \item $\lVert u_ix-x\rVert\to0$ and $\lVert xu_i-x\rVert\to0$ for every
 $x\in A$.
\end{enumerate}

\ResultHeading{B.30}{B 30.}
Let $X$ be a Borel space, and let $\mathcal B$ be the set of Borel
subsets of $X$. A positive measure on $X$ is a map
$\mu\colon\mathcal B\to[0,+\infty]$ such that:

\begin{enumerate}[label=\arabic*\textsuperscript{o}]
 \item if $X_1,X_2,\ldots$ are pairwise disjoint elements of $\mathcal B$,
 then
 \[
  \mu(X_1\cup X_2\cup\cdots)
  =\mu(X_1)+\mu(X_2)+\cdots;
 \]
 \item $X$ is the union of a sequence of Borel subsets
 $Y_1,Y_2,\ldots$ such that $\mu(Y_i)<+\infty$ for every $i$.
\end{enumerate}
A positive measure $\mu$ on $X$ is said to be standard if there exists a
Borel subset $N$ of $X$ such that $\mu(N)=0$ and $X\setminus N$ is a
standard Borel space.

\ResultHeading{B.31}{B 31.}
A topological space is said to be separable if its topology has a
countable base. When a space is metrizable, saying that it is separable
amounts to saying that it has a countable dense subset. A compact
metrizable space is separable.

\ResultHeading{B.32}{B 32.}
In a topological space, a set is called a $G_\delta$ if it is the
intersection of a countable family of open sets.

\ResultHeading{B.33}{B 33.}
A character of a commutative algebra is a nonzero homomorphism from that
algebra into the base field.

\ResultHeading{B.34}{B 34.}
Let $V$ be a finite-dimensional complex vector space, let $G$ be a
compact group, and let $\pi$ be a continuous homomorphism from $G$ into
the linear group of $V$. There exists on $V$ a scalar product invariant
under $\pi(G)$ which makes $V$ a Hilbert space
\cite[p.~70, Proposition~1]{appBref6}.

% Source PDF physical page 369
\BookBackMatter{appendix-b-special-bibliography}%
{Special Bibliography for Appendix B}

\begin{DixmierBibliography}{99}
\bibitem[1]{appBref1}
N.~Bourbaki, \emph{Algèbre}, chap.~VIII (Act. Sc. Ind., no.~1261,
Hermann, Paris, 1958).

\bibitem[2]{appBref2}
N.~Bourbaki, \emph{Topologie générale}, chap.~IX, 2nd ed. (Act. Sc.
Ind., no.~1045, Hermann, Paris, 1958).

\bibitem[3]{appBref3}
N.~Bourbaki, \emph{Topologie générale}, chap.~X, 2nd ed. (Act. Sc.
Ind., no.~1084, Hermann, Paris, 1961).

\bibitem[4]{appBref4}
N.~Bourbaki, \emph{Espaces vectoriels topologiques}, chap.~I--II,
2nd ed. (Act. Sc. Ind., no.~1189, Hermann, Paris, 1960).

\bibitem[5]{appBref5}
N.~Bourbaki, \emph{Espaces vectoriels topologiques}, chap.~III--IV--V
(Act. Sc. Ind., no.~1229, Hermann, Paris, 1955).

\bibitem[6]{appBref6}
N.~Bourbaki, \emph{Intégration}, chap.~VII--VIII (Act. Sc. Ind.,
no.~1306, Hermann, Paris, 1963).

\bibitem[7]{appBref7}
J.~Dixmier, Dual et quasi-dual d'une algèbre de Banach involutive
(Trans. Amer. Math. Soc., vol.~104, 1962, pp.~278--283).

\bibitem[8]{appBref8}
R.~Godement, Sur la théorie des représentations unitaires (Ann. of Math.,
vol.~53, 1951, pp.~68--124).

\bibitem[9]{appBref9}
M.~Hervé, Sur les représentations intégrales à l'aide des points extrémaux
dans un ensemble compact convexe métrisable (C. R. Acad. Sc., vol.~253,
1961, pp.~366--368).

\bibitem[10]{appBref10}
G.~W.~Mackey, Borel structure in groups and their duals (Trans. Amer.
Math. Soc., vol.~85, 1957, pp.~134--165).

\bibitem[11]{appBref11}
C.~E.~Rickart, \emph{General Theory of Banach Algebras} (The University
Series in Higher Mathematics, D. Van Nostrand, Princeton, 1960).

\bibitem[12]{appBref12}
B.~Sz.-Nagy, \emph{Spektraldarstellung linearer Transformationen des
Hilbertschen Raumes} (Erg. der Math., Springer-Verlag, Berlin, 1942).

\bibitem[13]{appBref13}
A.~Weil, \emph{L'intégration dans les groupes topologiques et ses
applications}, 2nd ed. (Act. Sc. Ind., no.~1145, Hermann, Paris, 1953).
\end{DixmierBibliography}

% Source PDF physical page 370
\BookBackMatter{bibliography}{Bibliography}

The first part of this bibliography extends approximately through 1963. The articles are arranged alphabetically by author. The second part extends approximately from 1964 to 1968. The articles are arranged by year. For some time it has become increasingly artificial to separate work devoted to von Neumann algebras from work devoted to $C^{*}$-algebras. The division made between the present book and the second edition of my book on von Neumann algebras (Cahiers Scientifiques, fasc.~XXV) is therefore rather arbitrary.

Some articles in theoretical physics related to $C^{*}$-algebras have been noted, but no attempt has been made to provide a complete list.

\begin{DixmierBibliography}{999}

\bibitem[A v N]{refAvN}
J.~Dixmier, \emph{Les alg\`{e}bres d'op\'{e}rateurs dans l'espace hilbertien (Alg\`{e}bres de von Neumann)}, 2nd ed., Cahiers Scientifiques, fasc.~XXV, Gauthier-Villars, Paris, 1969.

\bibitem[1]{ref1}
R. Arens, On a theorem of Gelfand and Neumark (Proc. Nat. Acad. Sc. U.S. A., vol.~32, 1946, pp.~237--239).

\bibitem[2]{ref2}
A. Badrikian, Structure de certains groupes localement compacts (Cahiers Rhodaniens, vol.~5, 1953, pp.~27--51).

\bibitem[3]{ref3}
H. F. Bohnenblust and S. Karlin, Geometrical properties of the unit sphere of Banach algebras (Ann. Math., vol.~62, 1955, pp.~217--229).

\bibitem[4]{ref4}
W. F. Darsow, Positive definite functions and states (Ann. Math., vol.~60, 1954, pp.~447--453).

\bibitem[5]{ref5}
J. Dixmier, Sur les repr\'{e}sentations unitaires des groupes de Lie alg\'{e}briques (Ann. Inst. Fourier, vol.~7, 1957, pp.~315--328).

\bibitem[6]{ref6}
J. Dixmier, Sur les repr\'{e}sentations unitaires des groupes de Lie nilpotents. IV (Can. J. Math., vol.~11, 1959, pp.~321--344).

\bibitem[7]{ref7}
J. Dixmier, Sur les repr\'{e}sentations unitaires des groupes de Lie nilpotents. V (Bull. Soc. Math. Fr., vol.~87, 1959, pp.~65--79).

\bibitem[8]{ref8}
J. Dixmier, Sur les repr\'{e}sentations unitaires des groupes de Lie nilpotents. VI (Can. J. Math., vol.~12, 1960, pp.~324--352).

\bibitem[9]{ref9}
J. Dixmier, Repr\'{e}sentations int\'{e}grables du groupe de De Sitter (Bull. Soc. Math. Fr., vol.~89, 1961, pp.~9--41).

\bibitem[10]{ref10}
J. Dixmier, Sur les repr\'{e}sentations unitaires des groupes de Lie r\'{e}solubles (Math. J. Okayama Univ., vol.~11, 1962, pp.~1--18).

\bibitem[11]{ref11}
J. Dixmier, Sur les $C^{*}$-alg\`{e}bres (Bull. Soc. Math. Fr., vol.~88, 1960, pp.~95--112).

\bibitem[12]{ref12}
J. Dixmier, Sur les structures bor\'{e}liennes du spectre d'une $C^{*}$-alg\`{e}bre (Publ. Inst. Hautes \'{E}tudes Sc., no.~6, 1960, pp.~297--303).

\bibitem[13]{ref13}
J. Dixmier, Points isol\'{e}s dans le dual d'un groupe localement compact (Bull. Sc. Math., vol.~85, 1961, pp.~91--96).

\bibitem[14]{ref14}
J. Dixmier, Points s\'{e}par\'{e}s dans le spectre d'une $C^{*}$-alg\`{e}bre (Acta Sc. Math., vol.~22, 1961, pp.~115--128).

\bibitem[15]{ref15}
J. Dixmier, Dual et quasi-dual d'une alg\`{e}bre de Banach involutive (Trans. Amer. Math. Soc., vol.~104, 1962, pp.~278--283).

\bibitem[16]{ref16}
J. Dixmier, Traces sur les $C^{*}$-alg\`{e}bres (Ann. Inst. Fourier, vol.~13, 1963, pp.~219--262).

\bibitem[17]{ref17}
J. Dixmier, Quasi-dual d'un id\'{e}al dans une $C^{*}$-alg\`{e}bre (Bull. Sc. Math., vol.~87, 1963, pp.~7--11).

\bibitem[18]{ref18}
J. Dixmier, Champs continus d'espaces hilbertiens et de $C^{*}$-alg\`{e}bres. II (J. Math. pures et appl., vol.~42, 1963, pp.~1--20).

\bibitem[19]{ref19}
J. Dixmier and A. Douady, Champs continus d'espaces hilbertiens et de $C^{*}$-alg\`{e}bres (Bull. Soc. Math. Fr., vol.~91, 1963, pp.~227--284).

\bibitem[20]{ref20}
E. G. Effros, A decomposition theory for representations of $C^{*}$-algebras (Trans. Amer. Math. Soc., vol.~107, 1963, pp.~83--106).

\bibitem[21]{ref21}
E. G. Effros, Order ideals in a $C^{*}$-algebra and its dual (Duke Math. J., vol.~30, 1963, pp.~391--412).

\bibitem[22]{ref22}
L. Ehrenpreis and F. I. Mautner, Some properties of the Fourier transform on semi-simple Lie groups. I (Ann. Math., vol.~61, 1955, pp.~406--439).

% Source PDF physical page 371

\bibitem[23]{ref23}
L. Ehrenpreis and F. I. Mautner, Some properties of the Fourier transform on semi-simple Lie groups. II (Trans. Amer. Math. Soc., vol.~84, 1957, pp.~1--55).

\bibitem[24]{ref24}
L. Ehrenpreis and F. I. Mautner, Some properties of the Fourier transform on semi-simple Lie groups. III (Trans. Amer. Math. Soc., vol.~90, 1959, pp.~431--484).

\bibitem[25]{ref25}
J. Ernest, A decomposition theory for unitary representations of locally compact groups (Trans. Amer. Math. Soc., vol.~104, 1962, pp.~252--277).

\bibitem[26]{ref26}
J. M. G. Fell, Representations of weakly closed algebras (Math. Ann., vol.~133, 1957, pp.~118--126).

\bibitem[27]{ref27}
J. M. G. Fell, The dual spaces of $C^{*}$-algebras (Trans. Amer. Math. Soc., vol.~94, 1960, pp.~365--403).

\bibitem[28]{ref28}
J. M. G. Fell, $C^{*}$-algebras with smooth dual (Illinois J. Math., vol.~4, 1960, pp.~221--230).

\bibitem[29]{ref29}
J. M. G. Fell, The structure of algebras of operator fields (Acta Math., vol.~106, 1961, pp.~233--280).

\bibitem[30]{ref30}
J. M. G. Fell, A new proof that nilpotent groups are CCR (Proc. Amer. Math. Soc., vol.~13, 1962, pp.~93--99).

\bibitem[31]{ref31}
J. M. G. Fell, Weak containment and induced representations of groups (Can. J. Math., vol.~14, 1962, pp.~237--268).

\bibitem[32]{ref32}
J. M. G. Fell, Weak containment and induced representations of groups. II, \`{a} para\^{\i}tre.

\bibitem[33]{ref33}
J. M. G. Fell, Weak containment and Kronecker products of group representations, \`{a} para\^{\i}tre.

\bibitem[34]{ref34}
J. M. G. Fell and J. Feldman, Separable representations of rings of operators (Ann. Math., vol.~65, 1957, pp.~241--249).

\bibitem[35]{ref35}
J. M. G. Fell and E. Thoma, Einige Bemerkungen \"{u}ber vollsymmetrisch Banachsche Algebren (Arch. Math., vol.~12, 1961, pp.~69--70).

\bibitem[36]{ref36}
M. Fukamiya, On $B^{*}$-algebras (Proc. Jap. Acad., vol.~27, 1951, pp.~321--327).

\bibitem[37]{ref37}
M. Fukamiya, On a theorem of Gelfand and Neumark and the $B^{*}$-algebra (Kumamoto J. of Sc., vol.~1, 1952, pp.~17--22).

\bibitem[38]{ref38}
M. Fukamiya, M. Misonou, and Z. Takeda, On order and commutativity of $B^{*}$-algebras (T\^{o}hoku Math. J., vol.~6, 1954, pp.~89--93).

\bibitem[39]{ref39}
L. G\aa{}rding and A. Wightman, Representations of the anticommutation relations (Proc. Nat. Acad. Sc. U.S. A., vol.~40, 1954, pp.~617--621).

\bibitem[40]{ref40}
L. G\aa{}rding and A. Wightman, Representations of the commutation relations (Proc. Nat. Acad. Sc. U.S. A., vol.~40, 1954, pp.~622--626).

\bibitem[41]{ref41}
I. M. Gelfand and M. A. Naimark, On the imbedding of normed rings into the ring of operators in Hilbert space (Mat. Sb., vol.~12, 1943, pp.~197--213).

\bibitem[42]{ref42}
I. M. Gelfand and M. A. Naimark, Anneaux norm\'{e}s \`{a} involution et leurs repr\'{e}sentations (Izv. Akad. Nauk S.S.S.R., vol.~12, 1948, pp.~445--480) (in Russian).

\bibitem[43]{ref43}
I. M. Gelfand and D. Ra\"{\i}kov, Repr\'{e}sentations unitaires irr\'{e}ductibles des groupes localement compacts (Mat. Sb., vol.~13, 1943, pp.~301--316) (in Russian).

\bibitem[44]{ref44}
J. Glimm, On a certain class of operator algebras (Trans. Amer. Math. Soc., vol.~95, 1960, pp.~318--340).

\bibitem[45]{ref45}
J. Glimm, A Stone-Weierstrass theorem for $C^{*}$-algebras (Ann. Math., vol.~72, 1960, pp.~216--244).

\bibitem[46]{ref46}
J. Glimm, Type I $C^{*}$-algebras (Ann. Math., vol.~73, 1961, pp.~572--612).

\bibitem[47]{ref47}
J. Glimm, Families of induced representations (Pacific J. Math., vol.~12, 1962, pp.~885--911).

\bibitem[48]{ref48}
J. Glimm, What are $C^{*}$-algebras (Lecture at Columbia University, December 1961).

\bibitem[49]{ref49}
J. Glimm and R. V. Kadison, Unitary operators in $C^{*}$-algebras (Pacific J. Math., vol.~10, 1960, pp.~547--556).

% Source PDF physical page 372

\bibitem[50]{ref50}
R. Godement, Les fonctions de type positif et la th\'{e}orie des groupes (Trans. Amer. Math. Soc., vol.~63, 1948, pp.~1--84).

\bibitem[51]{ref51}
R. Godement, Sur les relations d'orthogonalit\'{e} de V. Bargmann (C. R. Acad. Sc., vol.~225, 1947, pp.~521--523 et 657--659).

\bibitem[52]{ref52}
R. Godement, Sur la th\'{e}orie des repr\'{e}sentations unitaires (Ann. Math., vol.~53, 1951, pp.~68--124).

\bibitem[53]{ref53}
R. Godement, M\'{e}moire sur la th\'{e}orie des caract\`{e}res dans les groupes localement compacts unimodulaires (J. Math. pures et appl., vol.~30, 1951, pp.~1--110).

\bibitem[54]{ref54}
R. Godement, Th\'{e}orie des caract\`{e}res. II. D\'{e}finition et propri\'{e}t\'{e}s g\'{e}n\'{e}rales des caract\`{e}res (Ann. Math., vol.~59, 1954, pp.~63--85).

\bibitem[55]{ref55}
A. Grothendieck, Un r\'{e}sultat sur le dual d'une $C^{*}$-alg\`{e}bre (J. Math. pures et appl., vol.~36, 1957, pp.~97--108).

\bibitem[56]{ref56}
A. Guichardet, Sur un probl\`{e}me pos\'{e} par G. W. Mackey (C. R. Acad. Sc., vol.~250, 1960, pp.~962--963).

\bibitem[57]{ref57}
A. Guichardet, Caract\'{e}res des alg\`{e}bres de Banach involutives (Ann. Inst. Fourier, vol.~13, 1962, pp.~1--81).

\bibitem[58]{ref58}
A. Gurevitch, Unitary representations in Hilbert space of a compact topological group (Mat. Sb., vol.~13, 1943, pp.~79--86).

\bibitem[59]{ref59}
Harish-Chandra, Representations of semi-simple Lie groups. III (Trans. Amer. Math. Soc., vol.~76, 1954, pp.~234--253).

\bibitem[60]{ref60}
Harish-Chandra, Representations of semi-simple Lie groups. VI (Amer. J. Math., vol.~78, 1956, pp.~564--628).

\bibitem[61]{ref61}
Harish-Chandra, Invariant eigendistributions on semi-simple Lie groups (Bull. Amer. Math. Soc., vol.~69, 1963, pp.~117--123).

\bibitem[62]{ref62}
R. Hirschfeld, Sur l'analyse harmonique dans les groupes localement compacts (C. R. Acad. Sc., vol.~246, 1958, pp.~1138--1140).

\bibitem[63]{ref63}
S. Ito, Positive definite functions on homogeneous spaces (Proc. Jap. Acad., vol.~26, 1950, pp.~17--28).

\bibitem[64]{ref64}
N. Jacobson, A topology for the set of primitive ideals in an arbitrary ring (Proc. Nat. Acad. Sc. U.S. A., vol.~31, 1945, pp.~333--338).

\bibitem[65]{ref65}
R. V. Kadison, Isometries of operator algebras (Ann. Math., vol.~54, 1951, pp.~325--338).

\bibitem[66]{ref66}
R. V. Kadison, A generalized Schwarz inequality and algebraic invariants for operator algebras (Ann. Math., vol.~56, 1952, pp.~494--503).

\bibitem[67]{ref67}
R. V. Kadison, On the orthogonalization of operator representations (Amer. J. Math., vol.~77, 1955, pp.~600--621).

\bibitem[68]{ref68}
R. V. Kadison, Operator algebras with a faithful weakly-closed representation (Ann. Math., vol.~64, 1956, pp.~175--181).

\bibitem[69]{ref69}
R. V. Kadison, Irreducible operator algebras (Proc. Nat. Acad. Sc. U.S. A., vol.~43, 1957, pp.~273--276).

\bibitem[70]{ref70}
R. V. Kadison, Unitary invariants for representations of operator algebras (Ann. Math., vol.~66, 1957, pp.~304--379).

\bibitem[71]{ref71}
R. V. Kadison, States and representations (Trans. Amer. Math. Soc., vol.~103, 1962, pp.~304--319).

\bibitem[72]{ref72}
R. V. Kadison and I. M. Singer, Some remarks on representations of connected groups (Proc. Nat. Acad. Sc. U.S. A., vol.~38, 1952, pp.~419--423).

\bibitem[73]{ref73}
R. V. Kadison and I. M. Singer, Extensions of pure states (Amer. J. Math., vol.~81, 1959, pp.~383--400).

\bibitem[74]{ref74}
I. Kaplansky, Normed algebras (Duke Math. J., vol.~16, 1949, pp.~399--418).

\bibitem[75]{ref75}
I. Kaplansky, Groups with representations of bounded degree (Can. J. Math., vol.~1, 1949, pp.~105--112).

\bibitem[76]{ref76}
I. Kaplansky, Group algebras in the large (T\^{o}hoku Math. J., vol.~3, 1951, pp.~249--256).

% Source PDF physical page 373

\bibitem[77]{ref77}
I. Kaplansky, The structure of certain operator algebras (Trans. Amer. Math. Soc., vol.~70, 1951, pp.~219--255).

\bibitem[78]{ref78}
I. Kaplansky, Representations of separable algebras (Duke Math. J., vol.~19, 1952, pp.~219--222).

\bibitem[79]{ref79}
I. Kaplansky, Functional analysis, Some aspects of analysis and probability, pp.~1--34 (John Wiley and Sons, New York, 1958).

\bibitem[80]{ref80}
J. L. Kelley and R. L. Vaught, The positive cone in Banach algebras (Trans. Amer. Math. Soc., vol.~74, 1953, pp.~44--55).

\bibitem[81]{ref81}
A. A. Kirillov, Repr\'{e}sentations unitaires des groupes de Lie nilpotents (Uspekhi Mat. Nauk, vol.~17, 1962, pp.~57--110) (in Russian).

\bibitem[82]{ref82}
P. Koosis, An irreducible unitary representation of a compact group is finite dimensional (Proc. Amer. Math. Soc., vol.~8, 1957, pp.~712--715).

\bibitem[83]{ref83}
R. A. Kunze, $L^{p}$-Fourier transforms on locally compact unimodular groups (Trans. Amer. Math. Soc., vol.~89, 1958, pp.~519--540).

\bibitem[84]{ref84}
R. A. Kunze and E. M. Stein, Uniformly bounded representations and harmonic analysis of the $2\times2$ real unimodular group (Amer. J. Math., vol.~82, 1960, pp.~1--62).

\bibitem[85]{ref85}
M. Kuranishi, On non connected maximally almost periodic groups (T\^{o}hoku Math. J., vol.~2, 1950, pp.~40--46).

\bibitem[86]{ref86}
L. H. Loomis, An introduction to abstract harmonic analysis (The University Series in higher mathematics, D. Van Nostrand, New York, 1953).

\bibitem[87]{ref87}
G. W. Mackey, Induced representations of groups (Amer. J. Math., vol.~73, 1951, pp.~576--592).

\bibitem[88]{ref88}
G. W. Mackey, Induced representations of locally compact groups. I (Ann. Math., vol.~55, 1952, pp.~101--139).

\bibitem[89]{ref89}
G. W. Mackey, Induced representations of locally compact groups. II. The Frobenius reciprocity theorem (Ann. Math., vol.~58, 1953, pp.~193--221).

\bibitem[90]{ref90}
G. W. Mackey, The theory of group representations (Notes mim\'{e}ographi\'{e}es, Universit\'{e} de Chicago, 1955).

\bibitem[91]{ref91}
G. W. Mackey, Borel structure in groups and their duals (Trans. Amer. Math. Soc., vol.~85, 1957, pp.~134--165).

\bibitem[92]{ref92}
G. W. Mackey, Unitary representations of group extensions (Acta Math., vol.~99, 1958, pp.~265--311).

\bibitem[93]{ref93}
G. W. Mackey, Induced representations and normal subgroups (Proc. of the International Symposium on linear spaces, Jerusalem, 1960).

\bibitem[94]{ref94}
S. Matsushima, Sur le th\'{e}or\`{e}me de Plancherel (Proc. Jap. Acad., vol.~30, 1954, pp.~557--561).

\bibitem[95]{ref95}
F. I. Mautner, Unitary representations of locally compact groups. I (Ann. Math., vol.~51, 1950, pp.~1--25).

\bibitem[96]{ref96}
F. I. Mautner, Unitary representations of locally compact groups. II (Ann. Math., vol.~52, 1950, pp.~528--556).

\bibitem[97]{ref97}
F. I. Mautner, Infinite-dimensional representations of certain groups (Proc. Amer. Math. Soc., vol.~1, 1950, pp.~582--584).

\bibitem[98]{ref98}
F. I. Mautner, The structure of the regular representation of certain discrete groups (Duke Math. J., vol.~17, 1950, pp.~437--441).

\bibitem[99]{ref99}
F. I. Mautner, The regular representation of a restricted direct product of finite groups (Trans. Amer. Math. Soc., vol.~70, 1951, pp.~531--548).

\bibitem[100]{ref100}
F. I. Mautner, Induced representations (Amer. J. Math., vol.~74, 1952, pp.~737--758).

\bibitem[101]{ref101}
F. I. Mautner, Note on the Fourier inversion formula on groups (Trans. Amer. Math. Soc., vol.~78, 1955, pp.~371--384).

\bibitem[102]{ref102}
D. Milman, On the theory of rings with involution (Dokl. Akad. Nauk S.S.S.R., vol.~76, 1951, pp.~349--352).

% Source PDF physical page 374

\bibitem[103]{ref103}
S. Murakami, Remarks on the structure of maximally almost periodic groups (Osaka Math. J., vol.~2, 1950, pp.~119--129).

\bibitem[104]{ref104}
L. Nachbin, On the finite dimensionality of every irreducible unitary representation of a compact group (Proc. Amer. Math. Soc., vol.~12, 1961, pp.~11--12).

\bibitem[105]{ref105}
M. A. Naimark, Anneaux \`{a} involution (Usp. mat. Nauk, vol.~3, 1948, pp.~52--145) (in Russian).

\bibitem[106]{ref106}
M. A. Naimark, Sur un probl\`{e}me de la th\'{e}orie des anneaux \`{a} involution (Usp. Mat. Nauk, vol.~6, 1951, pp.~160--164) (in Russian).

\bibitem[107]{ref107}
M. A. Naimark, Sur un analogue continu du lemme de Schur (Dokl. Akad. Nauk S.S.S.R., vol.~98, 1954, pp.~185--188) (in Russian).

\bibitem[108]{ref108}
M. A. Naimark, Anneaux norm\'{e}s (Moscou, 1956) (in Russian).

\bibitem[109]{ref109}
M. A. Naimark, Repr\'{e}sentations factorielles d'un groupe localement compact (Dokl. Akad. Nauk S.S.S.R., vol.~134, 1960, pp.~275--277) (in Russian). D\'{e}composition en repr\'{e}sentations factorielles des repr\'{e}sentations unitaires des groupes localement compacts (Sibirsk. Mat. \v{Z}., vol.~2, 1961, pp.~89--99) (in Russian).

\bibitem[110]{ref110}
M. A. Naimark and S. V. Fomin, Sommes directes continues d'espaces hilbertiens et quelques applications (Usp. Mat. Nauk, vol.~10, 1955, pp.~111--142) (in Russian).

\bibitem[111]{ref111}
M. Nakamura, The two-sided representation of an operator algebra (Proc. Jap. Acad., vol.~27, 1951, pp.~172--176).

\bibitem[112]{ref112}
M. Nakamura and T. Turumaru, Simple algebras of completely continuous operators (T\^{o}hoku Math. J., vol.~4, 1952, pp.~303--308).

\bibitem[113]{ref113}
M. Nakamura and T. Turumaru, On extensions of pure states of an abelian operator algebra (T\^{o}hoku Math. J., vol.~6, 1954, pp.~253--257).

\bibitem[114]{ref114}
M. Nakamura, M. Takesaki, and H. Umegaki, A remark on the expectation of operator algebras (Kodai Mat. Sem. Rep., vol.~12, 1960, pp.~82--90).

\bibitem[115]{ref115}
E. Nelson and W. F. Stinespring, Representation of elliptic operators in an enveloping algebra (Amer. J. Math., vol.~81, 1959, pp.~547--560).

\bibitem[116]{ref116}
T. Ogasawara, Finite dimensionality of certain Banach algebras (J. Sc. Hiroshima Univ., vol.~17, 1954, pp.~359--364).

\bibitem[117]{ref117}
T. Ogasawara, A theorem on operator algebras (J. Sc. Hiroshima Univ., vol.~18, 1955, pp.~307--309).

\bibitem[118]{ref118}
T. Ogasawara and K. Yoshinaga, A characterization of dual $B^{*}$-algebras (J. Sc. Hiroshima Univ., vol.~18, 1954, pp.~179--182).

\bibitem[119]{ref119}
T. Ogasawara and K. Yoshinaga, Weakly completely continuous Banach $^{*}$-algebras (J. Sc. Hiroshima Univ., vol.~18, 1954, pp.~15--36).

\bibitem[120]{ref120}
T. Ono, Note on a $B^{*}$-algebra (J. Math. Soc. Japan, vol.~11, 1959, pp.~146--158).

\bibitem[121]{ref121}
D. Ra\"{\i}kov, Sur quelques esp\'{e}ces de convergence des fonctions de type positif (Dokl. Akad. Nauk S.S.S.R., vol.~58, 1947, pp.~1279--1282) (in Russian).

\bibitem[122]{ref122}
C. Rickart, Banach algebras with an adjoint operation (Ann. Math., vol.~47, 1946, pp.~528--550).

\bibitem[123]{ref123}
C. Rickart, The uniqueness of norm problem in Banach algebras (Ann. Math., vol.~51, 1950, pp.~615--628).

\bibitem[124]{ref124}
C. Rickart, Representation of certain Banach algebras on Hilbert space (Duke Math. J., vol.~18, 1951, pp.~27--39).

\bibitem[125]{ref125}
C. Rickart, Spectral permanence for certain Banach algebras (Proc. Amer. Math. Soc., vol.~4, 1953, pp.~191--196).

\bibitem[126]{ref126}
C. Rickart, General theory of Banach algebras (The University series in Higher Mathematics, D. Van Nostrand, New York, 1960).

\bibitem[127]{ref127}
A. Rosenberg, The number of irreducible representations of simple rings with no minimal ideals (Amer. J. Math., vol.~75, 1953, pp.~523--530).

\bibitem[128]{ref128}
S. Sakai, On linear functionals of $W^{*}$-algebras (Proc. Jap. Acad., vol.~34, 1958, pp.~571--574).

% Source PDF physical page 375

\bibitem[129]{ref129}
S. Sakai, On some problems of $C^{*}$-algebras (T\^{o}hoku Math. J., vol.~11, 1959, pp.~453--455).

\bibitem[130]{ref130}
S. Sakai, On a conjecture of Kaplansky (T\^{o}hoku Math. J., vol.~12, 1960, pp.~31--33).

\bibitem[131]{ref131}
S. Sakai, The theory of $W^{*}$-algebras (Notes mim\'{e}ographi\'{e}es, Yale University, 1962).

\bibitem[132]{ref132}
J. A. Schatz, Revue de l'article \cite{ref37} dans Math. Reviews, vol.~14, 1953, pp.~884.

\bibitem[133]{ref133}
I. E. Segal, The group algebra of a locally compact group (Trans. Amer. Math. Soc., vol.~61, 1947, pp.~69--105).

\bibitem[134]{ref134}
I. E. Segal, Irreducible representations of operator algebras (Bull. Amer. Math. Soc., vol.~53, 1947, pp.~73--88).

\bibitem[135]{ref135}
I. E. Segal, Two-sided ideals in operator algebras (Ann. Math., vol.~50, 1949, pp.~856--865).

\bibitem[136]{ref136}
I. E. Segal, The two-sided regular representation of a unimodular locally compact group (Ann. Math., vol.~51, 1950, pp.~293--298).

\bibitem[137]{ref137}
I. E. Segal, An extension of Plancherel's formula to separable unimodular groups (Ann. Math., vol.~52, 1950, pp.~272--292).

\bibitem[138]{ref138}
I. E. Segal, A class of operator algebras which are determined by groups (Duke Math. J., vol.~18, 1951, pp.~221--265).

\bibitem[139]{ref139}
I. E. Segal, Foundations of the theory of dynamical systems of infinitely many degrees of freedom I. (Mat.-fys. Medd. Dan. Vid. Selsk., vol.~31, 1959, pp.~1--38).

\bibitem[140]{ref140}
I. E. Segal and J. von Neumann, A theorem on unitary representations of semi-simple Lie groups (Ann. Math., vol.~52, 1950, pp.~509--517).

\bibitem[141]{ref141}
S. Sherman, The second adjoint of a $C^{*}$-algebra (Proc. Intern. Congr. Math. Cambridge, vol.~1, 1950, pp.~470).

\bibitem[142]{ref142}
S. Sherman, Order in operator algebras (Amer. J. Math., vol.~73, 1951, pp.~227--232).

\bibitem[143]{ref143}
I. M. Singer, Uniformly continuous representations of Lie groups (Ann. Math., vol.~56, 1952, pp.~242--247).

\bibitem[144]{ref144}
W. F. Stinespring, Positive functions on $C^{*}$-algebras (Proc. Amer. Math. Soc., vol.~6, 1955, pp.~211--216).

\bibitem[145]{ref145}
W. F. Stinespring, Integration theorems for gages and duality for unimodular groups (Trans. Amer. Math. Soc., vol.~90, 1959, pp.~15--56).

\bibitem[146]{ref146}
W. F. Stinespring, A semi-simple matrix group is of type I (Proc. Amer. Math. Soc., vol.~9, 1958, pp.~965--967).

\bibitem[147]{ref147}
W. F. Stinespring, Integrability of Fourier transforms for unimodular Lie groups (Duke Math. J., vol.~26, 1959, pp.~123--131).

\bibitem[148]{ref148}
H. Sunouchi, An extension of the Plancherel formula to unimodular groups (T\^{o}hoku Math. J., vol.~4, 1952, pp.~216--230).

\bibitem[149]{ref149}
R. Takahashi, Sur les repr\'{e}sentations unitaires des groupes de Lorentz g\'{e}n\'{e}ralis\'{e}s (Bull. Soc. Math. Fr., vol.~91, 1963, pp.~289--433).

\bibitem[150]{ref150}
Z. Takeda, Conjugate spaces of operator algebras (Proc. Jap. Acad., vol.~30, 1954, pp.~90--95).

\bibitem[151]{ref151}
Z. Takeda, On the representation of operator algebras (Proc. Jap. Acad., vol.~30, 1954, pp.~299--304).

\bibitem[152]{ref152}
Z. Takeda, On the representation of operator algebras II (T\^{o}hoku Math. J., vol.~6, 1954, pp.~212--219).

\bibitem[153]{ref153}
Z. Takeda, Inductive limit and infinite direct product of operator algebras (T\^{o}hoku Math. J., vol.~7, 1955, pp.~67--86).

\bibitem[154]{ref154}
O. Takenouchi, Sur une classe de fonctions continues de type positif sur un groupe localement compact (Math. J. Okayama Univ., vol.~4, 1955, pp.~143--173).

\bibitem[155]{ref155}
O. Takenouchi, Sur la facteur repr\'{e}sentation d'un groupe de Lie r\'{e}soluble de type (E) (Math. J. Okayama Univ., vol.~7, 1957, pp.~151--161).

\bibitem[156]{ref156}
M. Takesaki, On the conjugate space of an operator algebra (T\^{o}hoku Math. J., vol.~10, 1958, pp.~194--203).

% Source PDF physical page 376

\bibitem[157]{ref157}
M. Takesaki, A note on the cross-norm of the direct product of operator algebras (Kodai Math. Sem. Rep., vol.~10, 1958, pp.~137--140).

\bibitem[158]{ref158}
M. Takesaki, A note on the direct product of operator algebras (Kodai Math. Sem. Rep., vol.~11, 1959, pp.~178--181).

\bibitem[159]{ref159}
M. Takesaki, On the non separability of singular representations of operator algebras (Kodai Math. Sem. Reports, vol.~12, 1960, pp.~102--108).

\bibitem[160]{ref160}
M. Takesaki, On some representations of $C^{*}$-algebras (T\^{o}hoku Math. J., vol.~15, 1963, pp.~79--95).

\bibitem[161]{ref161}
M. Tomita, Representations of operator algebras (Math. J. Okayama Univ., vol.~3, 1954, pp.~147--173).

\bibitem[162]{ref162}
M. Tomita, Harmonic analysis on locally compact groups\hfill\break
(Math. J. Okayama Univ., vol.~5, 1956, pp.~133--193).

\bibitem[163]{ref163}
M. Tomita, Spectral theory of operator algebras. I (Math. J. Okayama Univ., vol.~9, 1959, pp.~63--98).

\bibitem[164]{ref164}
M. Tomita, Spectral theory of operator algebras. II (Math. J. Okayama Univ., vol.~10, 1960, pp.~19--60).

\bibitem[165]{ref165}
J. Tomiyama, On the projection of norm one in $W^{*}$-algebras (Proc. Jap. Acad., vol.~33, 1957, pp.~608--612).

\bibitem[166]{ref166}
J. Tomiyama, On the projection of norm one in $W^{*}$-algebras. II (T\^{o}hoku Math. J., vol.~10, 1958, pp.~204--209).

\bibitem[167]{ref167}
J. Tomiyama, Topological representations of $C^{*}$-algebras (T\^{o}hoku Math. J., vol.~14, 1962, pp.~187--204).

\bibitem[168]{ref168}
J. Tomiyama, A characterization of $C^{*}$-algebras whose conjugate spaces are separable (T\^{o}hoku Math. J., vol.~15, 1963, pp.~96--102).

\bibitem[169]{ref169}
J. Tomiyama and M. Takesaki, Applications of fibre bundles to the certain class of $C^{*}$-algebras (T\^{o}hoku Math. J., vol.~13, 1961, pp.~498--523).

\bibitem[170]{ref170}
K. Tsuji, Representation theorems of operator algebras and their applications (Proc. Jap. Acad., vol.~31, 1955, pp.~272--277).

\bibitem[171]{ref171}
K. Tsuji, Harmonic analysis on locally compact groups (Bull. Kyushu Inst. Tech., vol.~2, 1956, pp.~16--32).

\bibitem[172]{ref172}
T. Turumaru, On the direct product of operator algebras. I (T\^{o}hoku Math. J., vol.~4, 1952, pp.~242--251).

\bibitem[173]{ref173}
T. Turumaru, On the direct product of operator algebras. II (T\^{o}hoku Math. J., vol.~5, 1953, pp.~1--7).

\bibitem[174]{ref174}
T. Turumaru, On the direct product of operator algebras. IV (T\^{o}hoku Math. J., vol.~8, 1956, pp.~281--285).

\bibitem[175]{ref175}
T. Turumaru, Crossed product of operator algebra (T\^{o}hoku Math. J., vol.~10, 1958, pp.~355--365).

\bibitem[176]{ref176}
H. Umegaki, On some representation theorems in an operator algebra. I (Proc. Jap. Acad., vol.~27, 1951, pp.~328--333).

\bibitem[177]{ref177}
H. Umegaki, On some representation theorems in an operator algebra. II (Proc. Jap. Acad., vol.~27, 1951, pp.~501--505).

\bibitem[178]{ref178}
H. Umegaki, On some representation theorems in an operator algebra. III (Proc. Jap. Acad., vol.~28, 1952, pp.~29--31).

\bibitem[179]{ref179}
H. Umegaki, Decomposition theorems of operator algebra and their applications (Jap. J. Math., vol.~22, 1953, pp.~27--50).

\bibitem[180]{ref180}
H. Umegaki, Note on irreducible decompositions of a positive linear functional (Kodai Math. Sem. Rep., vol.~6, 1954, pp.~25--32).

\bibitem[181]{ref181}
A. Weil, L'int\'{e}gration dans les groupes topologiques et ses applications, 2nd ed. (Act. Sc. Ind., no.~1145, Hermann, Paris, 1953).

\bibitem[182]{ref182}
A. Wightman and S. Schweber, Configuration space methods (Phys. Rev., vol.~98, 1955, pp.~812--837).

% Source PDF physical page 377

\bibitem[183]{ref183}
K. G. Wolfson, The algebra of bounded operators on Hilbert space (Duke Math. J., vol.~20, 1953, pp.~533--538).

\bibitem[184]{ref184}
A. Wulfsohn, Produit tensoriel de $C^{*}$-alg\`{e}bres (Bull. Sc. Math., vol.~87, 1963, pp.~13--21).

\bibitem[185]{ref185}
H. Yoshizawa, Some remarks on unitary representations of the free group (Osaka Math. J., vol.~3, 1951, pp.~55--63).

\bibitem[186]{ref186}
H. Yoshizawa, A proof of the Plancherel theorem (Proc. Jap. Acad., vol.~30, 1954, pp.~276--281).

\bibitem[187]{ref187}
H. Yoshizawa, On some types of convergence of positive definite functions (Osaka Math. J., vol.~1, 1949, pp.~90--94).

\medskip\noindent\textbf{Second Part.}\par

\medskip\noindent\textbf{1961.}\par

\bibitem[188]{ref188}
I. E. Segal, Foundations of the theory of dynamical systems of infinitely many degrees of freedom. II (Can. J. Math., vol.~13, pp.~1--18).

\bibitem[189]{ref189}
I. Vidav, Sur un syst\`{e}me d'axiomes caract\'{e}risant les alg\`{e}bres $C^{*}$ (Glasnik mat.-fiz. astr. Dru\v{s}tvo, vol.~16, pp.~189--192).

\medskip\noindent\textbf{1962.}\par

\bibitem[190]{ref190}
I. E. Segal, Mathematical characterization of the physical vacuum for a linear Bose-Einstein field (Ill. J. Math., vol.~6, pp.~500--523).

\medskip\noindent\textbf{1963.}\par

\bibitem[191]{ref191}
S. B. Cleveland, Homomorphisms of non commutative $^{*}$-algebras (Pacific J. Math., vol.~13, pp.~1097--1109).

\bibitem[192]{ref192}
E. Hewitt and K. A. Ross, Abstract harmonic analysis I (Springer, Berlin).

\bibitem[193]{ref193}
P. E. Miles, Order isomorphisms of $B^{*}$-algebras (Trans. Amer. Math. Soc., vol.~107, pp.~217--236).

\bibitem[194]{ref194}
M. A. Naimark, La structure des repr\'{e}sentations factorielles d'un groupe localement compact (Dokl. Akad. Nauk S.S.S.R., vol.~148, pp.~775--778, in Russian).

\bibitem[195]{ref195}
R. T. Prosser, On the ideal structure of operator algebras (Mem. Amer. Math. Soc., no.~45).

\bibitem[196]{ref196}
E. St{\o}rmer, Positive linear maps of operator algebras (Acta Math., vol.~110, pp.~233--278).

\bibitem[197]{ref197}
M. Takesaki, On the unitary equivalence among the components of decompositions of representations of involutive Banach algebras and the associated diagonal algebras (T\^{o}hoku Math. J., vol.~15, pp.~365--393).

\bibitem[198]{ref198}
D. Topping, Vector lattices of self-adjoint operators (Bull. Amer. Math. Soc., vol.~69, pp.~251--255).

\medskip\noindent\textbf{1964.}\par

\bibitem[199]{ref199}
J. Dixmier, Utilisation des facteurs hyperfinis dans la th\'{e}orie des $C^{*}$-alg\`{e}bres (C. R. Acad. Sc., vol.~258, pp.~4184--4187).

\bibitem[200]{ref200}
J. Dixmier, Traces sur les $C^{*}$-alg\`{e}bres. II (Bull. Sc. Math., vol.~88, pp.~39--57).

\bibitem[201]{ref201}
L. T. Gardner, A note on isomorphisms of $C^{*}$-algebras (Bull. Amer. Math. Soc., vol.~70, pp.~788--791).

% Source PDF physical page 378

\bibitem[202]{ref202}
A. Guichardet, Sur la d\'{e}composition des repr\'{e}sentations des $C^{*}$-alg\`{e}bres (C. R. Acad. Sc., vol.~258, pp.~768--770).

\bibitem[203]{ref203}
A. Guichardet, Caract\'{e}res et repr\'{e}sentations de produits tensoriels de $C^{*}$-alg\`{e}bres (Ann. scient. \'{E}c. Norm. Sup., vol.~81, pp.~189--206).

\bibitem[204]{ref204}
P. Miles, Derivations on $B^{*}$-algebras (Pacific J. Math., vol.~14, pp.~1359--1366).

\bibitem[205]{ref205}
A. E. Nussbaum, On the reduction of $C^{*}$-algebras (Bull. Amer. Math. Soc., vol.~70, pp.~567--573).

\bibitem[206]{ref206}
D. Shale and W. F. Stinespring, States of the Clifford algebra (Ann. of Math., vol.~80, pp.~365--381).

\bibitem[207]{ref207}
M. Takesaki, On the cross-norm of the direct product of $C^{*}$-algebras (T\^{o}hoku Math. J., vol.~16, pp.~111--122).

\bibitem[208]{ref208}
M. Takesaki, A complement to ``On the unitary equivalence among the components of decompositions of representations of involutive Banach algebras and the associated diagonal algebras'' (T\^{o}hoku Math. J., vol.~16, pp.~226--227).

\bibitem[209]{ref209}
E. Thoma, Ueber unit\"{a}re Darstellungen abz\"{a}hlbarer, diskreter Gruppen (Math. Ann., vol.~153, pp.~111--138).

\bibitem[210]{ref210}
N. T. Varopoulos, Sur les formes positives d'une alg\`{e}bre de Banach (C. R. Acad. Sc., vol.~258, pp.~2465--2467).

\bibitem[211]{ref211}
A. Wulfsohn, Le produit tensoriel de certaines $C^{*}$-alg\`{e}bres (C. R. Acad. Sc., vol.~258, pp.~6052--6054).

\medskip\noindent\textbf{1965.}\par

\bibitem[212]{ref212}
G. R. Allan, A note on $B^{*}$-algebras (Proc. Cambridge Phil. Soc., vol.~61, pp.~29--32).

\bibitem[213]{ref213}
H. J. Borchers, On the vacuum state in quantum field theory. II (Commun. math. Phys., vol.~1, pp.~57--79).

\bibitem[214]{ref214}
H. J. Borchers, Local rings and the connection of spin with statistics (Commun. math. Phys., vol.~1, pp.~281--307).

\bibitem[215]{ref215}
F. Combes, Relations entre formes positives sur une $C^{*}$-alg\`{e}bre (C. R. Acad. Sc., vol.~260, pp.~5435--5438).

\bibitem[216]{ref216}
F. Combes, Repr\'{e}sentations d'une $C^{*}$-alg\`{e}bre et formes lin\'{e}aires positives (C. R. Acad. Sc., vol.~260, pp.~5993--5996).

\bibitem[217]{ref217}
S. Doplicher, An algebraic spectrum condition (Commun. math. Phys., vol.~1, pp.~1--5).

\bibitem[218]{ref218}
E. G. Effros, Transformation groups and $C^{*}$-algebras (Ann. of Math., vol.~81, pp.~38--55).

\bibitem[219]{ref219}
J. Feldman, Borel sets of states and of representations (Michigan Math. J., vol.~12, pp.~363--366).

\bibitem[220]{ref220}
L. T. Gardner, On isomorphisms of $C^{*}$-algebras (Amer. J. Math., vol.~87, pp.~384--396).

\bibitem[221]{ref221}
A. Guichardet, Sur le produit tensoriel des $C^{*}$-alg\`{e}bres (Dokl. Akad. Nauk S.S.S.R., vol.~160, pp.~986--989 et vol.~168, p.~1231; in Russian).

\bibitem[222]{ref222}
A. Guichardet, Utilisation des sous-groupes distingu\'{e}s ouverts dans l'\'{e}tude des repr\'{e}sentations unitaires des groupes localement compacts (Compos. Math., vol.~17, pp.~1--35).

\bibitem[223]{ref223}
R. Haag and J. A. Swieca, When does a quantum field theory describe particles (Commun. math. Phys., vol.~1, pp.~308--320).

\bibitem[224]{ref224}
R. V. Kadison, Transformations of states in operator theory and dynamics (Topology, vol.~3, pp.~177--198).

\bibitem[225]{ref225}
D. Kastler, The $C^{*}$-algebra of a free boson field. I (Commun. math. Phys., vol.~1, pp.~14--48).

\bibitem[226]{ref226}
J. R. Ringrose, On subalgebras of a $C^{*}$-algebra (Pacific J. Math., vol.~15, pp.~1377--1382).

% Source PDF physical page 379

\bibitem[227]{ref227}
S. Sakai, On the central decomposition for positive functionals on $C^{*}$-algebras (Trans. Amer. Math. Soc., vol.~118, pp.~406--419).

\bibitem[228]{ref228}
E. St{\o}rmer, On the Jordan structure of $C^{*}$-algebras (Trans. Amer. Math. Soc., vol.~120, pp.~438--447).

\bibitem[229]{ref229}
D. M. Topping, Vector lattices of self-adjoint operators (Trans. Amer. Math. Soc., vol.~115, pp.~14--30).

\medskip\noindent\textbf{1966.}\par

\bibitem[230]{ref230}
W. B. Arveson, A theorem on the action of abelian unitary groups (Pacific J. Math., vol.~16, pp.~205--212).

\bibitem[231]{ref231}
E. Berkson, Some characterizations of $C^{*}$-algebras (Ill. J. Math., vol.~10, pp.~1--8).

\bibitem[232]{ref232}
F. F. Bonsall, J. Lindenstrauss, and R. R. Phelps, Extreme positive operators on algebras of functions (Math. Scand., vol.~18, pp.~161--182).

\bibitem[233]{ref233}
H. J. Borchers, Energy and momentum as observables in quantum field theory (Commun. math. Phys., vol.~2, pp.~49--54).

\bibitem[234]{ref234}
F. Combes, \'{E}tude des repr\'{e}sentations trac\'{e}es d'une $C^{*}$-alg\`{e}bre (C. R. Acad. Sc., vol.~262, s\'{e}rie A, pp.~114--117).

\bibitem[235]{ref235}
G. F. Dell'Antonio, On some groups of automorphisms of physical observables (Commun. math. Phys., vol.~2, pp.~384--397).

\bibitem[236]{ref236}
G. F. Dell'Antonio, S. Doplicher, and D. Ruelle, A theorem on canonical commutation and anticommutation relations (Commun. math. Phys., vol.~2, pp.~223--230).

\bibitem[237]{ref237}
S. Doplicher, D. Kastler, and D. W. Robinson, Covariance algebras in field theory and statistical mechanics (Commun. math. Phys., vol.~3, pp.~1--28).

\bibitem[238]{ref238}
J. Duncan, The continuity of the involution in Banach $^{*}$-algebras (J. London Math. Soc., vol.~41, pp.~701--706).

\bibitem[239]{ref239}
W. B. Glickfeld, A metric characterization of $C(X)$ and its generalization to $C^{*}$-algebras (Ill. J. Math., vol.~10, pp.~547--556).

\bibitem[240]{ref240}
V. Ja. Golodec, On factor-representations of anti-commutation relations (Dokl. Akad. Nauk S.S.S.R., vol.~167, pp.~19--22; in Russian).

\bibitem[241]{ref241}
A. Guichardet, Produits tensoriels infinis et repr\'{e}sentations des relations d'anticommutation (Ann. scient. \'{E}c. Norm. Sup., vol.~83, pp.~1--52).

\bibitem[242]{ref242}
A. Guichardet, Sur l'homologie et la cohomologie des alg\`{e}bres de Banach (C. R. Acad. Sc., vol.~262, s\'{e}rie A, pp.~38--41).

\bibitem[243]{ref243}
D. Kastler and D. W. Robinson, Invariant states in statistical mechanics (Commun. math. Phys., vol.~3, pp.~151--180).

\bibitem[244]{ref244}
G. Loupias and S. Miracle-Sol\'{e}, $C^{*}$-alg\`{e}bres des syst\`{e}mes canoniques. I (Commun. math. Phys., vol.~2, pp.~31--48).

\bibitem[245]{ref245}
G. Loupias and S. Miracle-Sol\'{e}, $C^{*}$-alg\`{e}bres des syst\`{e}mes canoniques. II (Ann. Inst. H. Poincar\'{e}, vol.~6, pp.~39--58).

\bibitem[246]{ref246}
T. Okayasu, On the tensor product of $C^{*}$-algebras (T\^{o}hoku Math. J., vol.~18, pp.~325--331).

\bibitem[247]{ref247}
T. Okayasu and M. Takesaki, Dual spaces of tensor products of $C^{*}$-algebras (T\^{o}hoku Math. J., vol.~18, pp.~332--337).

\bibitem[248]{ref248}
G. K. Pedersen, Measure theory for $C^{*}$-algebras (Math. Scand., vol.~19, pp.~131--145).

\bibitem[249]{ref249}
D. Rainjonneau, Existence des sommes dans certaines cat\'{e}gories d'al\-g\`{e}bres (C. R. Acad. Sc., vol.~262, s\'{e}rie A, pp.~283--285).

\bibitem[250]{ref250}
D. Ruelle, States of physical systems (Commun. math. Phys., vol.~3, pp.~133--150).

\bibitem[251]{ref251}
B. Russo, Linear mappings of operator algebras (Proc. Amer. Math. Soc., vol.~17, pp.~1019--1022).

\bibitem[252]{ref252}
B. Russo and H. A. Dye, A note on unitary operators in $C^{*}$-algebras (Duke Math. J., vol.~33, pp.~413--416).

% Source PDF physical page 380

\bibitem[253]{ref253}
S. Sakai, On a problem of Calkin (Amer. J. Math., vol.~88, pp.~935--941).

\bibitem[254]{ref254}
S. Sakai, On a characterization of type I $C^{*}$-algebras (Bull. Amer. Math. Soc., vol.~72, pp.~508--512).

\bibitem[255]{ref255}
S. Sakai, On pure states of $C^{*}$-algebras (Proc. Amer. Math. Soc., vol.~17, pp.~86--87).

\bibitem[256]{ref256}
N. B. Vasil'ev, $C^{*}$-algebras with finite dimensional irreducible representations (Usp. mat. Nauk, vol.~21, pp.~135--154; and Russian Math. Surveys, vol.~21, pp.~137--155).

\bibitem[257]{ref257}
G. Zeller-Meier, Produits crois\'{e}s d'une $C^{*}$-alg\`{e}bre par un groupe d'auto\-morphismes (C. R. Acad. Sc., vol.~263, s\'{e}rie A, pp.~20--23).

\medskip\noindent\textbf{1967.}\par

\bibitem[258]{ref258}
W. B. Arveson, Operator algebras and measure preserving automorphisms (Acta Math., vol.~118, pp.~95--109).

\bibitem[259]{ref259}
W. B. Arveson, An algebraic conjugacy invariant for measure-preserving transformations (Bull. Amer. Math. Soc., vol.~73, pp.~121--125).

\bibitem[260]{ref260}
S. K. Berberian and G. H. Orland, On the closure of the numerical range of an operator (Proc. Amer. Math. Soc., vol.~18, pp.~499--503).

\bibitem[261]{ref261}
M. Billik, Idempotent Reynolds operators (J. Math. Anal. Appl., vol.~18, pp.~135--144).

\bibitem[262]{ref262}
R. C. Busby, On structure spaces and extensions of $C^{*}$-algebras (J. functional anal., vol.~1, pp.~370--377).

\bibitem[263]{ref263}
J. M. Chaiken, Finite-particle representations and states of the canonical commutation relations (Ann. of Phys., vol.~42, pp.~23--80).

\bibitem[264]{ref264}
L. A. Coburn, The $C^{*}$-algebra generated by an isometry (Bull. Amer. Math. Soc., vol.~73, pp.~722--726).

\bibitem[265]{ref265}
F. Combes, Sur les \'{e}tats factoriels d'une $C^{*}$-alg\`{e}bre (C. R. Acad. Sc., vol.~265, s\'{e}rie A, pp.~736--739).

\bibitem[266]{ref266}
F. Combes, \'{E}tude des poids d\'{e}finis sur une $C^{*}$-alg\`{e}bre (C. R. Acad. Sc., vol.~265, s\'{e}rie A, pp.~340--343).

\bibitem[267]{ref267}
H. O. Cordes, On a class of $C^{*}$-algebras (Math. Ann., vol.~170, pp.~283--313).

\bibitem[268]{ref268}
C. Delaroche, Sur les centres des $C^{*}$-alg\`{e}bres (C. R. Acad. Sc., vol.~265, s\'{e}rie A, pp.~465--466).

\bibitem[269]{ref269}
C. Delaroche, Sur les centres des $C^{*}$-alg\`{e}bres (Bull. Sc. Math., vol.~91, pp.~105--112).

\bibitem[270]{ref270}
G. F. Dell'Antonio and S. Doplicher, Total number of particles and Fock representation (J. Math. Phys., vol.~8, pp.~663--666).

\bibitem[271]{ref271}
J. Dixmier, On some $C^{*}$-algebras considered by Glimm (J. functional anal., vol.~1, pp.~182--203).

\bibitem[272]{ref272}
J. Dixmier, Sur les automorphismes des alg\`{e}bres de Banach (C. R. Acad. Sc., vol.~264, pp.~729--731).

\bibitem[273]{ref273}
S. Doplicher, R. V. Kadison, D. Kastler, and D. W. Robinson, Asymptotically abelian systems (Commun. math. Phys., vol.~6, pp.~101--120).

\bibitem[274]{ref274}
E. G. Effros and E. St{\o}rmer, Jordan algebras of self-adjoint operators (Trans. Amer. Math. Soc., vol.~127, pp.~313--316).

\bibitem[275]{ref275}
E. G. Effros and F. Hahn, Locally compact transformation groups and $C^{*}$-algebras (Bull. Amer. Math. Soc., vol.~73, pp.~222--226).

\bibitem[276]{ref276}
E. G. Effros and F. Hahn, Locally compact transformation groups and $C^{*}$-algebras (Mem. Amer. Math. Soc., no.~75).

\bibitem[277]{ref277}
V. Ja. Golodec, Factor representations of type $\mathrm{II}_{1}$ for a Clifford algebra (Dokl. Akad. Nauk S.S.S.R., vol.~173, pp.~745--747; in Russian).

\bibitem[278]{ref278}
A. Guichardet, Produits tensoriels continus d'espaces et d'alg\`{e}bres de Banach (Commun. math. Phys., vol.~5, pp.~262--287).

\bibitem[279]{ref279}
R. Haag, N. M. Hugenholtz et M. Winnink, On the equilibrium states in quantum statistical mechanics (Commun. math. Phys., vol.~5, pp.~215--236).

% Source PDF physical page 381

\bibitem[280]{ref280}
H. Halpern, Finite sums of irreducible functional on $C^{*}$-algebras (Proc. Amer. Math. Soc., vol.~18, pp.~352--358).

\bibitem[281]{ref281}
B. E. Johnson, $AW^{*}$-algebras are $QW^{*}$-algebras (Pacific J. Math., vol.~23, pp.~97--99).

\bibitem[282]{ref282}
R. V. Kadison, The energy momentum spectrum of quantum fields (Commun. math. Phys., vol.~4, pp.~258--260).

\bibitem[283]{ref283}
D. A. Kajdan, Sur les relations entre l'espace dual d'un groupe et la structure de ses sous-groupes ferm\'{e}s (Funkc. nyj. Analiz., vol.~1, pp.~71--74; in Russian).

\bibitem[284]{ref284}
E. C. Lance, Automorphisms of postliminal $C^{*}$-algebras (Pacific J. Math., vol.~23, pp.~547--555).

\bibitem[285]{ref285}
O. Lanford and D. Ruelle, Integral representations of invariant states on $B^{*}$-algebras (J. math. Phys., vol.~8, pp.~1460--1463).

\bibitem[286]{ref286}
H. Leptin, Verallgemeinerte $L^{1}$-Algebren und projektive Darstellungen lokal kompakter Gruppen. I and II (Inventiones Math., vol.~3, pp.~257--281 and vol.~4, pp.~68--86).

\bibitem[287]{ref287}
H. Porta and J. T. Schwartz, Representations of the algebra of all operators in Hilbert space, and related analytic function algebras (Comm. pure appl. math., vol.~20, pp.~457--492).

\bibitem[288]{ref288}
D. W. Robinson, Statistical mechanics of quantum spin systems (Commun. math. Phys., vol.~6, pp.~151--160).

\bibitem[289]{ref289}
D. W. Robinson and D. Ruelle, Extremal invariant states (Ann. Inst. H. Poincar\'{e}, vol.~6, pp.~299--310).

\bibitem[290]{ref290}
D. Ruelle, A variational formulation of equilibrium statistical mechanics and the Gibbs phase rule (Commun. math. Phys., vol.~5, pp.~324--329).

\bibitem[291]{ref291}
K. Sait\^{o}, Non-commutative extension of Lusin's theorem (T\^{o}hoku Math. J., vol.~19, pp.~332--340).

\bibitem[292]{ref292}
S. Sakai, On type I $C^{*}$-algebras (Proc. Amer. Math. Soc., vol.~18, pp.~861--863).

\bibitem[293]{ref293}
E. St{\o}rmer, Two-sided ideals in $C^{*}$-algebras (Bull. Amer. Math. Soc., vol.~73, pp.~254--257).

\bibitem[294]{ref294}
E. St{\o}rmer, Large groups of automorphisms of $C^{*}$-algebras (Commun. math. Phys., vol.~5, pp.~1--22).

\bibitem[295]{ref295}
M. Takesaki, A duality in the representation theory of $C^{*}$-algebras (Ann. of Math., vol.~85, pp.~370--382).

\bibitem[296]{ref296}
M. Takesaki, Covariant representations of $C^{*}$-algebras and their locally compact automorphism groups (Acta Math., vol.~119, pp.~273--302).

\bibitem[297]{ref297}
D. Testard, Alg\'{e}bre de covariance (Ann. Inst. H. Poincar\'{e}, vol.~6, pp.~267--297).

\bibitem[298]{ref298}
M. Tomita, The second dual of a $C^{*}$-algebra (Mem. Fac. Sci. Kyushu Univ., vol.~21, pp.~185--193).

\bibitem[299]{ref299}
J. Tomiyama, Applications of Fubini type theorem to the tensor product of $C^{*}$-algebras (T\^{o}hoku Math. J., vol.~19, pp.~213--226).

\bibitem[300]{ref300}
B. J. Vowden, On the Gelfand-Neumark theorem (J. London Math. Soc., vol.~42, pp.~725--731).

\bibitem[301]{ref301}
G. Zeller-Meier, Repr\'{e}sentations fid\'{e}les des produits crois\'{e}s (C. R. Acad. Sc., vol.~264, s\'{e}rie A, pp.~679--682).

\bibitem[302]{ref302}
G. Zeller-Meier, Sur les automorphismes des alg\`{e}bres de Banach (C. R. Acad. Sc., vol.~264, s\'{e}rie A, pp.~1131--1132).

\medskip\noindent\textbf{1968.}\par

\bibitem[303]{ref303}
J. F. Aarnes, On the continuity of automorphic representations of groups (Commun. math. Phys., vol.~7, pp.~332--336).

\bibitem[304]{ref304}
E. Balslev, J. Manuceau, and A. Verbeure, Representations of anticommutation relations and Bogolioubov transformations (Commun. math. Phys., vol.~8, pp.~315--326).

% Source PDF physical page 382

\bibitem[305]{ref305}
E. Balslev and A. Verbeure, States on Clifford algebras (Commun. math. Phys., vol.~7, pp.~55--76).

\bibitem[306]{ref306}
H. Behncke, Structure of certain non normal operators (J. math. mech., vol.~18, pp.~103--107).

\bibitem[307]{ref307}
R. C. Busby, Double centralizers and extensions of $C^{*}$-algebras (Trans. Amer. Math. Soc., vol.~132, pp.~79--99).

\bibitem[308]{ref308}
J. M. Chaiken, Number operators for representations of the canonical commutation relations (Commun. math. Phys., vol.~8, pp.~164--184).

\bibitem[309]{ref309}
F. Combes, Poids sur une $C^{*}$-alg\`{e}bre (J. Math. pures et appl., vol.~47, pp.~57--100).

\bibitem[310]{ref310}
F. Combes, \'{E}l\'{e}ments semi-continus associ\'{e}s \`{a} une $C^{*}$-alg\`{e}bre (C. R. Acad. Sc., vol.~267, s\'{e}rie A, pp.~986--989).

\bibitem[311]{ref311}
E. B. Davies, On the Borel structure of $C^{*}$-algebras (with an appendix by R. V. Kadison) (Commun. math. Phys., vol.~8, pp.~147--163).

\bibitem[312]{ref312}
C. Delaroche, Sur les centres des $C^{*}$-alg\`{e}bres. II (Bull. Sc. Math., vol.~92, pp.~111--128).

\bibitem[313]{ref313}
J. Dixmier, Ideal center of a $C^{*}$-algebra (Duke Math. J., vol.~35, pp.~375--382).

\bibitem[314]{ref314}
J. Dixmier, Localization of a theorem of Glimm (Proc. Amer. Math. Soc., vol.~19, pp.~364--366).

\bibitem[315]{ref315}
S. Doplicher and D. Kastler, Ergodic states in a non commutative ergodic theory (Commun. math. Phys., vol.~7, pp.~1--20).

\bibitem[316]{ref316}
A. Guichardet, Sur un th\'{e}or\`{e}me de Sakai (C. R. Acad. Sc., vol.~266, s\'{e}rie A, pp.~974--975).

\bibitem[317]{ref317}
A. Guichardet, Alg\`{e}bres d'observables associ\'{e}es aux relations de commutation (A. Colin, Paris, Collection Intersciences).

\bibitem[318]{ref318}
A. Guichardet and A. Wulfsohn, Sur les produits tensoriels continus d'espaces hilbertiens (J. functional anal., vol.~2, pp.~371--377).

\bibitem[319]{ref319}
H. Leptin, Darstellungen verallgemeinerter $L^{1}$-Algebren\hfill\break
(Inventiones Math., vol.~5, pp.~192--215).

\bibitem[320]{ref320}
E. Luft, The two-sided closed ideals of the algebra of bounded linear operators of a Hilbert space (Czech Math. J., vol.~18, pp.~595--605).

\bibitem[321]{ref321}
J. Manuceau, \'{E}tude de quelques automorphismes de la $C^{*}$-alg\`{e}bre du champ de bosons libres (Ann. Inst. H. Poincar\'{e}, vol.~8, pp.~117--138).

\bibitem[322]{ref322}
J. Manuceau and A. Verbeure, Quasi-free states of the CCR-algebra and Bogoliubov transformations (Commun. math. Phys., vol.~9, pp.~293--302).

\bibitem[323]{ref323}
S. M. Newberger, R-algebras and locally convex spaces (Math. Ann., vol.~176, pp.~145--156).

\bibitem[324]{ref324}
T. W. Palmer, Characterizations of $C^{*}$-algebras (Bull. Amer. Math. Soc., vol.~74, pp.~538--540).

\bibitem[325]{ref325}
R. J. Plymen, A modification of Piron's axioms (Helv. Phys. Acta, vol.~41, pp.~69--74).

\bibitem[326]{ref326}
R. J. Plymen, $C^{*}$-algebras and Mackey's axioms (Commun. math. Phys., vol.~8, pp.~132--146).

\bibitem[327]{ref327}
G. Rideau, On some representations of the anticommutations relations (Commun. math. Phys., vol.~9, pp.~229--241).

\bibitem[328]{ref328}
J. E. Roberts and G. Roepstorff, Some basic concepts of algebraic quantum theory (Desy, Hamburg).

\bibitem[329]{ref329}
D. W. Robinson, Statistical mechanics of quantum spin systems. II (Commun. math. Phys., vol.~7, pp.~337--348).

\bibitem[330]{ref330}
S. Sakai, Derivations of simple $C^{*}$-algebras (J. functional anal., vol.~2, pp.~202--206).

\bibitem[331]{ref331}
E. St{\o}rmer, On partially ordered vector spaces and their duals, with applications to simplexes and $C^{*}$-algebras (Proc. London Math. Soc., vol.~18, pp.~245--265).

\bibitem[332]{ref332}
E. St{\o}rmer, Irreducible Jordan algebras of self-adjoint operators (Trans. Amer. Math. Soc., vol.~130, pp.~153--166).

% Source PDF physical page 383

\bibitem[333]{ref333}
E. St{\o}rmer, A characterization of pure states of $C^{*}$-algebras (Proc. Amer. Math. Soc., vol.~19, pp.~1100--1102).

\bibitem[334]{ref334}
R. F. Streater, Current commutation relations and continuous tensor products (Nuovo Cimento, vol.~53, pp.~487--495).

\bibitem[335]{ref335}
R. F. Streater, On certain non-relativistic quantized fields (Commun. math. Phys., vol.~7, pp.~93--98).

\bibitem[336]{ref336}
D. M. Topping, Transcendental quasinilpotents in operator algebras (J. functional anal., vol.~2, pp.~342--351).

\bibitem[337]{ref337}
W. Wils, D\'{e}sint\'{e}gration centrale des formes positives sur les $C^{*}$-alg\`{e}bres (C. R. Acad. Sc., vol.~267, s\'{e}rie A, pp.~810--812).

\bibitem[338]{ref338}
M. Winnink, An application of $C^{*}$-algebras to quantum statistical mechanics of systems in equilibrium (Th\'{e}se, Univ. de Groningen).

\bibitem[339]{ref339}
A. Wulfsohn, The primitive spectrum of a tensor product of $C^{*}$-algebras (Proc. Amer. Math. Soc., vol.~19, pp.~1094--1096).

\bibitem[340]{ref340}
K. Ylinen, Compact and finite-dimensional elements of normed algebras (Ann. Acad. Sc. Fennicae, s\'{e}rie A, 428).

\bibitem[341]{ref341}
G. Zeller-Meier, Produits crois\'{e}s d'une $C^{*}$-alg\`{e}bre par un groupe d'auto\-morphismes (J. Math. pures et appl., vol.~47, pp.~101--239).

\medskip\noindent\textbf{Forthcoming.}\par

\bibitem[342]{ref342}
C. A. Akemann, The general Stone-Weierstrass problem.

\bibitem[343]{ref343}
C. A. Akemann, Majorization of operators.

\bibitem[344]{ref344}
W. B. Arveson, A density theorem for operator algebras.

\bibitem[345]{ref345}
F. Combes, Sur les faces d'une $C^{*}$-alg\`{e}bre.

\bibitem[346]{ref346}
J. Dauns and K. H. Hofmann, Representations of rings by continuous sections.

\bibitem[347]{ref347}
E. B. Davies, Decomposition of traces on separable $C^{*}$-algebras.

\bibitem[348]{ref348}
E. B. Davies, The structure of $C^{*}$-algebras.

\bibitem[349]{ref349}
S. Doplicher, D. Kastler, and E. St{\o}rmer, Invariant states and asymptotic abelianness.

\bibitem[350]{ref350}
E. G. Effros, The canonical measures for a separable $C^{*}$-algebra.

\bibitem[351]{ref351}
J. M. G. Fell, An extension of Mackey's method to Banach $^{*}$-algebraic bundles.

\bibitem[352]{ref352}
R. Haag, D. Kastler, and L. Michel, Central decompositions of ergodic states.

\bibitem[353]{ref353}
E. J. Irwin, Concerning time-like asymptotic abelianness.

\bibitem[354]{ref354}
D. Kastler and J. C. T. Pool, Quasi-unitary algebras attached to temperature states in statistical mechanics; a comment on the work of Haag, Hugenholtz and Winnink.

\bibitem[355]{ref355}
E. T. Kehlet, On the monotone sequential closure of a $C^{*}$-algebra.

\bibitem[356]{ref356}
E. C. Lance, Automorphisms of certain operator algebras.

\bibitem[357]{ref357}
E. C. Lance, Inner automorphisms of UHF algebras.

\bibitem[358]{ref358}
O. E. Lanford and D. W. Robinson, Mean entropy of states in quantum statistical mechanics.

\bibitem[359]{ref359}
K. B. Laursen, Tensor products of Banach-$^{*}$-algebras.

\bibitem[360]{ref360}
J. Manuceau, F. Rocca, and D. Testard, On the product form of quasi-free states.

\bibitem[361]{ref361}
J. Manuceau and J. C. Trotin, On lattice spin systems.

\bibitem[362]{ref362}
G. K. Pedersen, Measure theory for $C^{*}$-algebras. II.

\bibitem[363]{ref363}
G. K. Pedersen, Measure theory for $C^{*}$-algebras. III.

\bibitem[364]{ref364}
G. K. Pedersen, On weak and monotone $\sigma$-closure of $C^{*}$-algebras.

\bibitem[365]{ref365}
G. K. Pedersen, A decomposition theorem for $C^{*}$-algebras.

\bibitem[366]{ref366}
R. T. Powers, Representations of the canonical anticommutation relations.

\bibitem[367]{ref367}
F. Rocca, M. Sirugue, and D. Testard, Translation invariant quasi-free states and Bogoliubov transformations.

% Source PDF physical page 384

\bibitem[368]{ref368}
F. Rocca, M. Sirugue, and D. Testard, Quasi-free states as equilibrium states under the Kubo-Martin-Schwinger boundary condition.

\bibitem[369]{ref369}
D. Ruelle, Some remarks on the ground state of infinite systems in statistical mechanics.

\bibitem[370]{ref370}
B. Russo, Isometries of the trace class.

\bibitem[371]{ref371}
S. Sakai, Derivations of uniformly hyperfinite $C^{*}$-algebras.

\bibitem[372]{ref372}
M. S. B. Smith, On the structure of Banach algebras.

\bibitem[373]{ref373}
E. St{\o}rmer, On extremal maps of operator algebras.

\bibitem[374]{ref374}
E. St{\o}rmer, $C^{*}$-homomorphisms of $C^{*}$-algebras.

\bibitem[375]{ref375}
E. St{\o}rmer, $C^{*}$-homomorphisms of $C^{*}$-algebras. II.

\bibitem[376]{ref376}
E. St{\o}rmer, Symmetric states of infinite tensor products of $C^{*}$-algebras.

\bibitem[377]{ref377}
E. St{\o}rmer, Positive linear maps and Jordan homomorphisms of $C^{*}$-algebras.

\bibitem[378]{ref378}
E. St{\o}rmer, States and invariant maps of operator algebras.

\bibitem[379]{ref379}
R. F. Streater, A continuum analogue of the lattice gas.

\bibitem[380]{ref380}
M. Takesaki, A liminal crossed product of a uniformly hyperfinite $C^{*}$-algebra by a compact abelian automorphism group.

\bibitem[381]{ref381}
D. M. Topping, UHF algebras are singly generated.

\bibitem[382]{ref382}
M. Weinless, Existence and uniqueness of the vacuum for linear quantized fields.

\end{DixmierBibliography}
% ===== ASSEMBLED TRANSLATION CHUNK 14: chunk_385_399.tex =====
% Source PDF physical page 385

% Source PDF physical page 386
\BookBackMatter{notation-index}{Index of Notation}

\begingroup
\setlength{\parindent}{0pt}
\setlength{\parskip}{0.08\baselineskip}
\def\DixIndexEntry#1#2{\par\noindent\hangindent=1.5em\hangafter=1 #1: #2.\par}
\def\DixR#1{\XRef{#1}{#1}}
\def\DixA#1{\XRef{A.#1}{A~#1}}
\def\DixB#1{\XRef{B.#1}{B~#1}}

\DixIndexEntry{$\mathbb C$}{\DixR{1.1.1}}
\DixIndexEntry{$\overline z$ ($z$ a complex number)}{\DixR{1.1.2}}
\DixIndexEntry{$\mathcal L(H)$ ($H$ a Hilbert space)}{\DixR{1.1.2}}
\DixIndexEntry{$\operatorname{Sp}_{A}x$, $\operatorname{Sp}x$ ($x$ an element of a unital algebra $A$)}{\DixR{1.1.5}}
\DixIndexEntry{$\operatorname{Sp}'_{A}x$, $\operatorname{Sp}'x$ ($x$ an element of an algebra $A$)}{\DixR{1.1.6}}
\DixIndexEntry{$M'$ ($M$ a subset of an algebra)}{\DixR{1.1.9}}
\DixIndexEntry{$f^{\ast}$ ($f$ a linear functional on an involutive algebra)}{\DixR{1.1.10}}
\DixIndexEntry{$\mathbb R$}{\DixR{1.3.9}}
\DixIndexEntry{$f(x)$ ($x$ a normal element of a $C^{*}$-algebra; $f$ a complex-valued function)}{\DixR{1.5.2}, \DixR{1.5.7}}
\DixIndexEntry{$x^{+}$, $x^{-}$, $\lvert x\rvert$ ($x$ a Hermitian element of a $C^{*}$-algebra)}{\DixR{1.5.7}}
\DixIndexEntry{$A^{+}$ ($A$ a $C^{*}$-algebra)}{\DixR{1.6.5}}
\DixIndexEntry{$x\geq y$, $y\leq x$ ($x,y$ elements of a $C^{*}$-algebra)}{\DixR{1.6.5}}
\DixIndexEntry{$f\geq g$, $g\leq f$ ($f,g$ linear functionals)}{\DixR{2.1.1}}
\DixIndexEntry{$H_{\pi}$ ($\pi$ a representation)}{\DixR{2.2.1}, \DixR{13.1.3}}
\DixIndexEntry{$\pi\simeq\pi'$ ($\pi,\pi'$ representations)}{\DixR{2.2.1}, \DixR{13.1.3}}
\DixIndexEntry{$\bigoplus_i\pi_i$, $\pi_1\oplus\cdots\oplus\pi_n$ ($\pi_i$ representations)}{\DixR{2.2.3}, \DixR{13.1.3}}
\DixIndexEntry{$c\pi$ ($c$ a cardinal; $\pi$ a representation)}{\DixR{2.2.3}, \DixR{13.1.3}}
\DixIndexEntry{$\pi_K$, $\pi_E$ ($\pi$ a representation; $K$ a vector subspace of $H_\pi$ invariant under $\pi$; $E$ the projection onto $K$)}{\DixR{2.2.4}, \DixR{13.1.3}}
\DixIndexEntry{$\pi'\leq\pi$, $\pi\geq\pi'$ ($\pi,\pi'$ representations)}{\DixR{2.2.4}, \DixR{13.1.3}}
\DixIndexEntry{$\bar\pi^{\mathrm o}$ ($\pi$ a representation)}{\DixR{2.2.8}}
\DixIndexEntry{$\widehat A$ ($A$ an involutive algebra)}{\DixR{2.3.2}}
\DixIndexEntry{$\omega_\xi$ ($\xi$ a vector in a Hilbert space)}{\DixR{2.4.2}}
\DixIndexEntry{$\pi_f$, $\xi_f$ ($f$ a positive functional)}{\DixR{2.4.4}}
\DixIndexEntry{$P(A)$ ($A$ an involutive normed algebra)}{\DixR{2.5.2}}
\DixIndexEntry{$\operatorname{Prim}(A)$ ($A$ an algebra)}{\DixR{2.9.7}}
\DixIndexEntry{$\widehat A_I$, $\widehat A^{I}$ ($A$ a $C^{*}$-algebra; $I$ a closed two-sided ideal)}{\DixR{2.11.2}}
\DixIndexEntry{$\operatorname{Prim}_{I}(A)$, $\operatorname{Prim}^{I}(A)$ ($A$ a $C^{*}$-algebra; $I$ a closed two-sided ideal)}{\DixR{2.11.5}}
\DixIndexEntry{$P_I(A)$, $P^{I}(A)$ ($A$ a $C^{*}$-algebra; $I$ a closed two-sided ideal)}{\DixR{2.11.8}}
\DixIndexEntry{$\operatorname{Prim}(\varphi)$, $\widehat\varphi$ ($\varphi$ an automorphism of an algebra)}{\DixR{3.1.9}}
\DixIndexEntry{$\operatorname{Rep}_n(A)$, $\operatorname{Irr}_n(A)$, $\widehat A_n$ ($A$ a $C^{*}$-algebra)}{\DixR{3.5.1}}
\DixIndexEntry{$\operatorname{Rep}(A)$, $\operatorname{Irr}(A)$ ($A$ a separable $C^{*}$-algebra)}{\DixR{3.8.1}}
\DixIndexEntry{$\mathcal{LC}(H)$ ($H$ a Hilbert space)}{\DixR{4.1.1}}
\DixIndexEntry{$\mathfrak m(A)$ ($A$ a $C^{*}$-algebra)}{\DixR{4.5.2}}
\DixIndexEntry{$J(A)$ ($A$ a $C^{*}$-algebra)}{\DixR{4.7.12}}
\DixIndexEntry{$\pi\perp\pi'$ ($\pi,\pi'$ representations)}{\DixR{5.2.2}}
\DixIndexEntry{$\pi\approx\pi'$ ($\pi,\pi'$ representations)}{\DixR{5.3.2}}
\DixIndexEntry{$\mathfrak m_f$ ($f$ a trace)}{\DixR{6.1.2}, \DixR{17.1.3}}
\DixIndexEntry{$\mathfrak n_s$ ($s$ a bitrace)}{\DixR{6.2.1}}
\DixIndexEntry{$N_s$, $A_s$, $H_s$, $J_s$, $\lambda_s$, $\rho_s$ ($s$ a bitrace)}{\DixR{6.2.2}}
\DixIndexEntry{$\mathfrak U_s$, $\mathfrak V_s$ ($s$ a bitrace)}{\DixR{6.2.3}}

% Source PDF physical page 387
\DixIndexEntry{$t_s$ ($s$ a bitrace)}{\DixR{6.2.4}}
\DixIndexEntry{$N_f$, $A_f$, $H_f$, $J_f$, $\lambda_f$, $\rho_f$, $\mathfrak U_f$, $\mathfrak V_f$, $t_f$ ($f$ a trace)}{\DixR{6.4.5}, \DixR{17.1.3}}
\DixIndexEntry{$\mathfrak m_s$ ($s$ a bitrace)}{\DixR{6.4.5}}
\DixIndexEntry{$\operatorname{Fac}_n(A)$, $\operatorname{Fac}(A)$ ($A$ a separable $C^{*}$-algebra)}{\DixR{7.1.3}}
\DixIndexEntry{$\check A$ ($A$ a separable $C^{*}$-algebra)}{\DixR{7.2.1}}
\DixIndexEntry{$\operatorname{FacI}_n(A)$, $\operatorname{FacI}(A)$, $\check A_{\mathrm I}$ ($A$ a separable $C^{*}$-algebra)}{\DixR{7.3.1}}
\DixIndexEntry{$\operatorname{Facf}_n(A)$, $\operatorname{Facf}(A)$, $\check A_f$ ($A$ a separable $C^{*}$-algebra)}{\DixR{7.4.1}}
\DixIndexEntry{$\displaystyle\int^{\oplus}\pi(\zeta)\,d\mu(\zeta)$ [$\zeta\mapsto\pi(\zeta)$ a measurable field of representations]}{\DixR{8.1.3}}
\DixIndexEntry{$\displaystyle\int_{\widehat A}^{\oplus}\zeta\,d\mu(\zeta)$ ($A$ a separable postliminal $C^{*}$-algebra; $\mu$ a positive measure on $\widehat A$)}{\DixR{8.6.3}}
\DixIndexEntry{$\theta_{\xi,\eta}$ ($\xi,\eta$ vectors in a Hilbert space)}{\DixR{10.6.1}}
\DixIndexEntry{$\alpha(H,\xi)$ ($H$ a Hilbert space; $\xi\in H$)}{\DixR{10.6.3}}
\DixIndexEntry{$\beta(A,p)$ ($A$ an elementary $C^{*}$-algebra; $p\in A$)}{\DixR{10.6.5}}
\DixIndexEntry{$\mathcal A(\mathcal H)$ ($\mathcal H$ a continuous field of Hilbert spaces)}{\DixR{10.7.2}}
\DixIndexEntry{$\alpha(\mathcal H,x)$ ($\mathcal H$ a continuous field of Hilbert spaces; $x$ a continuous vector field)}{\DixR{10.7.5}}
\DixIndexEntry{$\beta(\mathcal A,p)$ ($\mathcal A$ a continuous field of elementary $C^{*}$-algebras; $p$ a continuous field of projections)}{\DixR{10.7.5}}
\DixIndexEntry{$\mathcal U$}{\DixR{10.7.10}}
\DixIndexEntry{$\delta(\mathcal A)$ ($\mathcal A$ a continuous field of elementary $C^{*}$-algebras satisfying Fell's condition)}{\DixR{10.7.14}}
\DixIndexEntry{$\delta(A)$ ($A$ a separable continuous-trace $C^{*}$-algebra)}{\DixR{10.9.1}}
\DixIndexEntry{$E(A)$ ($A$ a $C^{*}$-algebra)}{\DixR{11.2.3}}
\DixIndexEntry{$g^{+}$, $g^{-}$ ($g$ a Hermitian linear functional)}{\DixR{12.3.5}}
\DixIndexEntry{$\pi\otimes\pi'$ ($\pi,\pi'$ representations)}{\DixR{13.1.5}}
\DixIndexEntry{$\overline\pi$ ($\pi$ a representation)}{\DixR{13.1.5}}
\DixIndexEntry{$1_H$ ($H$ a Hilbert space)}{\DixR{13.1.7}}
\DixIndexEntry{$M^1(G)$ ($G$ a locally compact group)}{\DixR{13.2.1}}
\DixIndexEntry{$\mathcal K(G)$ ($G$ a locally compact group)}{\DixR{13.2.2}}
\DixIndexEntry{$L^1(G)$ ($G$ a locally compact group)}{\DixR{13.2.2}}
\DixIndexEntry{$\check f$, $\widetilde f$, $f^{\ast}$, $f_s$, ${}_s f$ ($f$ a function on a group)}{\DixR{13.2.3}}
\DixIndexEntry{$\pi_\varphi$, $\xi_\varphi$ ($\varphi$ a continuous function of positive type)}{\DixR{13.4.6}}
\DixIndexEntry{$\mu\gg0$}{\DixR{13.7.1}, \DixR{13.7.6}}
\DixIndexEntry{$C^{*}(G)$ ($G$ a locally compact group)}{\DixR{13.9.1}}
\DixIndexEntry{$PP(G)$ ($G$ a topological group)}{\DixR{16.2.2}}
\DixIndexEntry{$\widehat G$ ($G$ a topological group)}{\DixR{13.1.4}, \DixR{18.1.1}}
\DixIndexEntry{$\operatorname{Rep}_n(G)$ ($G$ a topological group)}{\DixR{18.1.9}}
\DixIndexEntry{$\operatorname{Irr}_n(G)$ ($G$ a topological group)}{\DixR{18.1.10}}
\DixIndexEntry{$\widehat G_r$ ($G$ a locally compact group)}{\DixR{18.3.1}}
\DixIndexEntry{$\operatorname{Rep}(G)$, $\operatorname{Irr}(G)$ ($G$ a separable locally compact group)}{\DixR{18.5.1}}
\DixIndexEntry{$\operatorname{Fac}_n(G)$, $\operatorname{Fac}(G)$, $\operatorname{FacI}_n(G)$, $\operatorname{FacI}(G)$ ($G$ a separable locally compact group)}{\DixR{18.6.1}}
\DixIndexEntry{$\check G$ ($G$ a separable locally compact group)}{\DixR{18.6.2}}
\DixIndexEntry{$\check G_{\mathrm I}$ ($G$ a separable locally compact group)}{\DixR{18.6.3}}
\DixIndexEntry{$\operatorname{Facf}_n(G)$, $\operatorname{Facf}(G)$, $\check G_f$ ($G$ a separable locally compact group)}{\DixR{18.6.4}}
\DixIndexEntry{$\mathfrak m^{1/2}$ ($\mathfrak m$ an ideal)}{\DixA{9}}
\DixIndexEntry{$\mathfrak A_E$ ($\mathfrak A$ a von Neumann algebra; $E$ a projection)}{\DixA{15}}
\DixIndexEntry{$\displaystyle\prod_i\mathfrak A_i$ ($\mathfrak A_i$ von Neumann algebras)}{\DixA{16}}

% Source PDF physical page 388
\DixIndexEntry{$\mathfrak A_1\otimes\mathfrak A_2$ ($\mathfrak A_1,\mathfrak A_2$ von Neumann algebras)}{\DixA{18}}
\DixIndexEntry{$\operatorname{Tr}$}{\DixA{30}}
\DixIndexEntry{$E\sim F$ ($E,F$ projections)}{\DixA{41}}
\DixIndexEntry{$E\prec F$, $E\succ F$ ($E,F$ projections)}{\DixA{41}}
\DixIndexEntry{$\mathfrak U(A)$, $\mathfrak V(A)$ ($A$ a Hilbert algebra)}{\DixA{54}}
\DixIndexEntry{$\displaystyle\int^{\oplus}H(\zeta)\,d\mu(\zeta)$ [$\zeta\mapsto H(\zeta)$ a measurable field of Hilbert spaces]}{\DixA{73}}
\DixIndexEntry{$\displaystyle\int^{\oplus}V(\zeta)\,d\mu(\zeta)$ [$\zeta\mapsto V(\zeta)$ an isomorphism]}{\DixA{74}}
\DixIndexEntry{$\displaystyle\int^{\oplus}T(\zeta)\,d\mu(\zeta)$ [$\zeta\mapsto T(\zeta)$ a measurable field of operators]}{\DixA{77}}
\DixIndexEntry{$\displaystyle\int^{\oplus}\mathfrak A(\zeta)\,d\mu(\zeta)$ [$\zeta\mapsto\mathfrak A(\zeta)$ a measurable field of von Neumann algebras]}{\DixA{87}}
\DixIndexEntry{$\displaystyle\int^{\oplus}\Phi(\zeta)\,d\mu(\zeta)$ [$\zeta\mapsto\Phi(\zeta)$ a measurable field of isomorphisms]}{\DixA{91}}
\DixIndexEntry{$\displaystyle\int^{\oplus}f(\zeta)\,d\mu(\zeta)$ [$\zeta\mapsto f(\zeta)$ a measurable field of traces]}{\DixA{92}}
\DixIndexEntry{$\displaystyle\int^{\oplus}A(\zeta)\,d\mu(\zeta)$ [$\zeta\mapsto A(\zeta)$ a measurable field of Hilbert algebras]}{\DixA{94}}
\DixIndexEntry{$\rho(x)$ ($x$ an element of a Banach algebra)}{\DixB{1}}
\DixIndexEntry{$\mathcal F x$ ($x$ an element of a commutative Banach algebra)}{\DixB{3}}
\DixIndexEntry{$G_\delta$}{\DixB{32}}
\endgroup

% Source PDF physical page 389
\BookBackMatter{terminology-index}{Terminological Index}

\begingroup
\setlength{\parindent}{0pt}
\setlength{\parskip}{0.08\baselineskip}
\def\DixIndexEntry#1#2{\par\noindent\hangindent=1.5em\hangafter=1 #1: #2.\par}
\def\DixR#1{\XRef{#1}{#1}}
\def\DixA#1{\XRef{A.#1}{A~#1}}
\def\DixB#1{\XRef{B.#1}{B~#1}}

\DixIndexEntry{Abelian (projection)}{\DixA{35}}
\DixIndexEntry{Completed (Hilbert algebra)}{\DixA{56}}
\DixIndexEntry{Adjoint}{\DixR{1.1.1}}
\DixIndexEntry{Involutive Banach algebra}{\DixR{1.2.1}}
\DixIndexEntry{Von Neumann algebra}{\DixA{4}}
\DixIndexEntry{Decomposable von Neumann algebra}{\DixA{87}}
\DixIndexEntry{Von Neumann algebra generated by a set}{\DixA{11}}
\DixIndexEntry{Enveloping von Neumann algebra}{\DixR{12.1.4}}
\DixIndexEntry{Induced, reduced von Neumann algebra}{\DixA{15}}
\DixIndexEntry{Hilbert algebra}{\DixA{54}}
\DixIndexEntry{Hilbert algebra of a group}{\XRef{13.10.1}{14.1.1}\translatornote{The original index gives 14.1.1; the link is directed to 13.10.1, the relevant entry in the text.}}
\DixIndexEntry{Direct-sum Hilbert algebra}{\DixA{64}}
\DixIndexEntry{Involutive algebra}{\DixR{1.1.1}}
\DixIndexEntry{Involutive normed algebra}{\DixR{1.2.1}}
\DixIndexEntry{Involutive normed algebra obtained by adjoining an identity element}{\DixR{1.2.3}}
\DixIndexEntry{Algebraically equivalent (representations)}{\DixR{2.9.6}}
\DixIndexEntry{Ampliation}{\DixA{21}}
\DixIndexEntry{Antiliminal ($C^{*}$-algebra, group)}{\DixR{4.3.1}, \DixR{13.9.4}}
\DixIndexEntry{Borel map}{\DixB{19}}
\DixIndexEntry{Canonical map $\widehat A\to\operatorname{Prim}(A)$}{\DixR{2.9.7}}
\DixIndexEntry{Canonical map $P(A)\to\widehat A$}{\DixR{2.5.4}}
\DixIndexEntry{Self-adjoint (set)}{\DixR{1.1.1}}
\DixIndexEntry{Bi-invariant (subspace)}{\DixR{14.3.1}}
\DixIndexEntry{Bitrace}{\DixR{6.2.1}, \DixR{17.1.1}}
\DixIndexEntry{Bitrace associated with a trace}{\DixR{6.4.1}}
\DixIndexEntry{Bitrace of type I, II, \ldots}{\DixR{6.6.4}}
\DixIndexEntry{Maximal bitrace}{\DixR{6.3.3}}
\DixIndexEntry{Bounded relative to a Hilbert algebra (element)}{\DixA{56}}
\DixIndexEntry{$C^{*}$-algebra}{\DixR{1.3.1}}
\DixIndexEntry{Continuous-trace $C^{*}$-algebra}{\DixR{4.5.2}}
\DixIndexEntry{Generalized continuous-trace $C^{*}$-algebra}{\DixR{4.7.12}}
\DixIndexEntry{Conjugate $C^{*}$-algebra of a $C^{*}$-algebra}{\DixR{1.9.1}}
\DixIndexEntry{$C^{*}$-algebra defined by a continuous field of $C^{*}$-algebras}{\DixR{10.4.1}}
\DixIndexEntry{$C^{*}$-algebra defined by a continuous field of Hilbert spaces}{\DixR{10.9.1}}
\DixIndexEntry{Group $C^{*}$-algebra}{\DixR{13.9.1}}
\DixIndexEntry{Enveloping $C^{*}$-algebra}{\XRef{2.7.2}{2.7.7}\translatornote{The original index gives 2.7.7; the link is directed to 2.7.2, where the enveloping $C^{*}$-algebra is defined.}}
\DixIndexEntry{Character defined by a function, by a measure}{\DixR{17.2.3}}
\DixIndexEntry{Character defined by a distribution}{\DixR{17.2.7}}
\DixIndexEntry{Character of a commutative algebra}{\DixB{33}}
\DixIndexEntry{Character of a $C^{*}$-algebra}{\DixR{6.7.1}}
\DixIndexEntry{Character of a traceable factor representation}{\DixR{6.7.4}}
\DixIndexEntry{Character of a compact group}{\DixR{15.3.1}, \DixR{15.3.6}}

% Source PDF physical page 390
\DixIndexEntry{Character of a locally compact group}{\DixR{17.1.1}}
\DixIndexEntry{Normalized character of a compact group}{\DixR{15.3.1}}
\DixIndexEntry{Central (function)}{\DixR{15.3.2}}
\DixIndexEntry{Central (linear functional)}{\DixR{6.8.1}}
\DixIndexEntry{Central relative to a Hilbert algebra (element)}{\DixA{62}}
\DixIndexEntry{Ideal center associated with a representation}{\DixR{5.7.6}}
\DixIndexEntry{Canonical Borel field of Hilbert spaces over $\widehat A$}{\DixR{8.6.1}}
\DixIndexEntry{Borel field of Hilbert spaces}{\DixA{96}}
\DixIndexEntry{Constant field}{\DixR{10.1.4}, \DixA{71}}
\DixIndexEntry{Continuous field of $C^{*}$-algebras}{\DixR{10.3.1}}
\DixIndexEntry{Continuous field of $C^{*}$-algebras defined by a $C^{*}$-algebra}{\DixR{10.5.1}}
\DixIndexEntry{Continuous field of elementary $C^{*}$-algebras}{\DixR{10.3.1}}
\DixIndexEntry{Continuous field of elementary $C^{*}$-algebras associated with a continuous field of Hilbert spaces}{\DixR{10.7.2}}
\DixIndexEntry{Continuous field of Banach spaces}{\DixR{10.1.2}}
\DixIndexEntry{Continuous vector field}{\DixR{10.1.2}, \DixR{10.1.6}}
\DixIndexEntry{Square-integrable vector field}{\DixA{73}}
\DixIndexEntry{Measurable vector field}{\DixA{69}}
\DixIndexEntry{Induced field}{\DixR{10.1.7}}
\DixIndexEntry{Locally trivial field}{\DixR{10.1.8}}
\DixIndexEntry{Measurable field of von Neumann algebras}{\DixA{87}}
\DixIndexEntry{Measurable field of Hilbert algebras}{\DixA{93}}
\DixIndexEntry{Measurable field of representations}{\DixR{8.1.1}, \DixR{18.7.1}}
\DixIndexEntry{Measurable field of Hilbert spaces}{\DixA{69}}
\DixIndexEntry{Measurable field of traces}{\DixA{92}}
\DixIndexEntry{Measurable field of isomorphisms}{\DixA{91}}
\DixIndexEntry{Measurable field of operators}{\DixA{77}}
\DixIndexEntry{Field obtained by gluing}{\DixR{10.1.13}}
\DixIndexEntry{Trivial field}{\DixR{10.1.4}, \DixA{71}}
\DixIndexEntry{Measure class associated with a representation}{\DixR{8.4.3}}
\DixIndexEntry{Class of representations}{\DixR{2.2.1}, \DixR{13.1.3}}
\DixIndexEntry{Coefficient of a representation}{\DixR{13.1.1}}
\DixIndexEntry{Completely positive (map)}{\DixR{2.12.16}}
\DixIndexEntry{Fell's condition}{\DixR{10.5.7}}
\DixIndexEntry{Conjugate (Hilbert space)}{\DixR{2.2.8}}
\DixIndexEntry{Conjugate (representation)}{\DixR{13.1.5}}
\DixIndexEntry{Contained (representation)}{\DixR{2.2.4}, \DixR{13.1.3}}
\DixIndexEntry{Continuous (von Neumann algebra)}{\DixA{38}}
\DixIndexEntry{$C^{*}$-tensor product}{\DixR{2.12.15}}
\DixIndexEntry{$C^{*}$-seminorm}{\DixR{1.9.13}}
\DixIndexEntry{Polar decomposition of a linear functional}{\DixR{12.2.7}}
\DixIndexEntry{Enveloping polar decomposition of a linear functional}{\DixR{12.2.8}}
\DixIndexEntry{Central disintegration of a representation}{\DixR{8.4.3}}
\DixIndexEntry{Dimension of a representation}{\DixR{2.2.1}, \DixR{13.1.3}}
\DixIndexEntry{Dimension of a topological space}{\DixR{10.8.6}}
\DixIndexEntry{Formal dimension of a representation}{\XRef{14.3.4}{14.4.4}\translatornote{The original index gives 14.4.4; the link is directed to Definition 14.3.4.}}
\DixIndexEntry{Disjoint (representations)}{\DixR{5.2.2}, \DixR{13.1.4}}
\DixIndexEntry{Dual of a locally compact group}{\DixR{18.1.1}}
\DixIndexEntry{Dual ($C^{*}$-algebra)}{\DixR{4.7.20}}
\DixIndexEntry{Reduced dual}{\DixR{18.3.1}}
\DixIndexEntry{Elementary ($C^{*}$-algebra)}{\DixR{4.1.1}}
\DixIndexEntry{Trace element relative to a factor}{\DixA{32}}
\DixIndexEntry{Hilbert--Schmidt element in an elementary $C^{*}$-algebra}{\DixR{4.1.9}}

% Source PDF physical page 391
\DixIndexEntry{Hilbert--Schmidt element relative to a factor}{\DixA{32}}
\DixIndexEntry{Equivalent (representations)}{\DixR{2.2.1}, \DixR{13.1.3}}
\DixIndexEntry{Equivalent (projections)}{\DixA{41}}
\DixIndexEntry{Borel space}{\DixB{19}}
\DixIndexEntry{Quotient Borel space}{\DixB{19}}
\DixIndexEntry{Sum Borel space}{\DixB{19}}
\DixIndexEntry{Space of a representation}{\DixR{2.2.1}, \DixR{13.1.3}}
\DixIndexEntry{Model Hilbert space}{\DixR{3.5.1}}
\DixIndexEntry{State}{\DixR{2.1.1}}
\DixIndexEntry{Mutually singular (positive functionals)}{\DixR{12.3.2}}
\DixIndexEntry{Factor}{\DixA{6}}
\DixIndexEntry{Factor (representation)}{\DixR{5.2.6}, \DixR{13.1.4}}
\DixIndexEntry{Weak (topology)}{\DixA{1}}
\DixIndexEntry{Weakly contained}{\DixR{3.4.5}, \DixR{18.1.3}}
\DixIndexEntry{Faithful (trace)}{\DixA{28}}
\DixIndexEntry{Finite (von Neumann algebra)}{\DixA{31}}
\DixIndexEntry{Function of positive type}{\DixR{13.4.1}, \DixR{13.7.6}}
\DixIndexEntry{Function of positive type associated with a set of representations}{\DixR{13.4.6}}
\DixIndexEntry{Function of positive type associated with a representation}{\DixR{13.4.6}}
\DixIndexEntry{Function of positive type defined by a representation and a vector}{\DixR{13.4.6}}
\DixIndexEntry{Pure function of positive type}{\DixR{13.6.1}}
\DixIndexEntry{Almost-periodic function}{\DixR{16.2.2}}
\DixIndexEntry{Positive functional associated with a set of representations}{\DixR{2.4.2}}
\DixIndexEntry{Positive functional associated with a representation}{\DixR{2.4.2}}
\DixIndexEntry{Positive functional defined by a representation and a vector}{\DixR{2.4.2}}
\DixIndexEntry{Vector functional}{\DixR{2.4.2}}
\DixIndexEntry{Strong (topology)}{\DixA{1}}
\DixIndexEntry{Friedrichs extension}{\DixB{23}}
\DixIndexEntry{Compact group associated with a topological group}{\DixR{16.1.2}}
\DixIndexEntry{Unitary group}{\DixR{1.1.6}}
\DixIndexEntry{Hermitian (element)}{\DixR{1.1.3}}
\DixIndexEntry{Hermitian (linear functional)}{\DixR{1.1.10}}
\DixIndexEntry{Homogeneous ($C^{*}$-algebra)}{\DixR{10.9.3}}
\DixIndexEntry{Homogeneous (representation)}{\DixR{5.7.6}}
\DixIndexEntry{Induction of a von Neumann algebra}{\DixA{15}}
\DixIndexEntry{Embeddable in a compact group (group)}{\DixR{16.4.2}}
\DixIndexEntry{Direct integral of representations}{\DixR{8.1.3}, \DixR{18.7.5}}
\DixIndexEntry{Direct integral of Hilbert spaces}{\DixA{73}}
\DixIndexEntry{Involution}{\DixR{1.1.1}}
\DixIndexEntry{Irreducible (representation)}{\DixR{2.3.2}, \DixR{13.1.4}}
\DixIndexEntry{Isomorphism}{\DixR{1.1.7}, \DixR{1.2.4}, \DixR{10.1.3}, \DixA{19}, \DixA{70}}
\DixIndexEntry{Liminal ($C^{*}$-algebra, group)}{\DixR{4.2.1}, \DixR{13.9.4}}
\DixIndexEntry{Linear (topological group)}{\DixR{15.5.4}}
\DixIndexEntry{Plancherel measure}{\DixR{18.8.3}}
\DixIndexEntry{Measure of positive type}{\DixR{13.7.1}}
\DixIndexEntry{Positive measure on a Borel space}{\DixB{30}}
\DixIndexEntry{Standard measure}{\DixB{30}}
\DixIndexEntry{Minimal (projection)}{\DixA{34}}
\DixIndexEntry{Bounded element}{\DixR{13.8.1}}
\DixIndexEntry{Morphism}{\DixR{1.1.7}, \DixR{1.2.4}}
\DixIndexEntry{Mean of an almost-periodic function}{\DixR{16.3.2}}

% Source PDF physical page 392
\DixIndexEntry{Multiple of a representation}{\DixR{2.2.3}, \DixR{13.1.3}}
\DixIndexEntry{Multiplicity-free (representation)}{\DixR{5.4.5}, \DixR{13.1.4}}
\DixIndexEntry{Multiplicity $n$ (representation of)}{\DixR{5.4.8}, \DixR{13.1.4}}
\DixIndexEntry{Intertwining number}{\DixR{2.2.2}, \DixR{13.1.3}}
\DixIndexEntry{Nondegenerate (representation)}{\DixR{2.2.6}}
\DixIndexEntry{Normal (element)}{\DixR{1.1.3}}
\DixIndexEntry{Normal (morphism)}{\DixA{27}}
\DixIndexEntry{Normal (positive functional)}{\DixA{25}}
\DixIndexEntry{Normal (trace)}{\DixA{28}}
\DixIndexEntry{Norm (topology)}{\DixA{1}}
\DixIndexEntry{Trace-class operator}{\DixA{32}}
\DixIndexEntry{Decomposable operator}{\DixA{77}}
\DixIndexEntry{Intertwining operator}{\DixR{2.2.2}, \DixR{13.1.3}}
\DixIndexEntry{Hilbert--Schmidt operator}{\DixA{32}}
\DixIndexEntry{Diagonalizable operator}{\DixA{80}}
\DixIndexEntry{Positive and negative parts of a Hermitian linear functional}{\DixR{12.3.5}}
\DixIndexEntry{Positive (element of a $C^{*}$-algebra)}{\DixR{1.6.5}}
\DixIndexEntry{Positive (linear functional)}{\DixR{2.1.1}}
\DixIndexEntry{Postliminal ($C^{*}$-algebra, group)}{\DixR{4.3.1}, \DixR{13.9.4}}
\DixIndexEntry{Predual of a von Neumann algebra}{\DixA{23}}
\DixIndexEntry{Prime (two-sided ideal)}{\DixR{1.9.13}}
\DixIndexEntry{Primitive (two-sided ideal)}{\DixR{2.9.7}}
\DixIndexEntry{Product of von Neumann algebras}{\DixA{16}}
\DixIndexEntry{Product of $C^{*}$-algebras}{\DixR{1.3.3}}
\DixIndexEntry{Restricted product of $C^{*}$-algebras}{\DixR{1.9.14}}
\DixIndexEntry{Tensor product of von Neumann algebras}{\DixA{18}}
\DixIndexEntry{Tensor product of two representations}{\DixR{13.1.5}}
\DixIndexEntry{Projection}{\DixR{1.1.3}}
\DixIndexEntry{Characteristic projection of a Hilbert algebra}{\DixA{62}}
\DixIndexEntry{Central projection associated with a representation}{\XRef{14.2.4}{14.3.4}\translatornote{The original index gives 14.3.4; the link is directed to 14.2.4, where the term is introduced.}}
\DixIndexEntry{Canonical extension of a positive functional}{\DixR{2.1.5}}
\DixIndexEntry{Canonical extension of a representation}{\DixR{2.2.9}}
\DixIndexEntry{Canonical maximal extension of a bitrace}{\XRef{6.3.5}{6.3.9}\translatornote{The original index gives 6.3.9; the link is directed to Definition 6.3.5.}}
\DixIndexEntry{Properly infinite (von Neumann algebra)}{\DixA{31}}
\DixIndexEntry{Symmetric tensor power of a representation}{\XRef{13.11.9}{13.10.11}\translatornote{The original index gives 13.10.11; the link is directed to 13.11.9, where the term is defined.}}
\DixIndexEntry{Pure (positive functional)}{\DixR{2.5.2}}
\DixIndexEntry{Purely infinite (von Neumann algebra)}{\DixA{31}}
\DixIndexEntry{Quasi-dual}{\DixR{18.6.2}}
\DixIndexEntry{Quasi-equivalent (representations)}{\DixR{5.3.2}, \DixR{13.1.4}}
\DixIndexEntry{Quasi-equivalent (traced representations)}{\DixR{6.6.2}, \DixR{17.1.2}}
\DixIndexEntry{Quasi-spectrum}{\DixR{7.2.2}}
\DixIndexEntry{Rank of a continuous field of elementary $C^{*}$-algebras}{\DixR{10.8.3}}
\DixIndexEntry{Rank of an elementary $C^{*}$-algebra}{\DixR{4.1.9}}
\DixIndexEntry{Spectral radius}{\DixB{1}}
\DixIndexEntry{Real (representation)}{\XRef{13.11.11}{13.10.13}\translatornote{The original index gives 13.10.13; the link is directed to 13.11.11, where real representations are defined.}}
\DixIndexEntry{Regular (representation)}{\DixR{13.1.6}}
\DixIndexEntry{Right regular (representation)}{\DixR{13.1.6}}
\DixIndexEntry{Left regular (representation)}{\DixR{13.1.6}, \DixR{13.9.3}}
\DixIndexEntry{Square-integrable representation}{\XRef{14.1.3}{14.2.3}\translatornote{The original index gives 14.2.3; the link is directed to Definition 14.1.3.}}
\DixIndexEntry{Representation defined by a function of positive type}{\DixR{13.4.6}}
\DixIndexEntry{Representation defined by a positive functional}{\DixR{2.4.4}}
\DixIndexEntry{Representation of an involutive algebra}{\DixR{2.2.1}}

% Source PDF physical page 393
\DixIndexEntry{Integrable representation}{\XRef{14.5.2}{14.6.2}\translatornote{The original index gives 14.6.2; the link is directed to Definition 14.5.2.}}
\DixIndexEntry{Traced representation}{\DixR{6.6.1}, \DixR{17.1.2}}
\DixIndexEntry{Continuous unitary representation}{\DixR{13.1.1}}
\DixIndexEntry{Universal representation}{\DixR{2.7.6}}
\DixIndexEntry{Rich (sub-$C^{*}$-algebra)}{\DixR{11.1.1}}
\DixIndexEntry{Semifinite (von Neumann algebra)}{\DixA{31}}
\DixIndexEntry{Semifinite (trace)}{\DixR{6.1.1}}
\DixIndexEntry{Separable (continuous field)}{\DixR{10.2.1}}
\DixIndexEntry{Separable (topological space)}{\DixB{31}}
\DixIndexEntry{Separation of one set by another}{\DixR{11.1.7}}
\DixIndexEntry{Separating (vector)}{\DixA{14}}
\DixIndexEntry{Separated (point)}{\DixR{3.9.4}}
\DixIndexEntry{Direct sum of representations}{\DixR{2.2.3}, \DixR{13.1.3}}
\DixIndexEntry{Generated involutive subalgebra}{\DixR{1.1.8}}
\DixIndexEntry{Closed generated involutive subalgebra}{\DixR{1.2.5}}
\DixIndexEntry{Bi-invariant subspace associated with a representation}{\XRef{14.2.4}{14.3.4}\translatornote{The original index gives 14.3.4; the link is directed to 14.2.4, where the term is introduced.}}
\DixIndexEntry{Borel subspace}{\DixB{19}}
\DixIndexEntry{Essential subspace}{\DixR{2.2.6}}
\DixIndexEntry{Subrepresentation}{\DixR{2.2.4}}
\DixIndexEntry{Spectrum}{\DixR{3.1.5}}
\DixIndexEntry{Spectrum of a commutative Banach algebra}{\DixB{3}}
\DixIndexEntry{Standard (Borel space)}{\DixB{19}}
\DixIndexEntry{Borel structure defined by functions}{\DixB{19}}
\DixIndexEntry{Mackey Borel structure}{\DixR{3.8.2}, \DixR{7.2.2}, \DixR{18.5.3}, \DixR{18.6.2}}
\DixIndexEntry{Generated Borel structure}{\DixB{19}}
\DixIndexEntry{Finer Borel structure}{\DixB{19}}
\DixIndexEntry{Borel structure underlying a topology}{\DixB{19}}
\DixIndexEntry{Composition series of a $C^{*}$-algebra}{\DixR{4.3.2}}
\DixIndexEntry{Central support}{\DixA{10}}
\DixIndexEntry{Support of a normal positive functional}{\DixA{26}}
\DixIndexEntry{Support of a representation}{\DixR{3.4.6}}
\DixIndexEntry{Enveloping support of a positive functional}{\DixR{12.3.6}}
\DixIndexEntry{Character table}{\DixR{15.4.2}}
\DixIndexEntry{$T_0$-space}{\DixR{3.1.3}}
\DixIndexEntry{Jacobson topology}{\DixR{3.1.1}}
\DixIndexEntry{Topologically irreducible (representation)}{\DixR{2.3.2}, \DixR{13.1.4}}
\DixIndexEntry{Total (subset)}{\DixR{10.2.1}}
\DixIndexEntry{Cyclic for an algebra (vector)}{\DixA{14}}
\DixIndexEntry{Cyclic for a representation (vector)}{\DixR{2.2.5}, \DixR{13.1.3}}
\DixIndexEntry{Traceable (representation)}{\DixR{6.6.7}, \DixR{17.1.4}}
\DixIndexEntry{Trace}{\DixR{6.1.1}, \DixR{17.1.1}}
\DixIndexEntry{Trace associated with a bitrace}{\DixR{6.4.2}}
\DixIndexEntry{Trace of type I, II, \ldots}{\DixR{6.6.4}, \DixR{17.1.3}}
\DixIndexEntry{Traced (representation)}{\DixR{6.6.1}}
\DixIndexEntry{Natural trace}{\DixA{60}}
\DixIndexEntry{Gelfand transform}{\DixB{3}}
\DixIndexEntry{Trivial (representation)}{\DixR{13.1.7}}
\DixIndexEntry{Type I (von Neumann algebra)}{\DixA{35}}
\DixIndexEntry{Type $\mathrm I_n$ (von Neumann algebra)}{\DixA{47}}
\DixIndexEntry{Type I (involutive algebra)}{\DixR{5.5.1}}
\DixIndexEntry{Type I (group)}{\DixR{13.9.4}}
\DixIndexEntry{Type I (representation)}{\DixR{5.4.1}, \DixR{13.1.4}}

% Source PDF physical page 394
\DixIndexEntry{Type II, III (von Neumann algebra)}{\DixA{38}}
\DixIndexEntry{Type II, III (representation)}{\DixR{5.6.1}, \DixR{13.1.4}}
\DixIndexEntry{Type $\mathrm{II}_1$, $\mathrm{II}_\infty$ (von Neumann algebra)}{\DixA{40}}
\DixIndexEntry{Type $\mathrm{II}_1$, $\mathrm{II}_\infty$, III (representation)}{\DixR{5.6.7}, \DixR{13.1.4}}
\DixIndexEntry{Ultraweak (topology)}{\DixA{1}}
\DixIndexEntry{Ultrastrong (topology)}{\DixA{1}}
\DixIndexEntry{Unitary (element)}{\DixR{1.1.6}}
\DixIndexEntry{Approximate identity}{\DixB{29}}
\DixIndexEntry{Increasing directed approximate identity}{\DixR{1.7.1}}
\DixIndexEntry{Absolute value of a linear functional}{\DixR{12.2.7}, \DixR{12.2.8}}
\DixIndexEntry{Vector defined by a function of positive type}{\DixR{13.4.6}}
\DixIndexEntry{Vector defined by a positive functional}{\DixR{2.4.4}}
\endgroup

% Source PDF physical page 395

% Source PDF physical page 396

% Source PDF physical page 397

% Source PDF physical page 398

% Source PDF physical page 399


\end{document}
